\documentclass[oneside,10pt]{amsart}
\PassOptionsToPackage{hyperfootnotes=false}{hyperref}
\ifdefined\pdfvariable\pdfvariable suppressoptionalinfo 767\fi
\usepackage[no-math]{fontspec}
\setmainfont{Noto Serif}
\usepackage{amsmath,amssymb,amsthm}
\usepackage[all]{xy}
\usepackage{xstring,xparse,xcolor}
\definecolor{egalinkred}{rgb}{0.76,0.13,0.28}
\usepackage[
  linktocpage=true,
  colorlinks=true,
  hyperindex,
  hypertexnames=false,
  citecolor=blue,
  linkcolor=egalinkred
]{hyperref}
\usepackage[left=30mm,right=30mm,top=22mm,bottom=25mm]{geometry}
\usepackage{marginnote}
\usepackage[shortlabels]{enumitem}
\usepackage{longtable,multicol,microtype}

\setlength{\emergencystretch}{5em}
\raggedbottom

\hypersetup{
  pdftitle={Cơ sở của hình học đại số, EGA II — ấn bản tiếng Việt},
  pdfauthor={Alexander Grothendieck and Jean Dieudonné; AI translation and edition by OpenAI Codex},
  pdfsubject={EGA II bằng tiếng Việt}
}

% Statements. EGA numbers its statements in the text, so the environments
% are unnumbered and take the number as optional argument. For instance
% \begin{proposition}[2.1.3] starts with the bold heading "Proposition
% (2.1.3). ---" (in the language of the edition) and sets the statement in
% italics; definitions, examples, remarks and notation are set in roman;
% \begin{env}[2.1.4] starts a numbered paragraph with a bold "(2.1.4).".
\newtheoremstyle{egaitem}{}{}{\rmfamily}{}{\bfseries}{.}{ }{\thmnote{(#3)}}
\newtheoremstyle{egastatement}{}{}{\itshape}{}{\bfseries}{. --- }{ }{\thmname{#1}\thmnote{ (#3)}}
\newtheoremstyle{egadescription}{}{}{\rmfamily}{}{\bfseries}{. --- }{ }{\thmname{#1}\thmnote{ (#3)}}
\makeatletter
\newcommand*{\ega@statement}[3]{\theoremstyle{#1}\newtheorem*{#2}{#3}}
\ega@statement{egaitem}{env}{---}
\ega@statement{egastatement}{theorem}{Định lý}
\ega@statement{egastatement}{proposition}{Mệnh đề}
\ega@statement{egastatement}{lemma}{Bổ đề}
\ega@statement{egastatement}{corollary}{Hệ quả}
\ega@statement{egadescription}{definition}{Định nghĩa}
\ega@statement{egadescription}{example}{Ví dụ}
\ega@statement{egadescription}{remark}{Nhận xét}
\ega@statement{egadescription}{remarks}{Các nhận xét}
\ega@statement{egadescription}{notation}{Ký hiệu}
\makeatother
\renewcommand{\proofname}{Chứng minh}

% Notation of the EGA editions: the abbreviations the transcription uses for alphabets, symbols, operators and
% arrows. Written for these editions and dedicated to the public domain under CC0 1.0.
\makeatletter
% a command that may or may not exist already is (re)defined in one step
\newcommand{\ega@define}[2]{\providecommand{#1}{}\renewcommand{#1}{#2}}
\newcommand{\ega@operator}[2]{\providecommand{#1}{}\renewcommand{#1}{\operatorname{#2}}}

% alphabets: sheaves, bold letters, categories
\ega@define{\sh}{\mathcal}
\ega@define{\bb}{\mathbf}
\ega@define{\cat}{\mathtt}
\ega@define{\C}{\cat{C}}
\ega@define{\Set}{\cat{Set}}
\ega@define{\CHom}{\cat{Hom}}
\ega@define{\op}{^{\cat{op}}}
\ega@define{\GL}{\bb{GL}}
\ega@define{\bbGamma}{\boldsymbol{\Gamma}}

% letters and symbols
\renewcommand{\leq}{\leqslant}
\renewcommand{\geq}{\geqslant}
\ega@define{\setmin}{-}
\ega@define{\emp}{\varnothing}
\ega@define{\vphi}{\phi}
\ega@define{\kres}{k}
\ega@define{\rad}{\mathfrak{r}}
\ega@define{\nilrad}{\mathfrak{N}}
\ega@define{\HH}{\mathrm{H}}
\ega@define{\CHH}{\check{\mathrm{H}}}
\ega@define{\RR}{\mathrm{R}}
\ega@define{\LL}{\mathrm{L}}
\ega@define{\red}{\mathrm{red}}
\providecommand{\dual}{}\renewcommand{\dual}[1]{{#1}^{\vee}}
% the tilde of the sheaf associated with a module, written after it: M\supertilde
\ega@define{\supertilde}{{\,\widetilde{\,}\,}}

% operators
\ega@operator{\rg}{rg}
\ega@operator{\gr}{gr}
\ega@operator{\Gr}{Gr}
\ega@operator{\Hom}{Hom}
\ega@operator{\Proj}{Proj}
\ega@operator{\Tor}{Tor}
\ega@operator{\Ext}{Ext}
\ega@operator{\Supp}{Supp}
\ega@operator{\Ker}{Ker}
\ega@operator{\Im}{Im}
\ega@operator{\Coker}{Coker}
\ega@operator{\Spec}{Spec}
\ega@operator{\Spm}{Spm}
\ega@operator{\Spf}{Spf}
\ega@operator{\grad}{grad}
\ega@operator{\dimc}{dimc}
\ega@operator{\codim}{codim}
\ega@operator{\id}{id}
\ega@operator{\Der}{Der}
\ega@operator{\Diff}{Diff}
\ega@operator{\Hyp}{Hyp}
\ega@operator{\Tr}{Tr}

% arrows: long in displays, short in text, scripts and subscripts
\newcommand{\ega@arrow}[2]{\mathchoice{#1}{#2}{#2}{#2}}
\let\ega@shortmapsto\mapsto
\renewcommand{\to}{\ega@arrow{\longrightarrow}{\rightarrow}}
\providecommand{\from}{}\renewcommand{\from}{\ega@arrow{\longleftarrow}{\leftarrow}}
\renewcommand{\mapsto}{\ega@arrow{\longmapsto}{\ega@shortmapsto}}
\ega@define{\isoto}{\simeq}
\ega@define{\simto}{\xrightarrow{\sim}}

% wide accents
\newcommand{\ega@wide}[2]{\renewcommand{#1}[1]{#2{##1}}}
\ega@wide{\hat}{\widehat}
\ega@wide{\tilde}{\widetilde}
\makeatother
\def\Frac{\operatorname{Frac}}
\def\shHom{\mathcal{H}\!\mathit{om}}
\def\shProj{\mathcal{P}\!\mathit{roj}}
\def\shAut{\mathcal{A}\!\mathit{ut}}
% A centred row of asterisks.
\newcommand*{\sectionbreak}{\begin{center}***\end{center}}
\newcommand{\oldpage}[2][]{\marginnote{\normalfont\footnotesize\textbf{#1}~|~#2}\ignorespaces}

\let\EGAOriginalLabel\label
\renewcommand{\label}[1]{\phantomsection\EGAOriginalLabel{#1}}
\makeatletter
\def\EGARefNumberAux#1.#2\@nil{#2}
\newcommand{\EGARefNumber}[1]{\EGARefNumberAux#1\@nil}
\newcommand{\EGAMaybeRefLink}[2]{%
  \ifcsname r@#1\endcsname\hyperref[#1]{#2}\else#2\fi}
\NewDocumentCommand{\sref}{O{} m O{}}{%
  {\normalfont(\IfBlankF{#1}{\textbf{#1},~}%
  \EGAMaybeRefLink{#2}{\EGARefNumber{#2}}%
  \IfBlankF{#3}{,~#3})}}
\makeatother
% \eref{eq:chapter.number} refers to a numbered display, printing "(number)".
\makeatletter
\newcommand*{\ega@labelnumber}[1]{\StrBehind{#1}{.}}
\newcommand{\eref}[1]{\hyperref[#1]{\normalfont{(\ega@labelnumber{#1})}}}
\makeatother


\newcommand{\EGAResetVolumeCounters}{%
  \setcounter{section}{0}\setcounter{subsection}{0}\setcounter{equation}{0}%
  \setcounter{figure}{0}\setcounter{table}{0}\setcounter{footnote}{0}}
\newcommand{\EGAVolumeTitlePage}[3]{%
  \clearpage\EGAResetVolumeCounters\phantomsection\label{section:#1}%
  \pdfbookmark[0]{#2: #3}{volume.#1}\thispagestyle{empty}\vspace*{\fill}%
  \begin{center}
    {\Huge\bfseries #2\par}\vspace{1.5em}{\LARGE #3\par}\vspace{2em}
    {\large Alexander Grothendieck và Jean Dieudonné\par}\vspace{1.3em}
    {\large\bfseries Ấn bản tiếng Việt hiện hành\par}\vspace{1.3em}
  \end{center}
  \begin{minipage}{0.88\textwidth}\small
    Bản dịch, biên tập và dàn trang do OpenAI Codex — GPT-5.6 Sol, mức nỗ lực Ultra thực hiện. Không tuyên bố có biên tập hay thẩm định của con người.\par\medskip
    Nguồn: Alexander Grothendieck và Jean Dieudonné, Publications Mathématiques de l'IHÉS 8 (1961), tr. 5–222; chứng bản số NUMDAM PMIHES\_1961\_\_8\_\_5\_0.\par\medskip
    Ấn bản hiệu chỉnh.\par\medskip
    Ấn bản độc lập này bảo tồn chính xác quan hệ tác giả–nguồn; nó không hàm ý sự xác nhận của các tác giả, người kế quyền hay NUMDAM.
  \end{minipage}
  \vspace*{\fill}\clearpage}

\begin{document}
\EGAVolumeTitlePage{ega2}{EGA II}{Nghiên cứu toàn cục sơ cấp về một số lớp cấu xạ}

\begingroup\small
\noindent\textbf{Nguồn, quyền và tác giả}

Tác phẩm gốc của Alexandre Grothendieck, với sự cộng tác của Jean Dieudonné: EGA II, Publications Mathématiques de l’IHÉS 8 (1961), tr. 5–222, \href{https://numdam.org/item/PMIHES\_1961\_\_8\_\_5\_0/}{NUMDAM}. Bản tham khảo tiếng Anh và cấu trúc TeX được sử dụng có nguồn gốc từ bản dịch do con người thực hiện tại \href{https://github.com/ryankeleti/ega/tree/3d7afe28c99d1887ccdb85acf2ca1afd93fb1971}{ryankeleti/ega}, phiên bản 3d7afe28c99d1887ccdb85acf2ca1afd93fb1971. Các cộng tác viên được ghi tên gồm Mateo Carmona, Tim Holzschuh, Tim Hosgood, Ryan Keleti, Thiago Solovera \& Nery, Wayne Yang. \href{https://github.com/ryankeleti/ega/blob/3d7afe28c99d1887ccdb85acf2ca1afd93fb1971/README.md}{Ghi công theo từng mục} xác định các đóng góp của họ. Bản dịch tiếng Anh được kế thừa vẫn giữ nguồn gốc tác giả đó; bản dịch tiếng Việt và các phần bổ sung do AI thực hiện không được gán cho các cộng tác viên này. Ghi công không đồng nghĩa với việc được cho phép hay được bảo chứng.

Bản dịch tiếng Việt do AI tạo ra. Các hiệu chỉnh, biên tập và dàn trang cũng do AI thực hiện. Người xuất bản không dịch văn bản. Quyền tác giả và nguồn gốc của tác phẩm gốc cùng các đóng góp được kế thừa được bảo tồn, không gán cho AI. Công việc được ghi nhận gồm OpenAI Codex, GPT-5.6 Sol và GPT-6 Astra, mức nỗ lực Ultra.
\par\endgroup\clearpage
\begingroup\small
Quyền đối với các tác phẩm gốc và mọi đóng góp được bảo hộ trong các bản dịch và bản chép lại phái sinh này vẫn thuộc về các chủ thể quyền tương ứng. Dự án không đưa những tác phẩm hay đóng góp đó vào phạm vi công cộng và không cấp giấy phép đối với chúng. Các thông báo quyền của bên thứ ba và các sự cho phép đã có vẫn có hiệu lực.

Các ấn bản được cung cấp miễn phí cho mục đích phi thương mại, thuần túy giáo dục, trên giả định rằng các chủ thể quyền tương ứng sẽ không phản đối việc các bản dịch và bản chép lại được cung cấp miễn phí bên cạnh các bản gốc tiếng Pháp trên \href{https://www.numdam.org/}{NUMDAM}. Đây là một giả định, không phải tuyên bố đã được cho phép; việc truy cập miễn phí bản gốc tự nó không cho phép dịch, chuyển thể hoặc tái xuất bản.

Nếu bạn nắm giữ quyền đối với tài liệu ở đây và phản đối việc cung cấp tài liệu đó, hãy \href{https://github.com/KokunoYumeto/ega-fga-sga-vi/issues}{mở một issue trên GitHub}. Người xuất bản cam kết gỡ bỏ tài liệu bị phản đối theo yêu cầu của chủ thể quyền.

\par\endgroup\clearpage

\section*{Mục lục}

\begin{longtable}{@{}p{0.08\textwidth}p{0.86\textwidth}@{}}
  \hyperref[section:II.1]{\textsection1}. & Các cấu xạ affine.\\
  \hyperref[section:II.2]{\textsection2}. & Các phổ nguyên tố thuần nhất.\\
  \hyperref[section:II.3]{\textsection3}. & Phổ thuần nhất của một bó đại số phân bậc.\\
  \hyperref[section:II.4]{\textsection4}. & Các phân thớ xạ ảnh; các bó ample.\\
  \hyperref[section:II.5]{\textsection5}. & Các cấu xạ quasi-affine; các cấu xạ quasi-projective; các cấu xạ riêng; các cấu xạ xạ ảnh.\\
  \hyperref[section:II.6]{\textsection6}. & Các cấu xạ nguyên và các cấu xạ hữu hạn.\\
  \hyperref[section:II.7]{\textsection7}. & Các tiêu chuẩn định giá.\\
  \hyperref[section:II.8]{\textsection8}. & Các lược đồ thu được bằng phép nổ; các nón chiếu; bao đóng xạ ảnh.
\end{longtable}
\bigskip

\oldpage[II]{5}
Các lớp cấu xạ khác nhau được nghiên cứu trong chương này mà không vận dụng nhiều các phương pháp đối đồng điều; một nghiên cứu sâu hơn sử dụng các phương pháp ấy sẽ được thực hiện ở Chương~III, nơi ta sẽ chủ yếu dùng các mục~\textsection\textsection\hyperref[section:II.2]{2}, \hyperref[section:II.4]{4} và \hyperref[section:II.5]{5} của Chương~II.
Có thể bỏ qua mục~\textsection\hyperref[section:II.8]{8} trong lần đọc đầu tiên: mục này bổ sung một số nội dung cho hình thức luận được phát triển trong các mục~\textsection\textsection\hyperref[section:II.1]{1} đến \hyperref[section:II.3]{3}, vốn quy về những áp dụng đơn giản của hình thức luận ấy, và ta sẽ dùng mục này ít thường xuyên hơn các kết quả khác của chương.
\bigskip

\section{Các cấu xạ affine}
\label{section:II.1}

Phần lớn các kết quả trong mục này là những đối ứng ``toàn cục'' của các kết quả trong Chương~I, \textsection1; vì vậy, về bản chất chúng không mới và chỉ cung cấp một ngôn ngữ thuận tiện cho phần tiếp theo.
\bigskip

\subsection{Các $S$-tiền lược đồ và các $\mathcal{O}_S$-đại số}
\label{subsection:II.1.1}

\begin{env}[1.1.1]
\label{II.1.1.1}
Cho $S$ là một tiền lược đồ, $X$ là một $S$-tiền lược đồ và $f:X\to S$ là cấu xạ cấu trúc của nó.
Theo \sref[0]{0.4.2.4}, ảnh trực tiếp $f_*(\sh{O}_X)$ là một $\sh{O}_S$-đại số, mà ta sẽ
\oldpage[II]{6}
ký hiệu là $\sh{A}(X)$ khi điều đó không gây nhầm lẫn; nếu $U$ là một tập mở của $S$, thì
\[
  \sh{A}(f^{-1}(U))=\sh{A}(X)|U.
\]
Tương tự, với mọi $\sh{O}_X$-môđun $\sh{F}$ (tương ứng, mọi $\sh{O}_X$-đại số $\sh{B}$), ta viết $\sh{A}(\sh{F})$ (tương ứng, $\sh{A}(\sh{B})$) để chỉ ảnh trực tiếp $f_*(\sh{F})$ (tương ứng, $f_*(\sh{B})$); đây là một $\sh{A}(X)$-môđun (tương ứng, một $\sh{A}(X)$-đại số), chứ không chỉ là một $\sh{O}_S$-môđun (tương ứng, một $\sh{O}_S$-đại số).
\end{env}

\begin{env}[1.1.2]
\label{II.1.1.2}
Cho $Y$ là một $S$-tiền lược đồ thứ hai, $g:Y\to S$ là cấu xạ cấu trúc của nó và $h:X\to Y$ là một $S$-cấu xạ; khi đó ta có sơ đồ giao hoán
\[
\label{II.1.1.2.1}
  \xymatrix{
    X\ar[rr]^h\ar[rd]_f & &
    Y\ar[ld]^g\\
    & S.
  }
  \tag{1.1.2.1}
\]

Theo định nghĩa, $h=(\psi,\theta)$, trong đó $\theta:\sh{O}_Y\to h_*(\sh{O}_X)=\psi_*(\sh{O}_X)$ là một đồng cấu các bó vành; từ \sref[0]{0.4.2.2} suy ra một đồng cấu các $\sh{O}_S$-đại số $g_*(\theta):g_*(\sh{O}_Y)\to g_*(h_*(\sh{O}_X))=f_*(\sh{O}_X)$, nói cách khác, một đồng cấu các $\sh{O}_S$-đại số $\sh{A}(Y)\to\sh{A}(X)$, mà ta cũng ký hiệu là $\sh{A}(h)$.
Nếu $h':Y\to Z$ là một $S$-cấu xạ thứ hai, thì hiển nhiên $\sh{A}(h'\circ h)=\sh{A}(h)\circ\sh{A}(h')$.
Như vậy ta đã xác định một \emph{hàm tử phản biến $X\mapsto\sh{A}(X)$} trên phạm trù các $S$-tiền lược đồ, nhận giá trị trong phạm trù các $\sh{O}_S$-đại số.

Bây giờ, cho $\sh{F}$ là một $\sh{O}_X$-môđun, $\sh{G}$ là một $\sh{O}_Y$-môđun và $u:\sh{G}\to\sh{F}$ là một $h$-cấu xạ, nghĩa là, theo \sref[0]{0.4.4.1}, một đồng cấu các $\sh{O}_Y$-môđun $\sh{G}\to h_*(\sh{F})$.
Khi đó $g_*(u):g_*(\sh{G})\to g_*(h_*(\sh{F}))=f_*(\sh{F})$ là một đồng cấu $\sh{A}(\sh{G})\to\sh{A}(\sh{F})$ của các $\sh{O}_S$-môđun, mà ta ký hiệu là $\sh{A}(u)$; hơn nữa, cặp $(\sh{A}(h),\sh{A}(u))$ tạo thành một \emph{di-đồng cấu} từ $\sh{A}(Y)$-môđun $\sh{A}(\sh{G})$ đến $\sh{A}(X)$-môđun $\sh{A}(\sh{F})$.
\end{env}

\begin{env}[1.1.3]
\label{II.1.1.3}
Khi tiền lược đồ $S$ được cố định, ta có thể xem các cặp $(X,\sh{F})$, trong đó $X$ là một $S$-tiền lược đồ và $\sh{F}$ là một $\sh{O}_X$-môđun, như lập thành một \emph{phạm trù}, bằng cách định nghĩa một \emph{cấu xạ} $(X,\sh{F})\to(Y,\sh{G})$ là một cặp $(h,u)$, trong đó $h:X\to Y$ là một $S$-cấu xạ và $u:\sh{G}\to\sh{F}$ là một $h$-cấu xạ.
Khi đó có thể nói rằng $(X,\sh{F})\mapsto(\sh{A}(X),\sh{A}(\sh{F}))$ là một \emph{hàm tử phản biến}, nhận giá trị trong phạm trù mà các vật là những cặp gồm một $\sh{O}_S$-đại số và một môđun trên đại số ấy, còn các cấu xạ là các di-đồng cấu.
\end{env}

\subsection{Các tiền lược đồ affine trên một tiền lược đồ}
\label{subsection:II.1.2}

\begin{definition}[1.2.1]
\label{II.1.2.1}
Cho $X$ là một $S$-tiền lược đồ và $f:X\to S$ là cấu xạ cấu trúc của nó.
Ta nói $X$ là \emph{affine trên $S$} nếu tồn tại một phủ $(S_\alpha)$ của $S$ bởi các tập mở affine sao cho, với mọi $\alpha$, tiền lược đồ do $X$ cảm sinh trên tập mở $f^{-1}(S_\alpha)$ là affine.
\end{definition}

\begin{example}[1.2.2]
\label{II.1.2.2}
Mọi tiền lược đồ con đóng của $S$ đều là một $S$-tiền lược đồ affine trên $S$ (\sref[I]{I.4.2.3} và \sref[I]{I.4.2.4}).
\end{example}

\begin{remark}[1.2.3]
\label{II.1.2.3}
Một tiền lược đồ $X$ affine trên $S$ không nhất thiết là một lược đồ affine, như ví dụ $X=S$ cho thấy \sref{II.1.2.2}.
Mặt khác, nếu một lược đồ affine $X$ là một $S$-tiền lược đồ, thì $X$ không nhất thiết affine trên
\oldpage[II]{7}
$S$ (xem Ví dụ~\sref{II.1.3.3}).
Tuy nhiên, ta nhắc lại rằng, nếu $S$ là một \emph{lược đồ}, thì mọi $S$-tiền lược đồ đồng thời là một lược đồ affine đều affine trên $S$ \sref[I]{I.5.5.10}.
\end{remark}

\begin{proposition}[1.2.4]
\label{II.1.2.4}
Mọi $S$-tiền lược đồ affine trên $S$ đều tách trên $S$ (nói cách khác, đều là một $S$-lược đồ).
\end{proposition}

\begin{proof}
Điều này suy ra ngay từ \sref[I]{I.5.5.5} và \sref[I]{I.5.5.8}.
\end{proof}

\begin{proposition}[1.2.5]
\label{II.1.2.5}
Cho $X$ là một $S$-lược đồ affine trên $S$ và $f:X\to S$ là cấu xạ cấu trúc của nó.
Với mọi tập mở $U\subset S$, $f^{-1}(U)$ affine trên $U$.
\end{proposition}

\begin{proof}
Theo Định nghĩa~\sref{II.1.2.1}, ta có thể quy về trường hợp $S=\Spec(A)$ và $X=\Spec(B)$ đều affine; khi đó $f=({}^a\vphi,\widetilde{\vphi})$, trong đó $\vphi:A\to B$ là một đồng cấu.
Vì các $D(g)$, với $g\in A$, tạo thành một cơ sở của $S$, ta có thể quy về trường hợp $U=D(g)$; nhưng khi đó đã biết $f^{-1}(U)=D(\vphi(g))$ (\textbf{I},~1.2.2.2), từ đó suy ra mệnh đề.
\end{proof}

\begin{proposition}[1.2.6]
\label{II.1.2.6}
Cho $X$ là một $S$-lược đồ affine trên $S$ và $f:X\to S$ là cấu xạ cấu trúc của nó.
Với mọi $\sh{O}_X$-môđun quasi-coherent $\sh{F}$, $f_*(\sh{F})$ là một $\sh{O}_S$-môđun quasi-coherent.
\end{proposition}

\begin{proof}
Có tính đến Mệnh đề~\sref{II.1.2.4}, điều này suy ra từ \sref[I]{I.9.2.2}[a].
\end{proof}

Đặc biệt, $\sh{O}_S$-đại số $\sh{A}(X)=f_*(\sh{O}_X)$ là \emph{quasi-coherent}.

\begin{proposition}[1.2.7]
\label{II.1.2.7}
Cho $X$ là một $S$-lược đồ affine trên $S$.
Với mọi $S$-tiền lược đồ $Y$, ánh xạ $h\mapsto\sh{A}(h)$ từ tập $\Hom_S(Y,X)$ đến tập $\Hom(\sh{A}(X),\sh{A}(Y))$ \sref{II.1.1.2} là song ánh.
\end{proposition}

\begin{proof}
Cho $f:X\to S$ và $g:Y\to S$ là các cấu xạ cấu trúc.
Trước hết, giả sử $S=\Spec(A)$ và $X=\Spec(B)$ đều affine; ta phải chứng minh rằng, với mọi đồng cấu $\omega:f_*(\sh{O}_X)\to g_*(\sh{O}_Y)$ của các $\sh{O}_S$-đại số, tồn tại duy nhất một $S$-cấu xạ $h:Y\to X$ sao cho $\sh{A}(h)=\omega$.
Theo định nghĩa, với mọi tập mở $U\subset S$, $\omega$ xác định một đồng cấu $\omega_U=\Gamma(U,\omega):\Gamma(f^{-1}(U),\sh{O}_X)\to\Gamma(g^{-1}(U),\sh{O}_Y)$ của các $\Gamma(U,\sh{O}_S)$-đại số.
Đặc biệt, với $U=S$, ta thu được một đồng cấu $\vphi:\Gamma(X,\sh{O}_X)\to\Gamma(Y,\sh{O}_Y)$ của các $\Gamma(S,\sh{O}_S)$-đại số, tương ứng với một $S$-cấu xạ $h:Y\to X$ được xác định duy nhất, vì $X$ affine \sref[I]{I.2.2.4}.
Còn phải chứng minh $\sh{A}(h)=\omega$, hay nói cách khác, với mọi tập mở $U$ thuộc một cơ sở của $S$, $\omega_U$ trùng với đồng cấu đại số $\vphi_U$ tương ứng với $S$-cấu xạ $g^{-1}(U)\to f^{-1}(U)$ là hạn chế của $h$.
Ta có thể quy về trường hợp $U=D(\lambda)$, với $\lambda\in A$; khi đó, nếu $f=({}^a\rho,\widetilde{\rho})$, trong đó $\rho:A\to B$ là một đồng cấu vành, thì $f^{-1}(U)=D(\mu)$, với $\mu=\rho(\lambda)$, và $\Gamma(f^{-1}(U),\sh{O}_X)$ là vành phân thức $B_\mu$; sơ đồ
\[
  \xymatrix{
    B\ar[r]^\vphi\ar[d] &
    \Gamma(Y,\sh{O}_Y)\ar[d]\\
    B_\mu\ar[r]_{\vphi_U} &
    \Gamma(g^{-1}(U),\sh{O}_Y)
  }
\]
giao hoán, và sơ đồ tương tự trong đó $\vphi_U$ được thay bởi $\omega_U$ cũng giao hoán; khi đó đẳng thức $\vphi_U=\omega_U$ suy ra từ tính chất phổ dụng của các vành phân thức \sref[0]{0.1.2.4}.

Bây giờ ta chuyển sang trường hợp tổng quát; cho $(S_\alpha)$ là một phủ của $S$ bởi các tập mở affine
\oldpage[II]{8}
sao cho các $f^{-1}(S_\alpha)$ đều affine.
Khi đó mọi đồng cấu $\omega:\sh{A}(X)\to\sh{A}(Y)$ của các $\sh{O}_S$-đại số đều cho, bằng phép hạn chế, một họ đồng cấu
\[
  \omega_\alpha:\sh{A}(f^{-1}(S_\alpha))\to\sh{A}(g^{-1}(S_\alpha))
\]
của các $\sh{O}_{S_\alpha}$-đại số, do đó, theo phần trên, cho một họ các $S_\alpha$-cấu xạ $h_\alpha:g^{-1}(S_\alpha)\to f^{-1}(S_\alpha)$.
Còn phải thấy rằng, với mọi tập mở affine $U$ thuộc một cơ sở của $S_\alpha\cap S_\beta$, các hạn chế của $h_\alpha$ và $h_\beta$ trên $g^{-1}(U)$ trùng nhau; điều này hiển nhiên vì, theo phần trên, cả hai hạn chế đều tương ứng với đồng cấu $\sh{A}(X)|U\to\sh{A}(Y)|U$ là hạn chế của $\omega$.
\end{proof}

\begin{corollary}[1.2.8]
\label{II.1.2.8}
Cho $X$ và $Y$ là hai $S$-lược đồ affine trên $S$.
Một $S$-cấu xạ $h:Y\to X$ là đẳng cấu khi và chỉ khi $\sh{A}(h):\sh{A}(X)\to\sh{A}(Y)$ là đẳng cấu.
\end{corollary}

\begin{proof}
Điều này suy ra ngay từ Mệnh đề~\sref{II.1.2.7} và tính hàm tử của $\sh{A}(X)$.
\end{proof}

\subsection{Tiền lược đồ affine trên $S$ liên kết với một $\mathcal{O}_S$-đại số}
\label{subsection:II.1.3}
\label{II.1.3}

\begin{proposition}[1.3.1]
\label{II.1.3.1}
Cho $S$ là một tiền lược đồ.
Với mọi $\sh{O}_S$-đại số quasi-coherent $\sh{B}$, tồn tại một tiền lược đồ $X$ affine trên $S$, được xác định sai khác một $S$-đẳng cấu duy nhất, sao cho $\sh{A}(X)=\sh{B}$.
\end{proposition}

\begin{proof}
Tính duy nhất suy ra từ Hệ quả~\sref{II.1.2.8}; ta chứng minh sự tồn tại của $X$.
Với mọi tập mở affine $U\subset S$, cho $X_U$ là tiền lược đồ $\Spec(\Gamma(U,\sh{B}))$; vì $\Gamma(U,\sh{B})$ là một $\Gamma(U,\sh{O}_S)$-đại số, $X_U$ là một $S$-tiền lược đồ \sref[I]{I.1.6.1}.
Hơn nữa, vì $\sh{B}$ là quasi-coherent, $\sh{O}_S$-đại số $\sh{A}(X_U)$ được đồng nhất chính tắc với $\sh{B}|U$ (\sref[I]{I.1.3.7}, \sref[I]{I.1.3.13} và \sref[I]{I.1.6.3}).
Cho $V$ là một tập mở affine thứ hai của $S$, và cho $X_{U,V}$ là tiền lược đồ do $X_U$ cảm sinh trên $f_U^{-1}(U\cap V)$, trong đó $f_U$ ký hiệu cấu xạ cấu trúc $X_U\to S$; $X_{U,V}$ và $X_{V,U}$ đều affine trên $U\cap V$ \sref{II.1.2.5}, và theo định nghĩa, $\sh{A}(X_{U,V})$ cùng $\sh{A}(X_{V,U})$ được đồng nhất chính tắc với $\sh{B}|(U\cap V)$.
Do đó, theo \sref{II.1.2.8}, tồn tại một $S$-đẳng cấu chính tắc $\theta_{U,V}:X_{V,U}\to X_{U,V}$; hơn nữa, nếu $W$ là một tập mở affine thứ ba của $S$, và nếu $\theta_{U,V}'$, $\theta_{V,W}'$ và $\theta_{U,W}'$ lần lượt là các hạn chế của $\theta_{U,V}$, $\theta_{V,W}$ và $\theta_{U,W}$ trên các ảnh ngược của $U\cap V\cap W$ trong $X_V$, $X_W$ và $X_W$ qua các cấu xạ cấu trúc, thì $\theta_{U,V}'\circ\theta_{V,W}'=\theta_{U,W}'$.
Vì vậy, tồn tại một tiền lược đồ $X$, một phủ $(T_U)$ của $X$ bởi các tập mở affine, và với mỗi $U$, một đẳng cấu $\vphi_U:X_U\to T_U$, sao cho $\vphi_U$ đưa $f_U^{-1}(U\cap V)$ lên $T_U\cap T_V$ và $\theta_{U,V}=\vphi_U^{-1}\circ\vphi_V$ \sref[I]{I.2.3.1}.
Cấu xạ $g_U=f_U\circ\vphi_U^{-1}$ trang bị cho $T_U$ một cấu trúc $S$-tiền lược đồ, và các cấu xạ $g_U$ cùng $g_V$ trùng nhau trên $T_U\cap T_V$; do đó $X$ là một $S$-tiền lược đồ.
Hơn nữa, theo định nghĩa, hiển nhiên $X$ affine trên $S$ và $\sh{A}(T_U)=\sh{B}|U$; do đó $\sh{A}(X)=\sh{B}$.
\end{proof}

Ta nói $S$-lược đồ $X$ được xác định như trên là \emph{liên kết với $\sh{O}_S$-đại số $\sh{B}$}, hay là \emph{phổ của $\sh{B}$}, và ký hiệu nó là $\Spec(\sh{B})$.

\begin{corollary}[1.3.2]
\label{II.1.3.2}
Cho $X$ là một tiền lược đồ affine trên $S$ và $f:X\to S$ là cấu xạ cấu trúc của nó.
Với mọi tập mở affine $U\subset S$, tiền lược đồ cảm sinh trên $f^{-1}(U)$ là lược đồ affine của vành $\Gamma(U,\sh{A}(X))$.
\end{corollary}

\begin{proof}
\oldpage[II]{9}
Vì theo các Mệnh đề~\sref{II.1.2.6} và \sref{II.1.3.1}, ta có thể giả sử $X$ liên kết với một $\sh{O}_S$-đại số, nên hệ quả suy ra từ phép dựng $X$ được mô tả trong Mệnh đề~\sref{II.1.3.1}.
\end{proof}

\begin{example}[1.3.3]
\label{II.1.3.3}
Cho $S$ là mặt phẳng affine trên một trường $K$, trong đó điểm $0$ được nhân đôi \sref[I]{I.5.5.11}; với các ký hiệu của \sref[I]{I.5.5.11}, $S$ là hợp của hai tập mở affine $Y_1$ và $Y_2$; nếu $f$ là phép nhúng mở $Y_1\to S$, thì $f^{-1}(Y_2)=Y_1\cap Y_2$ không phải là một tập mở affine của $Y_1$ (\emph{loc. cit.}), điều này cho một ví dụ về một lược đồ affine không affine trên $S$.
\end{example}

\begin{corollary}[1.3.4]
\label{II.1.3.4}
Cho $S$ là một lược đồ affine; một $S$-tiền lược đồ $X$ affine trên $S$ khi và chỉ khi $X$ là một lược đồ affine.
\end{corollary}

\begin{corollary}[1.3.5]
\label{II.1.3.5}
Cho $X$ là một tiền lược đồ affine trên một tiền lược đồ $S$, và $Y$ là một $X$-tiền lược đồ.
$Y$ affine trên $S$ khi và chỉ khi $Y$ affine trên $X$.
\end{corollary}

\begin{proof}
Ta quy ngay về trường hợp $S$ là một lược đồ affine, và khi đó $X$ cũng là một lược đồ affine \sref{II.1.3.4}; cả hai điều kiện trong phát biểu đều diễn đạt rằng $Y$ là một lược đồ affine \sref{II.1.3.4}.
\end{proof}

\begin{env}[1.3.6]
\label{II.1.3.6}
Cho $X$ là một tiền lược đồ affine trên $S$.
Việc cho một tiền lược đồ $Y$ affine \emph{trên $X$}, theo Hệ quả~\sref{II.1.3.5}, tương đương với việc cho một tiền lược đồ $Y$ affine \emph{trên $S$} và một $S$-cấu xạ $g:Y\to X$; nói cách khác (Mệnh đề~\sref{II.1.3.1} và \sref{II.1.2.7}), điều đó tương đương với việc cho một $\sh{O}_S$-đại số quasi-coherent $\sh{B}$ và một đồng cấu $\sh{A}(X)\to\sh{B}$ của các $\sh{O}_S$-đại số (có thể xem đồng cấu này là xác định trên $\sh{B}$ một cấu trúc $\sh{A}(X)$-đại số).
Nếu $f:X\to S$ là cấu xạ cấu trúc, thì khi đó $\sh{B}=f_*(g_*(\sh{O}_Y))$.
\end{env}

\begin{corollary}[1.3.7]
\label{II.1.3.7}
Cho $X$ là một tiền lược đồ affine trên $S$; $X$ thuộc kiểu hữu hạn trên $S$ khi và chỉ khi $\sh{O}_S$-đại số quasi-coherent $\sh{A}(X)$ thuộc kiểu hữu hạn \sref[I]{I.9.6.2}.
\end{corollary}

\begin{proof}
Theo định nghĩa \sref[I]{I.9.6.2}, ta có thể quy về trường hợp $S$ affine; khi đó $X$ là một lược đồ affine \sref{II.1.3.4}, và nếu $S=\Spec(A)$ cùng $X=\Spec(B)$, thì $\sh{A}(X)$ là $\sh{O}_S$-đại số $\widetilde{B}$; vì $\Gamma(S,\widetilde{B})=B$, hệ quả suy ra từ \sref[I]{I.9.6.2} và \sref[I]{I.6.3.3}.
\end{proof}

\begin{corollary}[1.3.8]
\label{II.1.3.8}
Cho $X$ là một tiền lược đồ affine trên $S$; $X$ rút gọn khi và chỉ khi $\sh{O}_S$-đại số quasi-coherent $\sh{A}(X)$ rút gọn \sref[0]{0.4.1.4}.
\end{corollary}

\begin{proof}
Vấn đề có tính địa phương trên $S$; theo Hệ quả~\sref{II.1.3.2}, kết quả suy ra từ \sref[I]{I.5.1.1} và \sref[I]{I.5.1.4}.
\end{proof}

\subsection{Các bó quasi-coherent trên một tiền lược đồ affine trên $S$}
\label{subsection:II.1.4}

\begin{proposition}[1.4.1]
\label{II.1.4.1}
Cho $X$ là một tiền lược đồ affine trên $S$, $Y$ là một $S$-tiền lược đồ, và $\sh{F}$ (tương ứng, $\sh{G}$) là một $\sh{O}_X$-môđun (tương ứng, một $\sh{O}_Y$-môđun) quasi-coherent.
Khi đó, ánh xạ $(h,u)\mapsto(\sh{A}(h),\sh{A}(u))$ từ tập các cấu xạ $(Y,\sh{G})\to(X,\sh{F})$ đến tập các di-đồng cấu $(\sh{A}(X),\sh{A}(\sh{F}))\to(\sh{A}(Y),\sh{A}(\sh{G}))$ (\sref{II.1.1.2} và \sref{II.1.1.3}) là song ánh.
\end{proposition}

\begin{proof}
Chứng minh đi đúng theo cùng phương pháp như chứng minh của Mệnh đề~\sref{II.1.2.7}, sử dụng \sref[I]{I.2.2.5} và \sref[I]{I.2.2.4}; các chi tiết được dành cho người đọc.
\end{proof}

\begin{corollary}[1.4.2]
\label{II.1.4.2}
Nếu, ngoài các giả thiết của Mệnh đề~\sref{II.1.4.1}, ta giả sử $Y$ affine trên $S$, thì $(h,u)$ là một đẳng cấu khi và chỉ khi $(\sh{A}(h),\sh{A}(u))$ là một di-đẳng cấu.
\end{corollary}
\oldpage[II]{10}
\begin{proposition}[1.4.3]
\label{II.1.4.3}
Với mỗi cặp $(\sh{B},\sh{M})$ gồm một $\sh{O}_S$-đại số quasi-coherent $\sh{B}$ và một $\sh{B}$-môđun quasi-coherent $\sh{M}$ \emph{(được xét như một $\sh{O}_S$-môđun hoặc như một $\sh{B}$-môđun, hai cách này là như nhau~\sref[I]{I.9.6.1})}, tồn tại một cặp $(X,\sh{F})$ gồm một tiền lược đồ $X$ affine trên $S$ và một $\sh{O}_X$-môđun quasi-coherent $\sh{F}$, sao cho $\sh{A}(X)=\sh{B}$ và $\sh{A}(\sh{F})=\sh{M}$; hơn nữa, cặp này được xác định duy nhất đến một đẳng cấu duy nhất.
\end{proposition}

\begin{proof}
Tính duy nhất suy ra từ \sref{II.1.4.1} và \sref{II.1.4.2}; sự tồn tại được chứng minh như trong \sref{II.1.3.1}, và một lần nữa ta để các chi tiết cho độc giả.
\end{proof}

Ta ký hiệu bởi $\widetilde{\sh{M}}$ $\sh{O}_X$-môđun $\sh{F}$ và nói rằng nó \emph{liên kết} với $\sh{B}$-môđun quasi-coherent $\sh{M}$; với mọi tập mở affine $U\subset S$, $\widetilde{\sh{M}}|p^{-1}(U)$ (trong đó $p$ là cấu xạ cấu trúc $X\to S$) đẳng cấu chính tắc với $(\Gamma(U,\sh{M}))\supertilde$.

\begin{corollary}[1.4.4]
\label{II.1.4.4}
Trong phạm trù các $\sh{B}$-môđun quasi-coherent, $\sh{M}\mapsto\widetilde{\sh{M}}$ là một hàm tử đồng biến, cộng tính và khớp, giao hoán với các giới hạn quy nạp và các tổng trực tiếp.
\end{corollary}

\begin{proof}
Ta quy ngay về trường hợp $S$ affine, và khi đó hệ quả suy ra từ \sref[I]{I.1.3.5}, \sref[I]{I.1.3.9} và \sref[I]{I.1.3.11}.
\end{proof}

\begin{corollary}[1.4.5]
\label{II.1.4.5}
Dưới các giả thiết của \sref{II.1.4.3}, điều kiện cần và đủ để $\widetilde{\sh{M}}$ là một $\sh{O}_X$-môđun hữu hạn sinh là $\sh{M}$ là một $\sh{B}$-môđun hữu hạn sinh.
\end{corollary}

\begin{proof}
Ta quy ngay về trường hợp $S=\Spec(A)$ là một lược đồ affine.
Khi đó $\sh{B}=\widetilde{B}$, trong đó $B$ là một $A$-đại số, và $\sh{M}=\widetilde{M}$, trong đó $M$ là một $B$-môđun \sref[I]{I.1.3.13}; \emph{trên tiền lược đồ $X$}, $\sh{O}_X$ liên kết với vành $B$ còn $\widetilde{\sh{M}}$ liên kết với $B$-môđun $M$; vì vậy, điều kiện cần và đủ để $\widetilde{\sh{M}}$ hữu hạn sinh là $M$ hữu hạn sinh \sref[I]{I.1.3.13}, từ đó suy ra khẳng định.
\end{proof}

\begin{proposition}[1.4.6]
\label{II.1.4.6}
Cho $Y$ là một tiền lược đồ affine trên $S$, còn $X$ và $X'$ là hai tiền lược đồ affine trên $Y$ \emph{(do đó cũng affine trên $S$~\sref{II.1.3.5})}.
Đặt $\sh{B}=\sh{A}(Y)$, $\sh{A}=\sh{A}(X)$ và $\sh{A}'=\sh{A}(X')$.
Khi đó $X\times_Y X'$ affine trên $Y$ \emph{(do đó cũng affine trên $S$)}, và $\sh{A}(X\times_Y X')$ đồng nhất chính tắc với $\sh{A}\otimes_\sh{B}\sh{A}'$.
\end{proposition}

\begin{proof}
Thật vậy, theo \sref[I]{I.9.6.1}, $\sh{A}\otimes_\sh{B}\sh{A}'$ là một $\sh{B}$-đại số quasi-coherent, do đó cũng là một $\sh{O}_S$-đại số quasi-coherent \sref[I]{I.9.6.1}; cho $Z$ là phổ của $\sh{A}\otimes_\sh{B}\sh{A}'$ \sref{II.1.3.1}.
Các $\sh{B}$-đồng cấu chính tắc $\sh{A}\to\sh{A}\otimes_\sh{B}\sh{A}'$ và $\sh{A}'\to\sh{A}\otimes_\sh{B}\sh{A}'$ tương ứng, theo \sref{II.1.2.7}, với các $Y$-cấu xạ $p:Z\to X$ và $p':Z\to X'$.
Để thấy bộ ba $(Z,p,p')$ là một tích $X\times_Y X'$, ta có thể giới hạn ở trường hợp $S$ là một lược đồ affine có vành $C$ \sref[I]{I.3.2.6.4}.
Nhưng khi đó $Y$, $X$ và $X'$ là các lược đồ affine \sref{II.1.3.4}, với các vành $B$, $A$ và $A'$ là các $C$-đại số sao cho $\sh{B}=\widetilde{B}$, $\sh{A}=\widetilde{A}$ và $\sh{A}'=\widetilde{A'}$.
Theo \sref[I]{I.1.3.13}, $\sh{A}\otimes_\sh{B}\sh{A}'$ đồng nhất chính tắc với $\sh{O}_S$-đại số $(A\otimes_B A')\supertilde$; do đó vành $A(Z)$ đồng nhất với $A\otimes_B A'$, còn các cấu xạ $p$, $p'$ tương ứng với các đồng cấu chính tắc $A\to A\otimes_B A'$ và $A'\to A\otimes_B A'$.
Mệnh đề khi đó suy ra từ \sref[I]{I.3.2.2}.
\end{proof}

\begin{corollary}[1.4.7]
\label{II.1.4.7}
Cho $\sh{F}$ (tương ứng, $\sh{F}'$) là một $\sh{O}_X$-môđun quasi-coherent (tương ứng, một $\sh{O}_{X'}$-môđun quasi-coherent); khi đó $\sh{A}(\sh{F}\otimes_Y\sh{F}')$ đồng nhất chính tắc với $\sh{A}(\sh{F})\otimes_{\sh{A}(Y)}\sh{A}(\sh{F}')$.
\end{corollary}

\begin{proof}
Ta biết $\sh{F}\otimes_Y\sh{F}'$ là quasi-coherent trên $X\times_Y X'$ \sref[I]{I.9.1.2}.
Cho $g:Y\to S$, $f:X\to Y$ và $f':X'\to Y$ là các cấu xạ cấu trúc, sao cho cấu xạ cấu trúc
\oldpage[II]{11}
$h:Z\to S$ bằng $g\circ f\circ p$ và cũng bằng $g\circ f'\circ p'$.
Ta định nghĩa một đồng cấu chính tắc
\[
  \sh{A}(\sh{F})\otimes_{\sh{A}(Y)}\sh{A}(\sh{F}')
  \to
  \sh{A}(\sh{F}\otimes_Y\sh{F}')
\]
như sau: với mọi tập mở $U\subset S$, ta có các đồng cấu chính tắc
\[
  \begin{aligned}
  \Gamma(f^{-1}(g^{-1}(U)),\sh{F})
  &\longrightarrow\Gamma(h^{-1}(U),p^*(\sh{F})),
  \\[-2pt]
  \Gamma(f^{\prime-1}(g^{-1}(U)),\sh{F}')
  &\longrightarrow\Gamma(h^{-1}(U),p^{\prime*}(\sh{F}')).
  \end{aligned}
\]
Theo \sref[0]{0.4.4.3}, từ đó ta suy ra một đồng cấu chính tắc
\[
  \begin{aligned}
  &\Gamma(f^{-1}(g^{-1}(U)),\sh{F})
    \otimes_{\Gamma(g^{-1}(U),\sh{O}_Y)}
    \Gamma(f^{\prime-1}(g^{-1}(U)),\sh{F}')\longrightarrow\\
  &\hspace{17em}
    \Gamma(h^{-1}(U),p^*(\sh{F}))
    \otimes_{\Gamma(h^{-1}(U),\sh{O}_Z)}
    \Gamma(h^{-1}(U),p^{\prime*}(\sh{F}')).
  \end{aligned}
\]

Để thấy đây quả là một đẳng cấu của các $\sh{A}(Z)$-môđun, ta có thể giới hạn ở trường hợp $S$ (và do đó $X$, $X'$, $Y$ và $X\times_Y X'$) đều là các lược đồ affine, và, với các ký hiệu của \sref{II.1.4.6}, $\sh{F}=\widetilde{M}$, $\sh{F}'=\widetilde{M'}$, trong đó $M$ (tương ứng, $M'$) là một $A$-môđun (tương ứng, một $A'$-môđun).
Khi đó $\sh{F}\otimes_Y\sh{F}'$ đồng nhất với bó trên $X\times_Y X'$ liên kết với $(A\otimes_B A')$-môđun $M\otimes_B M'$ \sref[I]{I.9.1.3}, và hệ quả suy ra từ sự đồng nhất chính tắc giữa các $\sh{O}_S$-môđun $(M\otimes_B M')\supertilde$ và $\widetilde{M}\otimes_{\widetilde{B}}\widetilde{M'}$ (trong đó $M$, $M'$ và $B$ được xét như các $C$-môđun) (\sref[I]{I.1.3.13} và \sref[I]{I.1.6.3}).
\end{proof}

Đặc biệt, áp dụng \sref{II.1.4.7} cho trường hợp $X=Y$ và $\sh{F}'=\sh{O}_{X'}$, ta thấy $\sh{A}'$-môđun $\sh{A}(f^{\prime*}(\sh{F}))$ đồng nhất với $\sh{A}(\sh{F})\otimes_\sh{B}\sh{A}'$.

\begin{env}[1.4.8]
\label{II.1.4.8}
Đặc biệt, khi $X=X'=Y$ (với $X$ affine trên $S$), ta thấy rằng, nếu $\sh{F}$ và $\sh{G}$ là hai $\sh{O}_X$-môđun quasi-coherent, thì
\[
\label{II.1.4.8.1}
  \sh{A}(\sh{F}\otimes_{\sh{O}_X}\sh{G})
  =\sh{A}(\sh{F})\otimes_{\sh{A}(X)}\sh{A}(\sh{G})
  \tag{1.4.8.1}
\]
sai khác bởi một đẳng cấu chính tắc có tính hàm tử.
Nếu thêm rằng $\sh{F}$ có một trình bày hữu hạn, thì từ \sref[I]{I.1.6.3} và \sref[I]{I.1.3.12} suy ra
\[
\label{II.1.4.8.2}
  \sh{A}(\shHom_X(\sh{F},\sh{G}))
  =\shHom_{\sh{A}(X)}(\sh{A}(\sh{F}),\sh{A}(\sh{G}))
  \tag{1.4.8.2}
\]
sai khác bởi một đẳng cấu chính tắc.
\end{env}

\begin{remark}[1.4.9]
\label{II.1.4.9}
Nếu $X$ và $X'$ là hai tiền lược đồ affine trên $S$, thì tổng rời $X\sqcup X'$ cũng affine trên $S$, vì tổng rời của hai lược đồ affine là một lược đồ affine.
\end{remark}

\begin{proposition}[1.4.10]
\label{II.1.4.10}
Cho $S$ là một tiền lược đồ, $\sh{B}$ là một $\sh{O}_S$-đại số quasi-coherent và $X=\Spec(\sh{B})$.
Với mọi bó iđêan quasi-coherent $\sh{J}$ của $\sh{B}$, $\widetilde{\sh{J}}$ là một bó iđêan quasi-coherent của $\sh{O}_X$, và tiền lược đồ con đóng $Y$ của $X$ được xác định bởi $\widetilde{\sh{J}}$ đẳng cấu chính tắc với $\Spec(\sh{B}/\sh{J})$.
\end{proposition}

\begin{proof}
Từ \sref[I]{I.4.2.3} suy ra ngay rằng $Y$ affine trên $S$; do \sref{II.1.3.1}, ta quy về trường hợp $S$ affine, và khi đó mệnh đề suy ra ngay từ \sref[I]{I.4.1.2}.
\end{proof}

Ta cũng có thể phát biểu kết quả của \sref{II.1.4.10} như sau: nếu $h:\sh{B}\to\sh{B}'$ là một đồng cấu \emph{toàn ánh} của các $\sh{O}_S$-đại số quasi-coherent, thì ${}^a h:\Spec(\sh{B}')\to\Spec(\sh{B})$ là một \emph{phép nhúng đóng}.

\begin{proposition}[1.4.11]
\label{II.1.4.11}
\oldpage[II]{12}
Cho $S$ là một tiền lược đồ, $\sh{B}$ là một $\sh{O}_S$-đại số quasi-coherent và $X=\Spec(\sh{B})$.
Với mọi bó iđêan quasi-coherent $\sh{K}$ của $\sh{O}_S$, ta có (ký hiệu bởi $f$ cấu xạ cấu trúc $X\to S$)
\[
  f^*(\sh{K})\sh{O}_X=(\sh{K}\sh{B})\supertilde
\]
sai khác bởi một đẳng cấu chính tắc.
\end{proposition}

\begin{proof}
Vì vấn đề mang tính địa phương trên $S$, ta quy về trường hợp $S=\Spec(A)$ affine, và trong trường hợp này mệnh đề chính là \sref[I]{I.1.6.9}.
\end{proof}

\subsection{Đổi tiền lược đồ cơ sở}
\label{subsection:II.1.5}

\begin{proposition}[1.5.1]
\label{II.1.5.1}
Cho $X$ là một tiền lược đồ affine trên $S$.
Với mọi phép mở rộng $g:S'\to S$ của tiền lược đồ cơ sở, $X'=X_{(S')}=X\times_S S'$ affine trên $S'$.
\end{proposition}

\begin{proof}
Nếu $f'$ là phép chiếu $X'\to S'$, chỉ cần chứng minh rằng $f^{\prime-1}(U')$ là một tập mở affine với mọi tập mở affine $U'$ của $S'$ sao cho $g(U')$ được chứa trong một tập mở affine $U$ của $S$ \sref{II.1.2.1}; do đó ta có thể giới hạn ở trường hợp $S$ và $S'$ affine, và chỉ cần chứng minh $X'$ khi ấy là một lược đồ affine \sref{II.1.3.4}.
Nhưng khi đó $X$ là một lược đồ affine \sref{II.1.3.4}; nếu $A$, $A'$ và $B$ lần lượt là các vành của $S$, $S'$ và $X$, ta biết $X'$ là lược đồ affine có vành $A'\otimes_A B$ \sref[I]{I.3.2.2}.
\end{proof}

\begin{corollary}[1.5.2]
\label{II.1.5.2}
Dưới các giả thiết của \sref{II.1.5.1}, cho $f:X\to S$ là cấu xạ cấu trúc, còn $f':X'\to S'$ và $g':X'\to X$ là các phép chiếu, sao cho biểu đồ
\[
  \xymatrix{
    X\ar[d] &
    X'\ar[l]_{g'}\ar[d]^{f'}\\
    S &
    S'\ar[l]_g
  }
\]
giao hoán.
Với mọi $\sh{O}_X$-môđun quasi-coherent $\sh{F}$, tồn tại một đẳng cấu chính tắc của các $\sh{O}_{S'}$-môđun
\[
\label{II.1.5.2.1}
  u:g^*(f_*(\sh{F}))\isoto f'_*(g^{\prime*}(\sh{F})).
  \tag{1.5.2.1}
\]
Đặc biệt, tồn tại một đẳng cấu chính tắc từ $\sh{A}(X')$ lên $g^*(\sh{A}(X))$.
\end{corollary}

\begin{proof}
Để định nghĩa $u$, chỉ cần định nghĩa một đồng cấu
\[
  v:f_*(\sh{F})\to
  g_*(f'_*(g^{\prime*}(\sh{F})))
  =f_*(g'_*(g^{\prime*}(\sh{F})))
\]
và lấy $u=v^\sharp$ \sref[0]{0.4.4.3}.
Ta lấy $v=f_*(\rho)$, trong đó $\rho$ là đồng cấu chính tắc $\sh{F}\to g'_*(g^{\prime*}(\sh{F}))$ \sref[0]{0.4.4.3}.
Để chứng minh $u$ là một đẳng cấu, ta có thể giới hạn ở trường hợp $S$ và $S'$, do đó cả $X$ và $X'$, đều affine; với các ký hiệu của \sref{II.1.5.1}, khi đó ta có $\sh{F}=\widetilde{M}$, trong đó $M$ là một $B$-môđun.
Ta thấy ngay $g^*(f_*(\sh{F}))$ và $f'_*(g^{\prime*}(\sh{F}))$ đều bằng $\sh{O}_{S'}$-môđun liên kết với $A'$-môđun $A'\otimes_A M$ (trong đó $M$ được xét như một $A$-môđun), và $u$ là đồng cấu liên kết với ánh xạ đồng nhất (\sref[I]{I.1.6.3}, \sref[I]{I.1.6.5} và \sref[I]{I.1.6.7}).
\end{proof}

\begin{remark}[1.5.3]
\label{II.1.5.3}
Không được cho rằng \sref{II.1.5.2} vẫn đúng khi không còn giả sử $X$ affine trên $S$, ngay cả khi $S'=\Spec(\kres(s))$ ($s\in S$) và $S'\to S$ là cấu xạ chính tắc \sref[I]{I.2.4.5}---trường hợp mà $X'$ chính là \emph{thớ $f^{-1}(s)$} \sref[I]{I.3.6.2}.
Nói cách khác, khi $X$ không affine trên $S$, phép
\oldpage[II]{13}
«ảnh trực tiếp của các bó quasi-coherent» không giao hoán với phép «lấy các thớ».
Tuy nhiên, trong Chương~III \sref[III]{III.4.2.4}, ta sẽ thấy một kết quả theo hướng này, có tính «tiệm cận», đúng cho các bó \emph{coherent} trên $X$ khi $f$ riêng (\hyperref[subsection:II.5.4]{5.4}) và $S$ là Noether.
\end{remark}

\begin{corollary}[1.5.4]
\label{II.1.5.4}
Với mọi tiền lược đồ $X$ affine trên $S$ và mọi $s\in S$, thớ $f^{-1}(s)$ (trong đó $f$ ký hiệu cấu xạ cấu trúc $X\to S$) là một lược đồ affine.
\end{corollary}

\begin{proof}
Chỉ cần áp dụng \sref{II.1.5.1} cho $S'=\Spec(\kres(s))$ và dùng \sref{II.1.3.4}.
\end{proof}

\begin{corollary}[1.5.5]
\label{II.1.5.5}
Cho $X$ là một $S$-tiền lược đồ và $S'$ là một tiền lược đồ affine trên $S$; khi đó $X'=X_{(S')}$ là một tiền lược đồ affine trên $X$.
Hơn nữa, nếu $f:X\to S$ là cấu xạ cấu trúc, thì tồn tại một đẳng cấu chính tắc của các $\sh{O}_X$-đại số
\[
  \sh{A}(X')\isoto f^*(\sh{A}(S')),
\]
và, với mọi $\sh{A}(S')$-môđun quasi-coherent $\sh{M}$, một di-đẳng cấu chính tắc
\[
  f^*(\sh{M})\isoto\sh{A}(f^{\prime*}(\widetilde{\sh{M}})),
\]
trong đó $f'=f_{(S')}$ ký hiệu cấu xạ cấu trúc $X'\to S'$.
\end{corollary}

\begin{proof}
Chỉ cần hoán đổi vai trò của $X$ và $S'$ trong \sref{II.1.5.1} và \sref{II.1.5.2}.
\end{proof}

\begin{env}[1.5.6]
\label{II.1.5.6}
Bây giờ cho $S$ và $S'$ là hai tiền lược đồ, $q:S'\to S$ là một cấu xạ, $\sh{B}$ (tương ứng, $\sh{B}'$) là một $\sh{O}_S$-đại số quasi-coherent (tương ứng, một $\sh{O}_{S'}$-đại số quasi-coherent), và $u:\sh{B}\to\sh{B}'$ là một $q$-cấu xạ (tức là một đồng cấu $\sh{B}\to q_*(\sh{B}')$ của các $\sh{O}_S$-đại số).
Nếu $X=\Spec(\sh{B})$ và $X'=\Spec(\sh{B}')$, ta thu được chính tắc một cấu xạ
\[
  v=\Spec(u):X'\to X
\]
sao cho biểu đồ
\[
\label{II.1.5.6.1}
  \xymatrix{
    X'\ar[r]^-{v}\ar[d] & X\ar[d]\\
    S'\ar[r]_-{q} & S
  }
  \tag{1.5.6.1}
\]
giao hoán (các mũi tên thẳng đứng là các cấu xạ cấu trúc).

Thật vậy, việc cho $u$ tương đương với việc cho một đồng cấu của các $\sh{O}_{S'}$-đại số quasi-coherent $u^\sharp:q^*(\sh{B})\to\sh{B}'$ \sref[0]{0.4.4.3}; do đó đồng cấu này xác định chính tắc một $S'$-cấu xạ
\[
  w:\Spec(\sh{B}')\to\Spec(q^*(\sh{B}))
\]
sao cho $\sh{A}(w)=u^\sharp$ \sref{II.1.2.7}.
Mặt khác, từ \sref{II.1.5.2} suy ra $\Spec(q^*(\sh{B}))$ đồng nhất chính tắc với $X\times_S S'$; cấu xạ $v$ là hợp thành
\[
  X'\xrightarrow{w}X\times_S S'\xrightarrow{p_1}X
\]
của $w$ với phép chiếu thứ nhất, và tính giao hoán của \sref{II.1.5.6.1} suy ra từ các định nghĩa.

Cho $U$ (tương ứng, $U'$) là một tập mở affine của $S$ (tương ứng, của $S'$) sao cho $q(U')\subset U$, còn $A=\Gamma(U,\sh{O}_S)$ và $A'=\Gamma(U',\sh{O}_{S'})$ là các vành tương ứng, $B=\Gamma(U,\sh{B})$ và $B'=\Gamma(U',\sh{B}')$.
Hạn chế của $u$, như một $(q|U')$-cấu xạ $\sh{B}|U\to\sh{B}'|U'$, tương ứng với một di-đồng cấu đại số $B\to B'$; nếu $V$ và $V'$ lần lượt là các ảnh ngược của $U$ và $U'$ trong $X$ và $X'$ qua các cấu xạ cấu trúc, thì cấu xạ $V'\to V$, là hạn chế của $v$, tương ứng theo \sref[I]{I.1.7.3} với di-đồng cấu trên.
\end{env}

\begin{env}[1.5.7]
\label{II.1.5.7}
Dưới cùng các giả thiết như trong \sref{II.1.5.6}, cho $\sh{M}$ là một $\sh{B}$-môđun quasi-coherent; khi đó tồn tại một đẳng cấu chính tắc của các $\sh{O}_{X'}$-môđun
\[
\label{II.1.5.7.1}
  v^*(\widetilde{\sh{M}})
  \isoto
  (q^*(\sh{M})\otimes_{q^*(\sh{B})}\sh{B}')\supertilde.
  \tag{1.5.7.1}
\]
\oldpage[II]{14}
Thật vậy, đẳng cấu chính tắc \sref{II.1.5.2.1} cho một đẳng cấu chính tắc từ $p_1^*(\widetilde{\sh{M}})$ lên bó trên $\Spec(q^*(\sh{B}))$ liên kết với $q^*(\sh{B})$-môđun $q^*(\sh{M})$, rồi chỉ cần áp dụng \sref{II.1.4.7}.
\end{env}

\subsection{Các cấu xạ affine}
\label{subsection:II.1.6}

\begin{env}[1.6.1]
\label{II.1.6.1}
Ta nói một cấu xạ $f:X\to Y$ của các tiền lược đồ là \emph{affine} nếu nó xác định $X$ như một tiền lược đồ affine trên $Y$.
Trong ngôn ngữ này, các tính chất của những tiền lược đồ affine trên một tiền lược đồ khác được phát biểu như sau:
\end{env}

\begin{proposition}[1.6.2]
\label{II.1.6.2}
\medskip\noindent
\begin{enumerate}
  \item[{\rm(i)}] Một phép nhúng đóng là affine.
  \item[{\rm(ii)}] Hợp thành của hai cấu xạ affine là affine.
  \item[{\rm(iii)}] Nếu $f:X\to Y$ là một $S$-cấu xạ affine, thì $f_{(S')}:X_{(S')}\to Y_{(S')}$ là affine với mọi phép đổi cơ sở $S'\to S$.
  \item[{\rm(iv)}] Nếu $f:X\to Y$ và $f':X'\to Y'$ là hai $S$-cấu xạ affine, thì
    \[
      f\times_S f':X\times_S X'\to Y\times_S Y'
    \]
    là affine.
  \item[{\rm(v)}] Nếu $f:X\to Y$ và $g:Y\to Z$ là hai cấu xạ sao cho $g\circ f$ affine và $g$ tách, thì $f$ affine.
  \item[{\rm(vi)}] Nếu $f$ affine, thì $f_\red$ cũng vậy.
\end{enumerate}
\end{proposition}

\begin{proof}
Theo \sref[I]{I.5.5.12}, chỉ cần chứng minh (i), (ii) và (iii).
Nhưng (i) chính là Ví dụ~\sref{II.1.2.2}, còn (ii) chính là \sref{II.1.3.5}; cuối cùng, (iii) suy ra từ \sref{II.1.5.1}, vì $X_{(S')}$ đồng nhất với tích $X\times_Y Y_{(S')}$ \sref[I]{I.3.3.11}.
\end{proof}

\begin{corollary}[1.6.3]
\label{II.1.6.3}
Nếu $X$ là một lược đồ affine và $Y$ là một lược đồ, thì mọi cấu xạ $X\to Y$ đều affine.
\end{corollary}

\begin{proposition}[1.6.4]
\label{II.1.6.4}
Cho $Y$ là một tiền lược đồ Noether địa phương và $f:X\to Y$ là một cấu xạ kiểu hữu hạn.
Điều kiện cần và đủ để $f$ affine là $f_\red$ affine.
\end{proposition}

\begin{proof}
Theo \sref{II.1.6.2}[vi], ta chỉ cần chứng minh tính đủ của điều kiện.
Chỉ cần chứng minh rằng nếu $Y$ affine và Noether thì $X$ affine; nhưng khi đó $Y_\red$ affine, nên theo giả thiết $X_\red$ cũng affine.
Mặt khác, $X$ là Noether, do đó kết luận suy ra từ \sref[I]{I.6.1.7}.
\end{proof}

\subsection{Phân thớ vectơ liên kết với một bó môđun}
\label{subsection:II.1.7}

\begin{env}[1.7.1]
\label{II.1.7.1}
Cho $A$ là một vành và $E$ là một $A$-môđun.
Nhắc lại rằng \emph{đại số đối xứng} trên $E$, ký hiệu $\bb{S}(E)$ (hoặc $\bb{S}_A(E)$), là đại số thương của đại số tensor $\bb{T}(E)$ theo iđêan hai phía sinh bởi các phần tử $x\otimes y-y\otimes x$, trong đó $x$ và $y$ chạy khắp $E$.
Đại số $\bb{S}(E)$ được đặc trưng bởi tính chất phổ dụng sau: nếu $\sigma$ là ánh xạ chính tắc $E\to\bb{S}(E)$ (thu được bằng cách hợp thành $E\to\bb{T}(E)$ với ánh xạ chính tắc $\bb{T}(E)\to\bb{S}(E)$), thì mọi ánh xạ $A$-tuyến tính $E\to B$, trong đó $B$ là một $A$-đại số \emph{giao hoán}, đều phân tích duy nhất thành
\[
  E\xrightarrow{\sigma}\bb{S}(E)\xrightarrow{g}B,
\]
trong đó $g$ là một $A$-đồng cấu \emph{đại số}.
Từ đặc trưng này suy ra ngay rằng, với hai $A$-môđun $E$ và $F$, ta có
\[
  \bb{S}(E\oplus F)=\bb{S}(E)\otimes\bb{S}(F)
\]
\oldpage[II]{15}
sai khác một đẳng cấu chính tắc; hơn nữa, $\bb{S}(E)$ là một hàm tử hiệp biến theo $E$, từ phạm trù các $A$-môđun tới phạm trù các $A$-đại số giao hoán; cuối cùng, đặc trưng trên cũng cho thấy rằng nếu $E=\varinjlim E_\lambda$, thì $\bb{S}(E)=\varinjlim\bb{S}(E_\lambda)$ sai khác một đẳng cấu chính tắc.
Theo lối nói lạm dụng, một tích $\sigma(x_1)\sigma(x_2)\cdots\sigma(x_n)$, trong đó $x_i\in E$, thường được ký hiệu là $x_1x_2\cdots x_n$ nếu không gây nhầm lẫn.
Đại số $\bb{S}(E)$ là \emph{phân bậc}, trong đó $\bb{S}_n(E)$ là tập hợp các tổ hợp tuyến tính của những tích gồm $n$ phần tử của $E$ ($n\geq0$); đại số $\bb{S}(A)$ đẳng cấu chính tắc với đại số đa thức $A[T]$ theo một biến bất định, còn đại số $\bb{S}(A^n)$ đẳng cấu chính tắc với đại số đa thức $A[T_1,\dots,T_n]$ theo $n$ biến bất định.
\end{env}

\begin{env}[1.7.2]
\label{II.1.7.2}
Cho $\vphi:A\to B$ là một đồng cấu vành.
Nếu $F$ là một $B$-môđun, thì ánh xạ chính tắc $F\to\bb{S}(F)$ cho một ánh xạ chính tắc $F_{[\vphi]}\to\bb{S}(F)_{[\vphi]}$, do đó ánh xạ này phân tích thành $F_{[\vphi]}\to\bb{S}(F_{[\vphi]})\to\bb{S}(F)_{[\vphi]}$; đồng cấu chính tắc $\bb{S}(F_{[\vphi]})\to\bb{S}(F)_{[\vphi]}$ là toàn ánh, nhưng không nhất thiết song ánh.
Nếu $E$ là một $A$-môđun, thì mỗi di-đồng cấu $E\to F$ (nghĩa là mỗi $A$-đồng cấu $E\to F_{[\vphi]}$) do đó cho một cách chính tắc một $A$-đồng cấu đại số $\bb{S}(E)\to\bb{S}(F_{[\vphi]})\to\bb{S}(F)_{[\vphi]}$, nghĩa là một di-đồng cấu đại số $\bb{S}(E)\to\bb{S}(F)$.

Với cùng các ký hiệu, đối với mỗi $A$-môđun $E$, $\bb{S}(E\otimes_A B)$ được đồng nhất chính tắc với đại số $\bb{S}(E)\otimes_A B$; điều này suy ra ngay từ tính chất phổ dụng của $\bb{S}(E)$ \sref{II.1.7.1}.
\end{env}

\begin{env}[1.7.3]
\label{II.1.7.3}
Cho $R$ là một tập con nhân của vành $A$; áp dụng \sref{II.1.7.2} cho vành $B=R^{-1}A$ và nhớ rằng $R^{-1}E=E\otimes_A R^{-1}A$, ta thấy $\bb{S}(R^{-1}E)=R^{-1}\bb{S}(E)$ sai khác một đẳng cấu chính tắc.
Hơn nữa, nếu $R'\supset R$ là một tập con nhân thứ hai của $A$, thì sơ đồ
\[
  \xymatrix{
    {R^{-1}E}\ar[r]\ar[d] & {{R'}^{-1}E}\ar[d]\\
    {\bb{S}(R^{-1}E)}\ar[r] & {\bb{S}({R'}^{-1}E)}
  }
\]
giao hoán.
\end{env}

\begin{env}[1.7.4]
\label{II.1.7.4}
Bây giờ cho $(S,\sh{A})$ là một không gian vành hóa, và cho $\sh{E}$ là một $\sh{A}$-môđun trên $S$.
Nếu với mỗi tập mở $U\subset S$ ta gắn $\Gamma(U,\sh{A})$-đại số $\bb{S}(\Gamma(U,\sh{E}))$, thì nhờ tính hàm tử của $\bb{S}$ \sref{II.1.7.2}, ta thu được một tiền bó đại số; bó liên kết, được ký hiệu là $\bb{S}(\sh{E})$ hoặc $\bb{S}_\sh{A}(\sh{E})$, được gọi là \emph{$\sh{A}$-đại số đối xứng của $\sh{A}$-môđun $\sh{E}$}.
Từ \sref{II.1.7.1} suy ra ngay rằng $\bb{S}(\sh{E})$ là nghiệm của một bài toán phổ dụng: mỗi đồng cấu các $\sh{A}$-môđun $\sh{E}\to\sh{B}$, trong đó $\sh{B}$ là một $\sh{A}$-đại số giao hoán, phân tích duy nhất thành $\sh{E}\to\bb{S}(\sh{E})\to\sh{B}$, trong đó mũi tên thứ hai là một đồng cấu các $\sh{A}$-đại số.
Do đó có một tương ứng song ánh giữa các đồng cấu $\sh{E}\to\sh{B}$ của $\sh{A}$-môđun và các đồng cấu $\bb{S}(\sh{E})\to\sh{B}$ của $\sh{A}$-đại số.
Đặc biệt, mỗi đồng cấu $u:\sh{E}\to\sh{F}$ của các $\sh{A}$-môđun xác định một đồng cấu $\bb{S}(u):\bb{S}(\sh{E})\to\bb{S}(\sh{F})$ của các $\sh{A}$-đại số, và vì thế $\bb{S}(\sh{E})$ là một hàm tử hiệp biến theo $\sh{E}$.
\oldpage[II]{16}
Theo \sref{II.1.7.2} và vì $\bb{S}$ giao hoán với các giới hạn quy nạp, ta có $(\bb{S}(\sh{E}))_s=\bb{S}(\sh{E}_s)$ đối với mọi điểm $s\in S$.
Nếu $\sh{E}$ và $\sh{F}$ là hai $\sh{A}$-môđun, thì $\bb{S}(\sh{E}\oplus\sh{F})$ được đồng nhất chính tắc với $\bb{S}(\sh{E})\otimes_\sh{A}\bb{S}(\sh{F})$, như thấy được trên các tiền bó tương ứng.

Cũng lưu ý rằng $\bb{S}(\sh{E})$ là một $\sh{A}$-đại số phân bậc, là tổng trực tiếp vô hạn của các $\bb{S}_n(\sh{E})$, trong đó $\sh{A}$-môđun $\bb{S}_n(\sh{E})$ là bó liên kết với tiền bó $U\mapsto\bb{S}_n(\Gamma(U,\sh{E}))$.
Đặc biệt, lấy $\sh{E}=\sh{A}$, ta thấy $\bb{S}_\sh{A}(\sh{A})$ được đồng nhất với $\sh{A}[T]=\sh{A}\otimes_\bb{Z}\bb{Z}[T]$ ($T$ là một biến bất định và $\bb{Z}$ được xem như một bó hằng).
\end{env}

\begin{env}[1.7.5]
\label{II.1.7.5}
Cho $(T,\sh{B})$ là một không gian vành hóa thứ hai và $f$ là một cấu xạ $(S,\sh{A})\to(T,\sh{B})$.
Nếu $\sh{F}$ là một $\sh{B}$-môđun, thì $\bb{S}(f^*(\sh{F}))$ được đồng nhất chính tắc với $f^*(\bb{S}(\sh{F}))$; thật vậy, nếu $f=(\psi,\theta)$, thì theo định nghĩa \sref[0]{0.4.3.1},
\[
  \bb{S}(f^*(\sh{F}))=\bb{S}(\psi^*(\sh{F})\otimes_{\psi^*(\sh{B})}\sh{A})=\bb{S}(\psi^*(\sh{F}))\otimes_{\psi^*(\sh{B})}\sh{A}
\]
\sref{II.1.7.2}; đối với mỗi tập mở $U$ của $S$ và mỗi tiết diện $h$ của $\bb{S}(\psi^*(\sh{F}))$ trên $U$, trong một lân cận $V$ của mỗi điểm $s\in U$, $h$ trùng với một phần tử của $\bb{S}(\Gamma(V,\psi^*(\sh{F})))$; nếu viện đến định nghĩa của $\psi^*(\sh{F})$ \sref[0]{0.3.7.1} và lưu ý rằng mỗi phần tử của $\bb{S}(E)$, với $E$ là một môđun, là tổ hợp tuyến tính của một số hữu hạn các tích những phần tử của $E$, thì ta thấy tồn tại một lân cận $W$ của $\psi(s)$ trong $T$, một tiết diện $h'$ của $\bb{S}(\sh{F})$ trên $W$ và một lân cận $V'\subset V\cap\psi^{-1}(W)$ của $s$ sao cho $h$ trùng với $t\mapsto h'(\psi(t))$ trên $V'$; từ đó suy ra khẳng định.
\end{env}

\begin{proposition}[1.7.6]
\label{II.1.7.6}
Cho $A$ là một vành, $S=\Spec(A)$ là phổ nguyên tố của nó, và $\sh{E}=\widetilde{M}$ là $\sh{O}_S$-môđun liên kết với một $A$-môđun $M$; khi đó $\sh{O}_S$-đại số $\bb{S}(\sh{E})$ liên kết với $A$-đại số $\bb{S}(M)$.
\end{proposition}

\begin{proof}
Thật vậy, với mọi $f\in A$, $\bb{S}(M_f)=(\bb{S}(M))_f$ \sref{II.1.7.3}, nên mệnh đề suy ra từ các định nghĩa \sref[I]{I.1.3.4}.
\end{proof}

\begin{corollary}[1.7.7]
\label{II.1.7.7}
Nếu $S$ là một tiền lược đồ và $\sh{E}$ là một $\sh{O}_S$-môđun quasi-coherent, thì $\sh{O}_S$-đại số $\bb{S}(\sh{E})$ là quasi-coherent.
Hơn nữa, nếu $\sh{E}$ hữu hạn sinh, thì mỗi $\sh{O}_S$-môđun $\bb{S}_n(\sh{E})$ đều hữu hạn sinh.
\end{corollary}

\begin{proof}
Khẳng định thứ nhất là hệ quả trực tiếp của \sref{II.1.7.6} và \sref[I]{I.1.4.1}; khẳng định thứ hai suy ra từ việc, nếu $E$ là một $A$-môđun hữu hạn sinh, thì $\bb{S}_n(E)$ là một $A$-môđun hữu hạn sinh; sau đó áp dụng \sref[I]{I.1.3.13}.
\end{proof}

\begin{definition}[1.7.8]
\label{II.1.7.8}
Cho $\sh{E}$ là một $\sh{O}_S$-môđun quasi-coherent.
Ta gọi phổ \sref{II.1.3.1} của $\sh{O}_S$-đại số quasi-coherent $\bb{S}(\sh{E})$ là \emph{phân thớ vectơ trên $S$ xác định bởi $\sh{E}$}, và ký hiệu nó là $\bb{V}(\sh{E})$.
\end{definition}

Theo \sref{II.1.2.7}, đối với mỗi $S$-tiền lược đồ $X$, có một tương ứng song ánh chính tắc giữa các $S$-cấu xạ $X\to\bb{V}(\sh{E})$ và các đồng cấu \emph{$\sh{O}_S$-đại số $\bb{S}(\sh{E})\to\sh{A}(X)$}, do đó cũng giữa các $S$-cấu xạ này và các đồng cấu \emph{$\sh{O}_S$-môđun $\sh{E}\to\sh{A}(X)=f_*(\sh{O}_X)$} (trong đó $f$ là cấu xạ cấu trúc $X\to S$).
Đặc biệt:
\begin{env}[1.7.9]
\label{II.1.7.9}
Lấy $X$ là tiền lược đồ con của $S$ cảm sinh trên một \emph{tập mở $U\subset S$}.
Khi đó, các $S$-cấu xạ $U\to\bb{V}(\sh{E})$ chính là các $U$-tiết diện \sref[I]{I.2.5.5} của $U$-tiền lược đồ do $\bb{V}(\sh{E})$ cảm sinh trên tập mở $p^{-1}(U)$ (trong đó $p$ là cấu xạ cấu trúc $\bb{V}(\sh{E})\to S$).
Theo điều vừa thấy, các $U$-tiết diện này tương ứng song ánh với các đồng cấu $\sh{O}_S$-môđun $\sh{E}\to j_*(\sh{O}_S|U)$ (trong đó $j$ là phép nhúng chính tắc $U\to S$), hoặc,
\oldpage[II]{17}
tương đương, theo \sref[0]{0.4.4.3}, với các đồng cấu $(\sh{O}_S|U)$-môđun $j^*(\sh{E})=\sh{E}|U\to\sh{O}_S|U$.
Hơn nữa, hiển nhiên hạn chế lên một tập mở $U'\subset U$ của một $S$-cấu xạ $U\to\bb{V}(\sh{E})$ tương ứng với hạn chế lên $U'$ của đồng cấu tương ứng $\sh{E}|U\to\sh{O}_S|U$.
Ta kết luận rằng \emph{bó các mầm của những $S$-tiết diện của $\bb{V}(\sh{E})$ được đồng nhất chính tắc với đối ngẫu $\dual{\sh{E}}$ của $\sh{E}$}.

Đặc biệt, lấy $X=U=S$, đồng cấu \emph{không} $\sh{E}\to\sh{O}_S$ tương ứng với một $S$-tiết diện chính tắc của $\bb{V}(\sh{E})$, gọi là \emph{$S$-tiết diện không} (xem \sref{II.8.3.3}).
\end{env}

\begin{env}[1.7.10]
\label{II.1.7.10}
Bây giờ lấy $X$ là phổ $\{\xi\}$ của một trường $K$; khi đó cấu xạ cấu trúc $f:X\to S$ tương ứng với một đơn cấu $\kres(s)\to K$, trong đó $s=f(\xi)$ \sref[I]{I.2.4.6}; vì thế, các $S$-cấu xạ $\{\xi\}\to\bb{V}(\sh{E})$ chính là các \emph{điểm hình học của $\bb{V}(\sh{E})$ nhận giá trị trong mở rộng $K$ của $\kres(s)$} \sref[I]{I.3.4.5}, những điểm có ảnh thuộc $p^{-1}(s)$.
Theo \sref{II.1.7.8}, tập hợp các điểm này, còn có thể gọi là \emph{thớ hình học hữu tỉ trên $K$ của $\bb{V}(\sh{E})$ bên trên điểm $s$}, được đồng nhất với tập hợp các đồng cấu $\sh{O}_S$-môđun $\sh{E}\to f_*(\sh{O}_X)$, hoặc tương đương \sref[0]{0.4.4.3}, với tập hợp các đồng cấu $\sh{O}_X$-môđun $f^*(\sh{E})\to\sh{O}_X=K$.
Nhưng theo định nghĩa \sref[0]{0.4.3.1}, ta có $f^*(\sh{E})=\sh{E}_s\otimes_{\sh{O}_s}K=\sh{E}^s\otimes_{\kres(s)}K$, với $\sh{E}^s=\sh{E}_s/\mathfrak{m}_s\sh{E}_s$; do đó thớ hình học của $\bb{V}(\sh{E})$, hữu tỉ trên $K$ và nằm bên trên $s$, được đồng nhất với \emph{đối ngẫu của $K$-không gian vectơ $\sh{E}^s\otimes_{\kres(s)}K$}; nếu $\sh{E}^s$ hoặc $K$ có chiều hữu hạn trên $\kres(s)$, thì đối ngẫu này cũng được đồng nhất với $\dual{(\sh{E}^s)}\otimes_{\kres(s)}K$, trong đó $\dual{(\sh{E}^s)}$ ký hiệu đối ngẫu của $\kres(s)$-không gian vectơ $\sh{E}^s$.
\end{env}

\begin{proposition}[1.7.11]
\label{II.1.7.11}
\medskip\noindent
\begin{enumerate}
  \item[{\rm(i)}] $\bb{V}(\sh{E})$ là một hàm tử phản biến theo $\sh{E}$, từ phạm trù các $\sh{O}_S$-môđun quasi-coherent tới phạm trù các $S$-lược đồ affine.
  \item[{\rm(ii)}] Nếu $\sh{E}$ là một $\sh{O}_S$-môđun hữu hạn sinh, thì $\bb{V}(\sh{E})$ có kiểu hữu hạn trên $S$.
  \item[{\rm(iii)}] Nếu $\sh{E}$ và $\sh{F}$ là hai $\sh{O}_S$-môđun quasi-coherent, thì $\bb{V}(\sh{E}\oplus\sh{F})$ được đồng nhất chính tắc với $\bb{V}(\sh{E})\times_S\bb{V}(\sh{F})$.
  \item[{\rm(iv)}] Cho $g:S'\to S$ là một cấu xạ; đối với mỗi $\sh{O}_S$-môđun quasi-coherent $\sh{E}$, $\bb{V}(g^*(\sh{E}))$ được đồng nhất chính tắc với $\bb{V}(\sh{E})_{(S')}=\bb{V}(\sh{E})\times_S S'$.
  \item[{\rm(v)}] Một đồng cấu toàn ánh $\sh{E}\to\sh{F}$ của các $\sh{O}_S$-môđun quasi-coherent tương ứng với một phép nhúng đóng $\bb{V}(\sh{F})\to\bb{V}(\sh{E})$.
\end{enumerate}
\end{proposition}

\begin{proof}
(i) là hệ quả trực tiếp của \sref{II.1.2.7}, khi lưu ý rằng mỗi đồng cấu $\sh{O}_S$-môđun $\sh{E}\to\sh{F}$ xác định chính tắc một đồng cấu $\sh{O}_S$-đại số $\bb{S}(\sh{E})\to\bb{S}(\sh{F})$.
(ii) suy ra ngay từ các định nghĩa \sref[I]{I.6.3.1} và từ việc, nếu $E$ là một $A$-môđun hữu hạn sinh, thì $\bb{S}(E)$ là một $A$-đại số hữu hạn sinh.
Để chứng minh (iii), chỉ cần xuất phát từ đẳng cấu chính tắc $\bb{S}(\sh{E}\oplus\sh{F})\isoto\bb{S}(\sh{E})\otimes_{\sh{O}_S}\bb{S}(\sh{F})$ \sref{II.1.7.4} rồi áp dụng \sref{II.1.4.6}.
Tương tự, để chứng minh (iv), chỉ cần xuất phát từ đẳng cấu chính tắc $\bb{S}(g^*(\sh{E}))\isoto g^*(\bb{S}(\sh{E}))$ \sref{II.1.7.5} rồi áp dụng \sref{II.1.5.2}.
Cuối cùng, để thiết lập (v), chỉ cần nhận xét rằng nếu đồng cấu $\sh{E}\to\sh{F}$ là toàn ánh, thì đồng cấu $\sh{O}_S$-đại số tương ứng $\bb{S}(\sh{E})\to\bb{S}(\sh{F})$ cũng toàn ánh, và kết luận suy ra từ \sref{II.1.4.10}.
\end{proof}

\begin{env}[1.7.12]
\label{II.1.7.12}
Đặc biệt, lấy $\sh{E}=\sh{O}_S$; tiền lược đồ $\bb{V}(\sh{O}_S)$ là $S$-lược đồ affine, phổ của $\sh{O}_S$-đại số $\bb{S}(\sh{O}_S)$, được đồng nhất với $\sh{O}_S$-đại số $\sh{O}_S[T]=\sh{O}_S\otimes_\bb{Z}\bb{Z}[T]$
\oldpage[II]{18}
($T$ là một biến bất định); điều này hiển nhiên khi $S=\Spec(\bb{Z})$, theo \sref{II.1.7.6}, và từ đó chuyển sang trường hợp tổng quát bằng cách xét cấu xạ cấu trúc $S\to\Spec(\bb{Z})$ và dùng \sref{II.1.7.11}[iv].
Do kết quả này, ta còn đặt $\bb{V}(\sh{O}_S)=S[T]$, và do đó có công thức
\[
\label{II.1.7.12.1}
  S[T]=S\otimes_\bb{Z}\bb{Z}[T].
  \tag{1.7.12.1}
\]

Sự đồng nhất bó các mầm của những $S$-tiết diện của $S[T]$ với $\sh{O}_S$, đã thấy trong \sref[I]{I.3.3.15}, lại xuất hiện ở đây trong một bối cảnh tổng quát hơn như một trường hợp riêng của \sref{II.1.7.9}.
\end{env}

\begin{env}[1.7.13]
\label{II.1.7.13}
Đối với mỗi $S$-tiền lược đồ $X$, ta đã thấy trong \sref{II.1.7.8} rằng $\Hom_S(X,S[T])$ được đồng nhất chính tắc với $\Hom_{\sh{O}_S}(\sh{O}_S,\sh{A}(X))$, mà tập sau đẳng cấu chính tắc với $\Gamma(S,\sh{A}(X))$, và do đó được trang bị một cấu trúc vành; hơn nữa, với mỗi $S$-cấu xạ $h:X\to Y$ ứng với một đồng cấu vành $\Gamma(\sh{A}(h)):\Gamma(S,\sh{A}(Y))\to\Gamma(S,\sh{A}(X))$ \sref{II.1.1.2}.
Khi trang bị cho $\Hom_S(X,S[T])$ cấu trúc vành vừa định nghĩa, ta có thể xem $\Hom_S(X,S[T])$ như một \emph{hàm tử phản biến theo $X$}, từ phạm trù các $S$-tiền lược đồ tới phạm trù các vành.
Mặt khác, $\Hom_S(X,\bb{V}(\sh{E}))$ cũng được đồng nhất với $\Hom_{\sh{O}_S}(\sh{E},\sh{A}(X))$ (trong đó $\sh{A}(X)$ được xem như một \emph{$\sh{O}_S$-môđun}); do đó, ta có thể trang bị chính tắc cho tập này một cấu trúc \emph{môđun} trên vành $\Hom_S(X,S[T])$, và như trên, ta thấy rằng cặp
\[
  (\Hom_S(X,S[T]),\Hom_S(X,\bb{V}(\sh{E})))
\]
là một hàm tử phản biến theo $X$, nhận giá trị trong phạm trù mà các vật là những cặp $(A,M)$ gồm một vành $A$ và một $A$-môđun $M$, còn các cấu xạ là các di-đồng cấu.

Ta sẽ diễn giải những sự kiện này bằng cách nói rằng $S[T]$ là một \emph{$S$-lược đồ vành} và $\bb{V}(\sh{E})$ là một \emph{$S$-lược đồ môđun} trên $S$-lược đồ vành $S[T]$ (xem Chương~0, \textsection8).
\end{env}

\begin{env}[1.7.14]
\label{II.1.7.14}
Ta sẽ thấy rằng cấu trúc $S$-lược đồ môđun được xác định trên $S$-lược đồ $\bb{V}(\sh{E})$ cho phép dựng lại $\sh{O}_S$-môđun $\sh{E}$ sai khác một đẳng cấu duy nhất: để làm điều đó, ta sẽ chứng minh rằng $\sh{E}$ đẳng cấu chính tắc với một $\sh{O}_S$-môđun con của $\bb{S}(\sh{E})=\sh{A}(\bb{V}(\sh{E}))$, được xác định nhờ cấu trúc này.
Thật vậy \sref{II.1.7.4}, tập $\Hom_{\sh{O}_S}(\bb{S}(\sh{E}),\sh{A}(X))$ các đồng cấu \emph{$\sh{O}_S$-đại số} được đồng nhất chính tắc với $\Hom_{\sh{O}_S}(\sh{E},\sh{A}(X))$, tập các đồng cấu \emph{$\sh{O}_S$-môđun}: nếu $h$ và $h'$ là hai phần tử của tập sau, $s_i$ ($1\leq i\leq n$) là các tiết diện của $\sh{E}$ trên một tập mở $U\subset S$, và $t$ là một tiết diện của $\sh{A}(X)$ trên $U$, thì theo định nghĩa ta có
\[
  (h+h')(s_1s_2\cdots s_n)=\prod_{i=1}^n(h(s_i)+h'(s_i))
\]
và
\[
  (t\cdot h)(s_1s_2\cdots s_n)=t^n\prod_{i=1}^n h(s_i).
\]

Như vậy, nếu $z$ là một tiết diện của $\bb{S}(\sh{E})$ trên $U$, thì $h\mapsto h(z)$ là một ánh xạ từ $\Hom_S(X,\bb{V}(\sh{E}))=\Hom_{\sh{O}_S}(\bb{S}(\sh{E}),\sh{A}(X))$ vào $\Gamma(U,\sh{A}(X))$.
Ta sẽ
\oldpage[II]{19}
chứng minh rằng $\sh{E}$ được đồng nhất với môđun con của $\bb{S}(\sh{E})$ sao cho \emph{đối với mọi tập mở $U\subset S$, mọi tiết diện $z$ trên $U$ của $\sh{O}_S$-môđun con này và mọi $S$-tiền lược đồ $X$, ánh xạ $h\mapsto h(z)$ từ $\Hom_{\sh{O}_S|U}(\bb{S}(\sh{E})|U,\sh{A}(X)|U)$ vào $\Gamma(U,\sh{A}(X))$ là một đồng cấu các $\Gamma(U,\sh{A}(X))$-môđun}.

Thật vậy, hiển nhiên $\sh{E}$ có tính chất này; để chứng minh điều đảo lại, ta có thể quy về việc chứng minh rằng, khi $S=\Spec(A)$ và $\sh{E}=\widetilde{M}$, một tiết diện $z$ của $\bb{S}(\sh{E})$ trên $S$ có tính chất nêu trên (với $U=S$) nhất thiết là một tiết diện của $\sh{E}$; khi đó ta có $z=\sum_{n=0}^{\infty}z_n$, trong đó $z_n\in\bb{S}_n(M)$, và cần chứng minh rằng $z_n=0$ với $n\neq1$.
Đặt $B=\bb{S}(M)$ và lấy $X$ là tiền lược đồ $\Spec(B[T])$, trong đó $T$ là một biến bất định.
Tập $\Hom_{\sh{O}_S}(\bb{S}(\sh{E}),\sh{A}(X))$ được đồng nhất với tập các đồng cấu vành $h:B\to B[T]$ \sref[I]{I.1.3.13}, và theo điều đã thấy ở trên, ta có $(T\cdot h)(z)=\sum_{n=0}^{\infty}T^n h(z_n)$; giả thiết về $z$ kéo theo $\sum_{n=0}^{\infty}T^n h(z_n)=T\cdot\sum_{n=0}^{\infty}h(z_n)$ đối với mọi đồng cấu $h$.
Đặc biệt, lấy $h$ là phép nhúng chính tắc, ta được $\sum_{n=0}^{\infty}T^n z_n=T\cdot\sum_{n=0}^{\infty}z_n$, và điều này quả thực kéo theo kết luận $z_n=0$ với $n\neq1$.
\end{env}

\begin{proposition}[1.7.15]
\label{II.1.7.15}
Cho $Y$ là một tiền lược đồ có không gian nền Noether, hoặc là một lược đồ quasi-compact.
Mỗi $Y$-lược đồ affine $X$ có kiểu hữu hạn trên $Y$ đều $Y$-đẳng cấu với một $Y$-lược đồ con đóng của một $Y$-lược đồ dạng $\bb{V}(\sh{E})$, trong đó $\sh{E}$ là một $\sh{O}_Y$-môđun quasi-coherent hữu hạn sinh.
\end{proposition}

\begin{proof}
$\sh{O}_Y$-đại số quasi-coherent $\sh{A}(X)$ có kiểu hữu hạn \sref{II.1.3.7}.
Các giả thiết kéo theo rằng $\sh{A}(X)$ được sinh bởi một $\sh{O}_Y$-môđun con quasi-coherent hữu hạn sinh $\sh{E}$ \sref[I]{I.9.6.5}; theo định nghĩa, điều này kéo theo đồng cấu chính tắc $\bb{S}(\sh{E})\to\sh{A}(X)$, vốn mở rộng chính tắc phép nhúng $\sh{E}\to\sh{A}(X)$, là \emph{toàn ánh}; khi đó kết luận suy ra từ \sref{II.1.4.10}.
\end{proof}

\section{Các phổ nguyên tố thuần nhất}
\label{section:II.2}

\subsection{Đại cương về các vành và môđun phân bậc}
\label{subsection:II.2.1}

\begin{notation}[2.1.1]
\label{II.2.1.1}
Cho một vành $S$ \emph{phân bậc dương}; ta ký hiệu bởi $S_n$ phần của $S$ gồm các phần tử thuần nhất bậc $n$ ($n\geq0$), và bởi $S_+$ tổng (trực tiếp) của các $S_n$ với $n>0$; ta có $1\in S_0$, $S_0$ là một vành con của $S$, $S_+$ là một iđêan phân bậc của $S$, và $S$ là tổng trực tiếp của $S_0$ và $S_+$.
Nếu $M$ là một môđun \emph{phân bậc} trên $S$ (với các bậc dương hoặc âm), ta cũng ký hiệu bởi $M_n$ $S_0$-môđun gồm các phần tử thuần nhất bậc $n$ của $M$ (với $n\in\bb{Z}$).

Đối với mỗi số nguyên $d>0$, ta ký hiệu bởi $S^{(d)}$ tổng trực tiếp của các $S_{nd}$; khi xem các phần tử của $S_{nd}$ là thuần nhất bậc $n$, các $S_{nd}$ xác định trên $S^{(d)}$ một cấu trúc vành phân bậc.

Đối với mỗi số nguyên $k$ sao cho $0\leq k\leq d-1$, ta ký hiệu bởi $M^{(d,k)}$ tổng trực tiếp
\oldpage[II]{20}
của các $M_{nd+k}$ ($n\in\bb{Z}$); đó là một $S^{(d)}$-môđun phân bậc khi ta xem các phần tử của $M_{nd+k}$ là thuần nhất bậc $n$.
Ta viết $M^{(d)}$ thay cho $M^{(d,0)}$.

Với các ký hiệu trên, đối với mỗi số nguyên $n$ (dương hoặc âm), ta ký hiệu bởi $M(n)$ $S$-môđun phân bậc được xác định bởi $(M(n))_k=M_{n+k}$ với mọi $k\in\bb{Z}$.
Đặc biệt, $S(n)$ là một $S$-môđun phân bậc sao cho $(S(n))_k=S_{n+k}$, với quy ước $S_n=0$ khi $n<0$.
Ta nói một $S$-môđun phân bậc $M$ là \emph{tự do} nếu nó đẳng cấu, với tư cách \emph{môđun phân bậc}, với một tổng trực tiếp các môđun dạng $S(n)$; vì $S(n)$ là một $S$-môđun cyclic sinh bởi phần tử $1$ của $S$, được xem như một phần tử bậc $-n$, điều này tương đương với việc $M$ có một \emph{cơ sở} trên $S$ gồm các phần tử \emph{thuần nhất}.

Ta nói một $S$-môđun phân bậc $M$ \emph{có một trình bày hữu hạn} nếu tồn tại một dãy khớp $P\to Q\to M\to0$, trong đó $P$ và $Q$ là các tổng trực tiếp hữu hạn của những môđun dạng $S(n)$ và các đồng cấu đều có bậc $0$ (xem \sref{II.2.1.2}).
\end{notation}

\begin{env}[2.1.2]
\label{II.2.1.2}
Cho $M$ và $N$ là hai $S$-môđun phân bậc; ta xác định trên $M\otimes_S N$ một cấu trúc $S$-môđun \emph{phân bậc} như sau.
Trên tích tenxơ $M\otimes_\bb{Z}N$, ta có thể xác định một cấu trúc $\bb{Z}$-môđun phân bậc (trong đó $\bb{Z}$ được phân bậc bởi $\bb{Z}_0=\bb{Z}$, $\bb{Z}_n=0$ khi $n\neq0$) bằng cách đặt
\[
  (M\otimes_\bb{Z}N)_q
  =\bigoplus_{m+n=q}M_m\otimes_\bb{Z}N_n.
\]
Vì $M$ và $N$ lần lượt là tổng trực tiếp của các $M_m$ và các $N_n$, ta biết rằng $M\otimes_\bb{Z}N$ được đồng nhất chính tắc với tổng trực tiếp của tất cả các $M_m\otimes_\bb{Z}N_n$.
Khi đó, ta có $M\otimes_S N=(M\otimes_\bb{Z}N)/P$, trong đó $P$ là $\bb{Z}$-môđun con của $M\otimes_\bb{Z}N$ sinh bởi các phần tử $(xs)\otimes y-x\otimes(sy)$ với $x\in M$, $y\in N$ và $s\in S$; hiển nhiên $P$ là một $\bb{Z}$-môđun con \emph{phân bậc} của $M\otimes_\bb{Z}N$, và ta thấy ngay rằng khi chuyển sang môđun thương, ta thật sự thu được một cấu trúc $S$-môđun phân bậc trên $M\otimes_S N$.

Đối với hai $S$-môđun phân bậc $M$ và $N$, nhắc lại rằng một đồng cấu $u:M\to N$ của các $S$-môđun được gọi là \emph{có bậc $k$} nếu $u(M_j)\subset N_{j+k}$ với mọi $j\in\bb{Z}$.
Nếu $H_n$ ký hiệu tập hợp tất cả các đồng cấu bậc $n$ từ $M$ vào $N$, thì ta ký hiệu bởi $\Hom_S(M,N)$ \emph{tổng} (trực tiếp) của các $H_n$ ($n\in\bb{Z}$) trong $S$-môđun $H$ gồm tất cả các đồng cấu (của $S$-môđun) từ $M$ vào $N$; nói chung, $\Hom_S(M,N)$ không bằng môđun sau cùng này.
Tuy nhiên, ta có $H=\Hom_S(M,N)$ khi $M$ \emph{hữu hạn sinh}; thật vậy, khi đó có thể giả sử $M$ được sinh bởi một số hữu hạn phần tử thuần nhất $x_i$ ($1\leq i\leq n$), và mỗi đồng cấu $u\in H$ được viết duy nhất dưới dạng $\sum_{k\in\bb{Z}}u_k$, trong đó, với mỗi $k$, $u_k(x_i)$ bằng thành phần thuần nhất bậc $k+\deg(x_i)$ của $u(x_i)$ ($1\leq i\leq n$), kéo theo $u_k=0$ trừ một số hữu hạn chỉ số; theo định nghĩa, $u_k\in H_k$, từ đó suy ra kết luận.

Các phần tử bậc $0$ của $\Hom_S(M,N)$ được gọi là \emph{các đồng cấu của các $S$-môđun phân bậc}.
Hiển nhiên $S_mH_n\subset H_{m+n}$, do đó các $H_n$ xác định trên $\Hom_S(M,N)$ một cấu trúc $S$-môđun phân bậc.

Từ các định nghĩa này suy ra ngay rằng
\[
\label{II.2.1.2.1}
  M(m)\otimes_S N(n)=(M\otimes_S N)(m+n),
\tag{2.1.2.1}
\]
\[
\label{II.2.1.2.2}
  \Hom_S(M(m),N(n))=(\Hom_S(M,N))(n-m),
\tag{2.1.2.2}
\]
đối với hai $S$-môđun phân bậc $M$ và $N$.
\oldpage[II]{21}
Cho $S$ và $S'$ là hai vành phân bậc; một \emph{đồng cấu các vành phân bậc} $\vphi:S\to S'$ là một đồng cấu vành sao cho $\vphi(S_n)\subset S_n'$ với mọi $n\in\bb{Z}$ (nói cách khác, $\vphi$ phải là một đồng cấu \emph{bậc $0$} của các $\bb{Z}$-môđun phân bậc).
Dữ kiện về một đồng cấu như vậy xác định trên $S'$ một cấu trúc $S$-môđun \emph{phân bậc}; khi được trang bị cấu trúc này cùng cấu trúc vành phân bậc của nó, $S'$ được gọi là một \emph{$S$-đại số phân bậc}.

Nếu $M$ là một $S$-môđun phân bậc, thì tích tenxơ $M\otimes_S S'$ của các $S$-môđun \emph{phân bậc} được trang bị một cách tự nhiên cấu trúc $S'$-môđun \emph{phân bậc}, với phân bậc được xác định như trên.
\end{env}

\begin{lemma}[2.1.3]
\label{II.2.1.3}
Cho $S$ là một vành phân bậc dương.
Điều kiện cần và đủ để một tập con $E$ của $S_+$, gồm các phần tử thuần nhất, sinh $S_+$ như một $S$-môđun là $E$ sinh $S$ như một $S_0$-đại số.
\end{lemma}

\begin{proof}
Điều kiện hiển nhiên là đủ; ta chứng minh nó là cần.
Ký hiệu $E_n$ (tương ứng, $E^n$) là tập hợp các phần tử của $E$ có bậc bằng $n$ (tương ứng, có bậc $\leq n$); chỉ cần chứng minh, bằng quy nạp theo $n>0$, rằng $S_n$ là $S_0$-môđun sinh bởi các phần tử bậc $n$ vốn là tích của các phần tử thuộc $E^n$.
Điều này hiển nhiên với $n=1$ nhờ giả thiết; giả thiết còn cho thấy rằng
\[
  S_n=\sum_{p=0}^{n-1}S_pE_{n-p},
\]
và phép quy nạp khi đó là trực tiếp.
\end{proof}

\begin{corollary}[2.1.4]
\label{II.2.1.4}
Điều kiện cần và đủ để $S_+$ là một iđêan hữu hạn sinh là $S$ là một $S_0$-đại số hữu hạn sinh.
\end{corollary}

\begin{proof}
Ta luôn có thể giả sử rằng một hệ sinh hữu hạn của $S_0$-đại số $S$ (tương ứng, của $S$-iđêan $S_+$) gồm các phần tử thuần nhất, bằng cách thay mỗi phần tử sinh đang xét bởi các thành phần thuần nhất của nó.
\end{proof}

\begin{corollary}[2.1.5]
\label{II.2.1.5}
Điều kiện cần và đủ để $S$ là Noether là $S_0$ là Noether và $S$ là một $S_0$-đại số hữu hạn sinh.
\end{corollary}

\begin{proof}
Điều kiện hiển nhiên là đủ; nó là cần vì $S_0$ đẳng cấu với $S/S_+$ và $S_+$ phải là một iđêan hữu hạn sinh \sref{II.2.1.4}.
\end{proof}

\begin{lemma}[2.1.6]
\label{II.2.1.6}
Cho $S$ là một vành phân bậc dương và là một $S_0$-đại số hữu hạn sinh.
Cho $M$ là một $S$-môđun phân bậc hữu hạn sinh.
Khi đó:
\begin{enumerate}
  \item[{\rm(i)}] Các $M_n$ là những $S_0$-môđun hữu hạn sinh, và tồn tại một số nguyên $n_0$ sao cho $M_n=0$ với $n\leq n_0$.
  \item[{\rm(ii)}] Tồn tại một số nguyên $n_1$ và một số nguyên $h>0$ sao cho, với mọi số nguyên $n\geq n_1$, ta có $M_{n+h}=S_hM_n$.
  \item[{\rm(iii)}] Với mọi cặp số nguyên $(d,k)$ sao cho $d>0$ và $0\leq k\leq d-1$, $M^{(d,k)}$ là một $S^{(d)}$-môđun hữu hạn sinh.
  \item[{\rm(iv)}] Với mọi số nguyên $d>0$, $S^{(d)}$ là một $S_0$-đại số hữu hạn sinh.
  \item[{\rm(v)}] Tồn tại một số nguyên $h>0$ sao cho $S_{mh}=(S_h)^m$ với mọi $m>0$.
  \item[{\rm(vi)}] Với mọi số nguyên $n>0$, tồn tại một số nguyên $m_0$ sao cho $S_m\subset S_+^n$ với mọi $m\geq m_0$.
\end{enumerate}
\end{lemma}

\begin{proof}
Ta có thể giả sử $S$ được sinh (như một $S_0$-đại số) bởi các phần tử thuần nhất $f_i$, có bậc $h_i$ ($1\leq i\leq r$), và $M$ được sinh (như một $S$-môđun) bởi các phần tử thuần nhất $x_j$ có bậc $k_j$ ($1\leq j\leq s$).
Hiển nhiên $M_n$ gồm các tổ hợp
\oldpage[II]{22}
tuyến tính, với hệ số trong $S_0$, của các phần tử $f_1^{\alpha_1}\cdots f_r^{\alpha_r}x_j$ sao cho các $\alpha_i$ là những số nguyên $\geq0$ thỏa mãn
\[
  k_j+\sum_i\alpha_i h_i=n.
\]
Với mỗi $j$, chỉ có hữu hạn hệ $(\alpha_i)$ thỏa mãn phương trình này, vì các $h_i$ đều $>0$; từ đó suy ra khẳng định thứ nhất của~(i), còn khẳng định thứ hai là hiển nhiên.

Mặt khác, cho $h$ là bội chung nhỏ nhất của các $h_i$ và đặt $g_i=f_i^{h/h_i}$ ($1\leq i\leq r$), sao cho tất cả các $g_i$ đều có bậc $h$; gọi $z_\mu$ là các phần tử của $M$ có dạng $f_1^{\alpha_1}\cdots f_r^{\alpha_r}x_j$, với $0\leq\alpha_i<h/h_i$ đối với $1\leq i\leq r$.
Các phần tử này có số lượng hữu hạn; gọi $n_1$ là bậc lớn nhất trong các bậc của chúng.
Hiển nhiên, với $n\geq n_1$, mỗi phần tử của $M_{n+h}$ là một tổ hợp tuyến tính của các $z_\mu$, trong đó các hệ số là những đơn thức bậc $>0$ theo các $g_i$; do đó ta có $M_{n+h}=S_hM_n$, điều này chứng minh~(ii).
Cũng theo cách ấy, ta thấy (với mọi $d>0$) rằng một phần tử của $M^{(d,k)}$ là một tổ hợp tuyến tính, với hệ số trong $S_0$, của các phần tử dạng $g^d f_1^{\alpha_1}\cdots f_r^{\alpha_r}x_j$, với $0\leq\alpha_i<d$, trong đó $g$ là một phần tử thuần nhất của $S$; từ đó suy ra~(iii).
(iv) khi đó suy ra từ~(iii) và Bổ đề~\sref{II.2.1.3}, bằng cách lấy $M=S_+$, vì $(S_+)^{(d)}=(S^{(d)})_+$.
Khẳng định~(v) suy ra từ~(ii) bằng cách lấy $M=S$.
Cuối cùng, với một $n$ cho trước, chỉ có hữu hạn hệ $(\alpha_i)$ sao cho $\alpha_i\geq0$ và $\sum_i\alpha_i<n$; nếu chọn $m_0$ lớn hơn nghiêm ngặt giá trị lớn nhất của $\sum_i\alpha_i h_i$ trong các hệ này, thì $S_m\subset S_+^n$ với $m\geq m_0$, điều này chứng minh~(vi).
\end{proof}

\begin{corollary}[2.1.7]
\label{II.2.1.7}
Nếu $S$ là Noether, thì $S^{(d)}$ cũng vậy với mọi số nguyên $d>0$.
\end{corollary}

\begin{proof}
Điều này suy ra từ \sref{II.2.1.5} và \sref{II.2.1.6}[iv].
\end{proof}

\begin{env}[2.1.8]
\label{II.2.1.8}
Cho $\mathfrak{p}$ là một iđêan nguyên tố \emph{phân bậc} của vành phân bậc $S$; do đó $\mathfrak{p}$ là tổng trực tiếp của các nhóm con $\mathfrak{p}_n=\mathfrak{p}\cap S_n$.
Giả sử \emph{$\mathfrak{p}$ không chứa $S_+$}.
Khi đó, nếu $f\in S_+$ không thuộc $\mathfrak{p}$, quan hệ $f^nx\in\mathfrak{p}$ tương đương với $x\in\mathfrak{p}$; đặc biệt, nếu $f\in S_d$ ($d>0$), thì với mọi $x\in S_{m-nd}$, quan hệ $f^nx\in\mathfrak{p}_m$ tương đương với $x\in\mathfrak{p}_{m-nd}$.
\end{env}

\begin{proposition}[2.1.9]
\label{II.2.1.9}
Cho $n_0$ là một số nguyên $>0$; với mỗi $n\geq n_0$, cho $\mathfrak{p}_n$ là một nhóm con của $S_n$.
Điều kiện cần và đủ để tồn tại một iđêan nguyên tố phân bậc $\mathfrak{p}$ của $S$, không chứa $S_+$ và sao cho $\mathfrak{p}\cap S_n=\mathfrak{p}_n$ với mọi $n\geq n_0$, là các điều kiện sau được thỏa mãn:
\begin{enumerate}
  \item[{\rm($1^\circ$)}] $S_m\mathfrak{p}_n\subset\mathfrak{p}_{m+n}$ với mọi $m\geq0$ và mọi $n\geq n_0$.
  \item[{\rm($2^\circ$)}] Với $m\geq n_0$, $n\geq n_0$, $f\in S_m$ và $g\in S_n$, quan hệ $fg\in\mathfrak{p}_{m+n}$ kéo theo $f\in\mathfrak{p}_m$ hoặc $g\in\mathfrak{p}_n$.
  \item[{\rm($3^\circ$)}] $\mathfrak{p}_n\neq S_n$ với ít nhất một $n\geq n_0$.
\end{enumerate}
Hơn nữa, khi đó iđêan nguyên tố phân bậc $\mathfrak{p}$ là duy nhất.
\end{proposition}

\begin{proof}
Hiển nhiên các điều kiện $1^\circ$ và $2^\circ$ là cần.
Hơn nữa, nếu $\mathfrak{p}\not\supset S_+$, thì tồn tại ít nhất một $k>0$ sao cho $\mathfrak{p}\cap S_k\neq S_k$; nếu $f\in S_k$ không thuộc $\mathfrak{p}$, quan hệ $\mathfrak{p}\cap S_n=S_n$ kéo theo $\mathfrak{p}\cap S_{n-mk}=S_{n-mk}$ theo \sref{II.2.1.8}.
Do đó, nếu $\mathfrak{p}\cap S_n=S_n$ kể từ một giá trị nào đó của $n$, thì ta sẽ có $\mathfrak{p}\supset S_+$, trái với giả thiết; điều này chứng minh $3^\circ$ là cần.
Ngược lại, giả sử các điều kiện $1^\circ$, $2^\circ$ và $3^\circ$ được thỏa mãn.
Nhận xét rằng, nếu với một số nguyên $d\geq n_0$, $f\in S_d$ không thuộc $\mathfrak{p}_d$, thì, nếu $\mathfrak{p}$ tồn tại, $\mathfrak{p}_m$, với $m<n_0$, nhất thiết bằng tập hợp các $x\in S_m$ sao cho $f^rx\in\mathfrak{p}_{m+rd}$ với mọi giá trị $r$, trừ một số hữu hạn.
Điều này đã chứng minh rằng, nếu $\mathfrak{p}$ tồn tại, thì nó là duy nhất.
Còn phải chứng minh rằng, nếu ta xác định các $\mathfrak{p}_m$ với $m<n_0$ bởi điều kiện trên, thì $\mathfrak{p}=\sum_{n=0}^{\infty}\mathfrak{p}_n$ là một iđêan nguyên tố.
Trước hết, nhận xét rằng, nhờ $2^\circ$, với $m\geq n_0$, $\mathfrak{p}_m$ cũng được xác định là tập hợp các $x\in S_m$ sao cho $f^rx\in\mathfrak{p}_{m+rd}$ với mọi giá trị $r$, trừ một số hữu hạn.
Như
\oldpage[II]{23}
vậy, nếu $g\in S_m$ và $x\in\mathfrak{p}_n$, thì $f^rgx\in\mathfrak{p}_{m+n+rd}$ với mọi giá trị $r$, trừ một số hữu hạn, do đó $gx\in\mathfrak{p}_{m+n}$, điều này chứng minh $\mathfrak{p}$ là một iđêan của $S$.
Để chứng minh iđêan này là nguyên tố, nói cách khác, vành $S/\mathfrak{p}$, được phân bậc bởi các nhóm con $S_n/\mathfrak{p}_n$, là một miền nguyên, chỉ cần (bằng cách xét các thành phần bậc cao nhất của hai phần tử thuộc $S/\mathfrak{p}$) chứng minh rằng, nếu $x\in S_m$ và $y\in S_n$ sao cho $x\notin\mathfrak{p}_m$ và $y\notin\mathfrak{p}_n$, thì $xy\notin\mathfrak{p}_{m+n}$.
Nếu không, với $r$ đủ lớn, ta sẽ có $f^{2r}xy\in\mathfrak{p}_{m+n+2rd}$; nhưng $f^ry\notin\mathfrak{p}_{n+rd}$ với mọi $r>0$.
Khi đó, từ $2^\circ$ suy ra rằng, với mọi giá trị $r$ trừ một số hữu hạn, ta có $f^rx\in\mathfrak{p}_{m+rd}$, và sẽ kết luận được $x\in\mathfrak{p}_m$, trái với giả thiết.
\end{proof}

\begin{env}[2.1.10]
\label{II.2.1.10}
Ta nói một tập con $\mathfrak{J}$ của $S_+$ là một \emph{iđêan của $S_+$} nếu nó là một iđêan của $S$, và $\mathfrak{J}$ là một \emph{iđêan nguyên tố phân bậc của $S_+$} nếu nó là giao của $S_+$ với một iđêan nguyên tố phân bậc của $S$ \emph{không chứa $S_+$} (iđêan nguyên tố này hơn nữa là duy nhất theo \sref{II.2.1.9}).
Nếu $\mathfrak{J}$ là một iđêan của $S_+$, thì \emph{căn của $\mathfrak{J}$ trong $S_+$} là tập hợp các phần tử của $S_+$ có một lũy thừa nào đó thuộc $\mathfrak{J}$, nói cách khác, là tập hợp
\[
  \rad_+(\mathfrak{J})=\rad(\mathfrak{J})\cap S_+.
\]
Đặc biệt, căn của $0$ trong $S_+$ còn được gọi là \emph{căn lũy linh} của $S_+$ và được ký hiệu bởi $\nilrad_+$: do đó, đây là tập hợp các phần tử lũy linh của $S_+$.
Nếu $\mathfrak{J}$ là một iđêan \emph{phân bậc} của $S_+$, thì căn $\rad_+(\mathfrak{J})$ của nó là một iđêan \emph{phân bậc}: bằng cách chuyển sang vành thương $S/\mathfrak{J}$, ta có thể quy về trường hợp $\mathfrak{J}=0$, và tất cả quy về việc chứng minh rằng, nếu $x=x_h+x_{h+1}+\cdots+x_k$ là lũy linh, thì các $x_i\in S_i$ ($1\leq h\leq i\leq k$) cũng lũy linh.
Có thể giả sử $x_k\neq0$, và khi đó thành phần bậc cao nhất của $x^n$ là $x_k^n$; vì vậy $x_k$ lũy linh, rồi ta lập luận quy nạp theo $k$.
Ta nói vành phân bậc $S$ là \emph{rút gọn cốt yếu} nếu $\nilrad_+=0$, nói cách khác, nếu $S_+$ không chứa phần tử lũy linh $\neq0$.
\end{env}

\begin{env}[2.1.11]
\label{II.2.1.11}
Nhận xét rằng, nếu trong vành phân bậc $S$, một phần tử $x$ là ước của không, thì thành phần bậc cao nhất của nó cũng vậy.
Ta nói một vành $S$ là \emph{nguyên cốt yếu} nếu vành $S_+$ (\emph{không có phần tử đơn vị}) không chứa ước của không và khác không; để có điều đó, chỉ cần không một phần tử thuần nhất khác không nào của $S_+$ là ước của không trong vành này.
Hiển nhiên, nếu $\mathfrak{p}$ là một iđêan nguyên tố phân bậc của $S_+$, thì $S/\mathfrak{p}$ là nguyên cốt yếu.

Cho $S$ là một vành phân bậc nguyên cốt yếu và $x_0\in S_0$; nếu khi đó tồn tại \emph{một} phần tử thuần nhất $f\neq0$ của $S_+$ sao cho $x_0f=0$, thì ta có $x_0S_+=0$, vì $(x_0g)f=(x_0f)g=0$ với mọi $g\in S_+$, và do đó giả thiết kéo theo $x_0g=0$.
Điều kiện cần và đủ để $S$ là một miền nguyên là $S_0$ là một miền nguyên và linh hóa tử của $S_+$ trong $S_0$ bằng $0$.
\end{env}

\subsection{Các vành phân thức của một vành phân bậc}
\label{subsection:II.2.2}

\begin{env}[2.2.1]
\label{II.2.2.1}
Cho $S$ là một vành phân bậc dương và $f$ là một phần tử \emph{thuần nhất} của $S$, có bậc $d>0$; khi đó vành phân thức $S'=S_f$ được phân bậc bằng cách lấy $S_n'$ là tập hợp các $x/f^k$, trong đó $x\in S_{n+kd}$ và $k\geq0$ (ở đây ta lưu ý rằng $n$ có thể nhận các giá trị âm tùy ý); ta ký hiệu bởi $S_{(f)}$ vành con $S_0'=(S_f)_0$ của $S'$ gồm các phần tử \emph{bậc $0$}.

Nếu $f\in S_d$, thì các đơn thức $(f/1)^h$ trong $S_f$ ($h$ là số nguyên dương hoặc âm) tạo thành một \emph{hệ tự do} trên vành $S_{(f)}$, còn tập hợp các tổ hợp tuyến tính của chúng chính là
\oldpage[II]{24}
vành $(S^{(d)})_f$, do đó đẳng cấu với $S_{(f)}[T,T^{-1}]=S_{(f)}\otimes_\bb{Z}\bb{Z}[T,T^{-1}]$ (trong đó $T$ là một biến bất định).
Thật vậy, nếu có một quan hệ
\[
  \sum_{h=-a}^{b}z_h(f/1)^h=0,
\]
với $z_h=x_h/f^m$, trong đó tất cả các $x_h$ đều thuộc $S_{md}$, thì theo định nghĩa, quan hệ này tương đương với sự tồn tại của một $k\geq a$ sao cho
\[
  \sum_{h=-a}^{b}f^{h+k}x_h=0;
\]
vì bậc của các số hạng trong tổng này đôi một khác nhau, ta có $f^{h+k}x_h=0$ với mọi $h$, từ đó $z_h=0$ với mọi $h$.

Nếu $M$ là một $S$-môđun phân bậc, thì $M'=M_f$ là một $S_f$-môđun phân bậc, trong đó $M_n'$ là tập hợp các $z/f^k$ với $z\in M_{n+kd}$ ($k\geq0$); ta ký hiệu bởi $M_{(f)}$ tập hợp các phần tử thuần nhất bậc $0$ của $M'$.
Hiển nhiên $M_{(f)}$ là một $S_{(f)}$-môđun và ta có
\[
  (M^{(d)})_f=M_{(f)}\otimes_{S_{(f)}}(S^{(d)})_f.
\]
\end{env}

\begin{lemma}[2.2.2]
\label{II.2.2.2}
Cho $d$ và $e$ là hai số nguyên $>0$, $f\in S_d$ và $g\in S_e$.
Tồn tại một đẳng cấu vành chính tắc
\[
  S_{(fg)}\isoto(S_{(f)})_{g^d/f^e};
\]
nếu ta đồng nhất chính tắc hai vành này, thì tồn tại một đẳng cấu môđun chính tắc
\[
  M_{(fg)}\isoto(M_{(f)})_{g^d/f^e}.
\]
\end{lemma}

\begin{proof}
Thật vậy, $fg$ chia $f^eg^d$, và phần tử sau chia $(fg)^{de}$; do đó các vành phân bậc $S_{fg}$ và $S_{f^eg^d}$ được đồng nhất chính tắc.
Mặt khác, $S_{f^eg^d}$ cũng được đồng nhất với $(S_{f^e})_{g^d/1}$ \sref[0]{0.1.4.6}; và vì $f^e/1$ khả nghịch trong $S_{f^e}$, $S_{f^eg^d}$ cũng được đồng nhất với $(S_{f^e})_{g^d/f^e}$.
Mà phần tử $g^d/f^e$ có bậc $0$ trong $S_{f^e}$; ta suy ra ngay rằng vành con của $(S_{f^e})_{g^d/f^e}$ gồm các phần tử bậc $0$ là $(S_{(f^e)})_{g^d/f^e}$, và vì hiển nhiên $S_{(f^e)}=S_{(f)}$, điều này chứng minh phần thứ nhất của bổ đề; phần thứ hai được chứng minh theo cùng một cách.
\end{proof}

\begin{env}[2.2.3]
\label{II.2.2.3}
Dưới các giả thiết của \sref{II.2.2.2}, hiển nhiên đồng cấu chính tắc $S_f\to S_{fg}$ \sref[0]{0.1.4.1}, biến $x/f^k$ thành $g^kx/(fg)^k$, có bậc $0$ và do đó, bằng phép hạn chế, cho một \emph{đồng cấu chính tắc $S_{(f)}\to S_{(fg)}$}, sao cho sơ đồ
\[
  \xymatrix{
    & S_{(f)}\ar[dl]\ar[dr]\\
    S_{(fg)}\ar[rr]^-{\sim} & &
    (S_{(f)})_{g^d/f^e}
  }
\]
giao hoán.
Theo cùng một cách, ta xác định một đồng cấu chính tắc $M_{(f)}\to M_{(fg)}$.
\end{env}

\begin{lemma}[2.2.4]
\label{II.2.2.4}
Nếu $f$ và $g$ là hai phần tử thuần nhất của $S_+$, thì vành $S_{(fg)}$ được sinh bởi hợp của các ảnh chính tắc của $S_{(f)}$ và $S_{(g)}$.
\end{lemma}

\begin{proof}
Theo \sref{II.2.2.2}, chỉ cần thấy rằng
\[
  1/(g^d/f^e)=f^{d+e}/(fg)^d
\]
thuộc ảnh chính tắc của $S_{(g)}$ trong $S_{(fg)}$, điều này hiển nhiên theo định nghĩa.
\end{proof}

\begin{proposition}[2.2.5]
\label{II.2.2.5}
Cho $d$ là một số nguyên $>0$ và $f\in S_d$.
Khi đó tồn tại một đẳng cấu vành chính tắc
\[
  S_{(f)}\isoto S^{(d)}/(f-1)S^{(d)};
\]
nếu ta đồng nhất hai vành này bằng đẳng cấu đó, thì tồn tại một đẳng cấu môđun chính tắc
\[
  M_{(f)}\isoto M^{(d)}/(f-1)M^{(d)}.
\]
\end{proposition}

\begin{proof}
Đẳng cấu thứ nhất trong các đẳng cấu này được xác định bằng cách cho $x/f^n$, trong đó $x\in S_{nd}$, tương ứng với phần tử $\overline{x}$, là lớp của $x$ modulo $(f-1)S^{(d)}$; ánh xạ này được xác định tốt vì ta có đồng dư
\[
  f^hx\equiv x\pmod{(f-1)S^{(d)}}
\]
với mọi $x\in S^{(d)}$; do đó, nếu $f^hx=0$ với một $h>0$ nào đó,
\oldpage[II]{25}
thì ta có $\overline{x}=0$.
Mặt khác, nếu $x\in S_{nd}$ thỏa $x=(f-1)y$, với
$y=y_{hd}+y_{(h+1)d}+\cdots+y_{kd}$, trong đó $y_{jd}\in S_{jd}$ và $y_{hd}\neq0$, thì nhất thiết $h=n$ và $x=-y_{hd}$, đồng thời có các hệ thức $y_{(j+1)d}=fy_{jd}$ với $h\leq j\leq k-1$ và $fy_{kd}=0$, cuối cùng suy ra $f^{k-n+1}x=0$.
Do đó, với mỗi lớp $\overline{x}\text{ mod. }(f-1)S^{(d)}$ của một phần tử $x\in S_{nd}$, ta có thể cho tương ứng phần tử $x/f^n$ của $S_{(f)}$, vì nhận xét trên cho thấy ánh xạ này được xác định tốt.
Hiển nhiên hai ánh xạ được định nghĩa như vậy là các đồng cấu vành nghịch đảo của nhau.
Đối với $M$, ta tiến hành hoàn toàn tương tự.
\end{proof}

\begin{corollary}[2.2.6]
\label{II.2.2.6}
Nếu $S$ là Noether, thì $S_{(f)}$ cũng là Noether đối với mọi $f$ thuần nhất có bậc $>0$.
\end{corollary}

\begin{proof}
Điều này suy ra ngay từ \sref{II.2.1.7} và \sref{II.2.2.5}.
\end{proof}

\begin{env}[2.2.7]
\label{II.2.2.7}
Cho $T$ là một tập con nhân của $S_+$ gồm các phần tử \emph{thuần nhất}; khi đó $T_0=T\cup\{1\}$ là một tập con nhân của $S$.
Vì các phần tử của $T_0$ là thuần nhất, vành $T_0^{-1}S$ vẫn được phân bậc theo cách hiển nhiên; ta ký hiệu bởi $S_{(T)}$ vành con của $T_0^{-1}S$ gồm các phần tử bậc $0$, tức là các phần tử dạng $x/h$, trong đó $h\in T$ và $x$ thuần nhất có cùng bậc với $h$.
Theo \sref[0]{0.1.4.5}, $T_0^{-1}S$ được đồng nhất chính tắc với giới hạn quy nạp của các vành $S_f$, khi $f$ chạy qua $T$ (đối với các đồng cấu chính tắc $S_f\to S_{fg}$); vì sự đồng nhất này bảo toàn bậc, nó đồng nhất $S_{(T)}$ với \emph{giới hạn quy nạp} của các $S_{(f)}$ với $f\in T$.
Đối với mọi $S$-môđun phân bậc $M$, ta định nghĩa tương tự môđun $M_{(T)}$ (trên vành $S_{(T)}$), gồm các phần tử bậc $0$ của $T_0^{-1}M$, và thấy rằng môđun này là giới hạn quy nạp của các $M_{(f)}$ với $f\in T$.

Nếu $\mathfrak{p}$ là một iđêan nguyên tố thuần nhất của $S_+$, ta lần lượt ký hiệu bởi $S_{(\mathfrak{p})}$ và $M_{(\mathfrak{p})}$ vành $S_{(T)}$ và môđun $M_{(T)}$, trong đó $T$ là tập các phần tử \emph{thuần nhất} của $S_+$ không thuộc $\mathfrak{p}$.
\end{env}

\subsection{Phổ nguyên tố thuần nhất của một vành phân bậc}
\label{subsection:II.2.3}

\begin{env}[2.3.1]
\label{II.2.3.1}
Cho một vành $S$ phân bậc dương, ta gọi \emph{phổ nguyên tố thuần nhất} của $S$, và ký hiệu là $\Proj(S)$, tập các iđêan nguyên tố thuần nhất của $S_+$ \sref{II.2.1.10}, hay tương đương, tập các iđêan nguyên tố thuần nhất của $S$ \emph{không chứa $S_+$}; ta sẽ định nghĩa một cấu trúc \emph{lược đồ} có tập nền là $\Proj(S)$.
\end{env}

\begin{env}[2.3.2]
\label{II.2.3.2}
Đối với mỗi tập con $E$ của $S$, gọi $V_+(E)$ là tập các iđêan nguyên tố thuần nhất của $S$ chứa $E$ và không chứa $S_+$; do đó, đây là tập con $V(E)\cap\Proj(S)$ của $\Spec(S)$.
Từ \sref[I]{I.1.1.2}, ta suy ra:
\[
\label{II.2.3.2.1}
  V_+(0)=\Proj(S),\qquad V_+(S)=V_+(S_+)=\emp,
\tag{2.3.2.1}
\]
\[
\label{II.2.3.2.2}
  V_+\big(\textstyle\bigcup_\lambda E_\lambda\big)
  =\textstyle\bigcap_\lambda V_+(E_\lambda),
\tag{2.3.2.2}
\]
\[
\label{II.2.3.2.3}
  V_+(EE')=V_+(E)\cup V_+(E').
\tag{2.3.2.3}
\]

Thay $E$ bằng iđêan phân bậc sinh bởi $E$ không làm thay đổi $V_+(E)$; hơn nữa, nếu $\mathfrak{J}$ là một iđêan phân bậc của $S$, thì
\[
\label{II.2.3.2.4}
  V_+(\mathfrak{J})
  =V_+\big(\textstyle\bigcup_{q\geq n}(\mathfrak{J}\cap S_q)\big)
\tag{2.3.2.4}
\]
\oldpage[II]{26}
đối với mọi $n>0$.
Thật vậy, nếu $\mathfrak{p}\in\Proj(S)$ chứa các phần tử thuần nhất bậc $\geq n$ của $\mathfrak{J}$, thì theo giả thiết tồn tại một phần tử thuần nhất $f\in S_d$ không thuộc $\mathfrak{p}$; do đó, với mọi $m\geq0$ và mọi $x\in S_m\cap\mathfrak{J}$, ta có $f^rx\in\mathfrak{J}\cap S_{m+rd}$ với mọi $r\geq0$; trừ một số hữu hạn giá trị của $r$, ta có $m+rd\geq n$, nên $f^rx\in\mathfrak{p}\cap S_{m+rd}$, kéo theo $x\in\mathfrak{p}\cap S_m$ \sref{II.2.1.8}.

Cuối cùng, đối với mọi iđêan phân bậc $\mathfrak{J}$ của $S$, ta có
\[
  V_+(\mathfrak{J})=V_+(\rad_+(\mathfrak{J})).
\tag{2.3.2.5}
\]
\end{env}

\begin{env}[2.3.3]
\label{II.2.3.3}
Theo định nghĩa, các $V_+(E)$ là những tập con đóng của $X=\Proj(S)$ đối với tô pô cảm sinh bởi tô pô phổ của $\Spec(S)$, mà ta cũng gọi là \emph{tô pô phổ} trên $X$.
Với mọi $f\in S$, đặt
\[
\label{II.2.3.3.1}
  D_+(f)=D(f)\cap\Proj(S)=\Proj(S)\setmin V_+(f)
\tag{2.3.3.1}
\]
và do đó, đối với hai phần tử $f$, $g$ bất kỳ của $S$ \sref[I]{I.1.1.9.1},
\[
\label{II.2.3.3.2}
  D_+(fg)=D_+(f)\cap D_+(g).
\tag{2.3.3.2}
\]
\end{env}

\begin{proposition}[2.3.4]
\label{II.2.3.4}
Khi $f$ chạy qua tập các phần tử thuần nhất của $S_+$, các $D_+(f)$ lập thành một cơ sở của tô pô trên $X=\Proj(S)$.
\end{proposition}

\begin{proof}
Thật vậy, từ \sref{II.2.3.2.2} và \sref{II.2.3.2.4} suy ra rằng mọi tập con đóng của $X$ là giao của các tập dạng $V_+(f)$, trong đó $f$ thuần nhất có bậc $>0$.
\end{proof}

\begin{env}[2.3.5]
\label{II.2.3.5}
Cho $f$ là một phần tử \emph{thuần nhất} của $S_+$, có bậc $d>0$.
Đối với mỗi iđêan nguyên tố thuần nhất $\mathfrak{p}$ của $S$ không chứa $f$, ta biết rằng tập các $x/f^n$, trong đó $x\in\mathfrak{p}$ và $n\geq0$, là một iđêan nguyên tố của vành phân thức $S_f$ \sref[0]{0.1.2.6}; vì vậy, giao của nó với $S_{(f)}$ là một iđêan nguyên tố của vành này, được ký hiệu là $\psi_f(\mathfrak{p})$: đó là tập các $x/f^n$ với $n\geq0$ và $x\in\mathfrak{p}\cap S_{nd}$.
Như vậy ta đã định nghĩa một ánh xạ
\[
  \psi_f:D_+(f)\to\Spec(S_{(f)});
\]
hơn nữa, nếu $g\in S_e$ là một phần tử thuần nhất thứ hai của $S_+$, ta có một biểu đồ giao hoán
\[
\label{II.2.3.5.1}
  \xymatrix{
    D_+(f) \ar[r]^{\psi_f}
    & \Spec(S_{(f)})
  \\D_+(fg) \ar[u] \ar[r]_{\psi_{fg}}
    & \Spec(S_{(fg)}) \ar[u]
  }
\tag{2.3.5.1}
\]
trong đó mũi tên dọc bên trái là phép nhúng, còn mũi tên dọc bên phải là ánh xạ ${}^a\!\omega_{fg,f}$ cảm sinh bởi đồng cấu chính tắc $\omega=\omega_{fg,f}:S_{(f)}\to S_{(fg)}$ \sref[I]{I.1.2.1}.
Thật vậy, nếu $x/f^n\in\omega^{-1}(\psi_{fg}(\mathfrak{p}))$, với $fg\notin\mathfrak{p}$, thì theo định nghĩa $g^nx/(fg)^n\in\psi_{fg}(\mathfrak{p})$, nên $g^nx\in\mathfrak{p}$ và do đó $x\in\mathfrak{p}$; chiều đảo lại là hiển nhiên.
\end{env}

\begin{proposition}[2.3.6]
\label{II.2.3.6}
Ánh xạ $\psi_f$ là một phép đồng phôi từ $D_+(f)$ lên $\Spec(S_{(f)})$.
\end{proposition}

\begin{proof}
Trước hết, $\psi_f$ liên tục; thật vậy, nếu $h\in S_{nd}$, thì $h/f^n\in\psi_f(\mathfrak{p})$ khi và chỉ khi $h\in\mathfrak{p}$, do đó $\psi_f^{-1}(D(h/f^n))=D_+(hf)$; khẳng định của ta suy ra từ \sref{II.2.3.3.2}.
Hơn nữa, các $D_+(hf)$, khi $h$ chạy qua các tập $S_{nd}$, lập thành một cơ sở của tô pô trên $D_+(f)$, theo \sref{II.2.3.4} và \sref{II.2.3.3.2}.
Điều
\oldpage[II]{27}
vừa trình bày, khi xét đến tiên đề $(T_0)$ đúng trong $D_+(f)$ và trong $\Spec(S_{(f)})$, chứng minh rằng $\psi_f$ là đơn ánh và ánh xạ ngược $\psi_f(D_+(f))\to D_+(f)$ là liên tục.
Cuối cùng, để thấy $\psi_f$ là toàn ánh, nhận xét rằng nếu $\mathfrak{q}_0$ là một iđêan nguyên tố của $S_{(f)}$ và, với mọi $n>0$, ta ký hiệu bởi $\mathfrak{p}_n$ tập các $x\in S_n$ sao cho $x^d/f^n\in\mathfrak{q}_0$, thì các $\mathfrak{p}_n$ thỏa các điều kiện của \sref{II.2.1.9}.
Thật vậy, nếu $x,y\in S_n$ thỏa $x^d/f^n,y^d/f^n\in\mathfrak{q}_0$, thì $(x+y)^{2d}/f^{2n}\in\mathfrak{q}_0$, từ đó $(x+y)^d/f^n\in\mathfrak{q}_0$ vì $\mathfrak{q}_0$ là nguyên tố.
Điều này chứng minh rằng các $\mathfrak{p}_n$ là các nhóm con của $S_n$, và việc kiểm tra các điều kiện còn lại của \sref{II.2.1.9} là trực tiếp, khi xét đến tính nguyên tố của $\mathfrak{q}_0$.
Nếu $\mathfrak{p}$ là iđêan nguyên tố thuần nhất của $S$ được định nghĩa như vậy, thì quả thật $\psi_f(\mathfrak{p})=\mathfrak{q}_0$, vì với $x\in S_{nd}$, các hệ thức $x/f^n\in\mathfrak{q}_0$ và $x^d/f^{nd}\in\mathfrak{q}_0$ là tương đương do $\mathfrak{q}_0$ là nguyên tố.
\end{proof}

\begin{corollary}[2.3.7]
\label{II.2.3.7}
Để $D_+(f)=\emp$, điều kiện cần và đủ là $f$ lũy linh.
\end{corollary}

\begin{proof}
Thật vậy, để $\Spec(S_{(f)})=\emp$, điều kiện cần và đủ là $S_{(f)}=0$, hay cũng vậy, $1=0$ trong $S_f$, mà theo định nghĩa điều này có nghĩa là $f$ lũy linh.
\end{proof}

\begin{corollary}[2.3.8]
\label{II.2.3.8}
Cho $E$ là một tập con của $S_+$.
Các điều kiện sau là tương đương:
\begin{enumerate}
  \item[\rm(a)] $V_+(E)=X=\Proj(S)$.
  \item[\rm(b)] Mọi phần tử của $E$ đều lũy linh.
  \item[\rm(c)] Các thành phần thuần nhất của mọi phần tử thuộc $E$ đều lũy linh.
\end{enumerate}
\end{corollary}

\begin{proof}
Hiển nhiên (c) kéo theo (b), và (b) kéo theo (a).
Nếu $\mathfrak{J}$ là iđêan phân bậc của $S$ sinh bởi $E$, thì điều kiện (a) tương đương với $V_+(\mathfrak{J})=X$; \emph{a fortiori}, (a) kéo theo rằng mọi phần tử thuần nhất $f\in\mathfrak{J}$ thỏa $V_+(f)=X$, nên $f$ lũy linh theo \sref{II.2.3.7}.
\end{proof}

\begin{corollary}[2.3.9]
\label{II.2.3.9}
Nếu $\mathfrak{J}$ là một iđêan phân bậc của $S_+$, thì $\rad_+(\mathfrak{J})$ là giao của các iđêan nguyên tố thuần nhất của $S_+$ chứa $\mathfrak{J}$.
\end{corollary}

\begin{proof}
Xét vành phân bậc $S/\mathfrak{J}$, ta có thể quy về trường hợp $\mathfrak{J}=0$.
Cần chứng minh rằng nếu $f\in S_+$ không lũy linh thì tồn tại một iđêan nguyên tố thuần nhất của $S$ không chứa $f$; nhưng ít nhất một thành phần thuần nhất của $f$ không lũy linh, nên ta có thể giả sử $f$ thuần nhất.
Khi đó khẳng định suy ra từ \sref{II.2.3.7}.
\end{proof}

\begin{env}[2.3.10]
\label{II.2.3.10}
Đối với mỗi tập con $Y$ của $X=\Proj(S)$, gọi $\mathfrak{j}_+(Y)$ là tập các $f\in S_+$ sao cho $Y\subset V_+(f)$; cũng có thể nói $\mathfrak{j}_+(Y)=\mathfrak{j}(Y)\cap S_+$.
Do đó, $\mathfrak{j}_+(Y)$ là một iđêan của $S_+$ bằng căn của nó trong $S_+$.
\end{env}

\begin{proposition}[2.3.11]
\label{II.2.3.11}
\begin{enumerate}
  \item[\rm(i)] Đối với mọi tập con $E$ của $S_+$, $\mathfrak{j}_+(V_+(E))$ là căn trong $S_+$ của iđêan phân bậc của $S_+$ sinh bởi $E$.
  \item[\rm(ii)] Đối với mọi tập con $Y$ của $X$, $V_+(\mathfrak{j}_+(Y))=\overline{Y}$, trong đó $\overline{Y}$ là bao đóng của $Y$ trong $X$.
\end{enumerate}
\end{proposition}

\begin{proof}
\medskip\noindent
\begin{enumerate}
  \item[\rm(i)] Nếu $\mathfrak{J}$ là iđêan phân bậc của $S_+$ sinh bởi $E$, thì $V_+(E)=V_+(\mathfrak{J})$, và khẳng định suy ra từ \sref{II.2.3.9}.
  \item[\rm(ii)] Vì $V_+(\mathfrak{J})=\bigcap_{f\in\mathfrak{J}}V_+(f)$, hệ thức $Y\subset V_+(\mathfrak{J})$ kéo theo $Y\subset V_+(f)$ đối với mọi $f\in\mathfrak{J}$, và do đó $\mathfrak{j}_+(Y)\supset\mathfrak{J}$, suy ra $V_+(\mathfrak{j}_+(Y))\subset V_+(\mathfrak{J})$; điều này chứng minh (ii) theo định nghĩa của các tập con đóng.
\end{enumerate}
\end{proof}

\begin{corollary}[2.3.12]
\label{II.2.3.12}
Các tập con đóng $Y$ của $X=\Proj(S)$ và các iđêan phân bậc của $S_+$ bằng căn của chính chúng trong $S_+$ tương ứng song ánh qua các ánh xạ đảo chiều bao hàm $Y\mapsto\mathfrak{j}_+(Y)$ và $\mathfrak{J}\mapsto V_+(\mathfrak{J})$; hợp $Y_1\cup Y_2$ của hai tập con đóng của $X$ thì ứng với
\oldpage[II]{28}
$\mathfrak{j}_+(Y_1)\cap\mathfrak{j}_+(Y_2)$, còn giao của một họ tùy ý $(Y_\lambda)$ các tập con đóng thì ứng với căn trong $S_+$ của tổng các $\mathfrak{j}_+(Y_\lambda)$.
\end{corollary}

\begin{corollary}[2.3.13]
\label{II.2.3.13}
Cho $\mathfrak{J}$ là một iđêan phân bậc của $S_+$.
Để $V_+(\mathfrak{J})=\emp$, điều kiện cần và đủ là mỗi phần tử của $S_+$ có một lũy thừa nằm trong $\mathfrak{J}$.
\end{corollary}

Hệ quả vừa rồi cũng có thể được phát biểu dưới một trong các dạng tương đương sau:

\begin{corollary}[2.3.14]
\label{II.2.3.14}
Cho $(f_\alpha)$ là một họ các phần tử thuần nhất của $S_+$.
Để các $D_+(f_\alpha)$ lập thành một phủ của $X=\Proj(S)$, điều kiện cần và đủ là mỗi phần tử của $S_+$ có một lũy thừa nằm trong iđêan sinh bởi các $f_\alpha$.
\end{corollary}

\begin{corollary}[2.3.15]
\label{II.2.3.15}
Cho $(f_\alpha)$ là một họ các phần tử thuần nhất của $S_+$ và $f$ là một phần tử của $S_+$.
Các hệ thức sau là tương đương:
\begin{enumerate}
  \item[\rm(a)] $D_+(f)\subset\bigcup_\alpha D_+(f_\alpha)$;
  \item[\rm(b)] $V_+(f)\supset\bigcap_\alpha V_+(f_\alpha)$;
  \item[\rm(c)] một lũy thừa của $f$ thuộc iđêan sinh bởi các $f_\alpha$.
\end{enumerate}
\end{corollary}

\begin{corollary}[2.3.16]
\label{II.2.3.16}
Để $X=\Proj(S)$ là rỗng, điều kiện cần và đủ là mọi phần tử của $S_+$ đều lũy linh.
\end{corollary}

\begin{corollary}[2.3.17]
\label{II.2.3.17}
Trong tương ứng song ánh mô tả ở \sref{II.2.3.12}, các tập con đóng \emph{bất khả quy} của $X$ tương ứng với các iđêan phân bậc \emph{nguyên tố} của $S_+$.
\end{corollary}

\begin{proof}
Thật vậy, nếu $Y=Y_1\cup Y_2$, trong đó $Y_1$ và $Y_2$ đóng và khác $Y$, thì
\[
  \mathfrak{j}_+(Y)=\mathfrak{j}_+(Y_1)\cap\mathfrak{j}_+(Y_2)
\]
và các iđêan $\mathfrak{j}_+(Y_1)$, $\mathfrak{j}_+(Y_2)$ đều khác $\mathfrak{j}_+(Y)$; do đó $\mathfrak{j}_+(Y)$ không nguyên tố.
Ngược lại, nếu $\mathfrak{J}$ là một iđêan phân bậc không nguyên tố của $S_+$, thì tồn tại hai phần tử $f$, $g$ của $S_+$ sao cho $f\notin\mathfrak{J}$, $g\notin\mathfrak{J}$ và $fg\in\mathfrak{J}$.
Khi đó $V_+(f)\not\supset V_+(\mathfrak{J})$, $V_+(g)\not\supset V_+(\mathfrak{J})$, nhưng $V_+(\mathfrak{J})\subset V_+(f)\cup V_+(g)$ theo \sref{II.2.3.2.3}; suy ra $V_+(\mathfrak{J})$ là hợp của các tập con đóng $V_+(f)\cap V_+(\mathfrak{J})$ và $V_+(g)\cap V_+(\mathfrak{J})$, cả hai đều khác $V_+(\mathfrak{J})$.
\end{proof}

\subsection{Cấu trúc lược đồ trên $\Proj(S)$}
\label{subsection:II.2.4}

\begin{env}[2.4.1]
\label{II.2.4.1}
Cho $f$, $g$ là hai phần tử thuần nhất của $S_+$; xét các lược đồ affine $Y_f=\Spec(S_{(f)})$, $Y_g=\Spec(S_{(g)})$ và $Y_{fg}=\Spec(S_{(fg)})$.
Theo \sref{II.2.2.2}, cấu xạ $w_{fg,f}=({}^a\!\omega_{fg,f},\widetilde{\omega}_{fg,f})$ từ $Y_{fg}$ vào $Y_f$, tương ứng với đồng cấu chính tắc $\omega_{fg,f}:S_{(f)}\to S_{(fg)}$, là một \emph{phép nhúng mở} \sref[I]{I.1.3.6}.
Nhờ phép đồng phôi ngược của $\psi_f:D_+(f)\to Y_f$ \sref{II.2.3.6}, ta có thể chuyển cấu trúc lược đồ affine của $Y_f$ sang $D_+(f)$; do tính giao hoán của biểu đồ \sref{II.2.3.5.1}, lược đồ affine $D_+(fg)$ khi đó được đồng nhất với lược đồ cảm sinh trên tập con mở $D_+(fg)$ của không gian nền của lược đồ affine $D_+(f)$.
Khi xét đến \sref{II.2.3.4}, hiển nhiên $X=\Proj(S)$ được trang bị một cấu trúc \emph{tiền lược đồ} duy nhất mà hạn chế lên mỗi $D_+(f)$ là lược đồ affine vừa được định nghĩa.
Hơn nữa:
\end{env}

\begin{proposition}[2.4.2]
\label{II.2.4.2}
Tiền lược đồ $\Proj(S)$ là một lược đồ.
\end{proposition}

\begin{proof}
Theo \sref[I]{I.5.5.6}, chỉ cần chứng minh rằng, với mọi $f$, $g$ thuần nhất trong $S_+$, $D_+(f)\cap D_+(g)=D_+(fg)$ là affine và vành của nó được sinh bởi các ảnh chính tắc của các vành của $D_+(f)$ và $D_+(g)$; điểm thứ nhất hiển nhiên theo định nghĩa, còn điểm thứ hai suy ra từ \sref{II.2.2.4}.
\end{proof}
\oldpage[II]{29}
Khi nói đến phổ nguyên tố thuần nhất $\Proj(S)$ như một \emph{lược đồ}, ta luôn hiểu đó là cấu trúc vừa được định nghĩa.

\begin{example}[2.4.3]
\label{II.2.4.3}
Lấy $S=K[T_1,T_2]$, trong đó $K$ là một trường, $T_1$, $T_2$ là hai biến bất định, và $S$ được phân bậc theo bậc toàn phần.
Từ \sref{II.2.3.14} suy ra $\Proj(S)$ là hợp của $D_+(T_1)$ và $D_+(T_2)$; ta thấy ngay rằng các lược đồ affine này đẳng cấu chính tắc với $\Spec(K[T])$, và $\Proj(S)$ thu được bằng cách dán hai lược đồ affine ấy như mô tả trong \sref[I]{I.2.3.2} (xem \sref{II.7.4.14}).
\end{example}

\begin{proposition}[2.4.4]
\label{II.2.4.4}
Cho $S$ là một vành phân bậc dương, và cho $X$ là lược đồ $\Proj(S)$.
\begin{enumerate}
  \item[\rm(i)] Nếu $\mathfrak{N}_+$ là căn lũy linh của $S_+$ \sref{II.2.1.10}, thì lược đồ $X_\red$ đẳng cấu chính tắc với $\Proj(S/\mathfrak{N}_+)$; đặc biệt, nếu $S$ cốt yếu rút gọn, thì $\Proj(S)$ rút gọn.
  \item[\rm(ii)] Giả sử $S$ cốt yếu rút gọn. Khi đó, để $X$ là nguyên, điều kiện cần và đủ là $S$ cốt yếu nguyên.
\end{enumerate}
\end{proposition}

\begin{proof}
\medskip\noindent
\begin{enumerate}
  \item[\rm(i)] Cho $\overline{S}$ là vành phân bậc $S/\mathfrak{N}_+$, và ký hiệu bởi $x\mapsto\overline{x}$ đồng cấu chính tắc $S\to\overline{S}$ có bậc $0$.
  Với mọi $f\in S_d$ ($d>0$), đồng cấu chính tắc $S_f\to\overline{S}_{\overline{f}}$ \sref[0]{0.1.5.1} là toàn ánh và có bậc $0$, nên khi hạn chế cho một đồng cấu toàn ánh $S_{(f)}\to\overline{S}_{(\overline{f})}$.
  Nếu giả sử $f\notin\mathfrak{N}_+$, ta kiểm tra ngay rằng $\overline{S}_{(\overline{f})}$ là rút gọn và hạt nhân của đồng cấu trên là căn lũy linh của $S_{(f)}$, nói cách khác,
  \[
    \overline{S}_{(\overline{f})}=(S_{(f)})_\red.
  \]
  Do đó, đồng cấu này tương ứng với một phép nhúng đóng $D_+(\overline{f})\to D_+(f)$ đồng nhất $D_+(\overline{f})$ với $(D_+(f))_\red$ \sref[I]{I.5.1.2}, và đặc biệt là một phép đồng phôi giữa các không gian nền của hai lược đồ affine này.
  Hơn nữa, nếu $g\notin\mathfrak{N}_+$ là một phần tử thuần nhất thứ hai của $S_+$, thì biểu đồ
  \[
    \xymatrix{
      S_{(f)} \ar[r] \ar[d]
      & \overline{S}_{(\overline{f})} \ar[d]
    \\S_{(fg)} \ar[r]
      & \overline{S}_{(\overline{fg})}
    }
  \]
  là giao hoán.
  Vì các $D_+(f)$, với $f$ thuần nhất có bậc $>0$ và $f\notin\mathfrak{N}_+$, lập thành một phủ của $X=\Proj(S)$ \sref{II.2.3.7}, ta thấy các cấu xạ $D_+(\overline{f})\to D_+(f)$ là những hạn chế của một phép nhúng đóng $\Proj(\overline{S})\to\Proj(S)$, đồng thời là một phép đồng phôi giữa các không gian nền; kết luận suy ra từ \sref[I]{I.5.1.2}.
  \item[\rm(ii)] Giả sử $S$ cốt yếu nguyên, nói cách khác $(0)$ là một iđêan nguyên tố thuần nhất của $S_+$ khác $S_+$; khi đó $X$ rút gọn theo (i) và bất khả quy theo \sref{II.2.3.17}.
  Ngược lại, giả sử $S$ cốt yếu rút gọn và $X$ nguyên; khi đó, với $f\neq0$ thuần nhất trong $S_+$, ta có $D_+(f)\neq\emp$ \sref{II.2.3.7}.
  Giả thiết $X$ bất khả quy kéo theo $D_+(f)\cap D_+(g)\neq\emp$ đối với $f$, $g$ thuần nhất và $\neq0$ trong $S_+$; do đó $fg\neq0$ theo \sref{II.2.3.3.2}, và suy ra $S_+$ không có ước của không, từ đó có khẳng định đầu tiên.
\end{enumerate}
\end{proof}

\begin{env}[2.4.5]
\label{II.2.4.5}
Cho một vành giao hoán $A$, nhắc lại rằng một vành phân bậc $S$ được gọi là một \emph{$A$-đại số phân bậc} nếu nó được trang bị một cấu trúc $A$-đại số sao cho mỗi nhóm con $S_n$ là một $A$-môđun; hơn nữa, để có điều đó, chỉ cần $S_0$ là
\oldpage[II]{30}
một $A$-đại số; nói cách khác, ta xác định trên $S$ một cấu trúc $A$-đại số phân bậc bằng cách xác định trên $S_0$ một cấu trúc $A$-đại số và đặt $\alpha\cdot x=(\alpha\cdot1)x$ với $\alpha\in A$ và $x\in S_n$.
\end{env}

\begin{proposition}[2.4.6]
\label{II.2.4.6}
Giả sử $S$ là một $A$-đại số phân bậc.
Khi đó, trên $X=\Proj(S)$, bó cấu trúc $\sh{O}_X$ là một bó các $A$-đại số ($A$ được xem như bó hằng trên $X$); nói cách khác, $X$ là một lược đồ trên $\Spec(A)$.
\end{proposition}

\begin{proof}
Chỉ cần lưu ý rằng, với mọi $f$ thuần nhất trong $S_+$, $S_{(f)}$ là một đại số trên $A$, và sơ đồ
\[
  \xymatrix{
    S_{(f)} \ar[rr] && S_{(fg)}
  \\ & A \ar[ul] \ar[ur] &
  }
\]
giao hoán với $f,g$ thuần nhất trong $S_+$.
\end{proof}

\begin{proposition}[2.4.7]
\label{II.2.4.7}
Cho $S$ là một vành phân bậc dương.
\begin{enumerate}
  \item[\rm(i)] Với mọi số nguyên $d>0$, tồn tại một đẳng cấu chính tắc từ lược đồ $\Proj(S)$ lên lược đồ $\Proj(S^{(d)})$.
  \item[\rm(ii)] Cho $S'$ là vành phân bậc sao cho $S'_0=\bb{Z}$ và $S'_n=S_n$ (được xem như một $\bb{Z}$-môđun) với $n>0$. Tồn tại một đẳng cấu chính tắc từ lược đồ $\Proj(S)$ lên lược đồ $\Proj(S')$.
\end{enumerate}
\end{proposition}

\begin{proof}
\medskip\noindent
\begin{enumerate}
  \item[\rm(i)] Trước hết, ta chứng minh ánh xạ $\mathfrak{p}\mapsto\mathfrak{p}\cap S^{(d)}$ là một song ánh từ tập $\Proj(S)$ lên tập $\Proj(S^{(d)})$.
  Thật vậy, cho một iđêan nguyên tố phân bậc $\mathfrak{p}'\in\Proj(S^{(d)})$, và đặt $\mathfrak{p}_{nd}=\mathfrak{p}'\cap S_{nd}$ ($n\geq0$).
  Với mọi $n>0$ không chia hết cho $d$, ta định nghĩa $\mathfrak{p}_n$ là tập các $x\in S_n$ sao cho $x^d\in\mathfrak{p}_{nd}$.
  Nếu $x,y\in\mathfrak{p}_n$, thì $(x+y)^{2d}\in\mathfrak{p}_{2nd}$, do đó $(x+y)^d\in\mathfrak{p}_{nd}$ vì $\mathfrak{p}'$ là nguyên tố.
  Hiển nhiên các $\mathfrak{p}_n$ được định nghĩa như vậy với mọi $n\geq0$ thỏa các điều kiện của \sref{II.2.1.9}; do đó tồn tại một iđêan nguyên tố duy nhất $\mathfrak{p}\in\Proj(S)$ sao cho $\mathfrak{p}\cap S^{(d)}=\mathfrak{p}'$.
  Vì với mọi $f$ thuần nhất trong $S_+$, ta có $V_+(f)=V_+(f^d)$ \sref{II.2.3.2.3}, nên song ánh trên là một \emph{đồng phôi} giữa các không gian tôpô.
  Cuối cùng, với cùng các ký hiệu, $S_{(f)}$ và $S^{(d)}_{(f^d)}$ được đồng nhất chính tắc \sref{II.2.2.2}; do đó $\Proj(S)$ và $\Proj(S^{(d)})$ được đồng nhất chính tắc như các \emph{lược đồ}.
  \item[\rm(ii)] Nếu với mỗi $\mathfrak{p}\in\Proj(S)$ ta cho tương ứng iđêan nguyên tố duy nhất $\mathfrak{p}'\in\Proj(S')$ sao cho $\mathfrak{p}'\cap S_n=\mathfrak{p}\cap S_n$ với mọi $n>0$ \sref{II.2.1.9}, thì rõ ràng ta xác định được một đồng phôi chính tắc $\Proj(S)\simto\Proj(S')$ giữa các không gian nền, vì $V_+(f)$ là cùng một tập đối với $S$ và $S'$ khi $f$ là một phần tử thuần nhất của $S_+$.
  Hơn nữa, vì $S_{(f)}=S'_{(f)}$, nên $\Proj(S)$ và $\Proj(S')$ được đồng nhất như các \emph{lược đồ}.
\end{enumerate}
\end{proof}

\begin{corollary}[2.4.8]
\label{II.2.4.8}
Nếu $S$ là một $A$-đại số phân bậc và $S'_A$ là $A$-đại số phân bậc sao cho $(S'_A)_0=A$ và $(S'_A)_n=S_n$ với $n>0$, thì tồn tại một đẳng cấu chính tắc từ $\Proj(S)$ lên $\Proj(S'_A)$.
\end{corollary}

\begin{proof}
Thật vậy, với các ký hiệu của \sref{II.2.4.7}[(ii)], hai lược đồ này đều đẳng cấu chính tắc với $\Proj(S')$.
\end{proof}

\subsection{Bó liên kết với một môđun phân bậc}
\label{subsection:II.2.5}

\begin{env}[2.5.1]
\label{II.2.5.1}
Cho $M$ là một $S$-môđun \emph{phân bậc}.
Với mọi $f$ thuần nhất trong $S_+$, $M_{(f)}$ là một $S_{(f)}$-môđun, và do đó tương ứng với nó là một bó quasi-coherent liên kết $(M_{(f)})\supertilde$ trên lược đồ affine $\Spec(S_{(f)})$, được đồng nhất với $D_+(f)$ \sref[I]{I.1.3.4}.
\end{env}
\oldpage[II]{31}
\begin{proposition}[2.5.2]
\label{II.2.5.2}
Trên $X=\Proj(S)$ tồn tại một và chỉ một $\sh{O}_X$-môđun quasi-coherent $\widetilde{M}$ sao cho, với mọi $f$ thuần nhất trong $S_+$, ta có $\Gamma(D_+(f),\widetilde{M})=M_{(f)}$, và đồng cấu hạn chế $\Gamma(D_+(f),\widetilde{M})\to\Gamma(D_+(fg),\widetilde{M})$, với $f,g$ thuần nhất trong $S_+$, là đồng cấu chính tắc $M_{(f)}\to M_{(fg)}$ \sref{II.2.2.3}.
\end{proposition}

\begin{proof}
Giả sử $f\in S_d$ và $g\in S_e$.
Vì $D_+(fg)$ được đồng nhất với phổ nguyên tố của $(S_{(f)})_{g^d/f^e}$ theo \sref{II.2.2.2}, nên hạn chế trên $D_+(fg)$ của bó $(M_{(f)})\supertilde$ trên $D_+(f)$ được đồng nhất chính tắc với bó liên kết với môđun $(M_{(f)})_{g^d/f^e}$ \sref[I]{I.1.3.6}, và do đó cũng với $(M_{(fg)})\supertilde$ \sref{II.2.2.2}.
Suy ra tồn tại một đẳng cấu chính tắc
\[
  \theta_{g,f}:(M_{(f)})\supertilde|D_+(fg)
  \simto(M_{(g)})\supertilde|D_+(fg)
\]
sao cho, nếu $h$ là một phần tử thuần nhất thứ ba của $S_+$, thì $\theta_{f,h}=\theta_{f,g}\circ\theta_{g,h}$ trên $D_+(fgh)$.
Do đó, theo \sref[0]{0.3.3.1}, trên $X$ tồn tại một $\sh{O}_X$-môđun quasi-coherent $\sh{F}$ và, với mỗi $f$ thuần nhất trong $S_+$, một đẳng cấu $\eta_f$ từ $\sh{F}|D_+(f)$ lên $(M_{(f)})\supertilde$ sao cho $\theta_{g,f}=\eta_g\circ\eta_f^{-1}$.
Nếu bây giờ ta xét bó $\sh{G}$ liên kết với tiền bó (trên cơ sở của tôpô của $X$ gồm các $D_+(f)$) được xác định bởi $D_+(f)\mapsto M_{(f)}$, với các đồng cấu chính tắc $M_{(f)}\to M_{(fg)}$ làm các đồng cấu hạn chế, thì những điều trên chứng minh rằng $\sh{F}$ và $\sh{G}$ đẳng cấu với nhau (có xét đến \sref[I]{I.1.3.7}); bó $\sh{G}$ được ký hiệu là $\widetilde{M}$ và quả thật thỏa các điều kiện của phát biểu.
Đặc biệt, ta có $\widetilde{S}=\sh{O}_X$.
\end{proof}

\begin{definition}[2.5.3]
\label{II.2.5.3}
Ta nói $\sh{O}_X$-môđun quasi-coherent $\widetilde{M}$ được định nghĩa trong \sref{II.2.5.2} là \emph{liên kết} với $S$-môđun phân bậc $M$.
\end{definition}

Nhắc lại rằng các $S$-môđun phân bậc tạo thành một phạm trù khi ta giới hạn các đồng cấu môđun phân bậc ở các đồng cấu \emph{bậc $0$}.
Với quy ước này:

\begin{proposition}[2.5.4]
\label{II.2.5.4}
Hàm tử $M\mapsto\widetilde{M}$ là một hàm tử đồng biến, cộng tính và khớp từ phạm trù các $S$-môđun phân bậc vào phạm trù các $\sh{O}_X$-môđun quasi-coherent, và giao hoán với các giới hạn quy nạp và các tổng trực tiếp.
\end{proposition}

\begin{proof}
Vì các tính chất này có tính địa phương, chỉ cần kiểm tra chúng trên các bó dạng $\widetilde{M}|D_+(f)=(M_{(f)})\supertilde$.
Thế nhưng hàm tử địa phương hóa $M\mapsto M_f$, hàm tử lấy phần bậc không $N\mapsto N_0$ (trong phạm trù các $S_f$-môđun phân bậc), và hàm tử bó liên kết $P\mapsto\widetilde{P}$ (trong phạm trù các $S_{(f)}$-môđun) đều có ba tính chất: tính khớp và tính giao hoán với các giới hạn quy nạp và các tổng trực tiếp \sref[I]{I.1.3.5} và \sref[I]{I.1.3.9}; từ đó suy ra mệnh đề.
\end{proof}

Ta ký hiệu bởi $\widetilde{u}$ đồng cấu $\widetilde{M}\to\widetilde{N}$ tương ứng với một đồng cấu bậc $0$, $u:M\to N$.
Từ \sref{II.2.5.4}, ta suy ra ngay rằng các kết quả của \sref[I]{I.1.3.9} và \sref[I]{I.1.3.10} vẫn đúng đối với các $S$-môđun phân bậc và các đồng cấu bậc $0$ (với nghĩa đã cho ở đây của $\widetilde{M}$), vì các chứng minh hoàn toàn hình thức.

\begin{proposition}[2.5.5]
\label{II.2.5.5}
Với mọi $\mathfrak{p}\in X=\Proj(S)$, ta có $\widetilde{M}_\mathfrak{p}=M_{(\mathfrak{p})}$.
\end{proposition}

\begin{proof}
Theo định nghĩa, $\widetilde{M}_\mathfrak{p}=\varinjlim\Gamma(D_+(f),\widetilde{M})$, trong đó $f$ chạy qua tập các phần tử thuần nhất của $S_+$ sao cho $f\notin\mathfrak{p}$.
Vì $\Gamma(D_+(f),\widetilde{M})=M_{(f)}$, mệnh đề suy ra từ định nghĩa của $M_{(\mathfrak{p})}$ \sref{II.2.2.7}.
\end{proof}
\oldpage[II]{32}
Đặc biệt, \emph{vành địa phương $(\sh{O}_X)_\mathfrak{p}$} chính là vành $S_{(\mathfrak{p})}$, tức tập các phân thức $x/f$, trong đó $f$ thuần nhất trong $S_+$ và không thuộc $\mathfrak{p}$, còn $x$ thuần nhất cùng bậc với $f$.

Đặc biệt hơn, nếu $S$ là \emph{cốt yếu nguyên}, sao cho $\Proj(S)=X$ là \emph{nguyên} \sref{II.2.4.4}, và nếu $\xi=(0)$ là điểm tổng quát của $X$, thì \emph{trường các hàm hữu tỉ $R(X)=\sh{O}_\xi$} là trường gồm các phần tử $fg^{-1}$, trong đó $f$ và $g$ thuần nhất cùng bậc trong $S_+$ và $g\neq0$.

\begin{proposition}[2.5.6]
\label{II.2.5.6}
Nếu, với mọi $z\in M$ và mọi $f$ thuần nhất trong $S_+$, tồn tại một lũy thừa của $f$ triệt tiêu $z$, thì $\widetilde{M}=0$.
Điều kiện đủ này cũng là điều kiện cần khi $S$, với tư cách một $S_0$-đại số, được sinh bởi tập $S_1$ các phần tử thuần nhất bậc $1$.
\end{proposition}

\begin{proof}
Điều kiện $\widetilde{M}=0$ tương đương với $M_{(f)}=0$ với mọi $f$ thuần nhất trong $S_+$.
Mặt khác, nếu $f\in S_d$, thì $M_{(f)}=0$ có nghĩa là, với mọi $z\in M$ thuần nhất có bậc là bội của $d$, tồn tại một lũy thừa $f^n$ sao cho $f^nz=0$.
Nếu $M_{(f)}=0$ với mọi $f\in S_1$, thì điều kiện trong phát biểu được thỏa đối với mọi $f\in S_1$; điều kiện ấy lại càng được thỏa đối với mọi $f$ thuần nhất trong $S_+$ khi $S_1$ sinh $S$, vì khi đó mọi phần tử thuần nhất của $S_+$ là một tổ hợp tuyến tính của các tích những phần tử thuộc $S_1$.
\end{proof}

\begin{proposition}[2.5.7]
\label{II.2.5.7}
Cho $d>0$ là một số nguyên và $f\in S_d$.
Khi đó, với mọi $n\in\bb{Z}$, $(\sh{O}_X|D_+(f))$-môđun $(S(nd))\supertilde|D_+(f)$ đẳng cấu chính tắc với $\sh{O}_X|D_+(f)$.
\end{proposition}

\begin{proof}
Thật vậy, phép nhân với phần tử khả nghịch $(f/1)^n$ của $S_f$ là một song ánh từ $S_{(f)}=(S_f)_0$ lên
\[
  (S_f)_{nd}=(S_f(nd))_0=(S(nd)_f)_0=S(nd)_{(f)}.
\]
Nói cách khác, các $S_{(f)}$-môđun $S_{(f)}$ và $S(nd)_{(f)}$ đẳng cấu chính tắc với nhau; từ đó suy ra mệnh đề.
\end{proof}

\begin{corollary}[2.5.8]
\label{II.2.5.8}
Trên tập con mở $U=\bigcup_{f\in S_d}D_+(f)$, hạn chế của $\sh{O}_X$-môđun $(S(nd))\supertilde$ là một $(\sh{O}_X|U)$-môđun khả nghịch \sref[0]{0.5.4.1}.
\end{corollary}

\begin{corollary}[2.5.9]
\label{II.2.5.9}
Nếu iđêan $S_+$ của $S$ được sinh bởi tập $S_1$ các phần tử thuần nhất bậc $1$, thì các $\sh{O}_X$-môđun $(S(n))\supertilde$ đều khả nghịch với mọi $n\in\bb{Z}$.
\end{corollary}

\begin{proof}
Chỉ cần lưu ý rằng $X=\bigcup_{f\in S_1}D_+(f)$ theo giả thiết \sref{II.2.3.14}, rồi áp dụng \sref{II.2.5.8} với $U=X$.
\end{proof}

\begin{env}[2.5.10]
\label{II.2.5.10}
Trong toàn bộ phần còn lại của mục này, ta đặt
\[
\label{II.2.5.10.1}
  \sh{O}_X(n)=(S(n))\supertilde
\tag{2.5.10.1}
\]
với mọi $n\in\bb{Z}$ và, với mọi tập con mở $U$ của $X$ cùng mọi $(\sh{O}_X|U)$-môđun $\sh{F}$,
\[
\label{II.2.5.10.2}
  \sh{F}(n)=\sh{F}\otimes_{\sh{O}_X|U}(\sh{O}_X(n)|U)
\tag{2.5.10.2}
\]
với mọi $n\in\bb{Z}$.
Nếu iđêan $S_+$ được sinh bởi $S_1$, thì hàm tử $\sh{F}\mapsto\sh{F}(n)$ là \emph{khớp} theo $\sh{F}$ với mọi $n\in\bb{Z}$, vì khi đó $\sh{O}_X(n)$ là một $\sh{O}_X$-môđun \emph{khả nghịch}.
\end{env}

\begin{env}[2.5.11]
\label{II.2.5.11}
Cho $M$ và $N$ là hai $S$-môđun phân bậc.
Với mọi $f\in S_d$ ($d>0$), ta định nghĩa một đồng cấu chính tắc có tính hàm tử của các $S_{(f)}$-môđun
\[
\label{II.2.5.11.1}
  \lambda_f:M_{(f)}\otimes_{S_{(f)}}N_{(f)}
  \to(M\otimes_S N)_{(f)}
\tag{2.5.11.1}
\]
\oldpage[II]{33}
bằng cách hợp thành đồng cấu $M_{(f)}\otimes_{S_{(f)}}N_{(f)}\to M_f\otimes_{S_f}N_f$ (cảm sinh từ các phép nhúng $M_{(f)}\to M_f$, $N_{(f)}\to N_f$ và $S_{(f)}\to S_f$) với đẳng cấu chính tắc $M_f\otimes_{S_f}N_f\simto(M\otimes_S N)_f$ \sref[0]{0.1.3.4}, đồng thời lưu ý rằng theo định nghĩa tích tenxơ của hai môđun phân bậc, đẳng cấu sau cùng này bảo toàn bậc.
Với $x\in M_{md}$ và $y\in N_{nd}$ ($m,n\geq0$), do đó ta có
\[
  \lambda_f((x/f^m)\otimes(y/f^n))=(x\otimes y)/f^{m+n}.
\]

Từ định nghĩa này suy ra ngay rằng, nếu $g\in S_e$ ($e>0$), thì sơ đồ
\[
  \xymatrix{
    M_{(f)}\otimes_{S_{(f)}}N_{(f)} \ar[r]^{\lambda_f} \ar[d]
    & (M\otimes_S N)_{(f)} \ar[d]
  \\M_{(fg)}\otimes_{S_{(fg)}}N_{(fg)} \ar[r]_{\lambda_{fg}}
    & (M\otimes_S N)_{(fg)}
  }
\]
(trong đó mũi tên thẳng đứng bên phải là đồng cấu chính tắc, còn mũi tên bên trái cảm sinh từ các đồng cấu chính tắc) là giao hoán.
Do đó, các $\lambda_f$ cảm sinh một đồng cấu chính tắc có tính hàm tử của các $\sh{O}_X$-môđun
\[
\label{II.2.5.11.2}
  \lambda:\widetilde{M}\otimes_{\sh{O}_X}\widetilde{N}
  \to(M\otimes_S N)\supertilde.
\tag{2.5.11.2}
\]

Đặc biệt, xét hai iđêan phân bậc $\mathfrak{J},\mathfrak{K}$ của $S$.
Vì $\widetilde{\mathfrak{J}}$ và $\widetilde{\mathfrak{K}}$ là các bó iđêan của $\sh{O}_X$, nên ta có một đồng cấu chính tắc $\widetilde{\mathfrak{J}}\otimes_{\sh{O}_X}\widetilde{\mathfrak{K}}\to\sh{O}_X$; khi đó sơ đồ
\[
\label{II.2.5.11.3}
  \xymatrix{
    \widetilde{\mathfrak{J}}\otimes_{\sh{O}_X}\widetilde{\mathfrak{K}} \ar[rr]^\lambda \ar[dr]
    &&(\mathfrak{J}\otimes_S\mathfrak{K})\supertilde \ar[dl]
  \\&\sh{O}_X&
  }
\tag{2.5.11.3}
\]
giao hoán.
Thật vậy, ta quy việc kiểm tra về từng tập mở $D_+(f)$ (với $f$ thuần nhất trong $S_+$), và điều này suy ra ngay từ định nghĩa \sref{II.2.5.11.1} của $\lambda_f$ cùng \sref[I]{I.1.3.13}.

Cuối cùng, lưu ý rằng nếu $M,N,P$ là ba $S$-môđun phân bậc, thì sơ đồ
\[
\label{II.2.5.11.4}
  \xymatrix{
    \widetilde{M}\otimes_{\sh{O}_X}\widetilde{N}\otimes_{\sh{O}_X}\widetilde{P} \ar[r]^{\lambda\otimes1} \ar[d]_{1\otimes\lambda}
    &(M\otimes_S N)\supertilde\otimes_{\sh{O}_X}\widetilde{P} \ar[d]^\lambda
  \\widetilde{M}\otimes_{\sh{O}_X}(N\otimes_S P)\supertilde \ar[r]_\lambda
    &(M\otimes_S N\otimes_S P)\supertilde
  }
\tag{2.5.11.4}
\]
\oldpage[II]{34}
giao hoán.
Một lần nữa, chỉ cần kiểm tra điều đó trên từng $D_+(f)$, và điều này suy ra ngay từ các định nghĩa cùng \sref[I]{I.1.3.13}.
\end{env}

\begin{env}[2.5.12]
\label{II.2.5.12}
Dưới các giả thiết của \sref{II.2.5.11}, ta định nghĩa một đồng cấu chính tắc có tính hàm tử của các $S_{(f)}$-môđun
\[
\label{II.2.5.12.1}
  \mu_f:(\Hom_S(M,N))_{(f)}
  \to\Hom_{S_{(f)}}(M_{(f)},N_{(f)})
\tag{2.5.12.1}
\]
bằng cách cho $u/f^n$, trong đó $u$ là một đồng cấu bậc $nd$, tương ứng với đồng cấu $M_{(f)}\to N_{(f)}$ biến $x/f^m$ ($x\in M_{md}$) thành $u(x)/f^{m+n}$.
Với $g\in S_e$ ($e>0$), ta lại có một sơ đồ giao hoán:
\[
  \xymatrix{
    (\Hom_S(M,N))_{(f)} \ar[r]^{\mu_f} \ar[d]
    & \Hom_{S_{(f)}}(M_{(f)},N_{(f)}) \ar[d]
  \\(\Hom_S(M,N))_{(fg)} \ar[r]_{\mu_{fg}}
    & \Hom_{S_{(fg)}}(M_{(fg)},N_{(fg)})
  }
\]
(trong đó mũi tên bên trái là đồng cấu chính tắc, còn mũi tên bên phải cảm sinh từ các đồng cấu chính tắc).
Một lần nữa, có xét đến \sref[I]{I.1.3.8}, ta kết luận rằng các $\mu_f$ xác định một đồng cấu chính tắc có tính hàm tử của các $\sh{O}_X$-môđun
\[
\label{II.2.5.12.2}
  \mu:(\Hom_S(M,N))\supertilde
  \to\shHom_{\sh{O}_X}(\widetilde{M},\widetilde{N}).
\tag{2.5.12.2}
\]
\end{env}

\begin{proposition}[2.5.13]
\label{II.2.5.13}
Giả sử iđêan $S_+$ được sinh bởi $S_1$.
Khi đó, đồng cấu $\lambda$ \sref{II.2.5.11.2} là một đẳng cấu; đồng cấu $\mu$ \sref{II.2.5.12.2} cũng vậy khi $S$-môđun phân bậc $M$ có một trình bày hữu hạn \sref{II.2.1.1}.
\end{proposition}

\begin{proof}
Vì $X$ là hợp của các $D_+(f)$ với $f\in S_1$ \sref{II.2.3.14}, ta quy về việc chứng minh rằng $\lambda_f$ và $\mu_f$ là các đẳng cấu, dưới các giả thiết đang xét, khi $f$ thuần nhất và \emph{có bậc $1$}.
Khi đó, ta có thể định nghĩa một ánh xạ $\bb{Z}$-song tuyến tính $M_m\times N_n\to M_{(f)}\otimes_{S_{(f)}}N_{(f)}$ bằng cách cho $(x,y)$ tương ứng với phần tử $(x/f^m)\otimes(y/f^n)$ (khi $m<0$, ta viết $x/f^m$ để chỉ $f^{-m}x/1$).
Các ánh xạ này xác định một ánh xạ $\bb{Z}$-tuyến tính $M\otimes_{\bb{Z}}N\to M_{(f)}\otimes_{S_{(f)}}N_{(f)}$ và, nếu $s\in S_q$, ánh xạ này biến $(sx)\otimes y$ thành
\[
  (s/f^q)((x/f^m)\otimes(y/f^n))
\]
với $x\in M_m$ và $y\in N_n$.
Do đó ta thu được một di-đồng cấu môđun
\[
  \gamma_f:M\otimes_S N\to M_{(f)}\otimes_{S_{(f)}}N_{(f)}
\]
ứng với đồng cấu chính tắc $S\to S_{(f)}$ biến $s\in S_q$ thành $s/f^q$.
Giả sử thêm rằng, với một phần tử $\sum_i(x_i\otimes y_i)$ của $M\otimes_S N$, trong đó $x_i,y_i$ thuần nhất có các bậc lần lượt là $m_i,n_i$, ta có
\[
  f^r\sum_i(x_i\otimes y_i)=0,
\]
nói cách khác, $\sum_i(f^rx_i\otimes y_i)=0$.
Từ \sref[0]{0.1.3.4}, ta suy ra
\[
  \sum_i(f^rx_i/f^{m_i+r})\otimes(y_i/f^{n_i})=0,
\]
tức là $\gamma_f(\sum_i(x_i\otimes y_i))=0$.
Vì vậy, $\gamma_f$ phân tích thành
\[
  M\otimes_S N\to(M\otimes_S N)_f
  \xrightarrow{\gamma'_f}M_{(f)}\otimes_{S_{(f)}}N_{(f)}.
\]
Nếu $\lambda'_f$ là hạn chế của $\gamma'_f$
\oldpage[II]{35}
trên $(M\otimes_S N)_{(f)}$, thì ta kiểm tra ngay rằng $\lambda_f$ và $\lambda'_f$ là các $S_{(f)}$-đồng cấu nghịch đảo của nhau, từ đó suy ra phần thứ nhất của mệnh đề.

Để chứng minh phần thứ hai, giả sử $M$ là đối hạch của một đồng cấu $P\to Q$ giữa các $S$-môđun phân bậc, trong đó $P$ và $Q$ là các tổng trực tiếp của một số hữu hạn môđun dạng $S(n)$; dùng tính khớp trái của $\Hom_S(L,N)$ theo $L$ và tính khớp của $M_{(f)}$ theo $M$, ta quy ngay về việc chứng minh rằng $\mu_f$ là một đẳng cấu khi $M=S(n)$.
Với mọi phần tử thuần nhất $z$ của $N$, gọi $u_z$ là đồng cấu từ $S(n)$ vào $N$ sao cho $u_z(1)=z$; ta thấy ngay rằng $\eta:z\mapsto u_z$ là một đẳng cấu bậc $0$ từ $N(-n)$ lên $\Hom_S(S(n),N)$.
Tương ứng với nó là một đẳng cấu
\[
  \eta_f:(N(-n))_{(f)}\to(\Hom_S(S(n),N))_{(f)}.
\]

Mặt khác, gọi $\eta'_f$ là đẳng cấu $N_{(f)}\to\Hom_{S_{(f)}}(S(n)_{(f)},N_{(f)})$ gửi mỗi $z'\in N_{(f)}$ tới đồng cấu $v_{z'}$ sao cho
\[
  v_{z'}(s/f^k)=sz'/f^{n+k}
\]
với $s\in S_{n+k}=(S(n))_k$.
Ta dễ dàng nhận thấy rằng ánh xạ hợp thành
\[
  (N(-n))_{(f)}
  \xrightarrow{\eta_f}(\Hom_S(S(n),N))_{(f)}
  \xrightarrow{\mu_f}\Hom_{S_{(f)}}(S(n)_{(f)},N_{(f)})
  \xrightarrow{{\eta'_f}^{-1}}N_{(f)}
\]
là đẳng cấu $z/f^h\mapsto z/f^{h-n}$ từ $(N(-n))_{(f)}$ lên $N_{(f)}$; do đó $\mu_f$ là một đẳng cấu.
\end{proof}

Nếu iđêan $S_+$ được sinh bởi $S_1$, thì từ \sref{II.2.5.13} suy ra rằng, với mọi iđêan phân bậc $\mathfrak{J}$ của $S$ và mọi $S$-môđun phân bậc $M$, ta có
\[
\label{II.2.5.13.1}
  \widetilde{\mathfrak{J}}\cdot\widetilde{M}
  =(\mathfrak{J}\cdot M)\supertilde
\tag{2.5.13.1}
\]
sai khác một đẳng cấu chính tắc; thật vậy, điều này suy ra từ tính giao hoán của sơ đồ
\[
  \xymatrix{
    \widetilde{\mathfrak{J}}\otimes_{\sh{O}_X}\widetilde{M} \ar[rr]^\lambda \ar[dr]
    &&(\mathfrak{J}\otimes_S M)\supertilde \ar[dl]
  \\&\widetilde{M}&
  }
\]
mà ta kiểm tra như đối với sơ đồ \sref{II.2.5.11.3}.

\begin{corollary}[2.5.14]
\label{II.2.5.14}
Giả sử $S$ được sinh bởi $S_1$.
Với mọi $m,n\in\bb{Z}$, khi đó ta có:
\[
\label{II.2.5.14.1}
  \sh{O}_X(m)\otimes_{\sh{O}_X}\sh{O}_X(n)=\sh{O}_X(m+n)
\tag{2.5.14.1}
\]
\[
\label{II.2.5.14.2}
  \sh{O}_X(n)=(\sh{O}_X(1))^{\otimes n}
\tag{2.5.14.2}
\]
sai khác các đẳng cấu chính tắc.
\end{corollary}

\begin{proof}
Công thức thứ nhất suy ra từ \sref{II.2.5.13} và từ sự tồn tại của đẳng cấu chính tắc bậc $0$
\[
  S(m)\otimes_S S(n)\simto S(m+n),
\]
gửi phần tử $1\otimes1$ (trong đó phần tử $1$ thứ nhất thuộc $(S(m))_{-m}$ và phần tử thứ hai thuộc $(S(n))_{-n}$) tới phần tử $1\in(S(m+n))_{-(m+n)}$.
Sau đó chỉ cần chứng minh công thức thứ hai khi $n=-1$; theo \sref{II.2.5.13}, điều này quy về việc nhận thấy rằng $\Hom_S(S(1),S)$ đẳng cấu chính tắc với $S(-1)$, điều được kiểm tra ngay bằng cách trở lại các định nghĩa \sref{II.2.1.2} và nhớ rằng $S(1)$ là một $S$-môđun cyclic.
\end{proof}
\oldpage[II]{36}
\begin{corollary}[2.5.15]
\label{II.2.5.15}
Giả sử $S$ được sinh bởi $S_1$.
Khi đó, với mọi $S$-môđun phân bậc $M$ và mọi $n\in\bb{Z}$, ta có
\[
\label{II.2.5.15.1}
  (M(n))\supertilde=\widetilde{M}(n)
\tag{2.5.15.1}
\]
sai khác một đẳng cấu chính tắc.
\end{corollary}

\begin{proof}
Điều này suy ra từ các định nghĩa \sref{II.2.5.10.2} và \sref{II.2.5.10.1}, từ Mệnh đề~\sref{II.2.5.13}, và từ sự tồn tại của một đẳng cấu chính tắc $M(n)\simto M\otimes_S S(n)$ bậc $0$, gửi mỗi $z\in(M(n))_h=M_{n+h}$ tới
\[
  z\otimes1\in M_{n+h}\otimes(S(n))_{-n}
  \subset(M\otimes_S S(n))_h.
\]
\end{proof}

\begin{env}[2.5.16]
\label{II.2.5.16}
Ký hiệu bởi $S'$ vành phân bậc sao cho $S'_0=\bb{Z}$ và $S'_n=S_n$ với $n>0$.
Khi đó, nếu $f\in S_d$ ($d>0$), ta có $(S(n))_{(f)}=(S'(n))_{(f)}$ với mọi $n\in\bb{Z}$, vì một phần tử của $(S'(n))_{(f)}$ có dạng $x/f^k$, với $x\in S'_{n+kd}$ ($k>0$), và ta luôn có thể chọn $k$ sao cho $n+kd\neq0$.
Vì $X=\Proj(S)$ và $X'=\Proj(S')$ được đồng nhất chính tắc \sref{II.2.4.7}[(ii)], ta thấy rằng, với mọi $n\in\bb{Z}$, $\sh{O}_X(n)$ và $\sh{O}_{X'}(n)$ là ảnh của nhau qua phép đồng nhất ấy.

Mặt khác, lưu ý rằng, với mọi $d>0$ và mọi $n\in\bb{Z}$, ta có
\[
  (S^{(d)}(n))_h=S_{(n+h)d}=(S(nd))_{hd}
\]
Do đó, với $f\in S_d$, ta có $(S^{(d)}(n))_{(f)}=(S(nd))_{(f)}$.
Ta biết rằng các lược đồ $X=\Proj(S)$ và $X^{(d)}=\Proj(S^{(d)})$ được đồng nhất chính tắc \sref{II.2.4.7}[(i)].
Những điều trên cho thấy rằng, nếu $S_0$-đại số $S^{(d)}$ được sinh bởi $S_d$, thì $\sh{O}_X(nd)$ và $\sh{O}_{X^{(d)}}(n)$ là ảnh của nhau qua phép đồng nhất này, với mọi $n\in\bb{Z}$.
\end{env}

\begin{proposition}[2.5.17]
\label{II.2.5.17}
Cho $d>0$ là một số nguyên, và đặt $U=\bigcup_{f\in S_d}D_+(f)$.
Khi đó, với mọi số nguyên $n$, hạn chế lên $U$ của đồng cấu chính tắc
\[
  \sh{O}_X(nd)\otimes_{\sh{O}_X}\sh{O}_X(-nd)\to\sh{O}_X
\]
là một đẳng cấu.
\end{proposition}

\begin{proof}
Theo \sref{II.2.5.16}, ta có thể quy về trường hợp $d=1$, và khi đó kết luận suy ra từ chứng minh của \sref{II.2.5.13}.
\end{proof}

\subsection{$S$-môđun phân bậc liên kết với một bó trên $\operatorname{Proj}(S)$}
\label{subsection:II.2.6}

\emph{Trong toàn bộ mục này, ta giả sử iđêan $S_+$ của $S$ được sinh bởi tập $S_1$ các phần tử thuần nhất bậc $1$.}

\begin{env}[2.6.1]
\label{II.2.6.1}
$\sh{O}_X$-môđun $\sh{O}_X(1)$ khi đó là \emph{khả nghịch} \sref{II.2.5.9}; với mọi $\sh{O}_X$-môđun $\sh{F}$, ta đặt \sref[0]{0.5.4.6}
\[
\label{II.2.6.1.1}
  \Gamma_\bullet(\sh{F})
  =\Gamma_\bullet(\sh{O}_X(1),\sh{F})
  =\bigoplus_{n\in\bb{Z}}\Gamma(X,\sh{F}(n))
\tag{2.6.1.1}
\]
có xét đến \sref{II.2.5.14.2}.
Nhắc lại \sref[0]{0.5.4.6} rằng $\Gamma_\bullet(\sh{O}_X)$ được trang bị cấu trúc \emph{vành phân bậc}, còn $\Gamma_\bullet(\sh{F})$ được trang bị cấu trúc \emph{$\Gamma_\bullet(\sh{O}_X)$-môđun phân bậc}.

Vì $\sh{O}_X(n)$ tự do địa phương, $\Gamma_\bullet(\sh{F})$ là một hàm tử hiệp biến cộng tính \emph{khớp trái} theo $\sh{F}$; đặc biệt, nếu $\sh{J}$ là một bó iđêan của $\sh{O}_X$, thì $\Gamma_\bullet(\sh{J})$ được đồng nhất chính tắc với một \emph{iđêan phân bậc} của $\Gamma_\bullet(\sh{O}_X)$.
\end{env}

\begin{env}[2.6.2]
\label{II.2.6.2}
Cho $M$ là một $S$-môđun phân bậc.
Với mọi $f\in S_d$ ($d>0$), $x\mapsto x/1$ là một đồng cấu nhóm abel $M_0\to M_{(f)}$; vì $M_{(f)}$ được đồng nhất chính tắc
\oldpage[II]{37}
với $\Gamma(D_+(f),\widetilde{M})$, do đó ta thu được một đồng cấu nhóm abel
\[
  \alpha_0^f:M_0\to\Gamma(D_+(f),\widetilde{M}).
\]
Hiển nhiên, với mọi $g\in S_e$ ($e>0$), sơ đồ
\[
  \xymatrix{
    & \Gamma(D_+(f),\widetilde{M}) \ar[dd]
  \\M_0 \ar[ur]^{\alpha_0^f} \ar[dr]_{\alpha_0^{fg}}
  \\& \Gamma(D_+(fg),\widetilde{M})
  }
\]
giao hoán; điều này có nghĩa là, với mọi $x\in M_0$, các tiết diện $\alpha_0^f(x)$ và $\alpha_0^g(x)$ của $\widetilde{M}$ trùng nhau trên $D_+(f)\cap D_+(g)$, và do đó tồn tại một tiết diện duy nhất $\alpha_0(x)\in\Gamma(X,\widetilde{M})$ mà hạn chế của nó lên mỗi $D_+(f)$ là $\alpha_0^f(x)$.
Như vậy, không dùng giả thiết $S$ được sinh bởi $S_1$, ta đã định nghĩa một đồng cấu nhóm abel
\[
\label{II.2.6.2.1}
  \alpha_0:M_0\to\Gamma(X,\widetilde{M}).
\tag{2.6.2.1}
\]

Áp dụng kết quả này cho $S$-môđun phân bậc $M(n)$, với mỗi $n\in\bb{Z}$, ta thu được, với mỗi $n\in\bb{Z}$, một đồng cấu nhóm abel
\[
\label{II.2.6.2.2}
  \alpha_n:M_n=(M(n))_0\to\Gamma(X,\widetilde{M}(n))
\tag{2.6.2.2}
\]
có xét đến \sref{II.2.5.15}; từ đó ta thu được một đồng cấu có tính hàm tử (bậc $0$) giữa các nhóm abel phân bậc
\[
\label{II.2.6.2.3}
  \alpha:M\to\Gamma_\bullet(\widetilde{M})
\tag{2.6.2.3}
\]
(cũng được ký hiệu bởi $\alpha_M$), mà trên mỗi $M_n$ trùng với $\alpha_n$.

Đặc biệt, nếu lấy $M=S$, thì ta kiểm tra ngay, có xét đến định nghĩa của phép nhân trong $\Gamma_\bullet(\sh{O}_X)$ \sref[0]{0.5.4.6}, rằng $\alpha:S\to\Gamma_\bullet(\sh{O}_X)$ là một đồng cấu các vành phân bậc và, với mọi $S$-môđun phân bậc $M$, \sref{II.2.6.2.3} là một di-đồng cấu của các môđun phân bậc.
\end{env}

\begin{proposition}[2.6.3]
\label{II.2.6.3}
Với mọi $f\in S_d$ ($d>0$), $D_+(f)$ trùng với tập các $\mathfrak{p}\in X$ tại đó tiết diện $\alpha_d(f)$ của $\sh{O}_X(d)$ không triệt tiêu \sref[0]{0.5.5.2}.
\end{proposition}

\begin{proof}
Vì theo giả thiết $X=\bigcup_{g\in S_1}D_+(g)$, chỉ cần kiểm tra rằng, với mọi $g\in S_1$, tập các $\mathfrak{p}\in D_+(g)$ tại đó $\alpha_d(f)$ không triệt tiêu trùng với $D_+(fg)$.
Thật vậy, theo định nghĩa, hạn chế của $\alpha_d(f)$ lên $D_+(g)$ là tiết diện tương ứng với phần tử $f/1$ của $(S(d))_{(g)}$.
Qua đẳng cấu chính tắc $(S(d))_{(g)}\simto S_{(g)}$ \sref{II.2.5.7}, tiết diện này của $\sh{O}_X(d)$ trên $D_+(g)$ được đồng nhất với tiết diện của $\sh{O}_X$ trên $D_+(g)$ tương ứng với phần tử $f/g^d$ của $S_{(g)}$.
Nói rằng tiết diện này triệt tiêu tại $\mathfrak{p}\in D_+(g)$ có nghĩa là $f/g^d\in\mathfrak{q}$, trong đó $\mathfrak{q}$ là iđêan nguyên tố của $S_{(g)}$ tương ứng với $\mathfrak{p}$ \sref{II.2.3.6}; theo định nghĩa, điều này có nghĩa là $f\in\mathfrak{p}$, từ đó suy ra mệnh đề.
\end{proof}

\begin{env}[2.6.4]
\label{II.2.6.4}
Bây giờ cho $\sh{F}$ là một $\sh{O}_X$-môđun, và đặt $M=\Gamma_\bullet(\sh{F})$.
Nhờ sự tồn tại của đồng cấu các vành phân bậc $\alpha:S\to\Gamma_\bullet(\sh{O}_X)$, ta có thể xem $M$ như một $S$-môđun phân bậc.
Với mọi $f\in S_d$ ($d>0$), từ \sref{II.2.6.3} suy ra rằng hạn chế lên $D_+(f)$ của tiết diện $\alpha_d(f)$ của $\sh{O}_X(d)$ là khả nghịch; do đó, với mọi $n>0$, hạn chế lên $D_+(f)$ của tiết diện $\alpha_{nd}(f^n)$ của $\sh{O}_X(nd)$ cũng khả nghịch.
Cho $z\in M_{nd}=\Gamma(X,\sh{F}(nd))$ ($n>0$); nếu tồn tại một số nguyên $k\geq0$ sao cho hạn chế lên $D_+(f)$ của $f^kz$, tức là
\oldpage[II]{38}
tiết diện $(z|D_+(f))(\alpha_{kd}(f^k)|D_+(f))$ của $\sh{F}((n+k)d)$, bằng không, thì theo nhận xét trên ta cũng có $z|D_+(f)=0$.
Điều này cho thấy ta định nghĩa được một $S_{(f)}$-đồng cấu $\beta_f:M_{(f)}\to\Gamma(D_+(f),\sh{F})$ bằng cách gửi phần tử $z/f^n$ tới tiết diện
\[
  (z|D_+(f))(\alpha_{nd}(f^n)|D_+(f))^{-1}
\]
của $\sh{F}$ trên $D_+(f)$.
Hơn nữa, ta kiểm tra ngay rằng sơ đồ
\[
\label{II.2.6.4.1}
  \xymatrix{
    M_{(f)} \ar[r]^{\beta_f} \ar[d]
    & \Gamma(D_+(f),\sh{F}) \ar[d]
  \\M_{(fg)} \ar[r]_{\beta_{fg}}
    & \Gamma(D_+(fg),\sh{F})
  }
\tag{2.6.4.1}
\]
giao hoán với $g\in S_e$ ($e>0$).
Nhắc lại rằng $M_{(f)}$ được đồng nhất chính tắc với $\Gamma(D_+(f),\widetilde{M})$ và các $D_+(f)$ lập thành một cơ sở cho tôpô của $X$ \sref{II.2.3.4}, ta thấy các $\beta_f$ xuất phát từ một đồng cấu chính tắc duy nhất của các $\sh{O}_X$-môđun
\[
\label{II.2.6.4.2}
  \beta:(\Gamma_\bullet(\sh{F}))\supertilde\to\sh{F}
\tag{2.6.4.2}
\]
(cũng được ký hiệu bởi $\beta_{\sh{F}}$), hiển nhiên có tính hàm tử.
\end{env}

\begin{proposition}[2.6.5]
\label{II.2.6.5}
Cho $M$ là một $S$-môđun phân bậc và $\sh{F}$ là một $\sh{O}_X$-môđun; khi đó các đồng cấu hợp thành
\[
\label{II.2.6.5.1}
  \widetilde{M}
  \xrightarrow{\widetilde{\alpha}}(\Gamma_\bullet(\widetilde{M}))\supertilde
  \xrightarrow{\beta}\widetilde{M}
\tag{2.6.5.1}
\]
\[
\label{II.2.6.5.2}
  \Gamma_\bullet(\sh{F})
  \xrightarrow{\alpha}\Gamma_\bullet((\Gamma_\bullet(\sh{F}))\supertilde)
  \xrightarrow{\Gamma_\bullet(\beta)}\Gamma_\bullet(\sh{F})
\tag{2.6.5.2}
\]
là các đẳng cấu đồng nhất.
\end{proposition}

\begin{proof}
Việc kiểm tra \sref{II.2.6.5.1} mang tính địa phương: trên một tập mở $D_+(f)$, nó suy ra ngay từ các định nghĩa và từ việc $\beta$, khi áp dụng cho các bó quasi-coherent, được xác định bởi tác động của nó trên các tiết diện trên $D_+(f)$ \sref[I]{I.1.3.8}.
Việc kiểm tra \sref{II.2.6.5.2} được thực hiện riêng theo từng bậc: nếu đặt $M=\Gamma_\bullet(\sh{F})$, thì $M_n=\Gamma(X,\sh{F}(n))$ và
\[
  (\Gamma_\bullet(\widetilde{M}))_n
  =\Gamma(X,\widetilde{M}(n))
  =\Gamma(X,(M(n))\supertilde).
\]
Nếu $f\in S_1$ và $z\in M_n$, thì $\alpha_n^f(z)$ là phần tử $z/1$ của $(M(n))_{(f)}$, bằng $(f/1)^n(z/f^n)$; qua $\beta_f$, phần tử ấy tương ứng với tiết diện
\[
  \Big((\alpha_1(f))^n|D_+(f)\Big)
  \Big((z|D_+(f))((\alpha_1(f))^n|D_+(f))^{-1}\Big)
\]
trên $D_+(f)$, tức là hạn chế của $z$ lên $D_+(f)$, qua đó xác minh \sref{II.2.6.5.2}.
\end{proof}

\subsection{Các điều kiện hữu hạn}
\label{subsection:II.2.7}

\begin{proposition}[2.7.1]
\label{II.2.7.1}
\begin{enumerate}
  \item[{\rm(i)}] Nếu $S$ là một vành phân bậc Noether, thì $X=\Proj(S)$ là một lược đồ Noether.
  \item[{\rm(ii)}] Nếu $S$ là một $A$-đại số phân bậc hữu hạn sinh, thì $X=\Proj(S)$ là một lược đồ kiểu hữu hạn trên $Y=\Spec(A)$.
\end{enumerate}
\end{proposition}
\oldpage[II]{39}
\begin{proof}
\medskip\noindent
\begin{enumerate}
  \item[{\rm(i)}] Nếu $S$ là Noether, iđêan $S_+$ có một hệ hữu hạn các phần tử sinh thuần nhất $f_i$ ($1\leq i\leq p$); do đó, theo \sref{II.2.3.14}, không gian nền $X$ là hợp của các $D_+(f_i)=\Spec(S_{(f_i)})$, và mọi việc quy về chứng minh rằng mỗi $S_{(f_i)}$ là Noether, điều này suy ra từ \sref{II.2.2.6}.
  \item[{\rm(ii)}] Giả thiết kéo theo rằng $S_0$ là một $A$-đại số hữu hạn sinh và $S$ là một $S_0$-đại số hữu hạn sinh; do đó $S_+$ là một iđêan hữu hạn sinh \sref{II.2.1.4}.
    Khi đó, như trong (i), ta quy về chứng minh rằng, với mọi $f\in S_d$, $S_{(f)}$ là một $A$-đại số hữu hạn sinh.
    Theo \sref{II.2.2.5}, chỉ cần chứng minh rằng $S^{(d)}$ là một $A$-đại số hữu hạn sinh, điều này suy ra từ \sref{II.2.1.6}.
\end{enumerate}
\end{proof}

\begin{env}[2.7.2]
\label{II.2.7.2}
Trong phần tiếp theo, ta xét các điều kiện hữu hạn sau đây đối với một $S$-môđun phân bậc $M$:
\begin{enumerate}
  \item[{\rm(TF)}] Tồn tại một số nguyên $n$ sao cho môđun con $\bigoplus_{k\geq n}M_k$ là một $S$-môđun hữu hạn sinh.
  \item[{\rm(TN)}] Tồn tại một số nguyên $n$ sao cho $M_k=0$ với $k\geq n$.
\end{enumerate}

Nếu $M$ thỏa mãn (TN), thì $M_{(f)}=0$ với mọi phần tử thuần nhất $f$ trong $S_+$; do đó $\widetilde{M}=0$.

Cho $M$ và $N$ là hai $S$-môđun phân bậc; ta nói một đồng cấu $u:M\to N$ bậc $0$ là (TN)-\emph{đơn ánh} (tương ứng, (TN)-\emph{toàn ánh}, (TN)-\emph{song ánh}) nếu tồn tại một số nguyên $n$ sao cho $u_k:M_k\to N_k$ là đơn ánh (tương ứng, toàn ánh, song ánh) với $k\geq n$.
Như vậy, nói $u$ là (TN)-đơn ánh (tương ứng, (TN)-toàn ánh) tương đương với nói rằng $\Ker u$ (tương ứng, $\Coker u$) thỏa mãn (TN).
Theo \sref{II.2.5.4}, nếu $u$ là (TN)-đơn ánh (tương ứng, (TN)-toàn ánh, (TN)-song ánh), thì $\widetilde{u}$ là đơn ánh (tương ứng, toàn ánh, song ánh); khi $u$ là (TN)-song ánh, ta còn nói $u$ là một (TN)-\emph{đẳng cấu}.
\end{env}

\begin{proposition}[2.7.3]
\label{II.2.7.3}
Cho $S$ là một vành phân bậc sao cho iđêan $S_+$ hữu hạn sinh, và cho $M$ là một $S$-môđun phân bậc.
\begin{enumerate}
  \item[{\rm(i)}] Nếu $M$ thỏa mãn điều kiện (TF), thì $\sh{O}_X$-môđun $\widetilde{M}$ hữu hạn sinh.
  \item[{\rm(ii)}] Giả sử $M$ thỏa mãn (TF); điều kiện cần và đủ để $\widetilde{M}=0$ là $M$ thỏa mãn (TN).
\end{enumerate}
\end{proposition}

\begin{proof}
Ta vừa thấy điều kiện (TN) kéo theo $\widetilde{M}=0$.
Nếu $M$ thỏa mãn (TF), thì môđun con phân bậc $M'=\bigoplus_{k\geq n}M_k$, vốn hữu hạn sinh theo giả thiết, có tính chất rằng $M/M'$ thỏa mãn (TN); do đó $(M/M')\supertilde=0$, và tính khớp của hàm tử $M\mapsto\widetilde{M}$ \sref{II.2.5.4} kéo theo $\widetilde{M}=\widetilde{M'}$; vì vậy, để chứng minh $\widetilde{M}$ hữu hạn sinh, ta quy về trường hợp $M$ \emph{hữu hạn sinh}.
Vì vấn đề mang tính địa phương, chỉ cần chứng minh rằng $M_{(f)}$ là một $S_{(f)}$-môđun hữu hạn sinh \sref[I]{I.1.3.9}; nhưng $M^{(d)}$ là một $S^{(d)}$-môđun hữu hạn sinh \sref{II.2.1.6}[iii], và khi đó khẳng định của ta suy ra từ \sref{II.2.2.5}.

Bây giờ giả sử $M$ thỏa mãn (TF) và $\widetilde{M}=0$; khi đó, với cùng các ký hiệu như trên, ta có $\widetilde{M'}=0$, và điều kiện (TN) đối với $M'$ tương đương với điều kiện (TN) đối với $M$; do đó, để chứng minh rằng $\widetilde{M}=0$ kéo theo $M$ thỏa mãn (TN), ta lại có thể quy về trường hợp $M$ được sinh bởi một số hữu hạn phần tử thuần nhất $x_i$ ($1\leq i\leq p$); mặt khác, cho $(f_j)_{1\leq j\leq q}$ là một hệ các phần tử sinh thuần nhất của iđêan $S_+$.
Theo giả thiết, $M_{(f_j)}=0$ với mọi $j$; do đó tồn tại một số nguyên $n$ sao cho $f_j^nx_i=0$ với mọi $i$ và $j$.
Gọi $n_j$ là bậc của $f_j$, và gọi $m$ là giá trị lớn nhất của $\sum_j r_jn_j$ trên tập hữu hạn các hệ $(r_j)$ gồm những số nguyên không âm và thỏa mãn $\sum_j r_j\leq nq$; khi đó hiển nhiên
\oldpage[II]{40}
rằng nếu $k>m$ thì $S_kx_i=0$ với mọi $i$;
nếu $h$ là bậc lớn nhất của các $x_i$, ta suy ra $M_k=0$ khi $k>h+m$, và như vậy chứng minh được hoàn tất.
\end{proof}

\begin{corollary}[2.7.4]
\label{II.2.7.4}
Cho $S$ là một vành phân bậc sao cho iđêan $S_+$ hữu hạn sinh;
điều kiện cần và đủ để $X=\Proj(S)=\emp$ là tồn tại $n$ sao cho $S_k=0$ với $k\geq n$.
\end{corollary}

\begin{proof}
Thật vậy, điều kiện $X=\emp$ tương đương với $\sh{O}_X=\widetilde{S}=0$, và $S$ là một $S$-môđun cyclic.
\end{proof}

\begin{theorem}[2.7.5]
\label{II.2.7.5}
Giả sử iđêan $S_+$ được sinh bởi một số hữu hạn các phần tử thuần nhất bậc~$1$;
đặt $X=\Proj(S)$.
Khi đó, đối với mọi $\sh{O}_X$-môđun quasi-coherent $\sh{F}$, đồng cấu chính tắc $\beta:(\Gamma_\bullet(\sh{F}))\supertilde\to\sh{F}$ \sref{II.2.6.4} là một đẳng cấu.
\end{theorem}

\begin{proof}
Thật vậy, nếu $S_+$ được sinh bởi một số hữu hạn các phần tử $f_i\in S_1$, thì $X$ là hợp của các không gian con $\Spec(S_{(f_i)})$ \sref{II.2.3.6}; các không gian này quasi-compact, nên $X$ quasi-compact;
hơn nữa, $X$ là một lược đồ \sref{II.2.4.2};
theo \sref[I]{I.9.3.2}, \sref{II.2.5.14.2} và \sref{II.2.6.3}, đối với mọi $f\in S_d$ ($d>0$), ta có một đẳng cấu chính tắc $(\Gamma_\bullet(\sh{F}))_{(\alpha_d(f))}\simto\Gamma(D_+(f),\sh{F})$;
mặt khác, theo định nghĩa, $(\Gamma_\bullet(\sh{F}))_{(\alpha_d(f))}$ (trong đó $\Gamma_\bullet(\sh{F})$ được xét như một $\Gamma_\bullet(\sh{O}_X)$-môđun) chính là $(\Gamma_\bullet(\sh{F}))_{(f)}$ (trong đó $\Gamma_\bullet(\sh{F})$ được xét như một $S$-môđun);
đối chiếu với định nghĩa \sref[I]{I.9.3.1} của đẳng cấu chính tắc vừa nêu, ta thấy nó trùng với đồng cấu $\beta_f$, từ đó suy ra định lý.
\end{proof}

\begin{remark}[2.7.6]
\label{II.2.7.6}
Nếu giả sử vành phân bậc $S$ là \emph{Noether}, thì điều kiện của \sref{II.2.7.5} được thỏa mãn \emph{ipso facto} ngay khi giả sử iđêan $S_+$ được sinh bởi tập $S_1$ các phần tử thuần nhất bậc~$1$.
\end{remark}

\begin{corollary}[2.7.7]
\label{II.2.7.7}
Dưới các giả thiết của \sref{II.2.7.5}, mọi $\sh{O}_X$-môđun quasi-coherent $\sh{F}$ đều đẳng cấu với một $\sh{O}_X$-môđun dạng $\widetilde{M}$, trong đó $M$ là một $S$-môđun phân bậc.
\end{corollary}

\begin{corollary}[2.7.8]
\label{II.2.7.8}
Dưới các giả thiết của \sref{II.2.7.5}, mọi $\sh{O}_X$-môđun quasi-coherent hữu hạn sinh $\sh{F}$ đều đẳng cấu với một $\sh{O}_X$-môđun dạng $\widetilde{N}$, trong đó $N$ là một $S$-môđun phân bậc hữu hạn sinh.
\end{corollary}

\begin{proof}
Ta có thể giả sử $\sh{F}=\widetilde{M}$, trong đó $M$ là một $S$-môđun phân bậc \sref{II.2.7.7}.
Cho $(f_\lambda)_{\lambda\in L}$ là một hệ sinh thuần nhất của $M$;
với mỗi tập con hữu hạn $H$ của $L$, gọi $M_H$ là môđun con phân bậc của $M$ được sinh bởi các $f_\lambda$ sao cho $\lambda\in H$;
rõ ràng $M$ là giới hạn quy nạp của các môđun con $M_H$, do đó $\sh{F}$ là giới hạn quy nạp của các $\sh{O}_X$-môđun con $\widetilde{M_H}$ \sref{II.2.5.4}.
Nhưng vì $\sh{F}$ hữu hạn sinh và không gian nền của $X$ là quasi-compact, từ \sref[0]{0.5.2.3} suy ra $\sh{F}=\widetilde{M_H}$ đối với một tập con hữu hạn $H$ nào đó của $L$.
\end{proof}

\begin{corollary}[2.7.9]
\label{II.2.7.9}
Dưới các giả thiết của \sref{II.2.7.5}, cho $\sh{F}$ là một $\sh{O}_X$-môđun quasi-coherent hữu hạn sinh.
Khi đó tồn tại một số nguyên $n_0$ sao cho, với mọi $n\geq n_0$, $\sh{F}(n)$ đẳng cấu với một môđun thương của một $\sh{O}_X$-môđun dạng $\sh{O}_X^k$ ($k>0$ phụ thuộc vào $n$), và do đó được sinh bởi một số hữu hạn các tiết diện của nó trên $X$ \sref[0]{0.5.1.1}.
\end{corollary}

\begin{proof}
Theo \sref{II.2.7.8}, ta có thể giả sử $\sh{F}=\widetilde{M}$, trong đó $M$ là môđun thương của một tổng trực tiếp hữu hạn các $S$-môđun dạng $S(m_i)$;
theo \sref{II.2.5.4}, do đó ta có thể quy về trường hợp $M=S(m)$, và khi ấy $\sh{F}(n)=(S(m+n))\supertilde=\sh{O}_X(m+n)$ \sref{II.2.5.15}.
Vì vậy, chỉ cần chứng minh
\oldpage[II]{41}
  \begin{lemma}[2.7.9.1]
  \label{II.2.7.9.1}
  Dưới các giả thiết của \sref{II.2.7.5}, với mọi $n\geq0$, tồn tại một số nguyên $k$ (phụ thuộc vào $n$) và một đồng cấu toàn ánh $\sh{O}_X^k\to\sh{O}_X(n)$.
  \end{lemma}

Thật vậy, theo \sref{II.2.7.2}, chỉ cần chứng minh rằng, với một $k$ thích hợp, tồn tại một đồng cấu (TN)-toàn ánh $u$, \emph{bậc~$0$}, từ $S$-môđun phân bậc tích $S^k$ vào $S(n)$.
Ta có $(S(n))_0=S_n$ và, theo giả thiết, $S_h=S_1^h$ với mọi $h>0$, nên $SS_n=S_n+S_{n+1}+\ldots$.
Vì $S_n$ là một $S_0$-môđun hữu hạn sinh (\sref{II.2.1.5} và \sref{II.2.1.6}[i]), hãy chọn một hệ sinh $(a_i)_{1\leq i\leq k}$ của môđun này;
đồng cấu $u$ gửi phần tử thứ $i$ của cơ sở chính tắc của $S^k$ tới $a_i$ ($1\leq i\leq k$);
vì khi đó $\Coker u$ đồng nhất với $(S(n))_{-n}+\ldots+(S(n))_{-1}$, nên $u$ đúng là đồng cấu cần tìm.
\end{proof}

\begin{corollary}[2.7.10]
\label{II.2.7.10}
Dưới các giả thiết của \sref{II.2.7.5}, cho $\sh{F}$ là một $\sh{O}_X$-môđun quasi-coherent hữu hạn sinh.
Tồn tại một số nguyên $n_0$ sao cho, với mọi $n\geq n_0$, $\sh{F}$ đẳng cấu với một môđun thương của một $\sh{O}_X$-môđun dạng $(\sh{O}_X(-n))^k$ ($k$ phụ thuộc vào $n$).
\end{corollary}

\begin{proposition}[2.7.11]
\label{II.2.7.11}
Giả sử các giả thiết của \sref{II.2.7.5} được thỏa mãn, và cho $M$ là một $S$-môđun phân bậc.
Khi đó:
\begin{enumerate}
  \item[\rm{(i)}] Đồng cấu chính tắc $\widetilde{\alpha}:\widetilde{M}\to(\Gamma_\bullet(\widetilde{M}))\supertilde$ là một đẳng cấu.
  \item[\rm{(ii)}] Cho $\sh{G}$ là một $\sh{O}_X$-môđun con quasi-coherent của $\widetilde{M}$, và cho $N$ là $S$-môđun con phân bậc của $M$ được cho bởi ảnh ngược của $\Gamma_\bullet(\sh{G})$ qua $\alpha$.
    Khi đó $\widetilde{N}=\sh{G}$ (trong đó, theo \sref{II.2.5.4}, $\widetilde{N}$ được đồng nhất với một $\sh{O}_X$-môđun con của $\widetilde{M}$).
\end{enumerate}
\end{proposition}

\begin{proof}
Vì $\beta:(\Gamma_\bullet(\widetilde{M}))\supertilde\to\widetilde{M}$ là một đẳng cấu theo \sref{II.2.7.5}, nên $\widetilde{\alpha}$ là đẳng cấu ngược theo \sref{II.2.6.5.1}, từ đó suy ra (i).
Cho $P$ là môđun con phân bậc $\alpha(M)$ của $\Gamma_\bullet(\widetilde{M})$;
vì hàm tử $M\mapsto\widetilde{M}$ là khớp \sref{II.2.5.4}, ảnh của $\widetilde{M}$ qua $\widetilde{\alpha}$ bằng $\widetilde{P}$, do đó, theo (i), $\widetilde{P}=(\Gamma_\bullet(\widetilde{M}))\supertilde$.
Đặt $Q=\Gamma_\bullet(\sh{G})\cap P$;
theo điều vừa nêu và \sref{II.2.5.4}, ta có $\widetilde{Q}=(\Gamma_\bullet(\sh{G}))\supertilde$; vì thế, theo \sref{II.2.7.5}, hạn chế của $\beta$ trên $\widetilde{Q}$ là một \emph{đẳng cấu} từ $\sh{O}_X$-môđun này lên $\sh{G}$.
Nhưng theo định nghĩa của $N$ và \sref{II.2.5.4}, hạn chế của đẳng cấu $\widetilde{\alpha}$ trên $\widetilde{N}$ là một đẳng cấu từ $\widetilde{N}$ lên $\widetilde{Q}$; kết luận suy ra từ \sref{II.2.6.5.1}.
\end{proof}

\subsection{Tính hàm tử}
\label{subsection:II.2.8}

\begin{env}[2.8.1]
\label{II.2.8.1}
Cho $S$ và $S'$ là hai vành phân bậc dương, và cho $\vphi:S'\to S$ là một đồng cấu các vành phân bậc.
Ta ký hiệu $G(\vphi)$ là tập mở của $X=\Proj(S)$ bù với $V_+(\vphi(S'_+))$; nói cách khác, đó là hợp của các $D_+(\vphi(f'))$, trong đó $f'$ chạy qua tập các phần tử thuần nhất của $S'_+$.
Do đó, hạn chế lên $G(\vphi)$ của ánh xạ liên tục ${}^a\vphi$ từ $\Spec(S)$ vào $\Spec(S')$ \sref[I]{I.1.2.1} là một ánh xạ liên tục từ $G(\vphi)$ vào $\Proj(S')$; do lạm dụng ngôn ngữ, ta vẫn ký hiệu ánh xạ này là ${}^a\vphi$.
Nếu $f'\in S'_+$ thuần nhất, ta có
\[
\label{II.2.8.1.1}
  {}^a\vphi^{-1}(D_+(f'))=D_+(\vphi(f'))
\tag{2.8.1.1}
\]
do ${}^a\vphi$ gửi $G(\vphi)$ vào $\Proj(S')$ và theo \sref[I]{I.1.2.2.2}.
Mặt khác, đồng cấu $\vphi$ xác định chính tắc (với cùng các ký hiệu) một đồng cấu các vành phân bậc $S'_{f'}\to S_f$; khi hạn chế vào các phần tử bậc~$0$, ta thu được
\oldpage[II]{42}
một đồng cấu $S'_{(f')}\to S_{(f)}$, mà ta ký hiệu là $\vphi_{(f)}$;
theo \sref[I]{I.1.6.1}, đồng cấu này tương ứng với một cấu xạ các lược đồ affine $({}^a\vphi_{(f)},\widetilde{\vphi}_{(f)}):\Spec(S_{(f)})\to\Spec(S'_{(f')})$.
Nếu ta đồng nhất chính tắc $\Spec(S_{(f)})$ với lược đồ mà $\Proj(S)$ cảm sinh trên $D_+(f)$ \sref{II.2.3.6}, thì ta xác định được một cấu xạ $\Phi_f:D_+(f)\to D_+(f')$, và ${}^a\vphi_{(f)}$ được đồng nhất với hạn chế của ${}^a\vphi$ trên $D_+(f)$.
Mặt khác, rõ ràng rằng nếu $g'$ là một phần tử thuần nhất thứ hai của $S'_+$ và $g=\vphi(g')$, thì biểu đồ
\[
  \xymatrix{
    D_+(f) \ar[r]^{\Phi_f}
    & D_+(f')
  \\D_+(fg) \ar[u] \ar[r]_{\Phi_{fg}}
    & D_+(f'g') \ar[u]
  }
\]
giao hoán, vì biểu đồ
\[
  \xymatrix{
    S'_{(f')} \ar[r]^{\vphi_{(f)}} \ar[d]_{\omega_{f'g',f'}}
    & S_{(f)} \ar[d]^{\omega_{fg,f}}
  \\S'_{(f'g')} \ar[r]_{\vphi_{(fg)}}
    & S_{(fg)}
  }
\]
giao hoán.
Từ định nghĩa của $G(\vphi)$ và \sref{II.2.3.3.2}, ta suy ra:
\end{env}

\begin{proposition}[2.8.2]
\label{II.2.8.2}
Cho một đồng cấu các vành phân bậc $\vphi:S'\to S$; tồn tại duy nhất một cấu xạ $({}^a\vphi,\widetilde{\vphi})$ từ tiền lược đồ cảm sinh $G(\vphi)$ vào $\Proj(S')$ (được gọi là \emph{liên kết với $\vphi$} và ký hiệu là $\Proj(\vphi)$), sao cho, với mọi phần tử thuần nhất $f'\in S'_+$, hạn chế của cấu xạ này trên $D_+(\vphi(f'))$ trùng với cấu xạ liên kết với đồng cấu $S'_{(f')}\to S_{(\vphi(f'))}$ tương ứng với $\vphi$.
\end{proposition}

\begin{proof}
Ta nhận thấy rằng, với các ký hiệu trên, nếu $f'\in S'_d$, thì biểu đồ
\[
\label{II.2.8.2.1}
  \xymatrix{
    S'_{(f')} \ar[r]^{\vphi_{(f)}} \ar[d]_{\sim}
    & S_{(f)} \ar[d]^{\sim}
  \\{S'}^{(d)}/(f'-1){S'}^{(d)} \ar[r]
    & S^{(d)}/(f-1)S^{(d)}
  }
\tag{2.8.2.1}
\]
giao hoán (các mũi tên dọc là những đẳng cấu của \sref{II.2.2.5}).
\end{proof}

\begin{corollary}[2.8.3]
\label{II.2.8.3}
\begin{enumerate}
  \item[\rm{(i)}] Cấu xạ $\Proj(\vphi)$ là affine.
  \item[\rm{(ii)}] Nếu $\Ker(\vphi)$ lũy linh (và đặc biệt, nếu $\vphi$ đơn ánh), thì cấu xạ $\Proj(\vphi)$ là trội.
\end{enumerate}
\end{corollary}

\begin{proof}
Khẳng định~(i) là hệ quả trực tiếp của \sref{II.2.8.2} và \sref{II.2.8.1.1}.
Mặt khác, nếu $\Ker(\vphi)$ lũy linh, thì với mọi $f'$ thuần nhất trong $S'_+$, ta thấy ngay $\Ker(\vphi_f)$ (với $f=\vphi(f')$) lũy linh, và do đó $\Ker(\vphi_{(f)})$ cũng lũy linh;
kết luận khi ấy suy ra từ \sref{II.2.8.2} và \sref[I]{I.1.2.7}.
\end{proof}
\oldpage[II]{43}
Ta lưu ý rằng, nói chung, có những cấu xạ từ $\Proj(S)$ vào $\Proj(S')$ không affine và do đó không phát sinh từ các đồng cấu vành phân bậc $S'\to S$;
một ví dụ là cấu xạ cấu trúc $\Proj(S)\to\Spec(A)$ khi $A$ là một trường (trong đó $\Spec(A)$ được đồng nhất với $\Proj(A[T])$, x.~\sref{II.3.1.7});
thật vậy, điều này suy ra từ \sref[I]{I.2.3.2}.

\begin{env}[2.8.4]
\label{II.2.8.4}
Cho $S''$ là một vành thứ ba phân bậc dương, cho $\vphi':S''\to S'$ là một đồng cấu các vành phân bậc, và đặt $\vphi''=\vphi\circ\vphi'$.
Theo \sref{II.2.8.1.1} và công thức ${}^a\vphi''=({}^a\vphi')\circ({}^a\vphi)$, ta thấy ngay $G(\vphi'')\subset G(\vphi)$ và, nếu $\Phi$, $\Phi'$ và $\Phi''$ lần lượt là các cấu xạ liên kết với $\vphi$, $\vphi'$ và $\vphi''$, thì $\Phi''=\Phi'\circ(\Phi|G(\vphi''))$.
\end{env}

\begin{env}[2.8.5]
\label{II.2.8.5}
Giả sử $S$ (tương ứng, $S'$) là một $A$-đại số phân bậc (tương ứng, một $A'$-đại số phân bậc), và cho $\psi:A'\to A$ là một đồng cấu vành sao cho biểu đồ
\[
\label{II.2.8.5.1}
  \xymatrix{
    A' \ar[r]^{\psi} \ar[d]
    & A \ar[d]
  \\S' \ar[r]_{\vphi}
    & S
  }
\tag{2.8.5.1}
\]
giao hoán.
Khi đó có thể xét $G(\vphi)$ và $\Proj(S')$ lần lượt như các lược đồ trên $\Spec(A)$ và $\Spec(A')$;
nếu $\Phi$ và $\Psi$ lần lượt là các cấu xạ liên kết với $\vphi$ và $\psi$, thì biểu đồ
\[
  \xymatrix{
    G(\vphi) \ar[r]^{\Phi} \ar[d]
    & \Proj(S') \ar[d]
  \\\Spec(A) \ar[r]_{\Psi}
    & \Spec(A')
  }
\]
giao hoán:
thật vậy, chỉ cần chứng minh điều này cho hạn chế của $\Phi$ trên $D_+(f)$, trong đó $f=\vphi(f')$ và $f'$ thuần nhất trong $S'_+$;
khi ấy, điều này suy ra từ tính giao hoán của biểu đồ
\[
  \xymatrix{
    A' \ar[r]^{\psi} \ar[d]
    & A \ar[d]
  \\S'_{(f')} \ar[r]_{\vphi_{(f)}}
    & S_{(f)}
  }
\]
\end{env}

\begin{env}[2.8.6]
\label{II.2.8.6}
Bây giờ cho $M$ là một $S$-môđun phân bậc, và xét $S'$-môđun $M_{[\vphi]}$, hiển nhiên cũng phân bậc.
Cho $f'$ thuần nhất trong $S'_+$, và đặt $f=\vphi(f')$;
ta biết \sref[0]{0.1.5.2} rằng có một đẳng cấu chính tắc $(M_{[\vphi]})_{f'}\simto(M_f)_{[\vphi_f]}$, và rõ ràng đẳng cấu này bảo toàn bậc, do đó có một đẳng cấu chính tắc $(M_{[\vphi]})_{(f')}\simto(M_{(f)})_{[\vphi_{(f)}]}$.
Đẳng cấu này tương ứng chính tắc với một đẳng cấu các bó $(M_{[\vphi]})\supertilde|D_+(f')\simto(\Phi_f)_*(\widetilde{M}|D_+(f))$ (\sref{II.2.5.2} và \sref[I]{I.1.6.3}).
Hơn nữa,
\oldpage[II]{44}
nếu $g'$ là một phần tử thuần nhất thứ hai của $S'_+$ và $g=\vphi(g')$, thì biểu đồ
\[
  \xymatrix{
    (M_{[\vphi]})_{(f')} \ar[r]^{\sim} \ar[d]
    & (M_{(f)})_{[\vphi_{(f)}]} \ar[d]
  \\(M_{[\vphi]})_{(f'g')} \ar[r]^{\sim}
    & (M_{(fg)})_{[\vphi_{(fg)}]}
  }
\]
giao hoán; từ đó ta suy ra ngay rằng đẳng cấu
\[
  (M_{[\vphi]})\supertilde|D_+(f'g')\simto(\Phi_{fg})_*(\widetilde{M}|D_+(fg))
\]
là hạn chế trên $D_+(f'g')$ của đẳng cấu $(M_{[\vphi]})\supertilde|D_+(f')\simto(\Phi_f)_*(\widetilde{M}|D_+(f))$.
Vì $\Phi_f$ là hạn chế trên $D_+(f)$ của cấu xạ $\Phi$, nên xét đến \sref{II.2.8.1.1} và đặt $X'=\Proj(S')$, ta thấy rằng:
\end{env}

\begin{proposition}[2.8.7]
\label{II.2.8.7}
Tồn tại một đẳng cấu chính tắc có tính hàm tử từ $\sh{O}_{X'}$-môđun $(M_{[\vphi]})\supertilde$ lên $\sh{O}_{X'}$-môđun $\Phi_*(\widetilde{M}|G(\vphi))$.
\end{proposition}

Từ đó ta suy ra ngay một ánh xạ chính tắc có tính hàm tử từ tập gồm các $\vphi$-đồng cấu $M'\to M$, đi từ một $S'$-môđun phân bậc vào $S$-môđun phân bậc $M$, đến tập các $\Phi$-đồng cấu $\widetilde{M'}\to\widetilde{M}|G(\vphi)$.
Với các ký hiệu của \sref{II.2.8.4}, nếu $M''$ là một $S''$-môđun phân bậc, thì hợp thành của một $\vphi$-đồng cấu $M'\to M$ và một $\vphi'$-đồng cấu $M''\to M'$ tương ứng chính tắc với hợp thành của $\widetilde{M'}|G(\vphi')\to\widetilde{M}|G(\vphi'')$ và $\widetilde{M''}\to\widetilde{M'}|G(\vphi')$.

\begin{proposition}[2.8.8]
\label{II.2.8.8}
Dưới các giả thiết của \sref{II.2.8.1}, cho $M'$ là một $S'$-môđun phân bậc.
Tồn tại một đồng cấu chính tắc có tính hàm tử $\nu$ từ $(\sh{O}_X|G(\vphi))$-môđun $\Phi^*(\widetilde{M'})$ vào $(\sh{O}_X|G(\vphi))$-môđun $(M'\otimes_{S'}S)\supertilde|G(\vphi)$.
Nếu iđêan $S'_+$ được sinh bởi $S'_1$, thì $\nu$ là một đẳng cấu.
\end{proposition}

\begin{proof}
Thật vậy, với $f'\in S'_d$ ($d>0$), ta xác định một đồng cấu chính tắc có tính hàm tử của các $S_{(f)}$-môđun (trong đó $f=\vphi(f')$)
\[
\label{II.2.8.8.1}
  \nu_f:M'_{(f')}\otimes_{S'_{(f')}}S_{(f)}\to(M'\otimes_{S'}S)_{(f)}
\tag{2.8.8.1}
\]
bằng cách hợp thành đồng cấu $M'_{(f')}\otimes_{S'_{(f')}}S_{(f)}\to M'_{f'}\otimes_{S'_{f'}}S_f$ với đẳng cấu chính tắc $M'_{f'}\otimes_{S'_{f'}}S_f\simto(M'\otimes_{S'}S)_f$ \sref[0]{0.1.5.4}, và lưu ý rằng đẳng cấu sau bảo toàn bậc.
Ta thấy ngay các $\nu_f$ tương thích với các toán tử hạn chế từ $D_+(f)$ xuống $D_+(fg)$ (với một phần tử thứ hai $g'\in S'_+$ và $g=\vphi(g')$), từ đó xác định được đồng cấu
\[
  \nu:\Phi^*(\widetilde{M'})\to(M'\otimes_{S'}S)\supertilde|G(\vphi)
\]
khi xét đến \sref[I]{I.1.6.5}.
Để chứng minh khẳng định thứ hai, chỉ cần chứng minh $\nu_f$ là một đẳng cấu với mọi $f'\in S'_1$, vì khi ấy $G(\vphi)$ là hợp của các $D_+(\vphi(f'))$.
Trước hết, ta xác định một ánh xạ $\bb{Z}$-song tuyến tính $M'_m\times S_n\to M'_{(f')}\otimes_{S'_{(f')}}S_{(f)}$ bằng cách gửi $(x',s)$ tới phần tử $(x'/{f'}^m)\otimes(s/f^n)$ (với quy ước rằng $x'/{f'}^m$ là ${f'}^{-m}x'/1$
\oldpage[II]{45}
khi $m<0$).
Ta nhận thấy, như trong chứng minh của \sref{II.2.5.13}, rằng ánh xạ này sinh ra một di-đồng cấu môđun
\[
  \eta_f: M'\otimes_{S'}S \to M'_{(f')}\otimes_{S'_{(f')}}S_{(f)}.
\]
Hơn nữa, nếu với $r>0$ ta có $f^r\sum_i(x'_i\otimes s_i)=0$, thì điều này cũng được viết là $\sum_i({f'}^rx'_i\otimes s_i)=0$, từ đó, theo \sref[0]{0.1.5.4}, $\sum_i({f'}^rx'_i/{f'}^{m_i+r})\otimes(s_i/f^{n_i})=0$, tức là $\eta_f(\sum_i x'_i\otimes s_i)=0$; điều này chứng tỏ $\eta_f$ phân tích thành $M'\otimes_{S'}S\to(M'\otimes_{S'}S)_f\xrightarrow{\eta'_f}M'_{(f')}\otimes_{S'_{(f')}}S_{(f)}$;
cuối cùng, ta kiểm tra rằng $\eta'_f$ và $\nu_f$ là hai đẳng cấu ngược nhau.

Đặc biệt, từ \sref{II.2.1.2.1} suy ra một đồng cấu chính tắc
\[
\label{II.2.8.8.2}
  \Phi^*(\sh{O}_{X'}(n)) \simto \sh{O}_X(n)|G(\vphi)
\tag{2.8.8.2}
\]
với mọi $n\in\bb{Z}$.
\end{proof}

\begin{env}[2.8.9]
\label{II.2.8.9}
Cho $A$ và $A'$ là hai vành, và cho $\psi:A'\to A$ là một đồng cấu vành xác định một cấu xạ $\Psi:\Spec(A)\to\Spec(A')$.
Cho $S'$ là một $A'$-đại số phân bậc dương, và đặt $S=S'\otimes_{A'}A$, hiển nhiên là một $A$-đại số được phân bậc bởi các $S'_n\otimes_{A'}A$;
khi đó ánh xạ $\vphi:s'\mapsto s'\otimes1$ là một đồng cấu các vành phân bậc làm cho sơ đồ \sref{II.2.8.5.1} giao hoán.
Vì ở đây $S_+$ là $A$-môđun được sinh bởi $\vphi(S'_+)$, ta có $G(\vphi)=\Proj(S)=X$;
do đó, đặt $X'=\Proj(S')$, ta được sơ đồ giao hoán
\[
\label{II.2.8.9.1}
  \xymatrix{
    X \ar[r]^{\Phi} \ar[d]_p
    & X' \ar[d]
  \\Y \ar[r]_{\Psi}
    & Y'
  }
\tag{2.8.9.1}
\]

Bây giờ cho $M'$ là một $S'$-môđun phân bậc, và đặt $M=M'\otimes_{A'}A=M'\otimes_{S'}S$.
Trong các điều kiện này:
\end{env}

\begin{proposition}[2.8.10]
\label{II.2.8.10}
Sơ đồ \sref{II.2.8.9.1} đồng nhất lược đồ $X$ với tích $X'\times_{Y'}Y$;
hơn nữa, đồng cấu chính tắc $\nu:\Phi^*(\widetilde{M'})\to\widetilde{M}$ \sref{II.2.8.8} là một đẳng cấu.
\end{proposition}

\begin{proof}
Khẳng định thứ nhất sẽ được chứng minh nếu ta chỉ ra rằng, với mọi $f'$ thuần nhất trong $S'_+$, đặt $f=\vphi(f')$, các hạn chế của $\Phi$ và $p$ trên $D_+(f)$ đồng nhất lược đồ này với tích $D_+(f')\times_{Y'}Y$ \sref[I]{I.3.2.6.2};
nói cách khác, theo \sref[I]{I.3.2.2}, chỉ cần chứng minh rằng $S_{(f)}$ được đồng nhất chính tắc với $S'_{(f')}\otimes_{A'}A$, điều này suy ra ngay từ sự tồn tại của đẳng cấu chính tắc bảo toàn bậc $S_f\simto S'_{f'}\otimes_{A'}A$ \sref[0]{0.1.5.4}.
Khẳng định thứ hai suy ra từ việc, theo điều trên, $M'_{(f')}\otimes_{S'_{(f')}}S_{(f)}$ được đồng nhất với $M'_{(f')}\otimes_{A'}A$, và môđun sau lại được đồng nhất với $M_{(f)}$, bởi vì $M_f$ được đồng nhất chính tắc với $M'_{f'}\otimes_{A'}A$ qua một đẳng cấu bảo toàn bậc.
\end{proof}

\begin{corollary}[2.8.11]
\label{II.2.8.11}
Với mọi số nguyên $n\in\bb{Z}$, $\widetilde{M}(n)$ được đồng nhất với $\Phi^*(\widetilde{M'}(n))=\widetilde{M'}(n)\otimes_{Y'}\sh{O}_Y$;
đặc biệt, $\sh{O}_X(n)=\Phi^*(\sh{O}_{X'}(n))=\sh{O}_{X'}(n)\otimes_{Y'}\sh{O}_Y$.
\end{corollary}

\begin{proof}
Điều này suy ra từ \sref{II.2.8.10} và \sref{II.2.5.15}.
\end{proof}
\oldpage[II]{46}

\begin{env}[2.8.12]
\label{II.2.8.12}
Dưới các giả thiết của \sref{II.2.8.9}, với $f'\in S'_d$ ($d>0$) và $f=\vphi(f')$, sơ đồ
\[
\label{II.2.8.12.1}
  \xymatrix{
    M'_{(f')} \ar[r]^-{\sim} \ar[d]
    & {M'}^{(d)}/(f'-1){M'}^{(d)} \ar[d]
  \\M_{(f)} \ar[r]^-{\sim}
    & M^{(d)}/(f-1)M^{(d)}
  }
  \tag{2.8.12.1}
\]
(xem \sref{II.2.2.5}) là giao hoán.
\end{env}

\begin{env}[2.8.13]
\label{II.2.8.13}
Giữ các ký hiệu và giả thiết của \sref{II.2.8.9}, và cho $\sh{F}'$ là một $\sh{O}_{X'}$-môđun;
nếu đặt $\sh{F}=\Phi^*(\sh{F}')$, thì với mọi $n\in\bb{Z}$, ta có $\sh{F}(n)\simeq\Phi^*(\sh{F}'(n))$ nhờ \sref{II.2.8.11} và \sref[0]{0.4.3.3}.
Do đó, theo \sref[0]{0.3.7.1}, ta có một đồng cấu chính tắc
\[
  \Gamma(\rho): \Gamma(X',\sh{F}'(n)) \to \Gamma(X,\sh{F}(n))
\]
sinh ra một di-đồng cấu chính tắc của các môđun phân bậc
\[
  \Gamma_\bullet(\sh{F}') \to \Gamma_\bullet(\sh{F}).
\]

Giả sử iđêan $S_+$ được sinh bởi $S_1$ và $\sh{F}'=\widetilde{M'}$, do đó $\sh{F}=\widetilde{M}$ với $M=M'\otimes_{A'}A$.
Nếu $f'$ thuần nhất trong $S'_+$ và $f=\vphi(f')$, ta đã thấy rằng $M_{(f)}=M'_{(f')}\otimes_{A'}A$, và vì vậy sơ đồ
\[
  \xymatrix{
    M'_0 \ar[r] \ar[d]
    & M'_{(f')} \ar[d]
    & =\Gamma(D_+(f'),\widetilde{M'})
  \\M_0 \ar[r]
    & M_{(f)}
    & =\Gamma(D_+(f),\widetilde{M})
  }
\]
là giao hoán;
từ nhận xét này và định nghĩa của đồng cấu $\alpha$ \sref{II.2.6.2}, ta suy ra ngay rằng sơ đồ
\[
\label{II.2.8.13.1}
  \xymatrix{
    M' \ar[r]^-{\alpha_{M'}} \ar[d]
    & \Gamma_\bullet(\widetilde{M'}) \ar[d]
  \\M \ar[r]_-{\alpha_M}
    & \Gamma_\bullet(\widetilde{M})
  }
\tag{2.8.13.1}
\]
là giao hoán.
Tương tự, sơ đồ
\[
\label{II.2.8.13.2}
  \xymatrix{
    (\Gamma_\bullet(\sh{F}'))\supertilde \ar[r]^-{\beta_{\sh{F}'}} \ar[d]
    & \sh{F}' \ar[d]
  \\(\Gamma_\bullet(\sh{F}))\supertilde \ar[r]_-{\beta_{\sh{F}}}
    & \sh{F}
  }
\tag{2.8.13.2}
\]
là giao hoán (mũi tên dọc bên phải là $\Phi$-cấu xạ chính tắc $\sh{F}'\to\Phi^*(\sh{F}')=\sh{F}$).
\end{env}
\oldpage[II]{47}

\begin{env}[2.8.14]
\label{II.2.8.14}
Vẫn giữ các ký hiệu và giả thiết của \sref{II.2.8.9}, cho $N'$ là một $S'$-môđun phân bậc khác, và đặt $N=N'\otimes_{A'}A$.
Hiển nhiên các di-đồng cấu chính tắc $M'\to M$ và $N'\to N$ sinh ra một di-đồng cấu $M'\otimes_{S'}N'\to M\otimes_S N$ (tương đối với đồng cấu vành chính tắc $S'\to S$), và do đó cũng sinh ra một $S$-đồng cấu bậc~$0$, $(M'\otimes_{S'}N')\otimes_{A'}A\to M\otimes_S N$, ứng với nó (có xét đến \sref{II.2.8.10}) là một đồng cấu của các $\sh{O}_X$-môđun
\[
\label{II.2.8.14.1}
  \Phi^*((M'\otimes_{S'}N')\supertilde) \to (M\otimes_S N)\supertilde.
\tag{2.8.14.1}
\]

Hơn nữa, ta kiểm tra ngay rằng sơ đồ
\[
\label{II.2.8.14.2}
  \xymatrix{
    \Phi^*(\widetilde{M'}\otimes_{\sh{O}_{X'}}\widetilde{N'}) \ar[r]^\sim \ar[d]_{\Phi^*(\lambda)}
    & \widetilde{M}\otimes_{\sh{O}_X}\widetilde{N} \ar[d]^\lambda
    & =\Phi^*(\widetilde{M'})\otimes_{\sh{O}_X}\Phi^*(\widetilde{N'})
  \\\Phi^*((M'\otimes_{S'}N')\supertilde) \ar[r]
    & (M\otimes_S N)\supertilde
  }
\tag{2.8.14.2}
\]
là giao hoán, trong đó dòng thứ nhất là đẳng cấu chính tắc \sref[0]{0.4.3.3}.
Nếu iđêan $S'_+$ được sinh bởi $S'_1$, thì hiển nhiên $S_+$ được sinh bởi $S_1$, và khi đó hai mũi tên dọc của \sref{II.2.8.14.2} là các đẳng cấu \sref{II.2.5.13};
vì vậy, \sref{II.2.8.14.1} cũng là một đẳng cấu.

Tương tự, ta có một di-đồng cấu chính tắc $\Hom_{S'}(M',N')\to\Hom_S(M,N)$ biến một đồng cấu $u'$ bậc~$k$ thành đồng cấu $u'\otimes1$, cũng có bậc~$k$;
từ đó, ta lại suy ra một $S$-đồng cấu bậc~$0$
\[
  (\Hom_{S'}(M',N'))\otimes_{A'}A \to \Hom_S(M,N)
\]
và do đó một đồng cấu của các $\sh{O}_X$-môđun
\[
\label{II.2.8.14.3}
  \Phi^*((\Hom_{S'}(M',N'))\supertilde) \to (\Hom_S(M,N))\supertilde.
\tag{2.8.14.3}
\]

Hơn nữa, sơ đồ
\[
  \xymatrix{
    \Phi^*((\Hom_{S'}(M',N'))\supertilde) \ar[r] \ar[d]_{\Phi^*(\mu)}
    & (\Hom_S(M,N))\supertilde \ar[d]^\mu
  \\\Phi^*(\shHom_{\sh{O}_{X'}}(\widetilde{M'},\widetilde{N'})) \ar[r]
    & \shHom_{\sh{O}_X}(\widetilde{M},\widetilde{N})
  }
\]
là giao hoán (dòng ngang thứ hai là đồng cấu chính tắc \sref[0]{0.4.4.6}).
\end{env}
\oldpage[II]{48}

\begin{env}[2.8.15]
\label{II.2.8.15}
Với các ký hiệu và giả thiết của \sref{II.2.8.1}, từ \sref{II.2.4.7} suy ra rằng, sai khác các đẳng cấu, cấu xạ $\Phi$ không thay đổi khi thay $S_0$ và $S'_0$ bằng $\bb{Z}$ và thay $\vphi_0$ bằng ánh xạ đồng nhất, rồi thay lần lượt $S$ và $S'$ bằng $S^{(d)}$ và ${S'}^{(d)}$ ($d>0$), và thay $\vphi$ bằng hạn chế $\vphi^{(d)}$ của nó trên $S^{(d)}$.
\end{env}

\subsection{Các lược đồ con đóng của một lược đồ $\operatorname{Proj}(S)$}
\label{subsection:II.2.9}

\begin{env}[2.9.1]
\label{II.2.9.1}
Nếu $\vphi:S\to S'$ là một đồng cấu các vành phân bậc, ta nói $\vphi$ là (TN)-\emph{toàn ánh} (tương ứng, (TN)-\emph{đơn ánh}, (TN)-\emph{song ánh}) nếu tồn tại một số nguyên $n$ sao cho, với $k\geq n$, $\vphi_k:S_k\to S'_k$ là \emph{toàn ánh} (tương ứng, \emph{đơn ánh}, \emph{song ánh}).
Khi đó, từ \sref{II.2.8.15} suy ra rằng việc nghiên cứu $\Phi$ có thể quy về trường hợp $\vphi$ toàn ánh (tương ứng, đơn ánh, song ánh).
Thay vì nói $\vphi$ là (TN)-song ánh, ta cũng nói rằng nó là một (TN)-\emph{đẳng cấu}.
\end{env}

\begin{proposition}[2.9.2]
\label{II.2.9.2}
Cho $S$ là một vành phân bậc dương, và cho $X=\Proj(S)$.
\begin{enumerate}
  \item[{\rm(i)}] Nếu $\vphi:S\to S'$ là một đồng cấu (TN)-toàn ánh của các vành phân bậc, thì cấu xạ tương ứng $\Phi$ \sref{II.2.8.1} được xác định trên toàn bộ $\Proj(S')$ và là một phép nhúng đóng của $\Proj(S')$ vào $X$.
    Nếu $\mathfrak{J}$ là hạt nhân của $\vphi$, thì lược đồ con đóng của $X$ liên kết với $\Phi$ được xác định bởi bó iđêan quasi-coherent $\widetilde{\mathfrak{J}}$ của $\sh{O}_X$.
  \item[{\rm(ii)}] Giả sử thêm rằng iđêan $S_+$ được sinh bởi một số hữu hạn phần tử thuần nhất bậc~$1$.
    Cho $X'$ là một lược đồ con đóng của $X$, được xác định bởi một bó iđêan quasi-coherent $\sh{J}$ của $\sh{O}_X$.
    Cho $\mathfrak{J}$ là iđêan phân bậc của $S$ bằng nghịch ảnh của $\Gamma_\bullet(\sh{J})$ qua đồng cấu chính tắc $\alpha:S\to\Gamma_\bullet(\sh{O}_X)$ \sref{II.2.6.2}, và đặt $S'=S/\mathfrak{J}$.
    Khi đó, $X'$ là lược đồ con liên kết với phép nhúng đóng $\Proj(S')\to X$ tương ứng với đồng cấu chính tắc của các vành phân bậc $S\to S'$.
\end{enumerate}
\end{proposition}

\begin{proof}
\medskip\noindent
\begin{enumerate}
  \item[{\rm(i)}] Ta có thể giả sử rằng $\vphi$ toàn ánh \sref{II.2.9.1}.
    Vì theo giả thiết $\vphi(S_+)$ sinh $S'_+$, ta có $G(\vphi)=\Proj(S')$.
    Mặt khác, khẳng định thứ hai có thể kiểm tra cục bộ trên $X$;
    do đó, cho $f$ là một phần tử thuần nhất của $S_+$, và đặt $f'=\vphi(f)$.
    Vì $\vphi$ là một đồng cấu toàn ánh của các vành phân bậc, ta kiểm tra ngay rằng $\vphi_{(f')}:S_{(f)}\to S'_{(f')}$ toàn ánh và hạt nhân của nó là $\mathfrak{J}_{(f)}$; điều này hoàn tất chứng minh của (i) \sref[I]{I.4.2.3}.
  \item[{\rm(ii)}] Theo (i), ta quy về việc kiểm tra rằng đồng cấu $\widetilde{j}:\widetilde{\mathfrak{J}}\to\sh{O}_X$ suy ra từ phép nhúng chính tắc $j:\mathfrak{J}\to S$ là một đẳng cấu từ $\widetilde{\mathfrak{J}}$ lên $\sh{J}$, điều này suy ra từ \sref{II.2.7.11}.
\end{enumerate}
\end{proof}

Ta lưu ý rằng $\mathfrak{J}$ là iđêan \emph{lớn nhất} trong số các iđêan phân bậc $\mathfrak{J}'$ của $S$ thỏa $\widetilde{j}(\widetilde{\mathfrak{J'}})=\sh{J}$, bởi vì khi trở lại các định nghĩa \sref{II.2.6.2}, ta kiểm tra ngay rằng hệ thức này kéo theo $\alpha(\mathfrak{J}')\subset\Gamma_\bullet(\sh{J})$.

\begin{corollary}[2.9.3]
\label{II.2.9.3}
Giả sử các giả thiết của \sref{II.2.9.2}[(i)] được thỏa mãn và, thêm nữa, iđêan $S_+$ được sinh bởi $S_1$;
khi đó $\Phi^*((S(n))\supertilde)$ đẳng cấu chính tắc với $(S'(n))\supertilde$ với mọi $n\in\bb{Z}$, và do đó $\Phi^*(\sh{F}(n))$ đẳng cấu chính tắc với $\Phi^*(\sh{F})(n)$ với mọi $\sh{O}_X$-môđun $\sh{F}$.
\end{corollary}

\begin{proof}
Đây là một trường hợp riêng của \sref{II.2.8.8}, có xét đến \sref{II.2.5.10.2}.
\end{proof}

\begin{corollary}[2.9.4]
\label{II.2.9.4}
Giả sử các giả thiết của \sref{II.2.9.2}[(ii)] được thỏa mãn.
Điều kiện cần và đủ để tiền lược đồ con đóng $X'$ của $X$ là nguyên là iđêan phân bậc $\mathfrak{J}$ nguyên tố trong $S$.
\end{corollary}
\oldpage[II]{49}

\begin{proof}
Vì $X'$ đẳng cấu với $\Proj(S/\mathfrak{J})$, điều kiện là đủ theo \sref{II.2.4.4}.
Để thấy điều kiện là cần, xét dãy khớp $0\to\sh{J}\to\sh{O}_X\to\sh{O}_X/\sh{J}$, dãy này cho dãy khớp
\[
  0 \to \Gamma_\bullet(\sh{J}) \to \Gamma_\bullet(\sh{O}_X) \to \Gamma_\bullet(\sh{O}_X/\sh{J}).
\]

Chỉ cần chứng minh rằng, nếu $f\in S_m$ và $g\in S_n$ sao cho ảnh của $\alpha_{n+m}(fg)$ trong $\Gamma_\bullet(\sh{O}_X/\sh{J})$ bằng không, thì một trong hai ảnh của $\alpha_m(f)$ và $\alpha_n(g)$ bằng không.
Theo định nghĩa, các ảnh này là những tiết diện của các $(\sh{O}_X/\sh{J})$-môđun khả nghịch $\sh{L}=(\sh{O}_X/\sh{J})(m)$ và $\sh{L}'=(\sh{O}_X/\sh{J})(n)$ trên lược đồ nguyên $X'$;
giả thiết kéo theo rằng tích của hai tiết diện này bằng không trong $\sh{L}\otimes\sh{L}'$ (\sref{II.2.9.3} và \sref{II.2.5.14.1}), do đó một trong hai tiết diện bằng không theo \sref[I]{I.7.4.4}.
\end{proof}

\begin{corollary}[2.9.5]
\label{II.2.9.5}
Cho $A$ là một vành, $M$ là một $A$-môđun, $S$ là một $A$-đại số phân bậc được sinh bởi tập $S_1$ các phần tử thuần nhất bậc~$1$, $u:M\to S_1$ là một đồng cấu toàn ánh của các $A$-môđun, và $\overline{u}:\bb{S}(M)\to S$ là đồng cấu (của các $A$-đại số) từ đại số đối xứng $\bb{S}(M)$ của $M$ vào $S$ mở rộng $u$.
Khi đó, cấu xạ tương ứng với $\overline{u}$ là một phép nhúng đóng của $\Proj(S)$ vào $\Proj(\bb{S}(M))$.
\end{corollary}

\begin{proof}
Thật vậy, theo giả thiết $\overline{u}$ toàn ánh, và chỉ cần áp dụng \sref{II.2.9.2}.
\end{proof}

\section{Phổ thuần nhất của một bó đại số phân bậc}
\label{section:II.3}

\subsection{Phổ thuần nhất của một $\sh{O}_Y$-đại số phân bậc quasi-coherent}
\label{subsection:II.3.1}
\label{II.3.1}

\begin{env}[3.1.1]
\label{II.3.1.1}
Cho $Y$ là một tiền lược đồ, $\sh{S}$ là một $\sh{O}_Y$-đại số phân bậc, và $\sh{M}$ là một $\sh{S}$-môđun phân bậc.
Nếu $\sh{S}$ \emph{quasi-coherent}, thì mỗi thành phần thuần nhất $\sh{S}_n$ của nó là một $\sh{O}_Y$-môđun \emph{quasi-coherent}, vì đó là ảnh của $\sh{S}$ qua một đồng cấu từ $\sh{S}$ vào chính nó (\sref[I]{I.1.3.8} và \sref[I]{I.1.3.9});
tương tự, nếu $\sh{M}$ quasi-coherent với tư cách một $\sh{O}_Y$-môđun, thì các thành phần thuần nhất $\sh{M}_n$ của nó cũng quasi-coherent, và ngược lại.
Với một số nguyên $d>0$, ta ký hiệu bởi $\sh{S}^{(d)}$ tổng trực tiếp của các $\sh{O}_Y$-môđun $\sh{S}_{nd}$ ($n\in\bb{Z}$); tổng này quasi-coherent nếu $\sh{S}$ quasi-coherent \sref[I]{I.1.3.9};
với mọi số nguyên $k$ sao cho $0\leq k\leq d-1$, ta ký hiệu bởi $\sh{M}^{(d,k)}$ (hoặc $\sh{M}^{(d)}$ khi $k=0$) tổng trực tiếp của các $\sh{M}_{nd+k}$ ($n\in\bb{Z}$); đó là một $\sh{S}^{(d)}$-môđun phân bậc, quasi-coherent khi $\sh{S}$ và $\sh{M}$ đều quasi-coherent \sref[I]{I.9.6.1}.
Ta ký hiệu bởi $\sh{M}(n)$ $\sh{S}$-môđun phân bậc thỏa $(\sh{M}(n))_k=\sh{M}_{n+k}$ với mọi $k\in\bb{Z}$;
nếu $\sh{S}$ và $\sh{M}$ quasi-coherent, thì $\sh{M}(n)$ là một $\sh{S}$-môđun phân bậc quasi-coherent \sref[I]{I.9.6.1}.

Ta nói $\sh{M}$ là một $\sh{S}$-môđun phân bậc \emph{hữu hạn sinh} (tương ứng, \emph{có một trình bày hữu hạn}) nếu, với mọi $y\in Y$, tồn tại một lân cận mở $U$ của $y$ và các số nguyên $n_i$ (tương ứng, các số nguyên $m_i$ và $n_j$) sao cho có một đồng cấu toàn ánh bậc~$0$, $\bigoplus_{i=1}^r(\sh{S}(n_i)|U)\to\sh{M}|U$ (tương ứng, sao cho $\sh{M}|U$ đẳng cấu với đối hạt nhân của một đồng cấu bậc~$0$, $\bigoplus_{i=1}^r(\sh{S}(m_i)|U)\to\bigoplus_{j=1}^s(\sh{S}(n_j)|U)$).

Cho $U$ là một tập mở affine của $Y$, và cho $A=\Gamma(U,\sh{O}_Y)$ là vành của nó;
theo giả thiết, $\sh{O}_Y|U$-đại số phân bậc $\sh{S}|U$ đẳng cấu với $\widetilde{S}$, trong đó $S=\Gamma(U,\sh{S})$ là một $A$-đại số
\oldpage[II]{50}
phân bậc \sref[I]{I.1.4.3};
đặt $X_U=\Proj(\Gamma(U,\sh{S}))$.
Cho $U'\subset U$ là một tập mở affine thứ hai của $Y$, có vành $A'$, và $j:U'\to U$ là phép nhúng chính tắc, tương ứng với đồng cấu hạn chế $A\to A'$;
ta có $\sh{S}|U'=j^*(\sh{S}|U)$, và do đó $S'=\Gamma(U',\sh{S})$ được đồng nhất chính tắc với $S\otimes_AA'$ \sref[I]{I.1.6.4}.
Từ \sref{II.2.8.10} suy ra $X_{U'}$ được đồng nhất chính tắc với $X_U\times_U U'$, và do đó cũng với $f_U^{-1}(U')$, trong đó $f_U$ ký hiệu cấu xạ cấu trúc $X_U\to U$ \sref[I]{I.4.4.1}.
Ta ký hiệu bởi $\sigma_{U',U}$ đẳng cấu chính tắc $f_U^{-1}(U')\simto X_{U'}$ vừa được xác định, và bởi $\rho_{U',U}$ phép nhúng mở $X_{U'}\to X_U$ thu được bằng cách hợp $\sigma_{U',U}^{-1}$ với phép nhúng chính tắc $f_U^{-1}(U')\to X_U$.
Hiển nhiên, nếu $U''\subset U'$ là một tập mở affine thứ ba của $Y$, thì $\rho_{U'',U}=\rho_{U',U}\circ\rho_{U'',U'}$.
\end{env}

\begin{proposition}[3.1.2]
\label{II.3.1.2}
Cho $Y$ là một tiền lược đồ.
Với mọi $\sh{O}_Y$-đại số phân bậc quasi-coherent theo các bậc dương $\sh{S}$, tồn tại một tiền lược đồ $X$ trên $Y$, duy nhất sai khác một $Y$-đẳng cấu, có tính chất sau:
nếu $f:X\to Y$ là cấu xạ cấu trúc, thì với mọi tập mở affine $U$ của $Y$, tồn tại một \emph{đẳng cấu} $\eta_U$ từ tiền lược đồ cảm sinh $f^{-1}(U)$ lên $X_U=\Proj(\Gamma(U,\sh{S}))$ sao cho, nếu $V$ là một tập mở affine thứ hai của $Y$ chứa trong $U$, thì biểu đồ
\[
\label{II.3.1.2.1}
  \xymatrix{
    f^{-1}(V) \ar[r]^{\eta_V} \ar[d]
    & X_V \ar[d]^{\rho_{V,U}}
  \\f^{-1}(U) \ar[r]_{\eta_U}
    & X_U
  }
\tag{3.1.2.1}
\]
giao hoán.
\end{proposition}

\begin{proof}
Với hai tập mở affine $U$ và $V$ của $Y$, cho $X_{U,V}$ là tiền lược đồ do $X_U$ cảm sinh trên $f_U^{-1}(U\cap V)$;
ta sẽ xác định một $Y$-đẳng cấu $\theta_{U,V}:X_{V,U}\simto X_{U,V}$.
Để làm điều đó, xét một tập mở affine $W\subset U\cap V$:
bằng cách hợp các đẳng cấu
\[
  f_U^{-1}(W)
  \xrightarrow{\sigma_{W,U}} X_W
  \xrightarrow{\sigma_{W,V}^{-1}} f_V^{-1}(W),
\]
ta thu được một đẳng cấu $\tau_W$, và kiểm tra ngay rằng, nếu $W'\subset W$ là một tập mở affine, thì $\tau_{W'}$ là hạn chế của $\tau_W$ trên $f_U^{-1}(W')$;
do đó các $\tau_W$ chính là những hạn chế của một $Y$-đẳng cấu $\theta_{V,U}$.
Hơn nữa, nếu $U$, $V$, $W$ là ba tập mở affine của $Y$, và nếu $\theta'_{U,V}$, $\theta'_{V,W}$, $\theta'_{U,W}$ lần lượt là các hạn chế của $\theta_{U,V}$, $\theta_{V,W}$, $\theta_{U,W}$ trên các ảnh ngược của $U\cap V\cap W$ trong $X_V$, $X_W$, $X_W$, thì từ các định nghĩa trên suy ra $\theta'_{U,V}\circ\theta'_{V,W}=\theta'_{U,W}$.
Vì vậy, sự tồn tại của một $X$ thỏa các tính chất đã nêu suy ra từ \sref[I]{I.2.3.1};
tính duy nhất sai khác một $Y$-đẳng cấu là hiển nhiên khi xét đến \sref{II.3.1.2.1}.
\end{proof}

\begin{env}[3.1.3]
\label{II.3.1.3}
Ta nói tiền lược đồ $X$ được xác định trong \sref{II.3.1.2} là \emph{phổ thuần nhất} của $\sh{O}_Y$-đại số phân bậc quasi-coherent $\sh{S}$, và ký hiệu nó là $\Proj(\sh{S})$.
Hiển nhiên $\Proj(\sh{S})$ \emph{tách trên $Y$} (\sref{II.2.4.2} và \sref[I]{I.5.5.5});
nếu $\sh{S}$ là một $\sh{O}_Y$-đại số \emph{hữu hạn sinh} \sref[I]{I.9.6.2}, thì $\Proj(\sh{S})$ \emph{thuộc kiểu hữu hạn} trên $Y$ (\sref{II.2.7.1}[(ii)] và \sref[I]{I.6.3.1}).

Nếu $f$ là cấu xạ cấu trúc $X\to Y$, thì hiển nhiên, với mọi tiền lược đồ do $Y$ cảm sinh trên một tập mở $U$ của $Y$, $f^{-1}(U)$ được đồng nhất với phổ thuần nhất $\Proj(\sh{S}|U)$.
\end{env}

\begin{proposition}[3.1.4]
\label{II.3.1.4}
Cho $f\in\Gamma(Y,\sh{S}_d)$, với $d>0$.
Tồn tại một phần mở $X_f$ của không gian nền của $X=\Proj(\sh{S})$ có tính chất sau:
với mọi tập mở affine $U$ của $Y$,
\oldpage[II]{51}
ta có $X_f\cap\vphi^{-1}(U)=D_+(f|U)$ trong $\vphi^{-1}(U)$, được đồng nhất với $X_U=\Proj(\Gamma(U,\sh{S}))$, trong đó $\vphi$ ký hiệu cấu xạ cấu trúc $X\to Y$.
Hơn nữa, tiền lược đồ cảm sinh trên $X_f$ là affine trên $Y$ và đẳng cấu chính tắc với $\Spec(\sh{S}^{(d)}/(f-1)\sh{S}^{(d)})$ \sref{II.1.3.1}.
\end{proposition}

\begin{proof}
Ta có $f|U\in\Gamma(U,\sh{S}_d)=(\Gamma(U,\sh{S}))_d$.
Nếu $U$ và $U'$ là hai tập mở affine của $Y$ sao cho $U'\subset U$, thì $f|U'$ là ảnh của $f|U$ qua đồng cấu hạn chế
\[
  \Gamma(U,\sh{S}) \to \Gamma(U',\sh{S});
\]
do đó $D_+(f|U')$ bằng (với các ký hiệu của \sref{II.3.1.1}) tiền lược đồ cảm sinh trên ảnh ngược $\rho_{U',U}^{-1}(D_+(f|U))$ trong $X_{U'}$ \sref{II.2.8.1};
từ đó suy ra khẳng định thứ nhất.
Hơn nữa, tiền lược đồ do $X_U$ cảm sinh trên $D_+(f|U)$ được đồng nhất chính tắc với $\Spec((\Gamma(U,\sh{S}))_{(f|U)})$, và các đồng nhất này tương thích với các đồng cấu hạn chế \sref{II.2.8.1};
khi đó khẳng định thứ hai suy ra từ \sref{II.2.2.5} và tính giao hoán của biểu đồ \sref{II.2.8.2.1}.
\end{proof}

Ta còn nói rằng $X_f$ (xét như một tập mở trong không gian nền $X$) là tập hợp các $x\in X$ tại đó $f$ \emph{không triệt tiêu}.

\begin{corollary}[3.1.5]
\label{II.3.1.5}
Nếu $f\in\Gamma(Y,\sh{S}_d)$ và $g\in\Gamma(Y,\sh{S}_e)$, thì
\[
\label{II.3.1.5.1}
  X_{fg}=X_f\cap X_g.
\tag{3.1.5.1}
\]
\end{corollary}

\begin{proof}
Thật vậy, chỉ cần xét giao của hai vế với một tập $\vphi^{-1}(U)$, trong đó $U$ là một tập mở affine của $Y$, rồi áp dụng công thức \sref{II.2.3.3.2}.
\end{proof}

\begin{corollary}[3.1.6]
\label{II.3.1.6}
Cho $(f_\alpha)$ là một họ các tiết diện của $\sh{S}$ trên $Y$ sao cho $f_\alpha\in\Gamma(Y,\sh{S}_{d_\alpha})$;
nếu bó iđêan của $\sh{S}$ sinh bởi họ này \sref[0]{0.5.1.1} chứa mọi $\sh{S}_n$ kể từ một bậc nào đó, thì không gian nền $X$ là hợp của các $X_{f_\alpha}$.
\end{corollary}

\begin{proof}
Thật vậy, với mọi tập mở affine $U$ của $Y$, $\vphi^{-1}(U)$ là hợp của các $X_{f_\alpha}\cap\vphi^{-1}(U)$ \sref{II.2.3.14}.
\end{proof}

\begin{corollary}[3.1.7]
\label{II.3.1.7}
Cho $\sh{A}$ là một $\sh{O}_Y$-đại số quasi-coherent;
đặt
\[
  \sh{S}=\sh{A}[T]=\sh{A}\otimes_{\bb{Z}}\bb{Z}[T],
\]
trong đó $T$ là một biến bất định ($\bb{Z}$ và $\bb{Z}[T]$ được xem như các bó hằng trên $Y$).
Khi đó $X=\Proj(\sh{S})$ được đồng nhất chính tắc với $\Spec(\sh{A})$.
Đặc biệt, $\Proj(\sh{O}_Y[T])$ được đồng nhất với $Y$.
\end{corollary}

\begin{proof}
Áp dụng \sref{II.3.1.6} cho tiết diện duy nhất $f\in\Gamma(Y,\sh{S})$ bằng $T$ tại mọi điểm của $Y$, ta thấy $X_f=X$.
Hơn nữa, ở đây $d=1$, và $\sh{S}^{(1)}/(f-1)\sh{S}^{(1)}=\sh{S}/(f-1)\sh{S}$ đẳng cấu chính tắc với $\sh{A}$, từ đó suy ra hệ quả \sref{II.3.1.4}.
\end{proof}

Cho $g\in\Gamma(Y,\sh{O}_Y)$;
nếu lấy $\sh{S}=\sh{O}_Y[T]$, thì $g\in\Gamma(Y,\sh{S}_0)$;
đặt
\[
  h=gT\in\Gamma(Y,\sh{S}_1).
\]
Nếu $X=\Proj(\sh{S})$, đồng nhất chính tắc được xác định trong \sref{II.3.1.7} đồng nhất $X_h$ với tập mở $Y_g$ của $Y$ (theo nghĩa của \sref[0]{0.5.5.2}):
thật vậy, ta có thể quy về trường hợp $Y=\Spec(A)$ là affine, và khi đó mọi thứ quy về (có xét đến \sref{II.2.2.5}) sự kiện rằng vành phân thức $A_g$ được đồng nhất chính tắc với $A[T]/(gT-1)A[T]$ \sref[0]{0.1.2.3}.

\begin{proposition}[3.1.8]
\label{II.3.1.8}
Cho $\sh{S}$ là một $\sh{O}_Y$-đại số phân bậc quasi-coherent theo các bậc dương.
Khi đó:
\begin{enumerate}
  \item[(i)] Với mọi $d>0$, tồn tại một $Y$-đẳng cấu chính tắc từ $\Proj(\sh{S})$ lên $\Proj(\sh{S}^{(d)})$.
\oldpage[II]{52}
  \item[(ii)] Cho $\sh{S}'$ là $\sh{O}_Y$-đại số phân bậc bằng tổng trực tiếp của $\sh{O}_Y$ và các $\sh{S}_n$ ($n>0$);
    khi đó $\Proj(\sh{S}')$ và $\Proj(\sh{S})$ đẳng cấu chính tắc trên $Y$.
  \item[(iii)] Cho $\sh{L}$ là một $\sh{O}_Y$-môđun khả nghịch \sref[0]{0.5.4.1}, và cho $\sh{S}_{(\sh{L})}$ là $\sh{O}_Y$-đại số phân bậc bằng tổng trực tiếp của các $\sh{S}_d\otimes\sh{L}^{\otimes d}$ ($d\geq0$);
    khi đó $\Proj(\sh{S})$ và $\Proj(\sh{S}_{(\sh{L})})$ đẳng cấu chính tắc trên $Y$.
\end{enumerate}
\end{proposition}

\begin{proof}
Trong mỗi trường hợp trong ba trường hợp trên, chỉ cần xác định đẳng cấu một cách địa phương trên $Y$, vì việc kiểm tra tính tương thích với các phép hạn chế từ một tập mở xuống một tập mở nhỏ hơn là hiển nhiên.
Do đó có thể giả sử $Y$ là affine, và khi ấy (i) suy ra từ \sref{II.2.4.7}[(i)], còn (ii) suy ra từ \sref{II.2.4.8}.
Đối với (iii), nếu giả sử thêm rằng $\sh{L}$ đẳng cấu với $\sh{O}_Y$ (điều này được phép vì vấn đề có tính địa phương trên $Y$), thì đẳng cấu giữa $\Proj(\sh{S})$ và $\Proj(\sh{S}_{(\sh{L})})$ là hiển nhiên;
để xác định một đẳng cấu \emph{chính tắc}, cho $Y=\Spec(A)$ và $\sh{S}=\widetilde{S}$, trong đó $S$ là một $A$-đại số phân bậc, và cho $c$ là một phần tử sinh của $A$-môđun tự do $L$ sao cho $\sh{L}=\widetilde{L}$;
khi đó, với mọi $n>0$, $x_n\mapsto x_n\otimes c^{\otimes n}$ là một $A$-đẳng cấu từ $S_n$ lên $S_n\otimes L^{\otimes n}$, và các $A$-đẳng cấu này xác định một $A$-đẳng cấu các đại số phân bậc $\vphi_c:S\to S_{(L)}=\bigoplus_{n\geq0}S_n\otimes L^{\otimes n}$.
Bây giờ cho $f\in S_+$ thuần nhất bậc $d$;
với mọi $x\in S_{nd}$, ta có
\[
  \frac{x\otimes c^{\otimes nd}}{(f\otimes c^{\otimes d})^n}
  =
  \frac{x\otimes(\varepsilon c)^{\otimes nd}}{(f\otimes(\varepsilon c)^{\otimes d})^n}
\]
với mọi phần tử khả nghịch $\varepsilon\in A$, điều này cho thấy đẳng cấu $S_{(f)}\to(S_{(L)})_{(f\otimes c^{\otimes d})}$ suy ra từ $\vphi_c$ là \emph{độc lập} với phần tử sinh $c$ đã chọn của $L$, và hoàn tất chứng minh.
\end{proof}

\begin{env}[3.1.9]
\label{II.3.1.9}
Nhắc lại (\sref[0]{0.4.1.3} và \sref[I]{I.1.3.14}) rằng, để $\sh{O}_Y$-đại số phân bậc quasi-coherent $\sh{S}$ được \emph{sinh bởi $\sh{O}_Y$-môđun $\sh{S}_1$}, điều kiện cần và đủ là tồn tại một phủ $(U_\alpha)$ của $Y$ bởi các tập mở affine sao cho đại số phân bậc $\Gamma(U_\alpha,\sh{S})$ trên $\Gamma(U_\alpha,\sh{S}_0)$ được sinh bởi tập $\Gamma(U_\alpha,\sh{S}_1)$ các phần tử thuần nhất bậc $1$ của nó.
Khi đó, với mọi tập mở $V$ của $Y$, $\sh{S}|V$ được sinh bởi $(\sh{O}_Y|V)$-môđun $\sh{S}_1|V$.
\end{env}

\begin{proposition}[3.1.10]
\label{II.3.1.10}
Giả sử tồn tại một phủ mở affine hữu hạn $(U_i)$ của $Y$ sao cho mỗi đại số phân bậc $\Gamma(U_i,\sh{S})$ hữu hạn sinh trên $\Gamma(U_i,\sh{O}_Y)$.
Khi đó tồn tại $d>0$ sao cho $\sh{S}^{(d)}$ được sinh bởi $\sh{S}_d$, trong đó $\sh{S}_d$ là một $\sh{O}_Y$-môđun hữu hạn sinh.
\end{proposition}

\begin{proof}
Thật vậy, từ \sref{II.2.1.6}[(v)] suy ra rằng, với mỗi $i$, tồn tại một số nguyên $m_i$ sao cho $\Gamma(U_i,\sh{S}_{nm_i})=(\Gamma(U_i,\sh{S}_{m_i}))^n$ với mọi $n>0$;
chỉ cần lấy $d$ là một bội chung của các $m_i$, có xét đến \sref{II.2.1.6}[(i)].
\end{proof}

\begin{corollary}[3.1.11]
\label{II.3.1.11}
Dưới các giả thiết của \sref{II.3.1.10}, $\Proj(\sh{S})$ đẳng cấu trên $Y$ với một phổ thuần nhất $\Proj(\sh{S}')$, trong đó $\sh{S}'$ là một $\sh{O}_Y$-đại số phân bậc sinh bởi $\sh{S}'_1$, và $\sh{S}'_1$ là một $\sh{O}_Y$-môđun hữu hạn sinh.
\end{corollary}

\begin{proof}
Thật vậy, chỉ cần lấy $\sh{S}'=\sh{S}^{(d)}$, trong đó $d$ được xác định bởi tính chất của \sref{II.3.1.10}, rồi áp dụng \sref{II.3.1.8}[(i)].
\end{proof}

\begin{env}[3.1.12]
\label{II.3.1.12}
Nếu $\sh{S}$ là một $\sh{O}_Y$-đại số phân bậc quasi-coherent theo các bậc dương, thì theo \sref[I]{I.5.1.1}, căn lũy linh $\sh{N}$ của nó là một $\sh{O}_Y$-môđun quasi-coherent;
ta nói $\sh{N}_+=\sh{N}\cap\sh{S}_+$ là \emph{căn lũy linh của $\sh{S}_+$};
đây là một $\sh{S}_0$-môđun phân bậc quasi-coherent, vì có thể quy ngay về trường hợp $Y$ là affine, và khi đó khẳng định suy ra từ \sref{II.2.1.10}.
Với mọi $y\in Y$, $(\sh{N}_+)_y$ vì thế là căn lũy linh của $(\sh{S}_+)_y=(\sh{S}_y)_+$ \sref[I]{I.5.1.1}.
Ta nói $\sh{O}_Y$-đại số phân bậc $\sh{S}$ là \emph{cốt yếu rút gọn} nếu $\sh{N}_+=0$, điều này tương đương
\oldpage[II]{53}
với việc nói rằng $\sh{S}_y$ là một $\sh{O}_y$-đại số phân bậc cốt yếu rút gọn với mọi $y\in Y$.
Với mọi $\sh{O}_Y$-đại số phân bậc $\sh{S}$, $\sh{S}/\sh{N}_+$ là cốt yếu rút gọn.

Ta nói $\sh{S}$ là \emph{nguyên} nếu $\sh{S}_y$ là một miền nguyên với mọi $y\in Y$, và hơn nữa $(\sh{S}_y)_+=(\sh{S}_+)_y\neq0$ với mọi $y\in Y$.
\end{env}

\begin{proposition}[3.1.13]
\label{II.3.1.13}
Cho $\sh{S}$ là một $\sh{O}_Y$-đại số phân bậc theo các bậc dương.
Nếu $X=\Proj(\sh{S})$, thì $Y$-lược đồ $X_\red$ đẳng cấu chính tắc với $\Proj(\sh{S}/\sh{N}_+)$;
đặc biệt, nếu $\sh{S}$ cốt yếu rút gọn, thì $X$ rút gọn.
\end{proposition}

\begin{proof}
Sự kiện $X'=\Proj(\sh{S}/\sh{N}_+)$ rút gọn suy ra ngay từ \sref{II.2.4.4}[(i)], vì tính chất này có tính địa phương;
hơn nữa, với mọi tập mở affine $U\subset Y$, $(\vphi')^{-1}(U)$ bằng $(\vphi^{-1}(U))_\red$ (trong đó $\vphi$ và $\vphi'$ lần lượt ký hiệu các cấu xạ cấu trúc $X\to Y$ và $X'\to Y$);
ta kiểm tra ngay rằng các $U$-cấu xạ chính tắc $(\vphi')^{-1}(U)\to\vphi^{-1}(U)$ tương thích với các phép hạn chế và do đó xác định một phép nhúng đóng $X'\to X$ là một đồng phôi của các không gian nền;
từ đó suy ra kết luận \sref[I]{I.5.1.2}.
\end{proof}

\begin{proposition}[3.1.14]
\label{II.3.1.14}
Cho $Y$ là một tiền lược đồ nguyên và $\sh{S}$ là một $\sh{O}_Y$-đại số phân bậc quasi-coherent sao cho $\sh{S}_0=\sh{O}_Y$.
\begin{enumerate}
  \item[(i)] Nếu $\sh{S}$ là nguyên \sref{II.3.1.12}, thì $X=\Proj(\sh{S})$ là nguyên và cấu xạ cấu trúc $\vphi:X\to Y$ là trội.
  \item[(ii)] Giả sử thêm rằng $\sh{S}$ cốt yếu rút gọn.
    Khi đó, ngược lại, nếu $X$ là nguyên và $\vphi$ là trội, thì $\sh{S}$ là nguyên.
\end{enumerate}
\end{proposition}

\begin{proof}
\medskip\noindent
\begin{enumerate}
  \item[(i)] Nếu $(U_\alpha)$ là một cơ sở của $Y$ gồm các tập mở affine khác rỗng, thì chỉ cần chứng minh mệnh đề khi thay $Y$ bởi một trong các $U_\alpha$ và thay $\sh{S}$ bởi $\sh{S}|U_\alpha$:
    thật vậy, một mặt, từ đó suy ra các không gian nền $\vphi^{-1}(U_\alpha)$ là những tập mở bất khả quy (do đó khác rỗng) của $X$ sao cho $\vphi^{-1}(U_\alpha)\cap\vphi^{-1}(U_\beta)\neq\emp$ với mọi cặp chỉ số (vì $U_\alpha\cap U_\beta$ chứa một $U_\gamma$), nên $X$ bất khả quy \sref[0]{0.2.1.4};
    mặt khác, $X$ rút gọn vì đây là một tính chất địa phương, nên $X$ quả thực là nguyên và $\vphi(X)$ trù mật trong $Y$.

    Do đó, giả sử $Y=\Spec(A)$, trong đó $A$ là một miền nguyên \sref[I]{I.5.1.4}, và $\sh{S}=\widetilde{S}$, trong đó $S$ là một $A$-đại số phân bậc;
    giả thiết nói rằng, với mọi $y\in Y$, $\widetilde{S}_y=S_y$ là một miền nguyên phân bậc sao cho $(S_y)_+\neq0$.
    Chỉ cần chứng minh rằng $S$ là một \emph{miền nguyên}, vì khi đó $S_+\neq0$ và ta có thể áp dụng \sref{II.2.4.4}[(ii)].
    Thật vậy, cho $f,g$ là hai phần tử khác $0$ của $S$ và giả sử $fg=0$;
    với mọi $y\in Y$, ta có $(f/1)(g/1)=0$ trong $S_y$, nên theo giả thiết, $f/1=0$ hoặc $g/1=0$.
    Chẳng hạn, giả sử $f/1=0$ trong $S_y$;
    điều này có nghĩa là tồn tại $a\in A$ sao cho $a\notin\mathfrak{j}_y$ và $af=0$;
    khi đó, với \emph{mọi} $z\in Y$, ta có $(a/1)(f/1)=0$ trong miền nguyên $S_z$, và vì $a/1\neq0$ (do $A$ là một miền nguyên), nên $f/1=0$, kéo theo $f=0$.
  \item[(ii)] Vì vấn đề có tính địa phương trên $Y$, một lần nữa có thể giả sử $Y=\Spec(A)$, với $A$ là một miền nguyên, và $\sh{S}=\widetilde{S}$.
    Theo giả thiết, với mọi $y\in Y$, $(S_y)_+$ không chứa phần tử lũy linh khác $0$ nào, và điều tương tự đúng với $(S_0)_y=A_y$ theo giả thiết;
    do đó $S_y$ là một vành rút gọn với mọi $y\in Y$, và trước hết suy ra chính $S$ là rút gọn \sref[I]{I.5.1.1}.
    Mặt khác, giả thiết $X$ là nguyên kéo theo $S$ cốt yếu nguyên \sref{II.2.4.4}[(ii)], và mọi thứ khi đó quy về việc chứng minh rằng linh hóa tử $\mathfrak{J}$ của $S_+$ trong $A=S_0$ bằng
\oldpage[II]{54}
$0$ \sref{II.2.1.11}.
Nếu không, ta sẽ có $(S_h)_+=0$ với một $h\neq0$ nào đó trong $\mathfrak{J}$, và do đó, theo \sref{II.3.1.1}, $\vphi^{-1}(D(h))=\emp$, nên $\vphi(X)$ sẽ không trù mật trong $Y$, trái với giả thiết (vì $D(h)\neq\emp$, do $h$ không lũy linh).
\end{enumerate}
\end{proof}

\subsection{Bó trên $\operatorname{Proj}(\mathcal{S})$ liên kết với một $\mathcal{S}$-môđun phân bậc}
\label{subsection:II.3.2}

\begin{env}[3.2.1]
\label{II.3.2.1}
Cho $Y$ là một tiền lược đồ, $\sh{S}$ là một $\sh{O}_Y$-đại số phân bậc quasi-coherent theo các bậc dương, và $\sh{M}$ là một $\sh{S}$-môđun phân bậc quasi-coherent (trên $(Y,\sh{O}_Y)$ hoặc trên không gian vành hóa $(Y,\sh{S})$, hai cách này là như nhau \sref[I]{I.9.6.1}).
Với các ký hiệu của \sref{II.3.1.1}, ta ký hiệu bởi $\widetilde{\sh{M}}_U$ $\sh{O}_{X_U}$-môđun quasi-coherent $(\Gamma(U,\sh{M}))\supertilde$;
với $U'\subset U$, $\Gamma(U',\sh{M})$ được đồng nhất chính tắc với $\Gamma(U,\sh{M})\otimes_AA'$ \sref[I]{I.1.6.4};
do đó ta có $\widetilde{\sh{M}}_{U'}=\rho_{U',U}^*(\widetilde{\sh{M}}_U)$ \sref{II.2.8.11}.
\end{env}

\begin{proposition}[3.2.2]
\label{II.3.2.2}
Trên $\Proj(\sh{S})=X$ tồn tại duy nhất một $\sh{O}_X$-môđun quasi-coherent $\widetilde{\sh{M}}$ sao cho, với mọi tập mở affine $U$ của $Y$, ta có $\eta_U^*((\Gamma(U,\sh{M}))\supertilde)=\widetilde{\sh{M}}|f^{-1}(U)$ (trong đó $\eta_U$ ký hiệu đẳng cấu $f^{-1}(U)\simto\Proj(\Gamma(U,\sh{S}))$), với $f$ là cấu xạ cấu trúc $X\to Y$.
\end{proposition}

\begin{proof}
Vì $\rho_{U',U}$ được đồng nhất với cấu xạ nhúng $f^{-1}(U')\to f^{-1}(U)$ \sref{II.3.1.2.1}, mệnh đề suy ra ngay từ hệ thức $\widetilde{\sh{M}}_{U'}=\rho_{U',U}^*(\widetilde{\sh{M}}_U)$ và nguyên lý dán các bó \sref[0]{0.3.3.1}.
\end{proof}

Ta nói $\widetilde{\sh{M}}$ là $\sh{O}_X$-môđun \emph{liên kết với} $\sh{S}$-môđun phân bậc quasi-coherent $\sh{M}$.

\begin{proposition}[3.2.3]
\label{II.3.2.3}
Cho $\sh{M}$ là một $\sh{S}$-môđun phân bậc quasi-coherent, và cho $f\in\Gamma(Y,\sh{S}_d)$ ($d>0$).
Nếu $\xi_f$ là đẳng cấu chính tắc từ $X_f$ lên $Y$-tiền lược đồ $Z_f=\Spec(\sh{S}^{(d)}/(f-1)\sh{S}^{(d)})$ \sref{II.3.1.4}, thì $(\xi_f)_*(\widetilde{\sh{M}}|X_f)$ là $\sh{O}_{Z_f}$-môđun $(\sh{M}^{(d)}/(f-1)\sh{M}^{(d)})\supertilde$ \sref{II.1.4.3}.
\end{proposition}

\begin{proof}
Vì vấn đề có tính địa phương trên $Y$, ta được quy ngay về \sref{II.2.2.5}, có xét đến tính giao hoán của biểu đồ \sref{II.2.8.12.1}.
\end{proof}

\begin{proposition}[3.2.4]
\label{II.3.2.4}
Phép gán $\sh{M}\mapsto\widetilde{\sh{M}}$ là một hàm tử đồng biến, cộng tính và khớp từ phạm trù các $\sh{S}$-môđun phân bậc quasi-coherent vào phạm trù các $\sh{O}_X$-môđun quasi-coherent, và giao hoán với các giới hạn quy nạp và các tổng trực tiếp.
\end{proposition}

\begin{proof}
Vì vấn đề có tính địa phương trên $Y$, nó quy về \sref[I]{I.1.3.11}, \sref[I]{I.1.3.9} và \sref{II.2.5.4}.
\end{proof}

Đặc biệt, nếu $\sh{N}$ là một $\sh{S}$-môđun con phân bậc quasi-coherent của $\sh{M}$, thì $\widetilde{\sh{N}}$ được đồng nhất chính tắc với một $\sh{O}_X$-môđun con quasi-coherent của $\widetilde{\sh{M}}$;
đặc biệt hơn, với mọi bó iđêan phân bậc quasi-coherent $\sh{J}$ của $\sh{S}$, $\widetilde{\sh{J}}$ là một bó iđêan quasi-coherent của $\sh{O}_X$.

Nếu $\sh{M}$ là một $\sh{S}$-môđun phân bậc quasi-coherent và $\sh{I}$ là một bó iđêan quasi-coherent của $\sh{O}_Y$, thì $\sh{I}\sh{M}$ là một $\sh{S}$-môđun con phân bậc quasi-coherent của $\sh{M}$, và ta có
\[
\label{II.3.2.4.1}
  (\sh{I}\sh{M})\supertilde=\sh{I}\cdot\widetilde{\sh{M}}
\tag{3.2.4.1}
\]
(trong đó vế phải có nghĩa được xác định trong \sref[0]{0.4.3.5}).
Thật vậy, chỉ cần kiểm tra công thức này khi $Y=\Spec(A)$ là affine, $\sh{S}=\widetilde{S}$, trong đó $S$ là một $A$-đại số phân bậc, và $\sh{M}=\widetilde{M}$,
\oldpage[II]{55}
trong đó $M$ là một $S$-môđun phân bậc, và $\sh{I}=\widetilde{\mathfrak{I}}$, trong đó $\mathfrak{I}$ là một iđêan của $A$.
Với mọi phần tử thuần nhất $f$ của $S_+$, hạn chế trên $D_+(f)=\Spec(S_{(f)})$ của vế thứ nhất trong \sref{II.3.2.4.1} liên kết với $(\mathfrak{I}M)_{(f)}=\mathfrak{I}\cdot M_{(f)}$, và hạn chế của vế thứ hai cũng vậy, theo \sref[I]{I.1.3.13} và \sref[I]{I.1.6.9}.

\begin{proposition}[3.2.5]
\label{II.3.2.5}
Cho $f\in\Gamma(Y,\sh{S}_d)$ ($d>0$).
Trên tập mở $X_f$, $(\sh{O}_X|X_f)$-môđun $(\sh{S}(nd))\supertilde|X_f$ đẳng cấu chính tắc với $\sh{O}_X|X_f$ với mọi $n\in\bb{Z}$.
Đặc biệt, nếu $\sh{O}_Y$-đại số $\sh{S}$ được sinh bởi $\sh{S}_1$ \sref{II.3.1.9}, thì các $\sh{O}_X$-môđun $(\sh{S}(n))\supertilde$ đều khả nghịch với mọi $n\in\bb{Z}$.
\end{proposition}

\begin{proof}
Thật vậy, với mọi tập mở affine $U$ của $Y$, trong \sref{II.2.5.7} ta đã định nghĩa một đẳng cấu chính tắc từ $(\sh{S}(nd))\supertilde|(X_f\cap\vphi^{-1}(U))$ lên $\sh{O}_X|(X_f\cap\vphi^{-1}(U))$, có xét đến \sref{II.3.1.4} (trong đó $\vphi$ là cấu xạ cấu trúc $X\to Y$);
ta kiểm tra ngay rằng các đẳng cấu này tương thích với phép hạn chế từ $U$ xuống một tập mở affine $U'\subset U$, từ đó suy ra khẳng định thứ nhất.
Để chứng minh khẳng định thứ hai, chỉ cần lưu ý rằng, nếu $\sh{S}$ được sinh bởi $\sh{S}_1$, thì tồn tại một phủ $(U_\alpha)$ của $Y$ bởi các tập mở affine sao cho $\Gamma(U_\alpha,\sh{S})$ được sinh bởi $\Gamma(U_\alpha,\sh{S})_1=\Gamma(U_\alpha,\sh{S}_1)$;
khi đó ta quy về áp dụng kết quả của \sref{II.2.5.9}, vì tính khả nghịch là một tính chất địa phương.
\end{proof}

Ta còn đặt, với mọi $n\in\bb{Z}$,
\[
\label{II.3.2.5.1}
  \sh{O}_X(n)=(\sh{S}(n))\supertilde
\tag{3.2.5.1}
\]
và, với mọi $\sh{O}_X$-môđun $\sh{F}$,
\[
\label{II.3.2.5.2}
  \sh{F}(n)=\sh{F}\otimes_{\sh{O}_X}\sh{O}_X(n).
\tag{3.2.5.2}
\]

Từ các định nghĩa này suy ra ngay rằng, với mọi tập mở $U$ của $Y$, ta có
\[
  ((\sh{S}|U)(n))\supertilde=\sh{O}_X(n)|f^{-1}(U),
\]
trong đó $f$ là cấu xạ cấu trúc $X\to Y$.

\begin{proposition}[3.2.6]
\label{II.3.2.6}
Cho $\sh{M}$ và $\sh{N}$ là hai $\sh{S}$-môđun phân bậc quasi-coherent.
Tồn tại một đồng cấu chính tắc có tính hàm tử (theo $\sh{M}$ và $\sh{N}$)
\[
\label{II.3.2.6.1}
  \lambda:\widetilde{\sh{M}}\otimes_{\sh{O}_X}\widetilde{\sh{N}}\to(\sh{M}\otimes_{\sh{S}}\sh{N})\supertilde
\tag{3.2.6.1}
\]
và một đồng cấu chính tắc có tính hàm tử (theo $\sh{M}$ và $\sh{N}$)
\[
\label{II.3.2.6.2}
  \mu:(\shHom_{\sh{S}}(\sh{M},\sh{N}))\supertilde\to\shHom_{\sh{O}_X}(\widetilde{\sh{M}},\widetilde{\sh{N}}).
\tag{3.2.6.2}
\]

Hơn nữa, nếu $\sh{S}$ được sinh bởi $\sh{S}_1$ \sref{II.3.1.9}, thì $\lambda$ là một đẳng cấu;
nếu thêm vào đó $\sh{M}$ có một trình bày hữu hạn \sref{II.3.1.1}, thì $\mu$ là một đẳng cấu.
\end{proposition}

\begin{proof}
Các đồng cấu $\lambda$ và $\mu$ đã được định nghĩa trong \sref{II.2.5.11.2} và \sref{II.2.5.12.2} khi $Y$ affine;
vì các định nghĩa này mang tính địa phương, chúng chuyển ngay sang trường hợp tổng quát đang xét, có xét đến \sref{II.2.8.14}.
\end{proof}

\begin{corollary}[3.2.7]
\label{II.3.2.7}
Nếu $\sh{S}$ được sinh bởi $\sh{S}_1$, thì, với mọi $m,n\in\bb{Z}$, ta có
\[
\label{II.3.2.7.1}
  \sh{O}_X(m)\otimes_{\sh{O}_X}\sh{O}_X(n)=\sh{O}_X(m+n)
\tag{3.2.7.1}
\]
\[
\label{II.3.2.7.2}
  \sh{O}_X(n)=(\sh{O}_X(1))^{\otimes n}
\tag{3.2.7.2}
\]
sai khác các đẳng cấu chính tắc.
\end{corollary}
\oldpage[II]{56}

\begin{corollary}[3.2.8]
\label{II.3.2.8}
Giả sử $\sh{S}$ được sinh bởi $\sh{S}_1$.
Với mọi $\sh{S}$-môđun phân bậc $\sh{M}$ và mọi $n\in\bb{Z}$, ta có
\[
\label{II.3.2.8.1}
  (\sh{M}(n))\supertilde=\widetilde{\sh{M}}(n)
\tag{3.2.8.1}
\]
sai khác một đẳng cấu chính tắc.
\end{corollary}

\begin{proof}
Điều này suy ra từ các tính chất tương ứng khi $Y$ affine (\sref{II.2.5.14} và \sref{II.2.5.15}), cũng như từ \sref{II.2.8.11}.
\end{proof}

\begin{remarks}[3.2.9]
\label{II.3.2.9}
\medskip\noindent
\begin{enumerate}
  \item[\rm{(i)}] Nếu $\sh{S}=\sh{A}[T]$, trong đó $\sh{A}$ là một $\sh{O}_Y$-đại số quasi-coherent \sref{II.3.1.7}, thì ta kiểm tra ngay rằng mọi $\sh{O}_X$-môđun khả nghịch $\sh{O}_X(n)$ đều đẳng cấu chính tắc với $\sh{O}_X$.

  Hơn nữa, cho $\sh{N}$ là một $\sh{A}$-môđun quasi-coherent, và đặt $\sh{M}=\sh{N}\otimes_{\sh{A}}\sh{A}[T]$.
  Khi đó, từ \sref{II.3.2.3} và \sref{II.3.1.7} suy ra rằng, qua phép đồng nhất chính tắc $X=\Proj(\sh{A}[T])$ với $X'=\Spec(\sh{A})$, $\sh{O}_X$-môđun $\widetilde{\sh{M}}$ được đồng nhất với $\sh{O}_{X'}$-môđun $\widetilde{\sh{N}}$ liên kết với $\sh{N}$ (theo nghĩa của \sref{II.1.4.3}).
  \item[\rm{(ii)}] Cho $\sh{S}$ là một $\sh{O}_Y$-đại số phân bậc tùy ý, và cho $\sh{S}'$ là $\sh{O}_Y$-đại số phân bậc sao cho $\sh{S}'_0=\sh{O}_Y$ và $\sh{S}'_n=\sh{S}_n$ với mọi $n>0$;
  khi đó đẳng cấu chính tắc từ $X=\Proj(\sh{S})$ lên $X'=\Proj(\sh{S}')$ \sref{II.3.1.8}[(ii)] đồng nhất $\sh{O}_X(n)$ với $\sh{O}_{X'}(n)$ với mọi $n\in\bb{Z}$.
  Điều này suy ra từ cùng mệnh đề trong trường hợp affine \sref{II.2.5.16} và từ việc các phép đồng nhất ấy, trên các tập mở affine của $Y$, giao hoán với các phép hạn chế.
  Tương tự, cho $X^{(d)}=\Proj(\sh{S}^{(d)})$;
  khi đó đẳng cấu chính tắc từ $X$ lên $X^{(d)}$ \sref{II.3.1.8}[(i)] đồng nhất $\sh{O}_X(nd)$ với $\sh{O}_{X^{(d)}}(n)$ với mọi $n\in\bb{Z}$.
\end{enumerate}
\end{remarks}

\begin{proposition}[3.2.10]
\label{II.3.2.10}
Cho $\sh{L}$ là một $\sh{O}_Y$-môđun khả nghịch, và cho $g$ là đẳng cấu chính tắc từ $X_{(\sh{L})}=\Proj(\sh{S}_{(\sh{L})})$ lên $X=\Proj(\sh{S})$ \sref{II.3.1.8}[(iii)].
Khi đó, với mọi $n\in\bb{Z}$, $g_*(\sh{O}_{X_{(\sh{L})}}(n))$ đẳng cấu chính tắc với $\sh{O}_X(n)\otimes_Y\sh{L}^{\otimes n}$.
\end{proposition}

\begin{proof}
Trước hết, giả sử $Y$ affine, có vành $A$, và $\sh{L}=\widetilde{L}$, trong đó $L$ là một $A$-môđun tự do hạng một.
Với các ký hiệu trong chứng minh của \sref{II.3.1.8}[(iii)], ta định nghĩa, với $f\in S_d$, một đẳng cấu từ $(S(n))_{(f)}\otimes_A L^{\otimes n}$ lên $(S_{(L)}(n))_{(f\otimes c^{\otimes d})}$, gửi $(x/f^k)\otimes c^{\otimes n}$, trong đó $x\in S_{kd+n}$, tới phần tử $(x\otimes c^{\otimes(n+kd)})/(f\otimes c^{\otimes d})^k$;
ta kiểm tra ngay rằng đẳng cấu này không phụ thuộc vào phần tử sinh $c$ đã chọn của $L$;
hơn nữa, các đẳng cấu được định nghĩa như vậy với mỗi $f\in S_+$ tương thích với các phép hạn chế $D_+(f)\to D_+(fg)$.
Cuối cùng, trong trường hợp tổng quát, khi trở lại các định nghĩa của \sref{II.3.1.1}, ta thấy dễ dàng rằng các đẳng cấu được định nghĩa như vậy trên mỗi tập mở affine $U$ của $Y$ tương thích với việc chuyển từ $U$ xuống một tập mở affine $U'\subset U$.
\end{proof}

\subsection{$\mathcal{S}$-môđun phân bậc liên kết với một bó trên $\operatorname{Proj}(\mathcal{S})$}
\label{subsection:II.3.3}

\emph{Trong toàn bộ mục này, ta giả sử $\sh{O}_Y$-đại số phân bậc $\sh{S}$ được sinh bởi $\sh{S}_1$ \sref{II.3.1.9}.}
Nhắc lại rằng, theo \sref{II.3.1.8}[(i)], hạn chế này không cốt yếu khi có các điều kiện hữu hạn của \sref{II.3.1.10}.

\begin{env}[3.3.1]
\label{II.3.3.1}
Cho $p$ là cấu xạ cấu trúc $X=\Proj(\sh{S})\to Y$.
Với mọi $\sh{O}_X$-môđun $\sh{F}$, ta đặt
\[
\label{II.3.3.1.1}
  \bbGamma_*(\sh{F})=\bigoplus_{n\in\bb{Z}}p_*(\sh{F}(n))
\tag{3.3.1.1}
\]
\oldpage[II]{57}
và, đặc biệt,
\[
\label{II.3.3.1.2}
  \bbGamma_*(\sh{O}_X)=\bigoplus_{n\in\bb{Z}}p_*(\sh{O}_X(n)).
\tag{3.3.1.2}
\]

Ta biết theo \sref[0]{0.4.2.2} rằng tồn tại một đồng cấu chính tắc
\[
  p_*(\sh{F})\otimes_{\sh{O}_Y}p_*(\sh{G})\to p_*(\sh{F}\otimes_{\sh{O}_X}\sh{G})
\]
với hai $\sh{O}_X$-môđun $\sh{F}$ và $\sh{G}$;
do đó, từ \sref{II.3.2.7.1} suy ra rằng $\bbGamma_*(\sh{O}_X)$ được trang bị một cấu trúc \emph{$\sh{O}_Y$-đại số phân bậc}, và \sref{II.3.2.5.2} cũng xác định trên $\bbGamma_*(\sh{F})$ một cấu trúc \emph{môđun phân bậc trên $\bbGamma_*(\sh{O}_X)$}.

Theo \sref{II.3.2.5} và tính khớp trái của hàm tử $p_*$ \sref[0]{0.4.2.1}, tương ứng $\sh{F}\mapsto\bbGamma_*(\sh{F})$ là một hàm tử hiệp biến cộng tính và \emph{khớp trái} từ phạm trù các $\sh{O}_X$-môđun vào phạm trù các $\sh{O}_Y$-môđun phân bậc (trong đó các cấu xạ là những đồng cấu bậc~$0$).
Đặc biệt, nếu $\sh{J}$ là một bó iđêan của $\sh{O}_X$, thì $\bbGamma_*(\sh{J})$ được đồng nhất với một \emph{bó iđêan phân bậc} của $\bbGamma_*(\sh{O}_X)$.
\end{env}

\begin{env}[3.3.2]
\label{II.3.3.2}
Cho $\sh{M}$ là một $\sh{S}$-môđun phân bậc quasi-coherent.
Với mọi tập mở affine $U$ của $Y$, trong \sref{II.2.6.2} ta đã định nghĩa một đồng cấu các nhóm abel
\[
  \alpha_{0,U}:\Gamma(U,\sh{M}_0)\to\Gamma(p^{-1}(U),\widetilde{\sh{M}}).
\]
Ta kiểm tra ngay rằng các đồng cấu này giao hoán với các phép hạn chế \sref{II.2.8.13.1}, và do đó xác định (mà không dùng giả thiết $\sh{S}$ được sinh bởi $\sh{S}_1$) một đồng cấu các bó nhóm abel
\[
\label{II.3.3.2.1}
  \alpha_0:\sh{M}_0\to p_*(\widetilde{\sh{M}}).
\tag{3.3.2.1}
\]

Áp dụng kết quả này cho mỗi $\sh{M}_n=(\sh{M}(n))_0$, và có xét đến \sref{II.3.2.8.1}, ta định nghĩa một đồng cấu các bó nhóm abel
\[
\label{II.3.3.2.2}
  \alpha_n:\sh{M}_n\to p_*(\widetilde{\sh{M}}(n))
  \qquad\text{với mọi }n\in\bb{Z},
\tag{3.3.2.2}
\]
từ đó thu được một đồng cấu có tính hàm tử (bậc~$0$) của các bó nhóm abel phân bậc
\[
\label{II.3.3.2.3}
  \alpha:\sh{M}\to\bbGamma_*(\widetilde{\sh{M}})
\tag{3.3.2.3}
\]
(cũng được ký hiệu là $\alpha_{\sh{M}}$).

Đặc biệt, lấy $\sh{M}=\sh{S}$, ta kiểm tra rằng $\alpha:\sh{S}\to\bbGamma_*(\sh{O}_X)$ là một đồng cấu các $\sh{O}_Y$-đại số phân bậc, và \sref{II.3.3.2.3} là một di-đồng cấu các môđun phân bậc tương đối với đồng cấu các đại số phân bậc ấy.

Ta còn lưu ý rằng ứng với mỗi $\alpha_n$ là một đồng cấu chính tắc của các $\sh{O}_X$-môđun \sref[0]{0.4.4.3}
\[
\label{II.3.3.2.4}
  \alpha_n^\sharp:p^*(\sh{M}_n)\to\widetilde{\sh{M}}(n).
\tag{3.3.2.4}
\]

Ta kiểm tra dễ dàng rằng đồng cấu này chính là đồng cấu tương ứng có tính hàm tử \sref{II.3.2.4} với đồng cấu chính tắc (bậc~$0$) của các $\sh{O}_Y$-môđun phân bậc
\[
\label{II.3.3.2.5}
  \sh{M}_n\otimes_{\sh{O}_Y}\sh{S}\to\sh{M}(n),
\tag{3.3.2.5}
\]
\oldpage[II]{58}
trong đó phép phân bậc của vế thứ hai được cảm sinh một cách tự nhiên từ phép phân bậc của $\sh{M}$.
Thật vậy, ta có thể giới hạn vào trường hợp $Y=\Spec(A)$ affine, $\sh{M}=\widetilde{M}$ và $\sh{S}=\widetilde{S}$, trong đó $A$-đại số phân bậc $S$ được sinh bởi $S_1$, sao cho khi $f$ chạy qua $S_1$, các $D_+(f)$ tạo thành một phủ của $X$.
Trở lại các định nghĩa của \sref{II.2.6.2}, có xét đến \sref[I]{I.1.6.7}, ta thấy rằng hạn chế trên $D_+(f)$ của đồng cấu \sref{II.3.3.2.4} tương ứng theo \sref[I]{I.1.3.8} với đồng cấu các $S_{(f)}$-môđun $M_n\otimes_A S_{(f)}\to(M(n))_{(f)}$ gửi $x\otimes1$, trong đó $x\in M_n$, tới phần tử $x/1$;
điều này chứng minh khẳng định.
\end{env}

\begin{proposition}[3.3.3]
\label{II.3.3.3}
Với mọi tiết diện $f\in\Gamma(Y,\sh{S}_d)$ ($d>0$), $X_f$ chính là tập các điểm của $X$ tại đó $\alpha_d(f)$ (được xét như một tiết diện của $\sh{O}_X(d)$) không triệt tiêu \sref[0]{0.5.5.2}.
\end{proposition}

\begin{proof}
($\alpha_d(f)$ là một tiết diện của $p_*(\sh{O}_X(d))$ trên $Y$, nhưng theo định nghĩa, một tiết diện như vậy cũng là một tiết diện của $\sh{O}_X(d)$ trên $X$ \sref[0]{0.4.2.1}.)
Định nghĩa của $X_f$ \sref{II.3.1.4} quy ta về trường hợp $Y$ affine, đã được xử lý trong \sref{II.2.6.3}.
\end{proof}

\begin{env}[3.3.4]
\label{II.3.3.4}
Từ nay, ngoài giả thiết ở đầu mục này, ta giả sử rằng, với mọi $\sh{O}_X$-môđun quasi-coherent $\sh{F}$, các $p_*(\sh{F}(n))$ đều quasi-coherent trên $Y$, và do đó $\bbGamma_*(\sh{F})=\bigoplus_{n\in\bb{Z}}p_*(\sh{F}(n))$ cũng là một $\sh{O}_Y$-môđun quasi-coherent (\sref[I]{I.1.4.1} và \sref[I]{I.1.3.9});
điều này luôn xảy ra, đặc biệt, nếu $X$ thuộc kiểu hữu hạn trên $Y$ \sref[I]{I.9.2.2}.
Suy ra $(\bbGamma_*(\sh{F}))\supertilde$ được định nghĩa và là một $\sh{O}_X$-môđun quasi-coherent.
Với mọi tập mở affine $U$ của $Y$, ta có (\sref[I]{I.1.3.9} và \sref[I]{I.2.5.4})
\begin{align*}
  \Big(\Gamma(U,\bigoplus_{n\in\bb{Z}}p_*(\sh{F}(n)))\Big)\supertilde
  &=\bigoplus_{n\in\bb{Z}}\Big(\Gamma(U,p_*(\sh{F}(n)))\Big)\supertilde
\\&=\bigoplus_{n\in\bb{Z}}\Big(\Gamma(p^{-1}(U),\sh{F}(n))\Big)\supertilde
\\&=\Big(\bigoplus_{n\in\bb{Z}}\Gamma(p^{-1}(U),\sh{F}(n))\Big)\supertilde
\\&=(\bbGamma_*(\sh{F}|p^{-1}(U)))\supertilde,
\end{align*}
và do đó, theo \sref{II.2.6.4}, ta có một đồng cấu chính tắc
\[
  \beta_U:\Big(\Gamma(U,\bigoplus_{n\in\bb{Z}}p_*(\sh{F}(n)))\Big)\supertilde\to\sh{F}|p^{-1}(U).
\]

Hơn nữa, tính giao hoán của \sref{II.2.8.13.2} cho thấy các đồng cấu này giao hoán với các phép hạn chế trên $Y$;
do đó ta thu được một đồng cấu chính tắc có tính hàm tử
\[
\label{II.3.3.4.1}
  \beta:(\bbGamma_*(\sh{F}))\supertilde\to\sh{F}
\tag{3.3.4.1}
\]
(cũng được ký hiệu là $\beta_{\sh{F}}$) đối với các $\sh{O}_X$-môđun quasi-coherent.
\end{env}

\begin{proposition}[3.3.5]
\label{II.3.3.5}
Cho $\sh{M}$ là một $\sh{S}$-môđun phân bậc quasi-coherent và $\sh{F}$ là một $\sh{O}_X$-môđun quasi-coherent;
các đồng cấu hợp thành
\[
\label{II.3.3.5.1}
  \widetilde{\sh{M}}\xrightarrow{\widetilde{\alpha}}(\bbGamma_*(\widetilde{\sh{M}}))\supertilde\xrightarrow{\beta}\widetilde{\sh{M}}
\tag{3.3.5.1}
\]
\[
\label{II.3.3.5.2}
  \bbGamma_*(\sh{F})\xrightarrow{\alpha}\bbGamma_*((\bbGamma_*(\sh{F}))\supertilde)\xrightarrow{\bbGamma_*(\beta)}\bbGamma_*(\sh{F})
\tag{3.3.5.2}
\]
là các đẳng cấu đồng nhất.
\end{proposition}

\begin{proof}
Vấn đề quả thật mang tính địa phương trên $Y$, và ta quy về \sref{II.2.6.5}.
\end{proof}
\subsection{Các điều kiện hữu hạn}
\label{subsection:II.3.4}

\oldpage[II]{59}
\begin{proposition}[3.4.1]
\label{II.3.4.1}
Cho $Y$ là một tiền lược đồ và $\sh{S}$ là một $\sh{O}_Y$-đại số quasi-coherent được sinh bởi $\sh{S}_1$ \sref{II.3.1.9};
giả sử thêm rằng $\sh{S}_1$ hữu hạn sinh.
Khi đó $X=\Proj(\sh{S})$ thuộc kiểu hữu hạn trên $Y$.
\end{proposition}

\begin{proof}
Thật vậy, vì vấn đề mang tính địa phương trên $Y$, ta có thể giả sử $Y$ affine, có vành $A$;
khi đó $\sh{S}=\widetilde{S}$, trong đó $S=\Gamma(Y,\sh{S})$, và theo giả thiết, $S$ là một $A$-đại số được sinh bởi $S_1=\Gamma(Y,\sh{S}_1)$, hơn nữa ta có thể giả sử $S_1$ là một $A$-môđun hữu hạn sinh (\sref[I]{I.1.3.9} và \sref[I]{I.1.3.12}).
Do đó, $S$ là một $A$-đại số phân bậc hữu hạn sinh, và ta quy về \sref{II.2.7.1}[(ii)].
\end{proof}

\begin{env}[3.4.2]
\label{II.3.4.2}
Cho $\sh{S}$ là một $\sh{O}_Y$-đại số phân bậc quasi-coherent;
đối với một $\sh{S}$-môđun phân bậc quasi-coherent $\sh{M}$, ta xét các điều kiện hữu hạn sau đây:
\begin{enumerate}
  \item[(\textbf{TF})] Tồn tại một số nguyên $n$ sao cho $\sh{S}$-môđun $\bigoplus_{k\geq n}\sh{M}_k$ hữu hạn sinh.
  \item[(\textbf{TN})] Tồn tại một số nguyên $n$ sao cho $\sh{M}_k=0$ với $k\geq n$.
\end{enumerate}

Nếu $\sh{M}$ thỏa mãn (\textbf{TN}), thì $\widetilde{\sh{M}}=0$, vì vấn đề mang tính địa phương trên $Y$ \sref{II.2.7.2}.

Cho $\sh{M}$ và $\sh{N}$ là hai $\sh{S}$-môđun phân bậc quasi-coherent;
ta nói một đồng cấu $u:\sh{M}\to\sh{N}$ bậc~$0$ là \emph{(\textbf{TN})-đơn ánh} (tương ứng, \emph{(\textbf{TN})-toàn ánh}, \emph{(\textbf{TN})-song ánh}) nếu tồn tại một số nguyên $n$ sao cho $u_k:\sh{M}_k\to\sh{N}_k$ là đơn ánh (tương ứng, toàn ánh, song ánh) với $k\geq n$;
khi đó $\widetilde{u}:\widetilde{\sh{M}}\to\widetilde{\sh{N}}$ là đơn ánh (tương ứng, toàn ánh, song ánh) theo \sref{II.2.7.2}, vì vấn đề mang tính địa phương trên $Y$, và có xét đến \sref[I]{I.1.3.9};
khi $u$ là (\textbf{TN})-song ánh, ta còn nói $u$ là một \emph{(\textbf{TN})-đẳng cấu}.
\end{env}

\begin{proposition}[3.4.3]
\label{II.3.4.3}
Cho $Y$ là một tiền lược đồ và $\sh{S}$ là một $\sh{O}_Y$-đại số phân bậc quasi-coherent được sinh bởi $\sh{S}_1$, trong đó $\sh{S}_1$ được giả sử hữu hạn sinh.
Cho $\sh{M}$ là một $\sh{S}$-môđun phân bậc quasi-coherent.
\begin{enumerate}
  \item[(i)] Nếu $\sh{M}$ thỏa mãn điều kiện (\textbf{TF}), thì $\widetilde{\sh{M}}$ hữu hạn sinh.
  \item[(ii)] Giả sử $\sh{M}$ thỏa mãn (\textbf{TF}); điều kiện cần và đủ để $\widetilde{\sh{M}}=0$ là $\sh{M}$ thỏa mãn (\textbf{TN}).
\end{enumerate}
\end{proposition}

\begin{proof}
Vì các vấn đề mang tính địa phương trên $Y$, ta quy về trường hợp $Y$ affine, có vành $A$, $\sh{S}=\widetilde{S}$, trong đó $S$ là một $A$-đại số phân bậc sao cho iđêan $S_+$ hữu hạn sinh, và $\sh{M}=\widetilde{M}$, trong đó $M$ là một $S$-môđun phân bậc;
khi đó mệnh đề suy ra từ \sref{II.2.7.3}.
\end{proof}

\begin{theorem}[3.4.4]
\label{II.3.4.4}
Cho $Y$ là một tiền lược đồ và $\sh{S}$ là một $\sh{O}_Y$-đại số phân bậc quasi-coherent được sinh bởi $\sh{S}_1$, trong đó $\sh{S}_1$ được giả sử hữu hạn sinh;
đặt $X=\Proj(\sh{S})$.
Với mọi $\sh{O}_X$-môđun quasi-coherent $\sh{F}$, đồng cấu chính tắc $\beta$ \sref{II.3.3.4} là một đẳng cấu.
\end{theorem}

\begin{proof}
Trước hết, lưu ý rằng $\beta$ được định nghĩa theo \sref{II.3.4.1}.
Để thấy $\beta$ là một đẳng cấu, ta quy về trường hợp $Y$ affine, có vành $A$, $\sh{S}=\widetilde{S}$, trong đó $S$ là một $A$-đại số phân bậc được sinh bởi $S_1$, và $S_1$ là một $A$-môđun hữu hạn sinh.
Khi đó chỉ cần áp dụng \sref{II.2.7.5}.
\end{proof}

\begin{corollary}[3.4.5]
\label{II.3.4.5}
Dưới các giả thiết của \sref{II.3.4.4}, mọi $\sh{O}_X$-môđun quasi-coherent $\sh{F}$ đều đẳng cấu với một $\sh{O}_X$-môđun dạng $\widetilde{\sh{M}}$, trong đó $\sh{M}$ là một $\sh{S}$-môđun phân bậc quasi-coherent.
Nếu thêm vào đó $\sh{F}$ hữu hạn sinh, và nếu giả sử $Y$ là một lược đồ quasi-compact hoặc không gian nền của $Y$ là Noether, thì ta có thể giả sử $\sh{M}$ hữu hạn sinh.
\end{corollary}
\oldpage[II]{60}
\begin{proof}
Khẳng định thứ nhất suy ra ngay từ \sref{II.3.4.4}, bằng cách lấy $\sh{M}=\bbGamma_*(\sh{F})$.
Để chứng minh khẳng định thứ hai, chỉ cần chứng minh rằng $\sh{M}$ là giới hạn quy nạp của các $\sh{S}$-môđun con \emph{phân bậc} hữu hạn sinh $\sh{N}_\lambda$: thật vậy, từ đó suy ra $\widetilde{\sh{M}}$ là giới hạn quy nạp của các $\widetilde{\sh{N}}_\lambda$ \sref{II.3.2.4}, và do đó $\sh{F}$ là giới hạn quy nạp của các $\beta(\widetilde{\sh{N}}_\lambda)$;
vì $X$ quasi-compact (\sref{II.3.4.1} và \sref[I]{I.6.3.1}) và $\sh{F}$ hữu hạn sinh, nên $\sh{F}$ tất yếu bằng một trong các $\beta(\widetilde{\sh{N}}_\lambda)$ \sref[0]{0.5.2.3}.

Để xác định các $\sh{N}_\lambda$ có giới hạn quy nạp là $\sh{M}$, chỉ cần xét, với mỗi $n\in\bb{Z}$, $\sh{O}_Y$-môđun quasi-coherent $\sh{M}_n$, vốn là giới hạn quy nạp của các $\sh{O}_Y$-môđun con hữu hạn sinh $\sh{M}_n^{(\mu_n)}$ của nó, do các giả thiết trên $Y$ \sref[I]{I.9.4.9};
hiển nhiên $\sh{P}_{\mu_n}=\sh{S}\cdot\sh{M}_n^{(\mu_n)}$ là một $\sh{S}$-môđun phân bậc hữu hạn sinh, và ta kiểm tra ngay rằng, khi lấy các $\sh{N}_\lambda$ là những tổng hữu hạn của các $\sh{S}$-môđun dạng $\sh{P}_{\mu_n}$, ta thu được các đối tượng cần tìm.
\end{proof}

\begin{corollary}[3.4.6]
\label{II.3.4.6}
Giả sử các giả thiết của \sref{II.3.4.4} được thỏa mãn, và thêm rằng không gian nền của $Y$ quasi-compact;
cho $\sh{F}$ là một $\sh{O}_X$-môđun quasi-coherent hữu hạn sinh.
Tồn tại một số nguyên $n_0$ sao cho, với $n\geq n_0$, đồng cấu chính tắc
\[
  \sigma:p^*(p_*(\sh{F}(n)))\to\sh{F}(n)
\]
\sref[0]{0.4.4.3} là toàn ánh.
\end{corollary}

\begin{proof}
Với mọi $y\in Y$, cho $U$ là một lân cận mở affine của $y$ trong $Y$.
Tồn tại một số nguyên $n_0(U)$ sao cho, với $n\geq n_0(U)$, $\sh{F}(n)|p^{-1}(U)$ được sinh bởi một số hữu hạn các tiết diện của nó trên $p^{-1}(U)$ \sref{II.2.7.9};
nhưng các tiết diện này là ảnh chính tắc của những tiết diện của $p^*(p_*(\sh{F}(n)))$ trên $p^{-1}(U)$ (\sref[0]{0.3.7.1} và \sref[0]{0.4.4.3}), nên $\sh{F}(n)|p^{-1}(U)$ bằng ảnh chính tắc của $p^*(p_*(\sh{F}(n)))|p^{-1}(U)$.
Cuối cùng, vì $Y$ quasi-compact, tồn tại một phủ hữu hạn của $Y$ bởi các tập mở affine $U_i$; lấy $n_0$ là số lớn nhất trong các $n_0(U_i)$ thì chứng minh hoàn tất.
\end{proof}

\begin{remarks}[3.4.7]
\label{II.3.4.7}
Nếu $p=(\psi,\theta):X\to Y$ là một cấu xạ các không gian vành hóa và $\sh{F}$ là một $\sh{O}_X$-môđun, thì việc đồng cấu chính tắc $\sigma:p^*(p_*(\sh{F}))\to\sh{F}$ là toàn ánh được diễn giải như sau \sref[0]{0.4.4.1}: với mọi $x\in X$ và mọi tiết diện $s$ của $\sh{F}$ trên một lân cận mở $V$ của $x$, tồn tại một lân cận mở $U$ của $p(x)$ trong $Y$, một số hữu hạn các tiết diện $t_i$ ($1\leq i\leq m$) của $\sh{F}$ trên $p^{-1}(U)$, một lân cận $W\subset V\cap p^{-1}(U)$ của $x$, và các tiết diện $a_i$ ($1\leq i\leq m$) của $\sh{O}_X$ trên $W$, sao cho
\[
  s|W=\sum_{i=1}^m a_i\cdot(t_i|W).
\]
Khi $Y$ là một \emph{lược đồ affine} và $p_*(\sh{F})$ \emph{quasi-coherent}, điều kiện này tương đương với việc $\sh{F}$ \emph{được sinh bởi các tiết diện của nó trên $X$} \sref[0]{0.5.5.1}: thật vậy, nếu $Y=\Spec(A)$, ta có thể giả sử $U=D(f)$, với $f\in A$;
tồn tại một số nguyên $n>0$ và các tiết diện $s_i$ của $\sh{F}$ trên $X$ sao cho $t_i$ là hạn chế lên $p^{-1}(U)$ của $s_i/g^n$, với $g=\theta(f)$ (bằng cách áp dụng \sref[I]{I.1.4.1} cho $p_*(\sh{F})$);
vì $g$ khả nghịch trên $p^{-1}(U)$, nên ta có
\[
  s|W=\sum_i b_i\cdot(s_i|W),
\]
với $b_i=a_i(g|W)^{-n}$, từ đó suy ra khẳng định.
Khi $Y$ affine, Hệ quả~\sref{II.3.4.6} do đó cho lại \sref{II.2.7.9}.

Từ đó suy ra rằng, khi $Y$ là một tiền lược đồ bất kỳ, ba điều kiện
\oldpage[II]{61}
sau đây, đối với một $\sh{O}_X$-môđun quasi-coherent $\sh{F}$ sao cho $p_*(\sh{F})$ \emph{quasi-coherent}, là tương đương:
\begin{enumerate}
  \item[a)] \emph{Đồng cấu chính tắc $\sigma:p^*(p_*(\sh{F}))\to\sh{F}$ là toàn ánh.}
  \item[b)] \emph{Tồn tại một $\sh{O}_Y$-môđun quasi-coherent $\sh{G}$ và một đồng cấu toàn ánh $p^*(\sh{G})\to\sh{F}$.}
  \item[c)] \emph{Với mọi tập mở affine $U$ của $Y$, $\sh{F}|p^{-1}(U)$ được sinh bởi các tiết diện của nó trên $p^{-1}(U)$.}
\end{enumerate}

Thật vậy, ta vừa chứng minh sự tương đương giữa \emph{a)} và \emph{c)}.
Mặt khác, hiển nhiên \emph{a)} kéo theo \emph{b)}, vì theo giả thiết $p_*(\sh{F})$ quasi-coherent.
Ngược lại, mọi đồng cấu $u:p^*(\sh{G})\to\sh{F}$ đều phân tích thành
\[
  p^*(\sh{G})\to p^*(p_*(\sh{F}))\xrightarrow{\sigma}\sh{F}
\]
\sref[0]{0.3.5.4.4}; nếu $u$ toàn ánh thì $\sigma$ cũng vậy, điều này chứng minh \emph{b)} kéo theo \emph{a)}.
\end{remarks}

\begin{corollary}[3.4.8]
\label{II.3.4.8}
Giả sử các giả thiết của \sref{II.3.4.4} được thỏa mãn, và giả sử thêm rằng $Y$ là một lược đồ quasi-compact hoặc không gian nền của $Y$ là Noether.
Cho $\sh{F}$ là một $\sh{O}_X$-môđun quasi-coherent hữu hạn sinh;
khi đó tồn tại một số nguyên $n_0$ sao cho, với $n\geq n_0$, $\sh{F}$ đẳng cấu với một môđun thương của một $\sh{O}_X$-môđun dạng $(p^*(\sh{G}))(-n)$, trong đó $\sh{G}$ là một $\sh{O}_Y$-môđun quasi-coherent hữu hạn sinh (phụ thuộc vào $n$).
\end{corollary}

\begin{proof}
Vì cấu xạ cấu trúc $X\to Y$ là tách và thuộc kiểu hữu hạn, nên $p_*(\sh{F}(n))$ quasi-coherent \sref[I]{I.9.2.2}[b)] và do đó là giới hạn quy nạp của các $\sh{O}_Y$-môđun con quasi-coherent hữu hạn sinh của nó, do giả thiết trên $Y$ \sref[I]{I.9.4.9}.
Từ \sref{II.3.4.6}, \sref[0]{0.4.3.2} và \sref[0]{0.5.2.3}, ta suy ra $\sh{F}(n)$ là ảnh chính tắc của một $\sh{O}_X$-môđun dạng $p^*(\sh{G})$, trong đó $\sh{G}$ là một $\sh{O}_Y$-môđun con quasi-coherent hữu hạn sinh của $p_*(\sh{F}(n))$;
khi đó hệ quả suy ra từ \sref{II.3.2.5.2} và \sref{II.3.2.7.1}.
\end{proof}

\subsection{Các tính chất hàm tử}
\label{subsection:II.3.5}

\begin{env}[3.5.1]
\label{II.3.5.1}
Cho $Y$ là một tiền lược đồ, và $\sh{S}$, $\sh{S}'$ là hai $\sh{O}_Y$-đại số phân bậc quasi-coherent theo các bậc dương;
đặt $X=\Proj(\sh{S})$ và $X'=\Proj(\sh{S}')$, và cho $p$, $p'$ là các cấu xạ cấu trúc từ $X$, $X'$ đến $Y$.
Cho $\vphi:\sh{S}'\to\sh{S}$ là một đồng cấu các $\sh{O}_Y$-đại số phân bậc.
Với mỗi tập mở affine $U$ của $Y$, đặt $S_U=\Gamma(U,\sh{S})$ và $S'_U=\Gamma(U,\sh{S}')$;
đồng cấu $\vphi$ xác định một đồng cấu $\vphi_U:S'_U\to S_U$ của các $A_U$-đại số phân bậc, trong đó $A_U=\Gamma(U,\sh{O}_Y)$.
Tương ứng với nó, trong $p^{-1}(U)$ có một tập con mở $G(\vphi_U)$ và một cấu xạ $\Phi_U:G(\vphi_U)\to p'^{-1}(U)$ \sref{II.2.8.1}.
Hơn nữa, nếu $V\subset U$ là một tập mở affine, thì biểu đồ
\[
\label{II.3.5.1.1}
  \xymatrix{
    S'_U \ar[r]^{\vphi_U} \ar[d]
    & S_U \ar[d]
  \\S'_V \ar[r]_{\vphi_V}
    & S_V
  }
\tag{3.5.1.1}
\]
giao hoán, và từ các định nghĩa \sref{II.2.8.1}, ta kiểm tra ngay rằng
\[
  G(\vphi_V)=G(\vphi_U)\cap p^{-1}(V)
\]
và $\Phi_V$ là hạn chế của $\Phi_U$ lên $G(\vphi_V)$.
Như vậy, ta đã xác định một tập con mở $G(\vphi)$ của $X$ sao cho $G(\vphi)\cap p^{-1}(U)=G(\vphi_U)$ với mọi tập mở affine $U\subset Y$, và một $Y$-cấu xạ
\oldpage[II]{62}
\emph{affine} $\Phi:G(\vphi)\to X'$, được gọi là \emph{liên kết} với $\vphi$ và được ký hiệu bởi $\Proj(\vphi)$.
Khi, với mọi $y\in Y$, tồn tại một lân cận affine $U$ của $y$ sao cho $\Gamma(U,\sh{O}_Y)$-môđun $\Gamma(U,\sh{S}_+)$ được sinh bởi $\vphi(\Gamma(U,\sh{S}'_+))$, thì $G(\vphi_U)=p^{-1}(U)$, và do đó $G(\vphi)=X$.
\end{env}

\begin{proposition}[3.5.2]
\label{II.3.5.2}
\medskip\noindent
\begin{enumerate}
  \item[(i)] Nếu $\sh{M}$ là một $\sh{S}$-môđun phân bậc quasi-coherent, thì tồn tại một đẳng cấu chính tắc có tính hàm tử từ $\sh{O}_{X'}$-môđun $(\sh{M}_{[\vphi]})\supertilde$ lên $\sh{O}_{X'}$-môđun $\Phi_*(\widetilde{\sh{M}}|G(\vphi))$.
  \item[(ii)] Nếu $\sh{M}'$ là một $\sh{S}'$-môđun phân bậc quasi-coherent, thì tồn tại một đồng cấu chính tắc có tính hàm tử $\nu$ từ $(\sh{O}_X|G(\vphi))$-môđun $\Phi^*(\widetilde{\sh{M}'})$ vào $(\sh{O}_X|G(\vphi))$-môđun $(\sh{M}'\otimes_{\sh{S}'}\sh{S})\supertilde|G(\vphi)$.
    Nếu $\sh{S}'$ được sinh bởi $\sh{S}'_1$, thì $\nu$ là một đẳng cấu.
\end{enumerate}
\end{proposition}

\begin{proof}
Thật vậy, các đồng cấu đang xét đã được xác định khi $Y$ affine (\sref{II.2.8.7} và \sref{II.2.8.8});
trong trường hợp tổng quát, chỉ cần kiểm tra rằng chúng tương thích với các phép hạn chế từ một tập mở affine của $Y$ xuống một tập mở nhỏ hơn, điều này suy ra ngay từ tính giao hoán của \sref{II.3.5.1.1}.
\end{proof}

Đặc biệt, với mọi $n\in\bb{Z}$, ta có một đồng cấu chính tắc
\[
\label{II.3.5.2.1}
  \Phi^*(\sh{O}_{X'}(n))\to\sh{O}_X(n)|G(\vphi).
\tag{3.5.2.1}
\]

\begin{proposition}[3.5.3]
\label{II.3.5.3}
Cho $Y$, $Y'$ là hai tiền lược đồ, $\psi:Y'\to Y$ là một cấu xạ, và $\sh{S}$ là một $\sh{O}_Y$-đại số phân bậc quasi-coherent; đặt $\sh{S}'=\psi^*(\sh{S})$.
Khi đó $Y'$-lược đồ $X'=\Proj(\sh{S}')$ được đồng nhất chính tắc với $\Proj(\sh{S})\times_Y Y'$.
Hơn nữa, nếu $\sh{M}$ là một $\sh{S}$-môđun phân bậc quasi-coherent, thì $\sh{O}_{X'}$-môđun $(\psi^*(\sh{M}))\supertilde$ được đồng nhất với $\widetilde{\sh{M}}\otimes_Y\sh{O}_{Y'}$.
\end{proposition}

\begin{proof}
Trước hết, lưu ý rằng $\psi^*(\sh{S})$ và $\psi^*(\sh{M})$ là các $\sh{O}_{Y'}$-môđun quasi-coherent, cũng như các thành phần thuần nhất của chúng \sref[0]{0.5.1.4}.
Cho $U$ là một tập mở affine của $Y$, $U'\subset\psi^{-1}(U)$ là một tập mở affine của $Y'$, và $A$, $A'$ lần lượt là các vành của $U$, $U'$;
khi đó $\sh{S}|U=\widetilde{S}$, trong đó $S$ là một $A$-đại số phân bậc, và $\sh{S}'|U'$ được đồng nhất với $(S\otimes_A A')\supertilde$ \sref[I]{I.1.6.5};
khẳng định thứ nhất khi đó suy ra từ \sref{II.2.8.10} và \sref[I]{I.3.2.6.2}, vì ta kiểm tra ngay rằng phép chiếu $\Proj(\sh{S}'|U')\to\Proj(\sh{S}|U)$ được xác định bởi sự đồng nhất trên tương thích với các phép hạn chế trên $U$, $U'$ và do đó xác định một cấu xạ $\Proj(\sh{S}')\to\Proj(\sh{S})$.
Bây giờ, cho
\begin{align*}
  q &: \Proj(\sh{S})\to Y,
\\q'&: \Proj(\sh{S}')\to Y'
\end{align*}
là các cấu xạ cấu trúc;
$q'^{-1}(U')$ khi đó được đồng nhất với $q^{-1}(U)\times_U U'$, và hai bó $(\psi^*(\sh{M}))\supertilde|q'^{-1}(U')$, $(\widetilde{\sh{M}}\otimes_Y\sh{O}_{Y'})|q'^{-1}(U')$ đều được đồng nhất chính tắc với $(M\otimes_A A')\supertilde$, trong đó đặt $M=\Gamma(U,\sh{M})$, theo \sref{II.2.8.10} và \sref[I]{I.1.6.5};
từ đó suy ra khẳng định thứ hai, vì ta cũng kiểm tra ngay tính tương thích của các sự đồng nhất trên với các phép hạn chế.
\end{proof}

\begin{corollary}[3.5.4]
\label{II.3.5.4}
Với các ký hiệu của \sref{II.3.5.3}, $\sh{O}_{X'}(n)$ được đồng nhất chính tắc với $\sh{O}_X(n)\otimes_Y\sh{O}_{Y'}$ với mọi $n\in\bb{Z}$ (trong đó $X=\Proj(\sh{S})$).
\end{corollary}

\begin{proof}
Thật vậy, với các ký hiệu của \sref{II.3.5.3}, hiển nhiên $\psi^*(\sh{S}(n))=\sh{S}'(n)$ với mọi $n\in\bb{Z}$.
\end{proof}

\begin{env}[3.5.5]
\label{II.3.5.5}
Giữ các ký hiệu trên, ký hiệu bởi $\Psi$ phép chiếu chính tắc $X'\to X$ và đặt $\sh{M}'=\psi^*(\sh{M})$;
ta giả sử thêm rằng $\sh{S}$ được sinh bởi $\sh{S}_1$
\oldpage[II]{63}
và $X$ thuộc kiểu hữu hạn trên $Y$;
khi đó suy ra $\sh{S}'$ được sinh bởi $\sh{S}'_1$ (như thấy được bằng cách quy về trường hợp $Y$, $Y'$ affine), và $X'$ thuộc kiểu hữu hạn trên $Y'$ \sref[I]{I.6.3.4}.
Cho $\sh{F}$ là một $\sh{O}_X$-môđun và đặt $\sh{F}'=\Psi^*(\sh{F})$;
khi đó từ \sref{II.3.5.4} và \sref[0]{0.4.3.3} suy ra $\sh{F}'(n)=\Psi^*(\sh{F}(n))$ với mọi $n\in\bb{Z}$.
Hơn nữa, ta xác định một $\psi$-cấu xạ chính tắc $\theta_n:q_*(\sh{F}(n))\to q'_*(\sh{F}'(n))$ như sau: xét tính giao hoán của biểu đồ
\[
  \xymatrix{
    X \ar[d]_{q}
    & X' \ar[l]_{\Psi} \ar[d]^{q'}
  \\Y
    & Y' \ar[l]^{\psi}
  }
\]
ta cần xác định một đồng cấu $q_*(\sh{F}(n))\to\psi_*(q'_*(\Psi^*(\sh{F}(n))))=q_*(\Psi_*(\Psi^*(\sh{F}(n))))$, và chỉ cần lấy đồng cấu $\theta_n=q_*(\rho_n)$, trong đó $\rho_n$ là đồng cấu chính tắc $\sh{F}(n)\to\Psi_*(\Psi^*(\sh{F}(n)))$ \sref[0]{0.4.4.3}.
Hiển nhiên, với mọi tập mở affine $U$ của $Y$ và mọi tập mở affine $U'$ của $Y'$ sao cho $U'\subset\psi^{-1}(U)$, đồng cấu $\theta_n$ cảm sinh trên các tiết diện đồng cấu chính tắc \sref[0]{0.3.7.2} $\Gamma(q^{-1}(U),\sh{F}(n))\to\Gamma(q'^{-1}(U'),\sh{F}'(n))$.
Tính giao hoán của \sref{II.2.8.13.2} khi đó cho thấy rằng, nếu $\sh{F}$ quasi-coherent, biểu đồ
\[
  \xymatrix{
    \sh{F} \ar[r]^{\rho}
    & \sh{F}'
  \\(\bbGamma_*(\sh{F}))\supertilde \ar[u]^{\beta_{\sh{F}}} \ar[r]_{\widetilde{\theta}}
    & (\bbGamma_*(\sh{F}'))\supertilde \ar[u]_{\beta_{\sh{F}'}}
  }
\]
giao hoán (mũi tên ngang trên là $\Psi$-cấu xạ chính tắc $\sh{F}\to\Psi^*(\sh{F})$).

Tương tự, tính giao hoán của \sref{II.2.8.13.1} cho thấy biểu đồ
\[
  \xymatrix{
    \bbGamma_*(\widetilde{\sh{M}}) \ar[r]^{\theta}
    & \bbGamma_*(\widetilde{\sh{M}'})
  \\ \sh{M} \ar[r]_{\rho} \ar[u]^{\alpha_{\sh{M}}}
    & \sh{M}' \ar[u]_{\alpha_{\sh{M}'}}
  }
\]
giao hoán (mũi tên ngang dưới là $\psi$-cấu xạ chính tắc $\sh{M}\to\psi^*(\sh{M})$).
\end{env}

\begin{env}[3.5.6]
\label{II.3.5.6}
Bây giờ, xét hai tiền lược đồ $Y$, $Y'$, một cấu xạ $g:Y'\to Y$, một $\sh{O}_Y$-đại số phân bậc quasi-coherent $\sh{S}$, một $\sh{O}_{Y'}$-đại số phân bậc quasi-coherent $\sh{S}'$, và một $g$-cấu xạ các đại số phân bậc $u:\sh{S}\to\sh{S}'$, nghĩa là một đồng cấu các $\sh{O}_Y$-đại số phân bậc $\sh{S}\to g_*(\sh{S}')$;
ta biết rằng điều này tương đương với việc cho một đồng cấu các $\sh{O}_{Y'}$-đại số phân bậc $u^\sharp:g^*(\sh{S})\to\sh{S}'$.
Từ $u^\sharp$, ta suy ra một cách chính tắc một $Y'$-cấu xạ $W=\Proj(u^\sharp):G(u^\sharp)\to\Proj(g^*(\sh{S}))$, trong đó $G(u^\sharp)$ là một tập mở của $X'=\Proj(\sh{S}')$ \sref{II.3.5.1}.
Mặt khác, $X''=\Proj(g^*(\sh{S}))$ được đồng nhất
\oldpage[II]{64}
chính tắc với $X\times_Y Y'$, trong đó đặt $X=\Proj(\sh{S})$ \sref{II.3.5.3};
bằng cách hợp thành $\Proj(u^\sharp)$ với phép chiếu thứ nhất $p:X\times_Y Y'\to X$, ta thu được một cấu xạ $v:G(u^\sharp)\to X$, được ký hiệu bởi $\Proj(u)$, sao cho biểu đồ
\[
  \xymatrix{
    G(u^\sharp) \ar[r]^{v} \ar[d]
    & X \ar[d]
  \\Y' \ar[r]_{g}
    & Y
  }
\]
giao hoán.

Hơn nữa, với mọi $\sh{S}$-môđun phân bậc quasi-coherent $\sh{M}$, ta có một $v$-cấu xạ chính tắc
\[
\label{II.3.5.6.1}
  \nu:\widetilde{\sh{M}}\to(g^*(\sh{M})\otimes_{g^*(\sh{S})}\sh{S}')\supertilde|G(u^\sharp).
\tag{3.5.6.1}
\]

Thật vậy, $\nu^\sharp$ thu được bằng cách hợp thành các đồng cấu
\[
  v^*(\widetilde{\sh{M}})
  =W^*(p^*(\widetilde{\sh{M}}))
  \to W^*((g^*(\sh{M}))\supertilde)
  \to (g^*(\sh{M})\otimes_{g^*(\sh{S})}\sh{S}')\supertilde|G(u^\sharp)
\]
trong đó mũi tên thứ nhất đến từ đẳng cấu \sref{II.3.5.3}, còn mũi tên thứ hai là đồng cấu \sref{II.3.5.2}[(ii)];
khi $\sh{S}$ được sinh bởi $\sh{S}_1$, từ \sref{II.3.5.2} suy ra $\nu^\sharp$ là một \emph{đẳng cấu}.

Là một trường hợp riêng của \sref{II.3.5.6.1}, với mọi $n\in\bb{Z}$, ta có một $v$-cấu xạ chính tắc
\[
\label{II.3.5.6.2}
  \nu:\sh{O}_X(n)\to\sh{O}_{X'}(n)|G(u^\sharp).
\tag{3.5.6.2}
\]
\end{env}

\subsection{Các tiền lược đồ con đóng của một tiền lược đồ $\Proj(\sh{S})$}
\label{subsection:II.3.6}

\begin{env}[3.6.1]
\label{II.3.6.1}
Cho $Y$ là một tiền lược đồ và $\vphi:\sh{S}\to\sh{S}'$ là một đồng cấu bậc~$0$ của các $\sh{O}_Y$-đại số phân bậc quasi-coherent.
Ta nói $\vphi$ là \emph{(\textbf{TN})-toàn ánh} (tương ứng, \emph{(\textbf{TN})-đơn ánh}, \emph{(\textbf{TN})-song ánh}) nếu tồn tại một số nguyên $n$ sao cho, với $k\geq n$, $\vphi_k:\sh{S}_k\to\sh{S}'_k$ là toàn ánh (tương ứng, đơn ánh, song ánh).
Khi điều này xảy ra, việc nghiên cứu cấu xạ tương ứng $\Phi:\Proj(\sh{S}')\to\Proj(\sh{S})$ được quy về trường hợp $\vphi$ toàn ánh (tương ứng, đơn ánh, song ánh).
Điều này được chứng minh như trong \sref{II.2.9.1} (là trường hợp riêng tương ứng với $Y$ affine), bằng cách sử dụng \sref{II.3.1.8}.
Thay vì nói $\vphi$ là (\textbf{TN})-song ánh, ta cũng nói nó là một \emph{(\textbf{TN})-đẳng cấu}.
\end{env}

\begin{proposition}[3.6.2]
\label{II.3.6.2}
Cho $Y$ là một tiền lược đồ và $\sh{S}$ là một $\sh{O}_Y$-đại số phân bậc quasi-coherent; đặt $X=\Proj(\sh{S})$.
\begin{enumerate}
  \item[(i)] Nếu $\vphi:\sh{S}\to\sh{S}'$ là một đồng cấu (\textbf{TN})-toàn ánh của các $\sh{O}_Y$-đại số phân bậc, thì cấu xạ tương ứng $\Phi=\Proj(\vphi)$ \sref{II.3.5.1} được xác định trên toàn bộ $\Proj(\sh{S}')$ và là một phép nhúng đóng của $\Proj(\sh{S}')$ vào $X$.
    Nếu $\sh{J}$ là hạt nhân của $\vphi$, thì tiền lược đồ con đóng của $X$ liên kết với $\Phi$ được xác định bởi bó iđêan quasi-coherent $\widetilde{\sh{J}}$ của $\sh{O}_X$.
  \item[(ii)] Giả sử thêm rằng $\sh{S}_0=\sh{O}_Y$, $\sh{S}$ được sinh bởi $\sh{S}_1$, và $\sh{S}_1$ hữu hạn sinh.
    Cho $X'$ là một tiền lược đồ con đóng của $X=\Proj(\sh{S})$, được xác định bởi một bó iđêan quasi-coherent $\sh{I}$ của $\sh{O}_X$.
    Cho $\sh{J}$
\oldpage[II]{65}

là bó iđêan phân bậc quasi-coherent của $\sh{S}$, ảnh ngược của $\bbGamma_*(\sh{I})$ qua đồng cấu chính tắc $\alpha:\sh{S}\to\bbGamma_*(\sh{O}_X)$ \sref{II.3.3.2}, và đặt $\sh{S}'=\sh{S}/\sh{J}$.
Khi đó $X'$ là tiền lược đồ con liên kết \sref[I]{I.4.2.1} với phép nhúng đóng $\Proj(\sh{S}')\to X$ tương ứng với đồng cấu chính tắc $\sh{S}\to\sh{S}'$ của các $\sh{O}_Y$-đại số phân bậc.
\end{enumerate}
\end{proposition}

\begin{proof}
\medskip\noindent
\begin{enumerate}
  \item[(i)] Ta có thể giả sử $\vphi$ toàn ánh \sref{II.3.6.1}.
    Khi đó, với mọi tập mở affine $U$ của $Y$, đồng cấu $\Gamma(U,\sh{S})\to\Gamma(U,\sh{S}')$ là toàn ánh \sref[I]{I.1.3.9}; do đó, theo \sref{II.3.5.1}, ta có $G(\vphi)=\Proj(\sh{S}')$.
    Ta được quy ngay về việc chứng minh mệnh đề khi $Y$ affine, và khi ấy kết quả suy ra từ \sref{II.2.9.2}[(i)].
  \item[(ii)] Ta được quy về việc chứng minh rằng đồng cấu $\widetilde{\sh{J}}\to\sh{O}_X$ suy ra từ đơn ánh chính tắc $\sh{J}\to\sh{S}$ là một đẳng cấu từ $\widetilde{\sh{J}}$ lên $\sh{I}$;
    vì vấn đề mang tính địa phương trên $Y$, ta có thể giả sử $Y$ affine, có vành $A$, và khi đó $\sh{S}=\widetilde{S}$, trong đó $S$ là một $A$-đại số phân bậc được sinh bởi $S_1$, còn $S_1$ hữu hạn sinh trên $A$.
    Chỉ cần áp dụng \sref{II.2.9.2}[(ii)].
\end{enumerate}
\end{proof}

\begin{corollary}[3.6.3]
\label{II.3.6.3}
Dưới các điều kiện của \sref{II.3.6.2}[(i)], và giả sử thêm rằng $\sh{S}$ được sinh bởi $\sh{S}_1$, $\Phi^*(\sh{O}_X(n))$ được đồng nhất chính tắc với $\sh{O}_{X'}(n)$ với mọi $n\in\bb{Z}$.
\end{corollary}

\begin{proof}
Ta đã xác định một đẳng cấu chính tắc như vậy khi $Y$ affine \sref{II.2.9.3};
trong trường hợp tổng quát, chỉ cần kiểm tra rằng các đẳng cấu được xác định như trên đối với mỗi tập mở affine $U$ của $Y$ tương thích với phép chuyển từ $U$ sang một tập mở affine $U'\subset U$, điều này là hiển nhiên.
\end{proof}

\begin{corollary}[3.6.4]
\label{II.3.6.4}
Cho $Y$ là một tiền lược đồ, $\sh{S}$ là một $\sh{O}_Y$-đại số phân bậc quasi-coherent được sinh bởi $\sh{S}_1$, $\sh{M}$ là một $\sh{O}_Y$-môđun quasi-coherent, $u:\sh{M}\to\sh{S}_1$ là một $\sh{O}_Y$-đồng cấu toàn ánh, và $\overline{u}:\bb{S}_{\sh{O}_Y}(\sh{M})\to\sh{S}$ là đồng cấu của các $\sh{O}_Y$-đại số phân bậc mở rộng $u$ \sref{II.1.7.4}.
Khi đó, cấu xạ tương ứng với $\overline{u}$ là một phép nhúng đóng của $\Proj(\sh{S})$ vào $\Proj(\bb{S}_{\sh{O}_Y}(\sh{M}))$.
\end{corollary}

\begin{proof}
Thật vậy, $\overline{u}$ toàn ánh theo giả thiết, và ta áp dụng \sref{II.3.6.2}[(i)].
\end{proof}

\subsection{Các cấu xạ từ một tiền lược đồ vào một phổ thuần nhất}
\label{subsection:II.3.7}

\begin{env}[3.7.1]
\label{II.3.7.1}
Cho $q:X\to Y$ là một cấu xạ các tiền lược đồ, $\sh{L}$ là một $\sh{O}_X$-môđun khả nghịch, và $\sh{S}$ là một $\sh{O}_Y$-đại số phân bậc quasi-coherent theo các bậc dương;
khi đó $q^*(\sh{S})$ là một $\sh{O}_X$-đại số phân bậc quasi-coherent theo các bậc dương.
Xét $\sh{O}_X$-đại số phân bậc quasi-coherent $\sh{S}'=\bigoplus_{n\geq0}\sh{L}^{\otimes n}$, và giả sử cho một $\sh{O}_X$-đồng cấu của các đại số phân bậc
\[
  \psi:q^*(\sh{S})\longrightarrow
  \sh{S}'=\bigoplus_{n\geq0}\sh{L}^{\otimes n},
\]
điều này cũng tương đương với việc cho một $q$-cấu xạ của các đại số phân bậc
\[
  \psi^\flat:\sh{S}\longrightarrow q_*(\sh{S}').
\]
Ta biết $\Proj(\sh{S}')$ được đồng nhất chính tắc với $X$ (\sref{II.3.1.7} và \sref{II.3.1.8}[(iii)]);
từ $\psi$, ta suy ra chính tắc một tập con mở $G(\psi)$ của $X$ và một $Y$-cấu xạ
\[
\label{II.3.7.1.1}
  r_{\sh{L},\psi}:G(\psi)\longrightarrow \Proj(\sh{S})=P.
\tag{3.7.1.1}
\]
\oldpage[II]{66}

mà ta gọi là cấu xạ \emph{liên kết với $\sh{L}$ và $\psi$};
nhắc lại từ \sref{II.3.5.6} rằng cấu xạ này thu được bằng cách hợp thành $Y$-cấu xạ
\[
  \tau=\Proj(\psi):G(\psi)\longrightarrow\Proj(q^*(\sh{S}))
\]
với phép chiếu thứ nhất $\pi:\Proj(q^*(\sh{S}))=P\times_YX\to P$.
\end{env}

\begin{env}[3.7.2]
\label{II.3.7.2}
Ta mô tả tường minh $r=r_{\sh{L},\psi}$ khi $Y=\Spec(A)$ affine, và do đó $\sh{S}=\widetilde{S}$, trong đó $S$ là một $A$-đại số phân bậc theo các bậc dương.
Trước hết, giả sử $X=\Spec(B)$ cũng affine và $\sh{L}=\widetilde{L}$, trong đó $L$ là một $B$-môđun tự do hạng~$1$.
Khi đó $q^*(\sh{S})=(S\otimes_A B)\supertilde$ \sref[I]{I.1.6.5};
nếu $c$ là một phần tử sinh của $L$, thì $\psi_n:q^*(\sh{S}_n)\to\sh{L}^{\otimes n}$ tương ứng với đồng cấu $w_n:s\otimes b\mapsto bv_n(s)c^{\otimes n}$ từ $S_n\otimes_A B$ vào $L^{\otimes n}$, trong đó $v_n:S_n\to B$ là một đồng cấu các $A$-môđun, và các $v_n$ hợp lại thành một đồng cấu đại số $S\to B$.
Cho $f\in S_d$ ($d>0$) và đặt $g=v_d(f)$;
ta có $\pi^{-1}(D_+(f))=D_+(f\otimes1)$ theo \sref{II.2.8.10} và theo sự đồng nhất $D_+(f)$ với $\Spec(S_{(f)})$ \sref{II.2.3.6};
mặt khác, công thức \sref{II.2.8.1.1} cho thấy (có xét đến sự đồng nhất chính tắc của $X$ với $\Proj(\sh{S}')$) rằng
\[
  \tau^{-1}(D_+(f\otimes1))=D(g),
\]
do đó
\[
\label{II.3.7.2.1}
  r^{-1}(D_+(f))=D(g).
\tag{3.7.2.1}
\]

Hơn nữa, cấu xạ $\tau=\Proj(\psi)$, khi hạn chế trên $D(g)$, tương ứng với đồng cấu gửi $(s\otimes1)/(f\otimes1)^n$ (với $s\in S_{nd}$) tới $v_{nd}(s)/g^n$ \sref{II.2.8.1}, còn phép chiếu $\pi$, khi hạn chế trên $D_+(f\otimes1)$, tương ứng với đồng cấu $s/f^n\mapsto(s\otimes1)/(f\otimes1)^n$;
do đó, hạn chế của $r$ trên $D(g)$ tương ứng với đồng cấu các $A$-đại số $\omega:S_{(f)}\to B_g$ sao cho $\omega(s/f^n)=v_{nd}(s)/g^n$ với $s\in S_{nd}$ ($n>0$).
Chuyển sang trường hợp $X$ tùy ý ($Y$ vẫn affine), có xét đến \sref{II.2.8.1}, ta thu được:
\end{env}

\begin{proposition}[3.7.3]
\label{II.3.7.3}
Nếu $Y=\Spec(A)$ affine và $\sh{S}=\widetilde{S}$, trong đó $S$ là một $A$-đại số phân bậc, thì với mọi $f\in S_d=\Gamma(Y,\sh{S}_d)$, ta có
\[
\label{II.3.7.3.1}
  r_{\sh{L},\psi}^{-1}(D_+(f))=X_{\psi^\flat(f)}
  \qquad\bigl(\text{trong đó }\psi^\flat(f)\in\Gamma(X,\sh{L}^{\otimes d})\bigr)
\tag{3.7.3.1}
\]
và hạn chế $X_{\psi^\flat(f)}\to D_+(f)=\Spec(S_{(f)})$ của $r_{\sh{L},\psi}$ tương ứng \sref[I]{I.2.2.4} với đồng cấu đại số
\[
\label{II.3.7.3.2}
  \psi_{(f)}^\flat:S_{(f)}\longrightarrow
  \Gamma(X_{\psi^\flat(f)},\sh{O}_X)
\tag{3.7.3.2}
\]
sao cho, với $s\in S_{nd}=\Gamma(Y,\sh{S}_{nd})$,
\[
\label{II.3.7.3.3}
  \psi_{(f)}^\flat(s/f^n)
  =(\psi^\flat(s)|X_{\psi^\flat(f)})
   (\psi^\flat(f)|X_{\psi^\flat(f)})^{-n}.
\tag{3.7.3.3}
\]
\end{proposition}

Ta nói $r_{\sh{L},\psi}$ \emph{được xác định khắp nơi} nếu $G(\psi)=X$.
Để điều này xảy ra, điều kiện cần và đủ hiển nhiên là $G(\psi)\cap q^{-1}(U)=q^{-1}(U)$ với mọi tập mở affine $U\subset Y$;
nói cách khác, vấn đề mang tính \emph{địa phương trên $Y$}.
Nếu $Y$ affine, thì $G(\psi)$ là hợp của các $r^{-1}(D_+(f))$, khi $f$ chạy qua các phần tử thuần nhất trong $S_+$ \sref{II.2.8.1};
theo \sref{II.3.7.3.1}, các $X_{\psi^\flat(f)}$ phải lập thành một \emph{phủ} của $X$;
nói cách khác:

\begin{corollary}[3.7.4]
\label{II.3.7.4}
Dưới các giả thiết của \sref{II.3.7.3}, điều kiện cần và đủ để $r_{\sh{L},\psi}$ được xác định khắp nơi là: với mọi $x\in X$, tồn tại một số nguyên $n>0$ và một tiết diện $s$ của $\sh{S}_n$ trên $Y$ sao cho, nếu đặt $t=\psi^\flat(s)\in\Gamma(X,\sh{L}^{\otimes n})$, thì $t(x)\neq0$.
\end{corollary}
\oldpage[II]{67}

Ta lưu ý rằng điều kiện này luôn được thỏa mãn nếu $\psi$ là (\textbf{TN})-toàn ánh.

Tương tự, vấn đề $r_{\sh{L},\psi}$ có \emph{trội} hay không mang tính địa phương trên $Y$, và ta có:

\begin{corollary}[3.7.5]
\label{II.3.7.5}
Dưới các giả thiết của \sref{II.3.7.3}, điều kiện cần và đủ để $r_{\sh{L},\psi}$ trội là: với mọi số nguyên $n>0$, mọi tiết diện $s\in S_n$ sao cho $\psi^\flat(s)\in\Gamma(X,\sh{L}^{\otimes n})$ lũy linh cục bộ đều tự nó lũy linh.
\end{corollary}

\begin{proof}
Thật vậy, ta phải biểu thị rằng $r_{\sh{L},\psi}^{-1}(D_+(s))$ chỉ rỗng khi $D_+(s)$ rỗng, và do đó hệ quả suy ra từ \sref{II.3.7.3.1} và \sref{II.2.3.7}.
\end{proof}

\begin{proposition}[3.7.6]
\label{II.3.7.6}
Cho $q:X\to Y$ là một cấu xạ, $\sh{L}$ là một $\sh{O}_X$-môđun khả nghịch, $\sh{S}$ và $\sh{S}'$ là hai $\sh{O}_Y$-đại số phân bậc quasi-coherent, $u:\sh{S}'\to\sh{S}$ là một đồng cấu các đại số phân bậc, $\psi:q^*(\sh{S})\to\bigoplus_{n\geq0}\sh{L}^{\otimes n}$ là một đồng cấu các đại số phân bậc, và $\psi'=\psi\circ q^*(u)$ là đồng cấu hợp thành.
Nếu $r_{\sh{L},\psi'}$ được xác định khắp nơi thì $r_{\sh{L},\psi}$ cũng vậy;
nếu $u$ là (\textbf{TN})-toàn ánh và $r_{\sh{L},\psi'}$ trội thì $r_{\sh{L},\psi}$ cũng trội;
ngược lại, nếu $u$ là (\textbf{TN})-đơn ánh và $r_{\sh{L},\psi}$ trội thì $r_{\sh{L},\psi'}$ cũng trội.
\end{proposition}

\begin{proof}
Thật vậy, ta có $G(\psi')\subset G(\psi)$ \sref{II.2.8.4}, suy ra khẳng định thứ nhất;
nếu $u$ là (\textbf{TN})-toàn ánh, thì $\Proj(u):\Proj(\sh{S})\to\Proj(\sh{S}')$ được xác định khắp nơi và là một phép nhúng đóng;
vì $r_{\sh{L},\psi'}$ là hợp thành của $\Proj(u)$ với hạn chế của $r_{\sh{L},\psi}$ trên $G(\psi')$, nên nếu $r_{\sh{L},\psi'}$ trội thì $r_{\sh{L},\psi}$ cũng trội.
Cuối cùng, nếu $u$ là (\textbf{TN})-đơn ánh, ta biết $\Proj(u)$ là một cấu xạ trội từ $G(u)$ vào $\Proj(\sh{S}')$ \sref{II.2.8.3};
vì $G(\psi')$ là ảnh ngược của $G(u)$ qua $r_{\sh{L},\psi}$, ta thấy rằng nếu $r_{\sh{L},\psi}$ trội thì $r_{\sh{L},\psi'}$ cũng trội.
\end{proof}

\begin{proposition}[3.7.7]
\label{II.3.7.7}
Cho $Y$ là một tiền lược đồ quasi-compact, $q:X\to Y$ là một cấu xạ quasi-compact, $\sh{L}$ là một $\sh{O}_X$-môđun khả nghịch, và $\sh{S}$ là một $\sh{O}_Y$-đại số phân bậc quasi-coherent, là giới hạn quy nạp lọc của một hệ quy nạp $(\sh{S}^\lambda)$ các $\sh{O}_Y$-đại số phân bậc quasi-coherent.
Cho $\vphi_\lambda:\sh{S}^\lambda\to\sh{S}$ là đồng cấu chính tắc, $\psi:q^*(\sh{S})\to\bigoplus_{n\geq0}\sh{L}^{\otimes n}$ là một đồng cấu các đại số phân bậc, và đặt $\psi_\lambda=\psi\circ q^*(\vphi_\lambda)$.
Điều kiện cần và đủ để $r_{\sh{L},\psi}$ được xác định khắp nơi là tồn tại một $\lambda$ sao cho $r_{\sh{L},\psi_\lambda}$ được xác định khắp nơi;
khi đó $r_{\sh{L},\psi_\mu}$ được xác định khắp nơi với mọi $\mu\geq\lambda$.
\end{proposition}

\begin{proof}
Điều kiện là đủ theo \sref{II.3.7.6}.
Ngược lại, giả sử $r_{\sh{L},\psi}$ được xác định khắp nơi;
ta có thể quy về trường hợp $Y$ affine, vì nếu với mọi tập mở affine $U\subset Y$ tồn tại $\lambda(U)$ sao cho hạn chế của $r_{\sh{L},\psi_{\lambda(U)}}$ trên $q^{-1}(U)$ được xác định khắp nơi, thì chỉ cần ($Y$ quasi-compact) phủ $Y$ bởi một số hữu hạn các tập mở affine $U_i$ và lấy $\lambda\geq\lambda(U_i)$ với mọi chỉ số $i$, theo \sref{II.3.7.6}.
Nếu $Y$ affine, giả thiết kéo theo rằng với mọi $x\in X$, có một tiết diện $s^{(x)}$ của một $S_n$ sao cho, nếu đặt $t^{(x)}=\psi^\flat(s^{(x)})$, thì $t^{(x)}(x)\neq0$ (trong đó $t^{(x)}$ được xem là một tiết diện của $\sh{L}^{\otimes n}$ trên $X$); do đó $t^{(x)}(z)\neq0$ với mọi $z$ trong một lân cận $V(x)$ của $x$.
Phủ $X$ bởi một số hữu hạn các $V(x_i)$ và gọi $s^{(i)}$ là các tiết diện tương ứng của $S$;
khi đó tồn tại một $\lambda$ sao cho mọi $s^{(i)}$ đều có dạng $\vphi_\lambda(s_\lambda^{(i)})$, với $s_\lambda^{(i)}\in S^\lambda$ cho mọi $i$;
do đó, từ \sref{II.3.7.4} suy ra $r_{\sh{L},\psi_\lambda}$ được xác định khắp nơi.
Khẳng định cuối cùng là hệ quả hiển nhiên của \sref{II.3.7.6}.
\end{proof}

\begin{corollary}[3.7.8]
\label{II.3.7.8}
Dưới các giả thiết của \sref{II.3.7.7}, nếu các $r_{\sh{L},\psi_\lambda}$ đều trội thì $r_{\sh{L},\psi}$ cũng trội;
mệnh đề đảo đúng khi các $\vphi_\lambda$ đơn ánh.
\end{corollary}
\oldpage[II]{68}

\begin{proof}
Khẳng định thứ hai là một trường hợp riêng của \sref{II.3.7.6};
mặt khác, để chứng minh $r_{\sh{L},\psi}$ trội, ta có thể chỉ xét trường hợp $Y$ affine;
nếu $s\in S$ sao cho $\psi^\flat(s)$ lũy linh cục bộ, vì có thể viết $s=\vphi_\lambda(s_\lambda)$ với ít nhất một $\lambda$, nên từ giả thiết và \sref{II.3.7.5} suy ra $s_\lambda$ lũy linh, do đó $s$ cũng lũy linh, và tiêu chuẩn của \sref{II.3.7.5} được áp dụng.
\end{proof}

\begin{remarks}[3.7.9]
\label{II.3.7.9}
\medskip\noindent
\begin{enumerate}
  \item[(i)] Với các ký hiệu của \sref{II.3.7.1}, và có xét đến \sref{II.3.2.10}, với mọi $n\in\bb{Z}$ ta có một đồng cấu chính tắc
    \[
    \label{II.3.7.9.1}
      \theta:r_{\sh{L},\psi}^*(\sh{O}_P(n))\longrightarrow
      \sh{L}^{\otimes n}
    \tag{3.7.9.1}
    \]
    được xác định một cách tổng quát trong \sref{II.3.5.6.2}.
    Ta thấy ngay rằng, dưới các điều kiện của \sref{II.3.7.3}, hạn chế của đồng cấu này trên $X_{\psi^\flat(f)}$ gửi $s/f^k$ ($s\in S_{n+kd}$) tới phần tử $(\psi^\flat(s)|X_{\psi^\flat(f)})(\psi^\flat(f)|X_{\psi^\flat(f)})^{-k}$.
  \item[(ii)] Cho $\sh{F}$ là một $\sh{O}_X$-môđun quasi-coherent, và giả sử $q$ quasi-compact và tách, sao cho với mọi $n\geq0$, $q_*(\sh{F}\otimes\sh{L}^{\otimes n})$ là một $\sh{O}_Y$-môđun quasi-coherent \sref[I]{I.9.2.2}.
    Cho $\sh{M}'=\bigoplus_{n\geq0}\sh{F}\otimes\sh{L}^{\otimes n}$, là một $\sh{S}'$-môđun phân bậc quasi-coherent, và xét ảnh trực tiếp của nó
    $\sh{M}=q_*(\sh{M}')=\bigoplus_{n\geq0}q_*(\sh{F}\otimes\sh{L}^{\otimes n})$ (đây là một $\sh{S}$-môđun quasi-coherent nhờ đồng cấu $\psi^\flat$).
    Ta sẽ thấy rằng có một đồng cấu chính tắc của các $\sh{O}_X$-môđun
    \[
    \label{II.3.7.9.2}
      \xi:r_{\sh{L},\psi}^*(\widetilde{\sh{M}})
      \longrightarrow\sh{F}|G(\psi).
    \tag{3.7.9.2}
    \]

    Thật vậy, trong \sref{II.3.5.6.1} ta đã xác định một đồng cấu chính tắc
    \[
    \label{II.3.7.9.3}
      r_{\sh{L},\psi}^*(\widetilde{\sh{M}})
      \longrightarrow
      (q^*(\sh{M})\otimes_{q^*(\sh{S})}\sh{S}')\supertilde|G(\psi),
    \tag{3.7.9.3}
    \]
    trong đó vế phải được xem là một bó quasi-coherent trên $\Proj(\sh{S}')$.
    Mặt khác, ta có một đồng cấu chính tắc
    \[
    \label{II.3.7.9.4}
      q^*(q_*(\sh{M}'))\otimes_{q^*(\sh{S})}\sh{S}'
      \longrightarrow\sh{M}'
    \tag{3.7.9.4}
    \]
    với mọi $\sh{S}'$-môđun phân bậc quasi-coherent $\sh{M}'$:
    với mọi tập mở $U$ của $X$, mọi tiết diện $t'$ của $q^*(q_*(\sh{M}'_h))$ trên $U$ và mọi tiết diện $b'$ của $\sh{S}'_k$ trên $U$, ta cho $t'\otimes b'$ tương ứng với tiết diện $b'\sigma(t')$ của $\sh{M}'_{h+k}$, trong đó $\sigma(t')$ là tiết diện của $\sh{M}'_h$ trên $U$ tương ứng chính tắc \sref[0]{0.4.4.3} với $t'$.
    Từ đó suy ra một đồng cấu chính tắc
    \[
    \label{II.3.7.9.5}
      (q^*(q_*(\sh{M}'))\otimes_{q^*(\sh{S})}\sh{S}')\supertilde|G(\psi)
      \longrightarrow\widetilde{\sh{M}'}|G(\psi),
    \tag{3.7.9.5}
    \]
    và cuối cùng, vì $\widetilde{\sh{M}'}$ được đồng nhất chính tắc với $\sh{F}$ \sref{II.3.2.9}[(i)], ta thu được đúng đồng cấu chính tắc cần tìm.

    Dưới các điều kiện của \sref{II.3.7.3}, hạn chế của đồng cấu này trên $X_{\psi^\flat(f)}$ được mô tả tường minh như sau:
    cho một tiết diện $t_{nd}$ của $\sh{F}\otimes\sh{L}^{\otimes nd}$ trên $X$; nếu $t'_{nd}$ là tiết diện $t_{nd}$ được xem như một tiết diện của $q_*(\sh{F}\otimes\sh{L}^{\otimes nd})$ trên $Y$, thì phần tử $t'_{nd}/f^n$ tương ứng với tiết diện $(t_{nd}|X_{\psi^\flat(f)})(\psi^\flat(f)|X_{\psi^\flat(f)})^{-n}$ của $\sh{F}$ trên $X_{\psi^\flat(f)}$.
\end{enumerate}
\end{remarks}
\oldpage[II]{69}

\subsection{Các tiêu chuẩn nhúng vào một phổ thuần nhất}
\label{subsection:II.3.8}

\begin{env}[3.8.1]
\label{II.3.8.1}
Với các ký hiệu của \sref{II.3.7.1}, vấn đề $r_{\sh{L},\psi}$ có phải là một phép nhúng (tương ứng, một phép nhúng mở, một phép nhúng đóng) hay không hiển nhiên mang tính \emph{địa phương trên $Y$}.
\end{env}

\begin{proposition}[3.8.2]
\label{II.3.8.2}
Dưới các giả thiết của \sref{II.3.7.3}, điều kiện cần và đủ để $r_{\sh{L},\psi}$ được xác định khắp nơi và là một phép nhúng là tồn tại một họ tiết diện $s_\alpha\in S_{n_\alpha}$ ($n_\alpha>0$) sao cho, nếu đặt $f_\alpha=\psi^\flat(s_\alpha)$, thì các điều kiện sau được thỏa mãn:
\begin{enumerate}
  \item[(i)] Các $X_{f_\alpha}$ lập thành một phủ của $X$.
  \item[(ii)] Các $X_{f_\alpha}$ là những tập mở affine.
  \item[(iii)] Với mọi $\alpha$ và mọi $t\in\Gamma(X_{f_\alpha},\sh{O}_X)$, tồn tại một số nguyên $m>0$ và một $s\in S_{mn_\alpha}$ sao cho $t=(\psi^\flat(s)|X_{f_\alpha})(f_\alpha|X_{f_\alpha})^{-m}$.
\end{enumerate}

Điều kiện cần và đủ để $r_{\sh{L},\psi}$ được xác định khắp nơi và là một phép nhúng mở là tồn tại một họ $(s_\alpha)$ thỏa mãn các điều kiện (i), (ii), (iii) và:
\begin{enumerate}
  \item[(iv)] Với mọi $n>0$ và mọi $s\in S_{nn_\alpha}$ sao cho $\psi^\flat(s)|X_{f_\alpha}=0$, tồn tại một số nguyên $k>0$ sao cho $s_\alpha^ks=0$.
\end{enumerate}

Điều kiện cần và đủ để $r_{\sh{L},\psi}$ được xác định khắp nơi và là một phép nhúng đóng là tồn tại một họ $(s_\alpha)$ thỏa mãn (i), (ii), (iii) và:
\begin{enumerate}
  \item[(v)] Các $D_+(s_\alpha)$ lập thành một phủ của $P=\Proj(S)$.
\end{enumerate}
\end{proposition}

\begin{proof}
Thật vậy, điều kiện cần và đủ để $r$ là một phép nhúng (tương ứng, một phép nhúng đóng) là tồn tại một phủ của $r(G(\psi))$ (tương ứng, của $P$) bởi các $D_+(s_\alpha)$ sao cho, nếu đặt $V_\alpha=r^{-1}(D_+(s_\alpha))$, thì hạn chế của $r$ trên $V_\alpha$ là một \emph{phép nhúng đóng} của $V_\alpha$ vào $D_+(s_\alpha)$ \sref[I]{I.4.2.4}.
Điều kiện (i) đồng thời biểu thị rằng $r$ được xác định khắp nơi và các $D_+(s_\alpha)$ phủ $r(X)$, theo \sref{II.3.7.3.1};
vì $D_+(s_\alpha)$ affine, các điều kiện (ii) và (iii) biểu thị rằng hạn chế của $r$ trên $X_{f_\alpha}$ là một phép nhúng đóng vào $D_+(s_\alpha)$, theo \sref[I]{I.4.2.3};
cuối cùng, vì (iii) và (iv) biểu thị rằng $\psi_{(s_\alpha)}^\flat$ song ánh (theo ký hiệu của \sref{II.3.7.3.2}), nên (ii), (iii) và (iv) biểu thị rằng hạn chế của $r$ trên $X_{f_\alpha}$ là một đẳng cấu lên $D_+(s_\alpha)$ với mọi $\alpha$;
do đó (i), (ii), (iii) và (iv) biểu thị rằng $r$ là một phép nhúng mở.
\end{proof}

\begin{corollary}[3.8.3]
\label{II.3.8.3}
Dưới các giả thiết của \sref{II.3.7.6}, nếu $r_{\sh{L},\psi'}$ được xác định khắp nơi và là một phép nhúng thì $r_{\sh{L},\psi}$ cũng vậy.
Nếu giả sử thêm rằng $u$ là (\textbf{TN})-toàn ánh và nếu $r_{\sh{L},\psi'}$ là một phép nhúng mở (tương ứng, đóng), thì $r_{\sh{L},\psi}$ cũng vậy.
\end{corollary}

\begin{proof}
Thật vậy, theo \sref{II.3.8.2}, có một họ $s'_\alpha\in S'_{n_\alpha}$ sao cho, nếu đặt $f_\alpha=\psi'^\flat(s'_\alpha)$, thì các điều kiện (i), (ii), (iii) được thỏa mãn.
Mặt khác, nếu đặt $s_\alpha=u(s'_\alpha)$ thì ta cũng có $f_\alpha=\psi^\flat(s_\alpha)$; và nếu $t=(\psi'^\flat(s')|X_{f_\alpha})(f_\alpha|X_{f_\alpha})^{-m}$ thì, khi đặt $s=u(s')$, ta cũng có $t=(\psi^\flat(s)|X_{f_\alpha})(f_\alpha|X_{f_\alpha})^{-m}$, suy ra khẳng định thứ nhất.
Khẳng định thứ hai suy ra ngay từ việc $\Proj(u)$ là một phép nhúng đóng.
\end{proof}

\begin{proposition}[3.8.4]
\label{II.3.8.4}
Giả sử các giả thiết của \sref{II.3.7.7} được thỏa mãn và, hơn nữa, $q:X\to Y$ là một cấu xạ thuộc kiểu hữu hạn.
Khi đó, điều kiện cần và đủ để $r_{\sh{L},\psi}$ được xác định khắp nơi và là một phép nhúng là tồn tại một $\lambda$ sao cho $r_{\sh{L},\psi_\lambda}$ được xác định khắp nơi và là một phép nhúng;
trong trường hợp ấy, $r_{\sh{L},\psi_\mu}$ được xác định khắp nơi và là một phép nhúng với mọi $\mu\geq\lambda$.
\end{proposition}
\oldpage[II]{70}
\begin{proof}
Xét đến \sref{II.3.8.3}, chỉ cần chứng minh rằng nếu $r_{\sh{L},\psi}$ được xác định ở mọi nơi và là một phép nhúng, thì $r_{\sh{L},\psi_\lambda}$ cũng có các tính chất ấy với ít nhất một $\lambda$.
Bằng cùng lập luận như trong \sref{II.3.7.7}, sử dụng tính quasi-compact của $Y$, trước hết ta quy về trường hợp $Y=\Spec(A)$ là affine.
Vì $X$ quasi-compact, \sref{II.3.8.2} cho thấy tồn tại một họ \emph{hữu hạn} $(s_i\in S_{n_i})$ các phần tử của $S$ thỏa mãn các điều kiện~(i), (ii) và (iii).
Cấu xạ $X_{f_i}\to Y$ (trong đó $f_i=\psi^\flat(s_i)$) thuộc kiểu hữu hạn: thật vậy, đó là một cấu xạ giữa các lược đồ affine, nên nó quasi-compact \sref[I]{I.6.6.1}, và nó thuộc kiểu hữu hạn địa phương vì $q$ thuộc kiểu hữu hạn \sref[I]{I.6.3.2}; khi đó kết luận suy ra từ \sref[I]{I.6.6.3}.
Do đó, vành $B_i$ của $X_{f_i}$ là một $A$-đại số thuộc kiểu hữu hạn \sref[I]{I.6.3.3};
cho $(t_{ij})$ là một họ sinh của đại số này.
Theo giả thiết, tồn tại các phần tử $s'_{ij}\in S_{m_{ij}n_i}$ sao cho
\[
  t_{ij} = (\psi^\flat(s'_{ij})|X_{f_i})(\psi^\flat(s_i)|X_{f_i})^{-m_{ij}}.
\]
Cũng theo giả thiết, tồn tại một $\lambda$ và các phần tử $s_{i\lambda}\in S_{n_i}^\lambda$, $s'_{ij\lambda}\in S_{m_{ij}n_i}^\lambda$ có ảnh qua $\vphi_\lambda$ lần lượt là $s_i$ và $s'_{ij}$;
hiển nhiên $r_{\sh{L},\psi_\lambda}$ thỏa mãn các điều kiện~(i), (ii) và (iii) của \sref{II.3.8.2}.
\end{proof}

\begin{proposition}[3.8.5]
\label{II.3.8.5}
Cho $Y$ là một lược đồ quasi-compact hoặc một tiền lược đồ có không gian nền Noether, $q:X\to Y$ là một cấu xạ \emph{thuộc kiểu hữu hạn}, $\sh{L}$ là một $\sh{O}_X$-môđun khả nghịch, $\sh{S}$ là một $\sh{O}_Y$-đại số phân bậc quasi-coherent, và $\psi:q^*(\sh{S})\to\bigoplus_{n\geq0}\sh{L}^{\otimes n}$ là một đồng cấu các đại số phân bậc.
Điều kiện cần và đủ để $r_{\sh{L},\psi}$ được xác định ở mọi nơi và là một phép nhúng là tồn tại một số nguyên $n>0$ và một $\sh{O}_Y$-môđun con quasi-coherent \emph{hữu hạn sinh} $\sh{E}$ của $\sh{S}_n$ sao cho
\begin{enumerate}
  \item[a)] đồng cấu $\psi_n\circ q^*(j_n):q^*(\sh{E})\to\sh{L}^{\otimes n}$ (trong đó $j_n:\sh{E}\to\sh{S}_n$ là phép nhúng chính tắc) là toàn ánh; và
  \item[b)] nếu ký hiệu bởi $\sh{S}'$ $\sh{O}_Y$-đại số con (phân bậc) của $\sh{S}$ sinh bởi $\sh{E}$, và bởi $\psi'$ đồng cấu $\psi\circ q^*(j')$ (trong đó $j'$ là phép nhúng $\sh{S}'$ vào $\sh{S}$), thì $r_{\sh{L},\psi'}$ được xác định ở mọi nơi và là một phép nhúng.
\end{enumerate}

Khi đó, mọi $\sh{O}_Y$-môđun con quasi-coherent $\sh{E}'$ của $\sh{S}_n$ chứa $\sh{E}$ cũng có cùng tính chất, và ảnh của $\otimes^k\sh{E}$ trong $\sh{S}_{kn}$ cũng vậy với mọi $k>0$.
\end{proposition}

\begin{proof}
Tính đủ của điều kiện và hai khẳng định cuối là các trường hợp riêng của \sref{II.3.8.3}, có xét đến đẳng cấu chính tắc giữa $\Proj(\sh{S})$ và $\Proj(\sh{S}^{(k)})$ \sref{II.3.1.8}.

Cho $(U_i)$ là một phủ mở hữu hạn của $Y$ bởi các tập mở affine và đặt $A_i=A(U_i)$.
Vì $q^{-1}(U_i)$ quasi-compact, giả thiết rằng $r_{\sh{L},\psi}$ được xác định ở mọi nơi và là một phép nhúng, cùng với \sref{II.3.8.2}, kéo theo sự tồn tại của một họ hữu hạn $(s_{ij}\in S_{n_{ij}}^{(i)})$ các phần tử của $S^{(i)}=\Gamma(U_i,\sh{S})$ thỏa mãn các điều kiện~(i), (ii) và (iii).
Như trong chứng minh của \sref{II.3.8.4}, ta thấy cấu xạ $X_{f_{ij}}\to U_i$ (trong đó $f_{ij}=\psi^\flat(s_{ij})$) thuộc kiểu hữu hạn, và do đó vành $B_{ij}$ của $X_{f_{ij}}$ là một $A_i$-đại số thuộc kiểu hữu hạn, có một hệ sinh gồm các phần tử dạng $(\psi^\flat(t_{ijk})|X_{f_{ij}})(f_{ij}|X_{f_{ij}})^{-m_{ijk}}$, trong đó $t_{ijk}\in S_{m_{ijk}n_{ij}}^{(i)}$.
Cho $n$ là một bội chung của các $m_{ijk}n_{ij}$;
thay (với mỗi $(i,j,k)$) $s_{ij}$ bằng một lũy thừa $s_{ij}^\rho$ sao cho $\rho n_{ij}=n$, và nhân $t_{ijk}$ với $s_{ij}^{\rho-m_{ijk}}$, ta có thể giả sử rằng, với mỗi $i$, mọi $s_{ij}$ và $t_{ijk}$ đều thuộc $S_n^{(i)}$ và $m_{ijk}=1$.
Cho $E_i$ là $A_i$-môđun con của $S_n^{(i)}$ sinh bởi các phần tử ấy (với $i$ cố định).
Tồn tại một $\sh{O}_Y$-môđun con coherent hữu hạn sinh $\sh{E}_i$ của $\sh{S}_n$ sao cho $\sh{E}_i|U_i=(E_i)\supertilde$ \sref[I]{I.9.4.7}.
\oldpage[II]{71}
Hiển nhiên $\sh{O}_Y$-môđun con $\sh{E}$ của $\sh{S}_n$ bằng tổng các $\sh{E}_i$ là đối tượng cần tìm (vì mỗi tiết diện $f_{ij}$ có tính chất sau: với mọi $x\in X_{f_{ij}}$, tồn tại một lân cận affine $V\subset X_{f_{ij}}$ của $x$ sao cho $f_{ij}|V$ là một cơ sở của $\Gamma(V,\sh{L}^{\otimes n})$).
\end{proof}

\section{Các phân thớ xạ ảnh; các bó ample}
\label{section:II.4}


\subsection{Định nghĩa các phân thớ xạ ảnh}
\label{subsection:II.4.1}

\begin{definition}[4.1.1]
\label{II.4.1.1}
Cho $Y$ là một tiền lược đồ, $\sh{E}$ là một $\sh{O}_Y$-môđun quasi-coherent, và $\bb{S}_{\sh{O}_Y}(\sh{E})$ là $\sh{O}_Y$-đại số đối xứng của $\sh{E}$ \sref{II.1.7.4}, vốn quasi-coherent \sref{II.1.7.7}.
Ta gọi $Y$-lược đồ $P=\Proj(\bb{S}_{\sh{O}_Y}(\sh{E}))$ là \emph{phân thớ xạ ảnh trên $Y$ do $\sh{E}$ xác định}, và ký hiệu nó bởi $\bb{P}(\sh{E})$.
$\sh{O}_P$-môđun $\sh{O}_P(1)$ được gọi là \emph{bó cơ bản trên $P$}.
\end{definition}

Khi $Y$ affine với vành $A$, ta có $\sh{E}=\widetilde{E}$, trong đó $E$ là một $A$-môđun, và khi đó ta viết $\bb{P}(E)$ thay cho $\bb{P}(\widetilde{E})$.

Khi lấy $\sh{E}=\sh{O}_Y^n$, ta viết $\bb{P}_Y^{n-1}$ thay cho $\bb{P}(\sh{E})$;
nếu thêm nữa $Y$ affine với vành $A$, ta cũng viết $\bb{P}_A^{n-1}$ thay cho $\bb{P}_Y^{n-1}$.
Vì $\bb{S}_{\sh{O}_Y}(\sh{O}_Y)$ được đồng nhất chính tắc với $\sh{O}_Y[T]$ \sref{II.1.7.4}, $\bb{P}_Y^0$ được đồng nhất chính tắc với $Y$ \sref{II.3.1.7};
Ví dụ~\sref{II.2.4.3} chính là $\bb{P}_K^1$.

\begin{env}[4.1.2]
\label{II.4.1.2}
Cho $\sh{E}$ và $\sh{F}$ là hai $\sh{O}_Y$-môđun quasi-coherent, và $u:\sh{E}\to\sh{F}$ là một đồng cấu $\sh{O}_Y$-môđun;
ứng với nó một cách chính tắc là một đồng cấu $\bb{S}(u):\bb{S}_{\sh{O}_Y}(\sh{E})\to\bb{S}_{\sh{O}_Y}(\sh{F})$ của các $\sh{O}_Y$-đại số phân bậc \sref{II.1.7.4}.
Nếu $u$ \emph{toàn ánh}, thì $\bb{S}(u)$ cũng toàn ánh, và do đó \sref{II.3.6.2}[(i)] $\Proj(\bb{S}(u))$ là một \emph{phép nhúng đóng} $\bb{P}(\sh{F})\to\bb{P}(\sh{E})$, mà ta ký hiệu bởi $\bb{P}(u)$.
Vì vậy, có thể nói rằng $\bb{P}(\sh{E})$ là một \emph{hàm tử phản biến} theo $\sh{E}$, với điều kiện chỉ lấy các đồng cấu \emph{toàn ánh} làm cấu xạ giữa các $\sh{O}_Y$-môđun quasi-coherent.

Vẫn giả sử $u$ toàn ánh, và đặt $P=\bb{P}(\sh{E})$, $Q=\bb{P}(\sh{F})$, $j=\bb{P}(u)$, ta có, sai khác một đẳng cấu,
\[
\label{II.4.1.2.1}
  j^*(\sh{O}_P(n)) = \sh{O}_Q(n)
  \qquad\mbox{với mọi $n\in\bb{Z}$}
  \tag{4.1.2.1}
\]
như suy ra từ \sref{II.3.6.3}.
\end{env}

\begin{env}[4.1.3]
\label{II.4.1.3}
Bây giờ cho $\psi:Y'\to Y$ là một cấu xạ và đặt $\sh{E}'=\psi^*(\sh{E})$;
khi đó $\bb{S}_{\sh{O}_{Y'}}(\sh{E}')=\psi^*(\bb{S}_{\sh{O}_Y}(\sh{E}))$ \sref{II.1.7.5};
do đó, theo \sref{II.3.5.3}, ta có
\[
\label{II.4.1.3.1}
  \bb{P}(\psi^*(\sh{E})) = \bb{P}(\sh{E})\times_Y Y'
  \tag{4.1.3.1}
\]
sai khác một đẳng cấu chính tắc;
hơn nữa, hiển nhiên ta có
\[
  \psi^*((\bb{S}_{\sh{O}_Y}(\sh{E}))(n)) = (\bb{S}_{\sh{O}_{Y'}}(\sh{E}'))(n)
\]
với mọi $n\in\bb{Z}$; do đó, đặt $P=\bb{P}(\sh{E})$ và $P'=\bb{P}(\sh{E}')$, theo \sref{II.3.5.4}, ta có, sai khác một đẳng cấu,
\[
\label{II.4.1.3.2}
  \sh{O}_{P'}(n) = \sh{O}_P(n)\otimes_Y\sh{O}_{Y'}
  \qquad\mbox{với mọi $n\in\bb{Z}$.}
  \tag{4.1.3.2}
\]
\end{env}
\oldpage[II]{72}
\begin{proposition}[4.1.4]
\label{II.4.1.4}
Cho $\sh{L}$ là một $\sh{O}_Y$-môđun khả nghịch.
Với mọi $\sh{O}_Y$-môđun quasi-coherent $\sh{E}$, tồn tại một $Y$-đẳng cấu chính tắc $i_\sh{L}:\bb{P}(\sh{E})\simto\bb{P}(\sh{E}\otimes\sh{L})$;
hơn nữa, nếu đặt $P=\bb{P}(\sh{E})$ và $Q=\bb{P}(\sh{E}\otimes\sh{L})$, thì $i_\sh{L}^*(\sh{O}_Q(n))$ đẳng cấu chính tắc với $\sh{O}_P(n)\otimes_Y\sh{L}^{\otimes n}$ với mọi $n\in\bb{Z}$.
\end{proposition}

\begin{proof}
Trước hết, nhận xét rằng nếu $A$ là một vành, $E$ là một $A$-môđun và $L$ là một $A$-môđun \emph{tự do hạng một}, thì ta định nghĩa chính tắc một đồng cấu các $A$-môđun
\[
  \bb{S}_n(E\otimes L) \to \bb{S}_n(E)\otimes L^{\otimes n}
\]
bằng cách cho $(x_1\otimes y_1)\ldots(x_n\otimes y_n)$ ứng với phần tử
\[
  (x_1x_2\ldots x_n)\otimes(y_1\otimes y_2\otimes\ldots\otimes y_n)
  \qquad\mbox{($x_i\in E$, $y_i\in L$, với $1\leq i\leq n$);}
\]
ta kiểm tra ngay (bằng cách quy về trường hợp $L=A$) rằng đồng cấu này thực ra là một đẳng cấu.
Từ đó suy ra một đẳng cấu chính tắc của các $A$-đại số phân bậc $\bb{S}_A(E\otimes L)\simto\bigoplus_{n\geq0}\bb{S}_n(E)\otimes L^{\otimes n}$.
Trở lại các điều kiện của \sref{II.4.1.4}, những nhận xét trên cho phép định nghĩa một đẳng cấu chính tắc của các $\sh{O}_Y$-đại số phân bậc
\[
\label{II.4.1.4.1}
  \bb{S}_{\sh{O}_Y}(\sh{E}\otimes_{\sh{O}_Y}\sh{L}) \simto \bigoplus_{n\geq0}\bb{S}_n(\sh{E})\otimes_{\sh{O}_Y}\sh{L}^{\otimes n}
  \tag{4.1.4.1}
\]
bằng cách định nghĩa đẳng cấu này như một đẳng cấu các tiền bó và có xét đến \sref{II.1.7.4}, \sref[I]{I.1.3.9} và \sref[I]{I.1.3.12}.
Khi đó mệnh đề suy ra từ \sref{II.3.1.8}[(iii)] và \sref{II.3.2.10}.
\end{proof}

\begin{env}[4.1.5]
\label{II.4.1.5}
Với các giả thiết của \sref{II.4.1.1}, đặt $P=\bb{P}(\sh{E})$ và ký hiệu bởi $p$ cấu xạ cấu trúc $P\to Y$.
Vì theo định nghĩa $\sh{E}=(\bb{S}_{\sh{O}_Y}(\sh{E}))_1$, ta có một đồng cấu chính tắc $\alpha_1:\sh{E}\to p_*(\sh{O}_P(1))$ \sref{II.3.3.2.2}, và do đó cũng có \sref[0]{0.4.4.3} một đồng cấu chính tắc
\[
\label{II.4.1.5.1}
  \alpha_1^\sharp: p^*(\sh{E}) \to \sh{O}_P(1).
  \tag{4.1.5.1}
\]
\end{env}

\begin{proposition}[4.1.6]
\label{II.4.1.6}
Đồng cấu chính tắc \sref{II.4.1.5.1} là toàn ánh.
\end{proposition}

\begin{proof}
Trong \sref{II.3.3.2}, ta đã thấy $\alpha_1^\sharp$ tương ứng có tính hàm tử với đồng cấu chính tắc $\sh{E}\otimes_{\sh{O}_Y}\bb{S}_{\sh{O}_Y}(\sh{E})\to(\bb{S}_{\sh{O}_Y}(\sh{E}))(1)$;
vì theo định nghĩa $\sh{E}$ sinh $\bb{S}_{\sh{O}_Y}(\sh{E})$, đồng cấu này toàn ánh, nên kết luận suy ra từ \sref{II.3.2.4}.
\end{proof}


\subsection{Các cấu xạ từ một tiền lược đồ vào một phân thớ xạ ảnh}
\label{subsection:II.4.2}

\begin{env}[4.2.1]
\label{II.4.2.1}
Giữ các ký hiệu của \sref{II.4.1.5}, cho $X$ là một $Y$-tiền lược đồ, $q:X\to Y$ là cấu xạ cấu trúc, và $r:X\to P$ là một $Y$-\emph{cấu xạ}, sao cho biểu đồ sau giao hoán:
\[
  \xymatrix{
    P \ar[d]_p & X \ar[l]_r \ar[dl]^q
  \\Y
  }
\]
\oldpage[II]{73}

Vì hàm tử $r^*$ khớp phải \sref[0]{0.4.3.1}, từ đồng cấu trong \sref{II.4.1.5.1}, vốn toàn ánh theo \sref{II.4.1.6}, ta suy ra một đồng cấu toàn ánh
\[
  r^*(\alpha_1^\sharp): r^*(p^*(\sh{E})) \to r^*(\sh{O}_P(1)).
\]

Nhưng $r^*(p^*(\sh{E}))=q^*(\sh{E})$, và $r^*(\sh{O}_P(1))$ đẳng cấu địa phương với $r^*(\sh{O}_P)=\sh{O}_X$; nói cách khác, nó là một $\sh{O}_X$-môđun \emph{khả nghịch} $\sh{L}_r$.
Như vậy, từ $r$ ta đã định nghĩa một đồng cấu $\sh{O}_X$-môđun chính tắc toàn ánh
\[
\label{II.4.2.1.1}
  \vphi_r:q^*(\sh{E}) \to \sh{L}_r.
  \tag{4.2.1.1}
\]

Khi $Y=\Spec(A)$ là affine và $\sh{E}=\widetilde{E}$, ta có thể mô tả đồng cấu này tường minh hơn như sau.
Với $f\in E$, trước hết từ \sref{II.2.6.3} suy ra
\[
\label{II.4.2.1.2}
  r^{-1}(D_+(f)) = X_{\vphi_r^\flat(f)}.
  \tag{4.2.1.2}
\]

Mặt khác, cho $V$ là một tập mở affine của $X$ chứa trong $r^{-1}(D_+(f))$, và $B$ là vành của nó, vốn là một $A$-đại số;
đặt $S=\bb{S}_A(E)$;
hạn chế của $r$ lên $V$ tương ứng với một đồng cấu $A$-đại số $\omega:S_{(f)}\to B$, và ta có $q^*(\sh{E})|V=(E\otimes_A B)\supertilde$ cùng $\sh{L}_r|V=\widetilde{L_r}$, trong đó $L_r=(S(1))_{(f)}\otimes_{S_{(f)}}B_{[\omega]}$ \sref[I]{I.1.6.5}.
Hạn chế của $\vphi_r$ lên $q^*(\sh{E})|V$ tương ứng với đồng cấu $B$-môđun $u:E\otimes_A B\to L_r$ biến $x\otimes1$ thành $(x/1)\otimes1=(f/1)\otimes\omega(x/f)$.
Mở rộng chính tắc của $\vphi_r$ thành một đồng cấu các $\sh{O}_X$-đại số
\[
  \psi_r: q^*(\bb{S}(\sh{E})) = \bb{S}(q^*(\sh{E})) \to \bb{S}(\sh{L}_r) = \bigoplus_{n\geq0}\sh{L}_r^{\otimes n}
\]
do đó có tính chất: hạn chế của nó lên $q^*(\bb{S}_n(\sh{E}))|V$ tương ứng với đồng cấu $\bb{S}_n(E\otimes_A B)=\bb{S}_n(E)\otimes_A B\to L_r^{\otimes n}$ biến $s\otimes1$ thành $(f/1)^{\otimes n}\otimes\omega(s/f^n)$.
\end{env}

\begin{env}[4.2.2]
\label{II.4.2.2}
Ngược lại, giả sử đã cho một cấu xạ $q:X\to Y$, một $\sh{O}_X$-môđun khả nghịch $\sh{L}$ và một $\sh{O}_Y$-môđun quasi-coherent $\sh{E}$;
với mỗi đồng cấu $\vphi:q^*(\sh{E})\to\sh{L}$, ứng một cách chính tắc một đồng cấu các $\sh{O}_X$-đại số quasi-coherent
\[
  \psi: \bb{S}(q^*(\sh{E})) = q^*(\bb{S}(\sh{E})) \to \bigoplus_{n\geq0}\sh{L}^{\otimes n},
\]
và do đó \sref{II.3.7.1} một $Y$-cấu xạ $r_{\sh{L},\psi}:G(\psi)\to\Proj(\bb{S}(\sh{E}))=\bb{P}(\sh{E})$, mà ta ký hiệu bởi $r_{\sh{L},\vphi}$.
Nếu $\vphi$ \emph{toàn ánh}, thì $\psi$ cũng toàn ánh, nên \sref{II.3.7.4} $r_{\sh{L},\vphi}$ được \emph{xác định ở mọi nơi}.
Hơn nữa, với các ký hiệu của \sref{II.4.2.1} và \sref{II.4.2.2}:
\end{env}

\begin{proposition}[4.2.3]
\label{II.4.2.3}
Cho một cấu xạ $q:X\to Y$ và một $\sh{O}_Y$-môđun quasi-coherent $\sh{E}$, các ánh xạ $r\mapsto(\sh{L}_r,\vphi_r)$ và $(\sh{L},\vphi)\mapsto r_{\sh{L},\vphi}$ thiết lập một tương ứng song ánh giữa tập các $Y$-cấu xạ $r:X\to\bb{P}(\sh{E})$ và tập các lớp tương đương của những cặp $(\sh{L},\vphi)$ gồm một $\sh{O}_X$-môđun khả nghịch $\sh{L}$ và một đồng cấu toàn ánh $\vphi:q^*(\sh{E})\to\sh{L}$, trong đó hai cặp $(\sh{L},\vphi)$ và $(\sh{L}',\vphi')$ tương đương nếu tồn tại một $\sh{O}_X$-đẳng cấu $\tau:\sh{L}\simto\sh{L}'$ sao cho $\vphi'=\tau\circ\vphi$.
\end{proposition}

\begin{proof}
Trước hết, xuất phát từ một $Y$-cấu xạ $r:X\to\bb{P}(\sh{E})$, ta lập $\sh{L}_r$ và $\vphi_r$ \sref{II.4.2.1}, rồi đặt $r'=r_{\sh{L}_r,\vphi_r}$;
từ \sref{II.4.2.1} và \sref{II.3.7.2} suy ra ngay rằng các cấu xạ $r$ và $r'$ đồng nhất (khi lấy phần tử $(f/1)\otimes1$ làm phần tử sinh của $\sh{L}_r$ để định nghĩa các đồng cấu $v_n$ trong \sref{II.3.7.2}).
Ngược lại, xuất phát từ một cặp $(\sh{L},\vphi)$, ta lập
\oldpage[II]{74}
$r''=r_{\sh{L},\vphi}$, rồi lập $\sh{L}_{r''}$ và $\vphi_{r''}$;
ta sẽ chứng minh rằng tồn tại một đẳng cấu chính tắc $\tau:\sh{L}_{r''}\simto\sh{L}$ sao cho $\vphi=\tau\circ\vphi_{r''}$;
để định nghĩa nó, ta có thể quy về trường hợp $Y=\Spec(A)$ và $X=\Spec(B)$ là affine, rồi (với các ký hiệu của \sref{II.4.2.1} và \sref{II.3.7.2}) cho mỗi phần tử $(x/1)\otimes1$ của $L_{r''}$ (với $x\in E$) ứng với phần tử $v_1(x)c$ của $L$.
Ta kiểm tra ngay rằng $\tau$ không phụ thuộc vào phần tử sinh $c$ đã chọn của $L$;
vì $v_1$ toàn ánh theo giả thiết, để chứng minh $\tau$ là một đẳng cấu, chỉ cần chứng minh rằng nếu $x/1=0$ trong $(S(1))_{(f)}$, thì $v_1(x)/1=0$ trong $B_g$;
nhưng đẳng thức thứ nhất có nghĩa là $f^nx=0$ trong $\bb{S}_{n+1}(E)$ với một $n$ nào đó, và từ đó suy ra $v_{n+1}(f^nx)=g^nv_1(x)=0$ trong $B$, nên có kết luận.
Cuối cùng, hiển nhiên với hai cặp tương đương $(\sh{L},\vphi)$ và $(\sh{L}',\vphi')$, ta có $r_{\sh{L},\vphi}=r_{\sh{L}',\vphi'}$.
\end{proof}

Đặc biệt, với $X=Y$:
\begin{theorem}[4.2.4]
\label{II.4.2.4}
Tập các $Y$-tiết diện của $\bb{P}(\sh{E})$ tương ứng song ánh chính tắc với tập các $\sh{O}_Y$-môđun con quasi-coherent $\sh{F}$ của $\sh{E}$ sao cho $\sh{E}/\sh{F}$ khả nghịch.
\end{theorem}

Ta lưu ý rằng tính chất này tương ứng với định nghĩa cổ điển của ``không gian xạ ảnh'' như tập các siêu phẳng của một không gian véc tơ (trường hợp cổ điển tương ứng với $Y=\Spec(K)$, trong đó $K$ là một trường, và $\sh{E}=\widetilde{E}$, với $E$ là một $K$-không gian véc tơ hữu hạn chiều; khi đó các $\sh{F}$ có tính chất nêu trong \sref{II.4.2.4} tương ứng với các siêu phẳng của $E$, và mặt khác ta biết rằng các $Y$-tiết diện của $\bb{P}(\sh{E})$ khi ấy chính là \emph{các điểm của $\bb{P}(\sh{E})$ hữu tỉ trên $K$} \sref[I]{I.3.4.5}).

\begin{remark}[4.2.5]
\label{II.4.2.5}
Vì có một tương ứng song ánh chính tắc giữa các $Y$-cấu xạ từ $X$ vào $P$ và các cấu xạ đồ thị của chúng, tức các $X$-tiết diện của $P\times_Y X$ \sref[I]{I.3.3.14}, ta thấy rằng, theo chiều ngược lại, \sref{II.4.2.3} có thể suy ra từ \sref{II.4.2.4}.
Ký hiệu bởi $\Hyp_Y(X,\sh{E})$ tập các $\sh{O}_X$-môđun con quasi-coherent $\sh{F}$ của $\sh{E}\otimes_Y\sh{O}_X=q^*(\sh{E})$ sao cho $q^*(\sh{E})/\sh{F}$ là một $\sh{O}_X$-môđun khả nghịch.
Nếu $g:X'\to X$ là một $Y$-cấu xạ, thì từ tính khớp phải của $g^*$ suy ra
\[
  g^*(q^*(\sh{E})/\sh{F})=g^*q^*(\sh{E})/g^*(\sh{F}),
\]
nên bó sau là khả nghịch. Do đó, $X\mapsto\Hyp_Y(X,\sh{E})$ là một \emph{hàm tử phản biến} trên phạm trù các $Y$-tiền lược đồ.
Vì vậy, ta có thể diễn giải Định lý~\sref{II.4.2.4} như việc xác định một \emph{đẳng cấu chính tắc} giữa hai hàm tử $X\mapsto\Hom_Y(X,\bb{P}(\sh{E}))$ và $X\mapsto\Hyp_Y(X,\sh{E})$, cả hai đều phản biến theo $X$ trong phạm trù các $Y$-tiền lược đồ.
Điều này cũng cho một đặc trưng của phân thớ xạ ảnh $P=\bb{P}(\sh{E})$ bằng \emph{tính chất phổ dụng} sau, gần với trực giác hình học hơn các phép dựng trong các mục~\textsection\textsection\hyperref[section:II.2]{2} và \hyperref[section:II.3]{3}:
với mọi cấu xạ $q:X\to Y$ và mọi $\sh{O}_X$-môđun khả nghịch $\sh{L}$ là một môđun thương của $\sh{E}\otimes_{\sh{O}_Y}\sh{O}_X$, tồn tại một $Y$-cấu xạ duy nhất $r:X\to P$ sao cho $\sh{L}=r^*(\sh{O}_P(1))$.

Về sau ta sẽ thấy, theo cùng cách ấy, có thể định nghĩa, trong số những đối tượng khác, các lược đồ ``Grassmann''.
\end{remark}

\begin{corollary}[4.2.6]
\label{II.4.2.6}
Giả sử mọi $\sh{O}_Y$-môđun khả nghịch đều tầm thường \sref[I]{I.2.4.8}.
Cho $V$ là nhóm $\Hom_{\sh{O}_Y}(\sh{E},\sh{O}_Y)$, được xem như một môđun trên vành $A=\Gamma(Y,\sh{O}_Y)$, và cho $V^*$ là tập con của $V$ gồm các đồng cấu toàn ánh.
Khi đó tập các $Y$-tiết diện của $\bb{P}(\sh{E})$ được đồng nhất chính tắc với $V^*/A^*$, trong đó $A^*$ là nhóm các đơn vị của $A$.
\end{corollary}
\oldpage[II]{75}
Đặc biệt:
\begin{enumerate}
  \item Hệ quả \sref{II.4.2.6} áp dụng khi $Y$ là một \emph{lược đồ địa phương} \sref[I]{I.2.4.8}.
    Cho $Y$ là một tiền lược đồ bất kỳ, $y$ là một điểm của $Y$, và $Y'=\Spec(\kres(y))$;
    khi đó thớ $p^{-1}(y)$ của $\bb{P}(\sh{E})$ được đồng nhất, theo \sref{II.4.1.3.1}, với $\bb{P}(\sh{E}^y)$, trong đó $\sh{E}^y=\sh{E}_y\otimes_{\sh{O}_y}\kres(y)=\sh{E}_y/\mathfrak{m}_y\sh{E}_y$ được xét như một không gian vectơ trên $\kres(y)$.
    Tổng quát hơn, nếu $K$ là một mở rộng của $\kres(y)$, thì $p^{-1}(y)\otimes_{\kres(y)}K$ được đồng nhất với $\bb{P}(\sh{E}^y\otimes_{\kres(y)}K)$.
    Do đó, Hệ quả~\sref{II.4.2.6} cho thấy tập hợp các \emph{điểm hình học của $\bb{P}(\sh{E})$ nhận giá trị trong mở rộng $K$ của $\kres(y)$} \sref[I]{I.3.4.5}, còn có thể gọi là \emph{thớ hình học hữu tỉ trên $K$ của $\bb{P}(\sh{E})$ bên trên điểm $y$}, được đồng nhất với \emph{không gian xạ ảnh} liên kết với \emph{đối ngẫu} của $K$-không gian vectơ $\sh{E}^y\otimes_{\kres(y)}K$.
  \item Giả sử $Y$ affine có vành $A$, và hơn nữa mọi $\sh{O}_Y$-môđun khả nghịch đều tầm thường;
    lấy thêm $\sh{E}=\sh{O}_Y^n$;
    khi đó, trong \sref{II.4.2.6}, $V$ được đồng nhất với $A^n$ \sref[I]{I.1.3.8}, còn $V^*$ được đồng nhất với tập hợp các hệ $(f_i)_{1\leq i\leq n}$ gồm những phần tử của $A$ sinh iđêan $A$;
    hai hệ như vậy xác định cùng một $Y$-tiết diện của $\bb{P}_Y^{n-1}=\bb{P}_A^{n-1}$, nói cách khác, \emph{cùng một điểm của $\bb{P}_A^{n-1}$ nhận giá trị trong $A$}, khi và chỉ khi hệ này thu được từ hệ kia bằng cách nhân với một phần tử khả nghịch của $A$.
\end{enumerate}

Các tính chất này biện minh cho thuật ngữ «phân thớ xạ ảnh» đối với $\bb{P}(\sh{E})$.
Ta lưu ý rằng định nghĩa «không gian xạ ảnh» thu được như vậy thực ra \emph{đối ngẫu} với định nghĩa cổ điển;
điều này bị áp đặt bởi yêu cầu phải định nghĩa được $\bb{P}(\sh{E})$ cho một $\sh{O}_Y$-môđun quasi-coherent \emph{bất kỳ} $\sh{E}$, không nhất thiết tự do địa phương.

\begin{remark}[4.2.7]
\label{II.4.2.7}
Trong Chương~V, ta sẽ thấy rằng, nếu $Y$ khác rỗng, Noether địa phương và liên thông, còn $\sh{E}$ tự do địa phương có hạng ít nhất hai, thì, đặt $P=\bb{P}(\sh{E})$, mọi $\sh{O}_P$-môđun khả nghịch $\sh{L}$ đều đẳng cấu với một $\sh{O}_P$-môđun dạng $\sh{L}'\otimes_{\sh{O}_Y}\sh{O}_P(m)$, trong đó $\sh{L}'$ là một $\sh{O}_Y$-môđun khả nghịch được xác định duy nhất sai khác một đẳng cấu, còn $m$ là một số nguyên được xác định duy nhất.
Nói cách khác, $\HH^1(P,\sh{O}_P^*)$ đẳng cấu với $\bb{Z}\times\HH^1(Y,\sh{O}_Y^*)$ \sref[0]{0.5.4.7}.
Ta cũng sẽ thấy (\sref[III]{III.2.1.15}, có xét đến \sref[0]{0.5.4.10}) rằng $p_*(\sh{L})=0$ nếu $m<0$, và $p_*(\sh{L})$ đẳng cấu với $\sh{L}'\otimes_{\sh{O}_Y}(\bb{S}_{\sh{O}_Y}(\sh{E}))_m$ nếu $m\geq0$.
Nếu $\sh{F}$ là một $\sh{O}_Y$-môđun quasi-coherent, thì mọi $Y$-cấu xạ $\bb{P}(\sh{E})\to\bb{P}(\sh{F})$ được xác định bởi việc cho một $\sh{O}_Y$-môđun khả nghịch $\sh{L}'$, một số nguyên $m\geq0$, và một $\sh{O}_Y$-đồng cấu $\psi:\sh{F}\to\sh{L}'\otimes_{\sh{O}_Y}(\bb{S}_{\sh{O}_Y}(\sh{E}))_m$ sao cho đồng cấu tương ứng $\psi^\sharp$ của các $\sh{O}_{\bb{P}(\sh{E})}$-môđun là toàn ánh.
Ta cũng sẽ thấy rằng, nếu $Y$-cấu xạ đang xét là một đẳng cấu, thì $m=1$ và $\sh{F}$ đẳng cấu với $\sh{E}\otimes_{\sh{O}_Y}\sh{L}'$ (mệnh đề đảo của \sref{II.4.1.4}).
Điều này sẽ cho phép xác định bó các mầm tự đẳng cấu của $\bb{P}(\sh{E})$ như bó thương của bó nhóm $\shAut(\sh{E})$ (đẳng cấu địa phương với $\GL(n,\sh{O}_Y)$ nếu $\sh{E}$ có hạng $n$) theo $\sh{O}_Y^*$.
\end{remark}

\begin{env}[4.2.8]
\label{II.4.2.8}
Giữ các ký hiệu của \sref{II.4.2.1}, cho $u:X'\to X$ là một cấu xạ;
nếu $Y$-cấu xạ $r:X\to P$ tương ứng với đồng cấu $\vphi:q^*(\sh{E})\to\sh{L}$, thì $Y$-cấu xạ $r\circ u$ tương ứng với $u^*(\vphi):u^*(q^*(\sh{E}))\to u^*(\sh{L})$, như suy ra ngay từ các định nghĩa.
\end{env}

\begin{env}[4.2.9]
\label{II.4.2.9}
Cho $\sh{E}$ và $\sh{F}$ là hai $\sh{O}_Y$-môđun quasi-coherent, $v:\sh{E}\to\sh{F}$ là một đồng cấu toàn ánh, và $j=\bb{P}(v)$ là phép nhúng đóng tương ứng $\bb{P}(\sh{F})\to\bb{P}(\sh{E})$ \sref{II.4.1.2}.
Nếu $Y$-cấu xạ $r:X\to\bb{P}(\sh{F})$ tương ứng với đồng cấu $\vphi:q^*(\sh{F})\to\sh{L}$, thì
\oldpage[II]{76}
$Y$-cấu xạ $j\circ r$ tương ứng với $q^*(\sh{E})\xrightarrow{q^*(v)}q^*(\sh{F})\xrightarrow{\vphi}\sh{L}$;
điều này cũng suy ra từ định nghĩa trong \sref{II.4.2.1}.
\end{env}

\begin{env}[4.2.10]
\label{II.4.2.10}
Cho $\psi:Y'\to Y$ là một cấu xạ, và đặt $\sh{E}'=\psi^*(\sh{E})$.
Nếu $Y$-cấu xạ $r:X\to P$ tương ứng với đồng cấu $\vphi:q^*(\sh{E})\to\sh{L}$, thì $Y'$-cấu xạ
\[
  r_{(Y')}:X_{(Y')}\to P'=\bb{P}(\sh{E}')
\]
tương ứng với $\vphi_{(Y')}:q_{(Y')}^*(\sh{E}')=q^*(\sh{E})\otimes_Y\sh{O}_{Y'}\to\sh{L}\otimes_Y\sh{O}_{Y'}$.
Thật vậy, theo \sref{II.4.1.3.1}, ta có biểu đồ giao hoán
\[
  \xymatrix{
    Y' \ar[d]
    & P'=P_{(Y')} \ar[l]_{p_{(Y')}} \ar[d]^{u}
    & X_{(Y')} \ar[l]_{r_{(Y')}} \ar[d]^{v}
  \\Y
    & P \ar[l]_{p}
    & X \ar[l]_{r}
  }
\]

Theo \sref{II.4.1.3.2}, ta có
\[
  (r_{(Y')})^*(\sh{O}_{P'}(1))=(r_{(Y')})^*(u^*(\sh{O}_P(1)))=v^*(r^*(\sh{O}_P(1)))=v^*(\sh{L})=\sh{L}\otimes_Y\sh{O}_{Y'};
\]
mặt khác, $u^*(\alpha_1^\sharp)$ chính là đồng cấu chính tắc $\alpha_1^\sharp:(p_{(Y')})^*(\sh{E}')\to\sh{O}_{P'}(1)$;
điều này thấy được bằng cách viết tường minh các đồng cấu chính tắc $\alpha_1^\sharp$ trong $P$ và $P'$ như ở \sref{II.4.1.6}.
Từ đó suy ra khẳng định.
\end{env}

\subsection{Cấu xạ Segre}
\label{subsection:II.4.3}

\begin{env}[4.3.1]
\label{II.4.3.1}
Cho $Y$ là một tiền lược đồ, $\sh{E}$ và $\sh{F}$ là hai $\sh{O}_Y$-môđun quasi-coherent;
đặt $P_1=\bb{P}(\sh{E})$ và $P_2=\bb{P}(\sh{F})$, và ký hiệu bởi $p_1:P_1\to Y$ và $p_2:P_2\to Y$ các cấu xạ cấu trúc.
Đặt $Q=P_1\times_Y P_2$, và cho $q_1:Q\to P_1$ cùng $q_2:Q\to P_2$ là các phép chiếu chính tắc;
$\sh{O}_Q$-môđun $\sh{L}=\sh{O}_{P_1}(1)\otimes_Y\sh{O}_{P_2}(1)=q_1^*(\sh{O}_{P_1}(1))\otimes_{\sh{O}_Q}q_2^*(\sh{O}_{P_2}(1))$ là khả nghịch vì là tích tenxơ của hai $\sh{O}_Q$-môđun khả nghịch \sref[0]{0.5.4.4}.
Mặt khác, nếu $r=p_1\circ q_1=p_2\circ q_2$ là cấu xạ cấu trúc $Q\to Y$, thì $r^*(\sh{E}\otimes_{\sh{O}_Y}\sh{F})=q_1^*(p_1^*(\sh{E}))\otimes_{\sh{O}_Q}q_2^*(p_2^*(\sh{F}))$ \sref[0]{0.4.3.3};
do đó, bằng cách lấy tích tenxơ, các đồng cấu chính tắc toàn ánh \sref{II.4.1.5.1}, $p_1^*(\sh{E})\to\sh{O}_{P_1}(1)$ và $p_2^*(\sh{F})\to\sh{O}_{P_2}(1)$, cho một đồng cấu chính tắc
\[
\label{II.4.3.1.1}
  s:r^*(\sh{E}\otimes_{\sh{O}_Y}\sh{F})\to\sh{L}
  \tag{4.3.1.1}
\]
hiển nhiên toàn ánh;
do đó, từ \sref{II.4.2.2}, ta suy ra một cấu xạ chính tắc, gọi là \emph{cấu xạ Segre}:
\[
\label{II.4.3.1.2}
  \varsigma:\bb{P}(\sh{E})\times_Y\bb{P}(\sh{F})\to\bb{P}(\sh{E}\otimes_{\sh{O}_Y}\sh{F}).
  \tag{4.3.1.2}
\]

Ta mô tả tường minh cấu xạ $\varsigma$ khi $Y=\Spec(A)$ là affine, $\sh{E}=\widetilde{E}$ và $\sh{F}=\widetilde{F}$, trong đó $E$ và $F$ là hai $A$-môđun, nên $\sh{E}\otimes_{\sh{O}_Y}\sh{F}=(E\otimes_A F)^\sim$ \sref[I]{I.1.3.12};
đặt $R=\bb{S}_A(E)$, $S=\bb{S}_A(F)$ và $T=\bb{S}_A(E\otimes_A F)$;
cho $f\in E$, $g\in F$, và xét tập mở affine
\[
  D_+(f)\times_YD_+(g)=\Spec(B)
\]
\oldpage[II]{77}
của $Q$, trong đó $B=R_{(f)}\otimes_A S_{(g)}$;
hạn chế của $\sh{L}$ trên tập mở affine này là $\widetilde{L}$, trong đó
\[
  L=(R(1))_{(f)}\otimes_A(S(1))_{(g)},
\]
và phần tử $c=(f/1)\otimes(g/1)$ là một phần tử sinh của $L$ khi môđun này được xét như một $B$-môđun tự do \sref{II.2.5.7}.
Đồng cấu \sref{II.4.3.1.1} tương ứng với đồng cấu
\[
  (x\otimes y)\otimes b\longmapsto b((x/1)\otimes(y/1))
\]
từ $(E\otimes_A F)\otimes_A B$ vào $L$.
Với các ký hiệu của \sref{II.3.7.2}, do đó ta có $v_1(x\otimes y)=(x/f)\otimes(y/g)$;
hạn chế của $\varsigma$ trên $D_+(f)\times_YD_+(g)$ là một cấu xạ từ lược đồ affine này vào $D_+(f\otimes g)$, tương ứng với đồng cấu vành $\omega:T_{(f\otimes g)}\to R_{(f)}\otimes_A S_{(g)}$ sao cho
\[
\label{II.4.3.1.3}
  \omega((x\otimes y)/(f\otimes g))=(x/f)\otimes(y/g)
  \tag{4.3.1.3}
\]
với $x\in E$ và $y\in F$.
\end{env}

\begin{env}[4.3.2]
\label{II.4.3.2}
Từ \sref{II.4.2.3} suy ra rằng ta có một đẳng cấu chính tắc
\[
\label{II.4.3.2.1}
  \tau:\varsigma^*(\sh{O}_P(1))\simto\sh{O}_{P_1}(1)\otimes_Y\sh{O}_{P_2}(1)
  \tag{4.3.2.1}
\]
trong đó ta đặt $P=\bb{P}(\sh{E}\otimes_{\sh{O}_Y}\sh{F})$.
Ta chứng minh rằng, với $x\in\Gamma(Y,\sh{E})$ và $y\in\Gamma(Y,\sh{F})$, ta có
\[
\label{II.4.3.2.2}
  \tau(\alpha_1(x\otimes y))=\alpha_1(x)\otimes\alpha_1(y).
  \tag{4.3.2.2}
\]

Thật vậy, ta quy về trường hợp $Y$ affine, và khi đó, với các ký hiệu của \sref{II.4.3.1} và \sref{II.2.6.2}, ta có $\alpha_1^{f\otimes g}(x\otimes y)=(x\otimes y)/1$ trong $(T(1))_{(f\otimes g)}$, $\alpha_1^f(x)=x/1$ trong $(R(1))_{(f)}$, và $\alpha_1^g(y)=y/1$ trong $(S(1))_{(g)}$.
Định nghĩa của $\tau$ trong \sref{II.4.2.3} và phép tính $v_1$ trong \sref{II.4.3.1} chứng minh ngay \sref{II.4.3.2.2}.
Từ đó suy ra công thức
\[
\label{II.4.3.2.3}
  \varsigma^{-1}(P_{x\otimes y})=(P_1)_x\times_Y(P_2)_y
  \tag{4.3.2.3}
\]
với các ký hiệu của \sref{II.3.1.4}.
Thật vậy, có xét đến \sref{II.3.3.3}, công thức \sref{II.4.3.2.2} quy việc cần chứng minh (bằng cách trở lại trường hợp affine, với sự trợ giúp của \sref[I]{I.3.2.7} và \sref[I]{I.3.2.3}) về bổ đề sau:

\begin{lemma}[4.3.2.4]
\label{II.4.3.2.4}
Cho $B$ và $B'$ là hai $A$-đại số, và đặt $Y=\Spec(A)$, $Z=\Spec(B)$, $Z'=\Spec(B')$;
với mọi $t\in B$ và $t'\in B'$, ta có $D(t\otimes t')=D(t)\times_YD(t')$.
\end{lemma}

\begin{proof}
Thật vậy, nếu $p:Z\times_YZ'\to Z$ và $p':Z\times_YZ'\to Z'$ là các phép chiếu chính tắc, thì từ \sref[I]{I.1.2.2.2} suy ra $p^{-1}(D(t))=D(t\otimes1)$ và ${p'}^{-1}(D(t'))=D(1\otimes t')$;
kết luận suy ra từ \sref[I]{I.3.2.7} và \sref[I]{I.1.1.9.1}, vì $(t\otimes1)(1\otimes t')=t\otimes t'$.
\end{proof}
\end{env}

\begin{proposition}[4.3.3]
\label{II.4.3.3}
Cấu xạ Segre là một phép nhúng đóng.
\end{proposition}

\begin{proof}
Vì vấn đề mang tính địa phương trên $Y$, ta quy về trường hợp $Y$ affine.
Với các ký hiệu của \sref{II.4.3.1} và \sref{II.4.3.2}, các $D_+(f\otimes g)$ khi đó lập thành một cơ sở của tô pô trên $P$, vì các $f\otimes g$ sinh $T$ khi $f$ chạy qua $E$ và $g$ chạy qua $F$.
Mặt khác, theo \sref{II.4.3.2.3}, ta có $\varsigma^{-1}(D_+(f\otimes g))=D_+(f)\times_YD_+(g)$.
Do đó, theo \sref[I]{I.4.2.4}, chỉ cần chứng minh rằng hạn chế của $\varsigma$ trên $D_+(f)\times_YD_+(g)$ là một phép nhúng đóng vào $D_+(f\otimes g)$.
Nhưng, với cùng các ký hiệu, công thức \sref{II.4.3.1.3} cho thấy $\omega$ \emph{toàn ánh}, điều này hoàn tất chứng minh.
\end{proof}

\begin{env}[4.3.4]
\label{II.4.3.4}
Cấu xạ Segre có tính hàm tử theo $\sh{E}$ và $\sh{F}$ khi ta chỉ xét các
\oldpage[II]{78}
đồng cấu \emph{toàn ánh} của các $\sh{O}_Y$-môđun quasi-coherent.
Thật vậy, cần chứng minh rằng, nếu $\sh{E}\to\sh{E}'$ là một $\sh{O}_Y$-đồng cấu toàn ánh, thì biểu đồ
\[
  \xymatrix{
    \bb{P}(\sh{E}')\times_Y\bb{P}(\sh{F}) \ar[r]^{j\times1} \ar[d]_{\varsigma}
    & \bb{P}(\sh{E})\times_Y\bb{P}(\sh{F}) \ar[d]^{\varsigma}
  \\ \bb{P}(\sh{E}'\otimes_{\sh{O}_Y}\sh{F}) \ar[r]
    & \bb{P}(\sh{E}\otimes_{\sh{O}_Y}\sh{F})
  }
\]
giao hoán, trong đó $j$ ký hiệu phép nhúng đóng chính tắc $\bb{P}(\sh{E}')\to\bb{P}(\sh{E})$.
Đặt $P'_1=\bb{P}(\sh{E}')$ và giữ các ký hiệu còn lại của \sref{II.4.3.1};
$j\times1$ là một phép nhúng đóng \sref[I]{I.4.3.1}, và, sai khác các đẳng cấu, ta có
\[
  (j\times1)^*(\sh{O}_{P_1}(1)\otimes\sh{O}_{P_2}(1))
  =j^*(\sh{O}_{P_1}(1))\otimes\sh{O}_{P_2}(1)
  =\sh{O}_{P'_1}(1)\otimes\sh{O}_{P_2}(1)
\]
theo \sref{II.4.1.2.1} và \sref[I]{I.9.1.5};
do đó khẳng định suy ra từ \sref{II.4.2.8} và \sref{II.4.2.9}.
\end{env}

\begin{env}[4.3.5]
\label{II.4.3.5}
Với các ký hiệu của \sref{II.4.3.1}, cho $\psi:Y'\to Y$ là một cấu xạ, và đặt $\sh{E}'=\psi^*(\sh{E})$, $\sh{F}'=\psi^*(\sh{F})$;
khi đó cấu xạ Segre $\bb{P}(\sh{E}')\times\bb{P}(\sh{F}')\to\bb{P}(\sh{E}'\otimes\sh{F}')$ được đồng nhất với $\varsigma_{(Y')}$.
Thật vậy, giữ các ký hiệu của \sref{II.4.3.1}, đặt thêm $P'_1=\bb{P}(\sh{E}')$ và $P'_2=\bb{P}(\sh{F}')$;
ta biết theo \sref{II.4.1.3.1} rằng $P'_i$ được đồng nhất với $(P_i)_{(Y')}$ ($i=1,2$), nên cấu xạ cấu trúc $P'_1\times P'_2\to Y'$ được đồng nhất với $r_{(Y')}$.
Mặt khác, $\sh{E}'\otimes\sh{F}'$ được đồng nhất với $\psi^*(\sh{E}\otimes\sh{F})$, nên $\bb{P}(\sh{E}'\otimes\sh{F}')$ được đồng nhất với $(\bb{P}(\sh{E}\otimes\sh{F}))_{(Y')}$ \sref{II.4.1.3.1}.
Cuối cùng, $\sh{O}_{P'_1}(1)\otimes_{Y'}\sh{O}_{P'_2}(1)=\sh{L}'$ được đồng nhất với $\sh{L}\otimes_Y\sh{O}_{Y'}$, theo \sref{II.4.1.3.2} và \sref[I]{I.9.1.11}.
Khi đó đồng cấu chính tắc $(r_{(Y')})^*(\sh{E}'\otimes\sh{F}')\to\sh{L}'$ được đồng nhất với $s_{(Y')}$, và khẳng định suy ra từ \sref{II.4.2.10}.
\end{env}

\begin{remark}[4.3.6]
\label{II.4.3.6}
Tiền lược đồ tổng của $\bb{P}(\sh{E})$ và $\bb{P}(\sh{F})$ cũng đẳng cấu chính tắc với một \emph{tiền lược đồ con đóng của $\bb{P}(\sh{E}\oplus\sh{F})$}.
Thật vậy, các đồng cấu toàn ánh $\sh{E}\oplus\sh{F}\to\sh{E}$ và $\sh{E}\oplus\sh{F}\to\sh{F}$ tương ứng với các phép nhúng đóng $\bb{P}(\sh{E})\to\bb{P}(\sh{E}\oplus\sh{F})$ và $\bb{P}(\sh{F})\to\bb{P}(\sh{E}\oplus\sh{F})$;
mọi việc quy về chứng minh rằng các không gian nền của những tiền lược đồ con đóng của $\bb{P}(\sh{E}\oplus\sh{F})$ thu được như vậy không có điểm chung.
Vì vấn đề mang tính địa phương trên $Y$, ta có thể dùng các ký hiệu của \sref{II.4.3.1};
nhưng $\bb{S}_n(E)$ và $\bb{S}_n(F)$ được đồng nhất với các môđun con của $\bb{S}_n(E\oplus F)$ có giao chỉ gồm $0$;
và nếu $\mathfrak{p}$ là một iđêan nguyên tố phân bậc của $\bb{S}(E)$ sao cho $\mathfrak{p}\cap\bb{S}_n(E)\neq\bb{S}_n(E)$ với mọi $n\geq0$, thì tương ứng với nó trong $\bb{S}(E\oplus F)$ là một iđêan nguyên tố phân bậc có vết trên $\bb{S}_n(E)$ bằng $\mathfrak{p}\cap\bb{S}_n(E)$, nhưng \emph{chứa} $\bb{S}_+(F)$, như thấy ngay;
do đó, hai điểm lần lượt thuộc $\Proj(\bb{S}(E))$ và $\Proj(\bb{S}(F))$ không thể có cùng ảnh trong $\Proj(\bb{S}(E\oplus F))$.
\end{remark}

\subsection{Các phép nhúng vào phân thớ xạ ảnh; các bó rất ample}
\label{subsection:II.4.4}

\begin{proposition}[4.4.1]
\label{II.4.4.1}
Cho $Y$ là một lược đồ quasi-compact, hoặc một tiền lược đồ có không gian nền là Noether, $q:X\to Y$ là một cấu xạ \emph{thuộc kiểu hữu hạn}, và $\sh{L}$ là một $\sh{O}_X$-môđun khả nghịch.
\begin{enumerate}
  \item[\rm{(i)}] Cho $\sh{S}$ là một $\sh{O}_Y$-đại số phân bậc quasi-coherent theo các bậc dương, và $\psi:q^*(\sh{S})\to\bigoplus_{n\geq0}\sh{L}^{\otimes n}$ là một đồng cấu các đại số phân bậc.
    Để $r_{\sh{L},\psi}$ được xác định khắp nơi và là một phép nhúng, điều cần và
\oldpage[II]{79}
đủ là tồn tại một số nguyên $n>0$ và một $\sh{O}_Y$-môđun con quasi-coherent \emph{hữu hạn sinh} $\sh{E}$ của $\sh{S}_n$ sao cho đồng cấu $\psi'=\psi_n\circ q^*(j):q^*(\sh{E})\to\sh{L}^{\otimes n}=\sh{L}'$ (trong đó $j$ là đơn ánh $\sh{E}\to\sh{S}_n$) toàn ánh và cấu xạ $r_{\sh{L}',\psi'}:X\to\bb{P}(\sh{E})$ là một phép nhúng.
  \item[\rm{(ii)}] Cho $\sh{F}$ là một $\sh{O}_Y$-môđun quasi-coherent và $\vphi:q^*(\sh{F})\to\sh{L}$ là một đồng cấu toàn ánh.
    Để cấu xạ $r_{\sh{L},\vphi}$ là một phép nhúng $X\to\bb{P}(\sh{F})$, điều cần và đủ là tồn tại một $\sh{O}_Y$-môđun con quasi-coherent \emph{hữu hạn sinh} $\sh{E}$ của $\sh{F}$ sao cho đồng cấu $\vphi'=\vphi\circ q^*(j):q^*(\sh{E})\to\sh{L}$ (trong đó $j:\sh{E}\to\sh{F}$ là đơn ánh chính tắc) toàn ánh và $r_{\sh{L},\vphi'}:X\to\bb{P}(\sh{E})$ là một phép nhúng.
\end{enumerate}
\end{proposition}

\begin{proof}
\medskip\noindent
\begin{enumerate}
  \item[\rm{(i)}] Việc $r_{\sh{L},\psi}$ được xác định khắp nơi và là một phép nhúng, theo \sref{II.3.8.5}, tương đương với sự tồn tại của $n>0$ và $\sh{E}$ sao cho, nếu $\sh{S}'$ là đại số con của $\sh{S}$ được sinh bởi $\sh{E}$, thì đồng cấu $q^*(\sh{E})\to\sh{L}^{\otimes n}$ toàn ánh và cấu xạ $r_{\sh{L}',\psi'}:X\to\Proj(\sh{S}')$ được xác định khắp nơi và là một phép nhúng.
    Mặt khác, ta có một đồng cấu chính tắc toàn ánh $\bb{S}(\sh{E})\to\sh{S}'$, tương ứng với một phép nhúng đóng $\Proj(\sh{S}')\to\bb{P}(\sh{E})$ \sref{II.3.6.2};
    từ đó suy ra kết luận.
  \item[\rm{(ii)}] Vì $\sh{F}$ là giới hạn quy nạp của các môđun con quasi-coherent hữu hạn sinh $\sh{E}_\lambda$ của nó \sref[I]{I.9.4.9}, nên $\bb{S}(\sh{F})$ là giới hạn quy nạp của các $\bb{S}(\sh{E}_\lambda)$;
    khi đó kết luận suy ra từ \sref{II.3.8.4}, bằng cách lưu ý rằng, trong chứng minh của \sref{II.3.8.4}, ta có thể lấy mọi $n_i$ bằng $1$:
    thật vậy (giả sử $Y$ affine), nếu $r=r_{\sh{L},\vphi}$ là một phép nhúng, thì $r(X)$ là một không gian con quasi-compact của $\bb{P}(\sh{F})$ có thể được phủ bởi một số hữu hạn các tập mở của $\bb{P}(\sh{F})$ dạng $D_+(f)$, với $f\in F$, sao cho $D_+(f)\cap r(X)$ đóng.
\end{enumerate}
\end{proof}

\begin{definition}[4.4.2]
\label{II.4.4.2}
Cho $Y$ là một tiền lược đồ và $q:X\to Y$ là một cấu xạ.
Ta nói một $\sh{O}_X$-môđun khả nghịch $\sh{L}$ là \emph{rất ample đối với $q$} (hoặc \emph{rất ample đối với $Y$}, hoặc đơn giản là \emph{rất ample} nếu không gây nhầm lẫn) nếu tồn tại một $\sh{O}_Y$-môđun quasi-coherent $\sh{E}$ và một $Y$-phép nhúng $i$ của $X$ vào $P=\bb{P}(\sh{E})$ sao cho $\sh{L}$ đẳng cấu với $i^*(\sh{O}_P(1))$.
\end{definition}

Theo \sref{II.4.2.3}, điều này tương đương với việc tồn tại một $\sh{O}_Y$-môđun quasi-coherent $\sh{E}$ và một đồng cấu \emph{toàn ánh} $\vphi:q^*(\sh{E})\to\sh{L}$ sao cho $r_{\sh{L},\vphi}:X\to\bb{P}(\sh{E})$ là một \emph{phép nhúng}.

Ta lưu ý rằng sự tồn tại của một $\sh{O}_X$-môđun rất ample đối với $Y$ kéo theo $q$ là \emph{tách} (\sref{II.3.1.3} và \sref[I]{I.5.5.1}[(i) và (ii)]).

\begin{corollary}[4.4.3]
\label{II.4.4.3}
Giả sử tồn tại một $\sh{O}_Y$-đại số phân bậc quasi-coherent $\sh{S}$ được sinh bởi $\sh{S}_1$, và một $Y$-phép nhúng $i:X\to P=\Proj(\sh{S})$ sao cho $\sh{L}$ đẳng cấu với $i^*(\sh{O}_P(1))$;
khi đó $\sh{L}$ rất ample đối với $q$.
\end{corollary}

\begin{proof}
Thật vậy, nếu $\sh{F}=\sh{S}_1$, thì đồng cấu chính tắc $\bb{S}(\sh{F})\to\sh{S}$ toàn ánh; do đó, bằng cách hợp thành phép nhúng đóng tương ứng $\Proj(\sh{S})\to\bb{P}(\sh{F})$ \sref{II.3.6.2} với phép nhúng $i$, ta thu được một phép nhúng $j:X\to\bb{P}(\sh{F})=P'$ sao cho $\sh{L}$ đẳng cấu với $j^*(\sh{O}_{P'}(1))$ \sref{II.3.6.3}.
\end{proof}

\begin{proposition}[4.4.4]
\label{II.4.4.4}
Cho $q:X\to Y$ là một cấu xạ quasi-compact và $\sh{L}$ là một $\sh{O}_X$-môđun khả nghịch.
Các tính chất sau là tương đương:
\begin{enumerate}
  \item[\rm{(a)}] $\sh{L}$ rất ample đối với $q$.
  \item[\rm{(b)}] $q_*(\sh{L})$ quasi-coherent, đồng cấu chính tắc $\sigma:q^*(q_*(\sh{L}))\to\sh{L}$ toàn ánh, và cấu xạ $r_{\sh{L},\sigma}:X\to\bb{P}(q_*(\sh{L}))$ là một phép nhúng.
\end{enumerate}
\end{proposition}

\begin{proof}
Vì $q$ quasi-compact, ta biết rằng $q_*(\sh{L})$ quasi-coherent nếu $q$ tách \sref[I]{I.9.2.2}.
\oldpage[II]{80}
Ta biết từ \sref{II.3.4.7} rằng sự tồn tại của một đồng cấu toàn ánh $\vphi:q^*(\sh{E})\to\sh{L}$ (với $\sh{E}$ là một $\sh{O}_Y$-môđun quasi-coherent) kéo theo $\sigma$ toàn ánh;
hơn nữa, ứng với phân tích $q^*(\sh{E})\to q^*(q_*(\sh{L}))\xrightarrow{\sigma}\sh{L}$ của $\vphi$ một cách chính tắc là một phân tích
\[
  q^*(\bb{S}(\sh{E}))\to q^*(\bb{S}(q_*(\sh{L})))\to\bigoplus_{n\geq0}\sh{L}^{\otimes n}
\]
do đó, theo \sref{II.3.8.3}, giả thiết rằng $r_{\sh{L},\vphi}$ là một phép nhúng kéo theo $j=r_{\sh{L},\sigma}$ cũng là một phép nhúng;
hơn nữa, theo \sref{II.4.2.4}, $\sh{L}$ đẳng cấu với $j^*(\sh{O}_{P'}(1))$, trong đó $P'=\bb{P}(q_*(\sh{L}))$.
Vậy (a) và (b) tương đương.
\end{proof}

\begin{corollary}[4.4.5]
\label{II.4.4.5}
Giả sử $q$ quasi-compact.
Điều kiện cần và đủ để $\sh{L}$ rất ample đối với $Y$ là tồn tại một phủ mở $(U_\alpha)$ của $Y$ sao cho $\sh{L}|_{q^{-1}(U_\alpha)}$ rất ample đối với $U_\alpha$, với mọi $\alpha$.
\end{corollary}

\begin{proof}
Thật vậy, điều kiện~(b) của \sref{II.4.4.4} có tính địa phương trên $Y$.
\end{proof}

\begin{proposition}[4.4.6]
\label{II.4.4.6}
Cho $Y$ là một lược đồ quasi-compact, hoặc một tiền lược đồ có không gian nền Noether, $q:X\to Y$ là một cấu xạ \emph{thuộc kiểu hữu hạn}, và $\sh{L}$ là một $\sh{O}_X$-môđun khả nghịch.
Khi đó các điều kiện~(a) và~(b) của \sref{II.4.4.4} tương đương với các điều kiện sau:
\begin{enumerate}
  \item[\rm{(a')}] Tồn tại một $\sh{O}_Y$-môđun quasi-coherent hữu hạn sinh $\sh{E}$ và một đồng cấu toàn ánh $\vphi:q^*(\sh{E})\to\sh{L}$ sao cho $r_{\sh{L},\vphi}$ là một phép nhúng.
  \item[\rm{(b')}] Tồn tại một $\sh{O}_Y$-môđun con quasi-coherent hữu hạn sinh $\sh{E}$ của $q_*(\sh{L})$ có các tính chất nêu trong~(a').
\end{enumerate}
\end{proposition}

\begin{proof}
Hiển nhiên (a') hoặc (b') kéo theo (a);
mặt khác, (a) kéo theo (a') theo \sref{II.4.4.1}, và tương tự, (b) kéo theo (b').
\end{proof}

\begin{corollary}[4.4.7]
\label{II.4.4.7}
Giả sử $Y$ là một lược đồ quasi-compact hoặc một tiền lược đồ Noether.
Nếu $\sh{L}$ rất ample đối với $q$, thì tồn tại một $\sh{O}_Y$-đại số phân bậc quasi-coherent $\sh{S}$ sao cho $\sh{S}_1$ hữu hạn sinh và sinh $\sh{S}$, cùng một $Y$-phép nhúng \emph{mở trội} $i:X\to P=\Proj(\sh{S})$ sao cho $\sh{L}$ đẳng cấu với $i^*(\sh{O}_P(1))$.
\end{corollary}

\begin{proof}
Thật vậy, điều kiện~(b') của \sref{II.4.4.6} được thỏa mãn;
khi đó cấu xạ cấu trúc $p:\bb{P}(\sh{E})=P'\to Y$ là tách và thuộc kiểu hữu hạn \sref{II.3.1.3}, nên $P'$ là một lược đồ quasi-compact (tương ứng, một tiền lược đồ Noether) nếu $Y$ là một lược đồ quasi-compact (tương ứng, một tiền lược đồ Noether).
Gọi $Z$ là bao đóng \sref[I]{I.9.5.10} của tiền lược đồ con $X'$ của $P'$ gắn với phép nhúng $j=r_{\sh{L},\vphi}$ từ $X$ vào $P'$;
hiển nhiên $j$ là hợp thành của một phép nhúng mở trội $i:X\to Z$ và phép nhúng chính tắc $Z\to P'$.
Nhưng $Z$ đồng nhất với một tiền lược đồ $\Proj(\sh{S})$, trong đó $\sh{S}$ là một $\sh{O}_Y$-đại số phân bậc, thương của $\bb{S}(\sh{E})$ theo một bó iđêan phân bậc quasi-coherent \sref{II.3.6.2}, và hiển nhiên $\sh{S}_1$ hữu hạn sinh và sinh $\sh{S}$;
hơn nữa, $\sh{O}_Z(1)$ là ảnh ngược của $\sh{O}_{P'}(1)$ qua phép nhúng chính tắc \sref{II.3.6.3}, do đó $\sh{L}=i^*(\sh{O}_Z(1))$.
\end{proof}

\begin{proposition}[4.4.8]
\label{II.4.4.8}
Cho $q:X\to Y$ là một cấu xạ, $\sh{L}$ là một $\sh{O}_X$-môđun rất ample đối với $q$, và $\sh{L}'$ là một $\sh{O}_X$-môđun khả nghịch, sao cho tồn tại một $\sh{O}_Y$-môđun quasi-coherent $\sh{E}'$ và một đồng cấu toàn ánh $q^*(\sh{E}')\to\sh{L}'$.
Khi đó $\sh{L}\otimes_{\sh{O}_X}\sh{L}'$ rất ample đối với $q$.
\end{proposition}

\begin{proof}
Giả thiết kéo theo sự tồn tại của một $Y$-cấu xạ $r':X\to P'=\bb{P}(\sh{E}')$ sao cho $\sh{L}'=r'^*(\sh{O}_{P'}(1))$ \sref{II.4.2.1}.
Theo giả thiết, tồn tại một $\sh{O}_Y$-môđun quasi-coherent $\sh{E}$ và một
\oldpage[II]{81}
$Y$-phép nhúng $r:X\to P=\bb{P}(\sh{E})$ sao cho $\sh{L}=r^*(\sh{O}_P(1))$.
Đặt $Q=\bb{P}(\sh{E}\otimes\sh{E}')$ và xét cấu xạ Segre $\varsigma:P\times_Y P'\to Q$ \sref{II.4.3.1}.
Vì $r$ là một phép nhúng, nên $(r,r')_Y:X\to P\times_Y P'$ cũng vậy \sref[I]{I.5.3.14};
nhưng vì $\varsigma$ là một phép nhúng \sref{II.4.3.3}, nên $r'':X\xrightarrow{(r,r')}P\times_Y P'\xrightarrow{\varsigma}Q$ cũng vậy.
Mặt khác, theo \sref{II.4.3.2.1}, $\varsigma^*(\sh{O}_Q(1))$ đẳng cấu với $\sh{O}_P(1)\otimes_Y\sh{O}_{P'}(1)$, do đó, theo \sref[I]{I.9.1.4}, $r''^*(\sh{O}_Q(1))$ đẳng cấu với $\sh{L}\otimes\sh{L}'$, điều này chứng minh mệnh đề.
\end{proof}

\begin{corollary}[4.4.9]
\label{II.4.4.9}
Cho $q:X\to Y$ là một cấu xạ.
\begin{enumerate}
  \item Cho $\sh{L}$ là một $\sh{O}_X$-môđun khả nghịch và $\sh{K}$ là một $\sh{O}_Y$-môđun khả nghịch.
  Điều kiện cần và đủ để $\sh{L}$ rất ample đối với $q$ là $\sh{L}\otimes q^*(\sh{K})$ cũng vậy.
  \item Nếu $\sh{L}$ và $\sh{L}'$ là hai $\sh{O}_X$-môđun rất ample đối với $q$, thì $\sh{L}\otimes\sh{L}'$ cũng vậy;
  đặc biệt, $\sh{L}^{\otimes n}$ rất ample đối với $q$ với mọi $n>0$.
\end{enumerate}
\end{corollary}

\begin{proof}
Khẳng định~(ii) là hệ quả trực tiếp của \sref{II.4.4.8}, cũng như tính cần thiết của điều kiện~(i);
ngược lại, nếu $\sh{L}\otimes q^*(\sh{K})$ rất ample, thì theo điều vừa chứng minh, $(\sh{L}\otimes q^*(\sh{K}))\otimes q^*(\sh{K}^{-1})$ cũng vậy, và $\sh{O}_X$-môđun sau đẳng cấu với $\sh{L}$ (\sref[0]{0.5.4.3} và \sref[0]{0.5.4.5}).
\end{proof}

\begin{proposition}[4.4.10]
\label{II.4.4.10}
\medskip\noindent
\begin{enumerate}
  \item[\rm{(i)}] Với mọi tiền lược đồ $Y$, mọi $\sh{O}_Y$-môđun khả nghịch $\sh{L}$ đều rất ample đối với cấu xạ đồng nhất $1_Y$.
  \item[\rm{(i \emph{bis})}] Cho $f:X\to Y$ là một cấu xạ và $j:X'\to X$ là một phép nhúng.
  Nếu $\sh{L}$ là một $\sh{O}_X$-môđun rất ample đối với $f$, thì $j^*(\sh{L})$ rất ample đối với $f\circ j$.
  \item[\rm{(ii)}] Cho $Z$ là một tiền lược đồ quasi-compact, $f:X\to Y$ là một cấu xạ thuộc kiểu hữu hạn, $g:Y\to Z$ là một cấu xạ quasi-compact, $\sh{L}$ là một $\sh{O}_X$-môđun rất ample đối với $f$, và $\sh{K}$ là một $\sh{O}_Y$-môđun rất ample đối với $g$.
  Khi đó tồn tại một số nguyên $n_0>0$ sao cho $\sh{L}\otimes f^*(\sh{K}^{\otimes n})$ rất ample đối với $g\circ f$ với mọi $n\geq n_0$.
  \item[\rm{(iii)}] Cho $f:X\to Y$ và $g:Y'\to Y$ là hai cấu xạ, và đặt $X'=X_{(Y')}$.
  Nếu $\sh{L}$ là một $\sh{O}_X$-môđun rất ample đối với $f$, thì $\sh{L}'=\sh{L}\otimes_Y\sh{O}_{Y'}$ là một $\sh{O}_{X'}$-môđun rất ample đối với $f_{(Y')}$.
  \item[\rm{(iv)}] Cho $f_i:X_i\to Y_i$ ($i=1,2$) là hai $S$-cấu xạ.
  Nếu $\sh{L}_i$ là một $\sh{O}_{X_i}$-môđun rất ample đối với $f_i$ ($i=1,2$), thì $\sh{L}_1\otimes_S\sh{L}_2$ rất ample đối với $f_1\times_S f_2$.
  \item[\rm{(v)}] Cho $f:X\to Y$ và $g:Y\to Z$ là hai cấu xạ.
  Nếu một $\sh{O}_X$-môđun $\sh{L}$ rất ample đối với $g\circ f$, thì $\sh{L}$ rất ample đối với $f$.
  \item[\rm{(vi)}] Cho $f:X\to Y$ là một cấu xạ và $j$ là phép nhúng chính tắc $X_\red\to X$.
  Nếu một $\sh{O}_X$-môđun $\sh{L}$ rất ample đối với $f$, thì $j^*(\sh{L})$ rất ample đối với $f_\red$.
\end{enumerate}
\end{proposition}

\begin{proof}
Tính chất~(i~\emph{bis}) suy ra ngay từ định nghĩa \sref{II.4.4.2}, và hiển nhiên (vi) suy ra một cách hình thức từ (i~\emph{bis}) và (v), bằng một lập luận theo sát lập luận trong chứng minh \sref[I]{I.5.5.12}, mà ta để lại cho bạn đọc.
Để chứng minh (v), như trong \sref[I]{I.5.5.12}, ta xét phân tích $X\xrightarrow{\Gamma_f}X\times_Z Y\xrightarrow{p_2}Y$, lưu ý rằng $p_2=(g\circ f)\times1_Y$.
Từ giả thiết cùng với (i) và (iv) suy ra $\sh{L}\otimes_{\sh{O}_Z}\sh{O}_Y$ rất ample đối với $p_2$;
mặt khác, $\sh{L}=\Gamma_f^*(\sh{L}\otimes_{\sh{O}_Z}\sh{O}_Y)$ \sref[I]{I.9.1.4}, và $\Gamma_f$ là một phép nhúng \sref[I]{I.5.3.11};
do đó ta có thể áp dụng (i~\emph{bis}).
\oldpage[II]{82}
Để chứng minh (i), ta áp dụng định nghĩa \sref{II.4.4.2} với $\sh{E}=\sh{L}$, lưu ý rằng khi đó $\bb{P}(\sh{E})$ đồng nhất với $Y$ \sref{II.4.1.4}.

Ta chứng minh (iii).
Tồn tại một $\sh{O}_Y$-môđun quasi-coherent $\sh{E}$ và một $Y$-phép nhúng $i:X\to\bb{P}(\sh{E})=P$ sao cho $\sh{L}=i^*(\sh{O}_P(1))$;
khi đó, nếu đặt $\sh{E}'=g^*(\sh{E})$, thì $\sh{E}'$ là một $\sh{O}_{Y'}$-môđun quasi-coherent, và ta có $P'=\bb{P}(\sh{E}')=P_{(Y')}$ \sref{II.4.1.3.1}; $i_{(Y')}$ là một phép nhúng từ $X_{(Y')}$ vào $P'$ \sref[I]{I.4.3.2}, và $\sh{L}'$ đẳng cấu với $(i_{(Y')})^*(\sh{O}_{P'}(1))$ \sref{II.4.2.10}.

Để chứng minh (iv), lưu ý rằng theo giả thiết, tồn tại một $Y_i$-phép nhúng $r_i:X_i\to P_i=\bb{P}(\sh{E}_i)$, trong đó $\sh{E}_i$ là một $\sh{O}_{Y_i}$-môđun quasi-coherent, và $\sh{L}_i=r_i^*(\sh{O}_{P_i}(1))$ ($i=1,2$);
$r_1\times_S r_2$ là một $S$-phép nhúng từ $X_1\times_S X_2$ vào $P_1\times_S P_2$ \sref[I]{I.4.3.1}, và ảnh ngược của $\sh{O}_{P_1}(1)\otimes_S\sh{O}_{P_2}(1)$ qua phép nhúng này là $\sh{L}_1\otimes_S\sh{L}_2$.
Mặt khác, đặt $T=Y_1\times_S Y_2$, và gọi $p_1,p_2$ lần lượt là các phép chiếu từ $T$ lên $Y_1,Y_2$.
Nếu đặt $P'_i=\bb{P}(p_i^*(\sh{E}_i))$ ($i=1,2$), thì theo \sref{II.4.1.3.1}, ta có $P'_i=P_i\times_{Y_i}T$, do đó
\[
  P'_1\times_T P'_2
  =(P_1\times_{Y_1}T)\times_T(P_2\times_{Y_2}T)
  =P_1\times_{Y_1}(T\times_{Y_2}P_2)
  =P_1\times_{Y_1}(Y_1\times_S P_2)
  =P_1\times_S P_2
\]
sai khác bởi một đẳng cấu chính tắc.
Tương tự, ta có $\sh{O}_{P'_i}(1)=\sh{O}_{P_i}(1)\otimes_{Y_i}\sh{O}_T$ \sref{II.4.1.3.2}, và một phép tính tương tự (đặc biệt dựa trên \sref[I]{I.9.1.9.1} và \sref[I]{I.9.1.2}) cho thấy, dưới phép đồng nhất trên, $\sh{O}_{P'_1}(1)\otimes_T\sh{O}_{P'_2}(1)$ đồng nhất với $\sh{O}_{P_1}(1)\otimes_S\sh{O}_{P_2}(1)$.
Do đó ta có thể xem $r_1\times_S r_2$ như một $T$-phép nhúng từ $X_1\times_S X_2$ vào $P'_1\times_T P'_2$, với ảnh ngược của $\sh{O}_{P'_1}(1)\otimes_T\sh{O}_{P'_2}(1)$ qua phép nhúng này là $\sh{L}_1\otimes_S\sh{L}_2$.
Sau đó ta kết thúc lập luận như trong \sref{II.4.4.8}, bằng cách dùng cấu xạ Segre.

Còn lại phải chứng minh (ii).
Trước hết, ta có thể thu hẹp về trường hợp $Z$ là một lược đồ affine, vì nói chung tồn tại một phủ hữu hạn $(U_i)$ của $Z$ bởi các tập mở affine;
nếu mệnh đề đã được chứng minh cho $\sh{K}|_{g^{-1}(U_i)}$, $\sh{L}|_{f^{-1}(g^{-1}(U_i))}$ và một số nguyên $n_i$, thì chỉ cần lấy $n_0$ là số lớn nhất trong các $n_i$ để chứng minh mệnh đề cho $\sh{K}$ và $\sh{L}$ \sref{II.4.4.5}.
Giả thiết kéo theo $f$ và $g$ là các cấu xạ tách, do đó $X$ và $Y$ là các \emph{lược đồ} quasi-compact.

Tồn tại một phép nhúng $r:X\to P=\bb{P}(\sh{E})$, trong đó $\sh{E}$ là một $\sh{O}_Y$-môđun quasi-coherent \emph{hữu hạn sinh}, và $\sh{L}=r^*(\sh{O}_P(1))$, theo \sref{II.4.4.6}.
Ta sẽ thấy rằng tồn tại một $\sh{O}_P$-môđun $\sh{M}$ rất ample đối với cấu xạ hợp thành $P\xrightarrow{h}Y\xrightarrow{g}Z$, sao cho $\sh{O}_P(1)$ đẳng cấu với $\sh{M}\otimes_Y\sh{K}^{\otimes(-m)}$ với một số nguyên $m$.
Với $n\geq m+1$, $\sh{O}_P(1)\otimes_Y\sh{K}^{\otimes n}$ do đó rất ample đối với $Z$, theo giả thiết và theo (iv) áp dụng cho các cấu xạ $h:P\to Y$ và $1_Y$;
vì $r$ là một phép nhúng và $\sh{L}\otimes f^*(\sh{K}^{\otimes n})=r^*(\sh{O}_P(1)\otimes_Y\sh{K}^{\otimes n})$, kết luận suy ra từ (i~\emph{bis}).
Để thiết lập khẳng định của ta về $\sh{O}_P(1)$, ta sẽ dùng bổ đề sau:

\begin{lemma}[4.4.10.1]
\label{II.4.4.10.1}
Cho $Z$ là một lược đồ quasi-compact, hoặc một tiền lược đồ có không gian nền Noether, $g:Y\to Z$ là một cấu xạ quasi-compact, $\sh{K}$ là một $\sh{O}_Y$-môđun khả nghịch rất ample đối với $g$, và $\sh{E}$ là một $\sh{O}_Y$-môđun quasi-coherent hữu hạn sinh.
Khi đó tồn tại một số nguyên $m_0$ sao cho, với mọi $m\geq m_0$, $\sh{E}$ đẳng cấu với một môđun thương của một $\sh{O}_Y$-môđun có dạng $g^*(\sh{F})\otimes\sh{K}^{\otimes(-m)}$, trong đó $\sh{F}$ là một $\sh{O}_Z$-môđun quasi-coherent hữu hạn sinh (phụ thuộc vào $m$).
\end{lemma}

Bổ đề này sẽ được chứng minh trong \sref{II.4.5.10.1};
bạn đọc có thể kiểm tra rằng \sref{II.4.4.10} không được dùng ở bất kỳ đâu trong \sref{subsection:II.4.5}.
\oldpage[II]{83}
Thừa nhận bổ đề này, tồn tại một phép nhúng đóng $j_1$ từ $P$ vào
\[
  P_1=\bb{P}(g^*(\sh{F})\otimes\sh{K}^{\otimes(-m)})
\]
sao cho $\sh{O}_P(1)$ đẳng cấu với $j_1^*(\sh{O}_{P_1}(1))$ \sref{II.4.1.2}.
Mặt khác, tồn tại một đẳng cấu từ $P_1$ lên $P_2=\bb{P}(g^*(\sh{F}))$, đồng nhất $\sh{O}_{P_1}(1)$ với $\sh{O}_{P_2}(1)\otimes_Y\sh{K}^{\otimes(-m)}$ \sref{II.4.1.4};
do đó ta có một phép nhúng đóng $j_2:P\to P_2$ sao cho $\sh{O}_P(1)$ đẳng cấu với $j_2^*(\sh{O}_{P_2}(1))\otimes_Y\sh{K}^{\otimes(-m)}$.
Cuối cùng, $P_2$ đồng nhất với $P_3\times_Z Y$, trong đó $P_3=\bb{P}(\sh{F})$, còn $\sh{O}_{P_2}(1)$ đồng nhất với $\sh{O}_{P_3}(1)\otimes_Z\sh{O}_Y$ \sref{II.4.1.3}.
Theo định nghĩa, $\sh{O}_{P_3}(1)$ rất ample đối với $Z$;
vì $\sh{K}$ cũng vậy, từ (iv) ta suy ra $\sh{O}_{P_2}(1)\otimes_Y\sh{K}$ rất ample đối với $Z$;
theo (i~\emph{bis}), $\sh{M}=j_2^*(\sh{O}_{P_2}(1)\otimes_Y\sh{K})$ cũng vậy, và $\sh{O}_P(1)$ đẳng cấu với $\sh{M}\otimes_Y\sh{K}^{\otimes(-m-1)}$, điều này hoàn tất chứng minh.
\end{proof}

\begin{proposition}[4.4.11]
\label{II.4.4.11}
Cho $f:X\to Y$ và $f':X'\to Y$ là hai cấu xạ, $X''$ là tiền lược đồ tổng $X\sqcup X'$, và $f''$ là cấu xạ $X''\to Y$ trùng với $f$ (tương ứng, $f'$) trên $X$ (tương ứng, $X'$).
Cho $\sh{L}$ (tương ứng, $\sh{L}'$) là một $\sh{O}_X$-môđun (tương ứng, một $\sh{O}_{X'}$-môđun) khả nghịch, và $\sh{L}''$ là $\sh{O}_{X''}$-môđun khả nghịch trùng với $\sh{L}$ trên $X$ và với $\sh{L}'$ trên $X'$.
Điều kiện cần và đủ để $\sh{L}''$ rất ample đối với $f''$ là $\sh{L}$ rất ample đối với $f$ và $\sh{L}'$ rất ample đối với $f'$.
\end{proposition}

\begin{proof}
Ta thu hẹp ngay về trường hợp $Y$ affine.
Nếu $\sh{L}''$ rất ample, thì $\sh{L}$ và $\sh{L}'$ cũng vậy theo \sref{II.4.4.10}[(i~\emph{bis})].
Ngược lại, nếu $\sh{L}$ và $\sh{L}'$ rất ample, thì từ định nghĩa \sref{II.4.4.2} và \sref{II.4.3.6} suy ra ngay $\sh{L}''$ rất ample.
\end{proof}

\subsection{Các bó ample}
\label{subsection:II.4.5}

\begin{env}[4.5.1]
\label{II.4.5.1}
Với một tiền lược đồ $X$ và một $\sh{O}_X$-môđun khả nghịch $\sh{L}$ đã cho, ta đặt, với mọi $\sh{O}_X$-môđun $\sh{F}$ (khi không thể có nhầm lẫn về $\sh{L}$), $\sh{F}(n)=\sh{F}\otimes_{\sh{O}_X}\sh{L}^{\otimes n}$ ($n\in\bb{Z}$);
mặt khác, ta đặt $S=\bigoplus_{n\geq0}\Gamma(X,\sh{L}^{\otimes n})$ (vành con phân bậc của vành $\Gamma_\bullet(\sh{L})$ được định nghĩa trong \sref[0]{0.5.4.6}).
Nếu xem $X$ như một $\bb{Z}$-tiền lược đồ và ký hiệu $p$ là cấu xạ cấu trúc $X\to\Spec(\bb{Z})$, thì có một tương ứng song ánh giữa các đồng cấu $p^*(\widetilde{S})\to\bigoplus_{n\geq0}\sh{L}^{\otimes n}$ của các $\sh{O}_X$-đại số phân bậc và các tự đồng cấu của vành phân bậc $S$ \sref[I]{I.2.2.5};
đồng cấu $\varepsilon:p^*(\widetilde{S})\to\bigoplus_{n\geq0}\sh{L}^{\otimes n}$ ứng như vậy với tự đẳng cấu \emph{đồng nhất} của $S$ được gọi là \emph{chính tắc}.
Ứng với nó là một cấu xạ $G(\varepsilon)\to\Proj(S)$ \sref{II.3.7.1}, cấu xạ này cũng được gọi là \emph{chính tắc}.
\end{env}

\begin{theorem}[4.5.2]
\label{II.4.5.2}
Cho $X$ là một lược đồ quasi-compact hoặc một tiền lược đồ có không gian nền Noether, $\sh{L}$ là một $\sh{O}_X$-môđun khả nghịch, và $S$ là vành phân bậc $\bigoplus_{n\geq0}\Gamma(X,\sh{L}^{\otimes n})$.
Các điều kiện sau tương đương:
\begin{enumerate}
  \item[\rm{(a)}] Khi $f$ chạy qua tập hợp các phần tử thuần nhất của $S_+$, các $X_f$ tạo thành một cơ sở của tôpô trên $X$.
  \item[\rm{(a')}] Khi $f$ chạy qua tập hợp các phần tử thuần nhất của $S_+$, những $X_f$ affine tạo thành một phủ của $X$.
  \item[\rm{(b)}] Cấu xạ chính tắc $G(\varepsilon)\to\Proj(S)$ \sref{II.4.5.1} được xác định khắp nơi và là một phép nhúng mở trội.
\oldpage[II]{84}
  \item[\rm{(b')}] Cấu xạ chính tắc $G(\varepsilon)\to\Proj(S)$ được xác định khắp nơi và là một phép đồng phôi từ không gian nền của $X$ lên một không gian con của $\Proj(S)$.
  \item[\rm{(c)}] Với mọi $\sh{O}_X$-môđun quasi-coherent $\sh{F}$, nếu ký hiệu $\sh{F}_n$ là $\sh{O}_X$-môđun con của $\sh{F}(n)$ được sinh bởi các tiết diện của $\sh{F}(n)$ trên $X$, thì $\sh{F}$ là tổng của các $\sh{O}_X$-môđun con $\sh{F}_n(-n)$, với các số nguyên $n>0$.
  \item[\rm{(c')}] Tính chất~\rm{(c)} đúng với mọi bó iđêan quasi-coherent của $\sh{O}_X$.
\end{enumerate}

Hơn nữa, nếu $(f_\alpha)$ là một họ các phần tử thuần nhất của $S_+$ sao cho các $X_{f_\alpha}$ affine, thì hạn chế lên $\bigcup_\alpha X_{f_\alpha}$ của cấu xạ chính tắc $X\to\Proj(S)$ là một đẳng cấu từ $\bigcup_\alpha X_{f_\alpha}$ lên $\bigcup_\alpha(\Proj(S))_{f_\alpha}$.
\end{theorem}

\begin{proof}
Hiển nhiên (b) kéo theo (b'), và (b') kéo theo (a) theo công thức \sref{II.3.7.3.1} (lưu ý rằng $\varepsilon^\flat$ là đồng nhất).
Điều kiện~(a) kéo theo (a'), vì mọi $x\in X$ đều có một lân cận affine $U$ sao cho $\sh{L}|U$ đẳng cấu với $\sh{O}_X|U$;
nếu $f\in\Gamma(X,\sh{L}^{\otimes n})$ thỏa mãn $x\in X_f\subset U$, thì $X_f$ cũng là tập hợp các $x'\in U$ sao cho $(f|U)(x')\neq0$, và do đó là một tập mở affine \sref[I]{I.1.3.6}.
Để chứng minh (a') kéo theo (b), chỉ cần chứng minh khẳng định cuối của định lý và kiểm tra thêm rằng, nếu $X=\bigcup_\alpha X_{f_\alpha}$, thì điều kiện~(iv) của \sref{II.3.8.2} được thỏa mãn.
Điểm sau suy ra ngay từ \sref[I]{I.9.3.1}[(i)].
Đối với khẳng định cuối của \sref{II.4.5.2}, vì $X_{f_\alpha}$ là ảnh ngược của $(\Proj(S))_{f_\alpha}$ qua $G(\varepsilon)\to\Proj(S)$, chỉ cần áp dụng \sref[I]{I.9.3.2}.
Vậy (a), (a'), (b) và (b') tương đương.

Để chứng minh (a') kéo theo (c), lưu ý rằng nếu $X_f$ affine (với $f\in S_k$), thì $\sh{F}|X_f$ được sinh bởi các tiết diện của nó trên $X_f$ \sref[I]{I.1.3.9};
mặt khác, theo \sref[I]{I.9.3.1}[(ii)], một tiết diện $s$ như vậy có dạng $(t|X_f)\otimes(f|X_f)^{-m}$, trong đó $t\in\Gamma(X,\sh{F}(km))$;
theo định nghĩa, $t$ cũng là một tiết diện của $\sh{F}_{km}$, nên $s$ đúng là một tiết diện của $\sh{F}_{km}(-km)$ trên $X_f$, điều này chứng minh (c).
Hiển nhiên (c) kéo theo (c'), và còn phải chứng minh (c') kéo theo (a).
Thật vậy, cho $U$ là một lân cận mở của $x\in X$, và cho $\sh{J}$ là một bó iđêan quasi-coherent của $\sh{O}_X$ xác định một tiền lược đồ con đóng của $X$ có không gian nền là $X\setmin U$ \sref[I]{I.5.2.1}.
Giả thiết~(c') kéo theo tồn tại một số nguyên $n>0$ và một tiết diện $f$ của $\sh{J}(n)$ trên $X$ sao cho $f(x)\neq0$.
Nhưng hiển nhiên $f\in S_n$ và $x\in X_f\subset U$, điều này chứng minh (a).
\end{proof}

Khi $X$ là một tiền lược đồ có không gian nền Noether, các điều kiện tương đương của \sref{II.4.5.2} kéo theo $X$ là một \emph{lược đồ}, vì theo \sref{II.4.5.2}[(b)], nó đẳng cấu với một tiền lược đồ con của lược đồ $\Proj(S)$.

\begin{definition}[4.5.3]
\label{II.4.5.3}
Ta nói một $\sh{O}_X$-môđun khả nghịch $\sh{L}$ là \emph{ample} nếu $X$ là một lược đồ quasi-compact và các điều kiện tương đương của \sref{II.4.5.2} được thỏa mãn.
\end{definition}

Từ tiêu chuẩn~(a) của \sref{II.4.5.2} suy ra ngay rằng, nếu $\sh{L}$ là một $\sh{O}_X$-môđun ample, thì với mọi tập mở quasi-compact $U$ của $X$, $\sh{L}|U$ là một $(\sh{O}_X|U)$-môđun ample.

Từ chứng minh \sref{II.4.5.2} suy ra các $X_f$ \emph{affine} đã tạo thành một cơ sở của tôpô trên $X$.
Hơn nữa:

\begin{corollary}[4.5.4]
\label{II.4.5.4}
Cho $\sh{L}$ là một $\sh{O}_X$-môđun ample.
Với mọi không gian con hữu hạn $Z$ của $X$ và mọi lân cận $U$ của $Z$, tồn tại một số nguyên $n$ và một $f\in\Gamma(X,\sh{L}^{\otimes n})$ sao cho $X_f$ là một lân cận affine của $Z$ nằm trong $U$.
\end{corollary}

\begin{proof}
\oldpage[II]{85}
Theo \sref{II.4.5.2}[(b)], ta chỉ cần chứng minh rằng, với mọi tập con hữu hạn $Z'$ của $\Proj(S)$ và mọi lân cận mở $U$ của $Z'$, tồn tại một phần tử thuần nhất $f\in S_+$ sao cho $Z'\subset(\Proj(S))_f\subset U$ \sref{II.2.4.1}.
Thật vậy, theo định nghĩa, tập đóng $Y$, là phần bù của $U$ trong $\Proj(S)$, có dạng $V_+(\mathfrak{I})$, trong đó $\mathfrak{I}$ là một iđêan phân bậc của $S$ không chứa $S_+$ \sref{II.2.3.2};
mặt khác, theo định nghĩa, các điểm của $Z'$ là những iđêan nguyên tố phân bậc $\mathfrak{p}_i$ của $S$ không chứa $\mathfrak{I}$ \sref{II.2.3.1}.
Do đó tồn tại một phần tử $f\in\mathfrak{I}$ không thuộc bất kỳ $\mathfrak{p}_i$ nào (Bourbaki, \emph{Alg. comm.}, Chương~II, \S1, số~1, Mệnh đề~2), và vì các $\mathfrak{p}_i$ là phân bậc, lập luận trong \emph{loc. cit.} cho thấy thậm chí có thể chọn $f$ thuần nhất;
phần tử này khi đó thỏa mãn điều cần chứng minh.
\end{proof}

\begin{proposition}[4.5.5]
\label{II.4.5.5}
Giả sử $X$ là một lược đồ quasi-compact, hoặc là một tiền lược đồ có không gian nền Noether.
Các điều kiện~(a) đến (c') của \sref{II.4.5.2} cũng tương đương với các điều kiện sau:
\begin{enumerate}
  \item[\rm{(d)}] Với mọi $\sh{O}_X$-môđun quasi-coherent hữu hạn sinh $\sh{F}$, tồn tại một số nguyên $n_0$ sao cho, với mọi $n\geq n_0$, $\sh{F}(n)$ được sinh bởi các tiết diện của nó trên $X$.
  \item[\rm{(d')}] Với mọi $\sh{O}_X$-môđun quasi-coherent hữu hạn sinh $\sh{F}$, tồn tại hai số nguyên $n>0$, $k>0$ sao cho $\sh{F}$ đẳng cấu với một môđun thương của $\sh{O}_X$-môđun $\sh{L}^{\otimes(-n)}\otimes\sh{O}_X^k$.
  \item[\rm{(d'')}] Tính chất~(d') được thỏa mãn đối với mọi bó iđêan quasi-coherent hữu hạn sinh của $\sh{O}_X$.
\end{enumerate}
\end{proposition}

\begin{proof}
Vì $X$ quasi-compact, nếu một $\sh{O}_X$-môđun quasi-coherent hữu hạn sinh $\sh{F}$ có tính chất $\sh{F}(n)$ (vốn hữu hạn sinh) được sinh bởi các tiết diện của nó trên $X$, thì $\sh{F}(n)$ được sinh bởi một số \emph{hữu hạn} các tiết diện ấy \sref[0]{0.5.2.3};
do đó (d) kéo theo (d'), và hiển nhiên (d') kéo theo (d'').
Vì mọi $\sh{O}_X$-môđun quasi-coherent $\sh{G}$ là giới hạn quy nạp của các $\sh{O}_X$-môđun con hữu hạn sinh của nó \sref[I]{I.9.4.9}, để kiểm tra điều kiện~(c') của \sref{II.4.5.2}, chỉ cần kiểm tra điều kiện ấy cho một bó iđêan quasi-coherent hữu hạn sinh của $\sh{O}_X$; do đó (d'') kéo theo (c').
Còn phải chứng minh rằng nếu $\sh{L}$ ample thì tính chất~(d) được thỏa mãn.
Xét một phủ hữu hạn của $X$ bởi các $X_{f_i}$ ($f_i\in S_{n_i}$), mà ta có thể giả sử là affine;
thay các $f_i$ bằng những lũy thừa thích hợp (điều này không làm thay đổi các $X_{f_i}$), ta có thể giả sử mọi $n_i$ đều bằng cùng một số nguyên $m$.
Theo giả thiết, bó $\sh{F}|X_{f_i}$ là hữu hạn sinh, nên nó được sinh bởi một số hữu hạn các tiết diện $h_{ij}$ trên $X_{f_i}$ \sref[I]{I.1.3.13};
do đó tồn tại một số nguyên $k_0$ sao cho tiết diện $h_{ij}\otimes f_i^{\otimes k_0}$ mở rộng được thành một tiết diện của $\sh{F}(k_0m)$ trên $X$ đối với mọi cặp $(i,j)$ \sref[I]{I.9.3.1}.
\emph{A fortiori}, các $h_{ij}\otimes f_i^{\otimes k}$ mở rộng được thành những tiết diện của $\sh{F}(km)$ trên $X$ với mọi $k\geq k_0$, và với các giá trị ấy của $k$, $\sh{F}(km)$ do đó được sinh bởi các tiết diện của nó trên $X$.
Với mọi $p$ sao cho $0<p<m$, $\sh{F}(p)$ cũng hữu hạn sinh; do đó tồn tại một số nguyên $k_p$ sao cho $\sh{F}(p)(km)=\sh{F}(p+km)$ được sinh bởi các tiết diện của nó trên $X$ với mọi $k\geq k_p$.
Chọn $n_0$ lớn hơn mọi $k_pm$, ta kết luận rằng $\sh{F}(n)$ được sinh bởi các tiết diện của nó trên $X$ với mọi $n\geq n_0$, vì mỗi $n$ như vậy viết được dưới dạng $n=km+p$, với $k\geq k_p$ và $0\leq p<m$.
\end{proof}

\begin{proposition}[4.5.6]
\label{II.4.5.6}
Cho $X$ là một lược đồ quasi-compact và $\sh{L}$ là một $\sh{O}_X$-môđun khả nghịch.
\begin{enumerate}
  \item[\rm{(i)}] Cho $n>0$ là một số nguyên. Điều kiện cần và đủ để $\sh{L}$ ample là $\sh{L}^{\otimes n}$ ample.
  \item[\rm{(ii)}] Cho $\sh{L}'$ là một $\sh{O}_X$-môđun khả nghịch sao cho, với mọi $x\in X$, tồn tại một số nguyên $n>0$
\oldpage[II]{86}
    và một tiết diện $s'$ của $\sh{L}'^{\otimes n}$ trên $X$ sao cho $s'(x)\neq0$.
    Khi đó, nếu $\sh{L}$ ample thì $\sh{L}\otimes\sh{L}'$ cũng ample.
\end{enumerate}
\end{proposition}

\begin{proof}
Tính chất~(i) là hệ quả hiển nhiên của tiêu chuẩn~(a) trong \sref{II.4.5.2}, vì $X_{f^{\otimes n}}=X_f$.
Mặt khác, nếu $\sh{L}$ ample thì, với mọi $x\in X$ và mọi lân cận $U$ của $x$, tồn tại $m>0$ và $f\in\Gamma(X,\sh{L}^{\otimes m})$ sao cho $x\in X_f\subset U$ \sref{II.4.5.2}[(a)];
nếu $f'\in\Gamma(X,\sh{L}'^{\otimes n})$ thỏa $f'(x)\neq0$, thì $s(x)\neq0$ đối với $s=f^{\otimes n}\otimes f'^{\otimes m}\in\Gamma(X,(\sh{L}\otimes\sh{L}')^{\otimes mn})$; do đó $x\in X_s\subset X_f\subset U$, điều này chứng minh rằng $\sh{L}\otimes\sh{L}'$ ample \sref{II.4.5.2}[(a)].
\end{proof}

\begin{corollary}[4.5.7]
\label{II.4.5.7}
Tích tenxơ của hai $\sh{O}_X$-môđun ample là ample.
\end{corollary}

\begin{corollary}[4.5.8]
\label{II.4.5.8}
Cho $\sh{L}$ là một $\sh{O}_X$-môđun ample và $\sh{L}'$ là một $\sh{O}_X$-môđun khả nghịch;
khi đó tồn tại một số nguyên $n_0>0$ sao cho $\sh{L}^{\otimes n}\otimes\sh{L}'$ ample và được sinh bởi các tiết diện của nó trên $X$ với mọi $n\geq n_0$.
\end{corollary}

\begin{proof}
Thật vậy, từ \sref{II.4.5.5} suy ra rằng tồn tại một số nguyên $m_0$ sao cho $\sh{L}^{\otimes m}\otimes\sh{L}'$ được sinh bởi các tiết diện của nó trên $X$ với mọi $m\geq m_0$;
do \sref{II.4.5.6}, ta có thể lấy $n_0=m_0+1$.
\end{proof}

\begin{remark}[4.5.9]
\label{II.4.5.9}
Cho $P=\HH^1(X,\sh{O}_X^\times)$ là nhóm các lớp $\sh{O}_X$-môđun khả nghịch \sref[0]{0.5.4.7}, và cho $P^+$ là tập con của $P$ gồm các lớp của những bó ample.
Giả sử $P^+$ \emph{khác rỗng}.
Khi đó, từ \sref{II.4.5.7} và \sref{II.4.5.8} suy ra
\[
  P^+ + P^+ \subset P^+
  \quad\text{và}\quad
  P^+ - P^+ = P.
\]
Nói cách khác, $P^+\cup\{0\}$ là tập các phần tử \emph{dương} trong $P$ đối với một cấu trúc \emph{tiền thứ tự} trên $P$, tương thích với cấu trúc nhóm của nó, và thậm chí mang tính \emph{Archimedes} do \sref{II.4.5.8}.
Vì thế đôi khi người ta nói ``bó dương'' thay cho bó ample và ``bó âm'' cho nghịch đảo của một bó như vậy (nhưng ta sẽ không dùng thuật ngữ này).
\end{remark}

\begin{proposition}[4.5.10]
\label{II.4.5.10}
Cho $Y$ là một lược đồ affine, $q:X\to Y$ là một cấu xạ tách quasi-compact và $\sh{L}$ là một $\sh{O}_X$-môđun khả nghịch.
\begin{enumerate}
  \item[\rm{(i)}] Nếu $\sh{L}$ rất ample đối với $q$, thì $\sh{L}$ ample.
  \item[\rm{(ii)}] Giả sử thêm rằng cấu xạ $q$ thuộc \emph{kiểu hữu hạn}.
    Điều kiện cần và đủ để $\sh{L}$ ample là nó có một trong các tính chất tương đương sau:
    \begin{enumerate}
      \item[\rm{(e)}] Tồn tại $n_0>0$ sao cho, với mọi số nguyên $n\geq n_0$, $\sh{L}^{\otimes n}$ rất ample đối với $q$.
      \item[\rm{(e')}] Tồn tại $n>0$ sao cho $\sh{L}^{\otimes n}$ rất ample đối với $q$.
    \end{enumerate}
\end{enumerate}
\end{proposition}

\begin{proof}
Khẳng định thứ nhất suy ra từ định nghĩa \sref{II.4.4.2} của một $\sh{O}_X$-môđun rất ample:
nếu $A$ là vành của $Y$, thì tồn tại một $A$-môđun $E$ và một đồng cấu toàn ánh
\[
  \psi:q^*((\bb{S}(E))^{\supertilde})\to\bigoplus_{n\geq0}\sh{L}^{\otimes n}
\]
sao cho $i=r_{\sh{L},\psi}$ là một phép nhúng được xác định khắp nơi $X\to P=\bb{P}(\widetilde{E})$ và $\sh{L}=i^*(\sh{O}_P(1))$;
vì các $D_+(f)$, với $f$ thuần nhất trong $(\bb{S}(E))_+$, lập thành một cơ sở của tôpô trên $P$, và vì $i^{-1}(D_+(f))=X_{\psi^\flat(f)}$ theo \sref{II.3.7.3.1}, ta thấy điều kiện~(a) của \sref{II.4.5.2} được thỏa mãn, do đó $\sh{L}$ ample.

Bây giờ ta chứng minh rằng nếu $q$ thuộc kiểu hữu hạn và $\sh{L}$ ample thì tính chất~(e) được thỏa mãn.
Trước hết, từ tiêu chuẩn~(b) của \sref{II.4.5.2} và \sref{II.4.4.1}[(i)] suy ra rằng tồn tại
\oldpage[II]{87}
một số nguyên $k_0>0$ sao cho $\sh{L}^{\otimes k_0}$ rất ample đối với $q$.
Mặt khác, do \sref{II.4.5.5}, tồn tại một số nguyên $m_0$ sao cho, với mọi $m\geq m_0$, $\sh{L}^{\otimes m}$ được sinh bởi các tiết diện của nó trên $X$.
Đặt $n_0=k_0+m_0$;
nếu $n\geq n_0$, ta có thể viết $n=k_0+m$ với $m\geq m_0$, do đó $\sh{L}^{\otimes n}=\sh{L}^{\otimes k_0}\otimes\sh{L}^{\otimes m}$.
Vì $\sh{L}^{\otimes m}$ được sinh bởi các tiết diện của nó trên $X$, từ \sref{II.4.4.8} và \sref{II.3.4.7} suy ra rằng $\sh{L}^{\otimes n}$ rất ample đối với $q$.
Cuối cùng, hiển nhiên (e) kéo theo (e'), và (e') kéo theo $\sh{L}$ ample do (i) và \sref{II.4.5.6}[(i)].

  \begin{env}[4.5.10.1]
  \label{II.4.5.10.1}
  \emph{[Chứng minh Bổ đề~\sref{II.4.4.10.1}].}
  Đặt $\sh{E}(n)=\sh{E}\otimes\sh{K}^{\otimes n}$;
  vì $g$ tách \sref{II.4.4.2}, lập luận của \sref{II.3.4.8} áp dụng được và quy việc chứng minh về chỗ chỉ ra rằng đồng cấu chính tắc $g^*(g_*(\sh{E}(n)))\to\sh{E}(n)$ là toàn ánh với mọi $n$ đủ lớn.
  Hơn nữa, vì $Z$ quasi-compact, lập luận của \sref{II.3.4.6} quy việc chứng minh về trường hợp $Z$ affine.
  Nhưng khi đó $\sh{K}$ ample do \sref{II.4.5.10}[(i)], và kết luận suy ra từ \sref{II.4.5.5}[(d')].
  \end{env}
\end{proof}

\begin{corollary}[4.5.11]
\label{II.4.5.11}
Cho $Y$ là một lược đồ affine, $q:X\to Y$ là một cấu xạ tách thuộc kiểu hữu hạn, $\sh{L}$ là một $\sh{O}_X$-môđun ample và $\sh{L}'$ là một $\sh{O}_X$-môđun khả nghịch.
Tồn tại một số nguyên $n_0$ sao cho, với mọi $n\geq n_0$, $\sh{L}^{\otimes n}\otimes\sh{L}'$ rất ample đối với $q$.
\end{corollary}

\begin{proof}
Tồn tại $m_0$ sao cho, với mọi $m\geq m_0$, $\sh{L}^{\otimes m}\otimes\sh{L}'$ được sinh bởi các tiết diện của nó trên $X$ \sref{II.4.5.8};
mặt khác, tồn tại $k_0$ sao cho $\sh{L}^{\otimes k}$ rất ample đối với $q$ với mọi $k\geq k_0$.
Do đó $\sh{L}^{\otimes(k+m_0)}\otimes\sh{L}'$ rất ample với mọi $k\geq k_0$ (\sref{II.4.4.8} và \sref{II.3.4.7}).
\end{proof}

\begin{remark}[4.5.12]
\label{II.4.5.12}
Ta không biết giả thiết rằng một $\sh{O}_X$-môđun $\sh{L}$ có $\sh{L}^{\otimes n}$ rất ample (đối với $q$) có kéo theo cùng kết luận ấy cho $\sh{L}^{\otimes(n+1)}$ hay không.
\end{remark}

\begin{proposition}[4.5.13]
\label{II.4.5.13}
Cho $X$ là một tiền lược đồ quasi-compact, $Z$ là một tiền lược đồ con đóng của $X$ được xác định bởi một bó iđêan quasi-coherent lũy linh $\sh{J}$ của $\sh{O}_X$, và $j$ là phép nhúng chính tắc $Z\to X$.
Điều kiện cần và đủ để một $\sh{O}_X$-môđun khả nghịch $\sh{L}$ ample là $\sh{L}'=j^*(\sh{L})$ là một $\sh{O}_Z$-môđun ample.
\end{proposition}

\begin{proof}
Điều kiện là \emph{cần}.
Thật vậy, với mọi tiết diện $f$ của $\sh{L}^{\otimes n}$ trên $X$, cho $f'$ là ảnh chính tắc $f\otimes1$ của nó, một tiết diện của $\sh{L}'^{\otimes n}=\sh{L}^{\otimes n}\otimes_{\sh{O}_X}(\sh{O}_X/\sh{J})$ trên không gian $Z$ (đồng nhất với $X$);
hiển nhiên $X_f=Z_{f'}$, nên tiêu chuẩn~(a) của \sref{II.4.5.2} cho thấy $\sh{L}'$ ample.

Để thấy điều kiện là \emph{đủ}, trước hết lưu ý rằng ta có thể quy về trường hợp $\sh{J}^2=0$ bằng cách xét dãy (hữu hạn) các tiền lược đồ $X_k=(X,\sh{O}_X/\sh{J}^{k+1})$, trong đó mỗi tiền lược đồ là một tiền lược đồ con đóng của tiền lược đồ kế tiếp, được xác định bởi một bó iđêan có bình phương bằng không.
Hơn nữa, $X$ là một lược đồ vì theo giả thiết $X_\red$ là một lược đồ (\sref{II.4.5.3} và \sref[I]{I.5.5.1}).
Tiêu chuẩn~(a) của \sref{II.4.5.2} cho thấy chỉ cần chứng minh
  \begin{lemma}[4.5.13.1]
  \label{II.4.5.13.1}
  Dưới các giả thiết của \sref{II.4.5.13}, giả sử thêm rằng $\sh{J}$ có bình phương bằng không;
  với $\sh{L}$ là một $\sh{O}_X$-môđun khả nghịch, cho $g$ là một tiết diện của $\sh{L}'^{\otimes n}$ trên $Z$ sao cho $Z_g$ affine.
  Khi đó tồn tại một số nguyên $m>0$ sao cho $g^{\otimes m}$ là ảnh chính tắc của một tiết diện $f$ của $\sh{L}^{\otimes nm}$ trên $X$.
  \end{lemma}

  Ta có dãy khớp các $\sh{O}_X$-môđun
  \[
    0\to\sh{J}(n)\to\sh{O}_X(n)=\sh{L}^{\otimes n}\to\sh{O}_Z(n)=\sh{L}'^{\otimes n}\to0
  \]
\oldpage[II]{88}
  vì $\sh{F}(n)$ là một hàm tử khớp theo $\sh{F}$;
  từ đó có dãy khớp đối đồng điều
  \[
    0\to\Gamma(X,\sh{J}(n))\to\Gamma(X,\sh{L}^{\otimes n})\to\Gamma(X,\sh{L}'^{\otimes n})\xrightarrow{\partial}\HH^1(X,\sh{J}(n))
  \]
  đặc biệt gửi $g$ tới một phần tử $\partial g\in\HH^1(X,\sh{J}(n))$.
  Lưu ý rằng, vì $\sh{J}^2=0$, $\sh{J}$ có thể được xem như một $\sh{O}_Z$-môđun quasi-coherent và, với mọi $k$, $\sh{L}'^{\otimes k}\otimes_{\sh{O}_Z}\sh{J}(n)=\sh{J}(n+k)$;
  do đó, với mọi tiết diện $s\in\Gamma(X,\sh{L}'^{\otimes k})$, phép nhân tenxơ với $s$ là một đồng cấu $\sh{J}(n)\to\sh{J}(n+k)$ của các $\sh{O}_Z$-môđun, và vì thế cảm sinh một đồng cấu $\HH^i(X,\sh{J}(n))\xrightarrow{s}\HH^i(X,\sh{J}(n+k))$ của các nhóm đối đồng điều.

  Với những điều ấy, ta sẽ chứng minh rằng
  \[
  \label{II.4.5.13.2}
    g^{\otimes m}\otimes\partial g=0
  \tag{4.5.13.2}
  \]
  với $m>0$ đủ lớn.
  Thật vậy, $Z_g$ là một \emph{tập mở affine} của $Z$, nên $\HH^1(Z_g,\sh{J}(n))=0$ khi $\sh{J}(n)$ được xem như một \emph{$\sh{O}_Z$-môđun} \sref[I]{I.5.1.9.2}.
  Đặc biệt, nếu đặt $g'=g|Z_g$ và xét ảnh của nó qua ánh xạ $\partial:\Gamma(Z_g,\sh{L}'^{\otimes n})\to\HH^1(Z_g,\sh{J}(n))$, thì $\partial g'=0$.
  Để viết tường minh hệ thức này, lưu ý rằng ở bậc~$1$, đối đồng điều của một bó nhóm abel trùng với đối đồng điều \v{C}ech của nó (G, II, 5.9);
  do đó, để tính $\partial g$, phải xét một phủ mở $(U_\alpha)$ đủ mịn của $X$, mà ta có thể giả sử là \emph{hữu hạn} và gồm các tập mở affine, chọn với mỗi $\alpha$ một tiết diện $g_\alpha\in\Gamma(U_\alpha,\sh{L}^{\otimes n})$ có ảnh chính tắc trong $\Gamma(U_\alpha,\sh{L}'^{\otimes n})$ bằng $g|U_\alpha$, rồi xét lớp của đối chu trình $(g_{\alpha|\beta}-g_{\beta|\alpha})$, trong đó $g_{\alpha|\beta}$ ký hiệu hạn chế của $g_\alpha$ lên $U_\alpha\cap U_\beta$ (đối chu trình này nhận giá trị trong $\sh{J}(n)$).
  Ta còn có thể giả sử $\partial g'$ được tính theo cùng cách, bằng phủ gồm các $U_\alpha\cap Z_g$ và các hạn chế $g_\alpha|(U_\alpha\cap Z_g)$ (nếu cần, thay $(U_\alpha)$ bằng một phủ mịn hơn);
  hệ thức $\partial g'=0$ khi đó có nghĩa là, với mỗi $\alpha$, tồn tại một tiết diện $h_\alpha\in\Gamma(U_\alpha\cap Z_g,\sh{J}(n))$ sao cho $(g_{\alpha|\beta}-g_{\beta|\alpha})|(U_\alpha\cap U_\beta\cap Z_g)=h_{\alpha|\beta}-h_{\beta|\alpha}$, trong đó $h_{\alpha|\beta}$ ký hiệu hạn chế của $h_\alpha$ lên $U_\alpha\cap U_\beta\cap Z_g$ (G, II, 5.11).
  Khi đó tồn tại một số nguyên $m>0$ sao cho $g^{\otimes m}\otimes h_\alpha$ là hạn chế lên $U_\alpha\cap Z_g$ của một tiết diện $t_\alpha\in\Gamma(U_\alpha,\sh{J}(n+nm))$ với mọi $\alpha$ \sref[I]{I.9.3.1};
  do đó $g^{\otimes m}\otimes(g_{\alpha|\beta}-g_{\beta|\alpha})=t_{\alpha|\beta}-t_{\beta|\alpha}$ với mọi cặp chỉ số, điều này chứng minh \sref{II.4.5.13.2}.

  Mặt khác, lưu ý rằng nếu $s\in\Gamma(X,\sh{O}_Z(p))$ và $t\in\Gamma(X,\sh{O}_Z(q))$, thì trong nhóm $\HH^1(X,\sh{J}(p+q))$ ta có
  \[
  \label{II.4.5.13.3}
    \partial(s\otimes t)=(\partial s)\otimes t+s\otimes(\partial t).
  \tag{4.5.13.3}
  \]

  Thật vậy, ta lại có thể tính hai vế bằng cách xét một phủ mở $(U_\alpha)$ của $X$ và, với mỗi $\alpha$, một tiết diện $s_\alpha\in\Gamma(U_\alpha,\sh{O}_X(p))$ (tương ứng, $t_\alpha\in\Gamma(U_\alpha,\sh{O}_X(q))$) có ảnh chính tắc trong $\Gamma(U_\alpha,\sh{O}_Z(p))$ (tương ứng, $\Gamma(U_\alpha,\sh{O}_Z(q))$) bằng $s|U_\alpha$ (tương ứng, $t|U_\alpha$);
  hệ thức~\sref{II.4.5.13.3} khi đó suy ra từ các hệ thức
  \[
    (s_{\alpha|\beta}\otimes t_{\alpha|\beta})-(s_{\beta|\alpha}\otimes t_{\beta|\alpha})
    =(s_{\alpha|\beta}-s_{\beta|\alpha})\otimes t_{\alpha|\beta}+s_{\beta|\alpha}\otimes(t_{\alpha|\beta}-t_{\beta|\alpha})
  \]
  với cùng các ký hiệu như trên.
  Do đó, bằng quy nạp theo $k$, ta có
  \[
  \label{II.4.5.13.4}
    \partial(g^{\otimes k})=(kg^{\otimes(k-1)})\otimes(\partial g)
  \tag{4.5.13.4}
  \]

\oldpage[II]{89}
  và từ \sref{II.4.5.13.2} cùng \sref{II.4.5.13.4}, ta kết luận rằng $\partial(g^{\otimes(m+1)})=0$;
  do đó $g^{\otimes(m+1)}$ là ảnh chính tắc của một tiết diện $f$ của $\sh{L}^{\otimes n(m+1)}$ trên $X$, điều này hoàn tất chứng minh \sref{II.4.5.13}.
\end{proof}

\begin{corollary}[4.5.14]
\label{II.4.5.14}
Cho $X$ là một lược đồ Noether và $j$ là phép nhúng chính tắc $X_\red\to X$.
Điều kiện cần và đủ để một $\sh{O}_X$-môđun khả nghịch $\sh{L}$ ample là $j^*(\sh{L})$ là một $\sh{O}_{X_\red}$-môđun ample.
\end{corollary}

\begin{proof}
Điều này suy ra từ \sref[I]{I.6.1.6}.
\end{proof}


\subsection{Các bó ample tương đối}
\label{subsection:II.4.6}

\begin{definition}[4.6.1]
\label{II.4.6.1}
Cho $f:X\to Y$ là một cấu xạ quasi-compact và $\sh{L}$ là một $\sh{O}_X$-môđun khả nghịch.
Ta nói $\sh{L}$ \emph{ample tương đối với $f$}, hoặc \emph{tương đối với $Y$}, hoặc \emph{$f$-ample}, hoặc \emph{$Y$-ample} (hay thậm chí đơn giản là \emph{ample}, nếu không thể gây nhầm lẫn với khái niệm được định nghĩa trong \sref{II.4.5.3}) nếu tồn tại một phủ $(U_\alpha)$ của $Y$ gồm các tập mở affine sao cho, đặt $X_\alpha=f^{-1}(U_\alpha)$, $\sh{L}|X_\alpha$ là một $\sh{O}_{X_\alpha}$-môđun ample với mọi $\alpha$.
\end{definition}

Sự tồn tại của một $\sh{O}_X$-môđun $f$-ample tất yếu kéo theo rằng $f$ \emph{tách} (\sref{II.4.5.3} và \sref[I]{I.5.5.5}).

\begin{proposition}[4.6.2]
\label{II.4.6.2}
Cho $f:X\to Y$ là một cấu xạ quasi-compact và $\sh{L}$ là một $\sh{O}_X$-môđun khả nghịch.
Nếu $\sh{L}$ rất ample tương đối với $f$, thì nó ample tương đối với $f$.
\end{proposition}

\begin{proof}
Điều này suy ra từ tính địa phương (trên $Y$) của khái niệm bó rất ample \sref{II.4.4.5}, từ định nghĩa \sref{II.4.6.1} và từ tiêu chuẩn~\sref{II.4.5.10}[(i)].
\end{proof}

\begin{proposition}[4.6.3]
\label{II.4.6.3}
Cho $f:X\to Y$ là một cấu xạ quasi-compact, $\sh{L}$ là một $\sh{O}_X$-môđun khả nghịch, và cho $\sh{S}$ là $\sh{O}_Y$-đại số phân bậc $\bigoplus_{n\geq0}f_*(\sh{L}^{\otimes n})$.
Các điều kiện sau tương đương:
\begin{enumerate}
  \item[\rm{(a)}] $\sh{L}$ là $f$-ample.
  \item[\rm{(b)}] $\sh{S}$ quasi-coherent, và đồng cấu chính tắc $\sigma:f^*(\sh{S})\to\bigoplus_{n\geq0}\sh{L}^{\otimes n}$ \sref[0]{0.4.4.3} có tính chất rằng $Y$-cấu xạ $r_{\sh{L},\sigma}:G(\sigma)\to\Proj(\sh{S})=P$ được xác định khắp nơi và là một phép nhúng mở trội.
  \item[\rm{(b')}] Cấu xạ $f$ tách, và $Y$-cấu xạ $r_{\sh{L},\sigma}$ được xác định khắp nơi và là một phép đồng phôi từ không gian nền của $X$ lên một không gian con của $\Proj(\sh{S})$.
\end{enumerate}

Hơn nữa, khi các điều kiện ấy được thỏa mãn, với mọi $n\in\bb{Z}$, đồng cấu chính tắc
\[
\label{II.4.6.3.1}
  r_{\sh{L},\sigma}^*(\sh{O}_P(n))\to\sh{L}^{\otimes n}
\tag{4.6.3.1}
\]
được định nghĩa trong \sref{II.3.7.9.1} là một đẳng cấu.

Cuối cùng, với mọi $\sh{O}_X$-môđun quasi-coherent $\sh{F}$, nếu đặt $\sh{M}=\bigoplus_{n\geq0}f_*(\sh{F}\otimes\sh{L}^{\otimes n})$, thì đồng cấu chính tắc
\[
\label{II.4.6.3.2}
  r_{\sh{L},\sigma}^*(\widetilde{\sh{M}})\to\sh{F}
\tag{4.6.3.2}
\]
được định nghĩa trong \sref{II.3.7.9.2} là một đẳng cấu.
\end{proposition}

\begin{proof}
Ta đã lưu ý rằng (a) kéo theo $f$ tách, nên $\sh{S}$ quasi-coherent \sref[I]{I.9.2.2}[(a)].
Vì tính chất $r_{\sh{L},\sigma}$ là một phép nhúng được xác định khắp nơi mang tính địa phương trên $Y$, để chứng minh (a) kéo theo (b), ta có thể giả sử $Y$ affine và $\sh{L}$ ample;
\oldpage[II]{90}
khẳng định khi đó suy ra từ \sref{II.4.5.2}[(b)].
Hiển nhiên (b) kéo theo (b');
cuối cùng, để chứng minh (b') kéo theo (a), chỉ cần xét một phủ của $Y$ bởi các tập mở affine $(U_\alpha)$ và áp dụng tiêu chuẩn~\sref{II.4.5.2}[(b')] cho từng bó $\sh{L}|f^{-1}(U_\alpha)$.

Đối với hai khẳng định cuối, ta dùng sự kiện rằng $\sigma^\flat$ ở đây là đồng nhất, cùng với sự mô tả tường minh của các đồng cấu \sref{II.3.7.9.1} và \sref{II.3.7.9.2};
từ đó suy ra ngay rằng \sref{II.4.6.3.1} là một đẳng cấu.
Đối với \sref{II.4.6.3.2}, ta có thể quy về trường hợp $Y$ affine, do đó $\sh{L}$ ample;
hiển nhiên đồng cấu \sref{II.4.6.3.2} là đơn ánh, và tiêu chuẩn~\sref{II.4.5.2}[(c)] cho thấy nó là toàn ánh, từ đó có kết luận.
\end{proof}

\begin{corollary}[4.6.4]
\label{II.4.6.4}
Cho $(U_\alpha)$ là một phủ mở của $Y$.
Để $\sh{L}$ ample đối với $Y$, điều kiện cần và đủ là $\sh{L}|f^{-1}(U_\alpha)$ ample đối với $U_\alpha$ với mọi $\alpha$.
\end{corollary}

\begin{proof}
Thật vậy, điều kiện~(b) mang tính địa phương trên $Y$.
\end{proof}

\begin{corollary}[4.6.5]
\label{II.4.6.5}
Cho $\sh{K}$ là một $\sh{O}_Y$-môđun khả nghịch.
Để $\sh{L}$ là $Y$-ample, điều kiện cần và đủ là $\sh{L}\otimes f^*(\sh{K})$ là $Y$-ample.
\end{corollary}

\begin{proof}
Đây là hệ quả hiển nhiên của \sref{II.4.6.4}, bằng cách chọn các $U_\alpha$ sao cho $\sh{K}|U_\alpha$ đẳng cấu với $\sh{O}_Y|U_\alpha$ với mọi $\alpha$.
\end{proof}

\begin{corollary}[4.6.6]
\label{II.4.6.6}
Giả sử $Y$ affine;
để $\sh{L}$ là $Y$-ample, điều kiện cần và đủ là $\sh{L}$ ample.
\end{corollary}

\begin{proof}
Đây là hệ quả trực tiếp của định nghĩa \sref{II.4.6.1} và các tiêu chuẩn \sref{II.4.6.3}[(b)] và \sref{II.4.5.2}[(b)], vì ở đây $\Proj(\sh{S})=\Proj(\Gamma(Y,\sh{S}))$ theo định nghĩa.
\end{proof}

\begin{corollary}[4.6.7]
\label{II.4.6.7}
Cho $f:X\to Y$ là một cấu xạ quasi-compact.
Giả sử tồn tại một $\sh{O}_Y$-môđun quasi-coherent $\sh{E}$ và một $Y$-cấu xạ $g:X\to P=\bb{P}(\sh{E})$ là một phép đồng phôi từ không gian nền của $X$ lên một không gian con của $P$;
khi đó $\sh{L}=g^*(\sh{O}_P(1))$ là $Y$-ample.
\end{corollary}

\begin{proof}
Ta có thể giả sử $Y$ affine;
khi đó hệ quả suy ra từ tiêu chuẩn~\sref{II.4.5.2}[(a)], công thức~\sref{II.3.7.3.1} và \sref{II.4.2.3}.
\end{proof}

\begin{proposition}[4.6.8]
\label{II.4.6.8}
Cho $X$ là một lược đồ quasi-compact hoặc một tiền lược đồ có không gian nền Noether, và $f:X\to Y$ là một cấu xạ tách quasi-compact.
Để một $\sh{O}_X$-môđun khả nghịch $\sh{L}$ là $f$-ample, điều kiện cần và đủ là một trong các điều kiện tương đương sau được thỏa mãn:
\begin{enumerate}
  \item[\rm{(c)}] Với mọi $\sh{O}_X$-môđun quasi-coherent hữu hạn sinh $\sh{F}$, tồn tại một số nguyên $n_0>0$ sao cho, với mọi $n\geq n_0$, đồng cấu chính tắc $\sigma:f^*(f_*(\sh{F}\otimes\sh{L}^{\otimes n}))\to\sh{F}\otimes\sh{L}^{\otimes n}$ là toàn ánh.
  \item[\rm{(c')}] Với mọi bó iđêan quasi-coherent hữu hạn sinh $\sh{J}$ của $\sh{O}_X$, tồn tại một số nguyên $n>0$ sao cho đồng cấu chính tắc $\sigma:f^*(f_*(\sh{J}\otimes\sh{L}^{\otimes n}))\to\sh{J}\otimes\sh{L}^{\otimes n}$ là toàn ánh.
\end{enumerate}
\end{proposition}

\begin{proof}
Vì $X$ quasi-compact, $f(X)$ cũng quasi-compact, nên tồn tại một phủ hữu hạn $(U_i)$ của $f(X)$ bởi các tập mở affine $U_i$ của $Y$.
Để chứng minh điều kiện~(c) khi $\sh{L}$ là $f$-ample, ta có thể thay $Y$ bằng các $U_i$ và $X$ bằng các $f^{-1}(U_i)$, bởi vì nếu với mỗi $i$ ta tìm được một số nguyên $n_i$ sao cho (c) đúng (đối với $U_i$, $f^{-1}(U_i)$ và $\sh{L}|f^{-1}(U_i)$) với mọi $n\geq n_i$, thì chỉ cần lấy $n_0$ là số lớn nhất trong các $n_i$ để có (c) đối với $Y$, $X$ và $\sh{L}$.
Nhưng khi $Y$ affine, điều kiện~(c) suy ra từ \sref{II.4.5.5}[(d)], có tính đến \sref{II.4.6.6}.
Hiển nhiên (c) kéo theo (c').
Cuối cùng, để chứng minh (c') kéo theo $\sh{L}$ là $f$-ample, ta lại có thể quy về trường hợp $Y$ affine:
thật vậy, mọi bó iđêan quasi-coherent hữu hạn sinh $\sh{J}_i$ của $\sh{O}_X|f^{-1}(U_i)$ là hạn chế của một bó iđêan quasi-coherent hữu hạn sinh của $\sh{O}_X$ \sref[I]{I.9.4.7}, và giả thiết~(c') kéo theo rằng
\oldpage[II]{91}
$\sh{J}_i\otimes(\sh{L}^{\otimes n}|f^{-1}(U_i))$ được sinh bởi các tiết diện của nó (có tính đến \sref[I]{I.9.2.2} và \sref{II.3.4.7});
do đó chỉ cần áp dụng tiêu chuẩn~\sref{II.4.5.5}[(d'')].
\end{proof}

\begin{proposition}[4.6.9]
\label{II.4.6.9}
Cho $f:X\to Y$ là một cấu xạ quasi-compact và $\sh{L}$ là một $\sh{O}_X$-môđun khả nghịch.
\begin{enumerate}
  \item[\rm{(i)}] Cho $n>0$ là một số nguyên.
    Để $\sh{L}$ là $f$-ample, điều kiện cần và đủ là $\sh{L}^{\otimes n}$ là $f$-ample.
  \item[\rm{(ii)}] Cho $\sh{L}'$ là một $\sh{O}_X$-môđun khả nghịch, và giả sử tồn tại một số nguyên $n>0$ sao cho đồng cấu chính tắc $\sigma:f^*(f_*(\sh{L}'^{\otimes n}))\to\sh{L}'^{\otimes n}$ là toàn ánh.
    Khi đó, nếu $\sh{L}$ là $f$-ample thì $\sh{L}\otimes\sh{L}'$ cũng vậy.
\end{enumerate}
\end{proposition}

\begin{proof}
Ta có thể quy ngay về trường hợp $Y$ affine, và khi đó mệnh đề là hệ quả trực tiếp của \sref{II.4.5.6}.
\end{proof}

\begin{corollary}[4.6.10]
\label{II.4.6.10}
Tích tenxơ của hai $\sh{O}_X$-môđun $f$-ample là $f$-ample.
\end{corollary}

\begin{proposition}[4.6.11]
\label{II.4.6.11}
Cho $Y$ là một tiền lược đồ quasi-compact, $f:X\to Y$ là một cấu xạ kiểu hữu hạn, và $\sh{L}$ là một $\sh{O}_X$-môđun khả nghịch.
Để $\sh{L}$ là $f$-ample, điều kiện cần và đủ là nó có một trong các tính chất tương đương sau:
\begin{enumerate}
  \item[\rm{(d)}] Tồn tại $n_0>0$ sao cho, với mọi số nguyên $n\geq n_0$, $\sh{L}^{\otimes n}$ rất ample đối với $f$.
  \item[\rm{(d')}] Tồn tại $n>0$ sao cho $\sh{L}^{\otimes n}$ rất ample đối với $f$.
\end{enumerate}
\end{proposition}

\begin{proof}
Nếu $\sh{L}$ ample đối với $f$, thì tồn tại một phủ \emph{hữu hạn} $(U_i)$ của $Y$ bởi các tập mở affine sao cho các $\sh{L}|f^{-1}(U_i)$ ample.
Từ \sref{II.4.5.10}, ta suy ra tồn tại một số nguyên $n_0$ sao cho $\sh{L}^{\otimes n}|f^{-1}(U_i)$ rất ample đối với $f^{-1}(U_i)\to U_i$ với mọi $n\geq n_0$ và mọi $i$, do đó $\sh{L}^{\otimes n}$ rất ample đối với $f$ \sref{II.4.4.5}.
Ngược lại, (d') đã kéo theo $\sh{L}^{\otimes n}$ là $f$-ample \sref{II.4.6.2}, nên $\sh{L}$ cũng vậy \sref{II.4.6.9}[(i)].
\end{proof}

\begin{corollary}[4.6.12]
\label{II.4.6.12}
Cho $Y$ là một tiền lược đồ quasi-compact, $f:X\to Y$ là một cấu xạ kiểu hữu hạn, và $\sh{L}$, $\sh{L}'$ là hai $\sh{O}_X$-môđun khả nghịch.
Nếu $\sh{L}$ là $f$-ample, thì tồn tại $n_0$ sao cho $\sh{L}^{\otimes n}\otimes\sh{L}'$ rất ample đối với $f$ với mọi $n\geq n_0$.
\end{corollary}

\begin{proof}
Ta lập luận như trong \sref{II.4.6.11}, dùng một phủ hữu hạn của $Y$ bởi các tập mở affine và \sref{II.4.5.11}.
\end{proof}

\begin{proposition}[4.6.13]
\label{II.4.6.13}
\begin{enumerate}
  \item[\rm{(i)}] Với mọi tiền lược đồ $Y$, mọi $\sh{O}_Y$-môđun khả nghịch $\sh{L}$ đều ample đối với cấu xạ đồng nhất $1_Y$.
  \item[\rm{(i \emph{bis})}] Cho $f:X\to Y$ là một cấu xạ quasi-compact và $j:X'\to X$ là một cấu xạ quasi-compact, đồng thời là một phép đồng phôi từ không gian nền của $X'$ lên một không gian con của $X$.
    Nếu $\sh{L}$ là một $\sh{O}_X$-môđun ample đối với $f$, thì $j^*(\sh{L})$ ample đối với $f\circ j$.
  \item[\rm{(ii)}] Cho $Z$ là một tiền lược đồ quasi-compact, $f:X\to Y$ và $g:Y\to Z$ là hai cấu xạ quasi-compact, $\sh{L}$ là một $\sh{O}_X$-môđun ample đối với $f$, và $\sh{K}$ là một $\sh{O}_Y$-môđun ample đối với $g$.
    Khi đó tồn tại một số nguyên $n_0>0$ sao cho $\sh{L}\otimes f^*(\sh{K}^{\otimes n})$ ample đối với $g\circ f$ với mọi $n\geq n_0$.
  \item[\rm{(iii)}] Cho $f:X\to Y$ là một cấu xạ quasi-compact và $g:Y'\to Y$ là một cấu xạ, và đặt $X'=X_{(Y')}$.
    Nếu $\sh{L}$ là một $\sh{O}_X$-môđun ample đối với $f$, thì $\sh{L}'=\sh{L}\otimes_Y\sh{O}_{Y'}$ là một $\sh{O}_{X'}$-môđun ample đối với $f_{(Y')}$.
  \item[\rm{(iv)}] Cho $f_i:X_i\to Y_i$ ($i=1,2$) là hai $S$-cấu xạ quasi-compact.
    Nếu $\sh{L}_i$ là một $\sh{O}_{X_i}$-môđun ample đối với $f_i$ ($i=1,2$), thì $\sh{L}_1\otimes_S\sh{L}_2$ ample đối với $f_1\times_S f_2$.
  \item[\rm{(v)}] Cho $f:X\to Y$ và $g:Y\to Z$ là hai cấu xạ sao cho $g\circ f$ quasi-compact.
    Nếu một
\oldpage[II]{92}
    $\sh{O}_X$-môđun $\sh{L}$ ample đối với $g\circ f$, và nếu $g$ tách hoặc không gian nền của $X$ là Noether địa phương, thì $\sh{L}$ ample đối với $f$.
  \item[\rm{(vi)}] Cho $f:X\to Y$ là một cấu xạ quasi-compact và $j$ là phép nhúng chính tắc $X_\red\to X$.
    Nếu $\sh{L}$ là một $\sh{O}_X$-môđun ample đối với $f$, thì $j^*(\sh{L})$ ample đối với $f_\red$.
\end{enumerate}
\end{proposition}

\begin{proof}
Trước hết, lưu ý rằng (v) và (vi) suy ra từ (i), (i~\emph{bis}) và (iv) bằng cùng lập luận như trong \sref{II.4.4.10}, dùng \sref{II.4.6.4} thay cho \sref{II.4.4.5};
ta để chi tiết của lập luận cho người đọc.
Khẳng định~(i) là hệ quả hiển nhiên của \sref{II.4.4.10}[(i)] và \sref{II.4.6.2}.
Để chứng minh (i~\emph{bis}), (iii) và (iv), ta sẽ dùng bổ đề sau:
  \begin{lemma}[4.6.13.1]
  \label{II.4.6.13.1}
  \begin{enumerate}
    \item[\rm{(i)}] Cho $u:Z\to S$ là một cấu xạ, $\sh{L}$ là một $\sh{O}_S$-môđun khả nghịch, $s$ là một tiết diện của $\sh{L}$ trên $S$, và $s'$ là tiết diện tương ứng chính tắc của $u^*(\sh{L})=\sh{L}'$ trên $Z$.
      Khi đó $Z_{s'}=u^{-1}(S_s)$.
    \item[\rm{(ii)}] Cho $Z$, $Z'$ là hai $S$-tiền lược đồ, $p$, $p'$ là các phép chiếu của $T=Z\times_S Z'$, $\sh{L}$ (tương ứng, $\sh{L}'$) là một $\sh{O}_Z$-môđun khả nghịch (tương ứng, một $\sh{O}_{Z'}$-môđun khả nghịch), $t$ (tương ứng, $t'$) là một tiết diện của $\sh{L}$ (tương ứng, $\sh{L}'$) trên $Z$ (tương ứng, $Z'$), và $s$ (tương ứng, $s'$) là tiết diện tương ứng của $p^*(\sh{L})$ (tương ứng, $p'^*(\sh{L}')$) trên $Z\times_S Z'$.
      Khi đó $T_{s\otimes s'}=Z_t\times_S Z'_{t'}$.
  \end{enumerate}
  \end{lemma}

  \begin{proof}
  Từ các định nghĩa suy ra rằng ta có thể quy về trường hợp mọi tiền lược đồ đang xét đều affine.
  Hơn nữa, trong (i), ta có thể giả sử $\sh{L}=\sh{O}_S$;
  khi đó khẳng định~(i) suy ra ngay từ \sref[I]{I.1.2.2.2}.
  Tương tự, trong (ii), ta có thể quy về trường hợp $\sh{L}=\sh{O}_Z$ và $\sh{L}'=\sh{O}_{Z'}$, và khi đó khẳng định quy về Bổ đề~\sref{II.4.3.2.4}.
  \end{proof}

Bây giờ ta chứng minh (i~\emph{bis}).
Ta có thể giả sử $Y$ affine \sref{II.4.6.4}, do đó $\sh{L}$ ample \sref{II.4.6.6};
khi $s$ chạy qua hợp của các tập $\Gamma(X,\sh{L}^{\otimes n})$ ($n>0$), các $X_s$ lập thành một cơ sở của tôpô trên $X$ \sref{II.4.5.2}[(a)], nên theo giả thiết, các $j^{-1}(X_s)$ lập thành một cơ sở của tôpô trên $X'$;
do đó, từ Bổ đề~\sref{II.4.6.13.1}[(i)] và \sref{II.4.5.2}[(a)], ta kết luận rằng $j^*(\sh{L})$ ample.

Tiếp theo ta chứng minh (iii).
Ta lại có thể giả sử $Y$ và $Y'$ affine \sref{II.4.6.4}, từ đó phép chiếu $h:X'\to X$ là affine \sref{II.1.5.5}.
Vì $\sh{L}$ ample \sref{II.4.6.6}, khi $s$ chạy qua tập các tiết diện của $\sh{L}^{\otimes n}$ ($n>0$) trên $X$ sao cho $X_s$ affine, thì các $X_s$ phủ $X$ \sref{II.4.5.2}[(a')], nên các $h^{-1}(X_s)$ affine \sref{II.1.2.5} và phủ $X'$;
do đó, lại từ Bổ đề~\sref{II.4.6.13.1}[(i)] và \sref{II.4.5.2}[(a')], ta suy ra rằng $\sh{L}'$ ample, vì cấu xạ $f_{(Y')}$ quasi-compact \sref[I]{I.6.6.4}[(iii)].

Để chứng minh (iv), trước hết lưu ý rằng $f_1\times_S f_2$ quasi-compact \sref[I]{I.6.6.4}[(iv)].
Ta còn có thể giả sử $S$, $Y_1$ và $Y_2$ affine (\sref{II.4.6.4} và \sref[I]{I.3.2.7}), do đó $\sh{L}_i$ ample ($i=1,2$) \sref{II.4.6.6}.
Các tập mở $(X_1)_{s_1}\times_S(X_2)_{s_2}$ lập thành một phủ của $X_1\times_S X_2$, trong đó $s_i$ chạy qua các tiết diện của $\sh{L}_i^{\otimes n_i}$ sao cho $(X_i)_{s_i}$ affine \sref{II.4.5.2}[(a')].
Hơn nữa, bằng cách thay $s_1$ và $s_2$ bởi các lũy thừa thích hợp, điều này không thay đổi các $(X_i)_{s_i}$, ta luôn có thể giả sử $n_1=n_2$.
Từ \sref{II.4.6.13.1}[(ii)] và \sref{II.4.5.2}[(a')], ta suy ra $\sh{L}_1\otimes_S\sh{L}_2$ ample, từ đó có khẳng định, vì $Y_1\times_S Y_2$ affine \sref{II.4.6.6}.

Còn lại phải chứng minh (ii).
Bằng cùng lập luận như trong \sref{II.4.4.10}, nhưng ở đây dùng \sref{II.4.6.4}, ta có thể quy về trường hợp $Z$ affine.
Vì khi đó $\sh{K}$ ample và $Y$ quasi-compact, tồn tại một số hữu hạn các tiết diện $s_i\in\Gamma(Y,\sh{K}^{\otimes k_i})$ sao cho các $Y_{s_i}$
\oldpage[II]{93}
là \emph{affine} và phủ $Y$ \sref{II.4.5.2}[(a')];
bằng cách thay các $s_i$ bởi những lũy thừa thích hợp, ta còn có thể giả sử mọi $k_i$ đều bằng cùng một số nguyên $k$.
Cho $s'_i$ là các tiết diện của $f^*(\sh{K}^{\otimes k})$ trên $X$ tương ứng chính tắc với các $s_i$, sao cho các $X_{s'_i}=f^{-1}(Y_{s_i})$ \sref{II.4.6.13.1}[(i)] phủ $X$.
Vì $\sh{L}|X_{s'_i}$ ample (\sref{II.4.6.4} và \sref{II.4.6.6}), với mỗi $i$ tồn tại một số hữu hạn các tiết diện $t_{ij}\in\Gamma(X,\sh{L}^{\otimes n_{ij}})$ sao cho các $X_{t_{ij}}$ affine, được chứa trong $X_{s'_i}$ và phủ $X_{s'_i}$ \sref{II.4.5.2}[(a')];
ta cũng có thể giả sử mọi $n_{ij}$ đều bằng cùng một số $n$.
Với các điều kiện đó, $X$ tách và quasi-compact, nên tồn tại một số nguyên $m>0$ và, với mọi $(i,j)$, một tiết diện
\[
  u_{ij}\in\Gamma(X,\sh{L}^{\otimes n}\otimes_X f^*(\sh{K}^{\otimes mk}))
\]
sao cho $t_{ij}\otimes s'^{\otimes m}_i$ là hạn chế của $u_{ij}$ lên $X_{s'_i}$ \sref[I]{I.9.3.1};
hơn nữa, $X_{u_{ij}}=X_{t_{ij}}$, nên các $X_{u_{ij}}$ affine và phủ $X$.
Ta còn có thể giả sử $m$ có dạng $nr$;
nếu đặt $n_0=rk$, thì từ \sref{II.4.5.2}[(a')] ta thấy rằng $\sh{L}\otimes_{\sh{O}_X}f^*(\sh{K}^{\otimes n_0})$ ample.
Hơn nữa, tồn tại $h_0>0$ sao cho $\sh{K}^{\otimes h}$ được sinh bởi các tiết diện của nó trên $Y$ với mọi $h\geq h_0$ \sref{II.4.5.5};
\emph{a fortiori}, $f^*(\sh{K}^{\otimes h})$ được sinh bởi các tiết diện của nó trên $X$ với mọi $h\geq h_0$, theo định nghĩa của các ảnh ngược (\sref[0]{0.3.7.1} và \sref{II.4.4.1}).
Do đó ta kết luận rằng $\sh{L}\otimes f^*(\sh{K}^{\otimes n_0+h})$ ample với mọi $h\geq h_0$ \sref{II.4.5.6}, điều này hoàn tất chứng minh.
\end{proof}

\begin{remark}[4.6.14]
\label{II.4.6.14}
Dưới các điều kiện của (ii), không nên cho rằng $\sh{L}\otimes f^*(\sh{K})$ ample đối với $g\circ f$;
thật vậy, vì $\sh{L}\otimes f^*(\sh{K}^{-1})$ cũng ample đối với $f$ \sref{II.4.6.5}, ta sẽ suy ra $\sh{L}$ ample đối với $g\circ f$;
đặc biệt, lấy $g$ là cấu xạ đồng nhất thì \emph{mọi} $\sh{O}_X$-môđun khả nghịch đều sẽ ample đối với $f$, điều này nói chung không đúng (xem \sref{II.5.1.6}, \sref{II.5.3.4}[(i)] và \sref{II.5.3.1}).
\end{remark}

\begin{proposition}[4.6.15]
\label{II.4.6.15}
Cho $f:X\to Y$ là một cấu xạ quasi-compact, $\sh{J}$ là một bó iđêan quasi-coherent lũy linh địa phương của $\sh{O}_X$, $Z$ là tiền lược đồ con đóng của $X$ được xác định bởi $\sh{J}$, và $j:Z\to X$ là phép nhúng chính tắc.
Để một $\sh{O}_X$-môđun khả nghịch $\sh{L}$ ample đối với $f$, điều kiện cần và đủ là $j^*(\sh{L})$ ample đối với $f\circ j$.
\end{proposition}

\begin{proof}
Vì vấn đề mang tính địa phương trên $Y$ \sref{II.4.6.4}, ta có thể giả sử $Y$ affine;
vì khi đó $X$ quasi-compact, ta có thể giả sử $\sh{J}$ lũy linh.
Có tính đến \sref{II.4.6.6}, mệnh đề khi đó không gì khác hơn \sref{II.4.5.13}.
\end{proof}

\begin{corollary}[4.6.16]
\label{II.4.6.16}
Cho $X$ là một tiền lược đồ Noether địa phương và $f:X\to Y$ là một cấu xạ quasi-compact.
Để một $\sh{O}_X$-môđun khả nghịch $\sh{L}$ ample đối với $f$, điều kiện cần và đủ là ảnh ngược $\sh{L}'$ của nó qua phép nhúng chính tắc $X_\red\to X$ ample đối với $f_\red$.
\end{corollary}

\begin{proof}
Ta đã thấy điều kiện là cần thiết \sref{II.4.6.13}[(vi)];
ngược lại, nếu điều kiện được thỏa mãn, thì để chứng minh $\sh{L}$ ample đối với $f$, ta có thể quy về trường hợp $Y$ affine \sref{II.4.6.4};
khi đó $Y_\red$ cũng affine, nên $\sh{L}'$ ample \sref{II.4.6.6}, và $\sh{L}$ cũng vậy theo \sref{II.4.5.13}, vì khi ấy $X$ là Noether và $X_\red$ là một tiền lược đồ con đóng của $X$ được xác định bởi một bó iđêan quasi-coherent lũy linh \sref[I]{I.6.1.6}.
\end{proof}

\begin{proposition}[4.6.17]
\label{II.4.6.17}
Với các ký hiệu và giả thiết của \sref{II.4.4.11}, để $\sh{L}''$ ample đối với $f''$, điều kiện cần và đủ là $\sh{L}$ ample đối với $f$ và $\sh{L}'$ ample đối với $f'$.
\end{proposition}
\oldpage[II]{94}
\begin{proof}
Tính cần thiết của điều kiện suy ra từ \sref{II.4.6.13}[(i~\emph{bis})].
Để thấy điều kiện là đủ, ta có thể quy về trường hợp $Y$ affine, và khi đó việc $\sh{L}''$ ample suy ra từ tiêu chuẩn~\sref{II.4.5.2}[(a)] áp dụng cho $\sh{L}$, $\sh{L}'$ và $\sh{L}''$, lưu ý rằng một tiết diện của $\sh{L}$ trên $X$ có thể được kéo dài (bằng $0$) thành một tiết diện của $\sh{L}''$ trên $X''$.
\end{proof}

\begin{proposition}[4.6.18]
\label{II.4.6.18}
Cho $Y$ là một tiền lược đồ quasi-compact, $\sh{S}$ là một $\sh{O}_Y$-đại số phân bậc quasi-coherent hữu hạn sinh, $X=\Proj(\sh{S})$, và $f:X\to Y$ là cấu xạ cấu trúc.
Khi đó $f$ thuộc kiểu hữu hạn, và tồn tại một số nguyên $d>0$ sao cho $\sh{O}_X(d)$ khả nghịch và $f$-ample.
\end{proposition}

\begin{proof}
Theo \sref{II.3.1.10}, tồn tại một số nguyên $d>0$ sao cho $\sh{S}^{(d)}$ được sinh bởi $\sh{S}_d$.
Ta biết rằng, dưới đẳng cấu chính tắc giữa $X$ và $X^{(d)}=\Proj(\sh{S}^{(d)})$, $\sh{O}_X(d)$ được đồng nhất với $\sh{O}_{X^{(d)}}(1)$ \sref{II.3.2.9}[(ii)].
Do đó ta thấy có thể quy về trường hợp $\sh{S}$ được sinh bởi $\sh{S}_1$;
khi đó mệnh đề suy ra từ \sref{II.4.4.3} và \sref{II.4.6.2} (có tính đến việc $f$ là một cấu xạ kiểu hữu hạn \sref{II.3.4.1}).
\end{proof}

\section{Các cấu xạ quasi-affine; các cấu xạ quasi-projective; các cấu xạ riêng; các cấu xạ xạ ảnh}
\label{section:II.5}

\subsection{Các cấu xạ quasi-affine}
\label{subsection:II.5.1}

\begin{definition}[5.1.1]
\label{II.5.1.1}
Ta gọi một lược đồ là quasi-affine nếu nó đẳng cấu với một lược đồ con cảm sinh trên một tập mở quasi-compact của một lược đồ affine.
Ta nói một cấu xạ $f:X\to Y$ là quasi-affine, hoặc $X$ (được xét như một $Y$-tiền lược đồ qua $f$) là một $Y$-lược đồ quasi-affine, nếu tồn tại một phủ $(U_\alpha)$ của $Y$ bởi các tập mở affine sao cho các $f^{-1}(U_\alpha)$ là những lược đồ quasi-affine.
\end{definition}

Hiển nhiên một cấu xạ quasi-affine là \emph{tách} (\sref[I]{I.5.5.5} và \sref[I]{I.5.5.8}) và \emph{quasi-compact} \sref[I]{I.6.6.1};
mọi cấu xạ affine hiển nhiên là quasi-affine.

Nhắc lại rằng, với mọi tiền lược đồ $X$, nếu đặt $A=\Gamma(X,\sh{O}_X)$, đồng cấu đồng nhất $A\to A=\Gamma(X,\sh{O}_X)$ xác định một cấu xạ $X\to\Spec(A)$, gọi là \emph{chính tắc} \sref[I]{I.2.2.4};
đây không gì khác hơn cấu xạ chính tắc được định nghĩa trong \sref{II.4.5.1} đối với trường hợp riêng $\sh{L}=\sh{O}_X$, nếu nhớ rằng $\Proj(A[T])$ được đồng nhất chính tắc với $\Spec(A)$ \sref{II.3.1.7}.

\begin{proposition}[5.1.2]
\label{II.5.1.2}
Cho $X$ là một lược đồ quasi-compact hoặc một tiền lược đồ có không gian nền Noether, và $A$ là vành $\Gamma(X,\sh{O}_X)$.
Các điều kiện sau là tương đương:
\begin{enumerate}
  \item[{\rm(a)}] $X$ là một lược đồ quasi-affine.
  \item[{\rm(b)}] Cấu xạ chính tắc $u:X\to\Spec(A)$ là một phép nhúng mở.
  \item[{\rm(b')}] Cấu xạ chính tắc $u:X\to\Spec(A)$ là một phép đồng phôi từ $X$ lên một không gian con của không gian nền của $\Spec(A)$.
  \item[{\rm(c)}] $\sh{O}_X$-môđun $\sh{O}_X$ rất ample đối với $u$ \sref{II.4.4.2}.
  \item[{\rm(c')}] $\sh{O}_X$-môđun $\sh{O}_X$ ample \sref{II.4.5.1}.
  \item[{\rm(d)}] Khi $f$ chạy qua $A$, các $X_f$ lập thành một cơ sở của tôpô trên $X$.
  \item[{\rm(d')}] Khi $f$ chạy qua $A$, những $X_f$ nào affine lập thành một phủ của $X$.
\oldpage[II]{95}
  \item[{\rm(e)}] Mọi $\sh{O}_X$-môđun quasi-coherent đều được sinh bởi các tiết diện của nó trên $X$.
  \item[{\rm(e')}] Mọi bó iđêan quasi-coherent hữu hạn sinh của $\sh{O}_X$ đều được sinh bởi các tiết diện của nó trên $X$.
\end{enumerate}
\end{proposition}

\begin{proof}
Hiển nhiên (b) kéo theo (a), và (a) kéo theo (c) nhờ tiêu chuẩn~\sref{II.4.4.4}[(b)] áp dụng cho cấu xạ đồng nhất (có xét đến nhận xét đứng trước phát biểu);
mặt khác, (c) kéo theo (c$'$) \sref{II.4.5.10}[(i)], và (c$'$), (b) cùng (b$'$) tương đương với nhau nhờ \sref{II.4.5.2}[(b)] và \sref{II.4.5.2}[(b$'$)].
Cuối cùng, (c$'$) tương đương với từng điều kiện (d), (d$'$), (e) và (e$'$), nhờ \sref{II.4.5.2}[(a)], \sref{II.4.5.2}[(a$'$)], \sref{II.4.5.2}[(c)] và \sref{II.4.5.5}[(d$''$)].
\end{proof}

Ngoài ra, ta nhận thấy rằng, với các ký hiệu trên, những $X_f$ affine lập thành một \emph{cơ sở} của tôpô trên $X$, và cấu xạ chính tắc $u$ là \emph{xạ trội} \sref{II.4.5.2}.

\begin{corollary}[5.1.3]
\label{II.5.1.3}
Cho $X$ là một tiền lược đồ quasi-compact.
Nếu tồn tại một cấu xạ $v:X\to Y$ từ $X$ vào một lược đồ affine $Y$, là một phép đồng phôi từ $X$ lên một không gian con mở của $Y$, thì $X$ là quasi-affine.
\end{corollary}

\begin{proof}
Tồn tại một họ $(g_\alpha)$ các tiết diện của $\sh{O}_Y$ trên $Y$ sao cho các $D(g_\alpha)$ lập thành một cơ sở của tôpô trên $v(X)$;
nếu $v=(\psi,\theta)$ và ta đặt $f_\alpha=\theta(g_\alpha)$, thì $X_{f_\alpha}=\psi^{-1}(D(g_\alpha))$ \sref[I]{I.2.2.4.1}; do đó các $X_{f_\alpha}$ lập thành một cơ sở của tôpô trên $X$, và tiêu chuẩn~\sref{II.5.1.2}[(d)] được thỏa mãn.
\end{proof}

\begin{corollary}[5.1.4]
\label{II.5.1.4}
Nếu $X$ là một lược đồ quasi-affine, thì \emph{mọi} $\sh{O}_X$-môđun khả nghịch đều rất ample (tương đối với cấu xạ chính tắc) và, \emph{a fortiori}, ample.
\end{corollary}

\begin{proof}
Thật vậy, một môđun như thế $\sh{L}$ được sinh bởi các tiết diện của nó trên $X$ \sref{II.5.1.2}[(e)]; do đó $\sh{L}\otimes\sh{O}_X=\sh{L}$ là rất ample \sref{II.4.4.8}.
\end{proof}

\begin{corollary}[5.1.5]
\label{II.5.1.5}
Cho $X$ là một tiền lược đồ quasi-compact.
Nếu tồn tại một $\sh{O}_X$-môđun khả nghịch $\sh{L}$ sao cho $\sh{L}$ và $\sh{L}^{-1}$ đều ample, thì $X$ là một lược đồ quasi-affine.
\end{corollary}

\begin{proof}
Thật vậy, $\sh{O}_X=\sh{L}\otimes\sh{L}^{-1}$ khi đó là ample \sref{II.4.5.7}.
\end{proof}

\begin{proposition}[5.1.6]
\label{II.5.1.6}
Cho $f:X\to Y$ là một cấu xạ quasi-compact.
Các điều kiện sau tương đương:
\begin{enumerate}
  \item[{\rm(a)}] Cấu xạ $f$ là quasi-affine.
  \item[{\rm(b)}] $\sh{O}_Y$-đại số $f_*(\sh{O}_X)=\sh{A}(X)$ là quasi-coherent, và cấu xạ chính tắc $X\to\Spec(\sh{A}(X))$ tương ứng với đồng cấu đồng nhất $\sh{A}(X)\to\sh{A}(X)$ \sref{II.1.2.7} là một phép nhúng mở.
  \item[{\rm(b')}] $\sh{O}_Y$-đại số $\sh{A}(X)$ là quasi-coherent, và cấu xạ chính tắc $X\to\Spec(\sh{A}(X))$ là một phép đồng phôi từ $X$ lên một không gian con của $\Spec(\sh{A}(X))$.
  \item[{\rm(c)}] $\sh{O}_X$-môđun $\sh{O}_X$ rất ample đối với $f$.
  \item[{\rm(c')}] $\sh{O}_X$-môđun $\sh{O}_X$ ample đối với $f$.
  \item[{\rm(d)}] Cấu xạ $f$ là tách và, với mọi $\sh{O}_X$-môđun quasi-coherent $\sh{F}$, đồng cấu chính tắc $\sigma:f^*(f_*(\sh{F}))\to\sh{F}$ \sref[0]{0.4.4.3} là toàn ánh.
\end{enumerate}

Hơn nữa, khi $f$ là quasi-affine, mọi $\sh{O}_X$-môđun khả nghịch $\sh{L}$ đều rất ample đối với $f$.
\end{proposition}

\begin{proof}
Sự tương đương giữa (a) và (c$'$) suy ra từ tính địa phương trên $Y$ của tính $f$-ample \sref{II.4.6.4}, từ Định nghĩa~\sref{II.5.1.1} và từ tiêu chuẩn~\sref{II.5.1.2}[(c$'$)].
Các tính chất còn lại có tính địa phương trên $Y$
\oldpage[II]{96}
và do đó suy ra ngay từ \sref{II.5.1.2} và \sref{II.5.1.4}, có xét đến việc $f_*(\sh{F})$ là quasi-coherent khi $f$ tách \sref[I]{I.9.2.2}[(a)].
\end{proof}

\begin{corollary}[5.1.7]
\label{II.5.1.7}
Cho $f:X\to Y$ là một cấu xạ quasi-affine.
Với mọi tập mở $U$ của $Y$, hạn chế $f^{-1}(U)\to U$ của $f$ là quasi-affine.
\end{corollary}

\begin{corollary}[5.1.8]
\label{II.5.1.8}
Cho $Y$ là một lược đồ affine và $f:X\to Y$ là một cấu xạ quasi-compact.
Điều kiện cần và đủ để $f$ là quasi-affine là $X$ là một lược đồ quasi-affine.
\end{corollary}

\begin{proof}
Đây là hệ quả trực tiếp của \sref{II.5.1.6} và \sref{II.4.6.6}.
\end{proof}

\begin{corollary}[5.1.9]
\label{II.5.1.9}
Cho $Y$ là một lược đồ quasi-compact hoặc một tiền lược đồ có không gian nền Noether, và $f:X\to Y$ là một cấu xạ \emph{thuộc kiểu hữu hạn}.
Nếu $f$ là quasi-affine, thì tồn tại một $\sh{O}_Y$-đại số con quasi-coherent $\sh{B}$ của $\sh{A}(X)=f_*(\sh{O}_X)$, \emph{hữu hạn sinh} \sref[I]{I.9.6.2}, sao cho cấu xạ $X\to\Spec(\sh{B})$ tương ứng với phép nhúng chính tắc $\sh{B}\to\sh{A}(X)$ \sref{II.1.2.7} là một phép nhúng.
Hơn nữa, mọi $\sh{O}_Y$-đại số con quasi-coherent hữu hạn sinh $\sh{B}'$ của $\sh{A}(X)$ chứa $\sh{B}$ đều có cùng tính chất ấy.
\end{corollary}

\begin{proof}
Thật vậy, $\sh{A}(X)$ là giới hạn quy nạp của các $\sh{O}_Y$-đại số con quasi-coherent hữu hạn sinh của nó \sref[I]{I.9.6.5};
kết quả khi đó là một trường hợp riêng của \sref{II.3.8.4}, có xét đến sự đồng nhất giữa $\Spec(\sh{A}(X))$ và $\Proj(\sh{A}(X)[T])$ \sref{II.3.1.7}.
\end{proof}

\begin{proposition}[5.1.10]
\label{II.5.1.10}
\medskip\noindent
\begin{enumerate}
  \item[{\rm(i)}] Mọi cấu xạ quasi-compact $X\to Y$ là một phép đồng phôi từ không gian nền của $X$ lên một không gian con của không gian nền của $Y$ (đặc biệt, mọi phép nhúng đóng) đều là quasi-affine.
  \item[{\rm(ii)}] Hợp thành của hai cấu xạ quasi-affine là quasi-affine.
  \item[{\rm(iii)}] Nếu $f:X\to Y$ là một $S$-cấu xạ quasi-affine, thì $f_{(S')}:X_{(S')}\to Y_{(S')}$ là một cấu xạ quasi-affine với mọi mở rộng $S'\to S$ của tiền lược đồ cơ sở.
  \item[{\rm(iv)}] Nếu $f:X\to Y$ và $g:X'\to Y'$ là hai $S$-cấu xạ quasi-affine, thì $f\times_S g$ là quasi-affine.
  \item[{\rm(v)}] Nếu $f:X\to Y$ và $g:Y\to Z$ là hai cấu xạ sao cho $g\circ f$ quasi-affine, và nếu $g$ tách hoặc không gian nền của $X$ là Noether địa phương, thì $f$ là quasi-affine.
  \item[{\rm(vi)}] Nếu $f$ là một cấu xạ quasi-affine, thì $f_\red$ cũng vậy.
\end{enumerate}
\end{proposition}

\begin{proof}
Có xét đến tiêu chuẩn~\sref{II.5.1.6}[(c$'$)], các khẳng định (i), (iii), (iv), (v) và (vi) lần lượt suy ra ngay từ \sref{II.4.6.13}[(i~\emph{bis})], \sref{II.4.6.13}[(iii)], \sref{II.4.6.13}[(iv)], \sref{II.4.6.13}[(v)] và \sref{II.4.6.13}[(vi)].
Để chứng minh (ii), ta có thể quy về trường hợp $Z$ affine, và khi đó khẳng định suy ra trực tiếp từ \sref{II.4.6.13}[(ii)], áp dụng cho $\sh{L}=\sh{O}_X$ và $\sh{K}=\sh{O}_Y$.
\end{proof}

\begin{remark}[5.1.11]
\label{II.5.1.11}
Cho $f:X\to Y$ và $g:Y\to Z$ là hai cấu xạ sao cho $X\times_Z Y$ là Noether địa phương.
Khi đó phép nhúng đồ thị $\Gamma_f:X\to X\times_Z Y$ là quasi-affine, vì nó quasi-compact \sref[I]{I.6.3.5}, và \sref[I]{I.5.5.12} cho thấy rằng, trong (v), kết luận vẫn đúng nếu bỏ giả thiết $g$ tách.
\end{remark}

\begin{proposition}[5.1.12]
\label{II.5.1.12}
Cho $f:X\to Y$ là một cấu xạ quasi-compact và $g:X'\to X$ là một cấu xạ quasi-affine.
Nếu $\sh{L}$ là một $\sh{O}_X$-môđun ample đối với $f$, thì $g^*(\sh{L})$ là một $\sh{O}_{X'}$-môđun ample đối với $f\circ g$.
\end{proposition}

\begin{proof}
Vì $\sh{O}_{X'}$ rất ample đối với $g$, và vấn đề có tính địa phương trên $Y$ \sref{II.4.6.4}, nên từ \sref{II.4.6.13}[(ii)] suy ra rằng tồn tại (khi $Y$ affine) một số nguyên $n$ sao cho
\[
  g^*(\sh{L}^{\otimes n})=(g^*(\sh{L}))^{\otimes n}
\]
là ample đối với $f\circ g$; do đó $g^*(\sh{L})$ ample đối với $f\circ g$ \sref{II.4.6.9}.
\end{proof}
\oldpage[II]{97}
\subsection{Tiêu chuẩn Serre}
\label{subsection:II.5.2}

\begin{theorem}[5.2.1]
\label{II.5.2.1}
\emph{(Tiêu chuẩn Serre).}
Cho $X$ là một lược đồ quasi-compact hoặc một tiền lược đồ có không gian nền Noether.
Các điều kiện sau tương đương:
\begin{enumerate}
  \item[{\rm(a)}] $X$ là một lược đồ affine.
  \item[{\rm(b)}] Tồn tại một họ các phần tử $f_\alpha\in A=\Gamma(X,\sh{O}_X)$ sao cho các $X_{f_\alpha}$ affine và iđêan sinh bởi các $f_\alpha$ trong $A$ bằng $A$.
  \item[{\rm(c)}] Hàm tử $\Gamma(X,\sh{F})$ khớp theo $\sh{F}$ trong phạm trù các $\sh{O}_X$-môđun quasi-coherent; nói cách khác, nếu
    \[
      0\to\sh{F}'\to\sh{F}\to\sh{F}''\to0
      \tag{*}
    \]
    là một dãy khớp các $\sh{O}_X$-môđun quasi-coherent, thì dãy
    \[
      0\to\Gamma(X,\sh{F}')\to\Gamma(X,\sh{F})\to\Gamma(X,\sh{F}'')\to0
    \]
    là khớp.
  \item[{\rm(c')}] Tính chất~{\rm(c)} đúng với mọi dãy khớp {\rm($*$)} các $\sh{O}_X$-môđun quasi-coherent sao cho $\sh{F}$ đẳng cấu với một $\sh{O}_X$-môđun con của một tích hữu hạn $\sh{O}_X^n$.
  \item[{\rm(d)}] $\HH^1(X,\sh{F})=0$ với mọi $\sh{O}_X$-môđun quasi-coherent $\sh{F}$.
  \item[{\rm(d$'$)}] $\HH^1(X,\sh{J})=0$ với mọi bó iđêan quasi-coherent $\sh{J}$ của $\sh{O}_X$.
\end{enumerate}
\end{theorem}

\begin{proof}
Hiển nhiên (a) kéo theo (b);
mặt khác, (b) kéo theo rằng các $X_{f_\alpha}$ phủ $X$, vì theo giả thiết tiết diện $1$ là một tổ hợp tuyến tính của các $f_\alpha$, và khi đó các $D(f_\alpha)$ phủ $\Spec(A)$.
Do đó, khẳng định cuối của \sref{II.4.5.2} kéo theo rằng $X\to\Spec(A)$ là một đẳng cấu.

Ta biết rằng (a) kéo theo (c) \sref[I]{I.1.3.11}, và (c) hiển nhiên kéo theo (c$'$).
Ta chứng minh rằng (c$'$) kéo theo (b).
Trước hết, (c$'$) kéo theo rằng, với mọi điểm \emph{đóng} $x\in X$ và mọi lân cận mở $U$ của $x$, tồn tại $f\in A$ sao cho $x\in X_f\subset U$.
Thật vậy, cho $\sh{J}$ (tương ứng, $\sh{J}'$) là bó iđêan quasi-coherent của $\sh{O}_X$ xác định tiền lược đồ con đóng rút gọn của $X$ có không gian nền là $X\setmin U$ (tương ứng, $(X\setmin U)\cup\{x\}$) \sref[I]{I.5.2.1};
hiển nhiên $\sh{J}'\subset\sh{J}$ và $\sh{J}''=\sh{J}/\sh{J}'$ là một $\sh{O}_X$-môđun quasi-coherent có giá $\{x\}$ và thỏa $\sh{J}_x''=\kres(x)$.
Giả thiết~(c$'$), áp dụng cho dãy khớp $0\to\sh{J}'\to\sh{J}\to\sh{J}''\to0$, cho thấy $\Gamma(X,\sh{J})\to\Gamma(X,\sh{J}'')$ là toàn ánh.
Do đó, tiết diện của $\sh{J}''$ có mầm tại $x$ bằng $1_x$ là ảnh của một tiết diện $f\in\Gamma(X,\sh{J})\subset\Gamma(X,\sh{O}_X)$, và theo định nghĩa ta có $f(x)=1_x$ và $f(y)=0$ trong $X\setmin U$, điều này chứng minh khẳng định trên.
Hơn nữa, nếu $U$ affine thì $X_f$ cũng affine \sref[I]{I.1.3.6}; do đó hợp của các $X_f$ affine ($f\in A$) là một tập mở $Z$ chứa \emph{mọi điểm đóng} của $X$;
vì $X$ là một không gian Kolmogorov quasi-compact, tất yếu $Z=X$ \sref[0]{0.2.1.3}.
Vì $X$ quasi-compact, tồn tại một số \emph{hữu hạn} các phần tử $f_i\in A$ ($1\leq i\leq n$) sao cho các $X_{f_i}$ affine và phủ $X$.
Xét đồng cấu $\sh{O}_X^n\to\sh{O}_X$ xác định bởi các tiết diện $f_i$ \sref[0]{0.5.1.1};
vì với mọi $x\in X$, ít nhất một trong các $(f_i)_x$ khả nghịch, nên đồng cấu này \emph{toàn ánh}, và do đó ta có một dãy khớp $0\to\sh{R}\to\sh{O}_X^n\to\sh{O}_X\to0$, trong đó $\sh{R}$ là một $\sh{O}_X$-môđun con quasi-coherent của $\sh{O}_X^n$.
Khi đó, từ

\oldpage[II]{98}
(c$'$) suy ra rằng đồng cấu tương ứng $\Gamma(X,\sh{O}_X^n)\to\Gamma(X,\sh{O}_X)$ là toàn ánh, điều này chứng minh (b).

Cuối cùng, (a) kéo theo (d) \sref[I]{I.5.1.9.2}, và (d) hiển nhiên kéo theo (d$'$).
Còn phải chứng minh rằng (d$'$) kéo theo (c$'$).
Thật vậy, nếu $\sh{F}'$ là một $\sh{O}_X$-môđun con quasi-coherent của $\sh{O}_X^n$, thì lọc $0\subset\sh{O}_X\subset\sh{O}_X^2\subset\ldots\subset\sh{O}_X^n$ cảm sinh trên $\sh{F}'$ một lọc gồm các $\sh{F}'_k=\sh{F}'\cap\sh{O}_X^k$ ($0\leq k\leq n$), là các $\sh{O}_X$-môđun quasi-coherent \sref[I]{I.4.1.1}, và $\sh{F}'_{k+1}/\sh{F}'_k$ đẳng cấu với một $\sh{O}_X$-môđun con quasi-coherent của $\sh{O}_X^{k+1}/\sh{O}_X^k=\sh{O}_X$, nghĩa là với một bó iđêan quasi-coherent của $\sh{O}_X$.
Do đó, giả thiết~(d$'$) kéo theo $\HH^1(X,\sh{F}'_{k+1}/\sh{F}'_k)=0$;
dãy khớp đối đồng điều $\HH^1(X,\sh{F}'_k)\to\HH^1(X,\sh{F}'_{k+1})\to\HH^1(X,\sh{F}'_{k+1}/\sh{F}'_k)=0$ khi đó cho phép chứng minh bằng quy nạp theo $k$ rằng $\HH^1(X,\sh{F}'_k)=0$ với mọi $k$.
\end{proof}

\begin{remark}[5.2.1.1]
\label{II.5.2.1.1}
Khi $X$ là một tiền lược đồ \emph{Noether}, trong phát biểu của (c$'$) và (d$'$), ta có thể thay ``quasi-coherent'' bằng ``coherent''.
Thật vậy, trong chứng minh (c$'$) kéo theo (b), $\sh{J}$ và $\sh{J}'$ khi đó là các bó iđêan \emph{coherent}, và mặt khác, mọi môđun con quasi-coherent của một môđun coherent đều coherent \sref[I]{I.6.1.1};
từ đó suy ra kết luận.
\end{remark}

\begin{corollary}[5.2.2]
\label{II.5.2.2}
Cho $f:X\to Y$ là một cấu xạ tách quasi-compact.
Các điều kiện sau tương đương:
\begin{enumerate}
  \item[{\rm(a)}] $f$ là một cấu xạ affine.
  \item[{\rm(b)}] Hàm tử $f_*$ khớp trong phạm trù các $\sh{O}_X$-môđun quasi-coherent.
  \item[{\rm(c)}] Với mọi $\sh{O}_X$-môđun quasi-coherent $\sh{F}$, ta có $\RR^1f_*(\sh{F})=0$.
  \item[{\rm(c')}] Với mọi bó iđêan quasi-coherent $\sh{J}$ của $\sh{O}_X$, ta có $\RR^1f_*(\sh{J})=0$.
\end{enumerate}
\end{corollary}

\begin{proof}
Vì tất cả các điều kiện này đều có tính địa phương trên $Y$ theo định nghĩa của hàm tử $\RR^1f_*$ (T,~3.7.3), nên ta có thể giả sử $Y$ affine.
Nếu $f$ affine thì $X$ affine, và tính chất~(b) chính là \sref[I]{I.1.6.4}.
Ngược lại, ta chứng minh rằng (b) kéo theo (a):
với mọi $\sh{O}_X$-môđun quasi-coherent $\sh{F}$, $f_*(\sh{F})$ là một $\sh{O}_Y$-môđun quasi-coherent \sref[I]{I.9.2.2}[(a)].
Theo giả thiết, hàm tử $f_*$ khớp theo $\sh{F}$, và hàm tử $\Gamma(Y,\sh{G})$ khớp theo $\sh{G}$ (trong phạm trù các $\sh{O}_Y$-môđun quasi-coherent), vì $Y$ affine \sref[I]{I.1.3.11};
do đó $\Gamma(Y,f_*(\sh{F}))=\Gamma(X,\sh{F})$ khớp theo $\sh{F}$, điều này chứng minh khẳng định nhờ \sref{II.5.2.1}[(c)].

Nếu $f$ affine, thì $f^{-1}(U)$ affine với mọi tập mở affine $U$ của $Y$ \sref{II.1.3.2}; do đó $\HH^1(f^{-1}(U),\sh{F})=0$ \sref{II.5.2.1}[(d)], và theo định nghĩa điều này kéo theo $\RR^1f_*(\sh{F})=0$.
Cuối cùng, giả sử điều kiện~(c$'$) được thỏa mãn;
dãy khớp gồm các hạng tử bậc thấp của dãy phổ Leray (G,~II,~4.17.1 và I,~4.5.1) đặc biệt cho ta dãy khớp
\[
  0\to\HH^1(Y,f_*(\sh{J}))\to\HH^1(X,\sh{J})\to\HH^0(Y,\RR^1f_*(\sh{J})).
\]
Vì $Y$ affine và $f_*(\sh{J})$ quasi-coherent \sref[I]{I.9.2.2}[(a)], ta có $\HH^1(Y,f_*(\sh{J}))=0$ \sref{II.5.2.1};
do đó giả thiết~(c$'$) kéo theo $\HH^1(X,\sh{J})=0$, và từ \sref{II.5.2.1} ta kết luận rằng $X$ là một lược đồ affine.
\end{proof}

\begin{corollary}[5.2.3]
\label{II.5.2.3}
Nếu $f:X\to Y$ là một cấu xạ affine, thì với mọi $\sh{O}_X$-môđun quasi-coherent $\sh{F}$, đồng cấu chính tắc $\HH^1(Y,f_*(\sh{F}))\to\HH^1(X,\sh{F})$ là song ánh.
\end{corollary}
\oldpage[II]{99}
\begin{proof}
Ta có dãy khớp
\[
  0\to\HH^1(Y,f_*(\sh{F}))\to\HH^1(X,\sh{F})\to\HH^0(Y,\RR^1f_*(\sh{F}))
\]
gồm các hạng tử bậc thấp của dãy phổ Leray, và kết luận suy ra từ \sref{II.5.2.2}.
\end{proof}

\begin{remark}[5.2.4]
\label{II.5.2.4}
Trong Chương~III,~\textsection1, ta sẽ chứng minh rằng nếu $X$ affine thì $\HH^i(X,\sh{F})=0$ với mọi $i>0$ và mọi $\sh{O}_X$-môđun quasi-coherent $\sh{F}$.
\end{remark}

\subsection{Các cấu xạ quasi-xạ ảnh}
\label{subsection:II.5.3}

\begin{definition}[5.3.1]
\label{II.5.3.1}
Ta nói một cấu xạ $f:X\to Y$ là \emph{quasi-xạ ảnh}, hoặc $X$ (được xét như một $Y$-tiền lược đồ qua $f$) \emph{quasi-xạ ảnh trên $Y$}, hoặc $X$ là một \emph{$Y$-lược đồ quasi-xạ ảnh}, nếu $f$ thuộc kiểu hữu hạn và tồn tại một $\sh{O}_X$-môđun khả nghịch $f$-ample.
\end{definition}

Ta lưu ý rằng khái niệm này \emph{không có tính địa phương trên $Y$}:
các phản ví dụ của Nagata~\cite{II-26} và Hironaka cho thấy rằng, ngay cả khi $X$ và $Y$ là các lược đồ đại số không kỳ dị trên một trường đóng đại số, mỗi điểm của $Y$ có thể có một lân cận affine $U$ sao cho $f^{-1}(U)$ quasi-xạ ảnh trên $U$, mà $f$ không quasi-xạ ảnh.

Ta lưu ý rằng một cấu xạ quasi-xạ ảnh nhất thiết là \emph{tách} \sref{II.4.6.1}.
Khi $Y$ quasi-compact, nói $f$ quasi-xạ ảnh tương đương với nói $f$ thuộc kiểu hữu hạn và tồn tại một $\sh{O}_X$-môđun \emph{rất ample} đối với $f$ (\sref{II.4.6.2} và \sref{II.4.6.11}).
Hơn nữa:

\begin{proposition}[5.3.2]
\label{II.5.3.2}
Cho $Y$ là một lược đồ quasi-compact hoặc một tiền lược đồ có không gian nền Noether, và cho $X$ là một $Y$-tiền lược đồ.
Các điều kiện sau tương đương:
\begin{enumerate}
  \item[{\rm(a)}] $X$ là một $Y$-lược đồ quasi-xạ ảnh.
  \item[{\rm(b)}] $X$ thuộc kiểu hữu hạn trên $Y$, và tồn tại một $\sh{O}_Y$-môđun quasi-coherent hữu hạn sinh $\sh{E}$ sao cho $X$ $Y$-đẳng cấu với một tiền lược đồ con của $\bb{P}(\sh{E})$.
  \item[{\rm(c)}] $X$ thuộc kiểu hữu hạn trên $Y$, và tồn tại một $\sh{O}_Y$-đại số phân bậc quasi-coherent $\sh{S}$ sao cho $\sh{S}_1$ hữu hạn sinh và sinh $\sh{S}$, và sao cho $X$ $Y$-đẳng cấu với một tiền lược đồ con cảm sinh trên một tập mở trù mật khắp nơi của $\Proj(\sh{S})$.
\end{enumerate}
\end{proposition}

\begin{proof}
Điều này suy ra ngay từ nhận xét trên và từ \sref{II.4.4.3}, \sref{II.4.4.6} cùng \sref{II.4.4.7}.
\end{proof}

Ta lưu ý rằng, khi $Y$ là một tiền lược đồ \emph{Noether}, trong các điều kiện~(b) và (c) của \sref{II.5.3.2}, có thể bỏ giả thiết $X$ thuộc kiểu hữu hạn trên $Y$, vì khi đó giả thiết ấy tự động được thỏa mãn \sref[I]{I.6.3.5}.

\begin{corollary}[5.3.3]
\label{II.5.3.3}
Cho $Y$ là một lược đồ quasi-compact có một $\sh{O}_Y$-môđun ample $\sh{L}$ \sref{II.4.5.3}.
Điều kiện cần và đủ để một $Y$-lược đồ $X$ là quasi-xạ ảnh là nó thuộc kiểu hữu hạn trên $Y$ và đẳng cấu với một $Y$-lược đồ con của một phân thớ xạ ảnh dạng $\bb{P}_Y^r$.
\end{corollary}

\begin{proof}
Thật vậy, nếu $\sh{E}$ là một $\sh{O}_Y$-môđun quasi-coherent hữu hạn sinh, thì $\sh{E}$ đẳng cấu với một môđun thương của $\sh{O}_Y$-môđun $\sh{L}^{\otimes(-n)}\otimes_{\sh{O}_Y}\sh{O}_Y^k$ \sref{II.4.5.5}; do đó $\bb{P}(\sh{E})$ đẳng cấu với một lược đồ con đóng của $\bb{P}_Y^{k-1}$ (\sref{II.4.1.2} và \sref{II.4.1.4}).
\end{proof}

\begin{proposition}[5.3.4]
\label{II.5.3.4}
\medskip\noindent
\begin{enumerate}
  \item[{\rm(i)}] Mọi cấu xạ quasi-affine thuộc kiểu hữu hạn (và đặc biệt, mọi phép nhúng quasi-compact hoặc mọi cấu xạ affine thuộc kiểu hữu hạn) đều quasi-xạ ảnh.
  \item[{\rm(ii)}] Nếu $f:X\to Y$ và $g:Y\to Z$ là quasi-xạ ảnh, và nếu $Z$ quasi-compact, thì $g\circ f$ quasi-xạ ảnh.
\oldpage[II]{100}
  \item[{\rm(iii)}] Nếu $f:X\to Y$ là một $S$-cấu xạ quasi-xạ ảnh, thì $f_{(S')}:X_{(S')}\to Y_{(S')}$ là quasi-xạ ảnh đối với mọi phép mở rộng $S'\to S$ của tiền lược đồ cơ sở.
  \item[{\rm(iv)}] Nếu $f:X\to Y$ và $g:X'\to Y'$ là hai $S$-cấu xạ quasi-xạ ảnh, thì $f\times_S g$ là quasi-xạ ảnh.
  \item[{\rm(v)}] Nếu $f:X\to Y$ và $g:Y\to Z$ là hai cấu xạ sao cho $g\circ f$ quasi-xạ ảnh, và nếu $g$ tách hoặc $X$ là Noether địa phương, thì $f$ quasi-xạ ảnh.
  \item[{\rm(vi)}] Nếu $f$ là một cấu xạ quasi-xạ ảnh, thì $f_\red$ cũng vậy.
\end{enumerate}
\end{proposition}

\begin{proof}
(i) suy ra từ \sref{II.5.1.6} và \sref{II.5.1.10}[(i)].
Các khẳng định khác là những hệ quả trực tiếp của Định nghĩa~\sref{II.5.3.1}, các tính chất của cấu xạ kiểu hữu hạn \sref[I]{I.6.3.4} và \sref{II.4.6.13}.
\end{proof}

\begin{remark}[5.3.5]
\label{II.5.3.5}
Ta lưu ý rằng có thể xảy ra $f_\red$ quasi-xạ ảnh mà $f$ không quasi-xạ ảnh, ngay cả khi giả sử $Y$ là phổ của một đại số có hạng hữu hạn trên $\bb{C}$ và $f$ là riêng.
\end{remark}

\begin{corollary}[5.3.6]
\label{II.5.3.6}
Nếu $X$ và $X'$ là hai $Y$-lược đồ quasi-xạ ảnh, thì $X\sqcup X'$ là một $Y$-lược đồ quasi-xạ ảnh.
\end{corollary}

\begin{proof}
Điều này suy ra từ \sref{II.4.6.18}.
\end{proof}

\subsection{Các cấu xạ riêng và các cấu xạ đóng phổ dụng}
\label{subsection:II.5.4}

\begin{definition}[5.4.1]
\label{II.5.4.1}
Ta nói một cấu xạ các tiền lược đồ $f:X\to Y$ là \emph{riêng} nếu nó thỏa mãn hai điều kiện sau:
\begin{enumerate}
  \item[(a)] $f$ tách và thuộc kiểu hữu hạn;
  \item[(b)] đối với mọi tiền lược đồ $Y'$ và mọi cấu xạ $Y'\to Y$, phép chiếu $f_{(Y')}:X\times_Y Y'\to Y'$ là một cấu xạ đóng \sref[I]{I.2.2.6}.
\end{enumerate}

Khi đó, ta cũng nói $X$ (được xét như một $Y$-tiền lược đồ có cấu xạ cấu trúc $f$) là riêng trên $Y$.
\end{definition}

Hiển nhiên các điều kiện~(a) và (b) có tính địa phương trên $Y$.
Để kiểm tra ảnh của một tập con đóng $Z$ của $X\times_Y Y'$ qua phép chiếu $q:X\times_Y Y'\to Y'$ là đóng trong $Y'$, chỉ cần thấy $q(Z)\cap U'$ đóng trong $U'$ đối với mọi tập mở affine $U'$ của $Y'$;
vì $q(Z)\cap U'=q(Z\cap q^{-1}(U'))$ và $q^{-1}(U')$ được đồng nhất với $X\times_Y U'$ \sref[I]{I.4.4.1}, ta thấy rằng, để kiểm tra điều kiện~(b) của Định nghĩa~\sref{II.5.4.1}, có thể chỉ xét trường hợp $Y'$ là một lược đồ \emph{affine}.
Sau này ta sẽ thấy \sref{II.5.6.3} rằng, nếu $Y$ là Noether địa phương, thì thậm chí có thể chỉ cần kiểm tra (b) khi $Y'$ thuộc kiểu hữu hạn trên $Y$.

Hiển nhiên mọi cấu xạ riêng đều \emph{đóng}.

\begin{proposition}[5.4.2]
\label{II.5.4.2}
\medskip\noindent
\begin{enumerate}
  \item[{\rm(i)}] Một phép nhúng đóng là một cấu xạ riêng.
  \item[{\rm(ii)}] Hợp thành của hai cấu xạ riêng là riêng.
  \item[{\rm(iii)}] Nếu $X$ và $Y$ là các $S$-tiền lược đồ và $f:X\to Y$ là một $S$-cấu xạ riêng, thì $f_{(S')}:X_{(S')}\to Y_{(S')}$ là riêng đối với mọi phép mở rộng $S'\to S$ của tiền lược đồ cơ sở.
  \item[{\rm(iv)}] Nếu $f:X\to Y$ và $g:X'\to Y'$ là hai $S$-cấu xạ riêng, thì $f\times_S g:X\times_S X'\to Y\times_S Y'$ là một $S$-cấu xạ riêng.
\end{enumerate}
\end{proposition}
\oldpage[II]{101}
\begin{proof}
Chỉ cần chứng minh (i), (ii) và (iii) \sref[I]{I.3.5.1}.
Trong mỗi trường hợp ấy, việc kiểm tra điều kiện~(a) của \sref{II.5.4.1} suy ra từ các kết quả trước (\sref[I]{I.5.5.1} và \sref[I]{I.6.3.4}); còn phải kiểm tra điều kiện~(b).
Điều này là trực tiếp trong trường hợp (i), vì nếu $X\to Y$ là một phép nhúng đóng thì $X\times_Y Y'\to Y\times_Y Y'=Y'$ cũng vậy (\sref[I]{I.4.3.2} và \sref[I]{I.3.3.3}).
Để chứng minh (ii), xét hai cấu xạ riêng $X\to Y$, $Y\to Z$, và một cấu xạ $Z'\to Z$.
Ta có thể viết $X\times_Z Z'=X\times_Y(Y\times_Z Z')$ \sref[I]{I.3.3.9.1}; do đó phép chiếu $X\times_Z Z'\to Z'$ phân tích thành $X\times_Y(Y\times_Z Z')\to Y\times_Z Z'\to Z'$.
Theo nhận xét mở đầu, (ii) suy ra từ việc hợp thành của hai cấu xạ đóng vẫn là một cấu xạ đóng.
Cuối cùng, đối với mọi cấu xạ $S'\to S$, $X_{(S')}$ được đồng nhất với $X\times_Y Y_{(S')}$ \sref[I]{I.3.3.11}; đối với mọi cấu xạ $Z\to Y_{(S')}$, ta có thể viết
\[
  X_{(S')}\times_{Y_{(S')}}Z=(X\times_Y Y_{(S')})\times_{Y_{(S')}}Z=X\times_Y Z;
\]
vì theo giả thiết $X\times_Y Z\to Z$ là đóng, điều này chứng minh (iii).
\end{proof}

\begin{corollary}[5.4.3]
\label{II.5.4.3}
Cho $f:X\to Y$, $g:Y\to Z$ là hai cấu xạ sao cho $g\circ f$ là riêng.
\begin{enumerate}
  \item[{\rm(i)}] Nếu $g$ tách, thì $f$ là riêng.
  \item[{\rm(ii)}] Nếu $g$ tách và thuộc kiểu hữu hạn, và nếu $f$ toàn ánh, thì $g$ là riêng.
\end{enumerate}
\end{corollary}

\begin{proof}
(i) suy ra từ \sref{II.5.4.2} bằng phương pháp tổng quát \sref[I]{I.5.5.12}.
Để chứng minh (ii), chỉ cần kiểm tra điều kiện~(b) của Định nghĩa~\sref{II.5.4.1}.
Đối với mọi cấu xạ $Z'\to Z$, biểu đồ
\[
  \xymatrix{
    X\times_Z Z'\ar[r]^{f\times1_{Z'}}\ar[dr]_p &
    Y\times_Z Z'\ar[d]^{p'}\\
    & Z'
  }
\]
(trong đó $p$ và $p'$ là các phép chiếu) giao hoán \sref[I]{I.3.2.1};
hơn nữa, $f\times1_{Z'}$ toàn ánh vì $f$ toàn ánh \sref[I]{I.3.5.2}, và $p$ là một cấu xạ đóng theo giả thiết.
Do đó, mỗi tập con đóng $F$ của $Y\times_Z Z'$ là ảnh qua $f\times1_{Z'}$ của một tập con đóng $E$ của $X\times_Z Z'$; vì vậy $p'(F)=p(E)$ đóng trong $Z'$ theo giả thiết, từ đó suy ra hệ quả.
\end{proof}

\begin{corollary}[5.4.4]
\label{II.5.4.4}
Nếu $X$ là một tiền lược đồ riêng trên $Y$, và $\sh{S}$ là một $\sh{O}_Y$-đại số quasi-coherent, thì mọi $Y$-cấu xạ $f:X\to\Proj(\sh{S})$ đều riêng (và, \emph{a fortiori}, đóng).
\end{corollary}

\begin{proof}
Thật vậy, cấu xạ cấu trúc $p:\Proj(\sh{S})\to Y$ là tách, và $p\circ f$ là riêng theo giả thiết.
\end{proof}

\begin{corollary}[5.4.5]
\label{II.5.4.5}
Cho $f:X\to Y$ là một cấu xạ tách thuộc kiểu hữu hạn.
Cho $(X_i)_{1\leq i\leq n}$ (tương ứng, $(Y_i)_{1\leq i\leq n}$) là một họ hữu hạn các tiền lược đồ con đóng của $X$ (tương ứng, của $Y$), và $j_i$ (tương ứng, $h_i$) là phép nhúng chính tắc $X_i\to X$ (tương ứng, $Y_i\to Y$).
Giả sử không gian nền của $X$ là hợp của các $X_i$, và đối với mỗi $i$, tồn tại một cấu xạ $f_i:X_i\to Y_i$ sao cho biểu đồ
\[
  \xymatrix{
    X_i\ar[r]^{f_i}\ar[d]_{j_i} &
    Y_i\ar[d]^{h_i}\\
    X\ar[r]_f &
    Y
  }
\]
giao hoán.
Khi đó, điều kiện cần và đủ để $f$ là riêng là mỗi $f_i$ đều riêng.
\end{corollary}
\oldpage[II]{102}
\begin{proof}
Nếu $f$ là riêng, thì $f\circ j_i$ cũng riêng, vì $j_i$ là một phép nhúng đóng \sref{II.5.4.2};
vì $h_i$ là một phép nhúng đóng, do đó là một cấu xạ tách, $f_i$ là riêng theo \sref{II.5.4.3}.
Ngược lại, giả sử mỗi $f_i$ đều riêng, và xét tiền lược đồ $Z$ bằng tổng của các $X_i$; cho $u$ là cấu xạ $Z\to X$ có hạn chế bằng $j_i$ trên mỗi $X_i$.
Hạn chế của $f\circ u$ trên mỗi $X_i$ bằng $f\circ j_i=h_i\circ f_i$, do đó là riêng, vì cả $h_i$ lẫn $f_i$ đều riêng \sref{II.5.4.2};
khi ấy, từ Định nghĩa~\sref{II.5.4.1} suy ra ngay $f\circ u$ là riêng.
Nhưng vì $u$ toàn ánh theo giả thiết, ta kết luận $f$ là riêng theo \sref{II.5.4.3}.
\end{proof}

\begin{corollary}[5.4.6]
\label{II.5.4.6}
Cho $f:X\to Y$ là một cấu xạ tách thuộc kiểu hữu hạn; điều kiện cần và đủ để $f$ là riêng là $f_\red:X_\red\to Y_\red$ riêng.
\end{corollary}

\begin{proof}
Đây là trường hợp riêng của \sref{II.5.4.5}, với $n=1$, $X_1=X_\red$ và $Y_1=Y_\red$ \sref[I]{I.5.1.5}.
\end{proof}

\begin{env}[5.4.7]
\label{II.5.4.7}
Nếu $X$ và $Y$ là các tiền lược đồ Noether và $f:X\to Y$ là một cấu xạ tách thuộc kiểu hữu hạn, thì để kiểm tra $f$ là riêng, ta có thể quy về trường hợp các cấu xạ \emph{trội} giữa các tiền lược đồ \emph{nguyên}.
Thật vậy, cho $X_i$ ($1\leq i\leq n$) là các thành phần bất khả quy của $X$ (chỉ có hữu hạn thành phần), và đối với mỗi $i$, xét tiền lược đồ con đóng rút gọn duy nhất của $X$ có $X_i$ làm không gian nền, mà ta vẫn ký hiệu là $X_i$ \sref[I]{I.5.2.1}.
Mặt khác, cho $Y_i$ là tiền lược đồ con đóng rút gọn duy nhất của $Y$ có $\overline{f(X_i)}$ làm không gian nền.
Nếu $g_i$ (tương ứng, $h_i$) là cấu xạ nhúng $X_i\to X$ (tương ứng, $Y_i\to Y$), ta suy ra $f\circ g_i=h_i\circ f_i$, trong đó $f_i$ là một cấu xạ trội $X_i\to Y_i$ \sref[I]{I.5.2.2};
khi ấy các điều kiện để áp dụng \sref{II.5.4.5} được thỏa mãn, và điều kiện cần và đủ để $f$ là riêng là mọi $f_i$ đều riêng.
\end{env}

\begin{corollary}[5.4.8]
\label{II.5.4.8}
Cho $X$ và $Y$ là hai $S$-tiền lược đồ tách thuộc kiểu hữu hạn trên $S$, và cho $f:X\to Y$ là một $S$-cấu xạ.
Điều kiện cần và đủ để $f$ là riêng là, đối với mọi $S$-tiền lược đồ $S'$, cấu xạ $f\times_S 1_{S'}:X\times_S S'\to Y\times_S S'$ đóng.
\end{corollary}

\begin{proof}
Trước hết lưu ý rằng, nếu $g:X\to S$ và $h:Y\to S$ là các cấu xạ cấu trúc, thì theo định nghĩa $g=h\circ f$; do đó $f$ tách và thuộc kiểu hữu hạn (\sref[I]{I.5.5.1} và \sref[I]{I.6.3.4}).
Nếu $f$ là riêng, thì $f\times_S 1_{S'}$ cũng riêng \sref{II.5.4.2}; \emph{a fortiori}, $f\times_S 1_{S'}$ đóng.
Ngược lại, giả sử điều kiện trong phát biểu được thỏa mãn và cho $Y'$ là một $Y$-tiền lược đồ;
$Y'$ cũng có thể được xét như một $S$-tiền lược đồ, và vì $Y\to S$ tách, $X\times_Y Y'$ được đồng nhất với một tiền lược đồ con đóng của $X\times_S Y'$ \sref[I]{I.5.4.2}.
Trong biểu đồ giao hoán
\[
  \xymatrix{
    X\times_Y Y'\ar[r]^{f\times1_{Y'}}\ar[d] &
    Y\times_Y Y'=Y'\ar[d]\\
    X\times_S Y'\ar[r]_{f\times_S1_{Y'}} &
    Y\times_S Y',
  }
\]
các mũi tên dọc là những phép nhúng đóng; từ đó suy ra ngay rằng, nếu $f\times_S1_{Y'}$ là một cấu xạ đóng, thì $f\times_Y1_{Y'}$ cũng vậy.
\end{proof}

\begin{remark}[5.4.9]
\label{II.5.4.9}
Ta nói một cấu xạ $f:X\to Y$ là \emph{đóng phổ dụng} nếu nó thỏa mãn điều kiện~(b) của Định nghĩa~\sref{II.5.4.1}.
Người đọc sẽ nhận thấy rằng,
\oldpage[II]{103}
trong các mệnh đề \sref{II.5.4.2} đến \sref{II.5.4.8}, ở mọi nơi ta có thể thay ``riêng'' bằng ``đóng phổ dụng'' mà không làm thay đổi tính đúng đắn của các kết quả (và trong các giả thiết của \sref{II.5.4.3}, \sref{II.5.4.5}, \sref{II.5.4.6} và \sref{II.5.4.8}, ta có thể bỏ các điều kiện hữu hạn).
\end{remark}

\begin{env}[5.4.10]
\label{II.5.4.10}
Cho $f:X\to Y$ là một cấu xạ kiểu hữu hạn.
Ta nói một tập con đóng $Z$ của $X$ là \emph{riêng trên $Y$} (hay \emph{$Y$-riêng}, hay \emph{riêng đối với $f$}) nếu hạn chế của $f$ trên một tiền lược đồ con đóng của $X$ có không gian nền là $Z$ \sref[I]{I.5.2.1} là riêng.
Vì hạn chế này khi đó là tách, từ \sref{II.5.4.6} và \sref[I]{I.5.5.1}[(vi)] suy ra tính chất trên \emph{không phụ thuộc} vào tiền lược đồ con đóng của $X$ có $Z$ làm không gian nền.
Nếu $g:X'\to X$ là một cấu xạ \emph{riêng}, thì $g^{-1}(Z)$ là một tập con \emph{riêng} của $X'$:
nếu $T$ là một tiền lược đồ con của $X$ có $Z$ làm không gian nền, chỉ cần lưu ý rằng hạn chế của $g$ trên tiền lược đồ con đóng $g^{-1}(T)$ của $X'$ là một cấu xạ riêng $g^{-1}(T)\to T$ theo \sref{II.5.4.2}[(iii)], rồi áp dụng \sref{II.5.4.2}[(ii)].
Mặt khác, nếu $X''$ là một $Y$-\emph{lược đồ} kiểu hữu hạn và $u:X\to X''$ là một $Y$-cấu xạ, thì $u(Z)$ là một tập con \emph{riêng} của $X''$;
thật vậy, lấy $T$ là tiền lược đồ con đóng rút gọn của $X$ có $Z$ làm không gian nền;
vì hạn chế của $f$ trên $T$ là riêng, hạn chế của $u$ trên $T$ cũng riêng \sref{II.5.4.3}[(i)], do đó $u(Z)$ đóng trong $X''$;
cho $T''$ là một tiền lược đồ con đóng của $X''$ có $u(Z)$ làm không gian nền \sref[I]{I.5.2.1}, sao cho $u|T$ phân tích thành $T\xrightarrow{v}T''\xrightarrow{j}X''$, trong đó $j$ là phép nhúng chính tắc \sref[I]{I.5.2.2}; khi đó $v$ là riêng theo \sref{II.5.4.3}[(i)] và toàn ánh do cách chọn $T''$;
nếu $g$ là hạn chế trên $T''$ của cấu xạ cấu trúc $X''\to Y$, thì $g$ tách và thuộc kiểu hữu hạn, và ta có $f|T=g\circ v$;
do đó từ \sref{II.5.4.3}[(ii)] suy ra $g$ là riêng, từ đó có khẳng định của ta.
\end{env}

Đặc biệt, từ các nhận xét này suy ra rằng, nếu $Z$ là một tập con $Y$-riêng của $X$, thì:
\begin{enumerate}
  \item đối với mọi tiền lược đồ con đóng $X'$ của $X$, $Z\cap X'$ là một tập con $Y$-riêng của $X'$; và
  \item nếu $X$ là một tiền lược đồ con của một $Y$-lược đồ kiểu hữu hạn $X''$, thì $Z$ cũng là một tập con $Y$-riêng của $X''$ (và đặc biệt, \emph{đóng trong $X''$}).
\end{enumerate}

\subsection{Các cấu xạ xạ ảnh}
\label{subsection:II.5.5}

\begin{proposition}[5.5.1]
\label{II.5.5.1}
Cho $X$ là một $Y$-tiền lược đồ.
Các điều kiện sau là tương đương:
\begin{enumerate}
  \item[{\rm(a)}] $X$ $Y$-đẳng cấu với một tiền lược đồ con \emph{đóng} của một phân thớ xạ ảnh $\bb{P}(\sh{E})$, trong đó $\sh{E}$ là một $\sh{O}_Y$-môđun quasi-coherent hữu hạn sinh.
  \item[{\rm(b)}] Tồn tại một $\sh{O}_Y$-đại số phân bậc quasi-coherent $\sh{S}$ sao cho $\sh{S}_1$ hữu hạn sinh và sinh ra $\sh{S}$, đồng thời $X$ $Y$-đẳng cấu với $\Proj(\sh{S})$.
\end{enumerate}
\end{proposition}

\begin{proof}
Điều kiện~(a) kéo theo (b) theo \sref{II.3.6.2}[(ii)]: nếu $\sh{J}$ là một bó phân bậc quasi-coherent các iđêan của $\bb{S}(\sh{E})$, thì $\sh{O}_Y$-đại số phân bậc quasi-coherent $\sh{S}=\bb{S}(\sh{E})/\sh{J}$ được sinh bởi $\sh{S}_1$, và $\sh{S}_1$, ảnh chính tắc của $\sh{E}$, là một $\sh{O}_Y$-môđun hữu hạn sinh.
Điều kiện~(b) kéo theo (a) theo \sref{II.3.6.2}, áp dụng cho trường hợp ánh xạ $\sh{M}\to\sh{S}_1$ là ánh xạ đồng nhất.
\end{proof}
\oldpage[II]{104}
\begin{definition}[5.5.2]
\label{II.5.5.2}
Ta nói một $Y$-tiền lược đồ $X$ là \emph{xạ ảnh} trên $Y$, hay là một $Y$-lược đồ xạ ảnh, nếu nó thỏa mãn các điều kiện tương đương~(a) và (b) của \sref{II.5.5.1}.
Ta nói một cấu xạ $f:X\to Y$ là xạ ảnh nếu nó làm cho $X$ thành một $Y$-lược đồ xạ ảnh.
\end{definition}

Hiển nhiên nếu $f:X\to Y$ là xạ ảnh, thì tồn tại một $\sh{O}_X$-môđun \emph{rất ample} đối với $f$ \sref{II.4.4.2}.

\begin{theorem}[5.5.3]
\label{II.5.5.3}
\medskip\noindent
\begin{enumerate}
  \item[{\rm(i)}] Mọi cấu xạ xạ ảnh đều quasi-xạ ảnh và riêng.
  \item[{\rm(ii)}] Ngược lại, cho $Y$ là một lược đồ quasi-compact hoặc một tiền lược đồ có không gian nền Noether; khi đó mọi cấu xạ $f:X\to Y$ vừa quasi-xạ ảnh vừa riêng đều xạ ảnh.
\end{enumerate}
\end{theorem}

\begin{proof}
\medskip\noindent
\begin{enumerate}
  \item[(i)] Hiển nhiên nếu $f:X\to Y$ là xạ ảnh, thì nó thuộc kiểu hữu hạn và quasi-xạ ảnh (do đó, đặc biệt, tách); mặt khác, từ \sref{II.5.5.1}[(b)] và \sref{II.3.5.3} suy ra ngay rằng, nếu $f$ là xạ ảnh, thì $f\times_Y 1_{Y'}:X\times_Y Y'\to Y'$ cũng xạ ảnh đối với mọi cấu xạ $Y'\to Y$.
    Vì vậy, để chứng minh $f$ đóng phổ dụng, chỉ cần chứng minh một cấu xạ xạ ảnh $f$ là \emph{đóng}.
    Vì vấn đề có tính địa phương trên $Y$, ta có thể giả sử $Y=\Spec(A)$, và do \sref{II.5.5.1}, $X=\Proj(S)$, trong đó $S$ là một $A$-đại số phân bậc được sinh bởi một số hữu hạn các phần tử của $S_1$.
    Đối với mọi $y\in Y$, thớ $f^{-1}(y)$ được đồng nhất với $\Proj(S)\times_Y\Spec(\kres(y))$ \sref[I]{I.3.6.1}, do đó với $\Proj(S\otimes_A\kres(y))$ \sref{II.2.8.10};
    vì vậy $f^{-1}(y)$ rỗng khi và chỉ khi $S\otimes_A\kres(y)$ thỏa mãn điều kiện~(TN) \sref{II.2.7.4}, nói cách khác, khi và chỉ khi $S_n\otimes_A\kres(y)=0$ với $n$ đủ lớn.
    Nhưng vì $(S_n)_y$ là một $\sh{O}_y$-môđun hữu hạn sinh, điều kiện trên có nghĩa là $(S_n)_y=0$ với $n$ đủ lớn, theo bổ đề Nakayama.
    Nếu $\mathfrak{a}_n$ là linh hóa tử trong $A$ của $A$-môđun $S_n$, điều kiện trên còn có nghĩa là $\mathfrak{a}_n\not\subset\mathfrak{j}_y$ với $n$ đủ lớn \sref[0]{0.1.7.4}.
    Mà theo giả thiết $S_nS_1=S_{n+1}$, nên $\mathfrak{a}_n\subset\mathfrak{a}_{n+1}$; nếu $\mathfrak{a}$ là hợp của các $\mathfrak{a}_n$, ta thấy $f(X)=V(\mathfrak{a})$, điều này chứng minh $f(X)$ đóng trong $Y$.
    Bây giờ, nếu $X'$ là một tập con đóng tùy ý của $X$, thì tồn tại một tiền lược đồ con đóng của $X$ có $X'$ làm không gian nền \sref[I]{I.5.2.1}, và theo \sref{II.5.5.1}[(a)], hiển nhiên cấu xạ $X'\to X\xrightarrow{f}Y$ là xạ ảnh; do đó $f(X')$ đóng trong $Y$.
  \item[(ii)] Giả thiết trên $Y$ và sự kiện $f$ quasi-xạ ảnh kéo theo sự tồn tại của một $\sh{O}_Y$-môđun quasi-coherent hữu hạn sinh $\sh{E}$ và một $Y$-phép nhúng $j:X\to\bb{P}(\sh{E})$ \sref{II.5.3.2}.
    Nhưng vì $f$ là riêng, $j$ là \emph{đóng} theo \sref{II.5.4.4}; do đó $f$ là xạ ảnh.
\end{enumerate}
\end{proof}

\begin{remark}[5.5.4]
\label{II.5.5.4}
\medskip\noindent
\begin{enumerate}
  \item[(i)] Cho $f:X\to Y$ là một cấu xạ sao cho: $1^\circ$ $f$ là riêng; $2^\circ$ tồn tại một $\sh{O}_X$-môđun $\sh{L}$ \emph{rất ample} đối với $f$; $3^\circ$ $\sh{O}_Y$-môđun quasi-coherent $\sh{E}=f_*(\sh{L})$ \emph{hữu hạn sinh}.
    Khi đó $f$ là một cấu xạ \emph{xạ ảnh}: thật vậy, theo \sref{II.4.4.4}, tồn tại một $Y$-phép nhúng $r:X\to\bb{P}(\sh{E})$, và vì $f$ là riêng, $r$ là một phép nhúng \emph{đóng} \sref{II.5.4.4}.
    Trong Chương~III, \textsection3, ta sẽ thấy rằng, khi $Y$ là \emph{Noether địa phương}, điều kiện thứ ba trên ($\sh{E}$ hữu hạn sinh) là hệ quả của hai điều kiện kia; do đó trong trường hợp này hai điều kiện đầu \emph{đặc trưng} các cấu xạ xạ ảnh, và nếu $Y$ quasi-compact thì có thể thay điều kiện thứ hai (sự tồn tại của một $\sh{O}_X$-môđun $\sh{L}$ rất ample đối với $f$) bằng giả thiết tồn tại một $\sh{O}_X$-môđun \emph{ample} đối với $f$ \sref{II.4.6.11}.
  \item[(ii)] Cho $Y$ là một lược đồ quasi-compact sao cho tồn tại một $\sh{O}_Y$-môđun ample.
    Điều kiện cần và đủ để một $Y$-lược đồ $X$ là \emph{xạ ảnh} là nó $Y$-đẳng cấu với một $Y$-lược đồ con \emph{đóng} của một phân thớ xạ ảnh dạng $\bb{P}_Y^r$.
    Điều kiện rõ ràng là đủ.
\oldpage[II]{105}
Ngược lại, nếu $X$ xạ ảnh trên $Y$, thì $X$ quasi-xạ ảnh; do đó tồn tại một $Y$-phép nhúng $j$ của $X$ vào một $\bb{P}_Y^r$ nào đó \sref{II.5.3.3}, và phép nhúng này là \emph{đóng} theo \sref{II.5.4.4} và \sref{II.5.5.3}.
  \item[{\rm(iii)}] Lập luận của \sref{II.5.5.3} cho thấy rằng, với mọi tiền lược đồ $Y$ và mọi số nguyên $r\geq0$, cấu xạ cấu trúc $\bb{P}_Y^r\to Y$ là \emph{toàn ánh}; thật vậy, nếu đặt $\sh{S}=\bb{S}_{\sh{O}_Y}(\sh{O}_Y^{r+1})$, thì hiển nhiên $\sh{S}_y=\bb{S}_{\kres(y)}(\kres(y)^{r+1})$ \sref{II.1.7.3}, và do đó $(\sh{S}_n)_y\neq0$ với mọi $y\in Y$ và mọi $n\geq0$.
  \item[{\rm(iv)}] Từ các ví dụ của Nagata~\cite{II-26} suy ra rằng có những cấu xạ riêng không quasi-xạ ảnh.
\end{enumerate}
\end{remark}

\begin{proposition}[5.5.5]
\label{II.5.5.5}
\medskip\noindent
\begin{enumerate}
  \item[{\rm(i)}] Mọi phép nhúng đóng đều là một cấu xạ xạ ảnh.
  \item[{\rm(ii)}] Nếu $f:X\to Y$ và $g:Y\to Z$ là các cấu xạ xạ ảnh, và nếu $Z$ là một lược đồ quasi-compact hoặc một tiền lược đồ có không gian nền Noether, thì $g\circ f$ là xạ ảnh.
  \item[{\rm(iii)}] Nếu $f:X\to Y$ là một $S$-cấu xạ xạ ảnh, thì $f_{(S')}:X_{(S')}\to Y_{(S')}$ là xạ ảnh với mọi phép mở rộng $S'\to S$ của tiền lược đồ cơ sở.
  \item[{\rm(iv)}] Nếu $f:X\to Y$ và $g:X'\to Y'$ là các $S$-cấu xạ xạ ảnh, thì $f\times_S g$ cũng vậy.
  \item[{\rm(v)}] Nếu $g\circ f$ là một cấu xạ xạ ảnh và $g$ tách, thì $f$ xạ ảnh.
  \item[{\rm(vi)}] Nếu $f$ xạ ảnh, thì $f_\red$ cũng vậy.
\end{enumerate}
\end{proposition}

\begin{proof}
(i) suy ra ngay từ \sref{II.3.1.7}.
Ở đây phải chứng minh riêng (iii) và (iv), vì hạn chế đặt lên $Z$ trong (ii) (xem \sref[I]{I.3.5.1}).
Để chứng minh (iii), ta quy về trường hợp $S=Y$ \sref[I]{I.3.3.11}, và khi đó khẳng định suy ra ngay từ \sref{II.5.5.1}[b] và \sref{II.3.5.3}.
Để chứng minh (iv), ta lập tức quy về trường hợp $X=\bb{P}(\sh{E})$ và $X'=\bb{P}(\sh{E}')$, trong đó $\sh{E}$ (tương ứng, $\sh{E}'$) là một $\sh{O}_Y$-môđun (tương ứng, một $\sh{O}_{Y'}$-môđun) quasi-coherent hữu hạn sinh.
Cho $p$ và $p'$ lần lượt là các phép chiếu chính tắc từ $T=Y\times_S Y'$ lên $Y$ và $Y'$; theo \sref{II.4.1.3.1}, ta có $\bb{P}(p^*(\sh{E}))=\bb{P}(\sh{E})\times_Y T$ và $\bb{P}(p^{\prime *}(\sh{E}'))=\bb{P}(\sh{E}')\times_{Y'}T$; do đó
\begin{align*}
  \bb{P}(p^*(\sh{E}))\times_T\bb{P}(p^{\prime *}(\sh{E}'))&=(\bb{P}(\sh{E})\times_Y T)\times_T(T\times_{Y'}\bb{P}(\sh{E}'))\\
                                               &=\bb{P}(\sh{E})\times_Y(T\times_{Y'}\bb{P}(\sh{E}'))=\bb{P}(\sh{E})\times_S\bb{P}(\sh{E}')
\end{align*}
khi thay $T$ bằng $Y\times_S Y'$ và dùng \sref[I]{I.3.3.9.1}.
Mà $p^*(\sh{E})$ và $p^{\prime *}(\sh{E}')$ đều hữu hạn sinh trên $T$ \sref[0]{0.5.2.4}, nên $p^*(\sh{E})\otimes_{\sh{O}_T}p^{\prime *}(\sh{E}')$ cũng vậy;
vì $\bb{P}(p^*(\sh{E}))\times_T\bb{P}(p^{\prime *}(\sh{E}'))$ được đồng nhất với một tiền lược đồ con đóng của $\bb{P}(p^*(\sh{E})\otimes_{\sh{O}_T}p^{\prime *}(\sh{E}'))$ \sref{II.4.3.3}, điều này hoàn tất chứng minh (iv).
Để chứng minh (v) và (vi), ta có thể áp dụng \sref[I]{I.5.5.13}, vì mọi tiền lược đồ con đóng của một $Y$-lược đồ xạ ảnh đều là một $Y$-lược đồ xạ ảnh theo \sref{II.5.5.1}[a].

Còn phải chứng minh (ii); nhờ giả thiết trên $Z$, điều này suy ra từ \sref{II.5.5.3}, \sref{II.5.3.4}[ii] và \sref{II.5.4.2}[ii].
\end{proof}

\begin{proposition}[5.5.6]
\label{II.5.5.6}
Nếu $X$ và $X'$ là hai $Y$-lược đồ xạ ảnh, thì $X\sqcup X'$ là một $Y$-lược đồ xạ ảnh.
\end{proposition}

\begin{proof}
Đây là hệ quả hiển nhiên của \sref{II.5.5.2} và \sref{II.4.3.6}.
\end{proof}

\begin{proposition}[5.5.7]
\label{II.5.5.7}
Cho $X$ là một $Y$-lược đồ xạ ảnh và $\sh{L}$ là một $\sh{O}_X$-môđun ample đối với $Y$;
với mọi tiết diện $f$ của $\sh{L}$ trên $X$, $X_f$ affine trên $Y$.
\end{proposition}
\oldpage[II]{106}
\begin{proof}
Vì vấn đề có tính địa phương trên $Y$, ta có thể giả sử $Y=\Spec(A)$; hơn nữa, $X_{f^{\otimes n}}=X_f$, nên bằng cách thay $\sh{L}$ bởi một $\sh{L}^{\otimes n}$ thích hợp, ta có thể giả sử $\sh{L}$ rất ample đối với cấu xạ cấu trúc $q:X\to Y$ \sref{II.4.6.11}.
Do đó đồng cấu chính tắc $\sigma:q^*(q_*(\sh{L}))\to\sh{L}$ là toàn ánh, và cấu xạ tương ứng
\[
  r=r_{\sh{L},\sigma}:X\to P=\bb{P}(q_*(\sh{L}))
\]
là một phép nhúng sao cho $\sh{L}=r^*(\sh{O}_P(1))$ \sref{II.4.4.4}; hơn nữa, vì $X$ riêng trên $Y$, phép nhúng $r$ là đóng \sref{II.5.4.4}.
Mặt khác, theo định nghĩa, $f\in\Gamma(Y,q_*(\sh{L}))$, và $\sigma^\flat$ là đồng nhất của $q_*(\sh{L})$; khi đó từ công thức \sref{II.3.7.3.1} suy ra $X_f=r^{-1}(D_+(f))$;
do đó $X_f$ là một tiền lược đồ con đóng của lược đồ affine $D_+(f)$, và vì thế là một lược đồ affine.
\end{proof}

Đặc biệt, lấy $Y=X$, ta thu được (có tính đến \sref{II.4.6.13}[i]) hệ quả sau, mà chứng minh trực tiếp cũng là hiển nhiên:
\begin{corollary}[5.5.8]
\label{II.5.5.8}
Cho $X$ là một tiền lược đồ và $\sh{L}$ là một $\sh{O}_X$-môđun khả nghịch.
Với mọi tiết diện $f$ của $\sh{L}$ trên $X$, $X_f$ affine trên $X$ (và do đó là một lược đồ affine khi $X$ là một lược đồ affine).
\end{corollary}

\subsection{Bổ đề Chow}
\label{subsection:II.5.6}

\begin{theorem}[5.6.1]
\label{II.5.6.1}
\emph{(Bổ đề Chow)}.
Cho $S$ là một tiền lược đồ và $X$ là một $S$-lược đồ thuộc kiểu hữu hạn.
Giả sử một trong các giả thiết sau được thỏa mãn:
\begin{enumerate}
  \item[{\rm(a)}] $S$ là Noether;
  \item[{\rm(b)}] $S$ là một lược đồ quasi-compact, và $X$ có hữu hạn thành phần bất khả quy.
\end{enumerate}
Trong các điều kiện ấy:
\begin{enumerate}
  \item[{\rm(i)}] tồn tại một $S$-lược đồ $X'$ \emph{quasi-xạ ảnh} trên $S$ và một $S$-cấu xạ $f:X'\to X$ vừa \emph{xạ ảnh} vừa \emph{toàn ánh};
  \item[{\rm(ii)}] có thể chọn $X'$ và $f$ sao cho tồn tại một tập mở $U\subset X$ mà $U'=f^{-1}(U)$ trù mật trong $X'$ và hạn chế của $f$ lên $U'$ là một đẳng cấu $U'\isoto U$;
  \item[{\rm(iii)}] nếu $X$ rút gọn (tương ứng, bất khả quy, nguyên), thì có thể giả sử $X'$ rút gọn (tương ứng, bất khả quy, nguyên).
\end{enumerate}
\end{theorem}

\begin{proof}
Chứng minh được tiến hành qua nhiều bước.
\begin{enumerate}
  \item[(A)] Trước hết, ta có thể quy về trường hợp $X$ \emph{bất khả quy}.
    Thật vậy, dưới giả thiết~(a), $X$ là Noether; do đó, dưới cả hai giả thiết, các thành phần bất khả quy $X_i$ của $X$ là hữu hạn.
    Nếu định lý đã được chứng minh cho mỗi tiền lược đồ con đóng rút gọn của $X$ có $X_i$ làm không gian nền, và nếu $X'_i$ cùng $f_i:X'_i\to X_i$ lần lượt là tiền lược đồ và cấu xạ tương ứng với $X_i$, thì tiền lược đồ $X'$ là \emph{tổng} của các $X'_i$, cùng cấu xạ $f:X'\to X$ có hạn chế trên mỗi $X'_i$ bằng $j_i\circ f_i$ (trong đó $j_i$ là phép nhúng chính tắc $X_i\to X$), thỏa mãn kết luận của định lý.
    Quả thật, hiển nhiên $X'$ rút gọn nếu mọi $X'_i$ đều rút gọn; mặt khác, (ii) được thỏa mãn khi lấy $U$ là hợp của các tập $U_i\cap\complement\left(\bigcup_{j\neq i}X_j\right)$.
    Cuối cùng, vì các $X'_i$ quasi-xạ ảnh trên $S$, nên $X'$ cũng vậy
\oldpage[II]{107}
    \sref{II.5.3.6}; tương tự, các cấu xạ $X'_i\to X$ là xạ ảnh theo \sref{II.5.5.5}[i] và \sref{II.5.5.5}[ii], nên $f$ xạ ảnh \sref{II.5.5.6}, và hiển nhiên toàn ánh theo định nghĩa.
  \item[(B)] Bây giờ giả sử $X$ \emph{bất khả quy}.
    Vì cấu xạ cấu trúc $r:X\to S$ thuộc kiểu hữu hạn, tồn tại một phủ hữu hạn $(S_i)$ của $S$ bởi các tập mở affine, và với mỗi $i$, một phủ hữu hạn $(T_{ij})$ của $r^{-1}(S_i)$ bởi các tập mở affine; cấu xạ $T_{ij}\to S_i$ là affine và thuộc kiểu hữu hạn, nên quasi-xạ ảnh \sref{II.5.3.4}[i];
    vì dưới mỗi giả thiết~(a), (b), phép nhúng $S_i\to S$ là quasi-compact, nó là một cấu xạ quasi-xạ ảnh \sref{II.5.3.4}[i], nên hạn chế của $r$ lên $T_{ij}$ là một cấu xạ quasi-xạ ảnh \sref{II.5.3.4}[ii].
    Ký hiệu các $T_{ij}$ bởi $U_k$ ($1\leq k\leq n$).
    Với mỗi chỉ số $k$, tồn tại một phép nhúng mở $\vphi_k:U_k\to P_k$, trong đó $P_k$ xạ ảnh trên $S$ (\sref{II.5.3.2} và \sref{II.5.5.2}).
    Đặt $U=\bigcap_k U_k$; vì $X$ bất khả quy và các $U_k$ khác rỗng, $U$ khác rỗng và do đó trù mật trong $X$; các hạn chế của $\vphi_k$ lên $U$ xác định một cấu xạ
    \[
      \vphi:U\to P=P_1\times_S P_2\times_S\cdots\times_S P_n
    \]
    sao cho các biểu đồ
    \[
    \label{II.5.6.1.1}
      \xymatrix{
        U\ar[r]^\vphi\ar[d]_{j_k} &
        P\ar[d]^{p_k}\\
        U_k\ar[r]_{\vphi_k} &
        P_k
      }
      \tag{5.6.1.1}
    \]
    giao hoán, trong đó $j_k$ là phép nhúng chính tắc $U\to U_k$ và $p_k$ là phép chiếu chính tắc $P\to P_k$.
    Nếu $j$ là phép nhúng chính tắc $U\to X$, thì cấu xạ $\psi=(j,\vphi)_S:U\to X\times_S P$ là một \emph{phép nhúng} \sref[I]{I.5.3.14}.
    Dưới giả thiết~(a), $X\times_S P$ là Noether địa phương (\sref{II.3.4.1}, \sref[I]{I.6.3.7} và \sref[I]{I.6.3.8});
    dưới giả thiết~(b), $X\times_S P$ là một lược đồ quasi-compact (\sref[I]{I.5.5.1} và \sref[I]{I.6.6.4});
    trong cả hai trường hợp, tồn tại bao đóng $X'$ trong $X\times_S P$ của tiền lược đồ con $Z$ liên kết với $\psi$ (và có không gian nền là $\psi(U)$), và $\psi$ phân tích thành
    \[
    \label{II.5.6.1.2}
      \psi:U\xrightarrow{\psi'}X'\xrightarrow{h}X\times_S P,
      \tag{5.6.1.2}
    \]
    trong đó $\psi'$ là một \emph{phép nhúng mở} và $h$ là một \emph{phép nhúng đóng} \sref[I]{I.9.5.10}.
    Cho $q_1:X\times_S P\to X$ và $q_2:X\times_S P\to P$ là các phép chiếu chính tắc; ta đặt
    \[
    \label{II.5.6.1.3}
      f:X'\xrightarrow{h}X\times_S P\xrightarrow{q_1}X,
      \tag{5.6.1.3}
    \]
    \[
    \label{II.5.6.1.4}
      g:X'\xrightarrow{h}X\times_S P\xrightarrow{q_2}P.
      \tag{5.6.1.4}
    \]
    Ta sẽ thấy $X'$ và $f$ thỏa mãn kết luận của định lý.
  \item[(C)] Trước hết, ta chứng minh rằng $f$ \emph{xạ ảnh} và \emph{toàn ánh}, và hạn chế của $f$ lên $U'=f^{-1}(U)$ là một \emph{đẳng cấu} từ $U'$ lên $U$.
    Vì các $P_k$ xạ ảnh trên $S$, $P$ cũng vậy \sref{II.5.5.5}[iv], nên $X\times_S P$ xạ ảnh trên $X$ \sref{II.5.5.5}[iii], và $X'$ cũng vậy, vì nó là một tiền lược đồ con đóng của $X\times_S P$.
    Mặt khác, ta có $f\circ\psi'=q_1\circ(h\circ\psi')=q_1\circ\psi=j$, nên $f(X')$ chứa tập mở trù mật $U$ của $X$; nhưng $f$ là một cấu xạ \emph{đóng} \sref{II.5.5.3}, do đó $f(X')=X$.
    Bây giờ nhận xét rằng $q_1^{-1}(U)=U\times_S P$ được cảm sinh trên một tập mở của $X\times_S P$, và theo định nghĩa, tiền lược đồ $U'=h^{-1}(U\times_S P)$ được $X'$ cảm sinh trên tập mở $U'$; vì thế nó là bao đóng \emph{tương đối trong}
\oldpage[II]{108}
$U\times_S P$ của tiền lược đồ $Z$ \sref[I]{I.9.5.8}.
Mà phép nhúng $\psi$ phân tích thành $U\xrightarrow{\Gamma_\vphi}U\times_S P\xrightarrow{j\times1}X\times_S P$, và vì $P$ tách trên $S$, cấu xạ đồ thị $\Gamma_\vphi$ là một phép nhúng đóng \sref[I]{I.5.4.3}; do đó $Z$ là một tiền lược đồ con \emph{đóng} của $U\times_S P$, suy ra $U'=Z$.
Hơn nữa, vì $\psi$ là một phép nhúng, hạn chế của $f$ lên $U'$ là một đẳng cấu lên $U$, nghịch đảo của $\psi'$; cuối cùng, theo định nghĩa của $X'$, $U'$ trù mật trong $X'$.
  \item[(D)] Bây giờ ta chứng minh rằng $g$ là một \emph{phép nhúng}; điều này sẽ cho thấy $X'$ \emph{quasi-xạ ảnh} trên $S$, vì $P$ xạ ảnh trên $S$.
    Đặt
    \begin{align*}
      V_k  &=\vphi_k(U_k) \quad(\text{tập mở của }P_k)\\
      W_k  &=p_k^{-1}(V_k)\quad(\text{tập mở của }P)\\
      U_k' &=f^{-1}(U_k)  \quad(\text{tập mở của }X')\\
      U_k''&=g^{-1}(W_k)  \quad(\text{tập mở của }X').
    \end{align*}
    Hiển nhiên các $U'_k$ tạo thành một phủ mở của $X'$; trước hết ta sẽ thấy các $U_k''$ cũng tạo thành một phủ mở của $X'$, bằng cách chứng minh $U_k'\subset U_k''$.
    Để làm điều đó, chỉ cần chứng minh biểu đồ
    \[
    \label{II.5.6.1.5}
      \xymatrix{
        U_k'\ar[r]^{g|U_k'}\ar[d]_{f|U_k'} &
        P\ar[d]^{p_k}\\
        U_k\ar[r]_{\vphi_k} &
        P_k
      }
      \tag{5.6.1.5}
    \]
    giao hoán.
    Mà tiền lược đồ $U_k'=h^{-1}(U_k\times_S P)$ được $X'$ cảm sinh trên tập mở $U_k'$, nên là bao đóng của $Z=U'\subset U_k'$ tương đối trong $U_k'$ \sref[I]{I.9.5.8}.
    Do đó, để chứng minh tính giao hoán của biểu đồ \sref{II.5.6.1.5}, vì $P_k$ là một $S$-lược đồ, chỉ cần chứng minh rằng khi hợp biểu đồ ấy với phép nhúng chính tắc $U'\to U_k'$ (hoặc, tương đương, nhờ đẳng cấu từ $U'$ lên $U$, với $\psi$), ta thu được một biểu đồ giao hoán \sref[I]{I.9.5.6}.
    Nhưng theo định nghĩa, biểu đồ thu được khi ấy chính là \sref{II.5.6.1.1}, từ đó suy ra khẳng định.

    Do đó các $W_k$ tạo thành một phủ mở của $g(X')$; để chứng minh $g$ là một phép nhúng, chỉ cần thấy rằng mỗi hạn chế $g|U_k''$ là một phép nhúng vào $W_k$ \sref[I]{I.4.2.4}.
    Để làm điều này, xét cấu xạ $u_k:W_k\xrightarrow{p_k}V_k\xrightarrow{\vphi_k^{-1}}U_k\to X$; vì $X$ tách trên $S$, cấu xạ đồ thị $\Gamma_{u_k}:W_k\to X\times_S W_k$ là một phép nhúng đóng \sref[I]{I.5.4.3}, nên đồ thị $T_k=\Gamma_{u_k}(W_k)$ là một tiền lược đồ con đóng của $X\times_S W_k$;
    nếu ta chứng minh cấu xạ $U'\to X\times_S W_k$ phân tích qua tiền lược đồ con này, thì cấu xạ đến $X\times_S W_k$ từ tiền lược đồ con do $X'$ cảm sinh trên tập mở $U_k''$ của $X'$ cũng phân tích qua đồ thị ấy, theo \sref[I]{I.9.5.8}.
    Vì hạn chế của $q_2$ lên $T_k$ là một đẳng cấu lên $W_k$, hạn chế của $g$ lên $U''_k$ sẽ là một phép nhúng vào $W_k$, và khẳng định của ta sẽ được chứng minh.
    Cho $v_k$ là phép nhúng chính tắc $U'\to X\times_S W_k$; ta phải chứng minh tồn tại một cấu xạ $w_k:U'\to W_k$ sao cho $v_k=\Gamma_{u_k}\circ w_k$.
    Theo định nghĩa của tích, chỉ cần chứng minh $q_1\circ v_k=u_k\circ q_2\circ v_k$ \sref[I]{I.3.2.1}, hoặc, khi hợp ở bên phải
\oldpage[II]{109}
với đẳng cấu $\psi':U\to U'$, chỉ cần chứng minh $q_1\circ\psi=u_k\circ q_2\circ\psi$.
Nhưng vì $q_1\circ\psi=j$ và $q_2\circ\psi=\vphi$, khẳng định của ta suy ra từ tính giao hoán của biểu đồ \sref{II.5.6.1.1}, có tính đến định nghĩa của $u_k$.
  \item[(E)] Hiển nhiên, vì $U$, và do đó $U'$, bất khả quy, nên $X'$ trong phép dựng trên cũng bất khả quy, và cấu xạ $f$ vì thế \emph{song hữu tỉ} \sref[I]{I.2.2.9}.
    Nếu thêm rằng $X$ rút gọn, thì $U'$ cũng rút gọn, và do đó $X'$ rút gọn \sref[I]{I.9.5.9}.
    Điều này hoàn tất chứng minh.
\end{enumerate}
\end{proof}

\begin{corollary}[5.6.2]
\label{II.5.6.2}
Giả sử một trong hai giả thiết \emph{(a), (b)} của \sref{II.5.6.1} được thỏa mãn.
Để $X$ riêng trên $S$, điều kiện cần và đủ là tồn tại một lược đồ $X'$ xạ ảnh trên $S$ và một $S$-cấu xạ toàn ánh $f:X'\to X$ (khi đó là xạ ảnh theo \sref{II.5.5.5}[v]).
Khi điều này xảy ra, còn có thể chọn $f$ sao cho tồn tại một tập mở trù mật $U$ của $X$ mà hạn chế của $f$ lên $f^{-1}(U)$ là một đẳng cấu $f^{-1}(U)\isoto U$ và $f^{-1}(U)$ trù mật trong $X'$.
Nếu thêm rằng $X$ bất khả quy (tương ứng, rút gọn), thì có thể giả sử $X'$ cũng vậy; khi $X$ và $X'$ bất khả quy, $f$ là một cấu xạ song hữu tỉ.
\end{corollary}

\begin{proof}
Điều kiện là đủ theo \sref{II.5.5.3} và \sref{II.5.4.3}[ii].
Điều kiện là cần vì, với các ký hiệu của \sref{II.5.6.1}, nếu $X$ riêng trên $S$, thì $X'$ riêng trên $S$, bởi nó xạ ảnh trên $X$ và do đó riêng trên $X$ \sref{II.5.5.3}, rồi khẳng định của ta suy ra từ \sref{II.5.4.2}[ii]; hơn nữa, vì $X'$ quasi-xạ ảnh trên $S$, nó xạ ảnh trên $S$ theo \sref{II.5.5.3}.
\end{proof}

\begin{corollary}[5.6.3]
\label{II.5.6.3}
Cho $S$ là một tiền lược đồ Noether địa phương, $X$ là một $S$-lược đồ thuộc kiểu hữu hạn trên $S$, và $f_0:X\to S$ là cấu xạ cấu trúc của nó.
Để $X$ riêng trên $S$, điều kiện cần và đủ là, với mọi cấu xạ \emph{thuộc kiểu hữu hạn} $S'\to S$, $(f_0)_{(S')}:X_{(S')}\to S'$ là một cấu xạ đóng.
Thậm chí, chỉ cần điều kiện này được thỏa mãn với mọi $S$-tiền lược đồ dạng $S'=S\otimes_\bb{Z}\bb{Z}[T_1,\ldots,T_n]$ (trong đó $T_1,\ldots,T_n$ là các phần tử bất định).
\end{corollary}

\begin{proof}
Điều kiện hiển nhiên là cần; ta chứng minh nó là đủ.
Vì vấn đề có tính địa phương trên $S$ và $S'$ \sref{II.5.4.1}, ta có thể giả sử $S$ và $S'$ affine và Noether.
Theo bổ đề Chow, tồn tại một $S$-lược đồ xạ ảnh $P$, một phép nhúng $j:X'\to P$ và một cấu xạ xạ ảnh, toàn ánh $f:X'\to X$, sao cho biểu đồ
\[
  \xymatrix{
    X\ar[d]_{f_0} &
    X'\ar[l]_f\ar[d]^j\\
    S &
    P\ar[l]_r
  }
\]
giao hoán.
Vì $P$ thuộc kiểu hữu hạn trên $S$, giả thiết thứ nhất kéo theo phép chiếu $q_2:X\times_S P\to P$ là một cấu xạ \emph{đóng}.
Cấu xạ $(f,j):X'\to X\times_S P$ phân tích thành phép nhúng đồ thị $\Gamma_j:X'\to X'\times_S P$, là phép nhúng đóng vì $P$ tách trên $S$, rồi đến phép đổi cơ sở $f\times1_P:X'\times_S P\to X\times_S P$.
Mà $f$, do xạ ảnh, nên riêng \sref{II.5.5.3}, vì thế $f\times1_P$ đóng; do đó $(f,j)$ đóng.
Vì $j=q_2\circ(f,j)$ và $j$ vốn đã là một phép nhúng, ta kết luận $j$ là một phép nhúng đóng, và do đó riêng \sref{II.5.4.2}[i].
Hơn nữa, cấu xạ cấu trúc $r:P\to S$ là xạ ảnh, nên riêng \sref{II.5.5.3}; do đó $f_0\circ f=r\circ j$ riêng \sref{II.5.4.2}[ii]; cuối cùng, vì $f$ toàn ánh, $f_0$ riêng theo \sref{II.5.4.3}.

Để chứng minh mệnh đề chỉ bằng giả thiết thứ hai, chỉ cần chứng minh giả thiết ấy kéo theo giả thiết thứ nhất.
Mà nếu $S'$ affine và thuộc kiểu hữu hạn trên $S=\Spec(A)$,
\oldpage[II]{110}
ta có $S'=\Spec(A[c_1,\ldots,c_n])$ \sref[I]{I.6.3.3}, do đó $S'$ đẳng cấu với một tiền lược đồ con đóng của $S''=\Spec(A[T_1,\ldots,T_n])$ (trong đó các $T_i$ là những biến bất định).
Trong sơ đồ giao hoán
\[
  \xymatrix{
    X\times_S S'\ar[r]^{1_X\times j}\ar[d]_{(f_0)_{(S')}} &
    X\times_S S''\ar[d]^{(f_0)_{(S'')}}\\
    S'\ar[r]_j &
    S''
  }
\]
cả $j$ lẫn $1_X\times j$ đều là những phép nhúng đóng \sref[I]{I.4.3.1}, và $(f_0)_{(S'')}$ là đóng theo giả thiết;
do đó $(f_0)_{(S')}$ cũng là đóng.
\end{proof}

\section{Các cấu xạ nguyên và các cấu xạ hữu hạn}
\label{section:II.6}

\subsection{Các tiền lược đồ nguyên trên một tiền lược đồ khác}
\label{subsection:II.6.1}

\begin{definition}[6.1.1]
\label{II.6.1.1}
Cho $X$ là một $S$-tiền lược đồ và $f:X\to S$ là cấu xạ cấu trúc của nó.
Ta nói $X$ \emph{nguyên trên $S$}, hoặc $f$ là một \emph{cấu xạ nguyên}, nếu tồn tại một phủ $(S_\alpha)$ của $S$ bởi các tập mở affine sao cho, với mọi $\alpha$, tiền lược đồ cảm sinh $f^{-1}(S_\alpha)$ là một lược đồ affine có vành $B_\alpha$ là một đại số nguyên trên vành $A_\alpha$ của $S_\alpha$.
Ta nói $X$ \emph{hữu hạn trên $S$}, hoặc $f$ là một \emph{cấu xạ hữu hạn}, nếu $X$ nguyên trên $S$ và thuộc kiểu hữu hạn trên $S$.
\end{definition}

Khi $S$ là affine, có vành $A$, ta cũng nói ``nguyên (tương ứng, hữu hạn) trên $A$'' thay cho ``nguyên (tương ứng, hữu hạn) trên $S$''.

\begin{env}[6.1.2]
\label{II.6.1.2}
Hiển nhiên, nếu $X$ nguyên trên $S$ thì $X$ \emph{affine} trên $S$.
Điều kiện cần và đủ để một tiền lược đồ $X$ affine trên $S$ là nguyên (tương ứng, hữu hạn) trên $S$ là $\sh{O}_S$-đại số quasi-coherent liên kết $\sh{A}(X)$ có tính chất sau: tồn tại một phủ $(S_\alpha)$ của $S$ bởi các tập mở affine sao cho, với mọi $\alpha$, $\Gamma(S_\alpha,\sh{A}(X))$ là một đại số nguyên (tương ứng, nguyên và hữu hạn sinh) trên $\Gamma(S_\alpha,\sh{O}_S)$.
Một $\sh{O}_S$-đại số quasi-coherent có tính chất này được gọi là \emph{nguyên} (tương ứng, \emph{hữu hạn}) trên $\sh{O}_S$.
Do đó, việc cho một tiền lược đồ nguyên (tương ứng, hữu hạn) trên $S$ tương đương \sref{II.1.3.1} với việc cho một $\sh{O}_S$-đại số quasi-coherent nguyên (tương ứng, hữu hạn) trên $\sh{O}_S$.
Ta lưu ý rằng một $\sh{O}_S$-đại số quasi-coherent $\sh{B}$ là hữu hạn khi và chỉ khi nó là một \emph{$\sh{O}_S$-môđun hữu hạn sinh} \sref[I]{I.1.3.9};
nói một cách tương đương, $\sh{B}$ là một $\sh{O}_S$-đại số \emph{nguyên} và \emph{hữu hạn sinh}, vì một đại số nguyên và hữu hạn sinh trên một vành $A$ là một $A$-môđun hữu hạn sinh.
\end{env}

\begin{proposition}[6.1.3]
\label{II.6.1.3}
Cho $S$ là một tiền lược đồ Noether địa phương.
Điều kiện cần và đủ để một tiền lược đồ $X$ affine trên $S$ là hữu hạn trên $S$ là $\sh{O}_S$-đại số $\sh{A}(X)$ coherent.
\end{proposition}
\begin{proof}
Xét nhận xét vừa rồi, điều này quy về việc lưu ý rằng, nếu $S$ là Noether địa phương, thì các $\sh{O}_S$-môđun quasi-coherent hữu hạn sinh chính là các $\sh{O}_S$-môđun coherent \sref[I]{I.1.5.1}.
\end{proof}

\begin{proposition}[6.1.4]
\label{II.6.1.4}
Cho $X$ là một tiền lược đồ nguyên (tương ứng, hữu hạn) trên $S$, và cho $f:X\to S$ là cấu xạ cấu trúc.
Khi đó, với mọi tập mở affine $U\subset S$ có vành $A$, $f^{-1}(U)$ là một lược đồ affine có vành $B$ là một đại số nguyên (tương ứng, hữu hạn) trên $A$.
\end{proposition}
\oldpage[II]{111}
\begin{proof}
Trước hết ta chứng minh bổ đề sau:

  \begin{lemma}[6.1.4.1]
  \label{II.6.1.4.1}
  Cho $A$ là một vành, $M$ là một $A$-môđun, và $(g_i)_{1\leq i\leq m}$ là một hệ hữu hạn các phần tử của $A$ sao cho các $D(g_i)$ ($1\leq i\leq m$) phủ $\Spec(A)$.
  Nếu, với mọi $i$, $M_{g_i}$ là một $A_{g_i}$-môđun hữu hạn sinh, thì $M$ là một $A$-môđun hữu hạn sinh.
  \end{lemma}

  \begin{proof}
  Thật vậy, ta có thể giả sử $M_{g_i}$ có một hệ sinh hữu hạn $(m_{ij}/g_i^n)$, với $m_{ij}\in M$, trong đó $n$ giống nhau đối với mọi chỉ số $i$.
  Ta chứng minh các $m_{ij}$ lập thành một hệ sinh của $M$.
  Cho $M'$ là $A$-môđun con của $M$ sinh bởi các phần tử ấy, và cho $m$ là một phần tử của $M$.
  Theo giả thiết, với mỗi $i$, tồn tại các $a_{ij}\in A$ và một số nguyên $p$ (không phụ thuộc vào $i$) sao cho, trong $M_{g_i}$, $m/1=(\sum_j a_{ij}m_{ij})/g_i^p$;
  điều này kéo theo tồn tại một số nguyên $r\geq p$ sao cho, với mọi $i$, ta có $g_i^rm\in M'$.
  Mà vì các $D(g_i^r)=D(g_i)$ phủ $\Spec(A)$, iđêan của $A$ sinh bởi các $g_i^r$ bằng $A$; nói cách khác, tồn tại các phần tử $a_i\in A$ sao cho $\sum_i a_ig_i^r=1$.
  Do đó $m=(\sum_i a_i g_i^r)m\in M'$, và bổ đề được chứng minh.
  \end{proof}

Ta đã biết \sref{II.1.3.2} rằng $f^{-1}(U)$ là affine.
Nếu $\vphi$ là đồng cấu $A\to B$ tương ứng với $f$, thì tồn tại một phủ hữu hạn của $U$ bởi các tập mở $D(g_i)$ (với $g_i\in A$) sao cho, nếu đặt $h_i=\vphi(g_i)$, thì $B_{h_i}$ là một đại số nguyên (tương ứng, nguyên và hữu hạn) trên $A_{g_i}$.
Thật vậy, tồn tại một phủ của $U$ bởi các tập mở affine $V_\alpha\subset U$ sao cho, nếu đặt $A_\alpha=A(V_\alpha)$ và $B_\alpha=A(f^{-1}(V_\alpha))$, thì $B_\alpha$ là một đại số nguyên (tương ứng, hữu hạn) trên $A_\alpha$.
Mỗi $x\in U$ thuộc một $V_\alpha$ nào đó; do đó tồn tại $g\in A$ sao cho $x\in D(g)\subset V_\alpha$.
Nếu $g_\alpha$ là ảnh của $g$ trong $A_\alpha$, thì $A(D(g))=A_g=(A_\alpha)_{g_\alpha}$;
đặt $h=\vphi(g)$, và cho $h_\alpha$ là ảnh của $g_\alpha$ trong $B_\alpha$;
ta có
\[
  A(D(h))=B_h=(B_\alpha)_{h_\alpha},
\]
và vì $B_\alpha$ nguyên (tương ứng, hữu hạn) trên $A_\alpha$, nên $(B_\alpha)_{h_\alpha}$ nguyên (tương ứng, hữu hạn) trên $(A_\alpha)_{g_\alpha}$.
Bây giờ chỉ cần dùng việc $U$ quasi-compact để thu được phủ cần tìm.

Trước hết, giả sử các $B_{h_i}$ nguyên và hữu hạn trên các $A_{g_i}$;
vì, với tư cách $A_{g_i}$-môđun, $B_{h_i}$ cũng được viết là $B_{g_i}$, Bổ đề~\sref{II.6.1.4.1} cho thấy trong trường hợp này $B$ là một $A$-môđun hữu hạn sinh.

Bây giờ chỉ giả sử mỗi $B_{h_i}$ nguyên trên $A_{g_i}$;
cho $b\in B$, và cho $C$ là $A$-đại số con của $B$ sinh bởi $b$.
Với mọi $i$, $C_{h_i}$ là đại số trên $A_{g_i}$ sinh bởi $b/1$ trong $B_{h_i}$;
từ giả thiết suy ra mỗi $C_{h_i}$ là một $A_{g_i}$-môđun hữu hạn sinh, nên theo \sref{II.6.1.4.1}, $C$ là một $A$-môđun hữu hạn sinh; điều này chứng minh $B$ nguyên trên $A$.
\end{proof}

\begin{proposition}[6.1.5]
\label{II.6.1.5}
\medskip\noindent
\begin{enumerate}
  \item[\rm{(i)}] Một phép nhúng đóng là hữu hạn (và, \emph{a fortiori}, nguyên).
  \item[\rm{(ii)}] Hợp thành của hai cấu xạ hữu hạn (tương ứng, nguyên) là hữu hạn (tương ứng, nguyên).
  \item[\rm{(iii)}] Nếu $f:X\to Y$ là một $S$-cấu xạ hữu hạn (tương ứng, nguyên), thì $f_{(S')}:X_{(S')}\to Y_{(S')}$ là hữu hạn (tương ứng, nguyên) đối với mọi mở rộng cơ sở $S'\to S$.
  \item[\rm{(iv)}] Nếu $f:X\to Y$ và $g:X'\to Y'$ là hai $S$-cấu xạ hữu hạn (tương ứng, nguyên), thì $f\times_S g:X\times_S X'\to Y\times_S Y'$ cũng vậy.
  \item[\rm{(v)}] Nếu $f:X\to Y$ và $g:Y\to Z$ là hai cấu xạ sao cho $g\circ f$ hữu hạn (tương ứng, nguyên), và nếu $g$ tách, thì $f$ hữu hạn (tương ứng, nguyên).
  \item[\rm{(vi)}] Nếu $f:X\to Y$ là một cấu xạ hữu hạn (tương ứng, nguyên), thì $f_\red$ cũng vậy.
\end{enumerate}
\end{proposition}
\oldpage[II]{112}
\begin{proof}
Theo \sref[I]{I.5.5.12}, chỉ cần chứng minh (i), (ii) và (iii).
Để chứng minh một phép nhúng đóng $X\to S$ là hữu hạn, ta có thể quy về trường hợp $S=\Spec(A)$; khi đó mọi việc quy về nhận xét rằng vành thương $A/\mathfrak{J}$ là một $A$-môđun cyclic.
Để chứng minh hợp thành của hai cấu xạ hữu hạn (tương ứng, nguyên) $X\to Y$, $Y\to Z$ là hữu hạn (tương ứng, nguyên), ta lại có thể giả sử $Z$ (và do đó $X$, $Y$ \sref{II.1.3.4}) là affine; khi đó khẳng định tương đương với việc nói rằng, nếu $B$ là một $A$-đại số hữu hạn (tương ứng, nguyên) và $C$ là một $B$-đại số hữu hạn (tương ứng, nguyên), thì $C$ là một $A$-đại số hữu hạn (tương ứng, nguyên), điều này là hiển nhiên.
Cuối cùng, để chứng minh (iii), ta có thể quy về trường hợp $S=Y$, vì $X_{(S')}$ được đồng nhất với $X\times_Y Y_{(S')}$ \sref[I]{I.3.3.11};
hơn nữa, ta có thể giả sử $S=\Spec(A)$ và $S'=\Spec(A')$.
Khi đó $X$ là affine với vành $B$ \sref{II.1.3.4}, còn $X_{(S')}$ là affine với vành $A'\otimes_A B$;
chỉ cần lưu ý rằng, nếu $B$ là một $A$-đại số hữu hạn (tương ứng, nguyên), thì $A'\otimes_A B$ là một $A'$-đại số hữu hạn (tương ứng, nguyên).
\end{proof}

Ta cũng lưu ý rằng, nếu $X$ và $Y$ là hai $S$-tiền lược đồ hữu hạn (tương ứng, nguyên) trên $S$, thì \emph{tổng} $X\sqcup Y$ của chúng là một tiền lược đồ hữu hạn (tương ứng, nguyên) trên $S$, vì điều này quy về việc thấy rằng, nếu $B$ và $C$ là hai $A$-đại số hữu hạn (tương ứng, nguyên) trên $A$, thì $B\times C$ cũng vậy.

\begin{corollary}[6.1.6]
\label{II.6.1.6}
Nếu $X$ là một tiền lược đồ nguyên (tương ứng, hữu hạn) trên $S$, thì với mọi tập mở $U\subset S$, $f^{-1}(U)$ nguyên (tương ứng, hữu hạn) trên $U$.
\end{corollary}

\begin{proof}
Đây là một trường hợp riêng của \sref{II.6.1.5}[(iii)].
\end{proof}

\begin{corollary}[6.1.7]
\label{II.6.1.7}
Cho $f:X\to Y$ là một cấu xạ hữu hạn.
Khi đó, với mọi $y\in Y$, thớ $f^{-1}(y)$ là một lược đồ đại số hữu hạn trên $\kres(y)$ và, \emph{a fortiori}, không gian nền của nó rời rạc và hữu hạn.
\end{corollary}

\begin{proof}
Thật vậy, với tư cách $\kres(y)$-tiền lược đồ, $f^{-1}(y)$ được đồng nhất với $X\times_Y\Spec(\kres(y))$ \sref[I]{I.3.6.1}, vốn hữu hạn trên $\Spec(\kres(y))$ \sref{II.6.1.5}[(iii)];
do đó nó là một lược đồ affine có vành là một đại số có hạng hữu hạn trên $\kres(y)$ \sref{II.6.1.4}.
Hệ quả khi đó suy ra từ \sref[I]{I.6.4.4}.
\end{proof}

\begin{corollary}[6.1.8]
\label{II.6.1.8}
Cho $X$, $S$ là hai tiền lược đồ nguyên, và cho $f:X\to S$ là một cấu xạ \emph{trội}.
Nếu $f$ nguyên (tương ứng, hữu hạn), thì trường $R(X)$ các hàm hữu tỉ trên $X$ là một mở rộng đại số (tương ứng, một mở rộng đại số bậc hữu hạn) của trường $R(S)$ các hàm hữu tỉ trên $S$.
\end{corollary}

\begin{proof}
Cho $s$ là điểm tổng quát của $S$;
$\kres(s)$-tiền lược đồ $f^{-1}(s)$ nguyên (tương ứng, hữu hạn) trên $\Spec(\kres(s))$ \sref{II.6.1.5}[(iii)] và, theo giả thiết, chứa điểm tổng quát $x$ của $X$;
vành địa phương tại $x$ trong $f^{-1}(s)$ bằng $\kres(x)$ \sref[I]{I.3.6.5}, nên là một vành địa phương của một đại số nguyên (tương ứng, hữu hạn) trên $\kres(s)$ \sref{II.6.1.4}; từ đó suy ra hệ quả.
\end{proof}

\begin{remark}[6.1.9]
\label{II.6.1.9}
Giả thiết $g$ \emph{tách} là thiết yếu đối với tính đúng đắn của \sref{II.6.1.5}[(v)]: thật vậy, nếu $Y$ không tách trên $Z$, thì đồng nhất $1_Y$ là cấu xạ hợp thành $Y\xrightarrow{\Delta_Y}Y\times_Z Y\xrightarrow{p_1}Y$, nhưng $\Delta_Y$ không phải là một cấu xạ nguyên, như suy ra từ \sref{II.6.1.10}:
\end{remark}

\begin{proposition}[6.1.10]
\label{II.6.1.10}
Mọi cấu xạ nguyên đều đóng phổ dụng.
\end{proposition}

\begin{proof}
Cho $f:X\to Y$ là một cấu xạ nguyên;
theo \sref{II.6.1.5}[(iii)], chỉ cần chứng minh $f$ là \emph{đóng}.
Cho $Z$ là một tập con đóng của $X$;
\oldpage[II]{113}
tồn tại một tiền lược đồ con của $X$ có không gian nền là $Z$ \sref[I]{I.5.2.1}; do đó, từ \sref{II.6.1.5}[(i) và (ii)] suy ra ta có thể quy về việc chứng minh $f(X)$ \emph{đóng} trong $Y$.
Theo \sref{II.6.1.5}[(vi)], ta có thể giả sử $X$ và $Y$ \emph{rút gọn};
hơn nữa, nếu $T$ là tiền lược đồ con đóng rút gọn của $Y$ có không gian nền là $\overline{f(X)}$ \sref[I]{I.5.2.1}, thì ta biết $f$ phân tích thành $X\xrightarrow{g}T\xrightarrow{j}Y$, trong đó $j$ là phép nhúng chính tắc \sref[I]{I.5.2.2}; vì $j$ tách \sref[I]{I.5.5.1}[(i)], từ \sref{II.6.1.5}[(v)] suy ra $g$ là một cấu xạ nguyên.
Do đó ta có thể giả sử $f(X)$ \emph{trù mật} trong $Y$.
Cuối cùng, vì vấn đề có tính địa phương trên $Y$, ta có thể quy về trường hợp $Y=\Spec(A)$.
Khi đó $X=\Spec(B)$, trong đó $B$ là một $A$-đại số nguyên trên $A$ \sref{II.6.1.4};
hơn nữa, $A$ rút gọn \sref[I]{I.5.1.4}, và giả thiết $f(X)$ trù mật trong $Y$ kéo theo đồng cấu $\vphi:A\to B$ tương ứng với $f$ là \emph{đơn ánh} \sref[I]{I.1.2.7}.
Trong các điều kiện này, nói $f(X)=Y$ có nghĩa là mỗi iđêan nguyên tố của $A$ là giao với $A$ của một iđêan nguyên tố của $B$; đây chính là định lý thứ nhất của Cohen--Seidenberg (\cite[t.~I, p.~257, th.~3]{I-13}).
\end{proof}

\begin{corollary}[6.1.11]
\label{II.6.1.11}
Mọi cấu xạ hữu hạn $f:X\to Y$ đều là xạ ảnh.
\end{corollary}

\begin{proof}
Thật vậy, vì $f$ là affine, $\sh{O}_X$ là một $\sh{O}_X$-môđun \emph{rất ample} đối với $f$ \sref{II.5.1.2};
hơn nữa, $f_*(\sh{O}_X)$ là một $\sh{O}_Y$-môđun quasi-coherent \emph{hữu hạn sinh} \sref{II.6.1.2};
cuối cùng, $f$ tách, thuộc kiểu hữu hạn và đóng phổ dụng \sref{II.6.1.10}, nên thỏa các điều kiện áp dụng của tiêu chuẩn~\sref{II.5.5.4}[(i)].
\end{proof}

\begin{proposition}[6.1.12]
\label{II.6.1.12}
Cho $f:X'\to X$ là một cấu xạ hữu hạn, và đặt $\sh{B}=f_*(\sh{O}_{X'})$ (đó là một $\sh{O}_X$-đại số quasi-coherent và một $\sh{O}_X$-môđun hữu hạn sinh).
Cho $\sh{F}'$ là một $\sh{O}_{X'}$-môđun quasi-coherent;
điều kiện cần và đủ để $\sh{F}'$ tự do địa phương hạng~$r$ là $f_*(\sh{F}')$ là một $\sh{B}$-môđun tự do địa phương hạng~$r$.
\end{proposition}

\begin{proof}
Hiển nhiên, nếu $f_*(\sh{F}')|U$ đẳng cấu với $\sh{B}^r|U$ (với $U\subset X$ mở), thì $\sh{F}'|f^{-1}(U)$ đẳng cấu với $\sh{O}_{X'}^r|f^{-1}(U)$ \sref{II.1.4.2}.
Ngược lại, giả sử $\sh{F}'$ tự do địa phương hạng~$r$, và ta chứng minh $f_*(\sh{F}')$ đẳng cấu địa phương với $\sh{B}^r$ với tư cách $\sh{B}$-môđun.
Cho $x$ là một điểm của $X$;
khi $U$ chạy qua một hệ cơ bản các lân cận affine của $x$, $f^{-1}(U)$ chạy qua một hệ cơ bản các lân cận affine \sref{II.1.2.5} của tập hữu hạn $f^{-1}(x)$, vì $f$ đóng \sref{II.6.1.10}.
Mệnh đề khi đó suy ra từ bổ đề sau:
\end{proof}

\begin{lemma}[6.1.12.1]
\label{II.6.1.12.1}
Cho $Y$ là một tiền lược đồ, $\sh{E}$ là một $\sh{O}_Y$-môđun tự do địa phương hạng~$r$, và $Z$ là một tập con hữu hạn của $Y$ nằm trong một tập mở affine $V$.
Khi đó tồn tại một lân cận $U\subset V$ của $Z$ sao cho $\sh{E}|U$ đẳng cấu với $\sh{O}_Y^r|U$.
\end{lemma}

\begin{proof}
Hiển nhiên ta có thể giả sử $Y$ affine;
với mọi $z_i\in Z$, trong bao đóng $\overline{\{z_i\}}$ tồn tại ít nhất một điểm đóng $z'_i$ \sref[0]{0.2.1.3};
nếu $Z'$ là tập các $z'_i$, thì mọi lân cận của $Z'$ đều là một lân cận của $Z$, nên ta có thể giả sử $Z$ đóng trong $Y$ và rời rạc.
Xét tiền lược đồ con đóng rút gọn của $Y$ có không gian nền là $Z$ \sref[I]{I.5.2.1}, và cho $j:Z\to Y$ là phép nhúng chính tắc;
$j^*(\sh{E})=\sh{E}\otimes_Y\sh{O}_Z$ tự do địa phương hạng~$r$ trên lược đồ rời rạc $Z$, nên đẳng cấu với $\sh{O}_Z^r$;
nói cách khác, tồn tại $r$ tiết diện $s_i$ ($1\leq i\leq r$) của $\sh{E}\otimes_Y\sh{O}_Z$ trên $Z$ sao cho đồng cấu $\sh{O}_Z^r\to\sh{E}\otimes_Y\sh{O}_Z$ được xác định bởi các tiết diện ấy là song ánh.
Nhưng $Y=\Spec(A)$ là affine, $Z$ được xác định bởi một iđêan $\mathfrak{J}$ của $A$, và ta có $\sh{E}=\widetilde{M}$, trong đó $M$ là một $A$-môđun;
các $s_i$ là những phần tử của $M\otimes_A(A/\mathfrak{J})$, nên là ảnh của $r$ phần tử $t_i\in M=\Gamma(Y,\sh{E})$.
\oldpage[II]{114}
Với mọi $z_j\in Z$, khi đó tồn tại một lân cận $V_j$ của $z_j$ sao cho các hạn chế của các $t_i$ lên $V_j$ xác định một đẳng cấu $\sh{O}_Y^r|V_j\to\sh{E}|V_j$ \sref[0]{0.5.5.4};
lân cận $U$ là hợp của các $V_j$ đáp ứng yêu cầu.
\end{proof}

\begin{proposition}[6.1.13]
\label{II.6.1.13}
Cho $g:X'\to X$ là một cấu xạ nguyên của các tiền lược đồ, $Y$ là một tiền lược đồ chuẩn tắc nguyên địa phương, và $f$ là một ánh xạ hữu tỉ từ $Y$ vào $X'$ sao cho $g\circ f$ là một ánh xạ hữu tỉ được xác định khắp nơi \sref[I]{I.7.2.1}.
Khi đó $f$ được xác định khắp nơi.
\end{proposition}

\begin{proof}
Nếu $f_1$ và $f_2$ là hai cấu xạ (từ các tập mở trù mật của $Y$ vào $X'$) thuộc lớp $f$, thì hiển nhiên $g\circ f_1$ và $g\circ f_2$ là các cấu xạ tương đương; điều này biện minh cho ký hiệu $g\circ f$ đối với lớp tương đương của chúng.
Ta cũng nhắc lại rằng, nếu giả sử thêm $Y$ \emph{Noether địa phương}, thì giả thiết $Y$ chuẩn tắc đã kéo theo $Y$ nguyên địa phương \sref[I]{I.6.1.13}.

Để chứng minh \sref{II.6.1.13}, trước hết ta lưu ý rằng vấn đề có tính địa phương trên $Y$, nên có thể giả sử trong lớp $g\circ f$ tồn tại một \emph{cấu xạ} $h:Y\to X$.
Xét ảnh ngược $Y'=X'_{(h)}=X'_{(Y)}$, và lưu ý rằng cấu xạ $g'=g_{(Y)}:Y'\to Y$ là \emph{nguyên} \sref{II.6.1.5}[(iii)].
Xét tương ứng giữa các ánh xạ hữu tỉ từ $Y$ vào $X'$ và các $Y$-tiết diện hữu tỉ của $Y'$ \sref[I]{I.7.1.2}, ta thấy mình được quy về việc chứng minh trường hợp riêng của \sref{II.6.1.13} trong đó $X=Y$; nói cách khác, hệ quả sau:
\end{proof}

\begin{corollary}[6.1.14]
\label{II.6.1.14}
Cho $X$ là một tiền lược đồ chuẩn tắc nguyên địa phương, $g:X'\to X$ là một cấu xạ \emph{nguyên}, và $f$ là một $X$-tiết diện hữu tỉ của $X'$.
Khi đó $f$ được xác định khắp nơi.
\end{corollary}

\begin{proof}
Vì vấn đề có tính địa phương trên $X$, ta có thể giả sử $X$ nguyên; khi đó $f$ được đồng nhất với một cấu xạ từ một tập mở $U$ của $X$ vào $X'$ \sref[I]{I.7.2.2}, là một $U$-tiết diện của $g^{-1}(U)$.
Vì $g$ tách, $f$ là một phép nhúng đóng của $U$ vào $g^{-1}(U)$ \sref[I]{I.5.4.6};
cho $Z$ là tiền lược đồ con đóng của $g^{-1}(U)$ liên kết với $f$ \sref[I]{I.4.2.1}; tiền lược đồ này đẳng cấu với $U$, nên nguyên;
cho $X_1$ là tiền lược đồ con rút gọn của $X'$ có không gian nền là bao đóng $\overline{Z}$ của $Z$ trong $X'$ \sref[I]{I.5.2.1};
khi đó $Z$ là một tiền lược đồ con cảm sinh trên một tập mở của $X_1$ \sref[I]{I.5.2.3}; vì $Z$ bất khả quy nên $X_1$ cũng bất khả quy, do đó nguyên.
Khi ấy có thể xem cấu xạ $f$ là một $X$-tiết diện hữu tỉ của $X_1$;
vì hạn chế của $g$ lên $X_1$ là một cấu xạ nguyên \sref{II.6.1.5}[(i) và (ii)], cuối cùng ta được quy về việc chứng minh \sref{II.6.1.14} trong trường hợp riêng $X'=X_1$; nói cách khác, hệ quả sau:
\end{proof}

\begin{corollary}[6.1.15]
\label{II.6.1.15}
Cho $X$ là một tiền lược đồ nguyên chuẩn tắc, $X'$ là một tiền lược đồ nguyên, và $g:X'\to X$ là một cấu xạ \emph{nguyên}.
Nếu tồn tại một $X$-tiết diện hữu tỉ $f$ của $X'$, thì $g$ là một đẳng cấu.
\end{corollary}

\begin{proof}
Vì vấn đề có tính địa phương trên $X$, ta có thể giả sử $X$ affine với miền nguyên $A$; khi đó $X'$ affine với vành $A'$ nguyên trên $A$ \sref{II.6.1.4}, và $A'$ là một miền nguyên;
hơn nữa, lập luận của \sref{II.6.1.14} cho thấy có một tập mở trù mật của $X$ đẳng cấu với một tập mở trù mật của $X'$, nên $A$ và $A'$ có cùng trường phân thức.
Mặt khác, theo \sref[I]{I.8.2.1.1} và giả thiết các $\sh{O}_x$ đóng nguyên, vành $A$ đóng nguyên;
do đó $A'=A$, hoàn tất chứng minh \sref{II.6.1.13}.
\end{proof}

\subsection{Các cấu xạ quasi-hữu hạn}
\label{subsection:II.6.2}

\begin{proposition}[6.2.1]
\label{II.6.2.1}
Cho $f:X\to Y$ là một cấu xạ thuộc kiểu hữu hạn địa phương và $x$ là một điểm của $X$.
Các điều kiện sau tương đương:
\oldpage[II]{115}
\begin{enumerate}
  \item[\rm{(a)}] Điểm $x$ cô lập trong thớ $f^{-1}(f(x))$ của nó.
  \item[\rm{(b)}] Vành $\sh{O}_x$ là một $\sh{O}_{f(x)}$-môđun quasi-hữu hạn \sref[0]{0.7.4.1}.
\end{enumerate}
\end{proposition}

\begin{proof}
Vì vấn đề hiển nhiên có tính địa phương trên $X$ và trên $Y$, ta có thể giả sử $X=\Spec(A)$ và $Y=\Spec(B)$ đều affine, trong đó $A$ là một $B$-đại số hữu hạn sinh \sref[I]{I.6.3.3}.
Hơn nữa, ta có thể thay $X$ bởi $X\times_Y\Spec(\sh{O}_{f(x)})$ mà không làm thay đổi thớ $f^{-1}(f(x))$ cũng như vành địa phương $\sh{O}_x$ \sref[I]{I.3.6.5};
do đó có thể giả sử $B$ là một vành địa phương (bằng $\sh{O}_{f(x)}$).
Nếu $\mathfrak{n}$ là iđêan tối đại của $B$, thì $f^{-1}(f(x))$ là một lược đồ affine có vành $A/\mathfrak{n}A$, thuộc kiểu hữu hạn trên $\kres(f(x))=B/\mathfrak{n}$ \sref[I]{I.6.4.11}.
Khi đó, nếu (a) được thỏa mãn, ta còn có thể giả sử $f^{-1}(f(x))$ chỉ gồm điểm $x$;
vì thế $A/\mathfrak{n}A$ có hạng hữu hạn trên $B/\mathfrak{n}$ \sref[I]{I.6.4.4}, nói cách khác, $A$ là một $B$-môđun quasi-hữu hạn.
Ngược lại, nếu đúng như vậy thì $f^{-1}(f(x))$ là một lược đồ affine Artin, nên rời rạc \sref[I]{I.6.4.4};
do đó $x$ cô lập trong thớ của nó, và điều này chứng tỏ (b) kéo theo (a).
\end{proof}

\begin{corollary}[6.2.2]
\label{II.6.2.2}
Cho $f:X\to Y$ là một cấu xạ thuộc kiểu hữu hạn.
Các điều kiện sau đây tương đương:
\begin{enumerate}
  \item[\rm{(a)}] Mọi điểm $x\in X$ đều cô lập trong thớ $f^{-1}(f(x))$ của nó (nói cách khác, không gian con $f^{-1}(f(x))$ là \emph{rời rạc}).
  \item[\rm{(b)}] Với mọi $x\in X$, tiền lược đồ $f^{-1}(f(x))$ là một $\kres(f(x))$-tiền lược đồ hữu hạn.
  \item[\rm{(c)}] Với mọi $x\in X$, vành $\sh{O}_x$ là một $\sh{O}_{f(x)}$-môđun quasi-hữu hạn.
\end{enumerate}
\end{corollary}

\begin{proof}
Sự tương đương giữa (a) và (c) suy ra từ \sref{II.6.2.1}.
Mặt khác, vì $f^{-1}(f(x))$ là một $\kres(f(x))$-tiền lược đồ đại số \sref[I]{I.6.4.11}, sự tương đương giữa (a) và (b) suy ra từ \sref[I]{I.6.4.4}.
\end{proof}

\begin{definition}[6.2.3]
\label{II.6.2.3}
Khi $f:X\to Y$ là một cấu xạ thuộc kiểu hữu hạn thỏa mãn các điều kiện tương đương của \sref{II.6.2.2}, ta nói $f$ là \emph{quasi-hữu hạn}, hay $X$ \emph{quasi-hữu hạn trên $Y$}.
\end{definition}

Hiển nhiên mọi cấu xạ \emph{hữu hạn} đều quasi-hữu hạn \sref{II.6.1.7}.

\begin{proposition}[6.2.4]
\label{II.6.2.4}
\medskip\noindent
\begin{enumerate}
  \item[\rm{(i)}] Một phép nhúng $X\to Y$ đóng, hoặc có $X$ là Noether, là một cấu xạ quasi-hữu hạn.
  \item[\rm{(ii)}] Nếu $f:X\to Y$ và $g:Y\to Z$ là các cấu xạ quasi-hữu hạn, thì $g\circ f$ quasi-hữu hạn.
  \item[\rm{(iii)}] Nếu $X$ và $Y$ là các $S$-tiền lược đồ và $f:X\to Y$ là một $S$-cấu xạ quasi-hữu hạn, thì $f_{(S')}:X_{(S')}\to Y_{(S')}$ quasi-hữu hạn đối với mọi phép mở rộng cơ sở $g:S'\to S$.
  \item[\rm{(iv)}] Nếu $f:X\to Y$ và $g:X'\to Y'$ là hai $S$-cấu xạ quasi-hữu hạn, thì
    \[
      f\times_S g:X\times_S X'\to Y\times_S Y'
    \]
    quasi-hữu hạn.
  \item[\rm{(v)}] Cho $f:X\to Y$ và $g:Y\to Z$ là hai cấu xạ sao cho $g\circ f$ quasi-hữu hạn; nếu thêm nữa $g$ tách, hoặc $X$ là Noether, hoặc $X\times_ZY$ là Noether địa phương, thì $f$ quasi-hữu hạn.
  \item[\rm{(vi)}] Nếu $f$ quasi-hữu hạn thì $f_\red$ cũng vậy.
\end{enumerate}
\end{proposition}

\begin{proof}
Nếu $f:X\to Y$ là một phép nhúng, thì mỗi thớ có nhiều nhất một điểm, và do đó khẳng định~(i) suy ra từ \sref[I]{I.6.3.4}[(i)] và \sref[I]{I.6.3.5}.
Để chứng minh (ii), trước hết lưu ý rằng $h=g\circ f$ thuộc kiểu hữu hạn \sref[I]{I.6.3.4}[(ii)];
hơn nữa, nếu $z=h(x)$ và $y=f(x)$, thì $y$ cô lập trong $g^{-1}(z)$, nên tồn tại một lân cận mở $V$ của $y$ trong $Y$ không gặp điểm nào của $g^{-1}(z)$ khác $y$;
do đó $f^{-1}(V)$ là một lân cận mở của $x$ không gặp bất kỳ
\oldpage[II]{116}
$f^{-1}(y')$ nào, trong đó $y'\neq y$ thuộc $g^{-1}(z)$;
vì $x$ cô lập trong $f^{-1}(y)$, nó cô lập trong $h^{-1}(z)=f^{-1}(g^{-1}(z))$.
Để chứng minh (iii), ta có thể quy về trường hợp $Y=S$ \sref[I]{I.3.3.11};
trước hết lại lưu ý rằng $f'=f_{(S')}$ thuộc kiểu hữu hạn \sref[I]{I.6.3.4}[(iii)];
mặt khác, nếu $x'\in X'=X_{(S')}$ và đặt $y'=f'(x')$, $y=g(y')$, thì ${f'}^{-1}(y')$ được đồng nhất với $f^{-1}(y)\otimes_{\kres(y)}\kres(y')$ \sref[I]{I.3.6.5};
vì theo giả thiết $f^{-1}(y)$ có hạng hữu hạn trên $\kres(y)$, nên ${f'}^{-1}(y')$ có hạng hữu hạn trên $\kres(y')$, do đó rời rạc.
Các khẳng định (iv), (v) và (vi) suy ra từ (i), (ii) và (iii) bằng phương pháp tổng quát \sref[I]{I.5.5.12}, trừ khi các giả thiết trong (v) khác với giả thiết ``$g$ tách'';
để xử lý trường hợp này, trước hết lưu ý rằng nếu $x$ cô lập trong $f^{-1}(g^{-1}(g(f(x))))$, thì nó \emph{a fortiori} cô lập trong $f^{-1}(f(x))$;
khi đó, việc $f$ thuộc kiểu hữu hạn suy ra từ \sref[I]{I.6.3.6}.
\end{proof}

\begin{proposition}[6.2.5]
\label{II.6.2.5}
Cho $A$ là một vành Noether địa phương \emph{đầy đủ}, $X$ là một $A$-lược đồ thuộc kiểu hữu hạn địa phương, và $x$ là một điểm của $X$ nằm trên điểm đóng $y$ của $Y=\Spec(A)$.
Giả sử $x$ cô lập trong thớ $f^{-1}(y)$ của nó (trong đó $f$ là cấu xạ cấu trúc $X\to Y$).
Khi đó $\sh{O}_x$ là một $A$-môđun hữu hạn sinh, và $X$ $Y$-đẳng cấu với tổng \sref[I]{I.3.1} của $X'=\Spec(\sh{O}_x)$ (là một $Y$-lược đồ hữu hạn) và một $A$-lược đồ $X''$.
\end{proposition}

\begin{proof}
Từ \sref{II.6.2.1} suy ra $\sh{O}_x$ là một $A$-môđun quasi-hữu hạn.
Vì $\sh{O}_x$ là Noether \sref[I]{I.6.3.7} và đồng cấu $A\to\sh{O}_x$ là địa phương, giả thiết $A$ \emph{đầy đủ} kéo theo $\sh{O}_x$ là một $A$-môđun \emph{hữu hạn sinh} \sref[0]{0.7.4.3}.
Cho $X'=\Spec(\sh{O}_x)$ là lược đồ địa phương của $X$ tại điểm $x$ \sref[I]{I.2.4.1}, và cho $g:X'\to X$ là cấu xạ chính tắc.
Vì hợp thành $X'\xrightarrow{g}X\xrightarrow{f}Y$ là hữu hạn \sref{II.6.1.1} và $f$ tách, nên $g$ hữu hạn \sref{II.6.1.5}[(v)], do đó $g(X')$ \emph{đóng} trong $X$ \sref{II.6.1.10};
mặt khác, vì $g$ thuộc kiểu hữu hạn, theo định nghĩa của $g$ và \sref[I]{I.6.5.4}, $g$ là một đẳng cấu địa phương tại điểm đóng $x'$ của $X'$;
nhưng vì $X'$ là lân cận mở duy nhất của $x'$, điều này có nghĩa $g$ là một phép nhúng mở, nên $g(X')$ cũng \emph{mở} trong $X$, và chứng minh hoàn tất.
\end{proof}

\begin{corollary}[6.2.6]
\label{II.6.2.6}
Cho $A$ là một vành Noether địa phương \emph{đầy đủ}, $Y=\Spec(A)$, và $f:X\to Y$ là một cấu xạ tách và quasi-hữu hạn.
Khi đó $X$ $Y$-đẳng cấu với một tổng $X'\sqcup X''$, trong đó $X'$ là một $Y$-lược đồ hữu hạn còn $X''$ là một $Y$-lược đồ quasi-hữu hạn sao cho, nếu $y$ là điểm đóng của $Y$, thì $X''\cap f^{-1}(y)=\emp$.
\end{corollary}

\begin{proof}
Thật vậy, theo giả thiết, thớ $f^{-1}(y)$ hữu hạn và rời rạc; do đó hệ quả suy ra từ \sref{II.6.2.5} bằng quy nạp theo số điểm của thớ này.
\end{proof}

\begin{remark}[6.2.7]
\label{II.6.2.7}
Trong Chương~V, ta sẽ thấy rằng nếu $Y$ \emph{Noether địa phương}, thì một cấu xạ tách và quasi-hữu hạn $X\to Y$ nhất thiết \emph{quasi-affine}.
\end{remark}

\subsection{Bao đóng nguyên của một tiền lược đồ}
\label{subsection:II.6.3}

\begin{proposition}[6.3.1]
\label{II.6.3.1}
Cho $(X,\sh{A})$ là một không gian vành hóa, $\sh{B}$ là một $\sh{A}$-đại số (giao hoán), và $f$ là một tiết diện của $\sh{B}$ trên $X$.
Các tính chất sau đây tương đương:
\begin{enumerate}
  \item[\rm{(a)}] $\sh{A}$-môđun con của $\sh{B}$ sinh bởi các $f^n$, với $n\geq0$ \sref[0]{0.5.1.1}, là hữu hạn sinh.
  \item[\rm{(b)}] Tồn tại một $\sh{A}$-đại số con $\sh{C}$ của $\sh{B}$ là một $\sh{A}$-môđun hữu hạn sinh và sao cho $f\in\Gamma(X,\sh{C})$.
  \item[\rm{(c)}] Với mọi $x\in X$, $f_x$ nguyên trên $\sh{A}_x$.
\end{enumerate}
\end{proposition}
\oldpage[II]{117}
\begin{proof}
Vì $\sh{A}$-môđun con của $\sh{B}$ sinh bởi các $f^n$ là một $\sh{A}$-đại số, hiển nhiên (a) kéo theo (b).
Mặt khác, (b) kéo theo rằng với mọi $x\in X$, $\sh{A}_x$-môđun $\sh{C}_x$ là hữu hạn sinh, từ đó suy ra mọi phần tử của đại số $\sh{C}_x$, và đặc biệt là $f_x$, đều nguyên trên $\sh{A}_x$.
Cuối cùng, nếu với một $x\in X$ nào đó ta có một hệ thức dạng
\[
  f_x^n+(a_1)_xf_x^{n-1}+\ldots+(a_n)_x=0
\]
trong đó các $a_i$ ($1\leq i\leq n$) là các tiết diện của $\sh{A}$ trên một lân cận mở $U$ của $x$, thì tiết diện $f^n|U+a_1\cdot f^{n-1}|U+\ldots+a_n$ bằng không trên một lân cận $V\subset U$ của $x$;
do đó suy ra ngay rằng mọi $f^k|V$ ($k\geq0$) đều là tổ hợp tuyến tính, với các hệ số trong $\Gamma(V,\sh{A})$, của các $f^j|V$ với $0\leq j\leq n-1$;
ta kết luận rằng (c) kéo theo (a).
\end{proof}

Khi các điều kiện tương đương của \sref{II.6.3.1} được thỏa mãn, ta nói tiết diện $f$ \emph{nguyên trên $\sh{A}$}.

\begin{corollary}[6.3.2]
\label{II.6.3.2}
Dưới các giả thiết của \sref{II.6.3.1}, tồn tại duy nhất một $\sh{A}$-môđun con $\sh{A}'$ của $\sh{B}$ sao cho, với mọi $x\in X$, $\sh{A}'_x$ là tập các mầm $f_x\in\sh{B}_x$ nguyên trên $\sh{A}_x$.
Với mọi tập mở $U\subset X$, các tiết diện của $\sh{A}'$ trên $U$ là những tiết diện $f\in\Gamma(U,\sh{B})$ nguyên trên $\sh{A}|U$.
\end{corollary}

\begin{proof}
Sự tồn tại của $\sh{A}'$ là hiển nhiên khi lấy $\Gamma(U,\sh{A}')$ là tập các $f\in\Gamma(U,\sh{B})$ sao cho $f_x$ nguyên trên $\sh{A}_x$ với mọi $x\in U$.
Khẳng định thứ hai suy ra ngay từ \sref{II.6.3.1}.
\end{proof}

Hiển nhiên $\sh{A}'$ là một $\sh{A}$-đại số con của $\sh{B}$;
ta nói nó là \emph{bao đóng nguyên của $\sh{A}$ trong $\sh{B}$}.

\begin{env}[6.3.3]
\label{II.6.3.3}
Cho $(X,\sh{A})$ và $(Y,\sh{B})$ là hai không gian vành hóa, và cho
\[
  g=(\psi,\theta):X\to Y
\]
là một cấu xạ.
Cho $\sh{C}$ (tương ứng, $\sh{D}$) là một $\sh{A}$-đại số (tương ứng, một $\sh{B}$-đại số), và cho
\[
  u:\sh{D}\to\sh{C}
\]
là một $g$-cấu xạ \sref[0]{0.4.4.1}.
Khi đó, nếu $\sh{A}'$ (tương ứng, $\sh{B}'$) là bao đóng nguyên của $\sh{A}$ (tương ứng, $\sh{B}$) trong $\sh{C}$ (tương ứng, $\sh{D}$), thì \emph{hạn chế} của $u$ lên $\sh{B}'$ là một $g$-cấu xạ
\[
  u':\sh{B}'\to\sh{A}'.
\]

Thật vậy, nếu $j$ là đơn ánh chính tắc $\sh{B}'\to\sh{D}$, chỉ cần chứng minh rằng
\[
  v=u^\sharp\circ g^*(j):g^*(\sh{B}')\to\sh{C}
\]
ánh xạ $g^*(\sh{B}')$ vào $\sh{A}'$.
Thật vậy, theo định nghĩa của $\sh{B}'$, một phần tử của $g^*(\sh{B}')_x=\sh{B}'_{\psi(x)}\otimes_{\sh{B}_{\psi(x)}}\sh{A}_x$ nguyên trên $\sh{A}_x$, và do đó ảnh của nó qua $v_x$ cũng vậy, điều này chứng minh khẳng định.
\end{env}

\begin{proposition}[6.3.4]
\label{II.6.3.4}
Cho $X$ là một tiền lược đồ và $\sh{A}$ là một $\sh{O}_X$-đại số quasi-coherent.
Khi đó bao đóng nguyên $\sh{O}'_X$ của $\sh{O}_X$ trong $\sh{A}$ là một $\sh{O}_X$-đại số quasi-coherent, và với mọi tập mở affine $U$ của $X$, $\Gamma(U,\sh{O}'_X)$ là bao đóng nguyên của $\Gamma(U,\sh{O}_X)$ trong $\Gamma(U,\sh{A})$.
\end{proposition}

\begin{proof}
Ta có thể quy về trường hợp $X=\Spec(B)$ affine và $\sh{A}=\widetilde{A}$, trong đó $A$ là một
\oldpage[II]{118}
$B$-đại số; cho $B'$ là bao đóng nguyên của $B$ trong $A$.
Mọi việc quy về chứng minh rằng, với mọi $x\in X$, một phần tử của $A_x$ nguyên trên $B_x$ nhất thiết thuộc $B'_x$; điều này suy ra từ việc, đối với một vành giao hoán $C$, hai phép lấy bao đóng nguyên trong một $C$-đại số và chuyển sang một vành các phân thức (đối với một tập con nhân của $C$) giao hoán với nhau \cite[t.~I, pp.~261 và 257]{I-13}.
\end{proof}

$X$-lược đồ $X'=\Spec(\sh{O}'_X)$ khi đó được gọi là \emph{bao đóng nguyên của $X$ đối với $\sh{A}$} (hay \emph{đối với $\Spec(\sh{A})$});
hiển nhiên $X'$ \emph{nguyên} trên $X$ \sref{II.6.1.2}.

Từ \sref{II.6.3.4} suy ra ngay rằng, nếu $f:X'\to X$ là cấu xạ cấu trúc, thì với mọi tập mở $U$ của $X$, $f^{-1}(U)$ là \emph{bao đóng nguyên của tiền lược đồ do $X$ cảm sinh trên $U$, đối với $\sh{A}|U$}.

\begin{env}[6.3.5]
\label{II.6.3.5}
Cho $X$ và $Y$ là hai tiền lược đồ, $f:X\to Y$ là một cấu xạ, $\sh{A}$ (tương ứng, $\sh{B}$) là một $\sh{O}_X$-đại số (tương ứng, một $\sh{O}_Y$-đại số) quasi-coherent, và cho $u:\sh{B}\to\sh{A}$ là một $f$-cấu xạ.
Ta đã thấy trong \sref{II.6.3.3} rằng từ đó suy ra một $f$-cấu xạ $u':\sh{O}'_Y\to\sh{O}'_X$, trong đó $\sh{O}'_X$ (tương ứng, $\sh{O}'_Y$) là bao đóng nguyên của $\sh{O}_X$ (tương ứng, $\sh{O}_Y$) trong $\sh{A}$ (tương ứng, $\sh{B}$).
Do đó, nếu $X'$ (tương ứng, $Y'$) là bao đóng nguyên của $X$ (tương ứng, $Y$) đối với $\sh{A}$ (tương ứng, $\sh{B}$), thì từ $u$ ta thu được một cách chính tắc một cấu xạ $f'=\Spec(u'):X'\to Y'$ \sref{II.1.5.6} làm cho biểu đồ
\[
\label{II.6.3.5.1}
  \xymatrix{
    X' \ar[r]^{f'} \ar[d]
    & Y' \ar[d]
  \\X \ar[r]_{f}
    & Y
  }
\tag{6.3.5.1}
\]
giao hoán.
\end{env}

\begin{env}[6.3.6]
\label{II.6.3.6}
Giả sử $X$ chỉ có một số \emph{hữu hạn} các thành phần bất khả quy $X_i$ ($1\leq i\leq r$), với các điểm tổng quát $\xi_i$, và đặc biệt xét bao đóng nguyên của $X$ đối với một $\sh{R}(X)$-đại số quasi-coherent $\sh{A}$ (dù xem là $\sh{O}_X$-đại số hay $\sh{R}(X)$-đại số, hai cách là như nhau).
Ta biết \sref[I]{I.7.3.5} rằng $\sh{A}$ là tổng trực tiếp của $r$ $\sh{O}_X$-đại số quasi-coherent $\sh{A}_i$, trong đó giá của $\sh{A}_i$ nằm trong $X_i$, và bó do $\sh{A}_i$ cảm sinh trên $X_i$ là một bó đơn có thớ $A_i$ là một đại số trên $\sh{O}_{\xi_i}$.
Khi đó, theo \sref{II.6.3.4}, hiển nhiên bao đóng nguyên $\sh{O}'_X$ của $\sh{O}_X$ trong $\sh{A}$ là \emph{tổng trực tiếp} của các bao đóng nguyên $\sh{O}_X^{(i)}$ của $\sh{O}_X$ trong từng $\sh{A}_i$, và do đó bao đóng nguyên $X'=\Spec(\sh{O}'_X)$ của $X$ đối với $\sh{A}$ là một $X$-lược đồ bằng tổng của các $\Spec(\sh{O}_X^{(i)})=X'_i$ ($1\leq i\leq r$).

Giả sử thêm rằng $\sh{O}_X$-đại số $\sh{A}$ \emph{rút gọn}, hay tương đương, mỗi đại số $A_i$ rút gọn và do đó có thể được xem là một đại số trên \emph{trường} $\kres(\xi_i)$ (bằng trường các hàm hữu tỉ của tiền lược đồ con rút gọn của $X$ có không gian nền là $X_i$);
khi đó, theo \sref{II.1.3.8}, mỗi $X'_i$ là một $X$-lược đồ \emph{rút gọn}, và $X'$ cũng là bao đóng nguyên của $X_\red$.
Giả sử thêm nữa mỗi đại số $A_i$ là \emph{tổng trực tiếp của một số hữu hạn các trường $K_{ij}$} ($1\leq j\leq s_i$);
nếu $\sh{K}_{ij}$ là đại số con của $\sh{A}_i$ tương ứng với $K_{ij}$, thì hiển nhiên $\sh{O}_X^{(i)}$ là \emph{tổng trực tiếp} của các bao đóng nguyên $\sh{O}_X^{(ij)}$ của $\sh{O}_X$ trong từng $\sh{K}_{ij}$.
Do đó $X'_i$ là tổng của các $X$-lược đồ $X'_{ij}=\Spec(\sh{O}_X^{(ij)})$ ($1\leq j\leq s_i$).
Hơn nữa, dưới các giả thiết và với các ký hiệu này:
\end{env}

\oldpage[II]{119}
\begin{proposition}[6.3.7]
\label{II.6.3.7}
Mỗi $X'_{ij}$ là một $X$-tiền lược đồ nguyên và chuẩn tắc, và trường các hàm hữu tỉ $R(X'_{ij})$ của nó được đồng nhất chính tắc với bao đóng đại số $K'_{ij}$ của $\kres(\xi_i)$ trong $K_{ij}$.
\end{proposition}

\begin{proof}
Theo những điều trên, ta có thể giả sử $X$ nguyên, do đó $r=1$ và $s_1=1$, sao cho đại số duy nhất $A_1$ là một trường $K$;
cho $\xi$ là điểm tổng quát của $X$, và cho $f:X'\to X$ là cấu xạ cấu trúc.
Với mọi tập mở affine khác rỗng $U$ của $X$, $f^{-1}(U)$ được đồng nhất với phổ của bao đóng nguyên $B'_U$ trong trường $K$ của vành nguyên $B_U=\Gamma(U,\sh{O}_X)$ \sref{II.6.3.4};
vì vành $B'_U$ nguyên và đóng nguyên, các vành địa phương tại các điểm thuộc phổ của nó cũng vậy, nên theo định nghĩa $f^{-1}(U)$ là một lược đồ \emph{nguyên} và \emph{chuẩn tắc} (\sref[0]{0.4.1.4} và \sref[I]{I.5.1.4}).
Hơn nữa, vì $(0)$ là iđêan nguyên tố duy nhất của $B'_U$ nằm trên iđêan nguyên tố $(0)$ của $B_U$ \cite[t.~I, p.~259]{I-13}, $f^{-1}(\xi)$ chỉ gồm một điểm $\xi'$, và $\kres(\xi')$ là trường phân thức $K'$ của $B'_U$, tức chính là bao đóng đại số của $\kres(\xi)$ trong $K$.
Cuối cùng, $X'$ bất khả quy, vì khi $U$ chạy qua tập các tập mở affine khác rỗng của $X$, các $f^{-1}(U)$ lập thành một phủ mở của $X'$ gồm các tập mở bất khả quy;
hơn nữa, giao $f^{-1}(U\cap V)$ của hai tập mở như vậy chứa $\xi'$, nên khác rỗng, và ta kết luận bằng \sref[0]{0.2.1.4}.
\end{proof}

\begin{corollary}[6.3.8]
\label{II.6.3.8}
Cho $X$ là một tiền lược đồ rút gọn chỉ có một số hữu hạn các thành phần bất khả quy $X_i$ ($1\leq i\leq r$), và cho $\xi_i$ là điểm tổng quát của $X_i$.
Bao đóng nguyên $X'$ của $X$ đối với $\sh{R}(X)$ là tổng của $r$ $X$-lược đồ $X'_i$ nguyên và chuẩn tắc.
Nếu $f:X'\to X$ là cấu xạ cấu trúc, thì $f^{-1}(\xi_i)$ chỉ gồm điểm tổng quát $\xi'_i$ của $X'_i$, và ta có $\kres(\xi'_i)=\kres(\xi_i)$; nói cách khác, $f$ song hữu tỉ.
\end{corollary}

Trong trường hợp này, ta nói $X'$ là \emph{chuẩn hóa} của tiền lược đồ rút gọn $X$;
lưu ý rằng vì $f$ song hữu tỉ và nguyên, nên nó \emph{toàn ánh} \sref{II.6.1.10}.
Điều kiện cần và đủ để có $X'=X$ là $X$ \emph{chuẩn tắc}.
Khi $X$ là một tiền lược đồ nguyên, từ \sref{II.6.3.8} suy ra chuẩn hóa $X'$ của nó là nguyên.

\begin{env}[6.3.9]
\label{II.6.3.9}
Cho $X$ và $Y$ là hai tiền lược đồ nguyên, $f:X\to Y$ là một cấu xạ trội, và đặt $L=R(X)$, $K=R(Y)$ lần lượt là trường các hàm hữu tỉ của $X$ và $Y$;
ứng với $f$ có một đơn ánh chính tắc $K\to L$, và khi đồng nhất $K$ (tương ứng, $L$) với bó đơn $\sh{R}(Y)$ (tương ứng, $\sh{R}(X)$), đơn ánh này là một $f$-cấu xạ.
Cho $K_1$ (tương ứng, $L_1$) là một mở rộng của $K$ (tương ứng, $L$), và giả sử đã cho một đơn cấu $K_1\to L_1$ sao cho biểu đồ
\[
  \xymatrix{
    K_1 \ar[r]
    & L_1
  \\K \ar[u] \ar[r]
    & L \ar[u]
  }
\]
giao hoán;
nếu $K_1$ (tương ứng, $L_1$) được xem là một bó đơn trên $Y$ (tương ứng, $X$), và do đó là một $\sh{R}(Y)$-đại số (tương ứng, một $\sh{R}(X)$-đại số), điều này có nghĩa $K_1\to L_1$ là một $f$-cấu xạ.
Khi đó, nếu $X'$ (tương ứng, $Y'$) là bao đóng nguyên của $X$ (tương ứng, $Y$) đối với $L_1$ (tương ứng, $K_1$), thì $X'$ (tương ứng, $Y'$) là một tiền lược đồ nguyên và chuẩn tắc \sref{II.6.3.6} có trường các hàm hữu tỉ được đồng nhất chính tắc với bao đóng đại số
\oldpage[II]{120}
$L'$ (tương ứng $K'$) của $L$ (tương ứng $K$) trong $L_1$ (tương ứng $K_1$), và tồn tại một cấu xạ chính tắc (tất yếu là trội) $f':X'\to Y'$ làm cho biểu đồ \sref{II.6.3.5.1} giao hoán.
Trường hợp quan trọng nhất là trường hợp lấy $L_1=L$, khi đó $K_1$ là một mở rộng của $K$ chứa trong $L$, và giả sử $X$ nguyên và \emph{chuẩn tắc}, do đó $X'=X$.
Từ những điều trên suy ra rằng, khi $X$ chuẩn tắc và $Y'$ là bao đóng nguyên của $Y$ đối với một trường $K_1\subset L=R(X)$, mọi cấu xạ trội $f:X\to Y$ đều phân tích thành
\[
  f:X\xrightarrow{f'}Y'\to Y
\]
trong đó $f'$ là trội;
hơn nữa, một khi đã cho đơn cấu $K_1\to L$, thì $f'$ tất yếu là duy nhất, như thấy được bằng cách quy về trường hợp $X$ và $Y$ affine.
Do đó, với dữ liệu gồm $Y$, $L$ và một $K$-đơn cấu $K_1\to L$, bao đóng nguyên $Y'$ của $Y$ đối với $K_1$ là nghiệm của một \emph{bài toán phổ dụng}.
\end{env}

\begin{remark}[6.3.10]
\label{II.6.3.10}
Xét lại các giả thiết của \sref{II.6.3.6}, đồng thời giả sử thêm rằng mỗi đại số $A_i$ có \emph{hạng hữu hạn} trên $\kres(\xi_i)$ (điều này kéo theo $A_i$ là tổng trực tiếp của một số hữu hạn trường);
trong một số trường hợp, ta có thể khẳng định rằng cấu xạ cấu trúc $X'\to X$ không chỉ nguyên mà còn \emph{hữu hạn}.
Ta chỉ xét trường hợp $X$ là \emph{rút gọn};
vì vấn đề có tính địa phương trên $X$, ta còn có thể giả sử $X$ là affine, có vành $C$, và $C$ chỉ có một số hữu hạn iđêan tối tiểu $\mathfrak{p}_i$ ($1\leq i\leq r$), sao cho các $C_i=C/\mathfrak{p}_i$ là các vành nguyên;
khi đó $X'$ sẽ hữu hạn trên $X$ nếu bao đóng nguyên của mỗi $C_i$ trong mọi mở rộng hữu hạn bậc của trường phân thức của nó là một \emph{$C$-môđun hữu hạn sinh} \sref{II.6.3.4}.
Ta biết điều kiện này luôn được thỏa mãn khi $C$ là một \emph{đại số hữu hạn sinh trên một trường} \cite[t.~I, p.~267, th.~9]{I-13}, hoặc \emph{trên $\bb{Z}$} \cite[I,p.~93, th.~3]{I-9}, hoặc \emph{trên một vành địa phương Noether đầy đủ} \cite[p.~298]{II-25}.
Suy ra $X'\to X$ là một cấu xạ hữu hạn khi $X$ là một lược đồ \emph{kiểu hữu hạn} trên một trường, trên $\bb{Z}$ hoặc trên một vành địa phương Noether đầy đủ.
\end{remark}

\subsection{Định thức của một tự đồng cấu của một $\mathcal{O}_X$-môđun}
\label{subsection:II.6.4}

\begin{env}[6.4.1]
\label{II.6.4.1}
Cho $A$ là một vành (giao hoán), $E$ là một $A$-môđun tự do hạng~$n$, và $u$ là một tự đồng cấu của $E$;
nhắc lại rằng, để định nghĩa \emph{đa thức đặc trưng} của $u$, ta xét tự đồng cấu $u\otimes1$ của $A[T]$-môđun tự do hạng~$n$, $E\otimes_A A[T]$ (với $T$ là một biến bất định), và đặt
\[
\label{II.6.4.1.1}
  P(u,T) = \det(T\cdot I-(u\otimes 1))
\tag{6.4.1.1}
\]
(với $I$ là tự đẳng cấu đồng nhất của $E\otimes_A A[T]$).
Ta có
\[
\label{II.6.4.1.2}
  P(u,T) = T^n - \sigma_1(u)T^{n-1} + \ldots + (-1)^n\sigma_n(u)
\tag{6.4.1.2}
\]
trong đó $\sigma_i(u)$ là một phần tử của $A$, bằng một đa thức thuần nhất bậc~$i$ (với các hệ số nguyên) theo các phần tử của ma trận của $u$ đối với một cơ sở tùy ý của $E$;
ta gọi các $\sigma_i(u)$ là các \emph{hàm đối xứng sơ cấp} của $u$;
đặc biệt, ta có $\sigma_1(u)=\Tr u$ và $\sigma_n(u)=\det u$.
Nhắc lại rằng, theo định lý Hamilton--Cayley, ta có
\[
\label{II.6.4.1.3}
  P(u,u) = u^n - \sigma_1(u)u^{n-1} + \ldots + (-1)^n\sigma_n(u) = 0
\tag{6.4.1.3}
\]
\oldpage[II]{121}
điều này cũng được viết thành
\[
\label{II.6.4.1.4}
  (\det u)\cdot I_E = uQ(u)
\tag{6.4.1.4}
\]
(với $I_E$ là tự đẳng cấu đồng nhất của $E$), trong đó
\[
\label{II.6.4.1.5}
  Q(u) = (-1)^{n+1}(u^{n-1}-\sigma_1(u)u^{n-2}+\ldots+(-1)^{n-1}\sigma_{n-1}(u)).
\tag{6.4.1.5}
\]

Cho $\vphi:A\to B$ là một đồng cấu vành làm cho $B$ thành một $A$-đại số;
xét $B$-môđun $E_{(B)}=E\otimes_A B$, là môđun tự do hạng~$n$, và mở rộng $u\otimes1$ của $u$ thành một tự đồng cấu của $E_{(B)}$;
hiển nhiên $\sigma_i(u\otimes1)=\vphi(\sigma_i(u))$ với mọi chỉ số $i$.
\end{env}

\begin{env}[6.4.2]
\label{II.6.4.2}
Bây giờ giả sử $A$ là một \emph{vành nguyên}, có trường phân thức $K$, và $E$ là một $A$-môđun \emph{hữu hạn sinh} (nhưng không nhất thiết tự do).
Gọi $n$ là \emph{hạng} của $E$, tức là hạng của $K$-môđun tự do $E\otimes_A K$;
với mỗi tự đồng cấu $u$ của $E$, ta có một cách chính tắc tự đồng cấu $u\otimes1$ của $E\otimes_A K$;
theo lối lạm dụng ngôn ngữ, ta vẫn gọi \emph{đa thức đặc trưng} của $u$, và ký hiệu là $P(u,T)$, đa thức $P(u\otimes1,T)$; các hệ số $\sigma_i(u\otimes1)$ của nó (thuộc $K$) vẫn được ký hiệu là $\sigma_i(u)$ và được gọi là các \emph{hàm đối xứng sơ cấp của $u$};
đặc biệt, theo định nghĩa, $\det u=\det(u\otimes1)$.
Với các ký hiệu này, các công thức từ \sref{II.6.4.1.3} đến \sref{II.6.4.1.5} có nghĩa và vẫn đúng, nếu hiểu $u^j$ là hợp thành
\[
  E\xrightarrow{x\mapsto x\otimes1}E\otimes_A K
   \xrightarrow{(u\otimes1)^j}E\otimes_A K.
\]

Nếu $F$ là môđun con xoắn của $E$ và $E_0=E/F$, thì $u(F)\subset F$ và, qua việc lấy thương, $u$ cảm sinh một tự đồng cấu $u_0$ của $E_0$;
hơn nữa, $E\otimes_A K$ được đồng nhất với $E_0\otimes_A K$, và $u\otimes1$ với $u_0\otimes1$, do đó $\sigma_i(u)=\sigma_i(u_0)$ với $1\leq i\leq n$.

Nếu $E$ là \emph{không xoắn}, thì $E$ được đồng nhất một cách chính tắc với một $A$-môđun con của $E\otimes_A K$, và hệ thức $u\otimes1=0$ tương đương với $u=0$.
Nếu $E$ là một $A$-môđun tự do, thì hai định nghĩa của $\sigma_i(u)$ cho trong \sref{II.6.4.1} và \sref{II.6.4.2} trùng nhau theo những điều trên, điều này biện minh cho các ký hiệu đã dùng.

Cũng lưu ý rằng, khi $E$ là một \emph{môđun xoắn}, hay nói cách khác khi $E_0=\{0\}$, đại số ngoài của $E_0$ thu về $K$, và định thức của tự đồng cấu duy nhất $u_0$ của $E_0$ \emph{bằng $1$}.
\end{env}

\begin{proposition}[6.4.3]
\label{II.6.4.3}
Cho $A$ là một vành nguyên, $E$ là một $A$-môđun hữu hạn sinh và $u$ là một tự đồng cấu của $E$;
khi đó các hàm đối xứng sơ cấp $\sigma_i(u)$ của $u$ (và đặc biệt $\det u$) là các phần tử của $K$ nguyên trên $A$.
\end{proposition}

\begin{proof}
Gọi $A'$ là bao đóng nguyên của $A$;
vì $A'[T]$ là một vành đóng nguyên \cite[p.~99]{II-24}, nên nó là bao đóng nguyên của $A[T]$ trong trường phân thức $K(T)$ của vành này.
Thay $u$ bởi $T\cdot I-u\otimes1$ và $A$ bởi $A[T]$, ta thấy có thể quy về việc chứng minh rằng $\det u$ nguyên trên $A$.
Nếu $n$ là hạng của $E$, thì $\det(u)=\det(\wedge^n u)$ và $(\wedge^n u)\otimes1=\wedge^n(u\otimes1)$, do đó có thể giả sử $n=1$.
Nhưng khi đó ánh xạ $u\mapsto\det u$ là một đồng cấu từ $A$-môđun $\Hom_A(E,E)$ vào $K$;
vì $E$ hữu hạn sinh, $\Hom_A(E,E)$ đẳng cấu với một $A$-môđun con của $A$-môđun hữu hạn sinh $E^m$ (nếu $E$ có một hệ gồm $m$ phần tử sinh); do đó các phần tử $\det u$ thuộc một $A$-môđun con hữu hạn sinh của $K$, và vì thế nguyên trên $A$.
\end{proof}
\oldpage[II]{122}
\begin{corollary}[6.4.4]
\label{II.6.4.4}
Với các giả thiết của \sref{II.6.4.3}, nếu giả sử thêm rằng $A$ là chuẩn tắc, thì các $\sigma_i(u)$ (và đặc biệt $\Tr u$ và $\det u$) thuộc $A$.
\end{corollary}

\begin{proposition}[6.4.5]
\label{II.6.4.5}
Cho $A$ là một vành nguyên, $E$ là một $A$-môđun hữu hạn sinh hạng~$n$, và $u$ là một tự đồng cấu của $E$ sao cho các $\sigma_i(u)$ thuộc $A$.
Để $u$ là một tự đẳng cấu của $E$, điều kiện cần là $\det u$ khả nghịch trong $A$;
điều kiện này là đủ nếu $E$ không xoắn.
\end{proposition}

\begin{proof}
Điều kiện là \emph{đủ} khi $E$ không xoắn, vì giả thiết cùng các công thức \sref{II.6.4.1.4} và \sref{II.6.4.1.5} (đúng trong $E$, chứ không chỉ trong $E\otimes_A K$, vì $E$ không xoắn) chứng minh rằng $(\det u)^{-1}Q(u)$ là nghịch đảo của $u$.

Điều kiện là \emph{cần}, vì nếu $u$ khả nghịch thì theo \sref{II.6.4.3}, $\det(u^{-1})$ thuộc bao đóng nguyên $A'$ của $A$ trong trường phân thức $K$ của nó và hiển nhiên là nghịch đảo của $\det u$ \emph{trong $A'$}.
Khẳng định của ta khi đó suy ra từ bổ đề sau:
\end{proof}

\begin{lemma}[6.4.5.1]
\label{II.6.4.5.1}
Cho $A$ là một vành con của một vành $A'$ sao cho $A'$ \emph{nguyên} trên $A$.
Nếu một phần tử $x\in A$ khả nghịch trong $A'$, thì nó cũng khả nghịch trong $A$.
\end{lemma}

\begin{proof}
Nếu không, $x$ sẽ thuộc một iđêan tối đại $\mathfrak{m}$ của $A$, và từ định lý Cohen--Seidenberg thứ nhất \cite[t.~I, p.~257, th.~3]{I-13} sẽ suy ra rằng tồn tại một iđêan tối đại $\mathfrak{m}'$ của $A'$ sao cho $\mathfrak{m}=\mathfrak{m}'\cap A$;
khi đó $x\in\mathfrak{m}'$, điều này vô lý.
\end{proof}

\begin{corollary}[6.4.6]
\label{II.6.4.6}
Cho $A$ là một vành nguyên và đóng nguyên, $E$ là một $A$-môđun không xoắn hữu hạn sinh, và $u$ là một tự đồng cấu của $E$.
Để $u$ là một tự đẳng cấu của $E$, điều kiện cần và đủ là $\det u$ khả nghịch trong $A$.
\end{corollary}

\begin{proof}
Điều này suy ra từ \sref{II.6.4.4} và \sref{II.6.4.5}.
\end{proof}

\begin{remark}[6.4.7]
\label{II.6.4.7}
Về sau ta sẽ cần một dạng tổng quát hóa của các kết quả trên.
Xét một vành Noether \emph{rút gọn} $A$;
gọi $\mathfrak{p}_\alpha$ ($1\leq\alpha\leq r$) là các iđêan tối tiểu của nó, $K_\alpha$ là trường phân thức của vành nguyên $A/\mathfrak{p}_\alpha$, và $K$ là vành phân thức toàn phần của $A$, tức là \emph{tổng trực tiếp} của các trường $K_\alpha$.
Cho $E$ là một $A$-môđun hữu hạn sinh, và \emph{giả sử} $E\otimes_A K$ là một $K$-môđun \emph{tự do} hạng~$n$ (ở đây điều này không còn là hệ quả của các giả thiết khác);
điều đó tương đương với việc mọi $K_\alpha$-không gian vectơ $E\otimes_A K_\alpha=E_\alpha$ đều có \emph{cùng số chiều~$n$};
nếu $u$ là một tự đồng cấu của $E$, ta vẫn đặt $P(u,T)=P(u\otimes1,T)$ và $\sigma_j(u)=\sigma_j(u\otimes1)$, và đặc biệt $\det u=\det(u\otimes1)$;
do đó các $\sigma_j(u)$ là các phần tử của $K$.
Hiển nhiên $E\otimes_A K$ là tổng trực tiếp của các $E_\alpha$, và các môđun này ổn định dưới $u\otimes1$;
hạn chế của $u\otimes1$ lên $E_\alpha$ chính là mở rộng $u_\alpha$ của $u$ lên $K_\alpha$-không gian vectơ này;
suy ra $\sigma_j(u)$ là phần tử của $K$ có các thành phần trong các $K_\alpha$ bằng các $\sigma_j(u_\alpha)$.
Vì bao đóng nguyên của $A$ trong $K$ là \emph{tổng trực tiếp} của các bao đóng nguyên của $A$ trong các $K_\alpha$, nên các $\sigma_j(u)$ \emph{nguyên} trên $A$.

\begin{lemma}[6.4.7.1]
\label{II.6.4.7.1}
$A$-đại số con của $K$ sinh bởi tất cả các phần tử $\sigma_j(u)$ ($1\leq j\leq n$), khi $u$ chạy qua $\Hom_A(E,E)$, là một $A$-môđun hữu hạn sinh.
\end{lemma}

\begin{proof}
Chỉ cần chứng minh rằng $A[T]$-đại số con của $K[T]$ sinh bởi các $P(u,T)$ là một $A[T]$-môđun hữu hạn sinh, vì nếu các $F_i(T)$ ($1\leq i\leq m$) tạo thành một hệ sinh của $A[T]$-môđun này, thì các hệ số của các $F_i(T)$ nguyên trên $A$, và do đó sinh ra một $A$-đại số là một $A$-môđun hữu hạn sinh \cite[t.~I, p.~255, th.~1]{I-13}.
Do đó có thể
\oldpage[II]{123}
thay $A$ bởi $A[T]$ (là vành Noether) và $E$ bởi $E\otimes_A A[T]=E'$, sao cho $E'\otimes_{A[T]}K[T]=E\otimes_A K[T]$ là một $K[T]$-môđun tự do hạng~$n$.
Trở lại các ký hiệu ban đầu, ta thấy chỉ cần chứng minh rằng $A$-môđun sinh bởi các phần tử $\det u$, khi $u$ chạy qua $\Hom_A(E,E)$, là hữu hạn sinh;
\emph{a fortiori} (vì mọi môđun con của một $A$-môđun hữu hạn sinh đều hữu hạn sinh), chỉ cần chứng minh rằng, khi $v$ chạy qua tập các tự đồng cấu của $\wedge^n E$, $A$-môđun sinh bởi các $\det v$ là hữu hạn sinh;
nói cách khác, ta lại có thể quy về trường hợp $n=1$.
Nhưng khi đó mệnh đề suy ra từ việc $\Hom_A(E,E)$ là một $A$-môđun hữu hạn sinh và $v\mapsto\det v$ là một đồng cấu từ $A$-môđun này vào $K$.
\end{proof}

Gọi $F$ là hạt nhân của đồng cấu chính tắc $E\to E\otimes_A K$, và đặt $E_0=E/F$;
như trên, ta thấy $E\otimes_A K$ được đồng nhất với $E_0\otimes_A K$, rằng $u(F)\subset F$, và nếu $u_0$ là tự đồng cấu của $E_0$ suy ra từ $u$ bằng cách lấy thương, thì $u\otimes1$ được đồng nhất với $u_0\otimes1$, và $\sigma_j(u)=\sigma_j(u_0)$ với mọi $j$.
\emph{Nếu $F=0$}, các công thức \sref{II.6.4.1.3} và \sref{II.6.4.1.5} vẫn có nghĩa và vẫn đúng khi $E$ được đồng nhất với một môđun con của $E\otimes_A K$, và các $u^j$ với các đồng cấu $E\to E\otimes_A K$;
do đó Mệnh đề~\sref{II.6.4.5} cũng mở rộng sang trường hợp này, với cùng một chứng minh.
\end{remark}

\begin{env}[6.4.8]
\label{II.6.4.8}
Cho $(X,\sh{A})$ là một không gian vành hóa, $\sh{E}$ là một $\sh{A}$-môđun \emph{tự do địa phương} (có hạng hữu hạn), và $u$ là một tự đồng cấu của $\sh{E}$.
Theo giả thiết, tồn tại một cơ sở $\mathfrak{B}$ của tôpô trên $X$ sao cho, với mọi $V\in\mathfrak{B}$, $\sh{E}|V$ đẳng cấu với một $\sh{A}^n|V$ nào đó (trong đó $n$ có thể thay đổi theo $V$).
Cho $u$ là một tự đồng cấu của $\sh{E}$;
với mọi $V\in\mathfrak{B}$, $u_V$ do đó là một tự đồng cấu của $\Gamma(V,\sh{A})$-môđun $\Gamma(V,\sh{E})$, môđun này tự do theo giả thiết;
vì thế định thức $\det u_V$ được xác định và thuộc $\Gamma(V,\sh{A})$.
Hơn nữa, nếu $e_1,\ldots,e_n$ tạo thành một cơ sở của $\Gamma(V,\sh{E})$, thì các hạn chế của chúng lên mọi tập mở $W\subset V$ tạo thành một cơ sở của $\Gamma(W,\sh{E})$ trên $\Gamma(W,\sh{A})$, do đó $\det u_W$ là hạn chế của $\det u_V$ lên $W$.
Vì thế tồn tại duy nhất một tiết diện của $\sh{A}$ trên $X$, ký hiệu là $\det u$ và gọi là \emph{định thức} của $u$, sao cho hạn chế của $\det u$ lên mọi $V\in\mathfrak{B}$ bằng $\det u_V$.
Rõ ràng, với mọi $x\in X$, ta có $(\det u)_x=\det u_x$;
với hai tự đồng cấu $u$ và $v$ của $\sh{E}$, ta có
\[
\label{II.6.4.8.1}
  \det(u\circ v) = (\det u)(\det v)
\tag{6.4.8.1}
\]
cũng như
\[
\label{II.6.4.8.2}
  \det(1_{\sh{E}}) = 1_{\sh{A}}
\tag{6.4.8.2}
\]
và nếu hạng của $\sh{E}$ là không đổi (điều này đúng \sref[0]{0.5.4.1} nếu $X$ liên thông) và bằng $n$, thì
\[
\label{II.6.4.8.3}
  \det(s\cdot u) = s^n\det u
\tag{6.4.8.3}
\]
với mọi $s\in\Gamma(X,\sh{A})$ (lưu ý rằng $\det(0)=0_{\sh{A}}$ nếu $n\geq1$, nhưng $\det(0)=1_{\sh{A}}$ nếu $n=0$).
Hơn nữa, để $u$ là một \emph{tự đẳng cấu} của $\sh{E}$, điều kiện cần và đủ là $\det u$ \emph{khả nghịch} trong $\Gamma(X,\sh{A})$.

Nếu hạng của $\sh{E}$ là không đổi, ta có thể định nghĩa theo cùng cách các hàm đối xứng sơ cấp $\sigma_i(u)$, là các phần tử của $\Gamma(X,\sh{A})$;
ta lại có các hệ thức từ \sref{II.6.4.1.3} đến \sref{II.6.4.1.5}.
\oldpage[II]{124}
Như vậy ta đã xác định một đồng cấu $\det:\Hom_{\sh{A}}(\sh{E},\sh{E})\to\Gamma(X,\sh{A})$ của các monoid nhân;
vì theo định nghĩa $\Hom_{\sh{A}}(\sh{E},\sh{E})=\Gamma(X,\shHom_{\sh{A}}(\sh{E},\sh{E}))$, ta thấy trong định nghĩa này có thể thay $X$ bởi một tập mở tùy ý $U$ của $X$, từ đó xác định ngay một đồng cấu $\det:\shHom_{\sh{A}}(\sh{E},\sh{E})\to\sh{A}$ của các bó monoid nhân.
Khi $\sh{E}$ có hạng không đổi, ta định nghĩa theo cùng cách các đồng cấu $\sigma_i:\shHom_{\sh{A}}(\sh{E},\sh{E})\to\sh{A}$ của các bó tập hợp;
với $i=1$, đồng cấu $\sigma_1=\Tr$ là một đồng cấu các $\sh{A}$-môđun.

Cho $(Y,\sh{B})$ là một không gian vành hóa thứ hai, và $f:(X,\sh{A})\to(Y,\sh{B})$ là một cấu xạ các không gian vành hóa;
nếu $\sh{F}$ là một $\sh{B}$-môđun tự do địa phương, thì $f^*(\sh{F})$ là một $\sh{A}$-môđun tự do địa phương (có cùng hạng với $\sh{F}$ nếu môđun sau có hạng không đổi) \sref[0]{0.5.4.5}.
Với mọi tự đồng cấu $v$ của $\sh{F}$, $f^*(v)$ là một tự đồng cấu của $f^*(\sh{F})$, và từ các định nghĩa suy ra ngay rằng $\det f^*(v)$ là tiết diện của $\sh{A}=f^*(\sh{B})$ trên $X$ tương ứng một cách chính tắc với $\det v\in\Gamma(Y,\sh{B})$.
Cũng có thể nói rằng đồng cấu $f^*(\det):f^*(\shHom_{\sh{B}}(\sh{F},\sh{F}))\to f^*(\sh{B})=\sh{A}$ là hợp thành
\[
  f^*(\shHom_{\sh{B}}(\sh{F},\sh{F}))
  \xrightarrow{\gamma^\sharp} \shHom_{\sh{A}}(f^*(\sh{F}),f^*(\sh{F}))
  \xrightarrow{\det} \sh{A}
\]
\sref[0]{0.4.4.6}.
Ta có các kết quả tương tự đối với các $\sigma_i$.
\end{env}

\begin{env}[6.4.9]
\label{II.6.4.9}
Bây giờ giả sử $X$ là một \emph{tiền lược đồ nguyên địa phương}, sao cho bó $\sh{R}(X)$ các hàm hữu tỉ trên $X$ là một bó các trường đơn địa phương \sref[I]{I.7.3.4}, và là quasi-coherent với tư cách một $\sh{O}_X$-môđun.
Nếu $\sh{E}$ là một $\sh{O}_X$-môđun quasi-coherent \emph{hữu hạn sinh}, thì $\sh{E}'=\sh{E}\otimes_{\sh{O}_X}\sh{R}(X)$ là một $\sh{R}(X)$-môđun tự do địa phương \sref[I]{I.7.3.6};
với mọi tự đồng cấu $u$ của $\sh{E}$, $u\otimes1_{\sh{R}(X)}$ khi đó là một tự đồng cấu của $\sh{E}'$, và $\det(u\otimes1)$ là một tiết diện của $\sh{R}(X)$ trên $X$, vẫn được gọi là \emph{định thức} của $u$ và vẫn được ký hiệu là $\det u$.
Từ \sref{II.6.4.3} suy ra $\det u$ là một tiết diện của \emph{bao đóng nguyên} của $\sh{O}_X$ trong $\sh{R}(X)$ \sref{II.6.3.2};
nếu thêm rằng $X$ \emph{chuẩn tắc}, thì $\det u$ là một \emph{tiết diện của $\sh{O}_X$} trên $X$, và nếu còn giả sử $\sh{E}$ \emph{không xoắn}, thì theo \sref{II.6.4.6}, để $u$ là một tự đẳng cấu của $\sh{E}$, điều kiện cần và đủ là $\det u$ \emph{khả nghịch}.
Các công thức từ \sref{II.6.4.8.1} đến \sref{II.6.4.8.3} vẫn đúng;
từ đồng cấu $u\mapsto\det u$, áp dụng cho các môđun tiết diện của $\shHom_{\sh{O}_X}(\sh{E},\sh{E})$, ta suy ra một đồng cấu bó $\det:\shHom_{\sh{O}_X}(\sh{E},\sh{E})\to\sh{R}(X)$, nhận giá trị trong $\sh{O}_X$ khi $X$ chuẩn tắc.
Ta có các định nghĩa và kết quả tương tự cho các hàm đối xứng sơ cấp khác $\sigma_j(u)$, nếu $\sh{E}'$ có \emph{hạng không đổi};
nếu thêm rằng $X$ \emph{chuẩn tắc}, thì các $\sigma_j(u)$ là các \emph{tiết diện của $\sh{O}_X$} trên $X$.

Cuối cùng, cho $X$ và $Y$ là hai tiền lược đồ nguyên, và $f:X\to Y$ là một cấu xạ \emph{trội}.
Ta biết rằng tồn tại một đồng cấu chính tắc $f^*(\sh{R}(Y))\to\sh{R}(X)$ \sref[I]{I.7.3.8}\footnote{Xem \emph{Errata et Addenda}, danh sách 1.}, từ đó, với mọi $\sh{O}_Y$-môđun quasi-coherent hữu hạn sinh $\sh{F}$, suy ra một đồng cấu chính tắc $\theta:f^*(\sh{F}\otimes_{\sh{O}_Y}\sh{R}(Y))\to f^*(\sh{F})\otimes_{\sh{O}_X}\sh{R}(X)$.
\oldpage[II]{125}
Nếu $v$ là một tự đồng cấu của $\sh{F}$, thì $f^*(v\otimes1_{\sh{R}(Y)})$ là một tự đồng cấu của $f^*(\sh{F}\otimes_{\sh{O}_Y}\sh{R}(Y))$, và ta có một sơ đồ giao hoán
\[
  \xymatrix{
    f^*(\sh{F}\otimes_{\sh{O}_Y}\sh{R}(Y))
      \ar[r]^{f^*(v\otimes1)}
      \ar[d]_{\theta}
    & f^*(\sh{F}\otimes_{\sh{O}_Y}\sh{R}(Y))
      \ar[d]^{\theta}
  \\f^*(\sh{F})\otimes_{\sh{O}_X}\sh{R}(X)
      \ar[r]_{f^*(v)\otimes1}
    & f^*(\sh{F})\otimes_{\sh{O}_X}\sh{R}(X)
  }
\]

Từ đó dễ dàng suy ra rằng $\det f^*(v)$ là ảnh chính tắc, qua đồng cấu $f^*(\sh{R}(Y))\to\sh{R}(X)$, của tiết diện $\det v$ của $\sh{R}(Y)$;
thật vậy, ta quy ngay về trường hợp $X=\Spec(A)$ và $Y=\Spec(B)$ là affine, trong đó $A$ và $B$ là các miền nguyên có các trường phân thức tương ứng là $K$ và $L$, còn đồng cấu $B\to A$ là đơn ánh và do đó mở rộng thành một đơn cấu $L\to K$;
nếu $\sh{F}=\widetilde{M}$, trong đó $M$ là một $B$-môđun hữu hạn sinh, thì hạng của $M\otimes_B L$ trên $L$ \emph{bằng} hạng của $(M\otimes_B A)\otimes_A K$ trên $K$, và $\det((u\otimes1)\otimes1)$ là ảnh trong $K$ của $\det(u\otimes1)$ đối với mọi tự đồng cấu $u$ của $M$, từ đó suy ra kết luận.
\end{env}

\begin{env}[6.4.10]
\label{II.6.4.10}
Sau cùng, giả sử $X$ là một \emph{tiền lược đồ Noether địa phương rút gọn}, sao cho bó $\sh{R}(X)$ các hàm hữu tỉ trên $X$ vẫn là một $\sh{O}_X$-môđun quasi-coherent \sref[I]{I.7.3.4};
mặt khác, cho $\sh{E}$ là một $\sh{O}_X$-môđun \emph{coherent} sao cho $\sh{E}'=\sh{E}\otimes_{\sh{O}_X}\sh{R}(X)$ \emph{tự do địa phương} (có hạng hữu hạn).
Theo \sref{II.6.4.7}, nếu $\sh{E}'$ có hạng không đổi, thì với mọi tự đồng cấu $u$ của $\sh{E}$, ta có thể xác định các $\sigma_j(u)$, là những tiết diện của $\sh{R}(X)$ trên $X$.
Khi không còn giả sử $\sh{E}'$ có hạng không đổi, ta vẫn có thể xác định đồng cấu $\det:\shHom_{\sh{O}_X}(\sh{E},\sh{E})\to\sh{R}(X)$.
\end{env}

\subsection{Chuẩn của một bó khả nghịch}
\label{subsection:II.6.5}

\begin{env}[6.5.1]
\label{II.6.5.1}
Cho $(X,\sh{A})$ là một không gian vành hóa và $\sh{B}$ là một $\sh{A}$-đại số (giao hoán).
$\sh{A}$-môđun $\sh{B}$ được đồng nhất chính tắc với một $\sh{A}$-môđun con của $\shHom_{\sh{A}}(\sh{B},\sh{B})$, trong đó một tiết diện $f$ của $\sh{B}$ trên một tập mở $U$ của $X$ được đồng nhất với phép nhân bởi tiết diện ấy.
Nếu $(X,\sh{A})$ và $\sh{B}$ thỏa một trong các điều kiện liệt kê trong \sref{II.6.4.8}, \sref{II.6.4.9} hoặc \sref{II.6.4.10}, thì ta có thể xác định $\det(f)$ (và, trong một số trường hợp, các $\sigma_j(f)$), là những tiết diện của $\sh{A}$ hoặc của $\sh{R}(X)$ trên $U$; ta gọi chúng lần lượt là \emph{chuẩn} của $f$ và \emph{các hàm đối xứng sơ cấp} của $f$, và ký hiệu chuẩn là $N_{\sh{B}/\sh{A}}(f)$.
Ta sẽ giả sử \emph{một} trong các điều kiện sau được thỏa:
\begin{enumerate}
  \item[{\rm(I)}] \emph{$\sh{B}$ là một $\sh{A}$-môđun tự do địa phương (có hạng hữu hạn).}
  \item[{\rm(II)}] \emph{$(X,\sh{A})$ là một tiền lược đồ Noether địa phương rút gọn, $\sh{B}$ là một $\sh{A}$-môđun coherent sao cho $\sh{B}\otimes_{\sh{A}}\sh{R}(X)$ là một $\sh{R}(X)$-môđun tự do địa phương, và với mọi tiết diện $f\in\Gamma(U,\sh{B})$ trên một tập mở $U$ của $X$, chuẩn $N_{\sh{B}/\sh{A}}(f)$ là một tiết diện của $\sh{A}$ trên $U$.}
\end{enumerate}
\oldpage[II]{126}
Ta lưu ý rằng điều kiện cuối này tự động được thỏa khi tiền lược đồ Noether địa phương $X$ là \emph{chuẩn tắc} \sref{II.6.4.9}.

Mặt khác, giả thiết rằng $\sh{B}\otimes_{\sh{A}}\sh{R}(X)$ tự do địa phương còn có thể được diễn đạt như sau: ký hiệu bởi $X_\alpha$ các tiền lược đồ con đóng rút gọn của $X$ có không gian nền là các thành phần bất khả quy của $X$ \sref[I]{I.5.2.1}; do đó, chúng là các tiền lược đồ nguyên Noether địa phương.
Mỗi $x\in X$ thuộc một số hữu hạn các không gian con $X_\alpha$;
mặt khác, $\sh{B}\otimes_{\sh{A}}\sh{R}(X_\alpha)$ là một $\sh{R}(X_\alpha)$-môđun tự do địa phương có hạng không đổi~$k_\alpha$ \sref[I]{I.7.3.6};
nói rằng $\sh{B}\otimes_{\sh{A}}\sh{R}(X)$ là một $\sh{R}(X)$-môđun tự do địa phương có nghĩa là, với mọi $x\in X$, \emph{các hạng $k_\alpha$ tương ứng với những chỉ số thỏa $x\in X_\alpha$ đều bằng nhau}.
Quả vậy, vấn đề mang tính địa phương, và ta quy về trường hợp $X=\Spec(C)$, trong đó $C$ là một vành Noether rút gọn, còn $\sh{B}=\widetilde{D}$, trong đó $D$ là một $C$-đại số đồng thời là một $C$-môđun hữu hạn sinh;
nếu $\mathfrak{p}_i$ ($1\leq i\leq m$) là các iđêan nguyên tố tối tiểu của $C$, thì vành phân thức toàn phần $L$ của $C$ là tổng trực tiếp của các trường phân thức $K_i$ của các miền nguyên $C/\mathfrak{p}_i$, và $D\otimes_C L$ là tổng trực tiếp của các $D\otimes_C K_i$, từ đó suy ra kết luận.

Hiển nhiên, dưới giả thiết~(I) hoặc (II), bằng cách này ta xác định được một \emph{đồng cấu giữa các bó monoid nhân} $N_{\sh{B}/\sh{A}}:\sh{B}\to\sh{A}$, cũng được ký hiệu là $N$ nếu không gây nhầm lẫn, và gọi là đồng cấu \emph{chuẩn}.
Với hai tiết diện $f$, $g$ của $\sh{B}$ trên cùng một tập mở $U$, ta có
\[
\label{II.6.5.1.1}
  N_{\sh{B}/\sh{A}}(fg) = N_{\sh{B}/\sh{A}}(f)N_{\sh{B}/\sh{A}}(g)
\tag{6.5.1.1}
\]
đối với các tiết diện tương ứng của $\sh{A}$ trên $U$;
\[
\label{II.6.5.1.2}
  N_{\sh{B}/\sh{A}}(1_{\sh{B}}) = 1_{\sh{A}};
\tag{6.5.1.2}
\]
sau cùng, với mọi tiết diện $s$ của $\sh{A}$ trên $U$,
\[
\label{II.6.5.1.3}
  N_{\sh{B}/\sh{A}}(s\cdot1_{\sh{B}}) = s^n
\tag{6.5.1.3}
\]
nếu hạng của $\sh{B}$ không đổi và bằng $n$ (với $s=0_{\sh{A}}$, công thức này cho $N(0_{\sh{B}})=0_{\sh{A}}$ nếu $n\geq1$, và $N(0_{\sh{B}})=N(1_{\sh{B}})=1_{\sh{A}}$ nếu $n=0$).

Dưới giả thiết~(I), điều kiện cần và đủ để $f\in\Gamma(U,\sh{B})$ khả nghịch là $N(f)\in\Gamma(U,\sh{A})$ khả nghịch.
Dưới giả thiết~(II), điều kiện này vẫn cần;
nó là đủ (theo \sref{II.6.4.7}) khi giả sử $\sh{B}\to\sh{B}\otimes_{\sh{A}}\sh{R}(X)$ là \emph{đơn ánh} và điều kiện chặt hơn sau đây được thỏa:
\begin{enumerate}
  \item[{\rm(II~\emph{bis})}] \emph{$(X,\sh{A})$ là một tiền lược đồ Noether địa phương rút gọn, $\sh{B}$ là một $\sh{A}$-môđun coherent sao cho $\sh{B}\otimes_{\sh{A}}\sh{R}(X)$ là một $\sh{R}(X)$-môđun tự do địa phương, và với mọi tiết diện $f\in\Gamma(U,\sh{B})$ trên một tập mở $U$ sao cho $\sh{B}\otimes_{\sh{A}}\sh{R}(X)|U$ có hạng không đổi~$n$ trên $\sh{R}(X)|U$, các $\sigma_j(f)$ ($1\leq j\leq n$) là những tiết diện của $\sh{A}$ trên $U$.}
\end{enumerate}

(Ta lại lưu ý rằng điều kiện này được thỏa khi $X$ là \emph{chuẩn tắc}.)
\end{env}

\begin{env}[6.5.2]
\label{II.6.5.2}
Giả sử một trong các giả thiết~(I), (II) của \sref{II.6.5.1} được thỏa, và cho $\sh{L}'$ là một $\sh{B}$-môđun \emph{khả nghịch}.
Ta sẽ liên kết chính tắc với nó (sai khác một đẳng cấu duy nhất) một $\sh{A}$-môđun \emph{khả nghịch} $\sh{L}$ theo cách sau.
Ký hiệu bởi $\sh{A}^*$ (tương ứng, $\sh{B}^*$) bó con của $\sh{A}$ (tương ứng, của $\sh{B}$) sao cho $\Gamma(U,\sh{A}^*)$ (tương ứng, $\Gamma(U,\sh{B}^*)$) là tập hợp các
\oldpage[II]{127}
phần tử khả nghịch của $\Gamma(U,\sh{A})$ (tương ứng, của $\Gamma(U,\sh{B})$) với mọi tập mở $U\subset X$;
đó là các bó nhóm nhân, và hạn chế của $N_{\sh{B}/\sh{A}}$ lên $\sh{B}^*$ là một \emph{đồng cấu} $\sh{B}^*\to\sh{A}^*$ của các bó nhóm \sref{II.6.5.1}.
Cho $\mathfrak{L}$ là tập hợp các cặp $(U_\lambda,\nu_\lambda)$ có tính chất sau: $U_\lambda$ là một tập mở của $X$, và $\nu_\lambda$ là một đẳng cấu $\sh{L}'|U_\lambda\simto\sh{B}|U_\lambda$ của các $(\sh{B}|U_\lambda)$-môđun.
Theo giả thiết, các $U_\lambda$ lập thành một phủ của $X$;
với hai chỉ số bất kỳ $\lambda$ và $\mu$, đặt $\omega_{\lambda\mu}=(\nu_\lambda|U_\lambda\cap U_\mu)\circ(\nu_\mu|U_\lambda\cap U_\mu)^{-1}$; đây là một tự đẳng cấu của $\sh{B}|U_\lambda\cap U_\mu$, được đồng nhất chính tắc với một tiết diện của $\sh{B}^*$ trên $U_\lambda\cap U_\mu$, và $(\omega_{\lambda\mu})$ là một $1$-đối chu trình của phủ $\mathfrak{U}=(U_\lambda)$ nhận giá trị trong $\sh{B}^*$ \sref[0]{0.5.4.7}.
Việc $N_{\sh{B}/\sh{A}}:\sh{B}^*\to\sh{A}^*$ là một đồng cấu kéo theo rằng $(N_{\sh{B}/\sh{A}}\omega_{\lambda\mu})$ là một $1$-đối chu trình của $\mathfrak{U}$ nhận giá trị trong $\sh{A}^*$; vì vậy, tương ứng với nó (sai khác một đẳng cấu duy nhất) là một $\sh{A}$-môđun khả nghịch $\sh{L}$ \sref[0]{0.5.4.7}, mà ta ký hiệu là $N_{\sh{B}/\sh{A}}(\sh{L}')$ và gọi là \emph{chuẩn} của $\sh{B}$-môđun khả nghịch $\sh{L}'$.

Cho $\mathfrak{M}$ là một tập con của $\mathfrak{L}$ sao cho các $U_\lambda$ tương ứng vẫn lập thành một phủ của $X$, và ký hiệu phủ này là $\mathfrak{B}$;
hạn chế của đối chu trình $(\omega_{\lambda\mu})$ lên $\mathfrak{B}$ xác định một $1$-đối chu trình $(N_{\sh{B}/\sh{A}}\omega_{\lambda\mu})$ của $\mathfrak{B}$ nhận giá trị trong $\sh{A}^*$, là hạn chế của $1$-đối chu trình $(N_{\sh{B}/\sh{A}}\omega_{\lambda\mu})$ của $\mathfrak{U}$;
hiển nhiên có một đẳng cấu chính tắc từ $\sh{A}$-môđun khả nghịch được xác định bởi $1$-đối chu trình này của $\mathfrak{B}$ đến $N_{\sh{B}/\sh{A}}(\sh{L}')$, nhờ đó có thể xác định $\sh{A}$-môđun khả nghịch này bằng một phủ con bất kỳ của $\mathfrak{U}$.
Khả năng này cho thấy ngay rằng, nếu $\sh{L}'_1$ và $\sh{L}'_2$ là hai $\sh{B}$-môđun khả nghịch, thì theo \sref{II.6.5.1.1} và \sref{II.6.5.1.2},
\[
\label{II.6.5.2.1}
  N(\sh{L}'_1\otimes_{\sh{B}}\sh{L}'_2) = N(\sh{L}'_1)\otimes_{\sh{A}}N(\sh{L}'_2)
\tag{6.5.2.1}
\]
và
\[
\label{II.6.5.2.2}
  N_{\sh{B}/\sh{A}}(\sh{B}) = \sh{A}
\tag{6.5.2.2}
\]
cũng như
\[
\label{II.6.5.2.3}
  N({\sh{L}'}^{-1}) = (N(\sh{L}'))^{-1}
\tag{6.5.2.3}
\]
sai khác các đẳng cấu chính tắc.
Hơn nữa, từ \sref{II.6.5.1.3} suy ra rằng, nếu $\sh{L}$ là một $\sh{A}$-môđun khả nghịch và nếu $\sh{B}$ có hạng không đổi~$n$ trên $\sh{A}$ trong trường hợp~(I) (tương ứng, $\sh{B}\otimes_{\sh{A}}\sh{R}(X)$ có hạng không đổi~$n$ trên $\sh{R}(X)$ trong trường hợp~(II)), thì, sai khác một đẳng cấu chính tắc,
\[
\label{II.6.5.2.4}
  N_{\sh{B}/\sh{A}}(\sh{L}\otimes_{\sh{A}}\sh{B}) = \sh{L}^{\otimes n}.
\tag{6.5.2.4}
\]
\end{env}

\begin{env}[6.5.3]
\label{II.6.5.3}
Giờ ta hãy chứng minh rằng $N_{\sh{B}/\sh{A}}$ là một \emph{hàm tử hiệp biến} trên phạm trù các $\sh{B}$-môđun khả nghịch.
Thật vậy, cho $h':\sh{L}'_1\to\sh{L}'_2$ là một đồng cấu các $\sh{B}$-môđun khả nghịch, và cho $\mathfrak{B}=(U_\lambda)$ là một phủ mở của $X$ sao cho, với mọi $\lambda$, có các đẳng cấu $(\sh{B}|U_\lambda)$-môđun $\eta_\lambda^{(1)}:\sh{L}'_1|U_\lambda\simto\sh{B}|U_\lambda$ và $\eta_\lambda^{(2)}:\sh{L}'_2|U_\lambda\simto\sh{B}|U_\lambda$;
do đó, với mỗi $\lambda$, có một tự đồng cấu $h'_\lambda$ của $\sh{B}|U_\lambda$ sao cho $h'_\lambda\circ\eta_\lambda^{(1)}=\eta_\lambda^{(2)}\circ(h'|U_\lambda)$, và hiển nhiên có thể đồng nhất $h'_\lambda$ với một tiết diện của $\sh{B}$ trên $U_\lambda$ \sref[0]{0.5.1.1}.
Vì vậy, với mọi cặp chỉ số $(\lambda,\mu)$, các hạn chế lên $U_\lambda\cap U_\mu$ của $(\eta_\lambda^{(2)})^{-1}\circ h'_\lambda\circ\eta_\lambda^{(1)}$ và của $(\eta_\mu^{(2)})^{-1}\circ h'_\mu\circ\eta_\mu^{(1)}$ trùng nhau.
Từ đó, đối với các $1$-đối chu trình $(\omega_{\lambda\mu}^{(1)})$ và $(\omega_{\lambda\mu}^{(2)})$ nhận giá trị trong $\sh{B}^*$ tương ứng với $\sh{L}'_1$ và $\sh{L}'_2$, ta suy ra hệ thức
\[
  \omega_{\lambda\mu}^{(2)}h'_\mu = h'_\lambda\omega_{\lambda\mu}^{(1)}.
\]
\oldpage[II]{128}
Nếu đặt $h_\lambda=N(h'_\lambda)$, ta sẽ có các hệ thức tương tự
\[
  N(\omega_{\lambda\mu}^{(2)})h_\mu = h_\lambda N(\omega_{\lambda\mu}^{(1)})
\]
và do đó các $h_\lambda$ xác định một đồng cấu $N(\sh{L}'_1)\to N(\sh{L}'_2)$, mà ta ký hiệu là $N_{\sh{B}/\sh{A}}(h')$ hoặc $N(h')$.
Dưới giả thiết~(I), điều kiện cần và đủ để $h'$ là một \emph{đẳng cấu} là $N_{\sh{B}/\sh{A}}(h')$ là một đẳng cấu (vì đây là một vấn đề địa phương).
Dưới giả thiết~(II), điều kiện này vẫn cần;
nó là đủ nếu giả thiết~(II~\emph{bis}) được thỏa và $\sh{B}\to\sh{B}\otimes_{\sh{A}}\sh{R}(X)$ là đơn ánh.

Đặc biệt, lấy $\sh{L}'_1=\sh{B}$;
khi đó, các đồng cấu $\sh{B}\to\sh{L}'$ được đồng nhất \sref[0]{0.5.1.1} với các tiết diện của $\sh{L}'$ trên $X$, từ đó có một ánh xạ chính tắc
\[
  N_{\sh{B}/\sh{A}} : \Gamma(X,\sh{L}')\to\Gamma(X,N_{\sh{B}/\sh{A}}(\sh{L}')).
\]

Từ \sref{II.6.5.1.1} còn suy ra rằng, nếu $f'_1\in\Gamma(X,\sh{L}'_1)$ và $f'_2\in\Gamma(X,\sh{L}'_2)$, thì
\[
\label{II.6.5.3.1}
  N(f'_1\otimes f'_2) = N(f'_1)\otimes N(f'_2).
\tag{6.5.3.1}
\]

Với mọi $\sh{A}$-môđun khả nghịch $\sh{L}$ và mọi tiết diện $f\in\Gamma(X,\sh{L})$, khi xét đến \sref{II.6.5.2.4}, ta có
\[
\label{II.6.5.3.2}
  N_{\sh{B}/\sh{A}}(f\otimes 1_{\sh{B}}) = f^{\otimes n}
\tag{6.5.3.2}
\]
khi $\sh{B}$ có hạng không đổi~$n$ dưới giả thiết~(I) (tương ứng, khi $\sh{B}\otimes_{\sh{A}}\sh{R}(X)$ có hạng không đổi~$n$ dưới giả thiết~(II)).
Sau cùng, để đồng cấu $\sh{B}\to\sh{L}'$ tương ứng với một tiết diện $f'$ của $\sh{L}'$ trên $X$ là một đẳng cấu, điều kiện cần và đủ là $f'_x$ là một cơ sở của $\sh{L}'_x$ với mọi $x\in X$;
dưới giả thiết~(I), điều kiện này do đó tương đương với việc $(N(f'))_x$ là một cơ sở của $(N(\sh{L}'))_x$ với mọi $x$;
dưới giả thiết~(II), điều kiện này vẫn cần, và là đủ khi $\sh{B}$ thỏa giả thiết~(II~\emph{bis}) và $\sh{B}\to\sh{B}\otimes_{\sh{A}}\sh{R}(X)$ là đơn ánh.
\end{env}

\begin{env}[6.5.4]
\label{II.6.5.4}
Cho $(X,\sh{A})$ và $(X',\sh{A}')$ là hai không gian vành hóa, $f:X'\to X$ là một cấu xạ, $\sh{B}$ là một $\sh{A}$-đại số, và $\sh{B}'=f^*(\sh{B})$ là $\sh{A}'$-đại số ảnh ngược.
Giả sử một trong các giả thiết sau được thỏa:
\begin{enumerate}
  \item $\sh{B}$ thỏa giả thiết~(I) của \sref{II.6.5.1}.
  \item $(X,\sh{A})$ và $\sh{B}$ thỏa giả thiết~(II) của \sref{II.6.5.1}, $(X',\sh{A}')$ là một tiền lược đồ Noether địa phương rút gọn, và nếu ký hiệu bởi $X_\alpha$ và $X'_\beta$ các tiền lược đồ con đóng rút gọn tương ứng của $X$ và $X'$ có không gian nền là các thành phần bất khả quy của những không gian này, thì hạn chế của $f$ lên mỗi $X'_\beta$ là một cấu xạ \emph{trội} từ $X'_\beta$ vào một $X_\alpha$.
\end{enumerate}

Dưới các điều kiện này, $\sh{B}'$ lần lượt thỏa giả thiết~(I) hoặc giả thiết~(II) của \sref{II.6.5.1}; khẳng định thứ nhất là hiển nhiên.
Để chứng minh khẳng định thứ hai, chỉ cần thấy rằng, với mọi $x'\in X'$, hạng của các $\sh{B}'\otimes_{\sh{O}_{X'}}\sh{R}(X'_\beta)$ \emph{bằng nhau} đối với những $\beta$ thỏa $x'\in X'_\beta$.
Thật vậy, nếu hạn chế của $f$ lên $X'_\beta$ là một cấu xạ trội vào $X_\alpha$, thì hạng của $\sh{B}'\otimes_{\sh{O}_{X'}}\sh{R}(X'_\beta)$ bằng hạng của $\sh{B}\otimes_{\sh{O}_X}\sh{R}(X_\alpha)$ (điều này thấy ngay bằng cách quy về
\oldpage[II]{129}
trường hợp affine, như trong \sref{II.6.4.9}), từ đó suy ra khẳng định của ta nhờ giả thiết~(II) và \sref{II.6.5.1}.

Với những điều này, từ \sref{II.6.4.8}, \sref{II.6.4.9} và \sref{II.6.4.10} suy ra rằng, nếu $s$ là một tiết diện của $\sh{B}$ trên một tập mở $U\subset X$, còn $s'$ là tiết diện tương ứng của $\sh{B}'$ trên $f^{-1}(U)$, thì $N_{\sh{B}'/\sh{A}'}(s')$ là tiết diện của $\sh{A}'$ trên $f^{-1}(U)$ tương ứng với tiết diện $N_{\sh{B}/\sh{A}}(s)$ của $\sh{A}$ trên $U$.

Nếu $\sh{M}$ là một $\sh{B}$-môđun khả nghịch, thì từ những điều trên suy ra rằng, nếu $\sh{M}'=f^*(\sh{M})$ (là một $\sh{B}'$-môđun khả nghịch), ta có $N_{\sh{B}'/\sh{A}'}(\sh{M}')=f^*(N_{\sh{B}/\sh{A}}(\sh{M}))$, sai khác một đẳng cấu chính tắc.
\end{env}

\begin{env}[6.5.5]
\label{II.6.5.5}
Từ đây, giả sử $(X,\sh{A})$ là một \emph{tiền lược đồ}.
Việc cho một $\sh{A}$-đại số quasi-coherent $\sh{B}$, đồng thời là một \emph{$\sh{A}$-môđun hữu hạn sinh}, khi đó tương đương, như đã biết, với việc cho một cấu xạ \emph{hữu hạn} $g:X'\to X$ sao cho $g_*(\sh{O}_{X'})=\sh{B}$, được xác định sai khác một $X$-đẳng cấu (\sref{II.6.1.2} và \sref{II.1.3.1}).
Hơn nữa, việc cho một $\sh{O}_{X'}$-môđun quasi-coherent $\sh{F}'$ tương đương với việc cho một $\sh{B}$-môđun quasi-coherent $\sh{F}$ sao cho $g_*(\sh{F}')=\sh{F}$ \sref{II.1.4.3}, và điều kiện cần và đủ để $\sh{F}'$ khả nghịch là $\sh{F}$ khả nghịch \sref{II.6.1.12}.
Do đó, để chuyển các kết quả trên sang ngôn ngữ của cấu xạ hữu hạn $g$, ta phải giả sử hoặc $g_*(\sh{O}_{X'})$ là một $\sh{O}_X$-môđun \emph{tự do địa phương} (hữu hạn sinh), hoặc $(X,\sh{O}_X)$ và $g_*(\sh{O}_{X'})$ thỏa giả thiết~(II).
Với mọi $\sh{O}_{X'}$-môđun khả nghịch $\sh{L}'$, khi đó ta đặt
\[
\label{II.6.5.5.1}
  N_{X'/X}(\sh{L}') = N_{g_*(\sh{O}_{X'})/\sh{O}_X}(g_*(\sh{L}'))
\tag{6.5.5.1}
\]
và gọi đó là \emph{chuẩn} (đối với $g$) của $\sh{L}'$.
Tương tự, nếu $h':\sh{L}'_1\to\sh{L}'_2$ là một đồng cấu các $\sh{O}_{X'}$-môđun khả nghịch, ta đặt
\[
\label{II.6.5.5.2}
  N_{X'/X}(h') = N_{g_*(\sh{O}_{X'})/\sh{O}_X}(g_*(h')) : N_{X'/X}(\sh{L}'_1)\to N_{X'/X}(\sh{L}'_2).
\tag{6.5.5.2}
\]

Đặc biệt, với $\sh{L}'_1=\sh{O}_{X'}$, ta thu được một ánh xạ chính tắc
\[
\label{II.6.5.5.3}
  N_{X'/X} : \Gamma(X',\sh{L}') \to \Gamma(X,N_{X'/X}(\sh{L}')).
\tag{6.5.5.3}
\]

Ta sẽ để lại cho người đọc phần lớn các phép chuyển ngôn ngữ này và chỉ trình bày tường minh các kết quả sau:
\end{env}

\begin{proposition}[6.5.6]
\label{II.6.5.6}
Cho $g:X'\to X$ là một cấu xạ hữu hạn, và giả sử hoặc $g_*(\sh{O}_{X'})$ là một $\sh{O}_X$-môđun tự do địa phương, hoặc $(X,\sh{O}_X)$ và $g_*(\sh{O}_{X'})$ thỏa (II~\emph{bis}) (điều này đặc biệt xảy ra khi $X$ là Noether địa phương và chuẩn tắc).
Để một đồng cấu $h':\sh{L}'_1\to\sh{L}'_2$ của các $\sh{O}_{X'}$-môđun khả nghịch là một đẳng cấu, trong giả thiết thứ nhất, điều kiện cần và đủ là $N_{X'/X}(h')$ là một đẳng cấu;
trong giả thiết thứ hai, điều kiện này vẫn cần, và là đủ khi đồng cấu $g_*(\sh{O}_{X'})\to g_*(\sh{O}_{X'})\otimes_{\sh{O}_X}\sh{R}(X)$ là đơn ánh.
\end{proposition}

\begin{proof}
Ta lưu ý rằng ở đây đã dùng sự kiện: điều kiện cần và đủ để $g_*(h')$ là một đẳng cấu là $h'$ là một đẳng cấu \sref{II.1.4.2}.
\end{proof}

\begin{corollary}[6.5.7]
\label{II.6.5.7}
Cho $g:X'\to X$ là một cấu xạ hữu hạn, và giả sử hoặc $g_*(\sh{O}_{X'})$ là một $\sh{O}_X$-môđun tự do địa phương, hoặc $(X,\sh{O}_X)$ và $g_*(\sh{O}_{X'})$ thỏa (II~\emph{bis}) và $g_*(\sh{O}_{X'})\to g_*(\sh{O}_{X'})\otimes_{\sh{O}_X}\sh{R}(X)$ là đơn ánh.
Cho $\sh{L}'$ là một $\sh{O}_{X'}$-môđun khả nghịch, $f'$ là một tiết diện của $\sh{L}'$ trên $X'$, và $f=N_{X'/X}(f')$ là tiết diện của $\sh{L}=N_{X'/X}(\sh{L}')$ trên $X$ mà
\oldpage[II]{130}
tương ứng với nó \sref{II.6.5.5.3}.
Khi đó $g(X'\setminus X'_{f'})=X\setminus X_f$, và $X_f$ là tập mở lớn nhất $U$ của $X$ sao cho $g^{-1}(U)\subset X'_{f'}$.
\end{corollary}

\begin{proof}
Thật vậy, $g(X'\setminus X'_{f'})$ là đóng trong $X$ \sref{II.6.1.10};
do đó chỉ cần chứng minh khẳng định cuối.
Mà quan hệ $U\subset X_f$ tương đương với việc đồng cấu $\sh{O}_X|U\to\sh{L}|U$ được xác định bởi $f|U$ là một đẳng cấu.
Theo \sref{II.6.5.6}, điều này tương đương với việc đồng cấu $\sh{O}_{X'}|g^{-1}(U)\to\sh{L}'|g^{-1}(U)$ được xác định bởi $f'|g^{-1}(U)$ là một đẳng cấu, tức là tương đương với quan hệ $g^{-1}(U)\subset X'_{f'}$.
\end{proof}

\begin{proposition}[6.5.8]
\label{II.6.5.8}
Cho $g:X'\to X$ là một cấu xạ hữu hạn và $f:Y\to X$ là một cấu xạ;
gọi $Y'=X'_{(Y)}$, $g'=g_{(Y)}$ và $f'=f_{(X')}$, sao cho sơ đồ
\[
  \xymatrix{
    X' \ar[d]_{g}
    & Y' \ar[l]_{f'} \ar[d]^{g'}
  \\X
    & Y \ar[l]^{f}
  }
\]
giao hoán.

Giả sử $g_*(\sh{O}_{X'})$ là một $\sh{O}_X$-môđun tự do địa phương, hoặc $(X,\sh{O}_X)$ và $g_*(\sh{O}_{X'})$ thỏa mãn (II).
Giả sử thêm rằng $Y$ là một tiền lược đồ Noether địa phương và rút gọn, và hạn chế của $f$ lên mọi thành phần bất khả quy của $Y$ là một cấu xạ trội vào một thành phần bất khả quy của $X$.
Khi đó, với mọi $\sh{O}_{X'}$-môđun khả nghịch $\sh{L}'$, ta có
\[
  N_{Y'/Y}({f'}^*(\sh{L}')) = f^*(N_{X'/X}(\sh{L}'))
\]
sai khác một đẳng cấu chính tắc.
\end{proposition}

\begin{proof}
Lưu ý rằng $f^*(g_*(\sh{L}'))=g'_*({f'}^*(\sh{L}'))$ theo \sref{II.1.5.2}, và đặc biệt $g'_*(\sh{O}_{Y'})=f^*(g_*(\sh{O}_{X'}))$;
nếu $g_*(\sh{O}_{X'})$ tự do địa phương thì $g'_*(\sh{O}_{Y'})$ cũng vậy.
Kết luận khi đó suy ra từ các định nghĩa và \sref{II.6.5.4}.
\end{proof}

\begin{remark}[6.5.9]
\label{II.6.5.9}
Về sau, ta sẽ tổng quát hóa khái niệm chuẩn được phát triển ở trên và liên hệ nó với khái niệm ảnh trực tiếp của một ước.
\end{remark}


\subsection{Ứng dụng: các tiêu chuẩn về tính ample}
\label{subsection:II.6.6}

\begin{proposition}[6.6.1]
\label{II.6.6.1}
Cho $Y$ là một tiền lược đồ, $f:X\to Y$ là một cấu xạ quasi-compact và $g:X'\to X$ là một cấu xạ hữu hạn và toàn ánh.
Giả sử $g_*(\sh{O}_{X'})$ là một $\sh{O}_X$-môđun tự do địa phương, hoặc $(X,\sh{O}_X)$ và $g_*(\sh{O}_{X'})$ thỏa mãn (II~\emph{bis}).
Khi đó, với mọi $\sh{O}_{X'}$-môđun khả nghịch $\sh{L}'$ ample đối với $f\circ g$, $\sh{L}=N_{X'/X}(\sh{L}')$ ample đối với $f$.
\end{proposition}

\begin{proof}
Ta có thể giả sử $Y$ affine \sref{II.4.6.4}, và khi ấy, theo \sref{II.4.6.6}, mệnh đề tương đương với hệ quả sau:
\end{proof}

\begin{corollary}[6.6.2]
\label{II.6.6.2}
Cho $X$ là một tiền lược đồ quasi-compact và $g:X'\to X$ là một cấu xạ hữu hạn toàn ánh, sao cho $g_*(\sh{O}_{X'})$ là một $\sh{O}_X$-môđun tự do địa phương, hoặc $(X,\sh{O}_X)$ và $g_*(\sh{O}_{X'})$ thỏa mãn (II~\emph{bis}).
Khi đó, với mọi $\sh{O}_{X'}$-môđun ample $\sh{L}'$, $\sh{L}=N_{X'/X}(\sh{L}')$ là ample.
\end{corollary}

\begin{proof}
Trong giả thiết thứ hai, ta còn có thể giả sử đồng cấu chính tắc $g_*(\sh{O}_{X'})\to g_*(\sh{O}_{X'})\otimes_{\sh{O}_X}\sh{R}(X)$ là \emph{đơn ánh}.
Thật vậy, trong trường hợp ngược lại, cho $\sh{T}$
\oldpage[II]{131}
là hạt nhân của đồng cấu này; đó là một iđêan coherent của $\sh{B}=g_*(\sh{O}_{X'})$ \sref[I]{I.6.1.1}; đặt $X''=\Spec(\sh{B}/\sh{T})$;
do đó ta có sơ đồ giao hoán
\[
  \xymatrix{
    X'' \ar[rr]^{h} \ar[dr]_{g'}
    && X' \ar[dl]^{g}
  \\&X
  }
\]
trong đó $h$ là một phép nhúng đóng \sref{II.1.4.10}.
Hơn nữa, ta biết rằng giá của $\sh{T}$ là một tập đóng \sref[0]{0.5.2.2} không trù mật trong $X$ \sref[I]{I.7.4.6}; từ đó suy ra rằng, với điểm tổng quát $x$ của một thành phần bất khả quy của $X$, tồn tại một lân cận mở affine $U$ của $x$ sao cho $\sh{B}|U=(\sh{B}/\sh{T})|U$.
Vì theo giả thiết $g$ toàn ánh, ta suy ra $x\in g'(X'')$;
do đó $g'$ trội, và vì là một cấu xạ hữu hạn nên nó \emph{toàn ánh} \sref{II.6.1.10};
theo định nghĩa,
\[
  g'_*(\sh{O}_{X''})\otimes_{\sh{O}_X}\sh{R}(X)
  = (\sh{B}/\sh{T})\otimes_{\sh{O}_X}\sh{R}(X)
  = g_*(\sh{O}_{X'})\otimes_{\sh{O}_X}\sh{R}(X)
\]
do đó $(X,\sh{O}_X)$ và $g'_*(\sh{O}_{X''})$ thỏa mãn (II~\emph{bis}), và hơn nữa $g'_*(\sh{O}_{X''})\to g'_*(\sh{O}_{X''})\otimes_{\sh{O}_X}\sh{R}(X)$ là đơn ánh.
Cuối cùng, $h^*(\sh{L}')=\sh{L}''$ là một $\sh{O}_{X''}$-môđun ample \sref{II.4.6.13}[(i~\emph{bis})], và ta có $N_{X''/X}(\sh{L}'')=N_{X'/X}(\sh{L}')$.
Thật vậy, để xác định hai $\sh{O}_X$-môđun khả nghịch này, ta có thể dùng cùng một phủ mở affine $(U_\lambda)$ của $X$, sao cho các hạn chế của $g_*(\sh{L}')$ và $g'_*(\sh{L}'')$ lên $U_\lambda$ lần lượt đẳng cấu với $\sh{B}|U_\lambda$ và $(\sh{B}/\sh{T})|U_\lambda$.
Ta thấy ngay rằng với mỗi đẳng cấu $\eta_\lambda:g_*(\sh{L}')|U_\lambda\to\sh{B}|U_\lambda$ tương ứng một cách chính tắc một đẳng cấu
\[
  \eta'_\lambda : g'_*(\sh{L}'')|U_\lambda \to (\sh{B}/\sh{T})|U_\lambda
\]
sao cho, nếu $(\omega_{\lambda\mu})$ và $(\omega'_{\lambda\mu})$ là các $1$-đối chu trình tương ứng với các hệ đẳng cấu $(\eta_\lambda)$ và $(\eta'_\lambda)$ \sref{II.6.5.2}, thì $\omega'_{\lambda\mu}$ là ảnh chính tắc trong $\Gamma(U_\lambda\cap U_\mu,\sh{B}/\sh{T})$ của $\omega_{\lambda\mu}\in\Gamma(U_\lambda\cap U_\mu,\sh{B})$.
Theo định nghĩa của $\sh{T}$, do đó ta suy ra
\[
  N_{\sh{B}/\sh{A}}(\omega_{\lambda\mu}) = N_{(\sh{B}/\sh{T})/\sh{A}}(\omega'_{\lambda\mu})
\]
(trong đó $\sh{A}=\sh{O}_X$), và từ đó có đẳng thức đã nêu.

Vậy giả sử đồng cấu $g_*(\sh{O}_{X'})\to g_*(\sh{O}_{X'})\otimes_{\sh{O}_X}\sh{R}(X)$ là đơn ánh khi ta ở trong giả thiết~(II~\emph{bis}).
Chỉ cần chứng minh rằng, khi $f$ chạy qua các tiết diện của $\sh{L}^{\otimes n}$ (với $n>0$) trên $X$, các $X_f$ lập thành một cơ sở của tôpô trên $X$ \sref{II.4.5.2}.
Mà cho $x\in X$ và $U$ là một lân cận tùy ý của $x$;
vì $g^{-1}(x)$ hữu hạn \sref{II.6.1.7} và $\sh{L}'$ ample, tồn tại một số nguyên $n>0$ và một tiết diện $f'$ của ${\sh{L}'}^{\otimes n}$ trên $X'$ sao cho $X'_{f'}$ là một lân cận của $g^{-1}(x)$ nằm trong $g^{-1}(U)$ \sref{II.4.5.4}.
Vì
\[
  \sh{L}^{\otimes n} = N_{X'/X}({\sh{L}'}^{\otimes n}),
\]
chỉ cần lấy $f=N_{X'/X}(f')$;
thật vậy, khi ấy $X\setminus X_f=g(X'\setminus X'_{f'})$ \sref{II.6.5.7}, và do đó $x\in X_f\subset U$.
\end{proof}

\begin{corollary}[6.6.3]
\label{II.6.6.3}
Dưới các giả thiết của \sref{II.6.6.1}, để một $\sh{O}_X$-môđun khả nghịch $\sh{L}$ ample đối với $f$, điều kiện cần và đủ là $\sh{L}'=g^*(\sh{L})$ ample đối với $f\circ g$.
\end{corollary}

\begin{proof}
Điều kiện là cần vì $g$ affine \sref{II.5.1.12}.
Để chứng minh điều kiện là đủ, ta có thể giả sử $Y$ affine \sref{II.4.6.4}, do đó $X$ và $X'$ quasi-compact và $\sh{L}'$ ample \sref{II.4.6.6}; còn phải chứng minh rằng $\sh{L}$ ample. Mà tập các điểm
\oldpage[II]{132}
$x\in X$ thừa nhận một lân cận trên đó $g_*(\sh{O}_{X'})$ (tương ứng, $g_*(\sh{O}_{X'})\otimes_{\sh{O}_X}\sh{R}(X)$) có hạng cho trước~$n$, trong giả thiết thứ nhất (tương ứng, thứ hai), vừa mở vừa đóng trong $X$; do đó $X$ là tổng của một số hữu hạn các tiền lược đồ mở ấy, và vì thế ta có thể giả sử nó bằng một trong số đó \sref{II.4.6.17}.
Nhưng khi ấy $N_{X'/X}(\sh{L}')=\sh{L}^{\otimes n}$; do đó $\sh{L}^{\otimes n}$ ample theo \sref{II.6.6.2}, và vì thế $\sh{L}$ cũng ample \sref{II.4.5.6}.
\end{proof}

\begin{corollary}[6.6.4]
\label{II.6.6.4}
Giả sử các giả thiết của \sref{II.6.6.1} được thỏa mãn và giả sử thêm rằng $f:X\to Y$ là kiểu hữu hạn.
Khi đó, để $f$ quasi-xạ ảnh, điều kiện cần và đủ là $f\circ g$ quasi-xạ ảnh.
Nếu giả sử thêm rằng $Y$ là một lược đồ quasi-compact hoặc một tiền lược đồ có không gian nền Noether, thì để $f$ xạ ảnh, điều kiện cần và đủ là $f\circ g$ xạ ảnh.
\end{corollary}

\begin{proof}
Giả thiết kéo theo $f\circ g$ là kiểu hữu hạn.
Xét đến định nghĩa của các cấu xạ quasi-xạ ảnh \sref{II.5.3.1}, khẳng định thứ nhất suy ra từ \sref{II.6.6.1} và \sref{II.6.6.3}.
Xét đến kết quả này và \sref{II.5.5.3}[(ii)], còn phải chứng minh rằng, khi giả sử $f$ quasi-xạ ảnh, để $f$ riêng, điều kiện cần và đủ là $f\circ g$ riêng.
Mà khi ấy $f$ tách \sref{II.5.3.1} và là kiểu hữu hạn;
vì $g$ toàn ánh, khẳng định của ta suy ra từ \sref{II.5.4.2}[(ii)] và \sref{II.5.4.3}[(ii)].
\end{proof}

Đặc biệt:

\begin{corollary}[6.6.5]
\label{II.6.6.5}
Cho $X$ là một tiền lược đồ kiểu hữu hạn trên một trường $K$ và $K'$ là một mở rộng bậc hữu hạn của $K$.
Để $X$ xạ ảnh (tương ứng, quasi-xạ ảnh) trên $K$, điều kiện cần và đủ là $X'=X\otimes_K K'$ xạ ảnh (tương ứng, quasi-xạ ảnh) trên $K'$.
\end{corollary}

\begin{proof}
Quả vậy, điều kiện là cần (\sref{II.5.3.4}[(iii)] và \sref{II.5.5.5}[(iii)]).
Ngược lại, giả sử điều kiện được thỏa mãn, và cho $g:X'\to X$ là phép chiếu chính tắc.
Rõ ràng $g$ là một cấu xạ hữu hạn \sref{II.6.1.5}[(iii)] và toàn ánh \sref[I]{I.3.5.2}[(ii)].
Hơn nữa, $g_*(\sh{O}_{X'})$ là một $\sh{O}_X$-môđun tự do địa phương vì nó đẳng cấu với $\sh{O}_X\otimes_K K'$ \sref{II.1.5.2}.
Khi đó, từ giả thiết cùng \sref{II.6.1.11} và \sref{II.5.5.5}[(ii)] suy ra $X'$ xạ ảnh (tương ứng, quasi-xạ ảnh) \emph{trên $K$};
rồi từ \sref{II.6.6.4} suy ra $X$ xạ ảnh (tương ứng, quasi-xạ ảnh) trên $K$.
\end{proof}

Trong Chương~V, ta sẽ chứng minh rằng mệnh đề \sref{II.6.6.5} vẫn đúng khi $K'$ là một mở rộng \emph{tùy ý} của $K$.

Phần còn lại của mục này dành cho việc chứng minh tiêu chuẩn \sref{II.6.6.11}, một tinh chỉnh khá kỹ thuật của \sref{II.6.6.1};
có thể bỏ qua phần này trong lần đọc đầu tiên.

\begin{lemma}[6.6.6]
\label{II.6.6.6}
Cho $X$ là một tiền lược đồ Noether rút gọn và $\sh{E}$ là một $\sh{O}_X$-môđun coherent sao cho $\sh{E}\otimes_{\sh{O}_X}\sh{R}(X)$ là một $\sh{R}(X)$-môđun tự do địa phương có hạng không đổi~$n$.
Khi đó tồn tại một tiền lược đồ Noether rút gọn $Z$ và một cấu xạ hữu hạn song hữu tỉ $h:Z\to X$ có tính chất sau: các cấu xạ của các bó tập hợp $\sigma_i:\shHom_{\sh{O}_X}(\sh{E},\sh{E})\to\sh{R}(X)$ (với $1\leq i\leq n$) (xem \sref{II.6.4.10}) ánh xạ $\shHom_{\sh{O}_X}(\sh{E},\sh{E})$ vào $\sh{O}_X$-đại số coherent $h_*(\sh{O}_Z)$.
\end{lemma}

\begin{proof}
Xét một tập mở affine $U$ của $X$, có vành $A(U)=A$;
gọi $E=\Gamma(U,\sh{E})$ và $C_U$ là đại số con của $R(U)$ sinh bởi các $\sigma_i(u)$ khi $u$ chạy qua $\Hom_A(E,E)$;
ta đã thấy trong \sref{II.6.4.7.1} rằng $A$-đại số này là một \emph{$A$-môđun hữu hạn sinh}.
Hơn nữa, rõ ràng phép tạo các đại số $C_U$ giao hoán với các phép hạn chế từ một tập mở affine $U$ xuống một tập mở affine $U'\subset U$.
Như vậy, ta đã xác định một $\sh{O}_X$-đại số con \emph{hữu hạn} $\sh{C}$ của $\sh{R}(X)$ sao cho $\Gamma(U,\sh{C})=C_U$ với mọi tập mở affine $U$ của $X$. Ta
\oldpage[II]{133}
lấy $Z=\Spec(\sh{C})$ và lấy $h$ là cấu xạ cấu trúc, do đó là hữu hạn \sref{II.6.1.2};
vì $\sh{C}$ rút gọn, $Z$ là một tiền lược đồ Noether rút gọn \sref{II.1.3.8}.
Cuối cùng, theo định nghĩa, vành các phân thức toàn phần của $C_U$ là $R(U)$, và vì $C_U$ nằm trong bao đóng nguyên của $A(U)$ trong $R(U)$, có một tương ứng song ánh giữa các iđêan nguyên tố tối tiểu của $A(U)$ và các iđêan nguyên tố tối tiểu của $C_U$ \cite[t.~I, p.~259]{I-13}; điều này chứng minh $h$ song hữu tỉ và kết thúc chứng minh.
\end{proof}

\begin{corollary}[6.6.7]
\label{II.6.6.7}
Dưới các giả thiết của \sref{II.6.6.6}, cho $W$ là một tập mở của $X$ sao cho, với mọi $x\in W$, hoặc $X$ chuẩn tắc tại điểm $x$, hoặc $\sh{E}_x$ là một $\sh{O}_x$-môđun tự do.
Khi đó ta có thể giả sử $h$ được xác định sao cho hạn chế của $h$ lên $h^{-1}(W)$ là một đẳng cấu từ $h^{-1}(W)$ lên $W$.
\end{corollary}

\begin{proof}
Mỗi trong hai giả thiết kéo theo rằng, nếu $U\subset W$ là một tập mở affine thì, với các ký hiệu của \sref{II.6.6.6}, $(\sigma_i(u))_x\in A_x$ với mọi $x\in U$ \sref{II.6.4.3}, do đó $\sigma_i(u)\in A$, và kết luận suy ra từ định nghĩa của $h$ trong \sref{II.6.6.6}.
\end{proof}

\begin{env}[6.6.8]
\label{II.6.6.8}
Cho $X$ là một tiền lược đồ Noether rút gọn và $g:X'\to X$ là một cấu xạ hữu hạn toàn ánh, sao cho $\sh{B}=g_*(\sh{O}_{X'})$ là một $\sh{O}_X$-đại số \emph{coherent};
giả sử thêm rằng $\sh{B}\otimes_{\sh{O}_X}\sh{R}(X)$ là một $\sh{R}(X)$-môđun \emph{tự do địa phương có hạng không đổi~$n$}.
Khi đó ta có thể áp dụng Bổ đề~\sref{II.6.6.6}, lấy $\sh{E}=\sh{B}$; từ đó, với các ký hiệu của \sref{II.6.6.6}, ta thu được một đồng cấu các bó monoid nhân $\sigma_n:\shHom_{\sh{O}_X}(\sh{B},\sh{B})\to h_*(\sh{O}_Z)$; hợp thành đồng cấu này với đồng cấu chính tắc $\sh{B}\to\shHom_{\sh{O}_X}(\sh{B},\sh{B})$ \sref{II.6.5.1}, ta thu được một đồng cấu các bó monoid nhân:
\[
\label{II.6.6.8.1}
  N' : \sh{B} = g_*(\sh{O}_{X'}) \to h_*(\sh{O}_Z) = \sh{C}.
\tag{6.6.8.1}
\]

Với cách đặt đó, với mọi $\sh{O}_{X'}$-môđun khả nghịch $\sh{L}'$, $g_*(\sh{L}')$ là một $\sh{B}$-môđun khả nghịch \sref{II.6.1.12}, và phương pháp của \sref{II.6.5.2} cho phép gắn theo cách có tính hàm tử với $\sh{L}'$ một $\sh{C}$-môđun khả nghịch, mà ta ký hiệu là $N'(g_*(\sh{L}'))$.
\end{env}

\begin{lemma}[6.6.9]
\label{II.6.6.9}
Cho $X$ là một tiền lược đồ Noether rút gọn và $g:X'\to X$ là một cấu xạ hữu hạn toàn ánh sao cho $g_*(\sh{O}_{X'})\otimes_{\sh{O}_X}\sh{R}(X)$ là một $\sh{R}(X)$-môđun tự do địa phương có hạng không đổi.
Khi đó tồn tại một tiền lược đồ Noether rút gọn $Z$ và một cấu xạ hữu hạn song hữu tỉ $h:Z\to X$ có tính chất sau: với mọi $\sh{O}_{X'}$-môđun ample $\sh{L}'$, $\sh{O}_Z$-môđun khả nghịch $\sh{M}$ sao cho $h_*(\sh{M})=N'(g_*(\sh{L}'))$ (với các ký hiệu của \sref{II.6.6.8}) là ample.
\end{lemma}

\begin{proof}
Trước hết, giả sử đồng cấu $\sh{B}\to\sh{B}\otimes_{\sh{O}_X}\sh{R}(X)$ là \emph{đơn ánh}.
Xác định $Z$ và $h$ như trong \sref{II.6.6.6} (với $\sh{E}=g_*(\sh{O}_{X'})$).
Cho $z\in Z$;
ta phải chứng minh rằng tồn tại một số nguyên $m>0$ và một tiết diện $t$ của $\sh{M}^{\otimes m}$ trên $Z$ sao cho $Z_t$ là một tập mở affine chứa $z$ \sref{II.4.5.2}.
Gọi $x=h(z)$ và cho $U$ là một tập mở affine của $X$ chứa $x$;
khi đó $h^{-1}(U)$ là một lân cận mở affine của $z$, và chỉ cần tìm $t$ sao cho $z\in Z_t\subset h^{-1}(U)$, vì khi ấy $Z_t$ tất yếu affine \sref{II.5.5.8}.
Theo giả thiết, tồn tại một số nguyên $n>0$ và một tiết diện $s'$ của ${\sh{L}'}^{\otimes n}$ trên $X'$ sao cho
\[
\label{II.6.6.9.1}
  g^{-1}(x) \subset X'_{s'} \subset g^{-1}(U)
\tag{6.6.9.1}
\]
theo \sref{II.4.5.4}.
Theo định nghĩa, $s'$ cũng là một tiết diện của $g_*({\sh{L}'}^{\otimes n})$ trên $X$, và như trong \sref{II.6.5.2}, tương ứng với nó là một tiết diện
\[
  s=N'(s')
  \quad\text{của}\quad
  N'\bigl(g_*({\sh{L}'}^{\otimes n})\bigr)
  =N'(g_*(\sh{L}'))^{\otimes n}
\]
trên
\oldpage[II]{134}
$X$.
Ta sẽ chứng minh rằng, nếu $t$ là tiết diện $s$ được xem như một tiết diện của $\sh{M}^{\otimes n}$ trên $Z$, thì $t$ là tiết diện cần tìm.
Đặt
\[
\label{II.6.6.9.2}
  V = X\setminus g(X'\setminus X'_{s'})
\tag{6.6.9.2}
\]
là một tập mở của $X$ chứa $x$ và nằm trong $U$, theo \sref{II.6.6.9.1} và \sref{II.6.1.10}.
Ta sẽ chứng minh rằng
\[
\label{II.6.6.9.3}
  h^{-1}(V) \subset Z_t \subset h^{-1}(U),
\tag{6.6.9.3}
\]
điều này sẽ kết thúc chứng minh.
Điều đó tương đương với việc tập $T$ gồm các $y\in X$ sao cho $s_y$ khả nghịch chứa $V$ và nằm trong $U$.
Để làm điều này, trước hết xét một tập mở affine $W$ nằm trong $V$;
khi đó $g^{-1}(W)$ là một tập mở affine trong $X'$, và theo \sref{II.6.6.9.2}, $s'_{y'}$ khả nghịch với mọi $y'\in g^{-1}(W)$;
nhờ các giả thiết về $X$ và $\sh{B}$, ta có thể áp dụng các kết quả của \sref{II.6.4.7}, và thấy rằng nếu $y=g(y')$ thì $s_y$ khả nghịch;
nói cách khác, $V\subset T$.
Mặt khác, cũng từ \sref{II.6.4.7} suy ra rằng, ngược lại, nếu $s_y$ khả nghịch thì $s'_{y'}$ cũng khả nghịch, điều này kéo theo $y'\in g^{-1}(U)$ theo \sref{II.6.6.9.1}, và do đó $y\in U$; vì thế $T\subset U$ trong trường hợp này.

Từ đó ta chuyển sang trường hợp tổng quát bằng cùng lập luận như trong \sref{II.6.6.2}, thay $X'$ bởi $X''$ sao cho $g'_*(\sh{O}_{X''})\to g'_*(\sh{O}_{X''})\otimes_{\sh{O}_X}\sh{R}(X)$ là đơn ánh, và thay $\sh{L}'$ bởi một $\sh{O}_{X''}$-môđun ample $\sh{L}''$ sao cho $N'(g_*(\sh{L}'))=N'(g'_*(\sh{L}''))$.
Như vậy Bổ đề~\sref{II.6.6.9} được chứng minh trong mọi trường hợp (với một cách chọn $h$ thích hợp).
\end{proof}

\begin{corollary}[6.6.10]
\label{II.6.6.10}
Giả sử các giả thiết của \sref{II.6.6.9} được thỏa mãn;
với mọi $\sh{O}_X$-môđun khả nghịch $\sh{L}$ sao cho $g^*(\sh{L})$ ample, $h^*(\sh{L})$ là ample.
\end{corollary}

\begin{proof}
Nếu đặt $\sh{L}'=g^*(\sh{L})$, thì $g_*(\sh{L}')=\sh{L}\otimes_{\sh{O}_X}\sh{B}$ \sref[0]{0.5.4.10}, do đó
\[
  N'(g_*(\sh{L}')) = (\sh{L}\otimes_{\sh{O}_X}\sh{C})^{\otimes n}
\]
(bằng cùng lập luận như đối với \sref{II.6.5.2.4}).
Do đó ta suy ra $\sh{M}=(h^*(\sh{L}))^{\otimes n}$, và vì $\sh{M}$ ample nên $h^*(\sh{L})$ cũng ample \sref{II.4.5.6}.
\end{proof}

\begin{proposition}[6.6.11]
\label{II.6.6.11}
Cho $Y$ là một lược đồ affine, $X$ là một tiền lược đồ Noether rút gọn, $f:X\to Y$ là một cấu xạ quasi-compact và $g:X'\to X$ là một cấu xạ hữu hạn toàn ánh.
Cho $W$ là một tập con mở của $X$ sao cho, với mọi $x\in W$, hoặc $X$ chuẩn tắc tại điểm $x$, hoặc tồn tại một lân cận mở $T\subset W$ của $x$ sao cho $(g_*(\sh{O}_{X'}))|T$ là một $(\sh{O}_X|T)$-môđun tự do địa phương.
Khi đó tồn tại một $Y$-tiền lược đồ rút gọn $Z$ và một $Y$-cấu xạ \emph{hữu hạn và song hữu tỉ} $h:Z\to X$ sao cho hạn chế của $h$ lên $h^{-1}(W)$ là một đẳng cấu từ $h^{-1}(W)$ lên $W$, đồng thời có tính chất sau: với mọi $\sh{O}_X$-môđun khả nghịch $\sh{L}$ sao cho $g^*(\sh{L})$ ample đối với $f\circ g$, $h^*(\sh{L})$ ample đối với $f\circ h$.
\end{proposition}

\begin{proof}
Vì $Y$ affine, $g^*(\sh{L})$ ample, và do đó cần chứng minh rằng, với một cách chọn $h$ thích hợp, $h^*(\sh{L})$ ample \sref{II.4.6.6}.
Ta sẽ chứng minh rằng có thể thay $g$ bởi một cấu xạ hữu hạn toàn ánh $g':X''\to X$ sao cho ${g'}^*(\sh{L})$ ample và $g'_*(\sh{O}_{X''})\otimes_{\sh{O}_X}\sh{R}(X)$ là một $\sh{R}(X)$-môđun \emph{tự do địa phương có hạng không đổi};
khi đó ta quy về các điều kiện của \sref{II.6.6.10}, và mệnh đề sẽ được chứng minh.

Để làm điều này, cho $\sh{B}=g_*(\sh{O}_{X'})$;
ký hiệu bởi $X_i$ (với $1\leq i\leq n$) các tiền lược đồ con đóng rút gọn của $X$ có không gian nền là các thành phần bất khả quy của $X$ \sref[I]{I.5.2.1};
chúng nguyên theo giả thiết.
Cho $X'_i$ là tiền lược đồ con đóng
\oldpage[II]{135}
$g^{-1}(X_i)$ của $X'$, và $g_i:X'_i\to X_i$ là hạn chế của cấu xạ $g$ lên $X'_i$, cấu xạ này hữu hạn \sref{II.6.1.5}[(iii)] và toàn ánh;
gọi $k_i$ là \emph{hạng} của $\sh{O}_{X_i}$-đại số $\sh{B}_i=(g_i)_*(\sh{O}_{X'_i})$.
Vì $\sh{B}_i\otimes_{\sh{O}_{X_i}}\sh{R}(X_i)$ là một tiền bó hằng \sref[I]{I.7.3.5}, nên hạng $k_i$ cũng là hạng của $(\sh{O}_X|U)$-đại số $\sh{B}|U$ đối với mọi tập mở $U$ của $X$ chỉ gặp thành phần bất khả quy $X_i$.
Nếu $T$ là một tập mở trong $X$ sao cho $\sh{B}|T$ đẳng cấu với $\sh{O}_X^m|T$, thì từ nhận xét trên suy ra các số $k_i$ bằng $m$ đối với mọi chỉ số $i$ sao cho $T\cap X_i\neq\emp$.
Gọi $U$ là tập mở của $X$ gồm các điểm mà tại một lân cận của chúng, $\sh{B}$ là một $\sh{O}_X$-môđun tự do địa phương, và gọi $U_j$ ($1\leq j\leq s$) là các thành phần liên thông của nó; các thành phần này mở trong $X$ và chỉ có hữu hạn (vì $U$ là Noether).
Ký hiệu bởi $V_j$ tiền lược đồ con đóng của $X'$ là \emph{bao đóng} của tiền lược đồ cảm sinh trên tập mở $g^{-1}(U_j)$ \sref[I]{I.9.5.11}.
Theo những điều trên, đối với mọi chỉ số $i$ sao cho $X_i\cap U_j\neq\emp$, các hạng $k_i$ đều bằng cùng một số nguyên $m_j$;
hơn nữa, cần lưu ý rằng cùng một $X_i$ không thể gặp hai $U_j$ có chỉ số khác nhau.
Gọi $i_\lambda$ là các chỉ số $i$ sao cho $X_i\cap U=\emp$.
Xét tích $k$ của tất cả các $k_i$, đặt $n_i=k/k_i$, và định nghĩa tiền lược đồ $X''$ như sau.
Với mỗi $j$ ($1\leq j\leq s$), xét $k/m_j$ tiền lược đồ đẳng cấu với $V_j$, và với mỗi $\lambda$, xét $k/k_{i_\lambda}=n_{i_\lambda}$ tiền lược đồ đẳng cấu với $X'_{i_\lambda}$;
$X''$ là \emph{tổng} của tất cả các tiền lược đồ ấy.
Ta định nghĩa một cấu xạ $g'':X''\to X'$ sao cho hạn chế của nó lên mỗi hạng tử của $X''$ là phép nhúng chính tắc;
hiển nhiên $g''$ là một cấu xạ hữu hạn trội, do đó toàn ánh (vì một cấu xạ hữu hạn là đóng \sref{II.6.1.10}).
Đặt $g'=g\circ g''$; đây là một cấu xạ hữu hạn toàn ánh $X''\to X$.
Ta có
\[
  {g'}^*(\sh{L})={g''}^*(g^*(\sh{L})),
\]
do đó ${g'}^*(\sh{L})$ là một $\sh{O}_{X''}$-môđun ample \sref{II.5.1.12}.
Khi ấy hiển nhiên, đối với tiền lược đồ mới $X''$ này, các hạng được định nghĩa như các $k_i$ đối với $X'$ đều \emph{bằng $k$};
xét đến \sref[I]{I.7.3.3}, suy ra ngay rằng với mọi tập mở affine $T$ của $X$,
\[
  (g'_*(\sh{O}_{X''})\otimes_{\sh{O}_X}\sh{R}(X))|T
\]
là một $(\sh{R}(X)|T)$-môđun đẳng cấu với $(\sh{R}(X)|T)^k$.
\end{proof}

\begin{corollary}[6.6.12]
\label{II.6.6.12}
Nếu trong phát biểu của \sref{II.6.6.11} ta có $W=X$, thì điều kiện cần và đủ để một $\sh{O}_X$-môđun khả nghịch $\sh{L}$ ample tương đối với $f$ là $g^*(\sh{L})$ ample tương đối với $f\circ g$.
\end{corollary}

\begin{remark}[6.6.13]
\label{II.6.6.13}
Trong Chương~III, ta sẽ thấy rằng nếu $Y$ là Noether, $f$ thuộc kiểu hữu hạn, và hạn chế của $f$ lên tiền lược đồ con đóng rút gọn của $X$ có không gian nền là $X\setminus W$ là \emph{riêng}, thì kết luận của \sref{II.6.6.12} vẫn còn đúng.
Nhưng trong Chương~V, ta sẽ đưa ra các ví dụ về những lược đồ đại số $X$ trên một trường $K$ (cấu xạ cấu trúc $X\to\Spec(K)$ không riêng) mà chuẩn hóa $X'$ của chúng là quasi-affine, song bản thân chúng không quasi-affine (thành thử $\sh{O}_X$ không ample, mặc dù $\sh{O}_{X'}$ ample \sref{II.5.1.2} và cấu xạ $X'\to X$ hữu hạn toàn ánh \sref{II.6.3.10}).
Trong mục tiếp theo, ta sẽ thấy tình huống này không thể xảy ra khi thay «quasi-affine» bằng «affine».
\end{remark}

\subsection{Định lý Chevalley}
\label{subsection:II.6.7}

Ta sẽ chứng minh (nhờ tiêu chuẩn Serre \sref{II.5.2.1}) định lý sau đây, mà C.~Chevalley đã chứng minh bằng các phương pháp khác trong trường hợp các lược đồ \emph{đại số}.
\oldpage[II]{136}
\begin{theorem}[6.7.1]
\label{II.6.7.1}
Cho $X$ là một lược đồ affine, $Y$ là một tiền lược đồ Noether, và $f:X\to Y$ là một cấu xạ hữu hạn toàn ánh.
Khi đó $Y$ là một lược đồ affine.
\end{theorem}

\begin{proof}
Hiển nhiên
\[
  f_\red:X_\red\to Y_\red
\]
là hữu hạn \sref{II.6.1.5}[(vi)];
vì $X_\red$ là một lược đồ affine, và nói $Y$ affine hay $Y_\red$ affine là như nhau do $Y$ là Noether \sref[I]{I.6.1.7}, nên ta có thể giả sử $Y$ \emph{rút gọn}.
Với mọi tập đóng $Y'$ của $Y$, khi ấy có duy nhất một tiền lược đồ con rút gọn của $Y$ có $Y'$ làm không gian nền \sref[I]{I.5.1.2};
ảnh ngược $f^{-1}(Y')$ của nó, đẳng cấu chính tắc với $X\times_Y Y'$ \sref[I]{I.4.4.1}, là affine vì là một tiền lược đồ con đóng của $X$, và hạn chế của $f$ lên $f^{-1}(Y')$, được đồng nhất với $f\times_Y1_{Y'}$, là một cấu xạ hữu hạn toàn ánh \sref{II.6.1.5}[(iii)].
Do nguyên lý quy nạp Noether \sref[0]{0.2.2.2}, ta có thể (xét đến \sref[I]{I.6.1.7}) quy về việc chứng minh định lý dưới giả thiết rằng, với \emph{mọi} tập đóng $Y'\neq Y$, mọi tiền lược đồ con đóng của $Y$ có $Y'$ làm không gian nền đều affine.
Từ đó suy ra rằng \emph{với mọi $\sh{O}_Y$-môđun coherent $\sh{F}$ có giá (đóng) $Z$ khác $Y$, ta có $\HH^1(Y,\sh{F})=0$}.
Thật vậy, tồn tại một tiền lược đồ con đóng $Y'$ của $Y$ có $Z$ làm không gian nền và sao cho, nếu $j:Y'\to Y$ là phép nhúng chính tắc, thì
\[
  \sh{F}=j_*(j^*(\sh{F}))
\]
\sref[I]{I.9.3.5};
do đó
\[
  \HH^1(Y,\sh{F})=\HH^1(Y',j^*(\sh{F}))=0
\]
theo \sref{II.5.2.3} và \sref[I]{I.5.1.9.2}.

Trước hết giả sử $Y$ không bất khả quy, và cho $Y'$ là một thành phần bất khả quy của $Y$, còn $Y''=Y\setminus Y'$;
ta vẫn ký hiệu bởi $Y'$ tiền lược đồ con đóng rút gọn của $Y$ có $Y'$ làm không gian nền, và bởi $j$ phép nhúng chính tắc $Y'\to Y$.
Cho $\sh{F}$ là một $\sh{O}_Y$-môđun coherent, và xét đồng cấu chính tắc
\[
  \rho:\sh{F}\to\sh{F}'=j_*(j^*(\sh{F}))
\]
\sref[0]{0.4.4.3};
$\sh{F}'$ là một $\sh{O}_Y$-môđun coherent theo \sref[0]{0.5.3.10} và \sref[0]{0.5.3.12}, vì
\[
  j_*(\sh{O}_{Y'})=\sh{O}_Y/\sh{J},
\]
trong đó $\sh{J}$ ký hiệu bó iđêan của $\sh{O}_Y$ xác định tiền lược đồ con $Y'$.
Do đó $\sh{G}=\Ker\rho$ và $\sh{K}=\Im\rho$ cũng là các $\sh{O}_Y$-môđun coherent \sref[0]{0.5.3.4};
nhưng theo định nghĩa, thớ $\sh{F}'_y$ của $\sh{F}'$ tại điểm tổng quát $y$ của $Y'$ bằng $\sh{F}_y$, vì $y$ là điểm trong của $Y'$ và do đó $\sh{J}_y=0$, bởi $Y$ rút gọn.
Suy ra $y$ không thuộc giá (đóng) của $\sh{G}$;
mặt khác, giá của $\sh{F}'$ (và huống chi giá của $\sh{K}$) chứa trong $Y'$;
nói cách khác, các giá của $\sh{G}$ và $\sh{K}$ đều \emph{khác $Y$}.
Suy ra
\[
  \HH^1(Y,\sh{G})=\HH^1(Y,\sh{K})=0,
\]
và dãy khớp đối đồng điều áp dụng cho dãy khớp
\[
  0\to\sh{G}\to\sh{F}\to\sh{K}\to0
\]
cho thấy $\HH^1(Y,\sh{F})=0$.
Khi đó kết luận suy ra từ tiêu chuẩn Serre \sref{II.5.2.1}.

Vậy giả sử $Y$ bất khả quy, và do đó \emph{nguyên}.
Ta cũng có thể giả sử $X$ \emph{nguyên}:
thật vậy, nếu ký hiệu bởi $X_i$ các tiền lược đồ con đóng rút gọn của $X$ có các thành phần bất khả quy của $X$ làm không gian nền \sref[I]{I.5.2.1}, và bởi $f_i$ hạn chế của $f$ lên $X_i$, thì ít nhất một trong các $f_i$ là trội, và vì đó là một cấu xạ hữu hạn \sref{II.6.1.5}, nên nó toàn ánh \sref{II.6.1.10};
mặt khác, vì $X_i$ là một lược đồ affine, ta thấy có thể thay $X$ bằng $X_i$ trong phát biểu.

\begin{lemma}[6.7.1.1]
\label{II.6.7.1.1}
Cho $X$, $Y$ là hai tiền lược đồ nguyên Noether, $x$ (tương ứng $y$) là điểm tổng quát của $X$ (tương ứng của $Y$), và $f:X\to Y$ là một cấu xạ hữu hạn toàn ánh.
Cho $\sh{L}$ là một $\sh{O}_X$-môđun khả nghịch sao cho tồn tại một lân cận mở affine $U$ của $y$ và một tiết diện $g\in\Gamma(X,\sh{L})$ sao cho
\oldpage[II]{137}
\[
  x\in X_g\subset f^{-1}(U).
\]
Khi đó tồn tại hai số nguyên $m,n>0$, một đồng cấu
\[
  u:\sh{O}_Y^m\to f_*(\sh{L}^{\otimes n}),
\]
và một lân cận mở $V$ của $y$ sao cho hạn chế $u|V$ là một đẳng cấu từ $\sh{O}_Y^m|V$ lên $f_*(\sh{L}^{\otimes n})|V$.
\end{lemma}

\begin{proof}
Gọi $C$ là vành (nguyên) của $U$, và $k=\sh{O}_y$ là trường phân thức của nó:
vì $f$ hữu hạn, $U'=f^{-1}(U)$ là affine \sref{II.1.3.2};
gọi $D$ là vành (nguyên) của nó, có trường phân thức $K=\sh{O}_x$;
theo giả thiết, $D$ là một $C$-môđun hữu hạn sinh \sref{II.6.1.4}, nên $K$ là một mở rộng có hạng hữu hạn của $k$.
Thớ
\[
  f^{-1}(y)=X\times_Y\Spec(\kres(y))=X\times_Y\Spec(\sh{O}_y)
\]
được đồng nhất với $\Spec(K)$ \sref[I]{I.3.6.6};
cho $s_i$ ($1\leq i\leq m$) là các phần tử của $D$ lập thành một cơ sở của $K$ trên $k$.
Tồn tại $n>0$ sao cho các tiết diện
\[
  (s_i|X_g)g^{\otimes n}
\]
của $\sh{L}^{\otimes n}$ trên $X_g$ kéo dài thành các tiết diện $b_i$ ($1\leq i\leq m$) của $\sh{L}^{\otimes n}$ trên $X$ \sref[I]{I.9.3.1}.
Theo định nghĩa, các $b_i$ cũng là các tiết diện của $f_*(\sh{L}^{\otimes n})$ trên $Y$, và vì thế xác định một đồng cấu
\[
  u:\sh{O}_Y^m\to f_*(\sh{L}^{\otimes n})
\]
\sref[0]{0.5.1.1};
ta sẽ chứng minh rằng $u$ đáp ứng yêu cầu.
Ta có
\[
  \sh{L}^{\otimes n}|U'=\widetilde{M},
\]
trong đó $M$ là một $D$-môđun hữu hạn sinh;
nếu $\vphi$ là phép nhúng $C\to D$ tương ứng với cấu xạ $f^{-1}(U)\to U$ là hạn chế của $f$, thì $M_{[\vphi]}$ là một $C$-môđun hữu hạn sinh.
Do đó
\[
  f_*(\sh{L}^{\otimes n})|U=(M_{[\vphi]})\supertilde
\]
\sref[I]{I.1.6.3} là coherent, và vì $U$ là một tập mở affine tùy ý của $Y$, nên $f_*(\sh{L}^{\otimes n})$ coherent.
Hơn nữa, $u|U=\widetilde{\theta}$, trong đó $\theta$ là một $C$-đồng cấu
\[
  C^m\to M_{[\vphi]},
\]
và $u_y=\theta_y$ là đồng cấu
\[
  \theta\otimes1:k^m=C^m\otimes_Ck\to M_{[\vphi]}\otimes_Ck;
\]
theo định nghĩa, đồng cấu sau cùng này là một \emph{đẳng cấu}, vì các $(b_i)_x$ lập thành một cơ sở của $M_{[\vphi]}\otimes K$ trên $k$, còn theo giả thiết $g_x$ là một phần tử sinh của $(\sh{L}^{\otimes n})_x$.
Suy ra các giá của những $\sh{O}_Y$-môđun $\Ker u$ và $\Coker u$ không chứa $y$;
vì các $\sh{O}_Y$-môđun này coherent \sref[0]{0.5.3.3}, các giá của chúng đóng \sref[0]{0.5.2.2}, và bổ đề được chứng minh.
\end{proof}

Với điều đó, các giả thiết của Bổ đề \sref{II.6.7.1.1} được thỏa mãn trong trường hợp đang xét khi lấy $\sh{L}=\sh{O}_X$, vì $X$ affine \sref[I]{I.1.1.10};
ta đặt $\sh{A}=\sh{O}_Y$, $\sh{B}=f_*(\sh{O}_X)$.
Theo tiêu chuẩn Serre \sref{II.5.2.1.1}, chỉ cần chứng minh rằng với mọi $\sh{O}_Y$-môđun coherent $\sh{F}$, ta có $\HH^1(Y,\sh{F})=0$;
thậm chí chỉ cần chứng minh điều đó khi $\sh{F}\subset\sh{O}_Y$, điều này kéo theo $\sh{F}$ không xoắn vì $Y$ nguyên;
thực ra, ta sẽ chứng minh $\HH^1(Y,\sh{F})=0$ đối với \emph{mọi} $\sh{O}_Y$-môđun coherent \emph{không xoắn} $\sh{F}$.
Đồng cấu $u:\sh{A}^m\to\sh{B}$ xác định một đồng cấu
\[
  v:\sh{G}=\shHom_{\sh{A}}(\sh{B},\sh{F})
  \to
  \shHom_{\sh{A}}(\sh{A}^m,\sh{F})
  =\sh{F}^m.
\]

Trước hết chứng minh $v$ \emph{đơn ánh}:
theo giả thiết, $\sh{T}=\Coker u$ có giá không gặp $V$, nên là một $\sh{O}_Y$-môđun \emph{xoắn} \sref[I]{I.7.4.6};
dãy khớp
\[
  \sh{A}^m\to\sh{B}\to\sh{T}\to0
\]
cho, nhờ tính khớp trái của hàm tử $\shHom_{\sh{A}}$, dãy khớp
\[
  0\to\shHom_{\sh{A}}(\sh{T},\sh{F})
  \to\sh{G}\xrightarrow{v}\sh{F}^m.
\]
Nhưng vì $\sh{F}$ không xoắn, ta có
\[
  \shHom_{\sh{A}}(\sh{T},\sh{F})=0
\]
\sref[0]{0.5.2.6}, từ đó suy ra khẳng định.
Vậy ta có dãy khớp
\[
  0\to\sh{G}\to\sh{F}^m\to\Coker v\to0
\]
\oldpage[II]{138}
trong đó $\sh{G}$ và $\Coker v$ là các $\sh{O}_Y$-môđun coherent (\sref[0]{0.5.3.4} và \sref[0]{0.5.3.5});
theo dãy khớp đối đồng điều, chỉ cần chứng minh
\[
  \HH^1(Y,\sh{G})=\HH^1(Y,\Coker v)=0
\]
để suy ra
\[
  \HH^1(Y,\sh{F}^m)=(\HH^1(Y,\sh{F}))^m=0,
\]
và do đó $\HH^1(Y,\sh{F})=0$.
Nhưng hạn chế $v|V$ là một đẳng cấu, nên giá của $\Coker v$ khác $Y$, từ đó
\[
  \HH^1(Y,\Coker v)=0
\]
theo giả thiết ban đầu.
Mặt khác, $\sh{G}$ là một $\sh{B}$-môđun \emph{coherent} \sref[I]{I.9.6.4};
vì $X$ affine trên $Y$, tồn tại một $\sh{O}_X$-môđun quasi-coherent $\sh{K}$ sao cho $\sh{G}$ đẳng cấu với $f_*(\sh{K})$ \sref{II.1.4.3};
vì $X$ affine, ta có $\HH^1(X,\sh{K})=0$ \sref[I]{I.5.1.9.2}, nên cũng có $\HH^1(Y,\sh{G})=0$ theo \sref{II.5.2.3}; điều này hoàn tất chứng minh Định lý \sref{II.6.7.1}.
\end{proof}

\begin{corollary}[6.7.2]
\label{II.6.7.2}
Cho $X$ là một tiền lược đồ Noether và $(X_i)_{1\leq i\leq n}$ là một phủ hữu hạn của không gian $X$ gồm các tập đóng.
Điều kiện cần và đủ để $X$ affine là với mỗi $i$, tồn tại một tiền lược đồ con đóng của $X$, có $X_i$ làm không gian nền và affine.
\end{corollary}

\begin{proof}
Thật vậy, nếu điều đó đúng, gọi $X'$ là lược đồ \emph{tổng} của các $X_i$;
hiển nhiên $X'$ affine, và ta định nghĩa một cấu xạ toàn ánh $f:X'\to X$ bằng cách lấy hạn chế của $f$ lên $X_i$ là phép nhúng chính tắc.
Theo \sref{II.6.7.1}, chỉ còn phải kiểm tra $f$ \emph{hữu hạn}, và điều này đã được chứng minh trong \sref{II.6.1.5}.
\end{proof}

\begin{corollary}[6.7.3]
\label{II.6.7.3}
Điều kiện cần và đủ để một tiền lược đồ Noether $X$ affine là các tiền lược đồ con đóng rút gọn có những thành phần bất khả quy của $X$ làm không gian nền đều affine.
\end{corollary}

\section{Các tiêu chuẩn định giá}
\label{section:II.7}

Trong mục này, ta đưa ra các tiêu chuẩn định giá cho tính tách và tính riêng của một cấu xạ, nghĩa là các tiêu chuẩn có sự tham gia của một lược đồ phụ trợ biến thiên dạng $\Spec(A)$, trong đó $A$ là một vành định giá.
Với những giả thiết «Noether» thích hợp, trong các tiêu chuẩn này ta có thể giới hạn vào trường hợp $A$ là một vành định giá \emph{rời rạc}.
Đây hẳn là trường hợp duy nhất mà ta sẽ phải áp dụng về sau; trong trường hợp tổng quát, ta chỉ đưa các vành định giá tùy ý vào để nối với những trình bày cổ điển.

\subsection{Nhắc lại về các vành định giá}
\label{subsection:II.7.1}

\begin{env}[7.1.1]
\label{II.7.1.1}
Trong số các tính chất tương đương khác nhau đặc trưng cho các vành định giá, ta sẽ dùng tính chất sau:
một vành $A$ được gọi là một \emph{vành định giá} nếu nó là một vành nguyên không phải là trường, và nếu trong tập hợp các vành địa phương nằm trong trường phân thức $K$ của $A$ và khác $K$, vành $A$ là \emph{tối đại} đối với quan hệ chi phối \sref[I]{I.8.1.1}.
Nhắc lại rằng một vành định giá là \emph{đóng nguyên}.
Nếu $A$ là một vành định giá, thì $A_\mathfrak{p}$ cũng là một vành định giá với mọi iđêan nguyên tố $\mathfrak{p}\neq0$ của $A$.
\end{env}

\begin{env}[7.1.2]
\label{II.7.1.2}
Cho $K$ là một trường và $A$ là một vành con địa phương của $K$ không phải là trường;
\oldpage[II]{139}
khi đó tồn tại một vành định giá có $K$ làm trường phân thức và chi phối $A$ (\cite[p.~1-07, Bổ đề~2]{I-1}).

Mặt khác, cho $B$ là một vành định giá, $k$ là trường thặng dư của nó, $K$ là trường phân thức của nó, và $L$ là một mở rộng của $k$.
Khi đó tồn tại một vành định giá \emph{đầy đủ} $C$ chi phối $B$ và có trường thặng dư là $L$.
Thật vậy, $L$ là một mở rộng đại số của một mở rộng thuần siêu việt
\[
  L'=k(T_\mu)_{\mu\in M};
\]
ta biết rằng có thể mở rộng định giá của $K$ tương ứng với $B$ thành một định giá của
\[
  K'=K(T_\mu)_{\mu\in M}
\]
sao cho $L'$ là trường thặng dư của định giá này (\cite[p.~98]{II-24});
thay $B$ bởi vành đầy đủ hóa của vành thuộc định giá đã mở rộng ấy, ta thấy có thể giới hạn vào trường hợp $B$ đầy đủ và $L$ là một bao đóng đại số của $k$.
Nếu $\overline{K}$ là một bao đóng đại số của $K$, khi ấy ta có thể mở rộng lên $\overline{K}$ định giá xác định $B$, và trường thặng dư tương ứng là một bao đóng đại số của $k$, như thấy được bằng cách nâng lên $\overline{K}$ các hệ số của một đa thức monic thuộc $k[T]$.
Cuối cùng ta quy về trường hợp $L=k$, và khi đó chỉ cần lấy $C$ là vành đầy đủ hóa của $B$ để đáp ứng yêu cầu.
\end{env}

\begin{env}[7.1.3]
\label{II.7.1.3}
Cho $K$ là một trường và $A$ là một vành con của $K$;
bao đóng nguyên $A'$ của $A$ trong $K$ là giao của các vành định giá chứa $A$ và có $K$ làm trường phân thức (\cite[p.~51, Định lý~2]{I-11}).
Mệnh đề \sref{II.7.1.2} được biểu diễn dưới dạng hình học tương đương sau đây.
\end{env}

\begin{proposition}[7.1.4]
\label{II.7.1.4}
Cho $Y$ là một tiền lược đồ, $p:X\to Y$ là một cấu xạ, $x$ là một điểm của $X$, $y=p(x)$, và $y'\neq y$ là một chuyên biệt hóa \sref[0]{0.2.1.2} của $y$.
Khi đó tồn tại một lược đồ địa phương $Y'$, là phổ của một vành định giá, và một cấu xạ tách $f:Y'\to Y$ sao cho, nếu $a$ ký hiệu điểm đóng duy nhất của $Y'$ và $b$ là điểm tổng quát của $Y'$, thì
\[
  f(a)=y'
  \qquad\text{và}\qquad
  f(b)=y.
\]
Ngoài ra, có thể giả sử một trong hai tính chất bổ sung sau được thỏa mãn:
\begin{enumerate}
  \item[\rm{(i)}] $Y'$ là phổ của một vành định giá đầy đủ có trường thặng dư đóng đại số, và tồn tại một $\kres(y)$-đồng cấu $\kres(x)\to\kres(b)$.
  \item[\rm{(ii)}] Tồn tại một $\kres(y)$-đẳng cấu $\kres(x)\simto\kres(b)$.
\end{enumerate}
\end{proposition}

\begin{proof}
Cho $Y_1$ là tiền lược đồ con đóng rút gọn của $Y$ có $\overline{\{y\}}$ làm không gian nền \sref[I]{I.5.2.1}, và cho $X_1$ là tiền lược đồ con đóng ảnh ngược $p^{-1}(Y_1)$;
vì $y'\in\overline{\{y\}}$ theo giả thiết và $\kres(x)$ trong $X$ cũng chính là $\kres(x)$ trong $X_1$, ta có thể giả sử $Y$ \emph{nguyên}, có điểm tổng quát $y$.
Khi đó $\sh{O}_{y'}$ là một vành địa phương nguyên không phải là trường, trường phân thức của nó là $\sh{O}_y=\kres(y)$, và $\kres(x)$ là một mở rộng của $\kres(y)$.

Để thực hiện các điều kiện $f(a)=y'$ và $f(b)=y$, cũng như điều kiện bổ sung~(i), trước hết chọn một vành định giá $B$ chi phối $\sh{O}_{y'}$ và có $\kres(x)$ làm trường phân thức, rồi chọn một vành định giá đầy đủ $A'$ chi phối $B$ và có trường thặng dư đóng đại số.
Khi đó
\[
  \kres(x)=\Frac(B)\longrightarrow\Frac(A')=\kres(b)
\]
là $\kres(y)$-đồng cấu cần tìm.
Đối với điều kiện~(ii), lấy trực tiếp một vành định giá $A'$ chi phối $\sh{O}_{y'}$ và có $\kres(x)$ làm trường phân thức.
Trong cả hai trường hợp, đặt $Y'=\Spec(A')$; sự tồn tại của các vành đã nêu được bảo đảm bởi \sref{II.7.1.2}.
\end{proof}

\begin{env}[7.1.5]
\label{II.7.1.5}
Nhắc lại rằng một vành địa phương $A$ được gọi là \emph{có chiều~1} nếu tồn tại một iđêan nguyên tố khác iđêan tối đại $\mathfrak{m}$, và nếu mọi iđêan nguyên tố của $A$ khác $\mathfrak{m}$ đều là một iđêan nguyên tố \emph{tối tiểu};
khi $A$ \emph{nguyên}, điều này tương đương với việc $\mathfrak{m}$ và $(0)$ là các iđêan nguyên tố duy nhất và $\mathfrak{m}\neq(0)$;
nói cách khác, $Y=\Spec(A)$ chỉ gồm hai
\oldpage[II]{140}
các điểm $a$ và $b$: $a$ là điểm \emph{đóng} duy nhất, ta có $\mathfrak{j}_a=\mathfrak{m}$, và $\kres(a)=k$ là \emph{trường thặng dư} $k=A/\mathfrak{m}$;
$b$ là \emph{điểm tổng quát} của $Y$, $\mathfrak{j}_b=(0)$, tập $\{b\}$ là tập mở duy nhất của $Y$ khác $\emp$ và $Y$ (do đó là một tập mở \emph{trù mật khắp nơi}), và $\kres(b)=K$ là \emph{trường phân thức} của $A$.
\end{env}

\begin{env}[7.1.6]
\label{II.7.1.6}
Đối với một vành địa phương Noether $A$ có chiều~1, ta biết (\cite[pp.~2-08 và 17-01]{I-1}) rằng các điều kiện sau tương đương:
\begin{enumerate}
    \item[\rm{(a)}] $A$ chuẩn tắc;
    \item[\rm{(b)}] $A$ chính quy;
    \item[\rm{(c)}] $A$ là một vành định giá;
\end{enumerate}
hơn nữa, khi đó $A$ là một \emph{vành định giá rời rạc}.
Các Mệnh đề~\sref{II.7.1.2} và \sref{II.7.1.3} khi ấy có các mệnh đề tương tự sau đây đối với các vành định giá rời rạc:
\end{env}

\begin{proposition}[7.1.7]
\label{II.7.1.7}
Cho $A$ là một vành địa phương nguyên Noether không phải là một trường, $K$ là trường phân thức của nó, và $L$ là một mở rộng hữu hạn sinh của $K$;
khi đó tồn tại một vành định giá rời rạc chi phối $A$ và có $L$ làm trường phân thức.
\end{proposition}

\begin{proof}
Trước hết, giả sử $L=K$.
Gọi $\mathfrak{m}$ là iđêan tối đại của $A$, $(x_1,\ldots,x_n)$ là một hệ sinh gồm các phần tử khác không của $\mathfrak{m}$, và $B$ là vành con $A[x_2/x_1,\ldots,x_n/x_1]$ của $K$, vốn là Noether.
Hiển nhiên iđêan $\mathfrak{m}B$ của $B$ đồng nhất với iđêan chính $x_1B$;
nếu $\mathfrak{p}$ là một iđêan nguyên tố tối tiểu chứa $x_1B$, thì $\mathfrak{p}$ có độ cao~1 (\cite[t.~I, p.~277]{I-13});
nói cách khác, $B_\mathfrak{p}$ là một vành địa phương Noether \emph{có chiều~1};
hiển nhiên $\mathfrak{p}B_\mathfrak{p}\cap A$ là một iđêan của $A$ chứa $\mathfrak{m}$ và không chứa $1$, nên bằng $\mathfrak{m}$, và do đó $B_\mathfrak{p}$ \emph{chi phối} $A$ \sref[I]{I.8.1.1}.
Từ Định lý Krull--Akizuki (\cite[p.~293]{II-25}) suy ra rằng bao đóng nguyên $C$ của $B_\mathfrak{p}$ là một vành Noether (mặc dù $C$ không nhất thiết là một $B_\mathfrak{p}$-môđun hữu hạn sinh);
nếu $\mathfrak{n}$ là một iđêan tối đại của $C$, thì $C_\mathfrak{n}$ là một vành địa phương Noether chuẩn tắc có chiều~1 (\cite[p.~295]{II-25}), do đó là một vành định giá rời rạc chi phối $B_\mathfrak{p}$ và, \emph{a fortiori}, chi phối $A$.

Bây giờ, nếu $L$ là một mở rộng hữu hạn sinh của $K$, thì theo phần trên, ta có thể quy về trường hợp $A$ đã là một vành định giá rời rạc.
Gọi $w$ là một định giá của $K$ liên kết với $A$;
tồn tại một định giá rời rạc $w'$ của $L$ \emph{mở rộng} $w$: thật vậy, bằng quy nạp theo số phần tử sinh của $L$, ta quy về trường hợp $L=K(\alpha)$, và khi đó mệnh đề là cổ điển (\cite[p.~106]{II-24}).
\end{proof}

\begin{corollary}[7.1.8]
\label{II.7.1.8}
Cho $A$ là một vành nguyên Noether, $K$ là trường phân thức của nó, và $L$ là một mở rộng hữu hạn sinh của $K$.
Khi đó bao đóng nguyên của $A$ trong $L$ là giao của các vành định giá rời rạc có $L$ làm trường phân thức và chứa $A$.
\end{corollary}

\begin{proof}
Thật vậy, một vành định giá rời rạc như thế, vì chuẩn tắc, \emph{a fortiori} chứa mọi phần tử của $L$ nguyên trên $A$.
Do đó chỉ cần chứng minh rằng, nếu $x\in L$ không nguyên trên $A$, thì tồn tại một vành định giá rời rạc $C$ có $L$ làm trường phân thức, chứa $A$ và không chứa $x$.
Giả thiết về $x$ kéo theo $x\not\in B=A[1/x]$, hay nói cách khác, $1/x$ không khả nghịch trong vành Noether $B$.
Vì thế tồn tại một iđêan nguyên tố $\mathfrak{p}$ của $B$ chứa $1/x$.
Vành địa phương nguyên $B_\mathfrak{p}$ là Noether và được chứa trong $L$, vốn là một mở rộng hữu hạn sinh của trường phân thức của $B_\mathfrak{p}$ (trường sau chứa $K$).
Theo \sref{II.7.1.7}, do đó tồn tại một vành định giá rời rạc $C$ chi phối $B_\mathfrak{p}$ và có $L$ làm trường phân thức;
vì $1/x\in\mathfrak{p}B_\mathfrak{p}$ thuộc iđêan tối đại của $C$, ta có $x\not\in C$, điều này kết thúc chứng minh.
\end{proof}

Dạng hình học của \sref{II.7.1.7} như sau:
\oldpage[II]{141}
\begin{proposition}[7.1.9]
\label{II.7.1.9}
Cho $Y$ là một tiền lược đồ Noether địa phương, $p:X\to Y$ là một cấu xạ thuộc kiểu hữu hạn địa phương, $x$ là một điểm của $X$, $y=p(x)$, và $y'\neq y$ là một chuyên biệt hóa của $y$.
Khi đó tồn tại một lược đồ địa phương $Y'$, là phổ của một vành định giá rời rạc, một cấu xạ tách $f:Y'\to Y$, và một $Y$-ánh xạ hữu tỉ $g$ từ $Y'$ vào $X$ sao cho, nếu ký hiệu bởi $a$ điểm đóng của $Y'$ và bởi $b$ điểm tổng quát của $Y'$, thì $f(a)=y'$, $f(b)=y$, $g(b)=x$, và trong biểu đồ giao hoán
\[
\label{II.7.1.9.1}
    \xymatrix{
        & \kres(x) \ar[dl]_{\gamma}
    \\  \kres(b)
        & \kres(y) \ar[u]_{\pi} \ar[l]^{\vphi}
    }
\tag{7.1.9.1}
\]
(trong đó $\pi$, $\vphi$ và $\gamma$ lần lượt là các đồng cấu tương ứng với $p$, $f$ và $g$), đồng cấu $\gamma$ là song ánh.
\end{proposition}

\begin{proof}
Như trong \sref{II.7.1.4}, ta có thể quy về trường hợp $Y$ nguyên với điểm tổng quát $y$ (có xét đến \sref[I]{I.6.3.4}[iv]), và vì bài toán có tính địa phương trên $X$ và $Y$, có thể giả sử $p$ thuộc kiểu hữu hạn;
khi đó ta ở trong tình huống của \sref{II.7.1.4}, với tính chất bổ sung rằng $\kres(x)$ là một mở rộng \emph{hữu hạn sinh} của $\kres(y)$ \sref[I]{I.6.4.11} và $\sh{O}_{y'}$ là Noether;
điều này cho phép áp dụng \sref{II.7.1.7} và lấy $Y'=\Spec(A')$, trong đó $A'$ là một vành định giá rời rạc chi phối $\sh{O}_{y'}$ và có trường phân thức là $\kres(x)$.
Như vậy ta đã xác định một biểu đồ giao hoán \sref{II.7.1.9.1}, trong đó $\gamma$ là song ánh, còn $\pi$ và $\vphi$ tương ứng với các cấu xạ $p$ và $f$.
Hơn nữa, vì $X$ và $Y$ là Noether địa phương \sref[I]{I.6.6.2} và $Y'$ nguyên, tồn tại duy nhất một $Y$-ánh xạ hữu tỉ $g$ từ $Y'$ vào $X$ tương ứng với đẳng cấu $\gamma$ \sref[I]{I.7.1.15}, điều này kết thúc chứng minh.
\end{proof}

\subsection{Tiêu chuẩn định giá về tính tách}
\label{subsection:II.7.2}

\begin{proposition}[7.2.1]
\label{II.7.2.1}
Cho $X$ và $Y$ là hai tiền lược đồ, và $f:X\to Y$ là một cấu xạ quasi-compact.
Hai điều kiện sau tương đương:
\begin{enumerate}
    \item[\rm{(a)}] Cấu xạ $f$ là đóng.
    \item[\rm{(b)}] Với mọi $x\in X$ và mọi chuyên biệt hóa $y'\neq y$ của $y=f(x)$, tồn tại một chuyên biệt hóa $x'$ của $x$ sao cho $f(x')=y'$.
\end{enumerate}
\end{proposition}

\begin{proof}
Điều kiện~(b) diễn đạt rằng $f(\overline{\{x\}})=\overline{\{y\}}$, nên là hệ quả của~(a).
Để chứng minh (b) kéo theo (a), xét một tập con đóng $X'$ của không gian nền $X$;
đặt $Y'=\overline{f(X')}$, và ta chứng minh $Y'=f(X')$.
Xét các tiền lược đồ con đóng rút gọn của $X$ và $Y$ có không gian nền lần lượt là $X'$ và $Y'$ \sref[I]{I.5.2.1};
khi đó tồn tại một cấu xạ $f':X'\to Y'$ sao cho biểu đồ
\[
    \xymatrix{
        X' \ar[r]^{f'} \ar[d]
        & Y' \ar[d]
    \\  X \ar[r]_{f}
        & Y
    }
\]
giao hoán \sref[I]{I.5.2.2}, và vì $f$ quasi-compact nên $f'$ cũng vậy.
Do đó ta quy về chứng minh rằng, nếu $f$ là một cấu xạ quasi-compact và \emph{trội}, thì
\oldpage[II]{142}
điều kiện~(b) kéo theo $f(X)=Y$.
Thật vậy, lấy một điểm $y'$ của $Y$, và gọi $y$ là điểm tổng quát của một thành phần bất khả quy của $Y$ chứa $y'$;
theo~(b), chỉ cần chứng minh rằng $f^{-1}(y)$ không rỗng.
Nhưng ta biết tính chất này là hệ quả của việc $f$ quasi-compact và trội \sref[I]{I.6.6.5}.
\end{proof}

\begin{corollary}[7.2.2]
\label{II.7.2.2}
Cho $f:X\to Y$ là một phép nhúng quasi-compact.
Để không gian nền $X$ đóng trong $Y$, điều kiện cần và đủ là nó chứa mọi chuyên biệt hóa, trong $Y$, của từng điểm của nó.
\end{corollary}

\begin{proposition}[7.2.3]
\label{II.7.2.3}
Cho $Y$ là một tiền lược đồ (tương ứng, một tiền lược đồ Noether địa phương), và $f:X\to Y$ là một cấu xạ (tương ứng, một cấu xạ thuộc kiểu hữu hạn địa phương).
Các điều kiện sau tương đương:
\begin{enumerate}
    \item[\rm{(a)}] $f$ tách.
    \item[\rm{(b)}] Cấu xạ đường chéo $X\to X\times_Y X$ là quasi-compact và, đối với mọi $Y$-tiền lược đồ dạng $Y'=\Spec(A)$, trong đó $A$ là một vành định giá (tương ứng, một vành định giá rời rạc), hai $Y$-cấu xạ từ $Y'$ vào $X$ trùng nhau tại điểm tổng quát của $Y'$ thì bằng nhau.
    \item[\rm{(c)}] Cấu xạ đường chéo $X\to X\times_Y X$ là quasi-compact và, đối với mọi $Y$-tiền lược đồ dạng $Y'=\Spec(A)$, trong đó $A$ là một vành định giá (tương ứng, một vành định giá rời rạc), hai $Y'$-tiết diện của $X'=X_{(Y')}$ trùng nhau tại điểm tổng quát của $Y'$ thì bằng nhau.
\end{enumerate}
\end{proposition}

\begin{proof}
Sự tương đương giữa (b) và (c) suy ra từ tương ứng song ánh giữa các $Y$-cấu xạ từ $Y'$ vào $X$ và các $Y'$-tiết diện của $X'$ \sref[I]{I.3.3.14}.
Nếu $X$ tách trên $Y$, thì điều kiện~(b) được thỏa mãn theo \sref[I]{I.7.2.2.1}, vì $Y'$ nguyên.
Còn phải chứng minh rằng (b) kéo theo cấu xạ đường chéo $\Delta:X\to X\times_Y X$ là đóng, và điều này tương đương với việc chứng minh nó thỏa tiêu chuẩn của \sref{II.7.2.2}.
Vậy lấy $z$ là một điểm của đường chéo $\Delta(X)$, và $z'\neq z$ là một chuyên biệt hóa của $z$ trong $X\times_Y X$.
Khi đó tồn tại \sref{II.7.1.4} một vành định giá $A$ và một cấu xạ $f$ từ $Y'=\Spec(A)$ vào $X\times_Y X$ sao cho $f$ gửi điểm đóng $a$ của $Y'$ đến $z'$ và điểm tổng quát $b$ của $Y'$ đến $z$;
cấu xạ này làm cho $Y'$ thành một $(X\times_Y X)$-tiền lược đồ và, \emph{a fortiori}, một $Y$-tiền lược đồ.
Hợp $f$ với hai phép chiếu của $X\times_Y X$, ta thu được hai $Y$-cấu xạ $g_1$ và $g_2$ từ $Y'$ vào $X$ mà theo giả thiết trùng nhau tại điểm $b$;
do đó chúng bằng cùng một cấu xạ $g$, điều này có nghĩa \sref[I]{I.5.3.1} rằng $f$ phân tích thành $f=\Delta\circ g$, và vì thế $z'\in\Delta(X)$.
Nếu giả sử $Y$ Noether địa phương và $f$ thuộc kiểu hữu hạn địa phương, thì $X\times_Y X$ là Noether địa phương \sref[I]{I.6.6.7};
khi đó có thể lặp lại lập luận trên, giả sử $A$ là một vành định giá rời rạc, theo \sref{II.7.1.9}.
\end{proof}

\begin{remark}[7.2.4]
\label{II.7.2.4}
\begin{enumerate}
    \item[\rm{(i)}] Giả thiết rằng cấu xạ $\Delta$ quasi-compact luôn được thỏa mãn khi $Y$ Noether địa phương và $f$ thuộc kiểu hữu hạn địa phương, vì khi đó $X\times_Y X$ là Noether địa phương \sref[I]{I.6.6.4}[i].
        Trong trường hợp tổng quát, giả thiết ấy cũng có nghĩa rằng, đối với mọi phủ $(U_\alpha)$ của $X$ bởi các mở affine, các tập $U_\alpha\cap U_\beta$ đều \emph{quasi-compact}.
    \item[\rm{(ii)}] Để $f$ tách, \emph{chỉ cần} đòi hỏi điều kiện~(b) hoặc (c) đối với các vành định giá $A$ \emph{đầy đủ} có trường thặng dư \emph{đóng đại số};
        điều này suy ra từ chứng minh của \sref{II.7.2.3} và từ \sref{II.7.1.4}.
\end{enumerate}
\end{remark}
\oldpage[II]{143}

\subsection{Tiêu chuẩn định giá về tính riêng}
\label{subsection:II.7.3}

\begin{proposition}[7.3.1]
\label{II.7.3.1}
Cho $A$ là một vành định giá, $Y=\Spec(A)$, $b$ là điểm tổng quát của $Y$, $X$ là một \emph{lược đồ} nguyên, và $f:X\to Y$ là một cấu xạ \emph{đóng} sao cho $f^{-1}(b)$ chỉ gồm một điểm $x$ và đồng cấu tương ứng $\kres(b)\to\kres(x)$ là song ánh.
Khi đó $f$ là một đẳng cấu.
\end{proposition}

\begin{proof}
Vì $f$ đóng và trội, ta có $f(X)=Y$;
theo \sref[I]{I.4.2.2}, chỉ cần chứng minh rằng, với mọi $y'\neq b$ trong $Y$, tồn tại \emph{duy nhất một} điểm $x'$ sao cho $f(x')=y'$, và đồng cấu tương ứng $\sh{O}_{y'}\to\sh{O}_{x'}$ là song ánh, vì khi đó $f$ sẽ là một đồng phôi.
Thật vậy, nếu $f(x')=y'$, thì $\sh{O}_{x'}$ là một vành địa phương được chứa trong $K=\kres(x)=\kres(b)$ và chi phối $\sh{O}_{y'}$;
vành sau là vành địa phương $A_{y'}$, do đó là một vành định giá \sref{II.7.1.1} có $K$ làm trường phân thức.
Mặt khác, $\sh{O}_{x'}\neq K$, vì $x'$ không phải là điểm tổng quát của $X$ \sref[0]{0.2.1.3};
do đó ta kết luận $\sh{O}_{x'}=\sh{O}_{y'}$.
Vì $X$ là một lược đồ nguyên, đẳng thức $\sh{O}_{x'}=\sh{O}_{x''}$ kéo theo $x'=x''$ \sref[I]{I.8.2.2}, điều này kết thúc chứng minh.
\end{proof}

\begin{env}[7.3.2]
\label{II.7.3.2}
Cho $A$ là một vành định giá, $K$ là trường phân thức của nó, $Y=\Spec(A)$, và $b$ là điểm tổng quát của $Y$, sao cho $\sh{O}_b=\kres(b)$ bằng $K$;
cho $f:X\to Y$ là một cấu xạ.
Ta biết \sref[I]{I.7.1.4} rằng các \emph{$Y$-tiết diện hữu tỉ} của $X$ tương ứng song ánh với các \emph{mầm} tại $b$ của những $Y$-tiết diện được xác định trong các lân cận của $b$, từ đó có một ánh xạ chính tắc
\[
\label{II.7.3.2.1}
    \Gamma_\mathrm{rat}(X/Y) \to \Gamma(f^{-1}(b)/\Spec(K))
\tag{7.3.2.1}
\]
trong đó, theo định nghĩa \sref[I]{I.3.4.5}, các phần tử của $\Gamma(f^{-1}(b)/\Spec(K))$ được đồng nhất với các \emph{điểm của $f^{-1}(b)=X\otimes_A K$ hữu tỉ trên $K$}.
Khi $f$ \emph{tách}, từ \sref[I]{I.5.4.7} suy ra ánh xạ \sref{II.7.3.2.1} là \emph{đơn ánh}, vì $Y$ là một lược đồ nguyên.

Hợp \sref{II.7.3.2.1} với ánh xạ chính tắc $\Gamma(X/Y)\to\Gamma_\mathrm{rat}(X/Y)$ \sref[I]{I.7.1.2}, ta thu được một ánh xạ chính tắc
\[
\label{II.7.3.2.2}
    \Gamma(X/Y) \to \Gamma(f^{-1}(b)/\Spec(K)).
\tag{7.3.2.2}
\]
Khi $f$ \emph{tách}, ánh xạ này cũng \emph{đơn ánh} \sref[I]{I.5.4.7}.
\end{env}

\begin{proposition}[7.3.3]
\label{II.7.3.3}
Cho $A$ là một vành định giá có trường phân thức $K$, $Y=\Spec(A)$, $b$ là điểm tổng quát của $Y$, và $f:X\to Y$ là một cấu xạ \emph{tách} và \emph{đóng}.
Khi đó ánh xạ chính tắc \sref{II.7.3.2.2} là song ánh (điều này tương đương với việc nói rằng nó \emph{toàn ánh}, và kéo theo các $Y$-tiết diện hữu tỉ của $X$ được \emph{xác định khắp nơi}).
\end{proposition}

\begin{proof}
Vậy lấy $x$ là một điểm của $f^{-1}(b)$ \emph{hữu tỉ} trên $K$.
Vì $f$ tách, cấu xạ $f^{-1}(b)\to\Spec(K)$ tương ứng với $f$ cũng tách \sref[I]{I.5.5.1}[iv], và vì mọi tiết diện của $f^{-1}(b)$ là một phép nhúng đóng \sref[I]{I.5.4.6}, nên $\{x\}$ đóng \emph{trong $f^{-1}(b)$}.
Xét tiền lược đồ con đóng rút gọn $X'$ của $X$ có không gian nền là bao đóng $\overline{\{x\}}$ của $\{x\}$ \emph{trong $X$}.
Hiển nhiên hạn chế của $f$ lên $X'$ thỏa các giả thiết của \sref{II.7.3.1}, do đó là một \emph{đẳng cấu} từ $X'$ lên $Y$, mà đẳng cấu ngược của nó chính là $Y$-tiết diện cần tìm của $X$.
\end{proof}

\begin{env}[7.3.4]
\label{II.7.3.4}
Để phát biểu hai kết quả sau, ta sẽ dùng một thuật ngữ sẽ được biện minh và thảo luận trong Chương~IV: nếu $F$ là một tập con
\oldpage[II]{144}
của một tiền lược đồ $Y$, ta định nghĩa \emph{đối chiều} của $F$ trong $Y$, ký hiệu $\codim_Y F$, là cận dưới đúng của các số nguyên $\dim(\sh{O}_z)$ khi $z$ chạy qua $F$.
\end{env}

\begin{corollary}[7.3.5]
\label{II.7.3.5}
Cho $Y$ là một tiền lược đồ Noether địa phương rút gọn, và $N$ là tập các điểm $y\in Y$ tại đó $Y$ không chính quy \sref[0]{0.4.1.4};
giả sử $\codim_Y N\geq2$.
Cho $f:X\to Y$ là một cấu xạ thuộc kiểu hữu hạn, \emph{tách} và \emph{đóng}, và $g$ là một $Y$-tiết diện hữu tỉ của $X$;
nếu $Y'$ là tập các điểm của $Y$ tại đó $g$ không được xác định (một tập \emph{đóng} \sref[I]{I.7.2.1}), thì $\codim_Y Y'\geq2$.
\end{corollary}

\begin{proof}
Chỉ cần chứng minh rằng $g$ được xác định tại mọi điểm $z\in Y$ sao cho $\dim\sh{O}_z\leq1$.
Nếu $\dim\sh{O}_z=0$, thì $z$ là điểm tổng quát của một thành phần bất khả quy của $Y$ \sref[I]{I.1.1.14}, do đó thuộc mọi tập mở trù mật khắp nơi của $Y$, và đặc biệt thuộc miền xác định của $g$.
Vậy giả sử $\dim\sh{O}_z=1$; theo giả thiết, $\sh{O}_z$ khi đó là một vành địa phương Noether chính quy, và do đó, theo \sref{II.7.1.6}, là một \emph{vành định giá rời rạc}.
Đặt $Z=\Spec(\sh{O}_z)$;
vì $U=Y\setmin Y'$ trù mật khắp nơi, $U\cap Z$ không rỗng \sref[I]{I.2.4.2};
gọi $g'$ là ánh xạ hữu tỉ từ $Z$ vào $X$ cảm sinh bởi $g$ \sref[I]{I.7.2.8};
chỉ cần chứng minh $g'$ là một \emph{cấu xạ} \sref[I]{I.7.2.9}.
Thật vậy, có thể xem $g'$ là một $Z$-tiết diện hữu tỉ của $Z$-tiền lược đồ $f^{-1}(Z)=X\times_Y Z$;
hiển nhiên cấu xạ $f^{-1}(Z)\to Z$ tương ứng với $f$ là đóng, và từ \sref[I]{I.5.5.1}[i] suy ra nó tách;
do đó từ \sref{II.7.3.3} ta kết luận rằng $g'$ được xác định khắp nơi;
vì $Z$ rút gọn và $X$ tách trên $Y$, $g'$ là một cấu xạ \sref[I]{I.7.2.2}.
\end{proof}

\begin{corollary}[7.3.6]
\label{II.7.3.6}
Cho $S$ là một tiền lược đồ Noether địa phương, và $X$, $Y$ là hai $S$-tiền lược đồ;
giả sử $Y$ rút gọn và, hơn nữa, tập $N$ gồm các điểm $y\in Y$ tại đó $Y$ không chính quy thỏa $\codim_Y N\geq2$;
cuối cùng, giả sử cấu xạ cấu trúc $X\to S$ là riêng.
Cho $f$ là một $S$-ánh xạ hữu tỉ từ $Y$ vào $X$, và $Y'$ là tập các điểm của $Y$ tại đó $f$ không được xác định;
khi đó $\codim_Y Y'\geq2$.
\end{corollary}

\begin{proof}
Ta biết \sref[I]{I.7.1.2} rằng có thể đồng nhất các $S$-ánh xạ hữu tỉ từ $Y$ vào $X$ với các $Y$-tiết diện hữu tỉ của $X\times_S Y$;
vì cấu xạ cấu trúc $X\times_S Y\to Y$ là đóng \sref{II.5.4.1}, có thể áp dụng \sref{II.7.3.5}, từ đó suy ra hệ quả.
\end{proof}

\begin{remark}[7.3.7]
\label{II.7.3.7}
Các giả thiết về $Y$ trong \sref{II.7.3.5} và \sref{II.7.3.6} đặc biệt được thỏa mãn khi $Y$ \emph{chuẩn tắc} \sref[0]{0.4.1.4}, theo \sref{II.7.1.6}.
\end{remark}

Ta có thể đặc trưng các cấu xạ đóng phổ dụng (tương ứng, các cấu xạ riêng) bằng một mệnh đề đảo của \sref{II.7.3.3}:
\begin{theorem}[7.3.8]
\label{II.7.3.8}
Cho $Y$ là một tiền lược đồ (tương ứng, một tiền lược đồ Noether địa phương), và $f:X\to Y$ là một cấu xạ tách quasi-compact (tương ứng, một cấu xạ thuộc kiểu hữu hạn).
Các điều kiện sau tương đương:
\begin{enumerate}
    \item[\rm{(a)}] $f$ đóng phổ dụng (tương ứng, riêng).
    \item[\rm{(b)}] Đối với mọi $Y$-lược đồ dạng $Y'=\Spec(A)$, trong đó $A$ là một vành định giá (tương ứng, một vành định giá rời rạc) có trường phân thức $K$, ánh xạ chính tắc
        \[
            \Hom_Y(Y',X) \to \Hom_Y(\Spec(K),X)
        \]
        tương ứng với đơn ánh chính tắc $A\to K$ là toàn ánh (tương ứng, song ánh).
\oldpage[II]{145}
  \item[\rm{(c)}] Đối với mọi $Y$-lược đồ dạng $Y'=\Spec(A)$, trong đó $A$ là một vành định giá (tương ứng, một vành định giá rời rạc), ánh xạ chính tắc \sref{II.7.3.2.2} liên quan đến $Y'$-tiền lược đồ $X_{(Y')}$ là toàn ánh (tương ứng, song ánh).
\end{enumerate}
\end{theorem}

\begin{proof}
Sự tương đương giữa (b) và (c) suy ra ngay từ \sref[I]{I.3.3.14};
(a) kéo theo (c), bởi vì trong mỗi trường hợp, (a) kéo theo $f_{(Y')}$ tách \sref[I]{I.5.5.1}[iv] và đóng, và chỉ cần áp dụng \sref{II.7.3.3}.
Còn phải chứng minh rằng (b) kéo theo (a).

Trước hết, xét trường hợp $Y$ tùy ý còn $f$ tách và quasi-compact.
Nếu $f$ thỏa mãn điều kiện~(b), thì
\[
  f_{(Y'')}:X_{(Y'')}\longrightarrow Y''
\]
cũng thỏa mãn điều kiện ấy, trong đó $Y''$ là một $Y$-tiền lược đồ tùy ý, nhờ sự tương đương giữa (b) và (c) cùng với hệ thức
\[
  X_{(Y'')}\times_{Y''}Y'=X\times_Y Y'
\]
đối với mọi cấu xạ $Y'\to Y''$ \sref[I]{I.3.3.9.1};
hơn nữa, vì $f_{(Y'')}$ tách và quasi-compact khi $f$ có các tính chất ấy (\sref[I]{I.5.5.1}[iv] và \sref[I]{I.6.6.4}[iii]), ta quy về việc chứng minh rằng (b) kéo theo $f$ \emph{đóng}.
Để làm điều đó, chỉ cần kiểm tra điều kiện~(b) của \sref{II.7.2.1}.
Vậy lấy $x\in X$ và một chuyên biệt hóa $y'$ của $y=f(x)$, với $y'\ne y$;
theo \sref{II.7.1.4}, tồn tại một lược đồ $Y'$, là phổ của một vành định giá, và một cấu xạ tách $g:Y'\to Y$ sao cho, nếu $a$ ký hiệu điểm đóng và $b$ ký hiệu điểm tổng quát của $Y'$, thì $g(a)=y'$, $g(b)=y$, và tồn tại một $\kres(y)$-đồng cấu $\kres(x)\to\kres(b)$.
Đồng cấu cuối cùng này tương ứng chính tắc với một $Y$-cấu xạ $\Spec(\kres(b))\to X$ \sref[I]{I.2.4.6}; do đó, từ (b) suy ra tồn tại một $Y$-cấu xạ $h:Y'\to X$ tương ứng với cấu xạ vừa nêu.
Khi đó $h(b)=x$;
nếu đặt $h(a)=x'$, thì $x'$ là một chuyên biệt hóa của $x$, và
\[
  f(x')=f(h(a))=g(a)=y'.
\]

Bây giờ, nếu $Y$ Noether địa phương và $f$ thuộc kiểu hữu hạn, giả thiết~(b) trước hết kéo theo $f$ \emph{tách}, theo \sref{II.7.2.3}, vì cấu xạ đường chéo $X\to X\times_Y X$ là quasi-compact \sref{II.7.2.4}.
Hơn nữa, để chứng minh $f$ riêng, chỉ cần chứng minh rằng
\[
  f_{(Y'')}:X_{(Y'')}\longrightarrow Y''
\]
là \emph{đóng} đối với mọi $Y$-tiền lược đồ $Y''$ \emph{thuộc kiểu hữu hạn}, theo \sref{II.5.6.3}.
Vì khi đó $Y''$ Noether địa phương, ta có thể lặp lại lập luận của trường hợp thứ nhất, lấy $Y'$ là phổ của một vành định giá rời rạc và áp dụng \sref{II.7.1.9} thay cho \sref{II.7.1.4}.
\end{proof}

\begin{remarks}[7.3.9]
\label{II.7.3.9}
\begin{enumerate}
  \item[\rm{(i)}] Khi $Y$ là một tiền lược đồ tùy ý và $f$ là một cấu xạ tách, để $f$ đóng phổ dụng, \emph{chỉ cần} điều kiện~(b) hoặc điều kiện~(c) được thỏa mãn đối với các vành định giá $A$ \emph{đầy đủ} có trường thặng dư \emph{đóng đại số};
  thật vậy, điều này suy ra từ chứng minh trên và \sref{II.7.1.4}.

  \item[\rm{(ii)}] Từ tiêu chuẩn~(c) của \sref{II.7.3.8} suy ra một chứng minh mới, gần với các phương pháp cổ điển hơn, cho sự kiện rằng một cấu xạ xạ ảnh $X\to Y$ là đóng \sref{II.5.5.3}.
  Thật vậy, có thể giả sử $Y$ affine, và do đó đồng nhất $X$ với một tiền lược đồ con đóng của một phân thớ xạ ảnh $\bb{P}_Y^n$ \sref{II.5.3.3};
  để chứng minh $X\to Y$ đóng, chỉ cần kiểm tra điều ấy đối với cấu xạ cấu trúc $\bb{P}_Y^n\to Y$, và tiêu chuẩn~(c) của \sref{II.7.3.8}, cùng với \sref{II.4.1.3.1}, cho thấy ta quy về việc chứng minh sự kiện sau:
  \emph{nếu $Y$ là phổ của một vành định giá $A$ có trường phân thức $K$, thì mọi điểm của $\bb{P}_Y^n$ có giá trị trong $K$ đều xuất phát, bằng cách hạn chế tới điểm tổng quát của $Y$, từ một điểm của $\bb{P}_Y^n$ có giá trị trong $A$.}
  Mà mọi $\sh{O}_Y$-môđun khả nghịch đều tầm thường \sref[I]{I.2.4.8};
  do đó, từ \sref{II.4.2.6} suy ra rằng một điểm của $\bb{P}_Y^n$ có giá trị trong $K$ được đồng nhất với một lớp các hệ $(\zeta c_0,\zeta c_1,\ldots,\zeta c_n)$ trong $K$, trong đó $\zeta\ne0$ và các $c_i\in K$ không đồng thời bằng không.
  Bằng cách nhân các $c_i$ với một phần tử của $A$ có
\oldpage[II]{146}
  định giá thích hợp, ta có thể giả sử mọi $c_i$ đều thuộc $A$ và ít nhất một trong số chúng khả nghịch.
  Nhưng khi đó, theo \sref{II.4.2.6}, hệ $(c_0,\ldots,c_n)$ cũng xác định một điểm của $\bb{P}_Y^n$ có giá trị trong $A$, điều này chứng minh khẳng định.

  \item[\rm{(iii)}] Các tiêu chuẩn \sref{II.7.2.3} và \sref{II.7.3.8} đặc biệt thuận tiện khi ta xem việc cho một $Y$-tiền lược đồ $X$ tương đương với việc cho hàm tử gán cho mỗi $Y$-tiền lược đồ $Y'$ tập hợp
  \[
    X(Y')=\Hom_Y(Y',X);
  \]
  chẳng hạn, các tiêu chuẩn này sẽ cho phép ta chứng minh rằng, dưới một số điều kiện, các ``lược đồ Picard'' là riêng.
\end{enumerate}
\end{remarks}

\begin{corollary}[7.3.10]
\label{II.7.3.10}
Cho $Y$ là một lược đồ nguyên (tương ứng, một lược đồ nguyên Noether địa phương), $X$ là một lược đồ nguyên và $f:X\to Y$ là một cấu xạ trội.
\begin{enumerate}
  \item[\rm{(i)}] Nếu $f$ là quasi-compact và đóng phổ dụng, thì mọi vành định giá có trường phân thức là trường $R(X)$ của các hàm hữu tỉ trên $X$, đồng thời chi phối một vành địa phương của $Y$, cũng chi phối một vành địa phương của $X$.
  \item[\rm{(ii)}] Ngược lại, giả sử $f$ thuộc kiểu hữu hạn và tính chất nêu trong~(i) được thỏa mãn đối với mọi vành định giá (tương ứng, mọi vành định giá rời rạc) có $R(X)$ làm trường phân thức.
  Khi đó $f$ là một cấu xạ riêng.
\end{enumerate}
\end{corollary}

\begin{proof}
Trước hết, lưu ý rằng trong mọi trường hợp các giả thiết đều kéo theo $f$ là một cấu xạ tách \sref[I]{I.5.5.9}.
\begin{enumerate}
  \item[\rm{(i)}] Đặt $K=R(Y)$, $L=R(X)$; lấy một điểm $y$ của $Y$ và một vành định giá $A$ chi phối $\sh{O}_y$, có $L$ làm trường phân thức.
  Đồng cấu đơn ánh $\sh{O}_y\to A$ khi đó xác định một cấu xạ $h:Y'=\Spec(A)\to Y$ \sref[I]{I.2.4.4} sao cho $h(a)=y$, trong đó $a$ ký hiệu điểm đóng của $Y'$.
  Hơn nữa, nếu $\eta$ là điểm tổng quát của $Y$, cũng là điểm tổng quát của $\Spec(\sh{O}_y)$, thì $h(b)=\eta$, trong đó $b$ ký hiệu điểm tổng quát của $Y'$ (vì $K\subset L$ theo giả thiết).
  Nếu $\xi$ là điểm tổng quát của $X$, thì theo giả thiết $\kres(\xi)=\kres(b)=L$, do đó có một $Y$-cấu xạ $g:\Spec(L)\to X$ sao cho $g(b)=\xi$.
  Theo \sref{II.7.3.8}, $g$ mở rộng thành một $Y$-cấu xạ $g':Y'\to X$.
  Nếu $x=g'(a)$, hiển nhiên $A$ chi phối $\sh{O}_x$.

  \item[\rm{(ii)}] Vì bài toán có tính địa phương trên $Y$, có thể giả sử $Y$ là affine (tương ứng, affine và Noether).
  Vì $f$ thuộc kiểu hữu hạn, trong cả hai trường hợp đều có thể áp dụng bổ đề Chow \sref{II.5.6.1}.
  Do đó tồn tại một cấu xạ xạ ảnh $p:P\to Y$, một phép nhúng $j:X'\to P$ và một cấu xạ xạ ảnh, toàn ánh, song hữu tỉ $g:X'\to X$, với $X'$ nguyên, sao cho biểu đồ
  \[
    \xymatrix{
      P \ar[d]_p
      & X' \ar[l]_j \ar[d]^g
      \\ Y
      & X \ar[l]^f
    }
  \]
  giao hoán.
  Chỉ cần chứng minh rằng $j$ là một phép nhúng \emph{đóng}; khi đó $f\circ g=p\circ j$ là một cấu xạ xạ ảnh, do đó là một cấu xạ riêng, và vì $g$ toàn ánh nên $f$ cũng là một cấu xạ riêng \sref{II.5.4.3}.
  Gọi $Z$ là tiền lược đồ con đóng rút gọn của $P$ có không gian nền là $\overline{j(X')}$ \sref[I]{I.5.2.1}.
  Vì $X'$ nguyên, $j$ phân tích thành $i\circ h$, trong đó $i:Z\to P$ là phép nhúng chính tắc, $h:X'\to Z$ là một phép nhúng mở trội \sref[I]{I.5.2.3}, và $Z$ nguyên.
\oldpage[II]{147}
  Hơn nữa, $Z$ xạ ảnh trên $Y$; ta thấy rằng có thể thu hẹp về trường hợp $P$ \emph{nguyên}, còn $j$ \emph{trội} và \emph{song hữu tỉ}, và khi đó mọi việc quy về chứng minh $j$ \emph{toàn ánh}.
  Thật vậy, lấy $z\in P$.
  Khi đó $\sh{O}_z$ là một vành địa phương nguyên (tương ứng, nguyên và Noether) có trường phân thức
  \[
    L=R(P)=R(X')=R(X).
  \]
  Có thể chỉ xét trường hợp $z$ không phải là điểm tổng quát của $P$.
  Khi đó, theo \sref{II.7.1.2} và \sref{II.7.1.7}, tồn tại một vành định giá (tương ứng, một vành định giá rời rạc) $A$ chi phối $\sh{O}_z$ và có $L$ làm trường phân thức.
  \emph{A fortiori}, $A$ chi phối $\sh{O}_y$, với $y=p(z)$; theo giả thiết, tồn tại một $x\in X$ sao cho $A$ chi phối $\sh{O}_x$.
  Vì $g$ là một cấu xạ riêng, phần đầu của chứng minh cho thấy $A$ cũng chi phối $\sh{O}_{x'}$ với một $x'\in X'$.
  Suy ra $\sh{O}_z$ và $\sh{O}_{j(x')}=\sh{O}_{x'}$ liên hệ với nhau theo nghĩa của \sref[I]{I.8.1.4}; vì $P$ là một lược đồ, điều này kéo theo $z=j(x')$ \sref[I]{I.8.2.2}, và chứng minh kết thúc.
\end{enumerate}
\end{proof}

\begin{corollary}[7.3.11]
\label{II.7.3.11}
Cho $X$ và $Y$ là hai lược đồ nguyên, và $f:X\to Y$ là một cấu xạ trội, quasi-compact và đóng phổ dụng.
Giả sử thêm rằng $Y$ là affine, với vành (nguyên) $B$.
Khi đó $\Gamma(X,\sh{O}_X)$ đẳng cấu một cách chính tắc với một vành con của bao đóng nguyên của $B$ trong $R(X)$.
\end{corollary}

\begin{proof}
Thật vậy, theo \sref[I]{I.8.2.1.1}, $\Gamma(X,\sh{O}_X)$ được đồng nhất với giao của các $\sh{O}_x$ khi $x\in X$;
theo \sref{II.7.3.10}, \sref{II.7.1.2} và \sref{II.7.1.3}, $\Gamma(X,\sh{O}_X)$ do đó được chứa trong giao của các vành định giá chứa $B$ và có $R(X)$ làm trường phân thức;
kết luận khi đó suy ra từ \sref{II.7.1.3}.
\end{proof}

\begin{remarks}[7.3.12]
\label{II.7.3.12}
Dưới các giả thiết của \sref{II.7.3.11}, và khi giả sử $R(X)$ là một mở rộng hữu hạn sinh của $R(Y)$, trong nhiều trường hợp có thể kết luận rằng $\Gamma(X,\sh{O}_X)$ là một \emph{môđun hữu hạn sinh} trên vành $B=\Gamma(Y,\sh{O}_Y)$.
Điều này đúng, chẳng hạn, khi $B$ là một \emph{đại số hữu hạn sinh trên một trường}, vì khi đó ta biết rằng bao đóng nguyên của $B$ trong một mở rộng hữu hạn sinh của trường phân thức của nó là một $B$-môđun hữu hạn sinh (\cite[t.~I, p.~267, th.~9]{I-13});
kết luận khi đó suy ra từ \sref{II.7.3.11} và từ việc $B$ là Noether.

Đặc biệt, \emph{một lược đồ $X$ riêng và affine trên một trường $K$ là hữu hạn}.
Thật vậy, theo \sref{II.1.6.4}, \sref{II.5.4.6} và \sref[I]{I.6.4.4}[(c)], có thể quy về trường hợp $X$ \emph{rút gọn}.
Hơn nữa, chỉ cần chứng minh rằng mỗi tiền lược đồ con đóng của $X$ có không gian nền là một thành phần bất khả quy của $X$ (số các thành phần này là hữu hạn) đều hữu hạn trên $K$; do đó, xét đến \sref{II.5.4.5}, cuối cùng ta quy về trường hợp $X$ \emph{nguyên}.
Nhưng khi đó kết quả suy ra từ các nhận xét ở trên.

Trong chương~III, ta sẽ chứng minh lại mệnh đề trên bằng các phương pháp khác và như một hệ quả của những kết quả tổng quát hơn, bằng cách chỉ ra rằng, nếu $f:X\to Y$ là một cấu xạ riêng và $Y$ là Noether địa phương, thì $f_*(\sh{F})$ là coherent đối với mọi $\sh{O}_X$-môđun coherent $\sh{F}$ \sref[III]{III.4.4.2}.

Cuối cùng, lưu ý rằng tiêu chuẩn~\sref{II.7.3.10} được lấy làm \emph{định nghĩa} của các cấu xạ riêng trong hình học đại số cổ điển.
Chúng tôi chỉ nêu điều này ở đây như một nhận xét, vì tiêu chuẩn~\sref{II.7.3.8} dường như dễ áp dụng hơn trong mọi ứng dụng mà chúng tôi biết.
\end{remarks}

\oldpage[II]{148}

\subsection{Đường cong đại số và trường hàm chiều 1}
\label{subsection:II.7.4}

Mục đích của mục này là chỉ ra cách diễn đạt khái niệm cổ điển về đường cong đại số bằng ngôn ngữ lược đồ (như được trình bày, chẳng hạn, trong sách của C.~Chevalley \cite{II-23}).
Trong toàn bộ mục này, \emph{ta dùng $k$ để chỉ một trường; mọi lược đồ được xét đều là $k$-lược đồ thuộc kiểu hữu hạn, và mọi cấu xạ đều là $k$-cấu xạ}.

\begin{proposition}[7.4.1]
\label{II.7.4.1}
Cho $X$ là một tiền lược đồ thuộc kiểu hữu hạn trên $k$ (do đó là Noether);
gọi $x_i$ ($1\leq i\leq n$) là các điểm tổng quát của các thành phần bất khả quy $X_i$ của $X$, và đặt $K_i=\kres(x_i)$ ($1\leq i\leq n$).
Khi đó các điều kiện sau tương đương:
\begin{enumerate}
    \item[\rm{(a)}] Mỗi $K_i$ là một trường mở rộng của $k$ có bậc siêu việt bằng $1$.
    \item[\rm{(b)}] Với mọi điểm đóng $x$ của $X$, vành địa phương $\sh{O}_x$ có chiều $1$ \sref{II.7.1.5}.
    \item[\rm{(c)}] Các tập con đóng bất khả quy của $X$ khác với các $X_i$ chính là các điểm đóng của $X$.
\end{enumerate}
\end{proposition}

\begin{proof}
Vì $X$ là quasi-compact, mọi tập con đóng bất khả quy $F$ của $X$ đều chứa một điểm đóng \sref[0]{0.2.1.3}.
Theo \sref[I]{I.2.4.2}, có một tương ứng song ánh giữa các iđêan nguyên tố của $\sh{O}_x$ và các tập con đóng bất khả quy của $X$ chứa $x$ \sref[I]{I.1.1.14};
sự tương đương giữa (b) và (c) suy ra ngay từ đó.
Mặt khác, nếu $\mathfrak{p}_\alpha$ ($1\leq\alpha\leq r$) là các iđêan nguyên tố tối tiểu của vành địa phương Noether $\sh{O}_x$, thì các vành địa phương $\sh{O}_x/\mathfrak{p}_\alpha$ là các miền nguyên, và các trường phân thức của chúng chính là những $K_i$ sao cho $x\in X_i$.
Hơn nữa, ta biết (\cite[p.~4-06, th.~2]{I-1}) rằng chiều của một $k$-đại số địa phương nguyên thuộc kiểu hữu hạn bằng bậc siêu việt trên $k$ của trường phân thức của nó.
Cuối cùng, chiều của $\sh{O}_x$ là cận trên đúng của các chiều của $\sh{O}_x/\mathfrak{p}_\alpha$;
nhưng điều kiện~(a) kéo theo rằng các chiều này đều bằng $1$, nên (a) kéo theo (b);
ngược lại, nếu $\sh{O}_x$ có chiều $1$, thì không một $\mathfrak{p}_\alpha$ nào có thể bằng iđêan tối đại của $\sh{O}_x$, bởi nếu không thì $\sh{O}_x$ sẽ có chiều $0$;
do đó mỗi $\sh{O}_x/\mathfrak{p}_\alpha$ có chiều $1$, điều này chứng tỏ (b) kéo theo (a).
\end{proof}

Ta lưu ý rằng, dưới các điều kiện của \sref{II.7.4.1}, tập $X$ hoặc là \emph{rỗng} hoặc là \emph{vô hạn}, như suy ra ngay từ \sref[I]{I.6.4.4}.

\begin{definition}[7.4.2]
\label{II.7.4.2}
Một \emph{đường cong đại số trên $k$} là một \emph{lược đồ} đại số không rỗng trên $k$ thỏa các điều kiện của \sref{II.7.4.1}.
\end{definition}

Trong ngôn ngữ về chiều, sẽ được giới thiệu ở Chương~IV, điều này được diễn đạt bằng cách nói rằng một đường cong đại số trên $k$ là một $k$-lược đồ đại số không rỗng \emph{mà mọi thành phần bất khả quy đều có chiều $1$}.

Ta lưu ý rằng, nếu $X$ là một đường cong đại số trên $k$, thì các tiền lược đồ con đóng rút gọn $X_i$ ($1\leq i\leq n$) của $X$ có không gian nền là các thành phần bất khả quy của $X$ cũng là các đường cong đại số trên $k$.

\begin{corollary}[7.4.3]
\label{II.7.4.3}
Cho $X$ là một đường cong đại số bất khả quy.
Điểm không đóng duy nhất của $X$ là điểm tổng quát của nó.
Các tập con đóng của $X$ khác $X$ chính là các tập hữu hạn gồm các điểm đóng;
chúng cũng là những tập con duy nhất của $X$ không trù mật khắp nơi.
\end{corollary}

\begin{proof}
Nếu một điểm $x\in X$ không đóng, thì bao đóng của nó trong $X$ là một tập con đóng bất khả quy của $X$, do đó theo \sref{II.7.4.1} nhất thiết bằng toàn bộ $X$; vì vậy $x$ là điểm tổng quát của $X$.
Một tập con đóng $F$ của $X$ khác $X$ không thể chứa
\oldpage[II]{149}
điểm tổng quát của $X$, nên mọi điểm của nó đều đóng (trong $X$ và, \emph{a fortiori}, trong $F$);
bằng cách xét tiền lược đồ con đóng rút gọn của $X$ có $F$ làm không gian nền \sref[I]{I.5.2.1}, từ \sref[I]{I.6.2.2} suy ra rằng $F$ hữu hạn và rời rạc.
Do đó, bao đóng trong $X$ của mọi tập con vô hạn của $X$ nhất thiết bằng $X$.
\end{proof}

Khi $X$ là một đường cong đại số bất kỳ, bằng cách áp dụng \sref{II.7.4.3} cho các thành phần bất khả quy của $X$, ta thấy rằng những điểm không đóng duy nhất của $X$ là các điểm tổng quát của các thành phần đó.

\begin{corollary}[7.4.4]
\label{II.7.4.4}
Cho $X$ và $Y$ là hai đường cong đại số bất khả quy trên $k$, và $f:X\to Y$ là một $k$-cấu xạ.
Để $f$ trội, điều kiện cần và đủ là $f^{-1}(y)$ hữu hạn với mọi $y\in Y$.
\end{corollary}

\begin{proof}
Thật vậy, nếu $f$ không trội, thì theo \sref{II.7.4.3}, $f(X)$ nhất thiết là một tập con \emph{hữu hạn} của $Y$; do đó không thể có $f^{-1}(y)$ hữu hạn đối với mọi điểm của $Y$, vì nếu không thì $X$ sẽ hữu hạn, điều này là vô lý \sref{II.7.4.1}.
Ngược lại, nếu $f$ trội, thì với mọi $y\in Y$ khác điểm tổng quát $\eta$ của $Y$, tập $f^{-1}(y)$ đóng trong $X$, vì $\{y\}$ đóng trong $Y$ \sref{II.7.4.3};
mặt khác, theo giả thiết, $f^{-1}(y)$ không chứa điểm tổng quát $\xi$ của $X$, và do đó hữu hạn theo \sref{II.7.4.3}.
Cuối cùng, để thấy rằng khi $f$ trội thì $f^{-1}(\eta)$ hữu hạn, ta lưu ý rằng thớ $f^{-1}(\eta)$ là một lược đồ bất khả quy thuộc kiểu hữu hạn trên $\kres(\eta)$ và có điểm tổng quát $\xi$ (\sref[I]{I.6.3.9} và \sref[I]{I.6.4.11}).
Vì $\kres(\xi)$ và $\kres(\eta)$ là các mở rộng hữu hạn sinh của $k$, đều có bậc siêu việt $1$, nên $\kres(\xi)$ nhất thiết là một mở rộng hữu hạn của $\kres(\eta)$; do đó $\xi$ đóng trong $f^{-1}(\eta)$ \sref[I]{I.6.4.2}, và $f^{-1}(\eta)$ chỉ gồm điểm duy nhất $\xi$.
\end{proof}

Ta sẽ thấy ở Chương~III rằng một cấu xạ \emph{riêng} $f:X\to Y$ giữa các tiền lược đồ Noether, sao cho $f^{-1}(y)$ hữu hạn với mọi $y\in Y$, thì nhất thiết \emph{hữu hạn};
do đó từ \sref{II.7.4.4} suy ra rằng một cấu xạ trội riêng từ một đường cong đại số bất khả quy vào một đường cong đại số là \emph{hữu hạn}.

\begin{corollary}[7.4.5]
\label{II.7.4.5}
Cho $X$ là một đường cong đại số trên $k$.
Điều kiện cần và đủ để $X$ chính quy là $X$ chuẩn tắc, hay tương đương, các vành địa phương tại các điểm đóng của nó là các vành định giá rời rạc.
\end{corollary}

\begin{proof}
Điều này suy ra ngay từ điều kiện~(b) của \sref{II.7.4.1} và \sref{II.7.1.6}.
\end{proof}

\begin{corollary}[7.4.6]
\label{II.7.4.6}
Cho $X$ là một đường cong đại số rút gọn và $\sh{A}$ là một $\sh{R}(X)$-đại số coherent rút gọn;
khi đó bao đóng nguyên $X'$ của $X$ đối với $\sh{A}$ \sref{II.6.3.4} là một đường cong đại số chuẩn tắc, và cấu xạ chính tắc $X'\to X$ là hữu hạn.
\end{corollary}

\begin{proof}
Việc $X'\to X$ là hữu hạn suy ra từ \sref{II.6.3.10};
do đó $X'$ là một $k$-lược đồ đại số;
hơn nữa, nếu $x_i$ ($1\leq i\leq n$) là các điểm tổng quát của các thành phần bất khả quy của $X$, và $x'_j$ ($1\leq j\leq m$) là các điểm tổng quát của các thành phần bất khả quy của $X'$, thì mỗi $\kres(x'_j)$ là một mở rộng đại số hữu hạn của một trong các $\kres(x_i)$ \sref{II.6.3.6}, nên có bậc siêu việt $1$ trên $k$.
Vì thế $X'$ quả là một đường cong đại số trên $k$ và, hơn nữa, ta biết rằng $X'$ là tổng của một số hữu hạn lược đồ nguyên và chuẩn tắc (\sref{II.6.3.6} và \sref{II.6.3.7}).
\end{proof}

\begin{env}[7.4.7]
\label{II.7.4.7}
Ta nói một đường cong đại số $X$ trên $k$ là \emph{đầy đủ} nếu nó \emph{riêng} trên $k$.
\end{env}

\begin{corollary}[7.4.8]
\label{II.7.4.8}
Điều kiện cần và đủ để một đường cong đại số rút gọn $X$ trên $k$ đầy đủ là chuẩn hóa $X'$ của nó cũng đầy đủ.
\end{corollary}

\oldpage[II]{150}
\begin{proof}
Thật vậy, cấu xạ chính tắc $f:X'\to X$ là hữu hạn \sref{II.7.4.6}, do đó là cấu xạ riêng \sref{II.6.1.11} và toàn ánh \sref{II.6.3.8};
nếu $g:X\to\Spec(k)$ là cấu xạ cấu trúc, thì $g$ là cấu xạ riêng khi và chỉ khi $g\circ f$ là cấu xạ riêng, như suy ra từ \sref{II.5.4.2}[(ii)] và \sref{II.5.4.3}[(ii)], vì theo giả thiết $g$ là cấu xạ tách.
\end{proof}

\begin{proposition}[7.4.9]
\label{II.7.4.9}
Cho $X$ là một đường cong đại số \emph{chuẩn tắc} trên $k$ và $Y$ là một $k$-lược đồ đại số \emph{riêng}.
Khi đó mọi $k$-ánh xạ hữu tỉ từ $X$ vào $Y$ đều xác định khắp nơi, hay nói cách khác, là một cấu xạ.
\end{proposition}

\begin{proof}
Theo \sref{II.7.3.7}, tại những điểm $x\in X$ mà một ánh xạ như vậy không xác định, chiều của $\sh{O}_x$ phải là $\geq2$; do đó tập hợp các điểm ấy là rỗng;
khẳng định cuối cùng suy ra từ \sref[I]{I.7.2.3}.
\end{proof}

\begin{corollary}[7.4.10]
\label{II.7.4.10}
Một đường cong đại số chuẩn tắc trên $k$ là quasi-xạ ảnh trên $k$.
\end{corollary}

\begin{proof}
Vì $X$ là tổng của một số hữu hạn các đường cong đại số nguyên và chuẩn tắc \sref{II.6.3.8}, ta có thể quy về trường hợp $X$ nguyên \sref{II.5.3.6}.
Vì $X$ là quasi-compact, nó được phủ bởi một số hữu hạn các mở affine $U_i$ ($1\leq i\leq n$), và vì mỗi $U_i$ thuộc kiểu hữu hạn trên $k$, với mỗi $i$ tồn tại một số nguyên $n_i$ cùng một $k$-phép nhúng $f_i:U_i\to\bb{P}_k^{n_i}$ \sref{II.5.3.3} và \sref{II.5.3.4}[(i)].
Vì $U_i$ trù mật trong $X$, từ \sref{II.7.4.9} suy ra $f_i$ mở rộng thành một $k$-cấu xạ $g_i:X\to\bb{P}_k^{n_i}$; do đó ta thu được một $k$-cấu xạ $g=(g_1,\ldots,g_n)_k$ từ $X$ vào tích $P$ của các $\bb{P}_k^{n_i}$ trên $k$.
Hơn nữa, với mỗi $i$, vì hạn chế của $g_i$ trên $U_i$ là một phép nhúng, hạn chế của $g$ trên $U_i$ cũng là một phép nhúng \sref[I]{I.5.3.14}.
Vì các $U_i$ phủ $X$ và $g$ là một cấu xạ tách \sref[I]{I.5.5.1}[(v)], $g$ là một phép nhúng từ $X$ vào $P$ \sref[I]{I.8.2.8}.
Vì cấu xạ Segre \sref{II.4.3.3} cho một phép nhúng của $P$ vào $\bb{P}_k^N$, điều này hoàn tất chứng minh rằng $X$ là quasi-xạ ảnh.
\end{proof}

\begin{corollary}[7.4.11]
\label{II.7.4.11}
Mọi đường cong đại số chuẩn tắc $X$ đều đẳng cấu với lược đồ cảm sinh trên một mở trù mật khắp nơi của một đường cong đại số chuẩn tắc và đầy đủ $\widehat{X}$; đường cong $\widehat{X}$ này được xác định duy nhất đến một đẳng cấu duy nhất.
\end{corollary}

\begin{proof}
Nếu $X_1$ và $X_2$ là hai đường cong chuẩn tắc và đầy đủ, thì từ \sref{II.7.4.9} suy ra ngay rằng mọi đẳng cấu từ một mở $U_1$ trù mật trong $X_1$ lên một mở $U_2$ trù mật trong $X_2$ đều mở rộng duy nhất thành một đẳng cấu từ $X_1$ lên $X_2$;
do đó có khẳng định duy nhất.
Để chứng minh sự tồn tại của $\widehat{X}$, chỉ cần lưu ý rằng có thể xem $X$ là một lược đồ con của $\bb{P}_k^n$ \sref{II.7.4.10}.
Gọi $\overline{X}$ là bao đóng của $X$ trong $\bb{P}_k^n$ \sref[I]{I.9.5.11};
vì $X$ là lược đồ cảm sinh bởi $\overline{X}$ trên một mở trù mật trong $\overline{X}$ \sref[I]{I.9.5.10}, các điểm tổng quát $x_i$ của các thành phần bất khả quy của $X$ cũng là các điểm tổng quát của các thành phần bất khả quy của $\overline{X}$, và các $\kres(x_i)$ là như nhau đối với hai lược đồ; do đó, theo \sref{II.7.4.1}, $\overline{X}$ là một đường cong đại số trên $k$, rút gọn \sref[I]{I.9.5.9} và xạ ảnh trên $k$ \sref{II.5.5.1}, nên đầy đủ \sref{II.5.5.3}.
Ta lấy $\widehat{X}$ là \emph{chuẩn hóa} của $\overline{X}$, nó vẫn đầy đủ \sref{II.7.4.8};
hơn nữa, nếu $h:\widehat{X}\to\overline{X}$ là cấu xạ chính tắc, thì hạn chế của $h$ trên $h^{-1}(X)$ là một đẳng cấu lên $X$, vì $X$ chuẩn tắc \sref{II.6.3.4}; và vì $h^{-1}(X)$ chứa các điểm tổng quát của các thành phần bất khả quy của $\widehat{X}$ \sref{II.6.3.8}, nó trù mật trong $\widehat{X}$, điều này hoàn tất chứng minh.
\end{proof}

\oldpage[II]{151}
\begin{remark}[7.4.12]
\label{II.7.4.12}
Trong Chương~V, ta sẽ chứng minh rằng kết luận của \sref{II.7.4.10} vẫn đúng mà không cần giả sử đường cong là chuẩn tắc (thậm chí không cần rút gọn);
ta cũng sẽ chứng minh rằng điều kiện cần và đủ để một đường cong đại số (rút gọn hoặc không) là affine là các thành phần bất khả quy của nó (được lấy với cấu trúc rút gọn) không đầy đủ.
\end{remark}

\begin{corollary}[7.4.13]
\label{II.7.4.13}
Cho $X$ là một đường cong chuẩn tắc bất khả quy trên $k$, với trường hàm $K=R(X)$, và $Y$ là một đường cong nguyên đầy đủ trên $k$, với trường hàm $L=R(Y)$.
Khi đó có một tương ứng song ánh chính tắc giữa các $k$-cấu xạ \emph{trội} $X\to Y$ và các $k$-đơn cấu $L\to K$.
\end{corollary}

\begin{proof}
Theo \sref{II.7.4.9}, các $k$-ánh xạ hữu tỉ từ $X$ vào $Y$ được đồng nhất với các $k$-cấu xạ $u:X\to Y$.
Vì các cấu xạ trội $u:X\to Y$ được đặc trưng bởi điều kiện $u(x)=y$ (trong đó $x$ và $y$ lần lượt ký hiệu các điểm tổng quát của $X$ và $Y$), hệ quả suy ra từ các nhận xét này và từ \sref[I]{I.7.1.13}.
\end{proof}

\begin{env}[7.4.14]
\label{II.7.4.14}
Có thể làm chính xác hơn kết quả của \sref{II.7.4.13} trong trường hợp lấy $Y$ là \emph{đường thẳng xạ ảnh} $\bb{P}_k^1=\Proj(k[T_0,T])$, trong đó $T_0$ và $T$ là hai biến bất định.
Khi đó $Y$ là một lược đồ nguyên \sref{II.2.4.4}, và lược đồ cảm sinh trên tập mở $D_+(T_0)$ của $Y$ đẳng cấu với $\Spec(k[T])$ \sref{II.2.3.6}; do đó điểm tổng quát của $Y$ là iđêan $(0)$ của $k[T]$, và trường hàm của $Y$ là $k(T)$, điều này chứng tỏ $Y$ là một đường cong đại số đầy đủ trên $k$.
Hơn nữa, iđêan nguyên tố thuần nhất duy nhất của $S=k[T_0,T]$ chứa $T_0$ và khác $S_+$ là iđêan chính $(T_0)$; vì vậy, phần bù của $D_+(T_0)$ trong $Y=\bb{P}_k^1$ chỉ gồm \emph{một điểm đóng}, gọi là «điểm ở vô cực», mà ta ký hiệu là $\infty$ (về một nghiên cứu tổng quát các mối liên hệ giữa phân thớ vectơ và phân thớ xạ ảnh, xem \sref{subsection:II.8.4}).
Với các ký hiệu này:
\end{env}

\begin{corollary}[7.4.15]
\label{II.7.4.15}
Cho $X$ là một đường cong chuẩn tắc bất khả quy trên $k$, với trường hàm $K=R(X)$.
Khi đó có một tương ứng song ánh chính tắc giữa tập $K$ và tập các cấu xạ $u$ từ $X$ vào $\bb{P}_k^1$ khác với cấu xạ hằng có giá trị $\infty$.
Một cấu xạ như vậy là trội khi và chỉ khi phần tử tương ứng của $K$ siêu việt trên $k$.
\end{corollary}

\begin{proof}
Mệnh đề này suy ra ngay từ \sref{II.7.4.9} và bổ đề sau:
\begin{lemma}[7.4.15.1]
\label{II.7.4.15.1}
Cho $X$ là một tiền lược đồ nguyên trên $k$, và đặt $K=R(X)$ là trường hàm của nó.
Khi đó có một tương ứng song ánh chính tắc giữa tập $K$ và tập các ánh xạ hữu tỉ $u$ từ $X$ vào $\bb{P}_k^1$ khác với cấu xạ hằng có giá trị $\infty$.
Một ánh xạ hữu tỉ như vậy là trội khi và chỉ khi phần tử tương ứng của $K$ siêu việt trên $k$.
\end{lemma}
Trước hết, các ánh xạ hữu tỉ từ $X$ vào $\bb{P}_k^1$ tương ứng song ánh với các điểm của $\bb{P}_k^1$ có giá trị trong mở rộng $K$ của $k$ \sref[I]{I.7.1.12}.
Nếu một điểm như vậy nằm tại điểm tổng quát của $\bb{P}_k^1$ \sref[I]{I.3.4.5}, thì ánh xạ hữu tỉ tương ứng hiển nhiên là trội.
Trong trường hợp ngược lại, vì mọi điểm của $\bb{P}_k^1$ khác điểm tổng quát đều đóng \sref{II.7.4.3}, ảnh của miền xác định $U$ của $u$ qua cấu xạ duy nhất $U\to\bb{P}_k^1$ thuộc lớp $u$ \sref[I]{I.7.2.2} chỉ gồm một điểm đóng $y$ của $\bb{P}_k^1$, và do đó cấu xạ này (không nhất thiết được xác định trên toàn $X$) không trội;
theo cách nói lạm dụng, khi đó ta nói ánh xạ hữu tỉ $u$ là «hằng, có giá trị $y$».
Còn phải thiết lập tương ứng song ánh giữa các điểm của $\bb{P}_k^1$ có giá trị trong $K$ mà không nằm tại $\infty$ \sref[I]{I.3.4.5} và các phần tử của $K$, rồi kiểm tra rằng một điểm như vậy nằm tại điểm tổng quát của $\bb{P}_k^1$ khi và chỉ khi nó tương ứng với một
\oldpage[II]{152}
phần tử siêu việt trên $k$.
Điều này suy ra ngay từ \sref{II.4.2.6}[ví dụ~1].
\end{proof}

\begin{corollary}[7.4.16]
\label{II.7.4.16}
Cho $X$ và $Y$ là hai đường cong đại số trên $k$, chuẩn tắc, đầy đủ và bất khả quy;
đặt $K=R(X)$ và $L=R(Y)$ là các trường hàm tương ứng.
Khi đó có một tương ứng song ánh chính tắc giữa tập các $k$-đẳng cấu $X\simto Y$ và tập các $k$-đẳng cấu $L\simto K$.
\end{corollary}

\begin{proof}
Đây là hệ quả hiển nhiên của \sref{II.7.4.13}.
\end{proof}

\begin{env}[7.4.17]
\label{II.7.4.17}
Hệ quả \sref{II.7.4.16} cho thấy rằng một đường cong đại số trên $k$, chuẩn tắc, đầy đủ và bất khả quy, \emph{được trường hàm $K$ của nó xác định duy nhất đến một đẳng cấu duy nhất};
theo định nghĩa, $K$ là một mở rộng hữu hạn sinh của $k$, có bậc siêu việt $1$ (theo cách gọi cổ điển, đó là một \emph{trường hàm đại số một biến}).
Hơn nữa:
\end{env}

\begin{proposition}[7.4.18]
\label{II.7.4.18}
Với mọi mở rộng $K$ hữu hạn sinh của $k$ có bậc siêu việt $1$, tồn tại một đường cong đại số $X$ (được xác định duy nhất đến một đẳng cấu duy nhất), chuẩn tắc, đầy đủ và bất khả quy, sao cho $R(X)=K$.
Tập các vành địa phương của $X$ được đồng nhất \sref[I]{I.8.2.1} với tập gồm trường $K$ và các vành định giá chứa $k$ và có $K$ làm trường phân thức.
\end{proposition}

\begin{proof}
Thật vậy, $K$ là một mở rộng hữu hạn của một mở rộng siêu việt thuần $k(T)$ của $k$, mà mở rộng sau, như đã thấy, được đồng nhất với trường các hàm hữu tỉ của đường thẳng xạ ảnh $Y=\bb{P}_k^1$.
Cho $X$ là \emph{bao đóng nguyên} của $Y$ đối với $K$ \sref{II.6.3.4};
khi đó $X$ là một đường cong đại số chuẩn tắc trên $k$, có trường hàm $K$ \sref{II.6.3.7}, và đầy đủ vì cấu xạ $X\to Y$ là hữu hạn \sref{II.7.4.6}.
Các vành địa phương $\sh{O}_x$ của $X$ hoặc là trường $K$, khi $x$ là điểm tổng quát, hoặc là các vành định giá rời rạc chứa $k$ và có $K$ làm trường phân thức, khi $x$ khác điểm tổng quát \sref{II.7.4.5}.
Ngược lại, cho $A$ là một vành như vậy;
vì cấu xạ $X\to\Spec(k)$ là riêng và $A$ chi phối $k$, nên $A$ cũng chi phối một vành địa phương $\sh{O}_x$ của $X$ \sref{II.7.3.10};
do vành sau là một vành định giá có $K$ làm trường phân thức, nó nhất thiết bằng $A$.
\end{proof}

\begin{remarks}[7.4.19]
\label{II.7.4.19}
Từ \sref{II.7.4.16} và \sref{II.7.4.18} suy ra rằng \emph{việc cho một đường cong đại số trên $k$, chuẩn tắc, đầy đủ và bất khả quy, về bản chất tương đương với việc cho một mở rộng trường hữu hạn sinh $K/k$ có bậc siêu việt $1$}.
Ta lưu ý rằng, nếu $k'$ là một mở rộng của trường cơ sở $k$, thì $X\otimes_k k'$ vẫn là một đường cong đại số đầy đủ trên $k'$ \sref{II.5.4.2}[(iii)], nhưng nói chung nó sẽ không rút gọn cũng không bất khả quy.
Tuy nhiên, nó sẽ vừa rút gọn vừa bất khả quy nếu $K/k$ là một mở rộng tách được và $k$ đóng đại số trong $K$ (theo thuật ngữ cổ điển mà ta sẽ không dùng, điều này được diễn đạt bằng cách nói rằng $K$ là một «mở rộng chính quy» của $k$).
Nhưng ngay cả trong trường hợp này, $X\otimes_k k'$ vẫn có thể không chuẩn tắc.
Bạn đọc sẽ tìm thấy chi tiết về những vấn đề này trong Chương~IV.
\end{remarks}

\section{Các lược đồ thu được bằng phép nổ; các nón chiếu; bao đóng xạ ảnh}
\label{section:II.8}

\subsection{Các tiền lược đồ thu được bằng phép nổ}
\label{subsection:II.8.1}
\label{II.8.1}

\begin{env}[8.1.1]
\label{II.8.1.1}
Cho $Y$ là một tiền lược đồ và, với mỗi số nguyên $n\geq0$, cho $\sh{I}_n$ là một bó iđêan quasi-coherent của $\sh{O}_Y$; giả sử các điều kiện sau được thỏa mãn:
\[
\label{II.8.1.1.1}
  \sh{I}_0=\sh{O}_Y,\qquad \sh{I}_n\subset\sh{I}_m\quad\text{khi }m\leq n,
\tag{8.1.1.1}
\]
\[
\label{II.8.1.1.2}
  \sh{I}_m\sh{I}_n\subset\sh{I}_{m+n}\quad\text{với mọi }m,n.
\tag{8.1.1.2}
\]

\oldpage[II]{153}
Ta lưu ý rằng các giả thiết này kéo theo
\[
\label{II.8.1.1.3}
  \sh{I}_1^n\subset\sh{I}_n.
\tag{8.1.1.3}
\]

Đặt
\[
\label{II.8.1.1.4}
  \sh{S}=\bigoplus_{n\geq0}\sh{I}_n.
\tag{8.1.1.4}
\]

Từ \sref{II.8.1.1.1} và \sref{II.8.1.1.2} suy ra rằng $\sh{S}$ là một đại số $\sh{O}_Y$ phân bậc quasi-coherent, và do đó xác định một $Y$-lược đồ $X=\Proj(\sh{S})$.
Nếu $\sh{J}$ là một bó iđêan khả nghịch của $\sh{O}_Y$, thì $\sh{I}_n\otimes_{\sh{O}_Y}\sh{J}^{\otimes n}$ được đồng nhất chính tắc với $\sh{I}_n\sh{J}^n$.
Nếu sau đó ta thay các $\sh{I}_n$ bằng các $\sh{I}_n\sh{J}^n$ và, theo đó, thay $\sh{S}$ bằng đại số $\sh{O}_Y$ quasi-coherent $\sh{S}_{(\sh{J})}$, thì $X_{(\sh{J})}=\Proj(\sh{S}_{(\sh{J})})$ đẳng cấu chính tắc với $X$ \sref{II.3.1.8}.
\end{env}

\begin{env}[8.1.2]
\label{II.8.1.2}
Giả sử $Y$ là \emph{nguyên địa phương}; khi đó bó các hàm hữu tỉ $\sh{R}(Y)$ là một đại số $\sh{O}_Y$ quasi-coherent \sref[I]{I.7.3.7}.
Ta nói một $\sh{O}_Y$-môđun con $\sh{I}$ của $\sh{R}(Y)$ là một \emph{iđêan phân thức} của $\sh{R}(Y)$ nếu nó \emph{hữu hạn sinh} \sref[0]{0.5.2.1}.
Giả sử đã cho, với mỗi $n\geq0$, một iđêan phân thức quasi-coherent $\sh{I}_n$ của $\sh{R}(Y)$ sao cho $\sh{I}_0=\sh{O}_Y$ và điều kiện \sref{II.8.1.1.2} được thỏa mãn (nhưng không nhất thiết điều kiện thứ hai trong \sref{II.8.1.1.1});
khi đó, bằng công thức \sref{II.8.1.1.4}, ta lại có thể xác định một đại số $\sh{O}_Y$ phân bậc quasi-coherent và $Y$-lược đồ tương ứng $X=\Proj(\sh{S})$;
ta lại có một đẳng cấu chính tắc $X\xrightarrow{\sim}X_{(\sh{J})}$ với mọi iđêan phân thức \emph{khả nghịch} $\sh{J}$ của $\sh{R}(Y)$.
\end{env}

\begin{definition}[8.1.3]
\label{II.8.1.3}
Cho $Y$ là một tiền lược đồ (tương ứng, một tiền lược đồ nguyên địa phương), và cho $\sh{I}$ là một bó iđêan quasi-coherent của $\sh{O}_Y$ (tương ứng, một iđêan phân thức quasi-coherent của $\sh{R}(Y)$).
Ta nói $Y$-lược đồ
\[
  X=\Proj\left(\bigoplus_{n\geq0}\sh{I}^n\right)
\]
thu được bằng phép nổ iđêan $\sh{I}$, hoặc là tiền lược đồ thu được bằng phép nổ của $Y$ tương đối với $\sh{I}$.
Khi $\sh{I}$ là một bó iđêan quasi-coherent của $\sh{O}_Y$ và $Y'$ là tiền lược đồ con đóng của $Y$ được xác định bởi $\sh{I}$, ta cũng nói rằng $X$ là $Y$-lược đồ thu được bằng phép nổ $Y'$.
\end{definition}

Theo định nghĩa, $\sh{S}=\bigoplus_{n\geq0}\sh{I}^n$; do đó $\sh{S}$ được sinh bởi $\sh{S}_1=\sh{I}$;
nếu $\sh{I}$ là một $\sh{O}_Y$-môđun hữu hạn sinh, thì $X$ là \emph{xạ ảnh} trên $Y$ \sref{II.5.5.2}.
Không cần giả thiết nào về $\sh{I}$, $\sh{O}_X$-môđun $\sh{O}_X(1)$ là \emph{khả nghịch} \sref{II.3.2.5} và \emph{rất ample}, theo \sref{II.4.4.3} áp dụng cho cấu xạ cấu trúc $X\to Y$.

Ta lưu ý rằng, nếu $f:X\to Y$ là cấu xạ cấu trúc, thì hạn chế của $f$ lên $f^{-1}(Y\setminus Y')$ là một \emph{đẳng cấu} lên $Y\setminus Y'$ khi $\sh{I}$ là một \emph{iđêan của $\sh{O}_Y$} và $Y'$ là tiền lược đồ con đóng do nó xác định:
thật vậy, vì mệnh đề này có tính địa phương trên $Y$, chỉ cần giả sử $\sh{I}=\sh{O}_Y$, và khẳng định của ta suy ra từ \sref{II.3.1.7}.

Khi thay $\sh{I}$ bằng $\sh{I}^d$ ($d>0$), $Y$-lược đồ $X$ thu được bằng phép nổ tương đối với iđêan $\sh{I}$ được thay bằng một $Y$-lược đồ $X'$ đẳng cấu chính tắc với nó \sref{II.8.1.1};
tương tự, với mọi iđêan (tương ứng, iđêan phân thức) \emph{khả nghịch} $\sh{J}$, tiền lược đồ $X_{(\sh{J})}$ thu được bằng phép nổ tương đối với iđêan $\sh{I}\sh{J}$ đẳng cấu chính tắc với $X$ \sref{II.8.1.1}.

Đặc biệt, khi $\sh{I}$ là một iđêan (tương ứng, một iđêan phân thức) \emph{khả nghịch}, $Y$-lược đồ thu được bằng phép nổ $\sh{I}$ là \emph{đẳng cấu với $Y$} \sref{II.3.1.7}.

\begin{proposition}[8.1.4]
\label{II.8.1.4}
Cho $Y$ là một tiền lược đồ nguyên.
\begin{enumerate}
    \item[\rm{(i)}] Với mọi dãy $(\sh{I}_n)$ các iđêan phân thức quasi-coherent của $\sh{R}(Y)$ thỏa mãn \sref{II.8.1.1.2}
\oldpage[II]{154}
        và sao cho $\sh{I}_0=\sh{O}_Y$, $Y$-lược đồ $X=\Proj(\bigoplus_{n\geq0}\sh{I}_n)$ là nguyên, và cấu xạ cấu trúc $f:X\to Y$ là trội.
    \item[\rm{(ii)}] Cho $\sh{I}$ là một iđêan phân thức quasi-coherent của $\sh{R}(Y)$, và cho $X$ là $Y$-lược đồ thu được bằng phép nổ của $Y$ tương đối với $\sh{I}$.
        Nếu $\sh{I}\neq0$, thì cấu xạ cấu trúc $f:X\to Y$ là song hữu tỉ và toàn ánh.
\end{enumerate}
\end{proposition}

\begin{proof}
\medskip\noindent
\begin{enumerate}
    \item[\rm{(i)}] Điều này suy ra từ việc $\sh{S}=\bigoplus_{n\geq0}\sh{I}_n$ là một đại số $\sh{O}_Y$ nguyên (\sref{II.3.1.12} và \sref{II.3.1.14}), vì với mọi $y\in Y$, $\sh{O}_y$ là một vành nguyên \sref[I]{I.5.1.4}.
    \item[\rm{(ii)}] Theo (i), $X$ là nguyên; hơn nữa, nếu $x$ và $y$ là các điểm tổng quát của $X$ và $Y$ (tương ứng), thì $f(x)=y$, và còn phải chứng minh rằng $\kres(x)$ là một không gian vectơ có chiều $1$ trên $\kres(y)$.
        Nhưng $x$ cũng là điểm tổng quát của thớ $f^{-1}(y)$; nếu $\psi$ là cấu xạ chính tắc $Z\to Y$, trong đó $Z=\Spec(\kres(y))$, thì tiền lược đồ $f^{-1}(y)$ có thể được đồng nhất với $\Proj(\sh{S}')$, trong đó $\sh{S}'=\psi^*(\sh{S})$ \sref{II.3.5.3}.
        Nhưng rõ ràng $\sh{S}'=\bigoplus_{n\geq0}(\sh{I}_y)^n$, và vì $\sh{I}$ là một iđêan phân thức quasi-coherent của $\sh{R}(Y)$ khác không, $\sh{I}_y\neq0$ \sref[I]{I.7.3.6}, suy ra $\sh{I}_y=\kres(y)$;
        khi đó $\Proj(\sh{S}')$ có thể được đồng nhất với $\Spec(\kres(y))$ \sref{II.3.1.7}, từ đó suy ra kết luận.
\end{enumerate}
\end{proof}

Ta sẽ chứng minh một mệnh đề đảo của \sref{II.8.1.4} trong \sref[III]{III.2.3.8}.

\begin{env}[8.1.5]
\label{II.8.1.5}
Ta trở lại tình huống và các ký hiệu của \sref{II.8.1.1}.
Theo định nghĩa, các đồng cấu bao hàm $\sh{I}_{n+1}\to\sh{I}_n$ \sref{II.8.1.1.1} xác định, với mỗi $k\in\bb{Z}$, một đơn cấu bậc $0$ giữa các $\sh{S}$-môđun phân bậc
\[
\label{II.8.1.5.1}
  u_k:\sh{S}_+(k+1)\to\sh{S}(k);
\tag{8.1.5.1}
\]
vì $\sh{S}_+(k+1)$ và $\sh{S}(k+1)$ là \textbf{(TN)}-đẳng cấu với nhau một cách chính tắc, $u_k$ tương ứng chính tắc với một đơn cấu giữa các $\sh{O}_X$-môđun \sref{II.3.4.2}:
\[
\label{II.8.1.5.2}
  \widetilde{u}_k:\sh{O}_X(k+1)\to\sh{O}_X(k).
\tag{8.1.5.2}
\]

Ta cũng nhắc lại \sref{II.3.2.6} rằng ta đã định nghĩa các đồng cấu chính tắc
\[
\label{II.8.1.5.3}
  \lambda:\sh{O}_X(h)\otimes_{\sh{O}_X}\sh{O}_X(k)\to\sh{O}_X(h+k)
\tag{8.1.5.3}
\]
và vì sơ đồ
\[
  \xymatrix{
    \sh{S}(h)\otimes_{\sh{S}}\sh{S}(k)\otimes_{\sh{S}}\sh{S}(l)
      \ar[r]
      \ar[d]
  & \sh{S}(h+k)\otimes_{\sh{S}}\sh{S}(l)
      \ar[d]
  \\\sh{S}(h)\otimes_{\sh{S}}\sh{S}(k+l)
      \ar[r]
  & \sh{S}(h+k+l)
  }
\]
giao hoán, từ tính hàm tử của các $\lambda$ \sref{II.3.2.6} suy ra rằng các đồng cấu \sref{II.8.1.5.3} xác định trên
\[
\label{II.8.1.5.4}
  \sh{S}_X=\bigoplus_{n\in\bb{Z}}\sh{O}_X(n)
\tag{8.1.5.4}
\]
\oldpage[II]{155}
một cấu trúc đại số $\sh{O}_X$ phân bậc quasi-coherent.

Hơn nữa, sơ đồ
\[
  \xymatrix{
    \sh{S}(h)\otimes_{\sh{S}}\sh{S}_+(k+1)
      \ar[r]
      \ar[d]_{1\otimes u_k}
  & \sh{S}_+(h+k+1)
      \ar[d]^{u_{k+h}}
  \\\sh{S}(h)\otimes_{\sh{S}}\sh{S}(k)
      \ar[r]
  & \sh{S}(h+k)
  }
\]
giao hoán; tính hàm tử của các $\lambda$ khi đó cho thấy ta có một sơ đồ giao hoán
\[
\label{II.8.1.5.5}
  \xymatrix{
    \sh{O}_X(h)\otimes_{\sh{O}_X}\sh{O}_X(k+1)
      \ar[r]^{\lambda}
      \ar[d]_{1\otimes\widetilde{u}_k}
  & \sh{O}_X(h+k+1)
      \ar[d]^{\widetilde{u}_{k+h}}
  \\\sh{O}_X(h)\otimes_{\sh{O}_X}\sh{O}_X(k)
      \ar[r]^{\lambda}
  & \sh{O}_X(h+k)
  }
\tag{8.1.5.5}
\]
trong đó các mũi tên ngang là các đồng cấu chính tắc.
Do đó có thể nói rằng các $\widetilde{u}_k$ xác định một \emph{đơn cấu} (bậc $0$) \emph{giữa các $\sh{S}_X$-môđun phân bậc}
\[
\label{II.8.1.5.6}
  \widetilde{u}:\sh{S}_X(1)\to\sh{S}_X.
\tag{8.1.5.6}
\]
\end{env}

\begin{env}[8.1.6]
\label{II.8.1.6}
Vẫn giữ các ký hiệu của \sref{II.8.1.5}, ta nhận thấy rằng, với $n\geq0$, hợp thành $\widetilde{v}_n=\widetilde{u}_{n-1}\circ\widetilde{u}_{n-2}\circ\ldots\circ\widetilde{u}_0$ là một \emph{đơn cấu} $\sh{O}_X(n)\to\sh{O}_X$;
ta ký hiệu bởi $\sh{I}_{n,X}$ ảnh của nó, do đó đây là một iđêan quasi-coherent của $\sh{O}_X$, \emph{đẳng cấu} với $\sh{O}_X(n)$.
Hơn nữa, sơ đồ
\[
  \xymatrix{
    \sh{O}_X(m)\otimes_{\sh{O}_X}\sh{O}_X(n)
      \ar[r]^{\lambda}
      \ar[d]_{\widetilde{v}_m\otimes\widetilde{v}_n}
  & \sh{O}_X(m+n)
      \ar[d]^{\widetilde{v}_{m+n}}
  \\\sh{O}_X
      \ar[r]_{\id}
  & \sh{O}_X
  }
\]
giao hoán với $m\geq0$, $n\geq0$.
Từ đó suy ra các quan hệ bao hàm sau:
\[
\label{II.8.1.6.1}
  \sh{I}_{0,X}=\sh{O}_X,\quad \sh{I}_{n,X}\subset\sh{I}_{m,X}
  \qquad\mbox{với $0\leq m\leq n$;}
\tag{8.1.6.1}
\]
\[
\label{II.8.1.6.2}
  \sh{I}_{m,X}\sh{I}_{n,X}\subset\sh{I}_{m+n,X}
  \qquad\qquad\mbox{với $m\geq0$, $n\geq0$.}
\tag{8.1.6.2}
\]
\end{env}

\oldpage[II]{156}

\begin{proposition}[8.1.7]
\label{II.8.1.7}
Cho $Y$ là một tiền lược đồ, $\sh{I}$ là một bó iđêan quasi-coherent của $\sh{O}_Y$, và $X=\Proj\left(\bigoplus_{n\geq0}\sh{I}^n\right)$ là $Y$-lược đồ thu được bằng phép nổ $\sh{I}$.
Khi đó, với mọi $n>0$, ta có một đẳng cấu chính tắc
\[
\label{II.8.1.7.1}
  \sh{O}_X(n)\simto\sh{I}^n\sh{O}_X=\sh{I}_{n,X}
\tag{8.1.7.1}
\]
(xem \sref[0]{0.4.3.5}), và do đó $\sh{I}^n\sh{O}_X$ là một $\sh{O}_X$-môđun khả nghịch rất ample nếu $n>0$.
\end{proposition}

\begin{proof}
Khẳng định cuối cùng là hiển nhiên, vì $\sh{O}_X(1)$ là khả nghịch \sref{II.3.2.5} và rất ample đối với $Y$ theo định nghĩa (\sref{II.4.4.3} và \sref{II.4.4.9}).
Mặt khác, theo định nghĩa, ảnh của $v_n$ chính là $\sh{I}^n\sh{S}$, và \sref{II.8.1.7.1} khi đó suy ra từ tính khớp của hàm tử $\sh{M}\mapsto\widetilde{\sh{M}}$ \sref{II.3.2.4} và từ công thức~\sref{II.3.2.4.1}.
\end{proof}

\begin{corollary}[8.1.8]
\label{II.8.1.8}
Dưới các giả thiết của \sref{II.8.1.7}, nếu $f:X\to Y$ là cấu xạ cấu trúc và $Y'$ là tiền lược đồ con đóng của $Y$ được xác định bởi $\sh{I}$, thì tiền lược đồ con đóng $X'=f^{-1}(Y')$ của $X$ được xác định bởi $\sh{I}\sh{O}_X$ ($\sh{I}\sh{O}_X$ đẳng cấu chính tắc với $\sh{O}_X(1)$); từ đó ta thu được một dãy khớp ngắn chính tắc
\[
\label{II.8.1.8.1}
  0\to\sh{O}_X(1)\to\sh{O}_X\to\sh{O}_{X'}\to0.
\tag{8.1.8.1}
\]
\end{corollary}

\begin{proof}
Điều này suy ra từ \sref{II.8.1.7.1} và từ \sref[I]{I.4.4.5}.
\end{proof}

\begin{env}[8.1.9]
\label{II.8.1.9}
Dưới các giả thiết của \sref{II.8.1.7}, ta có thể mô tả chính xác hơn cấu trúc của các $\sh{I}_{n,X}$.
Thật vậy, xét đồng cấu
\[
  \widetilde{u}_{-1}:\sh{O}_X\to\sh{O}_X(-1),
\]
đồng cấu này tương ứng chính tắc với một tiết diện $s$ của $\sh{O}_X(-1)$ trên $X$, mà ta gọi là \emph{tiết diện chính tắc} (đối với $\sh{I}$) \sref[0]{0.5.1.1}.
Mặt khác, trong sơ đồ~\sref{II.8.1.5.5}, các mũi tên ngang là các đẳng cấu \sref{II.3.2.7}; thay $h$ bởi $k$ và $k$ bởi $-1$ trong sơ đồ này, ta được $\widetilde{u}_{k-1}=1_k\otimes\widetilde{u}_{-1}$ (trong đó $1_k$ ký hiệu ánh xạ đồng nhất trên $\sh{O}_X(k)$), nói cách khác, đồng cấu $\widetilde{u}_k$ chính là \emph{phép nhân tenxơ với tiết diện chính tắc $s$} (với mọi $k\in\bb{Z}$).
Khi đó, đồng cấu $\widetilde{u}$ \sref{II.8.1.5.6} cũng được hiểu theo cùng một cách.

Do đó, với mọi $n\geq0$, đồng cấu $\widetilde{v}_n:\sh{O}_X(n)\to\sh{O}_X$ chính là phép nhân tenxơ với $s^{\otimes n}$; ta suy ra:
\end{env}

\begin{corollary}[8.1.10]
\label{II.8.1.10}
Với các ký hiệu của \sref{II.8.1.8}, không gian nền của $X'$ là tập hợp các điểm $x\in X$ sao cho $s(x)=0$, trong đó $s$ ký hiệu tiết diện chính tắc của $\sh{O}_X(-1)$.
\end{corollary}

\begin{proof}
Thật vậy, nếu $c_x$ là một phần tử sinh của thớ $(\sh{O}_X(1))_x$ tại một điểm $x$, thì $s_x\otimes c_x$ được đồng nhất chính tắc với một phần tử sinh của thớ của $\sh{I}_{1,X}$ tại điểm $x$, và do đó là phần tử khả nghịch khi và chỉ khi $s_x\notin\mathfrak{m}_x(\sh{O}_X(-1))_x$, hay tương đương, khi và chỉ khi $s(x)\neq0$.
\end{proof}

\begin{proposition}[8.1.11]
\label{II.8.1.11}
Cho $Y$ là một tiền lược đồ nguyên, $\sh{I}$ là một iđêan phân thức quasi-coherent của $\sh{R}(Y)$, và $X$ là $Y$-lược đồ thu được bằng phép nổ $\sh{I}$.
Khi đó $\sh{I}\sh{O}_X$ là một $\sh{O}_X$-môđun khả nghịch rất ample đối với $Y$.
\end{proposition}

\begin{proof}
Vì vấn đề có tính địa phương trên $Y$ \sref{II.4.4.5}, ta có thể quy về trường hợp $Y=\Spec(A)$, trong đó $A$ là một miền nguyên có trường phân thức $K$, và $\sh{I}=\widetilde{\mathfrak{I}}$, với $\mathfrak{I}$ là một iđêan phân thức của $K$;
khi đó tồn tại một phần tử $a\neq0$ của $A$ sao cho $a\mathfrak{I}\subset A$.
Đặt $S=\bigoplus_{n\geq0}\mathfrak{I}^n$;
ánh xạ $x\mapsto ax$ là một $A$-đẳng cấu từ $\mathfrak{I}^{n+1}=(S(1))_n$ lên $a\mathfrak{I}^{n+1}=a\mathfrak{I}S_n\subset\mathfrak{I}^n=S_n$,
\oldpage[II]{157}
và do đó xác định một \textbf{(TN)}-đẳng cấu bậc $0$ giữa các $S$-môđun phân bậc $S_+(1)\to a\mathfrak{I}S$.
Mặt khác, $x\mapsto a^{-1}x$ là một đẳng cấu bậc $0$ giữa các $S$-môđun phân bậc $a\mathfrak{I}S\simto\mathfrak{I}S$.
Bằng cách hợp thành \sref{II.3.2.4}, ta thu được một đẳng cấu các $\sh{O}_X$-môđun $\sh{O}_X(1)\simto\sh{I}\sh{O}_X$ và, vì $S$ được sinh bởi $S_1=\mathfrak{I}$, $\sh{O}_X(1)$ là khả nghịch \sref{II.3.2.5} và rất ample (\sref{II.4.4.3} và \sref{II.4.4.9}), điều phải chứng minh.
\end{proof}

\subsection{Các kết quả sơ bộ về phép địa phương hóa trong các vành phân bậc}
\label{subsection:II.8.2}

\begin{env}[8.2.1]
\label{II.8.2.1}
Cho $S$ là một vành phân bậc mà hiện tại ta không giả sử tập trung ở các bậc không âm.
Đặt
\[
\label{II.8.2.1.1}
  S^\geq=\bigoplus_{n\geq0}S_n,
  \qquad
  S^\leq=\bigoplus_{n\leq0}S_n
\tag{8.2.1.1}
\]
là hai vành con phân bậc của $S$, lần lượt tập trung ở các bậc không âm và không dương.
Nếu $f$ là một phần tử thuần nhất bậc $d$ (dương hoặc âm) của $S$, thì vành các phân thức $S_f=S'$ lại được trang bị một cấu trúc vành phân bậc, trong đó $S'_n$ ($n\in\bb{Z}$) được lấy là tập hợp các $x/f^k$ với $x\in S_{n+kd}$ ($k\geq0$);
ngoài ra, ta đặt $S_{(f)}=S'_0$, và viết $S_f^\geq$ và $S_f^\leq$ thay cho $S^{\prime\geq}$ và $S^{\prime\leq}$, theo thứ tự đó.
Nếu $d>0$, ta có
\[
\label{II.8.2.1.2}
  (S^\geq)_f=S_f
\tag{8.2.1.2}
\]
vì, nếu $x\in S_{n+kd}$ và $n+kd<0$, thì có thể viết $x/f^k=xf^h/f^{h+k}$, và $n+(h+k)d>0$ khi $h$ dương và đủ lớn.
Do định nghĩa, suy ra rằng
\[
\label{II.8.2.1.3}
  (S^\geq)_{(f)}=(S_f^\geq)_0=S_{(f)}.
\tag{8.2.1.3}
\]

Đối với một $S$-môđun phân bậc $M$, ta đặt tương tự
\[
\label{II.8.2.1.4}
  M^\geq=\bigoplus_{n\geq0}M_n,
  \qquad
  M^\leq=\bigoplus_{n\leq0}M_n
\tag{8.2.1.4}
\]
lần lượt là một $S^\geq$-môđun phân bậc và một $S^\leq$-môđun phân bậc, và giao của chúng là $S_0$-môđun $M_0$.
Nếu $f\in S_d$, ta cũng định nghĩa $M_f$ là $S_f$-môđun phân bậc mà các phần tử bậc $n$ là các $z/f^k$ với $z\in M_{n+kd}$ ($k\geq0$);
ta ký hiệu bởi $M_{(f)}$ tập hợp các phần tử bậc $0$ của $M_f$, đó là một $S_{(f)}$-môđun, và viết $M_f^\geq$ và $M_f^\leq$ thay cho $(M_f)^\geq$ và $(M_f)^\leq$, theo thứ tự đó.
Nếu $d>0$, như trên ta thấy rằng
\[
\label{II.8.2.1.5}
  (M^\geq)_f=M_f
\tag{8.2.1.5}
\]
và
\[
\label{II.8.2.1.6}
  (M^\geq)_{(f)}=(M_f^\geq)_0=M_{(f)}.
\tag{8.2.1.6}
\]
\end{env}

\begin{env}[8.2.2]
\label{II.8.2.2}
Cho $\bb{z}$ là một biến bất định, mà ta sẽ gọi là \emph{biến thuần nhất hóa}.
Nếu $S$ là một vành phân bậc (với các bậc nguyên tùy ý), thì đại số đa thức\footnote{Không nên nhầm ở đây với việc dùng ký hiệu $\widehat{S}$ để chỉ vành hoàn thành tách của một vành.}
\[
\label{II.8.2.2.1}
  \widehat{S}=S[\bb{z}]
\tag{8.2.2.1}
\]
\oldpage[II]{158}
là một $S$-đại số phân bậc, trong đó ta định nghĩa bậc của $f\bb{z}^n$ ($n\geq0$), với $f$ thuần nhất, bởi
\[
\label{II.8.2.2.2}
  \deg(f\bb{z}^n)=n+\deg f.
\tag{8.2.2.2}
\]
\end{env}

\begin{lemma}[8.2.3]
\label{II.8.2.3}
\begin{enumerate}
  \item[\rm{(i)}] Có các đẳng cấu chính tắc giữa các vành (không phân bậc)
    \[
    \label{II.8.2.3.1}
      \widehat{S}_{(\bb{z})}
      \simto
      \widehat{S}/(\bb{z}-1)\widehat{S}
      \simto
      S.
    \tag{8.2.3.1}
    \]
  \item[\rm{(ii)}] Có một đẳng cấu chính tắc giữa các vành (không phân bậc)
    \[
    \label{II.8.2.3.2}
      \widehat{S}_{(f)} \simto S_f^\leq
    \tag{8.2.3.2}
    \]
    với mọi $f\in S_d$ sao cho $d>0$.
\end{enumerate}
\end{lemma}

\begin{proof}
Đẳng cấu thứ nhất trong \sref{II.8.2.3.1} đã được xác định trong \sref{II.2.2.5}, còn đẳng cấu thứ hai là hiển nhiên;
do đó, đẳng cấu $\widehat{S}_{(\bb{z})}\simto S$ được xác định như vậy cho $x\bb{z}^n/\bb{z}^{n+k}$, trong đó $\deg(x)=k$ với $k\geq-n$, tương ứng với phần tử $x$.
Đồng cấu \sref{II.8.2.3.2} cho $x\bb{z}^n/f^k$, trong đó $\deg(x)=kd-n$, tương ứng với phần tử $x/f^k$, có bậc $-n$ trong $S_f^\leq$, và một lần nữa hiển nhiên rằng đây quả là một đẳng cấu.
\end{proof}

\begin{env}[8.2.4]
\label{II.8.2.4}
Cho $M$ là một $S$-môđun phân bậc.
Hiển nhiên $S$-môđun
\[
\label{II.8.2.4.1}
  \widehat{M}=M\otimes_S\widehat{S}=M\otimes_SS[\bb{z}]
\tag{8.2.4.1}
\]
là tổng trực tiếp của các $S$-môđun $M\otimes_SS\bb{z}^n$, và do đó cũng là tổng trực tiếp của các nhóm Abel $M_k\otimes S\bb{z}^n$ ($k\in\bb{Z}$, $n\geq0$);
ta định nghĩa trên $\widehat{M}$ một cấu trúc $\widehat{S}$-môđun phân bậc bằng cách đặt
\[
\label{II.8.2.4.2}
  \deg(x\otimes\bb{z}^n)=n+\deg x
\tag{8.2.4.2}
\]
với mọi phần tử thuần nhất $x$ của $M$.
Ta dành cho độc giả việc chứng minh kết quả tương tự với \sref{II.8.2.3}:
\end{env}

\begin{lemma}[8.2.5]
\label{II.8.2.5}
\begin{enumerate}
  \item[\rm{(i)}] Có một di-đẳng cấu chính tắc giữa các môđun (không phân bậc)
    \[
    \label{II.8.2.5.1}
      \widehat{M}_{(\bb{z})} \simto M.
    \tag{8.2.5.1}
    \]
  \item[\rm{(ii)}] Với mọi $f\in S_d$ ($d>0$), có một di-đẳng cấu giữa các môđun (không phân bậc)
    \[
    \label{II.8.2.5.2}
      \widehat{M}_{(f)} \simto M_f^\leq.
    \tag{8.2.5.2}
    \]
\end{enumerate}
\end{lemma}

\begin{env}[8.2.6]
\label{II.8.2.6}
Cho $S$ là một vành phân bậc theo các bậc \emph{dương}, và xét trong $S$ dãy giảm dần các iđêan phân bậc
\[
\label{II.8.2.6.1}
  S_{[n]} = \bigoplus_{m\geq n} S_m
  \qquad\mbox{($n\geq0$)}
\tag{8.2.6.1}
\]
(đặc biệt, ta có $S_{[0]}=S$ và $S_{[1]}=S_+$).
Vì hiển nhiên $S_{[m]}S_{[n]}\subset S_{[m+n]}$, ta có thể định nghĩa một \emph{vành phân bậc} $S^\natural$ bằng cách đặt
\[
\label{II.8.2.6.2}
  S^\natural = \bigoplus_{n\geq0} S_n^\natural
  \quad
  \text{với}
  \quad
  S_n^\natural = S_{[n]}.
\tag{8.2.6.2}
\]
Khi đó, $S_0^\natural$ là vành $S$ được xét như một \emph{vành không phân bậc}, và do đó $S^\natural$ là một $S_0^\natural$-đại số.
Với mọi phần tử thuần nhất $f\in S_d$ ($d>0$), ta ký hiệu bởi $f^\natural$ phần tử $f$ được xét như thuộc $S_{[d]}=S_d^\natural$.
Với các ký hiệu này:
\end{env}

\oldpage[II]{159}
\begin{lemma}[8.2.7]
\label{II.8.2.7}
Cho $S$ là một vành phân bậc theo các bậc dương, và cho $f$ là một phần tử thuần nhất của $S_d$ ($d>0$).
Có các đẳng cấu chính tắc giữa các vành
\[
\label{II.8.2.7.1}
  S_f \simto \bigoplus_{n\in\bb{Z}} S(n)_{(f)}
\tag{8.2.7.1}
\]
\[
\label{II.8.2.7.2}
  (S_f^\geq)_{f/1} \simto S_f
\tag{8.2.7.2}
\]
\[
\label{II.8.2.7.3}
  S_{(f^\natural)}^\natural \simto S_f^\geq
\tag{8.2.7.3}
\]
trong đó hai đẳng cấu đầu là các đẳng cấu giữa các vành phân bậc.
\end{lemma}

\begin{proof}
Theo định nghĩa, hiển nhiên ta có $(S_f)_n=(S(n)_f)_0$, từ đó suy ra đẳng cấu \sref{II.8.2.7.1}, chính là ánh xạ đồng nhất.
Mặt khác, vì $f/1$ khả nghịch trong $S_f$, có một đẳng cấu chính tắc $S_f\simto(S_f^\geq)_{f/1}=(S_f)_{f/1}$, theo \sref{II.8.2.1.2} áp dụng cho $S_f$;
đẳng cấu ngược, theo định nghĩa, là đẳng cấu \sref{II.8.2.7.2}.
Cuối cùng, nếu $x=\sum_{m\geq n}y_m$ là một phần tử của $S_{[n]}$, với $n=kd$, thì phần tử $x/(f^\natural)^k$ tương ứng với phần tử $\sum_m y_m/f^k$ của $S_f^\geq$, và ta kiểm tra ngay rằng điều này xác định một đẳng cấu \sref{II.8.2.7.3}.
\end{proof}

\begin{env}[8.2.8]
\label{II.8.2.8}
Đối với một $S$-môđun phân bậc $M$, ta đặt tương tự, với mọi $n\in\bb{Z}$,
\[
\label{II.8.2.8.1}
  M_{[n]} = \bigoplus_{m\geq n} M_m
\tag{8.2.8.1}
\]
và vì $S_{[m]}M_{[n]} \subset M_{[m+n]}$ ($m\geq0$), ta có thể định nghĩa một $S^\natural$-môđun phân bậc $M^\natural$ bằng cách đặt
\[
\label{II.8.2.8.2}
  M^\natural = \bigoplus_{n\in\bb{Z}} M_n^\natural
  \quad
  \text{với}
  \quad
  M_n^\natural = M_{[n]}.
\tag{8.2.8.2}
\]
\end{env}

Ta lại dành cho độc giả việc chứng minh bổ đề sau:
\begin{lemma}[8.2.9]
\label{II.8.2.9}
Với các ký hiệu của \sref{II.8.2.7} và \sref{II.8.2.8}, có các di-đẳng cấu chính tắc giữa các môđun
\[
\label{II.8.2.9.1}
  M_f \simto \bigoplus_{n\in\bb{Z}} M(n)_{(f)}
\tag{8.2.9.1}
\]
\[
\label{II.8.2.9.2}
  (M_f^\geq)_{f/1} \simto M_f
\tag{8.2.9.2}
\]
\[
\label{II.8.2.9.3}
  M_{(f^\natural)}^\natural \simto M_f^\geq
\tag{8.2.9.3}
\]
trong đó hai di-đẳng cấu đầu là các di-đẳng cấu giữa các môđun phân bậc.
\end{lemma}

\begin{lemma}[8.2.10]
\label{II.8.2.10}
Cho $S$ là một vành phân bậc theo các bậc dương.
\begin{enumerate}
  \item[\rm{(i)}] Điều kiện cần và đủ để $S^\natural$ là một $S_0^\natural$-đại số hữu hạn sinh (tương ứng, một vành Noether) là $S$ phải là một $S_0$-đại số hữu hạn sinh (tương ứng, một vành Noether).
  \item[\rm{(ii)}] Điều kiện cần và đủ để $S_{n+1}^\natural = S_1^\natural S_n^\natural$ ($n\geq n_0$) là $S_{n+1}=S_1S_n$ ($n\geq n_0$).
  \item[\rm{(iii)}] Điều kiện cần và đủ để $S_n^\natural = (S_1^\natural)^n$ ($n\geq n_0$) là $S_n=S_1^n$ ($n\geq n_0$).
  \item[\rm{(iv)}] Nếu $(f_\alpha)$ là một tập hợp các phần tử thuần nhất của $S_+$ sao cho $S_+$ được chứa trong căn trong $S$ của iđêan của $S$ được sinh bởi các $f_\alpha$, thì $S_+^\natural$ được chứa trong căn trong $S^\natural$ của iđêan của $S^\natural$ được sinh bởi các $f_\alpha^\natural$.
\end{enumerate}
\end{lemma}

\begin{proof}
\medskip\noindent
\begin{enumerate}
  \item[\rm{(i)}] Nếu $S^\natural$ là một $S_0^\natural$-đại số hữu hạn sinh, thì $S_+=S_1^\natural$ là một môđun hữu hạn sinh trên $S=S_0^\natural$, theo \sref{II.2.1.6}[i], và do đó $S$ là một $S_0$-đại số hữu hạn sinh \sref{II.2.1.4};
    nếu $S^\natural$ là một vành Noether, thì $S_0^\natural=S$ cũng là một vành Noether \sref{II.2.1.5}.
    Ngược lại, nếu $S$ là một $S_0$-đại số
\oldpage[II]{160}
    hữu hạn sinh, thì theo \sref{II.2.1.6}[ii], ta biết rằng tồn tại $h>0$ và $m_0>0$ sao cho $S_{n+h}=S_hS_n$ với $n\geq m_0$;
    hiển nhiên ta có thể giả sử rằng $m_0\geq h$.
    Hơn nữa, các $S_m$ là các $S_0$-môđun hữu hạn sinh \sref{II.2.1.6}[i].
    Do đó, nếu $n\geq m_0+h$, thì $S_n^\natural = S_hS_{n-h}^\natural = S_h^\natural S_{n-h}^\natural$;
    còn nếu $m<m_0+h$, thì khi đặt $E = S_{m_0}+\ldots+S_{m_0+h-1}$, ta có
    \[
      S_m^\natural = S_m + \ldots + S_{m_0+h-1} + S_hE + S_h^2E + \ldots.
    \]
    Với $1\leq m\leq m_0$, đặt $G_m$ là hợp của các hệ sinh hữu hạn của các $S_0$-môđun $S_i$ với $m\leq i\leq m_0+h-1$, được xem như một tập con của $S_{[m]}$.
    Với $m_0+1\leq m\leq m_0+h-1$, đặt $G_m$ là hợp của các hệ sinh hữu hạn của các $S_0$-môđun $S_i$ với $m\leq i\leq m_0+h-1$ và của $S_hE$, được xem như một tập con của $S_{[m]}$.
    Rõ ràng $S_m^\natural=S_0^\natural G_m$ với $1\leq m\leq m_0+h-1$, và do đó hợp $G$ của các $G_m$ với $1\leq m\leq m_0+h-1$ là một hệ sinh của $S_0^\natural$-đại số $S^\natural$.
    Vì vậy, ta kết luận rằng nếu $S=S_0^\natural$ là một vành Noether thì $S^\natural$ cũng là một vành Noether.
  \item[\rm{(ii)}] Rõ ràng, nếu $S_{n+1}=S_1S_n$ với $n\geq n_0$, thì $S_{n+1}^\natural=S_1S_n^\natural$, và \emph{a fortiori} $S_{n+1}^\natural=S_1^\natural S_n^\natural$ với $n\geq n_0$.
    Ngược lại, đẳng thức cuối này có thể viết thành
    \[
      S_{n+1} + S_{n+2} + \ldots
      =
      (S_1 + S_2 + \ldots)(S_n + S_{n+1} + \ldots)
    \]
    và việc so sánh các hạng tử bậc $n+1$ (trong $S$) ở hai vế cho ta $S_{n+1}=S_1S_n$.
  \item[\rm{(iii)}] Nếu $S_n=S_1^n$ với $n\geq n_0$, thì $S_n^\natural=S_1^n+S_1^{n+1}+\ldots$;
    vì $S_1^\natural$ chứa $S_1+S_1^2+\ldots$, ta có $S_n^\natural\subset S_1^{\natural n}$, và do đó $S_n^\natural=S_1^{\natural n}$ với $n\geq n_0$.
    Ngược lại, các hạng tử duy nhất của $S_1^{\natural n}=(S_1+S_2+\ldots)^n$ có bậc $n$ trong $S$ là các hạng tử của $S_1^n$;
    do đó, đẳng thức $S_n^\natural=S_1^{\natural n}$ kéo theo $S_n=S_1^n$.
  \item[\rm{(iv)}] Chỉ cần chứng minh rằng nếu một phần tử $g\in S_{k+h}$ được xem như một phần tử của $S_k^\natural$ ($k>0$, $h\geq0$), thì tồn tại một số nguyên $n>0$ sao cho $g^n$ là một tổ hợp tuyến tính (trong $S_{kn}^\natural$) của các $f_\alpha^\natural$ với các hệ số trong $S^\natural$.
    Theo giả thiết, tồn tại một số nguyên $m_0$ sao cho, với $m\geq m_0$, ta có, \emph{trong $S$}, $g^m = \sum_\alpha c_{\alpha m}f_\alpha$, trong đó các chỉ số $\alpha$ xuất hiện ở đây \emph{không phụ thuộc vào $m$};
    hơn nữa, hiển nhiên ta có thể giả sử rằng các $c_{\alpha m}$ là thuần nhất, với
    \[
      \deg(c_{\alpha m}) = m(k+h)-\deg f_\alpha
    \]
    trong $S$.
    Do đó, lấy $m_0$ đủ lớn để bảo đảm rằng $km_0>\deg f_\alpha$ đối với mọi $f_\alpha$ xuất hiện trong $g^{m_0}$;
    với mọi $\alpha$, đặt $c'_{\alpha m}$ là phần tử $c_{\alpha m}$ được xem như có bậc $km-\deg f_\alpha$ \emph{trong $S^\natural$};
    khi đó, trong $S^\natural$, ta có $g^m = \sum_\alpha c'_{\alpha m}f_\alpha^\natural$, điều này hoàn tất chứng minh.
\end{enumerate}
\end{proof}

\begin{env}[8.2.11]
\label{II.8.2.11}
Xét $S_0$-đại số phân bậc
\[
\label{II.8.2.11.1}
  S^\natural \otimes_S S_0
  =
  S^\natural/S_+S^\natural
  =
  \bigoplus_{n\geq0} S_{[n]}/S_+S_{[n]}.
\tag{8.2.11.1}
\]

Vì $S_n$ là một $S_0$-môđun thương của $S_{[n]}/S_+S_{[n]}$, nên tồn tại một đồng cấu chính tắc giữa các $S_0$-đại số phân bậc
\[
\label{II.8.2.11.2}
  S^\natural \otimes_S S_0 \to S
\tag{8.2.11.2}
\]
đồng cấu này hiển nhiên \emph{toàn ánh}, và do đó, theo \sref{II.2.9.2}, tương ứng với một \emph{phép nhúng đóng} chính tắc
\[
\label{II.8.2.11.3}
  \Proj(S) \to \Proj(S^\natural \otimes_S S_0).
\tag{8.2.11.3}
\]
\end{env}

\oldpage[II]{161}
\begin{proposition}[8.2.12]
\label{II.8.2.12}
Cấu xạ chính tắc \sref{II.8.2.11.3} là song ánh.
Để đồng cấu \sref{II.8.2.11.2} là (TN)-song ánh, điều kiện cần và đủ là tồn tại một $n_0$ sao cho $S_{n+1}=S_1S_n$ với $n\geq n_0$.
Nếu điều kiện sau cùng này được thỏa mãn, thì \sref{II.8.2.11.3} là một đẳng cấu;
mệnh đề đảo đúng khi $S$ là một vành Noether.
\end{proposition}

\begin{proof}
Để chứng minh khẳng định thứ nhất, theo \sref{II.2.8.3}, chỉ cần chứng minh rằng hạt nhân $\mathfrak{I}$ của đồng cấu \sref{II.8.2.11.2} gồm các phần tử \emph{lũy linh}.
Nhưng nếu $f\in S_{[n]}$ là một phần tử có lớp modulo $S_+S_{[n]}$ thuộc hạt nhân này, thì điều đó kéo theo $f\in S_{[n+1]}$;
khi đó $f^{n+1}$, được xét như một phần tử của $S_{[n(n+1)]}$, cũng thuộc $S_+S_{[n(n+1)]}$, vì nó có thể được viết dưới dạng $f\cdot f^n$;
do đó, lớp của $f^{n+1}$ modulo $S_+S_{[n(n+1)]}$ bằng không, qua đó chứng minh khẳng định của ta.
Vì giả thiết $S_{n+1}=S_1S_n$ với $n\geq n_0$ tương đương với $S_{n+1}^\natural=S_1^\natural S_n^\natural$ với $n\geq n_0$ \sref{II.8.2.10}[ii], nên theo định nghĩa, giả thiết này tương đương với việc \sref{II.8.2.11.2} là (TN)-đơn ánh, và do đó là (TN)-song ánh, suy ra \sref{II.8.2.11.3} là một đẳng cấu theo \sref{II.2.9.1}.
Ngược lại, nếu \sref{II.8.2.11.3} là một đẳng cấu, thì bó $\widetilde{\mathfrak{I}}$ trên $\Proj(S^\natural\otimes_S S_0)$ bằng không \sref{II.2.9.2}[i];
vì $S^\natural\otimes_S S_0$ là một vành Noether, do là một vành thương của $S^\natural$ \sref{II.8.2.10}[i], từ \sref{II.2.7.3} suy ra rằng $\mathfrak{I}$ thỏa mãn điều kiện (TN), và do đó $S_{n+1}^\natural=S_1^\natural S_n^\natural$ với $n\geq n_0$, và theo \sref{II.8.2.10}[ii], điều này hoàn tất chứng minh.
\end{proof}

\begin{env}[8.2.13]
\label{II.8.2.13}
Bây giờ xét các đơn ánh chính tắc $(S_+)^n\to S_{[n]}$, chúng xác định một đồng cấu đơn ánh bậc không giữa các vành phân bậc
\[
\label{II.8.2.13.1}
  \bigoplus_{n\geq0} (S_+)^n \to S^\natural.
\tag{8.2.13.1}
\]
\end{env}

\begin{proposition}[8.2.14]
\label{II.8.2.14}
Để đồng cấu \sref{II.8.2.13.1} là một (TN)-đẳng cấu, điều kiện cần và đủ là tồn tại một $n_0$ sao cho $S_n=S_1^n$ với mọi $n\geq n_0$.
Khi điều này xảy ra, cấu xạ tương ứng với \sref{II.8.2.13.1} được xác định khắp nơi và là một đẳng cấu
\[
  \Proj(S^\natural) \simto \Proj(\bigoplus_{n\geq0}(S_+)^n);
\]
mệnh đề đảo đúng khi $S$ là một vành Noether.
\end{proposition}

\begin{proof}
Hai khẳng định đầu là hiển nhiên, theo \sref{II.8.2.10}[iii] và \sref{II.2.9.1}.
Khẳng định thứ ba sẽ là hệ quả của \sref{II.8.2.10}[i và iii] và bổ đề sau:
\begin{lemma}[8.2.14.1]
\label{II.8.2.14.1}
Cho $T$ là một vành phân bậc theo các bậc dương, đồng thời là một $T_0$-đại số hữu hạn sinh.
Nếu cấu xạ tương ứng với đồng cấu đơn ánh $\bigoplus_{n\geq0}T_1^n\to T$ được xác định khắp nơi và cũng là một đẳng cấu $\Proj(T)\to\Proj(\bigoplus_{n\geq0}T_1^n)$, thì tồn tại một $n_0$ sao cho $T_n=T_1^n$ với $n\geq n_0$.
\end{lemma}

Cho $g_i$ ($1\leq i\leq r$) là các phần tử sinh của $T_0$-môđun $T_1$.
Trước hết, giả thiết kéo theo rằng các $D_+(g_i)$ phủ $\Proj(T)$ \sref{II.2.8.1}.
Cho $(h_j)_{1\leq j\leq s}$ là một hệ các phần tử thuần nhất của $T_+$, với $\deg(h_j)=n_j$, sao cho cùng với các $g_i$, chúng tạo thành một hệ sinh của iđêan $T_+$, hay, một cách tương đương \sref{II.2.1.3}, một hệ sinh của $T$ với tư cách là một $T_0$-đại số;
nếu ta đặt $T'=\bigoplus_{n\geq0}T_1^n$, thì theo giả thiết, phần tử $h_j/g_i^{n_j}$ của vành $T_{(g_i)}$ phải thuộc vành con $T'_{(g_i)}$, và do đó tồn tại một số nguyên $k$ sao cho $T_1^k h_j\subset T_1^{k+n_j}$ với mọi $j$.
Do đó, bằng quy nạp theo số mũ $r$, ta suy ra rằng $T_1^k h_j^r \subset T'$ với mọi $r\geq1$, và, theo định nghĩa của các $h_j$, do đó ta có $T_1^k T\subset T'$.
Mặt khác, với mọi $j$, tồn tại một số nguyên $m_j$ sao cho $h_j^{m_j}$ thuộc iđêan của $T$ sinh bởi các $g_i$ \sref{II.2.3.14}, do đó $h_j^{m_j}\in T_1 T$, và
\oldpage[II]{162}
$h_j^{m_j k}\in T_1^k T\subset T'$.
Do đó tồn tại một số nguyên $m_0\geq k$ sao cho $h_j^m\in T_1^{m n_j}$ với $m\geq m_0$.
Vậy, nếu $q$ là số lớn nhất trong các số nguyên $n_j$, thì $n_0=q s m_0+k$ là số cần tìm.
Thật vậy, một phần tử của $T_n$, với $n\geq n_0$, là tổng các đơn thức thuộc $T_1^\alpha u$, trong đó $u$ là một tích các lũy thừa của các $h_j$;
nếu $\alpha\geq k$, thì từ trên suy ra rằng $T_1^\alpha u\subset T_1^n$;
trong trường hợp ngược lại, một trong các số mũ của các $h_j$ là $\geq m_0$, nên $u\in T_1^\beta v$, trong đó $\beta\geq k$ và $v$ lại là một tích các lũy thừa của các $h_j$;
khi đó ta có thể quy về trường hợp trước, và do đó suy ra rằng $T_1^\alpha u\subset T_1^n$ trong mọi trường hợp.
\end{proof}

\begin{remark}[8.2.15]
\label{II.8.2.15}
Điều kiện $S_n=S_1^n$ với $n\geq n_0$ hiển nhiên kéo theo $S_{n+1}=S_1S_n$ với $n\geq n_0$, nhưng điều ngược lại không nhất thiết đúng, ngay cả khi ta giả sử $S$ là một vành Noether.
Chẳng hạn, cho $K$ là một trường, $A=K[\bb{x}]$, và $B=K[\bb{y}]/\bb{y}^2K[\bb{y}]$, trong đó $\bb{x}$ và $\bb{y}$ là các biến bất định, với $\bb{x}$ được lấy có bậc 1 và $\bb{y}$ có bậc 2, và cho $S=A\otimes_K B$, do đó $S$ là một đại số phân bậc trên $K$ có một cơ sở gồm các phần tử $1$, $\bb{x}^n$ ($n\geq1$), và $\bb{x}^n\bb{y}$ ($n\geq0$).
Hiển nhiên rằng $S_{n+1}=S_1S_n$ với $n\geq2$, nhưng $S_1^n=K\bb{x}^n$ trong khi $S_n=K\bb{x}^n+K\bb{x}^{n-2}\bb{y}$ với $n\geq2$.
\end{remark}
\subsection{Các nón chiếu}
\label{subsection:II.8.3}

\begin{env}[8.3.1]
\label{II.8.3.1}
Cho $Y$ là một tiền lược đồ;
trong toàn bộ mục này, ta chỉ xét các \emph{$Y$-tiền lược đồ} và các \emph{$Y$-cấu xạ}.
Cho $\sh{S}$ là một đại số $\sh{O}_Y$ phân bậc quasi-coherent theo các bậc \emph{dương};
\emph{hơn nữa, ta giả sử rằng $\sh{S}_0=\sh{O}_Y$}.
Theo các ký hiệu được đưa vào trong \sref{II.8.2.2}, ta đặt
\[
\label{II.8.3.1.1}
  \widehat{\sh{S}} = \sh{S}[\bb{z}] = \sh{S} \otimes_{\sh{O}_Y} \sh{O}_Y[\bb{z}]
\tag{8.3.1.1}
\]
mà ta xét như một đại số $\sh{O}_Y$ phân bậc theo các bậc dương bằng cách định nghĩa các bậc như trong \sref{II.8.2.2.2}, sao cho, với mọi tập mở affine $U$ của $Y$, ta có
\[
  \Gamma(U,\widehat{\sh{S}}) = (\Gamma(U,\sh{S}))[\bb{z}].
\]
Trong phần tiếp theo, ta đặt
\[
\label{II.8.3.1.2}
  X = \Proj(\sh{S}),
  \quad
  C = \Spec(\sh{S}),
  \quad
  \widehat{C} = \Proj(\widehat{\sh{S}})
\tag{8.3.1.2}
\]
(trong định nghĩa của $C$, ta xét $\sh{S}$ như một đại số $\sh{O}_Y$ không phân bậc), và ta nói rằng $C$ (tương ứng, $\widehat{C}$) là \emph{nón affine} (tương ứng, \emph{nón xạ ảnh}) được xác định bởi $\sh{S}$;
ta cũng nói ``nón'' thay cho ``nón affine''.
Theo lối lạm dụng ngôn ngữ, ta còn nói rằng $C$ (tương ứng, $\widehat{C}$) là \emph{nón chiếu affine của $X$} (tương ứng, \emph{nón chiếu xạ ảnh của $X$}), ngầm hiểu rằng tiền lược đồ $X$ được cho dưới dạng $\Proj(\sh{S})$;
cuối cùng, ta nói rằng $\widehat{C}$ là \emph{bao đóng xạ ảnh} của $C$ (ngầm hiểu rằng dữ liệu $\sh{S}$ nằm trong cấu trúc của $C$).
\end{env}
\begin{proposition}[8.3.2]
\label{II.8.3.2}
Tồn tại các $Y$-cấu xạ chính tắc
\[
\label{II.8.3.2.1}
  Y \xrightarrow{\varepsilon} C \xrightarrow{i} \widehat{C}
\tag{8.3.2.1}
\]
\[
\label{II.8.3.2.2}
  X \xrightarrow{j} \widehat{C}
\tag{8.3.2.2}
\]

\oldpage[II]{163}
sao cho $\varepsilon$ và $j$ là các phép nhúng đóng, còn $i$ là một cấu xạ affine, đồng thời là một phép nhúng mở trội, và
\[
\label{II.8.3.2.3}
  i(C) = \widehat{C}\setmin j(X);
\tag{8.3.2.3}
\]
hơn nữa, $\widehat{C}$ là tiền lược đồ con đóng nhỏ nhất của $\widehat{C}$ chứa $i(C)$.
\end{proposition}

\begin{proof}
Để định nghĩa $i$, xét tập mở sau của $\widehat{C}$:
\[
\label{II.8.3.2.4}
  \widehat{C}_{\bb{z}} = \Spec(\widehat{\sh{S}}/(\bb{z}-1)\widehat{\sh{S}})
\tag{8.3.2.4}
\]
\sref{II.3.1.4}, trong đó $\bb{z}$ được đồng nhất chính tắc với một tiết diện của $\widehat{\sh{S}}$ trên $Y$.
Khi đó, đẳng cấu $i: C\simto\widehat{C}_{\bb{z}}$ tương ứng với đẳng cấu chính tắc \sref{II.8.2.3.1}
\[
  \widehat{\sh{S}}/(\bb{z}-1)\widehat{\sh{S}} \simto \sh{S}.
\]

Cấu xạ $\varepsilon$ tương ứng với đồng cấu tăng cường $\sh{S}\to\sh{S}_0=\sh{O}_Y$, có hạt nhân $\sh{S}_+$ \sref{II.1.2.7}, và vì đồng cấu này toàn ánh nên $\varepsilon$ là một phép nhúng đóng \sref{II.1.4.10}.
Cuối cùng, $j$ tương ứng \sref{II.3.5.1} với đồng cấu toàn ánh bậc không $\widehat{\sh{S}}\to\sh{S}$, đồng cấu này hạn chế thành đồng nhất trên $\sh{S}$ và bằng không trên $\bb{z}\widehat{\sh{S}}$, là hạt nhân của nó;
$j$ được xác định khắp nơi và là một phép nhúng đóng theo \sref{II.3.6.2}.

Để chứng minh các khẳng định còn lại của \sref{II.8.3.2}, hiển nhiên ta có thể giới hạn vào trường hợp $Y=\Spec(A)$ là affine và $\sh{S}=\widetilde{S}$, trong đó $S$ là một $A$-đại số phân bậc, do đó $\widehat{\sh{S}}=(\widehat{S})\supertilde$;
khi đó, các phần tử thuần nhất $f$ của $S_+$ có thể được đồng nhất với các tiết diện của $\widehat{\sh{S}}$ trên $Y$, và tập mở của $\widehat{C}$, được ký hiệu là $D_+(f)$ trong \sref{II.2.3.3}, cũng có thể được viết là $\widehat{C}_f$ \sref{II.3.1.4};
tương tự, tập mở của $C$ được ký hiệu là $D(f)$ trong \sref[I]{I.1.1.1} cũng có thể được viết là $C_f$ \sref[0]{0.5.5.2}.
Với các nhận xét này, từ \sref{II.2.3.14} và định nghĩa của $\widehat{S}$ suy ra rằng, trong trường hợp này, các tập mở $\widehat{C}_{\bb{z}}=i(C)$ và $\widehat{C}_f$ (với $f$ thuần nhất trong $S_+$) tạo thành một \emph{phủ} của $\widehat{C}$.
Hơn nữa, với các ký hiệu này,
\[
\label{II.8.3.2.5}
  i^{-1}(\widehat{C}_f) = C_f;
\tag{8.3.2.5}
\]
thật vậy, $\widehat{C}_f\cap i(C) = \widehat{C}_f\cap\widehat{C}_{\bb{z}} = \widehat{C}_{f\bb{z}} = \Spec(\widehat{S}_{(f\bb{z})})$.
Nhưng nếu $d=\deg(f)$, thì $\widehat{S}_{(f\bb{z})}$ đẳng cấu chính tắc với $(\widehat{S}_{(\bb{z})})_{f/\bb{z}^d}$ \sref{II.2.2.2}, và từ định nghĩa của đẳng cấu trong \sref{II.8.2.3.1} suy ra rằng ảnh của $(\widehat{S}_{(\bb{z})})_{f/\bb{z}^d}$ qua đẳng cấu tương ứng giữa các vành các phân thức chính là $S_f$.
Vì $C_f=\Spec(S_f)$, điều này chứng minh \sref{II.8.3.2.5} và đồng thời cho thấy cấu xạ $i$ là affine;
hơn nữa, hạn chế của $i$ lên $C_f$, được xem như một cấu xạ vào $\widehat{C}_f$, tương ứng \sref[I]{I.1.7.3} với đồng cấu chính tắc $\widehat{S}_{(f)}\to\widehat{S}_{(f\bb{z})}$, và theo những điều trên cùng với \sref{II.8.2.3.2}, ta có thể phát biểu kết quả sau:
\begin{env}[8.3.2.6]
\label{II.8.3.2.6}
Nếu $Y=\Spec(A)$ là affine và $\sh{S}=\widetilde{S}$, thì, với mọi $f$ thuần nhất trong $S_+$, $\widehat{C}_f$ được đồng nhất chính tắc với $\Spec(S_f^\leq)$, và cấu xạ $C_f\to\widehat{C}_f$ thu được bằng cách hạn chế $i$ khi đó tương ứng với đơn ánh chính tắc $S_f^\leq\to S_f$.
\end{env}

Bây giờ lưu ý rằng (khi $Y$ affine) phần bù của $\widehat{C}_{\bb{z}}$ trong $\widehat{C}=\Proj(\widehat{S})$
\oldpage[II]{164}
theo định nghĩa là tập hợp các iđêan nguyên tố phân bậc của $\widehat{S}$ chứa $\bb{z}$, và tập hợp này chính là $j(X)$ theo định nghĩa của $j$, điều đó chứng minh \sref{II.8.3.2.3}.

Cuối cùng, để chứng minh khẳng định cuối của \sref{II.8.3.2}, ta có thể giả sử $Y$ là affine.
Với các ký hiệu trên, lưu ý rằng trong vành $\widehat{S}$, $\bb{z}$ không phải là một ước của không;
vì bao đóng của $i(C)$ là $\widehat{C}$, chỉ cần chứng minh bổ đề sau:
\begin{lemma}[8.3.2.7]
\label{II.8.3.2.7}
Cho $T$ là một vành phân bậc theo các bậc dương, $Z=\Proj(T)$, và $g$ là một phần tử thuần nhất của $T$ có bậc $d>0$.
Nếu $g$ không phải là một ước của không trong $T$, thì $Z$ là tiền lược đồ con đóng nhỏ nhất của $Z$ chứa $Z_g=D_+(g)$.
\end{lemma}

Theo \sref[I]{I.4.1.9}, vấn đề có tính địa phương trên $Z$;
do đó, với mọi phần tử thuần nhất $h\in T_e$ ($e>0$), chỉ cần chứng minh rằng $Z_h$ là tiền lược đồ con đóng nhỏ nhất của $Z_h$ chứa $Z_{gh}$;
từ các định nghĩa và \sref[I]{I.4.3.2} suy ra rằng điều kiện này tương đương với việc đồng cấu chính tắc $T_{(h)}\to T_{(gh)}$ là \emph{đơn ánh}.
Nhưng đồng cấu này có thể được đồng nhất với đồng cấu chính tắc $T_{(h)}\to(T_{(h)})_{g^e/h^d}$ \sref{II.2.2.3}.
Vì $g^e$ không phải là một ước của không trong $T$, nên $g^e/h^d$ không phải là một ước của không trong $T_h$ (và càng không phải trong $T_{(h)}$), bởi vì hệ thức $(g^e/h^d)(t/h^m)=0$ (với $t\in T$ và $m>0$) kéo theo sự tồn tại của một $n>0$ sao cho $h^ng^et=0$, do đó $h^nt=0$, và vì thế $t/h^m=0$ trong $T_h$.
Điều này kết thúc chứng minh \sref[0]{0.1.2.2}.
\end{proof}
\begin{env}[8.3.3]
\label{II.8.3.3}
Ta sẽ thường đồng nhất nón affine $C$ với tiền lược đồ con được cảm sinh bởi nón xạ ảnh $\widehat{C}$ trên tập mở $i(C)$ thông qua phép nhúng mở $i$.
Tiền lược đồ con đóng của $C$ gắn với phép nhúng đóng $\varepsilon$ được gọi là \emph{tiền lược đồ các đỉnh} của $C$;
ta cũng nói rằng $\varepsilon$, vốn là một $Y$-tiết diện của $C$, là \emph{tiết diện đỉnh}, hay \emph{tiết diện không}, của $C$;
ta có thể đồng nhất $Y$ với tiền lược đồ các đỉnh của $C$ thông qua $\varepsilon$.
Hơn nữa, $i\circ\varepsilon$ là một $Y$-tiết diện của $\widehat{C}$, do đó cũng là một phép nhúng đóng \sref[I]{I.5.4.6}, tương ứng với đồng cấu chính tắc toàn ánh bậc không $\widehat{\sh{S}} = \sh{S}[\bb{z}] \to \sh{O}_Y[\bb{z}]$ \sref{II.3.1.7}, có hạt nhân $\sh{S}_+[\bb{z}] = \sh{S}_+\widehat{\sh{S}}$;
tiền lược đồ con của $\widehat{C}$ gắn với phép nhúng đóng này cũng được gọi là \emph{tiền lược đồ các đỉnh} của $\widehat{C}$, và $i\circ\varepsilon$ được gọi là \emph{tiết diện đỉnh} của $\widehat{C}$;
nó có thể được đồng nhất với $Y$ thông qua $i\circ\varepsilon$.
Cuối cùng, tiền lược đồ con đóng của $\widehat{C}$ gắn với $j$ được gọi là \emph{phần ở vô cực} của $\widehat{C}$, và có thể được đồng nhất với $X$ thông qua $j$.
\end{env}

\begin{env}[8.3.4]
\label{II.8.3.4}
Các tiền lược đồ con của $C$ (tương ứng, $\widehat{C}$) được cảm sinh trên các tập \emph{mở}
\[
\label{II.8.3.4.1}
  E = C \setmin \varepsilon(Y),
  \qquad
  \widehat{E} = \widehat{C} \setmin i(\varepsilon(Y))
\tag{8.3.4.1}
\]
được gọi lần lượt (do lạm dụng ngôn ngữ) là \emph{nón affine bỏ đỉnh} và \emph{nón xạ ảnh bỏ đỉnh} xác định bởi $\sh{S}$;
lưu ý rằng, bất chấp cách gọi này, \emph{$E$ không nhất thiết là affine} trên $Y$, và $\widehat{E}$ cũng không nhất thiết là xạ ảnh trên $Y$ \sref{II.8.4.3}.
Khi đồng nhất $C$ với $i(C)$, ta có, đối với các không gian nền,
\[
\label{II.8.3.4.2}
  C \cup \widehat{E} = \widehat{C},
  \qquad
  C \cap \widehat{E} = E
\tag{8.3.4.2}
\]
do đó $\widehat{C}$ có thể được xem là thu được bằng \emph{phép dán} các tiền lược đồ con mở $C$ và $\widehat{E}$;
hơn nữa, theo \sref{II.8.3.2.3},
\[
\label{II.8.3.4.3}
  E = \widehat{E} \setmin j(X).
\tag{8.3.4.3}
\]

\oldpage[II]{165}
Nếu $Y=\Spec(A)$ là affine, thì, với các ký hiệu của \sref{II.8.3.2},
\[
\label{II.8.3.4.4}
  E = \bigcup C_f,
  \qquad
  \widehat{E} = \bigcup \widehat{C}_f,
  \qquad
  C_f = C \cap \widehat{C}_f
\tag{8.3.4.4}
\]
trong đó $f$ chạy qua tập hợp các phần tử thuần nhất của $S_+$ (hoặc chỉ qua một tập con $M$ của tập hợp này, sao cho $M$ sinh một iđêan của $S_+$ có căn trong $S_+$ là chính $S_+$, hay tương đương, sao cho các $X_f$, với $f\in M$, phủ $X$ \sref{II.2.3.14}).
Khi đó, phép dán $C$ và $\widehat{C}_f$ dọc theo $C_f$ được xác định bởi các cấu xạ nhúng $C_f\to C$ và $C_f\to\widehat{C}_f$, mà, như ta đã thấy \sref{II.8.3.2.6}, tương ứng lần lượt với các đồng cấu chính tắc $S\to S_f$ và $S_f^\leq\to S_f$.
\end{env}
\begin{proposition}[8.3.5]
\label{II.8.3.5}
Với các ký hiệu của \sref{II.8.3.1} và \sref{II.8.3.4}, cấu xạ gắn \sref{II.3.5.1} với đơn ánh chính tắc $\vphi:\sh{S} \to \widehat{\sh{S}} = \sh{S}[\bb{z}]$ là một cấu xạ affine toàn ánh (được gọi là phép co rút chính tắc)
\[
\label{II.8.3.5.1}
  p:\widehat{E} \to X
\tag{8.3.5.1}
\]
sao cho
\[
\label{II.8.3.5.2}
  p \circ j = 1_X.
\tag{8.3.5.2}
\]
\end{proposition}

\begin{proof}
Để chứng minh mệnh đề, ta có thể giới hạn vào trường hợp $Y$ là affine.
Xét biểu thức \sref{II.8.3.4.4} của $\widehat{E}$, việc miền xác định $G(\vphi)$ của $p$ bằng $\widehat{E}$ sẽ suy ra từ khẳng định thứ nhất trong các khẳng định sau:
\begin{env}[8.3.5.3]
\label{II.8.3.5.3}
Nếu $Y=\Spec(A)$ là affine và $\sh{S}=\widetilde{S}$, thì, với mọi $f\in S_+$ thuần nhất,
\[
\label{II.8.3.5.4}
  p^{-1}(X_f) = \widehat{C}_f
\tag{8.3.5.4}
\]
và hạn chế của $p$ lên $\widehat{C}_f=\Spec(S_f^\leq)$, được xem như một cấu xạ từ $\widehat{C}_f$ đến $X_f$, tương ứng với đơn ánh chính tắc $S_{(f)}\to S_f^\leq$.
Nếu thêm $f\in S_1$, thì $\widehat{C}_f$ đẳng cấu với $X_f\otimes_{\bb{Z}}\bb{Z}[T]$ (trong đó $T$ là một biến bất định).
\end{env}

Thật vậy, công thức~\sref{II.8.3.5.4} chính là một trường hợp riêng của \sref{II.2.8.1.1}, và khẳng định thứ hai chính là định nghĩa của $\Proj(\vphi)$ khi $Y$ là affine \sref{II.2.8.1}.
Khi đó, công thức~\sref{II.8.3.5.2} và tính toàn ánh của $p$ suy ra từ việc hợp thành $\sh{S}\to\widehat{\sh{S}}\to\sh{S}$ của các đồng cấu chính tắc là đồng nhất trên $\sh{S}$.
Cuối cùng, khẳng định cuối của \sref{II.8.3.5.3} suy ra từ việc $S_f^\leq$ đẳng cấu với $S_{(f)}[T]$ khi $f\in S_1$ \sref{II.2.2.1}.
\end{proof}
\begin{corollary}[8.3.6]
\label{II.8.3.6}
Hạn chế
\[
\label{II.8.3.6.1}
  \pi: E \to X
\tag{8.3.6.1}
\]
của $p$ lên $E$ là một cấu xạ affine toàn ánh.
Nếu $Y$ là affine và $f$ thuần nhất trong $S_+$, thì
\[
\label{II.8.3.6.2}
  \pi^{-1}(X_f) = C_f
\tag{8.3.6.2}
\]
và hạn chế của $\pi$ lên $C_f$ tương ứng với đơn ánh chính tắc $S_{(f)}\to S_f$.
Nếu thêm $f\in S_1$, thì $C_f$ đẳng cấu với $X_f\otimes_{\bb{Z}}\bb{Z}[T,T^{-1}]$ (trong đó $T$ là một biến bất định).
\end{corollary}

\begin{proof}
Công thức~\sref{II.8.3.6.2} suy ra ngay từ \sref{II.8.3.5.3} và \sref{II.8.3.2.5}, và cho thấy tính toàn ánh của $\pi$;
ta đã thấy rằng phép nhúng $i$, khi hạn chế lên $C_f$, tương ứng
\oldpage[II]{166}
với đơn ánh $S_f^\leq\to S_f$ \sref{II.8.3.2}.
Cuối cùng, khẳng định cuối là hệ quả của việc, với $f\in S_1$, $S_f$ đẳng cấu với $S_{(f)}[T,T^{-1}]$ \sref{II.2.2.1}.
\end{proof}
\begin{remark}[8.3.7]
\label{II.8.3.7}
Khi $Y$ là affine, các phần tử của không gian nền của $E$ là các iđêan nguyên tố $\mathfrak{p}$ của $S$ (không nhất thiết phân bậc) không chứa $S_+$, theo định nghĩa của phép nhúng $\varepsilon$ \sref{II.8.3.2}.
Với một iđêan $\mathfrak{p}$ như vậy, các $\mathfrak{p}\cap S_n$ hiển nhiên thỏa mãn các điều kiện của \sref{II.2.1.9}, nên tồn tại duy nhất một iđêan nguyên tố \emph{phân bậc} $\mathfrak{q}$ của $S$ sao cho $\mathfrak{q}\cap S_n=\mathfrak{p}\cap S_n$ với mọi $n$;
ánh xạ $\pi:E\to X$ giữa các không gian nền khi đó được mô tả bởi hệ thức
\[
\label{II.8.3.7.1}
  \pi(\mathfrak{p}) = \mathfrak{q}.
\tag{8.3.7.1}
\]

Thật vậy, để chứng minh hệ thức này, chỉ cần xét một phần tử $f$ thuần nhất trong $S_+$ sao cho $\mathfrak{p}\in D(f)$, và nhận xét rằng $\mathfrak{q}_{(f)}$ là ảnh ngược của $\mathfrak{p}_f$ qua đơn ánh $S_{(f)}\to S_f$.
\end{remark}

\begin{corollary}[8.3.8]
\label{II.8.3.8}
Nếu $\sh{S}$ được sinh bởi $\sh{S}_1$, thì các cấu xạ $p$ và $\pi$ là kiểu hữu hạn;
với mọi $x\in X$, thớ $p^{-1}(x)$ đẳng cấu với $\Spec(\kres(x)[T])$, và thớ $\pi^{-1}(x)$ đẳng cấu với $\Spec(\kres(x)[T,T^{-1}])$.
\end{corollary}

\begin{proof}
Điều này suy ra ngay từ \sref{II.8.3.5} và \sref{II.8.3.6}, khi nhận xét rằng, nếu $Y$ là affine và $S$ được sinh bởi $S_1$, thì các $X_f$, với $f\in S_1$, tạo thành một phủ của $X$ \sref{II.2.3.14}.
\end{proof}
\begin{remark}[8.3.9]
\label{II.8.3.9}
Nón affine bỏ đỉnh tương ứng với $\sh{O}_Y$-đại số phân bậc $\sh{O}_Y[T]$ (trong đó $T$ là một biến bất định) được đồng nhất với $G_m=\Spec(\sh{O}_Y[T,T^{-1}])$, vì nó chính là $C_T$, như ta đã thấy trong \sref{II.8.3.2} (xem \sref{II.8.4.4} cho một kết quả tổng quát hơn).
Tiền lược đồ này được trang bị một cách chính tắc một cấu trúc \textquotedblleft\emph{$Y$-lược đồ nhóm giao hoán}\textquotedblright.
Khái niệm này sẽ được trình bày chi tiết về sau, nhưng ở đây có thể được tóm tắt ngắn gọn như sau.
Một $Y$-lược đồ nhóm là một $Y$-lược đồ $G$ được trang bị hai $Y$-cấu xạ, $p:G\times_Y G\to G$ và $s:G\to G$, thỏa mãn các điều kiện tương tự về mặt hình thức với các tiên đề của luật hợp thành và phép lấy nghịch đảo trong một nhóm: biểu đồ
\[
  \xymatrix{
    G \times G \times G
      \ar[r]^{p \times 1}
      \ar[d]_{1 \times p}
  & G \times G
      \ar[d]^{p}
  \\G \times G
      \ar[r]_{p}
  & G
  }
\]
phải giao hoán (\textquotedblleft tính kết hợp\textquotedblright), và phải có một điều kiện tương ứng với việc, trong các nhóm, các ánh xạ
\[
  (x,y)
  \mapsto
  (x,x^{-1},y)
  \mapsto
  (x,x^{-1}y)
  \mapsto
  x(x^{-1}y)
\]
và
\[
  (x,y)
  \mapsto
  (x,x^{-1},y)
  \mapsto
  (x,yx^{-1})
  \mapsto
  (yx^{-1})x
\]
đều phải rút gọn thành $(x,y)\mapsto y$;
chẳng hạn, dãy cấu xạ tương ứng với ánh xạ hợp thành thứ nhất là
\[
  G \times G
  \xrightarrow{(1,s) \times 1}
  G \times G \times G
  \xrightarrow{1 \times p}
  G\times G
  \xrightarrow{p}
  G
\]
và người đọc sẽ viết dãy thứ hai một cách tương tự.

\oldpage[II]{167}
Hiển nhiên \sref[I]{I.3.4.3}, việc cho một cấu trúc $Y$-lược đồ nhóm trên một $Y$-lược đồ $G$ tương đương với việc cho, đối với \emph{mọi} $Y$-tiền lược đồ $Z$, một cấu trúc \emph{nhóm} trên tập hợp $\Hom_Y(Z,G)$, sao cho, với mọi $Y$-cấu xạ $Z\to Z'$, ánh xạ tương ứng $\Hom_Y(Z',G)\to\Hom_Y(Z,G)$ là một đồng cấu nhóm.
Trong trường hợp riêng của $G_m$ đang xét, $\Hom_Y(Z,G_m)$ được đồng nhất với tập hợp các $Z$-tiết diện của $Z\times_Y G_m$ \sref[I]{I.3.3.14}, do đó với tập hợp các $Z$-tiết diện của $\Spec(\sh{O}_Z[T,T^{-1}])$;
cuối cùng, lập luận như trong \sref[I]{I.3.3.15} cho thấy tập hợp này được đồng nhất một cách chính tắc với tập hợp các phần tử \emph{khả nghịch} của vành $\Gamma(Z,\sh{O}_Z)$, và cấu trúc nhóm trên tập hợp này đến từ phép nhân trong vành $\Gamma(Z,\sh{O}_Z)$.
Người đọc có thể kiểm tra rằng các cấu xạ $p$ và $s$ ở trên được xác định như sau: theo \sref{II.1.2.7} và \sref{II.1.4.6}, chúng tương ứng với các đồng cấu $\sh{O}_Y$-đại số
\begin{align*}
  \pi: &\sh{O}_Y[T,T^{-1}] \to \sh{O}_Y[T,T^{-1},T',T'^{-1}]
\\\sigma: &\sh{O}_Y[T,T^{-1}] \to \sh{O}_Y[T,T^{-1}]
\end{align*}
và được xác định hoàn toàn bởi các dữ liệu $\pi(T)=TT'$ và $\sigma(T)=T^{-1}$.

Như vậy, $G_m$ có thể được xem là một \textquotedblleft\emph{miền toán tử phổ dụng}\textquotedblright\ cho mọi \emph{nón affine} $C=\Spec(\sh{S})$, trong đó $\sh{S}$ là một $\sh{O}_Y$-đại số quasi-coherent phân bậc với các bậc dương.
Điều này có nghĩa là ta có thể xác định một cách chính tắc một $Y$-cấu xạ $G_m\times_Y C\to C$ có các tính chất hình thức của một luật ngoài trên một tập hợp được trang bị một nhóm toán tử;
hoặc, như ở trên đối với các lược đồ nhóm, ta có thể cho, đối với mọi $Y$-tiền lược đồ $Z$, một luật ngoài trên $\Hom_Y(Z,C)$ có nhóm $\Hom_Y(Z,G_m)$ làm tập hợp các toán tử, với các tiên đề thông thường của tập hợp được trang bị một nhóm toán tử và một điều kiện tương thích đối với các $Y$-cấu xạ $Z\to Z'$.
Trong trường hợp đang xét, cấu xạ $G_m\times_Y C\to C$ được xác định bởi một đồng cấu $\sh{O}_Y$-đại số $\sh{S} \to \sh{S}\otimes_{\sh{O}_Y}\sh{O}_Y[T,T^{-1}] = \sh{S}[T,T^{-1}]$, gắn với mỗi tiết diện $s_n\in\Gamma(U,\sh{S}_n)$ (trong đó $U$ là một tập mở của $Y$) tiết diện $s_n T^n\in\Gamma(U,\sh{S}\otimes_{\sh{O}_Y}\sh{O}_Y[T,T^{-1}])$.

Ngược lại, giả sử cho một $\sh{O}_Y$-đại số quasi-coherent $\sh{S}$, \emph{tiên nghiệm chưa phân bậc}, và, trên $C=\Spec(\sh{S})$, một cấu trúc \textquotedblleft\emph{$Y$-lược đồ tập hợp được trang bị một nhóm toán tử}\textquotedblright\ có $Y$-lược đồ nhóm $G_m$ làm miền toán tử;
khi đó ta thu được một cách chính tắc một \emph{phân bậc} $\sh{O}_Y$-đại số trên $\sh{S}$.
Thật vậy, việc cho một $Y$-cấu xạ $G_m\times_Y C\to C$ tương đương với việc cho một đồng cấu $\sh{O}_Y$-đại số $\psi:\sh{S}\to\sh{S}[T,T^{-1}]$, có thể được viết thành $\psi=\sum_{n\in\bb{Z}}\psi_n T^n$, trong đó các $\psi_n:\sh{S}\to\sh{S}$ là các đồng cấu $\sh{O}_Y$-môđun (với $\psi_n(s)=0$ trừ một số hữu hạn chỉ số $n$ trên một lân cận của mỗi điểm, đối với mọi tiết diện $s\in\Gamma(U,\sh{S})$ trên một tập mở $U$ bất kỳ của $Y$).
Khi đó ta kiểm tra được rằng các tiên đề của tập hợp được trang bị một nhóm toán tử kéo theo việc các $\psi_n(\sh{S})=\sh{S}_n$ xác định một phân bậc (với các bậc nguyên tùy ý) $\sh{O}_Y$-đại số trên $\sh{S}$, với các $\psi_n$ là các phép chiếu tương ứng.
Như vậy ta thu được khái niệm cấu trúc \textquotedblleft\emph{nón affine}\textquotedblright\ trên mọi $Y$-lược đồ affine, được xác định theo cách \textquotedblleft hình học\textquotedblright\ mà không cần dùng đến một phân bậc có trước.
\oldpage[II]{168}
Ở đây chúng tôi không phát triển thêm quan điểm này, và dành cho người đọc việc phát biểu chính xác các định nghĩa và kết quả tương ứng với những chỉ dẫn trên.
\end{remark}
\subsection{Bao đóng xạ ảnh của một phân thớ vectơ}
\label{subsection:II.8.4}

\begin{env}[8.4.1]
\label{II.8.4.1}
Cho $Y$ là một tiền lược đồ và $\sh{E}$ là một $\sh{O}_Y$-môđun quasi-coherent.
Nếu lấy $\sh{S}$ là $\sh{O}_Y$-đại số phân bậc $\bb{S}_{\sh{O}_Y}(\sh{E})$, thì định nghĩa~\sref{II.8.3.1.1} cho thấy $\widehat{\sh{S}}$ được đồng nhất với $\bb{S}_{\sh{O}_Y}(\sh{E}\oplus\sh{O}_Y)$.
Vì nón affine $\Spec(\sh{S})$ xác định bởi $\sh{S}$ theo định nghĩa là $\bb{V}(\sh{E})$, còn $\Proj(\sh{S})$ theo định nghĩa là $\bb{P}(\sh{E})$, ta có:
\end{env}
\begin{proposition}[8.4.2]
\label{II.8.4.2}
Bao đóng xạ ảnh của một phân thớ vectơ $\bb{V}(\sh{E})$ trên $Y$ đẳng cấu chính tắc với $\bb{P}(\sh{E}\oplus\sh{O}_Y)$, và phần ở vô cực của lược đồ vừa nêu đẳng cấu chính tắc với $\bb{P}(\sh{E})$.
\end{proposition}
\begin{remark}[8.4.3]
\label{II.8.4.3}
Chẳng hạn, lấy $\sh{E}=\sh{O}_Y^r$ với $r\geq2$;
khi đó các nón bỏ đỉnh $E$ và $\widehat{E}$ xác định bởi $\sh{S}$ không affine cũng không xạ ảnh trên $Y$ nếu $Y\neq\emp$.
Khẳng định thứ hai là hiển nhiên, vì $\widehat{C}=\bb{P}(\sh{O}_Y^{r+1})$ là xạ ảnh trên $Y$, và các không gian nền của $E$ và $\widehat{E}$ là các tập mở không đóng trong $\widehat{C}$; do đó các phép nhúng chính tắc $E\to\widehat{C}$ và $\widehat{E}\to\widehat{C}$ không xạ ảnh \sref{II.5.5.3}, và kết luận suy ra từ \sref{II.5.5.5}[v].
Mặt khác, giả sử $Y=\Spec(A)$ là affine và, chẳng hạn, $r=2$, thì $C=\Spec(A[T_1,T_2])$, còn $E$ là tiền lược đồ cảm sinh bởi $C$ trên tập mở $D(T_1)\cup D(T_2)$;
nhưng ta đã thấy rằng tiền lược đồ này không affine \sref[I]{I.5.5.11};
\emph{a fortiori} $\widehat{E}$ không thể affine, vì $E$ là tập mở nơi tiết diện $\bb{z}$ trên $\widehat{E}$ không triệt tiêu \sref{II.8.3.2}.

Tuy nhiên:
\end{remark}
\begin{proposition}[8.4.4]
\label{II.8.4.4}
Nếu $\sh{L}$ là một $\sh{O}_Y$-môđun khả nghịch, thì tồn tại các đẳng cấu chính tắc cho các nón bỏ đỉnh $E$ và $\widehat{E}$ tương ứng với $C=\bb{V}(\sh{L})$:
\[
\label{eq:2.8.4.4.1}
  \Spec\left(
    \bigoplus_{n\in\bb{Z}}\sh{L}^{\otimes n}
  \right)
  \simto
  E
\tag{8.4.4.1}
\]
\[
\label{eq:2.8.4.4.2}
  \bb{V}(\sh{L}^{-1})
  \simto
  \widehat{E}.
\tag{8.4.4.2}
\]

Hơn nữa, tồn tại một đẳng cấu chính tắc từ bao đóng xạ ảnh của $\bb{V}(\sh{L})$ lên bao đóng xạ ảnh của $\bb{V}(\sh{L}^{-1})$, biến tiết diện không (tương ứng, phần ở vô cực) của đối tượng thứ nhất thành phần ở vô cực (tương ứng, tiết diện không) của đối tượng thứ hai.
\end{proposition}

\begin{proof}
Ở đây ta có $\sh{S}=\bigoplus_{n\geq0}\sh{L}^{\otimes n}$;
đơn ánh chính tắc
\[
  \sh{S} \to \bigoplus_{n\in\bb{Z}} \sh{L}^{\otimes n}
\]
xác định một cấu xạ trội chính tắc
\[
\label{eq:2.8.4.4.3}
  \Spec\left(
    \bigoplus_{n\in\bb{Z}} \sh{L}^{\otimes n}
  \right)
  \to
  \bb{V}(\sh{L})
  =
  \Spec\left(
    \bigoplus_{n\geq0} \sh{L}^{\otimes n}
  \right)
\tag{8.4.4.3}
\]
và chỉ cần chứng minh rằng cấu xạ này là một đẳng cấu từ lược đồ $\Spec(\bigoplus_{n\in\bb{Z}} \sh{L}^{\otimes n})$ lên $E$.
Vì vấn đề là địa phương trên $Y$, ta có thể giả sử $Y=\Spec(A)$ là affine
\oldpage[II]{169}
và $\sh{L}=\sh{O}_Y$, nên $\sh{S}=(A[T])\supertilde$ và $\bigoplus_{n\in\bb{Z}}\sh{L}^{\otimes n} = (A[T,T^{-1}])\supertilde$.
Nhưng $A[T,T^{-1}]$ là vành các phân thức $A[T]_T$ của $A[T]$, nên \eref{eq:2.8.4.4.3} đồng nhất $\Spec(\bigoplus_{n\in\bb{Z}}\sh{L}^{\otimes n})$ với tiền lược đồ cảm sinh bởi $C=\bb{V}(\sh{L})$ trên tập mở $D(T)$;
phần bù $V(T)$ của tập mở này trong $C$ là không gian nền của tiền lược đồ con đóng của $C$ xác định bởi iđêan $TA[T]$, chính là tiết diện không của $C$, do đó $E=D(T)$.

Mặt khác, đẳng cấu \eref{eq:2.8.4.4.2} sẽ suy ra từ khẳng định cuối, vì $\bb{V}(\sh{L}^{-1})$ là phần bù của phần ở vô cực trong bao đóng xạ ảnh của nó, còn $\widehat{E}$ là phần bù của tiết diện không trong bao đóng xạ ảnh của $C=\bb{V}(\sh{L})$.
Các bao đóng xạ ảnh này lần lượt là $\bb{P}(\sh{L}^{-1}\oplus\sh{O}_Y)$ và $\bb{P}(\sh{L}\oplus\sh{O}_Y)$;
nhưng ta có thể viết $\sh{L}\oplus\sh{O}_Y = \sh{L}\otimes(\sh{L}^{-1}\oplus\sh{O}_Y)$.
Sự tồn tại của đẳng cấu chính tắc cần tìm khi đó suy ra từ \sref{II.4.1.4}, và chỉ còn phải kiểm tra rằng đẳng cấu này hoán đổi các tiết diện không và các phần ở vô cực.
Để làm điều đó, ta có thể quy về trường hợp $Y=\Spec(A)$ là affine, $L=Ac$ và $L^{-1}=Ac'$, trong đó đẳng cấu chính tắc $L\otimes L^{-1}\to A$ gửi $c\otimes c'$ tới phần tử $1$ của $A$.
Khi đó $\bb{S}(L\oplus A)$ là tích tenxơ của $A[\bb{z}]$ với $\bigoplus_{n\geq0}Ac^{\otimes n}$, còn $\bb{S}(L^{-1}\oplus A)$ là tích tenxơ của $A[\bb{z}]$ với $\bigoplus_{n\geq0}Ac^{'\otimes n}$, và đẳng cấu xác định trong \sref{II.4.1.4} gửi $\bb{z}^h\otimes c^{'\otimes(n-h)}$ tới phần tử $\bb{z}^{n-h}\otimes c^{\otimes h}$.
Nhưng trong $\bb{P}(\sh{L}^{-1}\oplus\sh{O}_Y)$, phần ở vô cực là tập hợp các điểm nơi tiết diện $\bb{z}$ triệt tiêu, còn tiết diện không là tập hợp các điểm nơi tiết diện $c'$ triệt tiêu;
vì ta có các định nghĩa tương tự đối với $\bb{P}(\sh{L}\oplus\sh{O}_Y)$, kết luận suy ra ngay từ mô tả trên.
\end{proof}
\subsection{Tính hàm tử}
\label{subsection:II.8.5}

\begin{env}[8.5.1]
\label{II.8.5.1}
Cho $Y$ và $Y'$ là các tiền lược đồ, $q:Y'\to Y$ là một cấu xạ, và $\sh{S}$ (tương ứng, $\sh{S}'$) là một $\sh{O}_Y$-đại số quasi-coherent phân bậc theo các bậc \emph{dương} (tương ứng, một $\sh{O}_{Y'}$-đại số quasi-coherent phân bậc theo các bậc \emph{dương}).
Xét một $q$-cấu xạ của các đại số phân bậc
\[
\label{eq:2.8.5.1.1}
  \vphi: \sh{S} \to \sh{S}'.
\tag{8.5.1.1}
\]

Ta biết \sref{II.1.5.6} rằng nó tương ứng một cách chính tắc với một cấu xạ
\[
  \Phi = \Spec(\vphi): \Spec(\sh{S}') \to \Spec(\sh{S})
\]
sao cho sơ đồ
\[
\label{eq:2.8.5.1.2}
  \xymatrix{
    C'
      \ar[r]^{\Phi}
      \ar[d]
  & C
      \ar[d]
  \\Y'
      \ar[r]_{q}
  & Y
  }
\tag{8.5.1.2}
\]
giao hoán, trong đó ta viết $C=\Spec(\sh{S})$ và $C'=\Spec(\sh{S}')$.
\emph{Giả sử thêm rằng $\sh{S}_0=\sh{O}_Y$ và $\sh{S}'_0=\sh{O}_{Y'}$};
gọi $\varepsilon:Y\to C$ và $\varepsilon':Y'\to C'$ là các phép nhúng chính tắc \sref{II.8.3.2};
khi đó ta có sơ đồ giao hoán
\[
\label{eq:2.8.5.1.3}
  \xymatrix{
    Y'
      \ar[r]^{q}
      \ar[d]_{\varepsilon'}
  & Y
      \ar[d]^{\varepsilon}
  \\C'
      \ar[r]_{\Phi}
  & C
  }
\tag{8.5.1.3}
\]
\oldpage[II]{170}
tương ứng với sơ đồ
\[
  \xymatrix{
    \sh{S}
      \ar[r]^{\vphi}
      \ar[d]
  & \sh{S}'
      \ar[d]
  \\\sh{O}_Y
      \ar[r]
  & \sh{O}_{Y'}
  }
\]
trong đó các mũi tên thẳng đứng là các đồng cấu tăng cường; tính giao hoán do đó suy ra từ giả thiết rằng $\vphi$ là một đồng cấu của các đại số \emph{phân bậc}.
\end{env}
\begin{proposition}[8.5.2]
\label{II.8.5.2}
Nếu $E$ (tương ứng, $E'$) là nón affine bỏ đỉnh xác định bởi $\sh{S}$ (tương ứng, $\sh{S}'$), thì $\Phi^{-1}(E)\subset E'$;
nếu, hơn nữa, $\Proj(\vphi):G(\vphi)\to\Proj(\sh{S})$ được xác định khắp nơi (hay tương đương, nếu $G(\vphi)=\Proj(\sh{S}')$), thì $\Phi^{-1}(E)=E'$, và ngược lại.
\end{proposition}

\begin{proof}
Khẳng định thứ nhất suy ra từ tính giao hoán của \eref{eq:2.8.5.1.3}.
Để chứng minh khẳng định thứ hai, ta có thể quy về trường hợp $Y=\Spec(A)$ và $Y'=\Spec(A')$ là affine, $\sh{S}=\widetilde{S}$ và $\sh{S}'=\widetilde{S'}$.
Với mọi $f$ thuần nhất thuộc $S_+$, đặt $f'=\vphi(f)$, ta có $\Phi^{-1}(C_f)=C'_{f'}$ \sref[I]{I.2.2.4.1};
nói rằng $G(\vphi)=\Proj(S')$ có nghĩa là căn (trong $S'_+$) của iđêan sinh bởi các $f'=\vphi(f)$ là chính $S'_+$ (\sref{II.2.8.1} và \sref{II.2.3.14}), và điều này tương đương với việc các $C'_{f'}$ phủ $E'$ \sref{II.8.3.4.4}.
\end{proof}
\begin{env}[8.5.3]
\label{II.8.5.3}
$q$-cấu xạ $\vphi$ mở rộng một cách chính tắc thành một $q$-cấu xạ của các đại số phân bậc
\[
\label{II.8.5.3.1}
  \widehat{\vphi}: \widehat{\sh{S}} \to \widehat{\sh{S}'}
\tag{8.5.3.1}
\]
bằng cách đặt $\widehat{\vphi}(\bb{z})=\bb{z}$.
Từ đó suy ra một cấu xạ
\[
  \widehat{\Phi} = \Proj(\widehat{\vphi}) : G(\widehat{\vphi}) \to \widehat{C} = \Proj(\widehat{\sh{S}})
\]
sao cho sơ đồ
\[
  \xymatrix{
    G(\widehat{\vphi})
      \ar[r]^{\widehat{\Phi}}
      \ar[d]
  & \widehat{C}
      \ar[d]
  \\Y'
      \ar[r]_{q}
  & Y
  }
\]
giao hoán \sref{II.3.5.6}.
Từ các định nghĩa suy ra ngay rằng, nếu kí hiệu $i:C\to\widehat{C}$ và $i':C'\to\widehat{C'}$ là các phép nhúng mở chính tắc \sref{II.8.3.2}, thì $i'(C')\subset G(\widehat{\vphi})$, và sơ đồ
\[
\label{eq:2.8.5.3.2}
  \xymatrix{
    C'
      \ar[r]^{\Phi}
      \ar[d]_{i'}
  & C
      \ar[d]^{i}
  \\G(\widehat{\vphi})
      \ar[r]_{\widehat{\Phi}}
  & \widehat{C}
  }
\tag{8.5.3.2}
\]
giao hoán.
Cuối cùng, nếu đặt $X=\Proj(\sh{S})$ và $X'=\Proj(\sh{S}')$, và nếu $j:X\to\widehat{C}$ và $j':X'\to\widehat{C'}$ là các phép nhúng đóng chính tắc \sref{II.8.3.2}, thì từ định nghĩa của các phép nhúng này cũng suy ra rằng $j'(G(\vphi))\subset G(\widehat{\vphi})$, và sơ đồ
\oldpage[II]{171}
\[
\label{II.8.5.3.3}
  \xymatrix{
    G(\vphi)
      \ar[r]^{\Proj(\vphi)}
      \ar[d]_{j'}
  & X
      \ar[d]^{j}
  \\G(\widehat{\vphi})
      \ar[r]_{\widehat{\Phi}}
  & \widehat{C}
  }
\tag{8.5.3.3}
\]
giao hoán.
\end{env}
\begin{proposition}[8.5.4]
\label{II.8.5.4}
Nếu $\widehat{E}$ (tương ứng, $\widehat{E'}$) là nón xạ ảnh bỏ đỉnh xác định bởi $\sh{S}$ (tương ứng, bởi $\sh{S}'$), thì $\widehat{\Phi}^{-1}(\widehat{E}) \subset \widehat{E'}$;
hơn nữa, nếu $p:\widehat{E}\to X$ và $p':\widehat{E'}\to X'$ là các phép co rút chính tắc, thì $p'(\widehat{\Phi}^{-1}(\widehat{E})) \subset G(\vphi)$, và sơ đồ
\[
\label{II.8.5.4.1}
  \xymatrix{
    \widehat{\Phi}^{-1}(\widehat{E})
      \ar[r]^{\widehat{\Phi}}
      \ar[d]_{p'}
  & \widehat{E}
      \ar[d]^{p}
  \\G(\vphi)
      \ar[r]_{\Proj(\vphi)}
  & X
  }
\tag{8.5.4.1}
\]
giao hoán.
Nếu $\Proj(\vphi)$ được xác định khắp nơi, thì $\widehat{\Phi}$ cũng vậy, và ta có $\widehat{\Phi}^{-1}(\widehat{E}) = \widehat{E'}$.
\end{proposition}

\begin{proof}
Khẳng định thứ nhất suy ra từ tính giao hoán của các sơ đồ \eref{eq:2.8.5.1.3} và \eref{eq:2.8.5.3.2}, còn hai khẳng định tiếp theo suy ra từ định nghĩa của các phép co rút chính tắc \sref{II.8.3.5} và định nghĩa của $\widehat{\vphi}$.
Để thấy rằng $\widehat{\Phi}$ được xác định khắp nơi bất cứ khi nào $\Proj(\vphi)$ được xác định khắp nơi, ta có thể quy về trường hợp $Y=\Spec(A)$ và $Y'=\Spec(A')$ là affine, và $\sh{S}=\widetilde{S}$ cùng $\sh{S}'=\widetilde{S'}$;
giả thiết nói rằng, khi $f$ chạy qua tập hợp các phần tử thuần nhất của $S_+$, căn trong $S'_+$ của iđêan sinh trong $S'_+$ bởi các $\vphi(f)$ là chính $S'_+$;
do đó suy ra ngay rằng căn trong $(S'[\bb{z}])_+$ của iđêan sinh bởi $\bb{z}$ và các $\vphi(f)$ là chính $(S'[\bb{z}])_+$, từ đó có khẳng định;
điều này đồng thời cho thấy $\widehat{E'}$ là hợp của các $\widehat{C'}_{\vphi(f)}$, và do đó bằng $\widehat{\Phi}^{-1}(\widehat{E})$.
\end{proof}
\begin{corollary}[8.5.5]
\label{II.8.5.5}
Mỗi khi $\Proj(\vphi)$ được xác định khắp nơi, ảnh ngược qua $\widehat{\Phi}$ của không gian nền của phần ở vô cực (tương ứng, của tiền lược đồ các đỉnh) của $\widehat{C}$ là không gian nền của phần ở vô cực (tương ứng, của tiền lược đồ các đỉnh) của $\widehat{C'}$.
\end{corollary}

\begin{proof}
Điều này suy ra ngay từ \sref{II.8.5.4} và \sref{II.8.5.2}, khi tính đến các hệ thức \sref{II.8.3.4.1} và \sref{II.8.3.4.2}.
\end{proof}
\subsection{Một đẳng cấu chính tắc cho các nón bỏ đỉnh}
\label{subsection:II.8.6}

\begin{env}[8.6.1]
\label{II.8.6.1}
Cho $Y$ là một tiền lược đồ, $\sh{S}$ là một $\sh{O}_Y$-đại số quasi-coherent phân bậc theo các bậc dương \emph{sao cho $\sh{S}_0=\sh{O}_Y$}, và $X$ là $Y$-lược đồ $\Proj(\sh{S})$.
Ta sẽ áp dụng các kết quả của \sref{subsection:II.8.5} cho trường hợp $Y'=X$ và $q:X\to Y$ là cấu xạ cấu trúc;
đặt
\[
\label{II.8.6.1.1}
  \sh{S}_X=\bigoplus_{n\in\bb{Z}}\sh{O}_X(n)
  \tag{8.6.1.1}
\]
\oldpage[II]{172}
là một $\sh{O}_X$-đại số phân bậc quasi-coherent, trong đó phép nhân được xác định nhờ các đồng cấu chính tắc \sref{II.3.2.6.1}
\[
  \sh{O}_X(m)\otimes_{\sh{O}_X}\sh{O}_X(n)\to\sh{O}_X(m+n)
\]
mà tính kết hợp được bảo đảm bởi sơ đồ giao hoán trong \sref{II.2.5.11.4}.
Cho $\sh{S}'$ là $\sh{O}_X$-đại số con phân bậc quasi-coherent $\sh{S}_X^\geq=\bigoplus_{n\geq0}\sh{O}_X(n)$ của $\sh{S}_X$, với các bậc dương.

Cuối cùng, xét $q$-cấu xạ chính tắc
\[
\label{II.8.6.1.2}
  \alpha: \sh{S} \to \sh{S}_X^\geq
\tag{8.6.1.2}
\]
được xác định trong \sref{II.3.3.2.3} như một đồng cấu $\sh{S}\to q_*(\sh{S}_X)$, nhưng hiển nhiên gửi $\sh{S}$ vào $q_*(\sh{S}_X^\geq)$.
Đặt
\[
\label{II.8.6.1.3}
  C_X = \Spec(\sh{S}_X^\geq),
  \quad
  \widehat{C}_X = \Proj(\sh{S}_X^\geq[\bb{z}]),
  \quad
  X' = \Proj(\sh{S}_X^\geq)
\tag{8.6.1.3}
\]
và lần lượt ký hiệu các nón affine và xạ ảnh bỏ đỉnh tương ứng bởi $E_X$ và $\widehat{E}_X$;
ký hiệu các cấu xạ chính tắc bởi $\varepsilon_X:X\to C_X$, $i_X:C_X\to\widehat{C}_X$, $j_X:X'\to\widehat{C}_X$, $p_X:\widehat{E}_X\to X'$ và $\pi_X:E_X\to X'$, như đã được xác định trong \sref{subsection:II.8.3}.
\end{env}
\begin{proposition}[8.6.2]
\label{II.8.6.2}
Cấu xạ cấu trúc $u:X'\to X$ là một \emph{đẳng cấu}, và cấu xạ $\Proj(\alpha)$ được xác định khắp nơi và đồng nhất với $u$.
Cấu xạ $\Proj(\widehat{\alpha}):\widehat{C}_X\to\widehat{C}$ được xác định khắp nơi, và các hạn chế của nó lên $\widehat{E}_X$ và $E_X$ là các \emph{đẳng cấu} tương ứng lên $\widehat{E}$ và $E$.
Cuối cùng, nếu ta đồng nhất $X'$ với $X$ qua $u$, thì các cấu xạ $p_X$ và $\pi_X$ được đồng nhất với các cấu xạ cấu trúc của các $X$-tiền lược đồ $\widehat{E}_X$ và $E_X$.
\end{proposition}

\begin{proof}
Hiển nhiên ta có thể quy về trường hợp $Y=\Spec(A)$ là affine và $\sh{S}=\widetilde{S}$;
khi đó $X$ là hợp của các tập mở affine $X_f$, trong đó $f$ chạy qua tập các phần tử thuần nhất của $S_+$, và vành của mỗi $X_f$ là $S_{(f)}$.
Từ \sref{II.8.2.7.1} suy ra
\[
\label{II.8.6.2.1}
  \Gamma(X_f, \sh{S}_X^\geq) = S_f^\geq.
\tag{8.6.2.1}
\]

Do đó $u^{-1}(X_f)=\Proj(S_f^\geq)$.
Nhưng nếu $f\in S_d$ ($d>0$), thì $\Proj(S_f^\geq)$ đẳng cấu chính tắc với $\Proj((S_f^\geq)^{(d)})$ \sref{II.2.4.7}, và ta cũng biết rằng $(S_f^\geq)^{(d)}=(S^{(d)})_f^\geq$ có thể được đồng nhất với $S_{(f)}[T]$ \sref{II.2.2.1} qua ánh xạ $T\mapsto f/1$;
do đó, từ \sref{II.3.1.7}, ta kết luận rằng cấu xạ cấu trúc $u^{-1}(X_f)\to X_f$ là một đẳng cấu, từ đó suy ra khẳng định thứ nhất.
Để chứng minh khẳng định thứ hai, lưu ý rằng hạn chế $u^{-1}(X_f)\cap G(\alpha)\to X=\Proj(S)$ của $\Proj(\alpha)$ tương ứng với ánh xạ chính tắc $x\mapsto x/1$ từ $S$ đến $S_f^\geq$ \sref{II.2.6.2};
trước hết, từ đó ta suy ra $G(\alpha)=X'$, rồi, có tính đến đẳng thức $u^{-1}(X_f)=(u^{-1}(X_f))_{f/1}$, từ \sref{II.2.8.1.1} suy ra rằng ảnh của $u^{-1}(X_f)$ qua $\Proj(\alpha)$ nằm trong $X_f$, và hạn chế của $\Proj(\alpha)$ lên $u^{-1}(X_f)$, được xem như một cấu xạ đến $X_f=\Spec(S_{(f)})$, quả thực đồng nhất với hạn chế của $u$.
Cuối cùng, áp dụng \sref{II.8.3.5.4} cho $p_X$ thay cho $p$, ta thấy rằng $p_X^{-1}(u^{-1}(X_f)) = \Spec((S_f^\geq)_{f/1}^\leq)$, và theo \sref{II.8.5.4.1}, tập mở này là ảnh ngược qua $\Proj(\widehat{\alpha})$ của $p^{-1}(X_f)=\Spec(S_f^\leq)$ \sref{II.8.3.5.3}.
Có tính đến \sref{II.8.2.3.2}, hạn chế của $\Proj(\widehat{\alpha})$ lên $p_X^{-1}(u^{-1}(X_f))$ tương ứng với đẳng cấu nghịch đảo của \sref{II.8.2.7.2}, được hạn chế lên $S_f^\leq$, từ đó suy ra khẳng định thứ ba;
khẳng định cuối cùng là hiển nhiên theo định nghĩa.

\oldpage[II]{173}
Ta cũng lưu ý rằng từ sơ đồ giao hoán trong \eref{eq:2.8.5.3.2} suy ra rằng \emph{phép hạn chế lên $C_X$ của $\Proj(\widehat{\alpha})$ chính là cấu xạ $\Spec(\alpha)$}.
\end{proof}
\begin{corollary}[8.6.3]
\label{II.8.6.3}
Được xem như các $X$-lược đồ, $\widehat{E}_X$ đẳng cấu chính tắc với $\Spec(\sh{S}_X^\leq)$, và $E_X$ với $\Spec(\sh{S}_X)$.
\end{corollary}

\begin{proof}
Vì ta biết các cấu xạ $p_X$ và $\pi_X$ là affine (\sref{II.8.3.5} và \sref{II.8.3.6}), theo \sref{II.1.3.1}, chỉ cần chứng minh hệ quả trong trường hợp $Y=\Spec(A)$ là affine và $\sh{S}=\widetilde{S}$.
Khẳng định thứ nhất suy ra từ sự tồn tại của các đẳng cấu chính tắc \sref{II.8.2.7.2} $(S_f^\geq)_{f/1}^\leq \simto S_f^\leq$ và từ việc các đẳng cấu này tương thích với phép chuyển từ $f$ sang $fg$ (trong đó $f$ và $g$ thuần nhất trong $S_+$).
Tương tự, áp dụng \sref{II.8.3.6.2} cho $\pi_X$ thay cho $\pi$, ta thấy $\pi_X^{-1}(u^{-1}(X_f)) = \Spec((S_f^\geq)_{f/1})$ với $f$ thuần nhất trong $S_+$, và khẳng định thứ hai suy ra từ sự tồn tại của các đẳng cấu chính tắc \sref{II.8.2.7.2} $(S_f^\geq)_{f/1} \simto S_f$.

Do đó, có thể nói rằng $\widehat{C}_X$, được xem như một $X$-lược đồ, thu được bằng \emph{phép dán} các $X$-lược đồ affine $C_X=\Spec(\sh{S}_X^\geq)$ và $\widehat{E}_X=\Spec(\sh{S}_X^\leq)$ trên $X$, trong đó giao của hai $X$-lược đồ affine là tập mở $E_X=\Spec(\sh{S}_X)$.
\end{proof}
\begin{corollary}[8.6.4]
\label{II.8.6.4}
Giả sử $\sh{O}_X(1)$ là một $\sh{O}_X$-môđun khả nghịch, và $\sh{S}_X$ đẳng cấu với $\bigoplus_{n\in\bb{Z}}(\sh{O}_X(1))^{\otimes n}$ (điều này đặc biệt đúng mỗi khi $\sh{S}$ được sinh bởi $\sh{S}_1$ (\sref{II.3.2.5} và \sref{II.3.2.7})).
Khi đó nón xạ ảnh bỏ đỉnh $\widehat{E}$ có thể được đồng nhất với phân thớ vectơ hạng 1 $\bb{V}(\sh{O}_X(-1))$ trên $X$, còn nón affine bỏ đỉnh $E$ với tiền lược đồ con của phân thớ vectơ này được cảm sinh trên phần bù của tiết diện không.
Với cách đồng nhất này, phép co rút chính tắc $\widehat{E}\to X$ được đồng nhất với cấu xạ cấu trúc của $X$-lược đồ $\bb{V}(\sh{O}_X(-1))$.
Cuối cùng, tồn tại một $Y$-cấu xạ chính tắc $\bb{V}(\sh{O}_X(1))\to C$, mà hạn chế của nó trên phần bù của tiết diện không của $\bb{V}(\sh{O}_X(1))$ là một đẳng cấu từ phần bù này lên nón affine bỏ đỉnh $E$.
\end{corollary}

\begin{proof}
Nếu đặt $\sh{L}=\sh{O}_X(1)$, thì $\sh{S}_X^\geq$ đồng nhất với $\bb{S}_{\sh{O}_X}(\sh{L})$, do đó $\widehat{E}_X$ được đồng nhất chính tắc với $\bb{V}(\sh{L}^{-1})$, theo \sref{II.8.6.3}, và $C_X$ với $\bb{V}(\sh{L})$.
Cấu xạ $\bb{V}(\sh{L})\to C$ là hạn chế của $\Proj(\widehat{\alpha})$, và các khẳng định của hệ quả khi đó là các trường hợp riêng của \sref{II.8.6.2}.
\end{proof}

Ta lưu ý rằng ảnh ngược qua cấu xạ $\bb{V}(\sh{O}_X(1))\to C$ của không gian nền của tiền lược đồ các đỉnh của $C$ là không gian nền của tiết diện không của $\bb{V}(\sh{O}_X(1))$ \sref{II.8.5.5};
nhưng nói chung, các tiền lược đồ con tương ứng của $C$ và của $\bb{V}(\sh{O}_X(1))$ không đẳng cấu.
Vấn đề này sẽ được nghiên cứu dưới đây.
\subsection{Phép nổ các nón chiếu}
\label{subsection:II.8.7}

\begin{env}[8.7.1]
\label{II.8.7.1}
Dưới các điều kiện của \sref{II.8.6.1}, khi đặt $r=\Proj(\widehat{\alpha})$, ta có một biểu đồ giao hoán
\[
\label{II.8.7.1.1}
  \xymatrix{
    X
      \ar[r]^{i_X\circ\varepsilon_X}
      \ar[d]_{q}
  & \widehat{C}_X
      \ar[d]^{r}
  \\Y
      \ar[r]_{i\circ\varepsilon}
  & \widehat{C}
  }
\tag{8.7.1.1}
\]
\oldpage[II]{174}
theo \eref{eq:2.8.5.1.3} và \eref{eq:2.8.5.3.2};
hơn nữa, hạn chế của $r$ trên phần bù $\widehat{C}_X\setmin i_X(\varepsilon_X(X))$ của tiết diện không là một \emph{đẳng cấu} lên phần bù $\widehat{C}\setmin i(\varepsilon(Y))$ của tiết diện không, theo \sref{II.8.6.2}.
Nếu để đơn giản, ta giả sử $Y$ affine, $\sh{S}$ thuộc kiểu hữu hạn và được sinh bởi $\sh{S}_1$, thì $X$ xạ ảnh trên $Y$ và $\widehat{C}_X$ xạ ảnh trên $X$ \sref{II.5.5.1}; do đó $\widehat{C}_X$ xạ ảnh trên $Y$ \sref{II.5.5.5}[ii], và \emph{a fortiori} trên $\widehat{C}$ \sref{II.5.5.5}[v].
Khi đó ta có một $Y$-cấu xạ xạ ảnh $r:\widehat{C}_X\to\widehat{C}$ (mà hạn chế của nó trên $C_X$ là một $Y$-cấu xạ xạ ảnh $C_X\to C$), \emph{co $X$ vào $Y$} và cảm sinh một \emph{đẳng cấu} khi được hạn chế trên \emph{các phần bù của $X$ và của $Y$}.
Do đó, giữa $C_X$ và $C$ có một mối liên hệ tương tự mối liên hệ giữa một tiền lược đồ thu được bằng phép nổ và tiền lược đồ ban đầu \sref{II.8.1.3}.
Thật vậy, ta sẽ chứng minh rằng $C_X$ có thể được đồng nhất với phổ thuần nhất của một $\sh{O}_C$-đại số phân bậc.
\end{env}
\begin{env}[8.7.2]
\label{II.8.7.2}
Giữ nguyên các ký hiệu của \sref{II.8.6.1}, với mọi $n\geq0$, xét iđêan quasi-coherent
\[
\label{II.8.7.2.1}
  \sh{S}_{[n]} = \bigoplus_{m\geq n}\sh{S}_m
\tag{8.7.2.1}
\]
của $\sh{O}_Y$-đại số phân bậc $\sh{S}$.
Rõ ràng là
\[
\label{II.8.7.2.2}
  \sh{S}_{[0]} = \sh{S},
  \qquad
  \sh{S}_{[n]} \subset \sh{S}_{[m]}
  \qquad\qquad
  \mbox{với $m\leq n$}
\tag{8.7.2.2}
\]
\[
\label{II.8.7.2.3}
  \sh{S}_{[m]} \sh{S}_{[n]} \subset \sh{S}_{[m+n]}.
\tag{8.7.2.3}
\]

Xét $\sh{O}_C$-môđun liên kết với $\sh{S}_{[n]}$, là một iđêan quasi-coherent của $\sh{O}_C = \widetilde{\sh{S}}$ \sref{II.1.4.4}
\[
\label{II.8.7.2.4}
  \sh{I}_n = (\sh{S}_{[n]})\supertilde.
\tag{8.7.2.4}
\]

Từ \sref{II.8.7.2.2} và \sref{II.8.7.2.3}, bằng cách dùng \sref{II.1.4.4} và \sref{II.1.4.8.1}, ta suy ra các công thức tương tự
\[
\label{II.8.7.2.5}
  \sh{I}_{0} = \sh{O}_C,
  \qquad
  \sh{I}_{n} \subset \sh{I}_{m}
  \qquad
  \text{với }m\leq n
\tag{8.7.2.5}
\]
\[
\label{II.8.7.2.6}
  \sh{I}_{m} \sh{I}_{n} \subset \sh{I}_{m+n}.
\tag{8.7.2.6}
\]

Như vậy ta ở trong bối cảnh của \sref{II.8.1.1}, điều này dẫn ta đến việc đưa vào $\sh{O}_C$-đại số phân bậc quasi-coherent
\[
\label{II.8.7.2.7}
  \sh{S}^\natural
  =
  \bigoplus_{n\geq0}\sh{I}_n
  =
  \left(
    \bigoplus_{n\geq0} \sh{S}_{[n]}
  \right)\supertilde.
\tag{8.7.2.7}
\]
\end{env}
\begin{proposition}[8.7.3]
\label{II.8.7.3}
Tồn tại một $C$-đẳng cấu chính tắc
\[
\label{II.8.7.3.1}
  h: C_X \simto \Proj(\sh{S}^\natural).
\tag{8.7.3.1}
\]
\end{proposition}

\begin{proof}
Trước hết, giả sử $Y=\Spec(A)$ là affine, do đó $\sh{S}=\widetilde{S}$, với $S$ là một $A$-đại số phân bậc theo các bậc dương, và $C=\Spec(S)$.
Định nghĩa~\sref{II.8.2.6.2} khi đó cho thấy, với các ký hiệu của \sref{II.8.2.6}, rằng $\sh{S}^\natural=(S^\natural)\supertilde$.
Để định nghĩa \sref{II.8.7.3.1}, xét một phần tử thuần nhất $f\in S_d$ ($d>0$) và phần tử tương ứng $f^\natural\in S^\natural$ \sref{II.8.2.6};
$S$-đẳng cấu trong \sref{II.8.2.7.3} khi đó xác định một $C$-đẳng cấu
\[
\label{II.8.7.3.2}
  \Spec(S_f^\geq) \simto \Spec(S_{(f^\natural)}^\natural).
\tag{8.7.3.2}
\]

\oldpage[II]{175}
Nhưng, với các ký hiệu của \sref{II.8.6.2}, nếu $v:C_X\to X$ là cấu xạ cấu trúc, thì từ \sref{II.8.6.2.1} suy ra $v^{-1}(X_f)=\Spec(S_f^\geq)$.
Ta cũng có $\Spec(S_{(f^\natural)}^\natural)=D_+(f^\natural)$, nghĩa là \sref{II.8.7.3.2} xác định một đẳng cấu $v^{-1}(X_f)\to D_+(f^\natural)$.
Hơn nữa, nếu $g\in S_e$ ($e>0$), thì biểu đồ
\[
  \xymatrix{
    v^{-1}(X_{fg})
      \ar[r]^{\sim}
      \ar[d]
  & D_+(f^\natural g^\natural)
      \ar[d]
  \\v^{-1}(X_f)
      \ar[r]^{\sim}
  & D_+(f^\natural)
  }
\]
giao hoán, theo định nghĩa của đẳng cấu trong \sref{II.8.2.7.3}.
Cuối cùng, theo định nghĩa, $S_+$ được sinh bởi các phần tử thuần nhất $f$, nên từ \sref{II.8.2.10}[iv] và \sref{II.2.3.14} suy ra rằng các $D_+(f^\natural)$ tạo thành một phủ của $\Proj(S^\natural)$, và các $v^{-1}(X_f)$ tạo thành một phủ của $C_X$, vì các $X_f$ tạo thành một phủ của $X$;
như vậy, trong trường hợp này ta đã xác định đẳng cấu \sref{II.8.7.3.1}.

Để chứng minh \sref{II.8.7.3} trong trường hợp tổng quát, chỉ cần chứng minh rằng nếu $U$ và $U'$ là các tập mở affine của $Y$, tương ứng với các vành $A$ và $A'$, và sao cho $U'\subset U$, thì, khi đặt $\sh{S}|U=\widetilde{S}$ và $\sh{S}|U'=\widetilde{S'}$, biểu đồ
\[
\label{II.8.7.3.3}
  \xymatrix{
    C_{U'}
      \ar[r]
      \ar[d]
  & \Proj(S^{'\natural})
      \ar[d]
  \\C_U
      \ar[r]
  & \Proj(S^\natural)
  }
\tag{8.7.3.3}
\]
giao hoán.
Nhưng $S'$ được đồng nhất chính tắc với $S\otimes_A A'$, do đó $S^{'\natural}$ được đồng nhất chính tắc với
\[
  S^\natural \otimes_S S' = S^\natural \otimes_A A';
\]
vì vậy, $\Proj(S^{'\natural}) = \Proj(S^\natural)\times_U U'$ \sref{II.2.8.10};
tương tự, nếu $X=\Proj(S)$ và $X'=\Proj(S')$, thì $X'=X\times_U U'$ và $\sh{S}_{X'}=\sh{S}_X\otimes_{\sh{O}_U}\sh{O}_{U'}$ \sref{II.3.5.4}, hay tương đương, $\sh{S}_{X'}=j^*(\sh{S}_X)$, trong đó $j$ là phép chiếu $X'\to X$.
Khi đó, theo \sref{II.1.5.2}, ta có $C_{U'} = C_U\times_X X' = C_U\times_U U'$, và tính giao hoán của \sref{II.8.7.3.3} là hiển nhiên.
\end{proof}
\begin{remark}[8.7.4]
\label{II.8.7.4}
\begin{enumerate}
  \item[\rm{(i)}] Phần cuối của chứng minh \sref{II.8.7.3} có thể được tổng quát hóa ngay như sau.
    Cho $g:Y'\to Y$ là một cấu xạ, $\sh{S}'=g^*(\sh{S})$, và $X'=\Proj(\sh{S}')$;
    khi đó ta có biểu đồ giao hoán
    \[
    \label{II.8.7.4.1}
      \xymatrix{
        C_{X'}
          \ar[r]
          \ar[d]
      & \Proj(\sh{S}^{'\natural})
          \ar[d]
      \\C_X
          \ar[r]
      & \Proj(\sh{S}^\natural)
      }
    \tag{8.7.4.1}
    \]

    Bây giờ, cho $\vphi:\sh{S}''\to\sh{S}$ là một đồng cấu của các $\sh{O}_Y$-đại số phân bậc sao cho, nếu đặt $X''=\Proj(\sh{S}'')$, thì $u=\Proj(\vphi):X\to X''$ được xác định khắp nơi;
    ta cũng có
\oldpage[II]{176}
    một $Y$-cấu xạ $v:C\to C''$ (với $C''=\Spec(\sh{S}'')$) sao cho $\sh{A}(v)=\vphi$, và, vì $\vphi$ là một đồng cấu của các đại số phân bậc, $\vphi$ cảm sinh một $v$-cấu xạ của các đại số phân bậc $\psi:\sh{S}^{\prime\prime\natural}\to\sh{S}^\natural$ \sref{II.1.4.1}.
    Hơn nữa, từ \sref{II.8.2.10}[iv] và giả thiết về $\vphi$ suy ra rằng $\Proj(\psi)$ được xác định khắp nơi.
    Cuối cùng, xét đến \sref{II.3.5.6.1}, tồn tại một $u$-cấu xạ chính tắc $\sh{S}_{X''}\to\sh{S}_X$, từ đó theo \sref{II.1.5.6} có một cấu xạ $w:C_X\to C_{X''}$.
    Với các nhận xét đó, biểu đồ
    \[
    \label{II.8.7.4.2}
      \xymatrix{
        C_X
          \ar[r]^-{\sim}
          \ar[d]_{w}
      & \Proj(\sh{S}^\natural)
          \ar[d]^{\Proj(\psi)}
      \\C_{X''}
          \ar[r]^-{\sim}
      & \Proj(\sh{S}^{\prime\prime\natural})
      }
    \tag{8.7.4.2}
    \]
    giao hoán, như ta kiểm tra ngay được bằng cách quy về trường hợp $Y$ là affine.
  \item[\rm{(ii)}] Ta lưu ý rằng, theo \sref{II.8.7.2.5} và \sref{II.8.7.2.6}, ta có $\sh{I}_1^m\subset\sh{I}_m\subset\sh{I}_1$ với mọi $m>0$.
    Nhưng, theo định nghĩa, $\sh{I}_1=(\sh{S}_+)\supertilde$, nên $\sh{I}_1$ xác định tiền lược đồ con đóng $\varepsilon(Y)$ trong $C$ (\sref{II.1.4.10} và \sref{II.8.3.2});
    do đó ta kết luận rằng, với mọi $m>0$, \emph{giá của $\sh{O}_C/\sh{I}_m$ được chứa trong không gian nền của tiền lược đồ các đỉnh $\varepsilon(Y)$};
    trên ảnh ngược của nón affine bỏ đỉnh $E$, cấu xạ cấu trúc $\Proj(\sh{S}^\natural)\to C$ vì vậy hạn chế thành một \emph{đẳng cấu} (theo \sref{II.8.7.3} và \sref{II.8.7.1}).
    Hơn nữa, bằng cách đồng nhất chính tắc $C$ với một tập mở của $\widehat{C}$ \sref{II.8.3.3}, ta có thể mở rộng một cách hiển nhiên các iđêan $\sh{I}_m$ của $\sh{O}_C$ thành các iđêan $\sh{J}_m$ của $\sh{O}_{\widehat{C}}$, bằng cách yêu cầu chúng trùng với $\sh{O}_{\widehat{C}}$ trên tập mở $\widehat{E}$ của $\widehat{C}$.
    Nếu đặt $\sh{T}=\bigoplus_{n\geq 0}\sh{J}_n$, vốn là một $\sh{O}_{\widehat{C}}$-đại số phân bậc quasi-coherent, thì ta có thể mở rộng đẳng cấu \sref{II.8.7.3.1} thành một $\widehat{C}$-đẳng cấu
    \[
    \label{II.8.7.4.3}
      \widehat{C}_X \simto \Proj(\sh{T}).
    \tag{8.7.4.3}
    \]

    Thật vậy, trên $\widehat{E}$, từ những điều trên suy ra rằng $\Proj(\sh{T})$ được đồng nhất chính tắc với $\widehat{E}$, và do đó ta xác định đẳng cấu \sref{II.8.7.4.3} trên $\widehat{E}$ bằng cách yêu cầu nó trùng với đẳng cấu chính tắc $\widehat{E}_X\to\widehat{E}$ \sref{II.8.6.2};
    hiển nhiên đẳng cấu này và \sref{II.8.7.3.1} khi đó trùng nhau trên $E$.
\end{enumerate}
\end{remark}
\begin{corollary}[8.7.5]
\label{II.8.7.5}
Giả sử tồn tại một số $n_0>0$ sao cho
\[
\label{II.8.7.5.1}
  \sh{S}_{n+1} = \sh{S}_1\sh{S}_n
  \qquad
  \text{với }n\geq n_0.
\tag{8.7.5.1}
\]

Khi đó, tiền lược đồ các đỉnh của $C_X$ (đẳng cấu với $X$) là ảnh ngược, qua cấu xạ chính tắc $r:C_X\to C$, của tiền lược đồ các đỉnh của $C$ (đẳng cấu với $Y$).
Ngược lại, nếu tính chất này đúng, và nếu hơn nữa ta giả sử $Y$ là Noether và $\sh{S}$ thuộc kiểu hữu hạn, thì tồn tại một số $n_0>0$ sao cho \sref{II.8.7.5.1} đúng.
\end{corollary}

\begin{proof}
Vì khẳng định thứ nhất là địa phương trên $Y$, ta có thể giả sử $Y=\Spec(A)$ là affine, do đó $\sh{S}=\widetilde{S}$, với $S$ là một $A$-đại số phân bậc theo các bậc dương.
Khẳng định khi đó suy ra từ \sref{II.8.2.12}, vì $\Proj(S^\natural\otimes_S S_0) = C_X\times_C\varepsilon(Y)$ (theo phép đồng nhất trong \sref{II.8.7.3.1}), hay nói cách khác, vì tiền lược đồ này là ảnh ngược của $\varepsilon(Y)$ trong $C_X$ \sref[I]{I.4.4.1}.
Mệnh đề đảo cũng suy ra từ \sref{II.8.2.12} khi $Y$ là affine Noether và $S$ thuộc kiểu hữu hạn.
\oldpage[II]{177}
Nếu $Y$ là Noether (nhưng không nhất thiết affine) và $\sh{S}$ thuộc kiểu hữu hạn, thì tồn tại một phủ hữu hạn của $Y$ bởi các tập mở affine Noether $U_i$, và từ những điều trên ta suy ra rằng, với mọi $i$, tồn tại một số nguyên $n_i$ sao cho $\sh{S}_{n+1}|U_i = (\sh{S}_1|U_i)(\sh{S}_n|U_i)$ với $n\geq n_i$;
số lớn nhất trong các $n_i$ khi đó bảo đảm rằng \sref{II.8.7.5.1} đúng.
\end{proof}
\begin{env}[8.7.6]
\label{II.8.7.6}
Bây giờ xét $C$-tiền lược đồ $Z$ thu được bằng \emph{phép nổ} \emph{tiền lược đồ các đỉnh $\varepsilon(Y)$} trong nón affine $C$;
theo định nghĩa~\sref{II.8.1.3}, đó chính là tiền lược đồ $\Proj(\bigoplus_{n\geq0}\sh{S}_+^n)$;
đơn ánh chính tắc
\[
\label{II.8.7.6.1}
  \iota: \bigoplus_{n\geq0} \sh{S}_+^n \to \sh{S}^\natural
\tag{8.7.6.1}
\]
xác định (nhờ phép đồng nhất trong \sref{II.8.7.3}) một $C$-cấu xạ trội chính tắc
\[
\label{II.8.7.6.2}
  G(\iota)\to Z
  \tag{8.7.6.2}
\]
trong đó $G(\iota)$ là một tập mở của $C_X$ \sref{II.3.5.1};
lưu ý rằng có thể xảy ra $G(\iota)\neq C_X$, như ví dụ trong đó $Y=\Spec(K)$, với $K$ là một trường, và $\sh{S}=\widetilde{S}$, với $S=K[\bb{y}]$, trong đó $\bb{y}$ là một biến bất định \emph{có bậc 2}, cho thấy;
nếu $R_n$ ký hiệu tập $(S_+)^n$, được xem như một tập con của $S_{[n]}=S_n^\natural$, thì $S_+^\natural$ không phải là căn trong $S^\natural$ của iđêan được sinh bởi hợp của các $R_n$ (xem \sref{II.2.3.14}).
\end{env}
\begin{corollary}[8.7.7]
\label{II.8.7.7}
Giả sử tồn tại một số $n_0>0$ sao cho
\[
\label{II.8.7.7.1}
  \sh{S}_n=\sh{S}_1^n
  \quad
  \text{với }n\geq n_0.
  \tag{8.7.7.1}
\]

Khi đó cấu xạ chính tắc \sref{II.8.7.6.2} được xác định khắp nơi và là một đẳng cấu $C_X\simto Z$.
Ngược lại, nếu tính chất này đúng, và nếu hơn nữa ta giả sử $Y$ là Noether và $\sh{S}$ thuộc kiểu hữu hạn, thì tồn tại một số $n_0$ sao cho \sref{II.8.7.7.1} đúng.
\end{corollary}

\begin{proof}
Khẳng định thứ nhất là địa phương trên $Y$, nên suy ra từ \sref{II.8.2.14};
mệnh đề đảo cũng suy ra tương tự, bằng cách lập luận như trong \sref{II.8.7.5}.
\end{proof}
\begin{remark}[8.7.8]
\label{II.8.7.8}
Vì điều kiện \sref{II.8.7.7.1} kéo theo \sref{II.8.7.5.1}, ta thấy rằng, khi điều kiện ấy được thỏa mãn, không những $C_X$ được đồng nhất với tiền lược đồ thu được bằng phép nổ đỉnh (được đồng nhất với $Y$) của nón affine $C$, mà đỉnh (được đồng nhất với $X$) của $C_X$ cũng được đồng nhất với tiền lược đồ con đóng là ảnh ngược của đỉnh $Y$ của $C$.
Hơn nữa, giả thiết~\sref{II.8.7.7.1} kéo theo rằng, trên $X=\Proj(\sh{S})$, các $\sh{O}_X$-môđun $\sh{O}_X(n)$ là khả nghịch (\sref{II.3.2.5} và \sref{II.3.2.9}), và $\sh{O}_X(n)=\sh{L}^{\otimes n}$, với $\sh{L}=\sh{O}_X(1)$ (\sref{II.3.2.7} và \sref{II.3.2.9});
theo Định nghĩa~\sref{II.8.6.1.1}, do đó $C_X$ là \emph{phân thớ vectơ} $\bb{V}(\sh{L})$ trên $X$, và đỉnh của nó là \emph{tiết diện không} của phân thớ vectơ này.
\end{remark}
\subsection{Các bó ample và các phép co}
\label{subsection:II.8.8}

\begin{env}[8.8.1]
\label{II.8.8.1}
Cho $Y$ là một tiền lược đồ, $f:X\to Y$ là một cấu xạ \emph{tách} và \emph{quasi-compact}, và $\sh{L}$ là một $\sh{O}_X$-môđun khả nghịch \emph{ample đối với $f$}.
Xét $\sh{O}_Y$-đại số phân bậc theo các bậc dương
\[
\label{II.8.8.1.1}
  \sh{S} = \sh{O}_Y \oplus \bigoplus_{n\geq1}f_*(\sh{L}^{\otimes n})
\tag{8.8.1.1}
\]
\oldpage[II]{178}
vốn là quasi-coherent \sref[I]{I.9.2.2}[a].
Có một đồng cấu chính tắc của các $\sh{O}_X$-đại số phân bậc
\[
\label{II.8.8.1.2}
  \tau: f^*(\sh{S}) \to \bigoplus_{n\geq0}\sh{L}^{\otimes n}
\tag{8.8.1.2}
\]
mà, trên các bậc $\geq1$, trùng với đồng cấu chính tắc $\sigma:f^*(f_*(\sh{L}^{\otimes n})) \to \sh{L}^{\otimes n}$ \sref[0]{0.4.4.3}, và là đồng nhất ở bậc 0.
Khi đó, giả thiết $\sh{L}$ là $f$-ample kéo theo (\sref{II.4.6.3} và \sref{II.3.6.1}) rằng $Y$-cấu xạ tương ứng
\[
\label{II.8.8.1.3}
  r = r_{\sh{L},\tau} : X \to P = \Proj(\sh{S})
\tag{8.8.1.3}
\]
được xác định khắp nơi và là một \emph{phép nhúng mở trội}, đồng thời
\[
\label{II.8.8.1.4}
  r^*(\sh{O}_P(n)) = \sh{L}^{\otimes n}
  \qquad\qquad
  \mbox{với mọi $n\in\bb{Z}$.}
\tag{8.8.1.4}
\]
\end{env}
\begin{proposition}[8.8.2]
\label{II.8.8.2}
Cho $C=\Spec(\sh{S})$ là nón affine được xác định bởi $\sh{S}$;
nếu $\sh{L}$ là $f$-ample, thì tồn tại một $Y$-cấu xạ chính tắc
\[
\label{II.8.8.2.1}
  g : V = \bb{V}(\sh{L}) \to C
\tag{8.8.2.1}
\]
sao cho sơ đồ
\[
\label{II.8.8.2.2}
  \xymatrix{
    X
      \ar[r]^{j}
      \ar[d]_{f}
  & \bb{V}(\sh{L})
      \ar[r]^{\pi}
      \ar[d]^{g}
  & X
      \ar[d]^{f}
  \\Y
      \ar[r]_{\varepsilon}
  & C
      \ar[r]_{\psi}
  & Y
  }
  \tag{8.8.2.2}
\]
giao hoán, trong đó $\psi$ và $\pi$ là các cấu xạ cấu trúc, còn $j$ và $\varepsilon$ là các phép nhúng chính tắc lần lượt đưa $X$ và $Y$ lên tiết diện không của $\bb{V}(\sh{L})$ và tiền lược đồ các đỉnh của $C$.
Hơn nữa, hạn chế của $g$ trên $\bb{V}(\sh{L})\setmin j(X)$ là một phép nhúng mở
\[
\label{II.8.8.2.3}
  \bb{V}(\sh{L})\setmin j(X)\to E=C\setmin\varepsilon(Y)
  \tag{8.8.2.3}
\]
vào nón affine bỏ đỉnh $E$ ứng với $\sh{S}$.
\end{proposition}

\begin{proof}
Với các ký hiệu của \sref{II.8.8.1}, đặt $\sh{S}_P^\geq = \bigoplus_{n\geq0}\sh{O}_P(n)$ và $C_P=\Spec(\sh{S}_P^\geq)$.
Theo \sref{II.8.6.2}, có một cấu xạ chính tắc $h=\Spec(\alpha):C_P\to C$ sao cho sơ đồ
\[
\label{II.8.8.2.4}
  \xymatrix{
    C_P
      \ar[r]
      \ar[d]_{h}
  & P
      \ar[d]^{p}
  \\C
      \ar[r]_{\psi}
  & Y
  }
\tag{8.8.2.4}
\]
giao hoán; hơn nữa, nếu $\varepsilon_P:P\to C_P$ là phép nhúng chính tắc, thì sơ đồ
\[
\label{II.8.8.2.5}
  \xymatrix{
    P
      \ar[r]^{\varepsilon_P}
      \ar[d]_{p}
  & C_P
      \ar[d]^{h}
  \\Y
      \ar[r]_{\varepsilon}
  & C
  }
\tag{8.8.2.5}
\]
giao hoán \sref{II.8.7.1.1}, và cuối cùng, hạn chế của $h$ trên nón affine bỏ đỉnh $E_P$ là một \emph{đẳng cấu} $E_P\simto E$ \sref{II.8.6.2}.
Từ \sref{II.8.8.1.4} suy ra rằng
\[
  r^*(\sh{S}_P^\geq)=\bb{S}_{\sh{O}_X}(\sh{L})
\]
\oldpage[II]{179}
và do đó có một $P$-cấu xạ chính tắc $q:\bb{V}(\sh{L})\to C_P$, với sơ đồ giao hoán
\[
\label{II.8.8.2.6}
  \xymatrix{
    \bb{V}(\sh{L})
      \ar[r]^{\pi}
      \ar[d]_{q}
  & X
      \ar[d]^{r}
  \\C_P
      \ar[r]
  & P
  }
\tag{8.8.2.6}
\]
đồng nhất $\bb{V}(\sh{L})$ với tích $C_P\times_P X$ \sref{II.1.5.2};
vì $r$ là một phép nhúng mở, nên $q$ cũng vậy \sref[I]{I.4.3.2}.
Hơn nữa, hạn chế của $q$ trên $\bb{V}(\sh{L})\setmin j(X)$ đưa tiền lược đồ này vào $E_P$, theo \sref{II.8.5.2}, và sơ đồ
\[
\label{II.8.8.2.7}
  \xymatrix{
    X
      \ar[r]^{j}
      \ar[d]_{r}
  & \bb{V}(\sh{L})
      \ar[d]^{q}
  \\P
      \ar[r]_{\varepsilon_P}
  & C_P
  }
\tag{8.8.2.7}
\]
giao hoán (vì đây là một trường hợp đặc biệt của \eref{eq:2.8.5.1.3}).
Các khẳng định của \sref{II.8.8.2} suy ra ngay từ những điều này khi lấy $g$ là cấu xạ hợp $h\circ q$.
\end{proof}
\begin{remark}[8.8.3]
\label{II.8.8.3}
Giả sử thêm rằng $Y$ là một tiền lược đồ \emph{Noether} và $f$ là một cấu xạ \emph{riêng}.
Vì khi đó $r$ là \emph{riêng} \sref{II.5.4.4}, nên đóng, và đồng thời là một phép nhúng mở trội, $r$ nhất thiết là một \emph{đẳng cấu} $X\simto P$.
Hơn nữa, ta sẽ thấy trong Chương~III \sref[III]{III.2.3.5.1} rằng $\sh{S}$ khi đó nhất thiết là một $\sh{O}_Y$-đại số \emph{hữu hạn sinh}.
Suy ra $\sh{S}^\natural$ là một $\sh{S}_0^\natural$-đại số \emph{hữu hạn sinh} (\sref{II.8.2.10}[i] và \sref{II.8.7.2.7});
vì $C_P$ là $C$-đẳng cấu với $\Proj(\sh{S}^\natural)$ \sref{II.8.7.3}, ta thấy cấu xạ $h:C_P\to C$ là \emph{xạ ảnh};
vì cấu xạ $r$ là một đẳng cấu, $q:\bb{V}(\sh{L})\to C_P$ cũng vậy, và ta kết luận rằng cấu xạ $g:\bb{V}(\sh{L})\to C$ là \emph{xạ ảnh}.
Hơn nữa, vì hạn chế của $h$ trên $E_P$ là một đẳng cấu lên $E$ và $q$ là một đẳng cấu, nên hạn chế \sref{II.8.8.2.3} của $g$ là một \emph{đẳng cấu} $\bb{V}(\sh{L})\setmin j(X)\simto E$.

Nếu giả sử thêm rằng $\sh{L}$ là \emph{rất ample} đối với $f$, ta cũng sẽ thấy trong Chương~III \sref[III]{III.2.3.5.1} rằng tồn tại một số nguyên $n_0>0$ sao cho $\sh{S}_n=\sh{S}_1^n$ với $n\geq n_0$.
Khi đó, theo \sref{II.8.7.7}, ta sẽ kết luận rằng $\bb{V}(\sh{L})$ được đồng nhất với tiền lược đồ $Z$ thu được bằng phép nổ tiền lược đồ các đỉnh (được đồng nhất với $Y$) của nón affine $C$,
và \emph{tiết diện không} của $\bb{V}(\sh{L})$ (được đồng nhất với $X$) là \emph{ảnh ngược} của tiền lược đồ các đỉnh $Y$ của $C$.

Một phần các kết quả trên thực ra có thể được chứng minh mà không cần giả thiết Noether:
\end{remark}
\begin{corollary}[8.8.4]
\label{II.8.8.4}
Cho $Y$ là một tiền lược đồ (tương ứng, một lược đồ quasi-compact), $f:X\to Y$ là một cấu xạ riêng, và $\sh{L}$ là một $\sh{O}_X$-môđun khả nghịch ample đối với $f$.
Khi đó cấu xạ trong \sref{II.8.8.2.1} là riêng (tương ứng, xạ ảnh), và hạn chế \sref{II.8.8.2.3} của nó là một đẳng cấu.
\end{corollary}

\begin{proof}
Để chứng minh $g$ là riêng, ta có thể quy về trường hợp $Y$ affine, và khi đó chỉ cần xét trường hợp $Y$ là một lược đồ quasi-compact.
Lập luận như trong \sref{II.8.8.3} trước hết cho thấy $r$ là một \emph{đẳng cấu} $X\simto P$;
do đó $q$ cũng là một đẳng cấu và, vì hạn chế của $h$ trên $E_P$ là một đẳng cấu $E_P\simto E$, ta đã thấy rằng \sref{II.8.8.2.3} là một đẳng cấu.
Chỉ còn phải chứng minh rằng $g$ là \emph{xạ ảnh}.

Vì $f$ thuộc kiểu hữu hạn theo giả thiết, ta có thể áp dụng \sref{II.3.8.5} cho đồng cấu
\oldpage[II]{180}
$\tau$ trong \sref{II.8.8.1.2}:
tồn tại một số nguyên $d>0$ và một $\sh{O}_Y$-môđun con quasi-coherent hữu hạn sinh $\sh{E}$ của $\sh{S}_d$ sao cho, nếu $\sh{S}'$ là $\sh{O}_Y$-đại số con của $\sh{S}$ được sinh bởi $\sh{E}$ và $\tau'=\tau\circ f^*(\vphi)$ (trong đó $\vphi$ là đơn ánh chính tắc $\sh{S}'\to\sh{S}$), thì $r'=r_{\sh{L},\tau'}$ là một phép nhúng
\[
  X\to P'=\Proj(\sh{S}').
\]
Hơn nữa, vì $\vphi$ là đơn ánh, $r'$ cũng là một \emph{phép nhúng trội} \sref{II.3.7.6};
lập luận như đối với $r$ do đó cho thấy $r'$ là một \emph{phép nhúng đóng toàn ánh};
vì $r'$ phân tích thành $X\xrightarrow{r}\Proj(\sh{S})\xrightarrow{\Phi}\Proj(\sh{S}')$, trong đó $\Phi=\Proj(\vphi)$, ta kết luận rằng $\Phi$ cũng là một \emph{phép nhúng đóng toàn ánh}.
Nhưng điều này kéo theo $\Phi$ là một \emph{đẳng cấu};
thật vậy, ta có thể quy về trường hợp $Y=\Spec(A)$ là affine và $\sh{S}=\widetilde{S}$, $\sh{S}'=\widetilde{S'}$, trong đó $S$ là một $A$-đại số phân bậc và $S'$ là một đại số con phân bậc của $S$.
Với mọi phần tử thuần nhất $t\in S'$, $S'_{(t)}$ khi đó là một vành con của $S_{(t)}$;
trở lại định nghĩa của $\Proj(\vphi)$ \sref{II.2.8.1}, ta thấy chỉ cần chứng minh rằng, nếu $B'$ là một vành con của một vành $B$ và cấu xạ $\Spec(B)\to\Spec(B')$ ứng với đơn ánh chính tắc $B'\to B$ là một phép nhúng đóng, thì cấu xạ này nhất thiết là một \emph{đẳng cấu};
nhưng điều này suy ra từ \sref[I]{I.4.2.3}.
Hơn nữa, $\Phi^*(\sh{O}_{P'}(n))=\sh{O}_P(n)$ (\sref{II.3.5.2}[ii] và \sref{II.3.5.4}), do đó $r^{'*}(\sh{O}_{P'}(n))$ đẳng cấu với $\sh{L}^{\otimes n}$ \sref{II.4.6.3}.
Đặt $\sh{S}''=\sh{S}^{'(d)}$, sao cho theo \sref{II.3.1.8}[i], $X$ được đồng nhất chính tắc với $P''=\Proj(\sh{S}'')$, và $\sh{L}''=\sh{L}^{\otimes d}$ được đồng nhất với $\sh{O}_{P''}(1)$ \sref{II.3.2.9}[ii].

Bây giờ, nếu $C''=\Spec(\sh{S}'')$, thì $\sh{S}_{P''}^\geq=\bigoplus_{n\geq 0}\sh{O}_{P''}(n)$ được đồng nhất với $\bigoplus_{n\geq 0}\sh{L}^{\prime\prime\otimes n}$, và do đó $C_{P''}=\Spec(\sh{S}_{P''}^\geq)$ được đồng nhất với $\bb{V}(\sh{L}'')$;
ta cũng biết từ \sref{II.8.7.3} rằng $C_{P''}$ là $C''$-đẳng cấu với $\Proj(\sh{S}^{\prime\prime\natural})$;
theo định nghĩa của $\sh{S}''$, $\sh{S}^{\prime\prime\natural}$ được sinh bởi $\sh{S}_1^{\prime\prime\natural}$, và $\sh{S}_1^{\prime\prime\natural}$ thuộc kiểu hữu hạn trên $\sh{S}_0^{\prime\prime\natural}=\sh{S}''$ (\sref{II.8.2.10}[i và iii]), do đó $\Proj(\sh{S}^{\prime\prime\natural})$ là \emph{xạ ảnh} trên $C''$ \sref{II.5.5.1}.
Xét sơ đồ
\[
\label{II.8.8.4.1}
  \xymatrix{
    \bb{V}(\sh{L})
      \ar[r]^-{g}
      \ar[d]_{u}
  & \Spec(\sh{S}) = C
      \ar[d]^{v}
  \\\bb{V}(\sh{L}'')
      \ar[r]_{g''}
  & \Spec(\sh{S}'') = C''
  }
  \tag{8.8.4.1}
\]
trong đó $g$ và $g''$ ứng, theo \sref{II.1.5.6}, với các $f$-cấu xạ chính tắc
\[
  \sh{S}\to\bigoplus_{n\geq0}\sh{L}^{\otimes n}
  \quad\mbox{và}\quad
  \sh{S}''\to\bigoplus_{n\geq0}\sh{L}^{\prime\prime\otimes n}
\]
trong \sref{II.3.3.2.3} (xem \sref{II.8.8.5} dưới đây), còn $v$ và $u$ lần lượt ứng với các cấu xạ bao hàm $\sh{S}''\to\sh{S}$ và $\bigoplus_{n\geq0}\sh{L}^{\otimes nd}\to\bigoplus_{n\geq0}\sh{L}^{\otimes n}$;
theo \sref{II.3.3.2}, sơ đồ này hiển nhiên giao hoán.
Ta vừa thấy rằng $g''$ là một cấu xạ xạ ảnh;
mặt khác, $u$ là một cấu xạ \emph{hữu hạn}.
Thật vậy, vì vấn đề có tính địa phương trên $X$, ta có thể giả sử $X$ affine có vành $A$ và $\sh{L}=\sh{O}_X$;
khi đó mọi việc quy về nhận xét rằng vành $A[T]$ là một môđun hữu hạn sinh trên vành con $A[T^d]$ của nó (với $T$ là một biến bất định).
Vì $Y$ là một lược đồ quasi-compact và $C''$ affine trên $Y$, ta biết rằng $C''$ cũng là một lược đồ quasi-compact,
\oldpage[II]{181}
do đó $g''\circ u$ là một cấu xạ xạ ảnh \sref{II.5.5.5}[ii];
do tính giao hoán của \sref{II.8.8.4.1}, $v\circ g$ cũng là xạ ảnh và, vì $v$ là affine, nên tách, cuối cùng ta kết luận rằng $g$ là xạ ảnh \sref{II.5.5.5}[v].
\end{proof}
\begin{env}[8.8.5]
\label{II.8.8.5}
Xét lại tình huống trong \sref{II.8.8.1}.
Ta sẽ thấy rằng cấu xạ $g:\bb{V}(\sh{L})\to C$ còn có một định nghĩa khác, áp dụng được cho mọi $\sh{O}_X$-môđun khả nghịch $\sh{L}$ (không nhất thiết ample).
Để làm điều này, xét $f$-cấu xạ
\[
\label{II.8.8.5.1}
  \tau^\flat:\sh{S}\to\bigoplus_{n\geq0}\sh{L}^{\otimes n}
  \tag{8.8.5.1}
\]
ứng với cấu xạ $\tau$ trong \sref{II.8.8.1.2}.
Theo \sref{II.1.5.6}, điều này cảm sinh một cấu xạ $g':V\to C$ sao cho, nếu $\pi:V\to X$ và $\psi:C\to Y$ là các cấu xạ cấu trúc, thì các sơ đồ
\[
\label{II.8.8.5.2}
  \xymatrix{
    X
      \ar[d]_{f}
  & V
      \ar[l]^{\pi}
      \ar[d]^{g'}
  \\Y
  & C
      \ar[l]_{\psi}
  }
  \qquad
  \xymatrix{
    X
      \ar[r]^{j}
      \ar[d]_{f}
  & V
      \ar[d]^{g'}
  \\Y
      \ar[r]_{\varepsilon}
  & C
  }
\tag{8.8.5.2}
\]
giao hoán (\eref{eq:2.8.5.1.2} và \eref{eq:2.8.5.1.3}).
Ta sẽ chứng minh rằng (nếu giả sử $\sh{L}$ là $f$-ample) \emph{hai cấu xạ $g$ và $g'$ đồng nhất}.

Vì vấn đề có tính địa phương trên $Y$, ta có thể giả sử $Y=\Spec(A)$ là affine và, theo \sref{II.8.8.1.3}, đồng nhất $X$ với một tập mở của $P=\Proj(S)$, trong đó $S=A\oplus\bigoplus_{n\geq1}\Gamma(X,\sh{L}^{\otimes n})$;
khi đó, từ \sref{II.8.8.1.4}, ta suy ra $\Gamma(X,\sh{O}_P(n))=\Gamma(X,\sh{L}^{\otimes n})$ với mọi $n\in\bb{Z}$.
Theo định nghĩa của $h=\Spec(\alpha)$, trong đó $\alpha$ là $p$-cấu xạ chính tắc $\widetilde{S}\to\sh{S}_P^\geq$ \sref{II.8.6.1.2}, ta cần chứng minh rằng hạn chế lên $X$ của $\alpha^\sharp:p^*(\widetilde{S})\to\sh{S}_P^\geq$ đồng nhất với $\tau$.
Theo \sref[0]{0.4.4.3}, chỉ cần chứng minh rằng nếu ta hợp thành đồng cấu chính tắc $\alpha_n:S_n\to\Gamma(P,\sh{O}_P(n))$ với đồng cấu hạn chế $\Gamma(P,\sh{O}_P(n))\to\Gamma(X,\sh{O}_P(n))=\Gamma(X,\sh{L}^{\otimes n})$, thì ta thu được ánh xạ đồng nhất với mọi $n>0$;
nhưng điều này suy ra ngay từ định nghĩa của đại số $S$ và của $\alpha_n$ \sref{II.2.6.2}.
\end{env}
\begin{proposition}[8.8.6]
\label{II.8.8.6}
Giả sử (với các ký hiệu của \sref{II.8.8.5}) rằng, nếu viết $f=(f_0,\lambda)$, thì đồng cấu $\lambda:\sh{O}_Y\to f_*(\sh{O}_X)$ là song ánh;
khi đó:
\begin{enumerate}
  \item[\rm{(i)}] nếu viết $g=(g_0,\mu)$, thì $\mu:\sh{O}_C\to g_*(\sh{O}_V)$ là một đẳng cấu; và
  \item[\rm{(ii)}] nếu $X$ là nguyên (tương ứng, địa phương nguyên và chuẩn tắc), thì $C$ là nguyên (tương ứng, chuẩn tắc).
\end{enumerate}
\end{proposition}

\begin{proof}
Thật vậy, $f$-cấu xạ $\tau^\flat$ khi đó là một \emph{đẳng cấu}
\[
  \tau^\flat:\sh{S}=\psi_*(\sh{O}_C)\to f_*(\pi_*(\sh{O}_V))=\psi_*(g_*(\sh{O}_V))
\]
và $Y$-cấu xạ $g$ có thể được xem là cấu xạ mà đối với nó đồng cấu $\sh{A}(g)$ trong \sref{II.1.1.2} bằng $\tau^\flat$.
Để thấy rằng $\mu$ là một đẳng cấu các $\sh{O}_C$-môđun, theo \sref{II.1.4.2}, chỉ cần thấy rằng $\sh{A}(\mu):\psi_*(\sh{O}_C)\to\psi_*(g_*(\sh{O}_V))$ là một đẳng cấu.
Nhưng theo Định nghĩa~\sref{II.1.1.2}, ta có $\sh{A}(\mu)=\sh{A}(g)$, từ đó suy ra kết luận của (i).

Để chứng minh (ii), ta có thể quy về trường hợp $Y$ affine, do đó $\sh{S}=\widetilde{S}$, với
\oldpage[II]{182}
$S=\bigoplus_{n\geq0}\Gamma(X,\sh{L}^{\otimes n})$;
giả thiết rằng $X$ nguyên kéo theo vành $S$ là vành nguyên \sref[I]{I.7.4.4}, và do đó $C$ cũng nguyên \sref[I]{I.5.1.4}.
Để chứng minh $C$ chuẩn tắc, ta sẽ dùng bổ đề sau:

\begin{lemma}[8.8.6.1]
\label{II.8.8.6.1}
Cho $Z$ là một tiền lược đồ nguyên chuẩn tắc.
Khi đó vành $\Gamma(Z,\sh{O}_Z)$ là vành nguyên và đóng nguyên.
\end{lemma}

\begin{proof}
Từ \sref[I]{I.8.2.1.1} suy ra rằng $\Gamma(Z,\sh{O}_Z)$ là giao, trong trường hàm hữu tỉ $R(Z)$, của các vành đóng nguyên $\sh{O}_z$ với mọi $z\in Z$.
\end{proof}

Với nhận xét đó, trước hết ta chứng minh rằng $V$ là \emph{địa phương nguyên} và \emph{chuẩn tắc};
để làm điều này, ta có thể quy về trường hợp $X=\Spec(A)$ là affine, với $A$ là vành nguyên và đóng nguyên \sref{II.6.3.8}, và $\sh{L}=\sh{O}_X$.
Vì khi đó $V=\Spec(A[T])$, và $A[T]$ là vành nguyên và đóng nguyên \cite[p.~99]{II-24}, điều này chứng minh khẳng định của ta.
Với mọi tập mở affine $U$ của $C$, $g^{-1}(U)$ là quasi-compact vì cấu xạ $g$ là quasi-compact;
vì $V$ địa phương nguyên, các thành phần liên thông của $g^{-1}(U)$ là các tiền lược đồ nguyên mở trong $g^{-1}(U)$, nên có hữu hạn thành phần; hơn nữa, vì $V$ chuẩn tắc, các tiền lược đồ này cũng chuẩn tắc \sref{II.6.3.8}.
Khi đó $\Gamma(U,\sh{O}_C)$, bằng $\Gamma(g^{-1}(U),\sh{O}_V)$ theo (i), là tích trực tiếp hữu hạn của các vành nguyên và đóng nguyên \sref{II.8.8.6.1}, điều này chứng minh rằng $C$ chuẩn tắc \sref{II.6.3.4}.
\end{proof}
\subsection{Tiêu chuẩn về tính ample của Grauert: phát biểu}
\label{subsection:II.8.9}

Ta sẽ chứng minh rằng các tính chất được chứng minh trong \sref{II.8.8.2} \emph{đặc trưng} cho các môđun ample đối với $f$ trong số các $\sh{O}_X$-môđun, và chính xác hơn, chứng minh tiêu chuẩn sau:
\begin{theorem}[8.9.1]
\label{II.8.9.1}
\emph{(Tiêu chuẩn Grauert).}
Cho $Y$ là một tiền lược đồ, $p:X\to Y$ là một cấu xạ tách và quasi-compact, và $\sh{L}$ là một $\sh{O}_X$-môđun khả nghịch.
Để $\sh{L}$ ample đối với $p$, điều kiện cần và đủ là tồn tại một $Y$-tiền lược đồ $C$, một $Y$-tiết diện $\varepsilon:Y\to C$ của $C$, và một $Y$-cấu xạ $q:\bb{V}(\sh{L})\to C$, thỏa mãn các tính chất sau:
\begin{enumerate}
  \item[\rm{(i)}] sơ đồ
    \[
      \label{II.8.9.1.1}
        \xymatrix{
          X
            \ar[r]^{j}
            \ar[d]_{p}
        & \bb{V}(\sh{L})
            \ar[d]^{q}
        \\Y
            \ar[r]_{\varepsilon}
        & C
        }
      \tag{8.9.1.1}
    \]
    giao hoán, trong đó $j$ là tiết diện không của phân thớ vectơ $\bb{V}(\sh{L})$; và
  \item[\rm{(ii)}] hạn chế của $q$ lên $\bb{V}(\sh{L})\setmin j(X)$ là một phép nhúng mở quasi-compact
    \[
      \bb{V}(\sh{L})\setmin j(X)\to C
    \]
    có ảnh không giao với $\varepsilon(Y)$.
\end{enumerate}
\end{theorem}

Lưu ý rằng, nếu $C$ \emph{tách} trên $Y$, thì trong điều kiện~(ii) ta có thể bỏ giả thiết phép nhúng mở là quasi-compact;
để thấy rằng tính chất này (tính quasi-compact) thực ra là hệ quả của các điều kiện khác, ta có thể quy về trường hợp $Y$ affine, và khi đó khẳng định suy ra từ \sref[I]{I.5.5.1}[i] và \sref[I]{I.5.5.10}.
Ta cũng có thể bỏ
\oldpage[II]{183}
giả thiết ấy nếu giả sử $X$ là Noether, vì khi đó $V$ cũng là Noether, và khẳng định suy ra từ \sref[I]{I.6.3.5}.

\begin{corollary}[8.9.2]
\label{II.8.9.2}
Nếu cấu xạ $p:X\to Y$ là riêng, thì trong phát biểu của Định lý~\sref{II.8.9.1}, ta có thể giả sử $q$ là riêng và thay “phép nhúng mở” bằng “đẳng cấu”.
\end{corollary}

Một cách gợi hình hơn, ta có thể nói (mỗi khi $p:X\to Y$ là riêng) rằng \emph{$\sh{L}$ ample đối với $p$ khi và chỉ khi ta có thể “\emph{co}” tiết diện không của phân thớ vectơ $\bb{V}(\sh{L})$ về tiền lược đồ cơ sở $Y$.}
Một trường hợp riêng quan trọng là khi $Y$ là phổ của một trường và phép “co” được thực hiện bằng cách co tiết diện không của $\bb{V}(\sh{L})$ \emph{thành một điểm duy nhất}.

\begin{env}[8.9.3]
\label{II.8.9.3}
Tính cần thiết của các điều kiện trong Định lý \sref{II.8.9.1}, cũng như bản thân Hệ quả \sref{II.8.9.2}, đều suy ra ngay từ \sref{II.8.8.2} và \sref{II.8.8.4}.

Để chứng minh các điều kiện của \sref{II.8.9.1} là đủ, ta xét một tình huống tổng quát hơn đôi chút.
Với mục đích đó, đặt (theo các ký hiệu của \sref{II.8.8.2})
\[
  \sh{S}'=\bigoplus_{n\geq 0}\sh{L}^{\otimes n}
\]
và
\[
  V=\bb{V}(\sh{L})=\Spec(\sh{S}').
\]

Tiền lược đồ con đóng $j(X)$, tức tiết diện không của $\bb{V}(\sh{L})$, được xác định bởi bó iđêan quasi-coherent $\sh{J}=(\sh{S}'_+)\supertilde$ của $\sh{O}_V$ \sref{II.1.4.10}.
$\sh{O}_V$-môđun này là \emph{khả nghịch}, vì tính chất đó là địa phương trên $X$, và điều này quy về việc nhận xét rằng iđêan $TA[T]$ trong vành đa thức $A[T]$ là một $A[T]$-môđun tự do sinh bởi một phần tử.
Hơn nữa, hiển nhiên (một lần nữa vì vấn đề là địa phương trên $X$) rằng
\[
  \sh{L}=j^*(\sh{J})
\]
và
\[
  j_*(\sh{L})=\sh{J}/\sh{J}^2.
\]

Mặt khác, nếu
\[
  \pi:\bb{V}(\sh{L})\to X
\]
là cấu xạ cấu trúc, thì $\pi_*(\sh{J})=\sh{S}'_+$ và $\pi_*(\sh{J}/\sh{J}^2)=\sh{L}$;
do đó có các đồng cấu chính tắc $\sh{L}\to\pi_*(\sh{J})\to\sh{L}$, trong đó đồng cấu thứ nhất là đơn ánh chính tắc $\sh{L}\to\sh{S}'_+$, đồng cấu thứ hai là phép chiếu chính tắc từ $\sh{S}'_+$ lên $\sh{S}'_1=\sh{L}$, và hợp thành của chúng là đồng nhất.
Ta cũng có thể nhúng chính tắc $\pi_*(\sh{J}) = \sh{S}'_+ = \bigoplus_{n\geq1}\sh{L}^{\otimes n}$ vào \emph{tích} $\prod_{n\geq1}\sh{L}^{\otimes n} = \varprojlim_n \pi_*(\sh{J}/\sh{J}^{n+1})$ (vì $\pi_*(\sh{J}/\sh{J}^{n+1}) = \sh{L}\oplus\sh{L}^{\otimes 2}\oplus\ldots\oplus\sh{L}^{\otimes n}$), và do đó ta có các đồng cấu chính tắc
\[
\label{II.8.9.3.1}
  \sh{L}\to\varprojlim\pi_*(\sh{J}/\sh{J}^{n+1})\to\sh{L}
  \tag{8.9.3.1}
\]
có hợp thành là đồng nhất.

Với nhận xét đó, tổng quát hóa của \sref{II.8.9.1} mà ta sẽ chứng minh là như sau:
\end{env}
\begin{proposition}[8.9.4]
\label{II.8.9.4}
Cho $Y$ là một tiền lược đồ, $V$ là một $Y$-tiền lược đồ, và $X$ là một tiền lược đồ con đóng của $V$ được xác định bởi một iđêan $\sh{J}$ của $\sh{O}_V$, là một $\sh{O}_V$-môđun \emph{khả nghịch};
nếu $j:X\to V$ là
\oldpage[II]{184}
đơn ánh chính tắc, ta đặt $\sh{L} = j^*(\sh{J}) = \sh{J}\otimes_{\sh{O}_V}\sh{O}_X$, khi đó $j_*(\sh{L})=\sh{J}/\sh{J}^2$.
Giả sử cấu xạ cấu trúc $p:X\to Y$ là tách và quasi-compact, và các điều kiện sau được thỏa mãn:
\begin{enumerate}
  \item[\rm{(i)}] tồn tại một $Y$-cấu xạ $\pi:V\to X$ kiểu hữu hạn sao cho $\pi\circ j=1_X$, và do đó $\pi_*(\sh{J}/\sh{J}^2)=\sh{L}$;
  \item[\rm{(ii)}] tồn tại một đồng cấu $\sh{O}_X$-môđun $\vphi:\sh{L}\to\varprojlim\pi_*(\sh{J}/\sh{J}^{n+1})$ sao cho hợp thành
    \[
      \sh{L}
      \xrightarrow{\vphi}
      \varprojlim \pi_*(\sh{J}/\sh{J}^{n+1})
      \xrightarrow{\alpha}
      \pi_*(\sh{J}/\sh{J}^2) = \sh{L}
    \]
    (trong đó $\alpha$ là đồng cấu chính tắc) là đồng nhất;
  \item[\rm{(iii)}] tồn tại một $Y$-tiền lược đồ $C$, một $Y$-tiết diện $\varepsilon$ của $C$ và một $Y$-cấu xạ $q:V\to C$ sao cho sơ đồ
    \[
    \label{II.8.9.4.1}
      \xymatrix{
        X
          \ar[r]^{j}
          \ar[d]_{p}
      & V
          \ar[d]^{q}
      \\Y
          \ar[r]_{\varepsilon}
      & C
      }
    \tag{8.9.4.1}
    \]
    giao hoán; và
  \item[\rm{(iv)}] hạn chế của $q$ lên $W=V\setmin j(X)$ là một phép nhúng mở quasi-compact vào $C$, có ảnh không giao với $\varepsilon(Y)$.
\end{enumerate}
Khi đó $\sh{L}$ ample đối với $p$.
\end{proposition}


\subsection{Tiêu chuẩn về tính ample của Grauert: chứng minh}
\label{subsection:II.8.10}

\begin{lemma}[8.10.1]
\label{II.8.10.1}
Cho $\pi:V\to X$ là một cấu xạ, $j:X\to V$ là một $X$-tiết diện của $V$ đồng thời là một phép nhúng đóng, và $\sh{J}$ là một iđêan quasi-coherent của $\sh{O}_V$ xác định tiền lược đồ con đóng của $V$ ứng với $j$.
Khi đó các khẳng định sau đều đúng.
\begin{enumerate}
  \item[\rm{(i)}] Với mọi $n\geq0$, $\pi_*(\sh{O}_V/\sh{J}^{n+1})$ và $\pi_*(\sh{J}/\sh{J}^{n+1})$ là các $\sh{O}_X$-môđun quasi-coherent, và $\pi_*(\sh{O}_V/\sh{J})=\sh{O}_X$ và $\pi_*(\sh{J}/\sh{J}^2)=j^*(\sh{J})$.
  \item[\rm{(ii)}] Nếu $X=\{\xi\}=\Spec(k)$, trong đó $k$ là một trường, thì $\varprojlim\pi_*(\sh{O}_V/\sh{J}^{n+1})$ đẳng cấu với vành hoàn thành tách của vành địa phương $\sh{O}_{j(\xi)}$ đối với tôpô $\mathfrak{m}_{j(\xi)}$-adic.
  \item[\rm{(iii)}] Giả sử $\sh{J}$ là một $\sh{O}_V$-môđun khả nghịch (điều này kéo theo rằng
    \[
      \sh{L} = j^*(\sh{J}) = \pi_*(\sh{J}/\sh{J}^2)
    \]
    là một $\sh{O}_X$-môđun khả nghịch), và tồn tại một đồng cấu $\vphi:\sh{L}\to\varprojlim\pi_*(\sh{J}/\sh{J}^{n+1})$ sao cho hợp thành $\sh{L} \xrightarrow{\vphi} \varprojlim\pi_*(\sh{J}/\sh{J}^{n+1}) \xrightarrow{\alpha} \pi_*(\sh{J}/\sh{J}^2)$ (trong đó $\alpha$ là đồng cấu chính tắc) là đồng nhất.
    Nếu đặt $\sh{S}=\bigoplus_{n\geq0}\sh{L}^{\otimes n}$, thì $\vphi$ cảm sinh chính tắc một đẳng cấu $\sh{O}_X$-đại số từ vành hoàn thành $\widehat{\sh{S}}$ của $\sh{S}$ đối với lọc chính tắc của nó (vành hoàn thành này đẳng cấu với tích $\prod_{n\geq0}\sh{L}^{\otimes n}$) lên $\varprojlim\pi_*(\sh{O}_V/\sh{J}^{n+1})$.
\end{enumerate}
\end{lemma}

\begin{proof}
Trước hết, lưu ý rằng giá của $\sh{O}_V$-môđun $\sh{O}_V/\sh{J}^{n+1}$ là $j(X)$, và giá của $\sh{J}/\sh{J}^{n+1}$ được chứa trong $j(X)$.
Trong trường hợp (ii), $j(X)$ là một điểm đóng $j(\xi)$ của $V$,
\oldpage[II]{185}
và, theo định nghĩa, $\pi_*(\sh{O}_V/\sh{J}^{n+1})$ là thớ của $\sh{O}_V/\sh{J}^{n+1}$ tại điểm $j(\xi)$, tức là, nếu đặt $C=\sh{O}_{j(\xi)}$ và ký hiệu $\mathfrak{m}$ là iđêan cực đại của $C$, thì đó là $C$-môđun $C/\mathfrak{m}^{n+1}$;
khi đó khẳng định (ii) là hiển nhiên.

Để chứng minh (i), lưu ý rằng vấn đề mang tính địa phương trên $X$;
do đó ta có thể quy về trường hợp $X$ là affine.
Cho $U$ là một tập mở affine của $V$;
khi đó $j(X)\cap U$ là một tập mở affine của $j(X)$, nên $U_0=\pi(j(X)\cap U)$ là một tập mở affine của $X$ đẳng cấu với $j(X)\cap U$;
với mọi tập mở affine $W_0\subset U_0$ của $X$, $W=\pi^{-1}(W_0)\cap U$ là một tập mở affine của $V$, vì $X$ là một lược đồ \sref[I]{I.5.5.10};
đặc biệt, $U'=U\cap\pi^{-1}(U_0)$ là một tập mở affine của $V$, và hiển nhiên $\pi(U')=U_0$ và $j(U_0)=j(X)\cap U$.
Theo định nghĩa,
\[
  \Gamma(W_0,\pi_*(\sh{O}_V/\sh{J}^{n+1}))
  =\Gamma(\pi^{-1}(W_0),\sh{O}_V/\sh{J}^{n+1});
\]
nhưng vì mỗi điểm của $\pi^{-1}(W_0)$ không thuộc $j(W_0)$ đều có một lân cận mở trong $\pi^{-1}(W_0)$ không giao với $j(X)$, mà trên đó $\sh{O}_V/\sh{J}^{n+1}$ bằng không, nên hiển nhiên các tiết diện của $\sh{O}_V/\sh{J}^{n+1}$ trên $\pi^{-1}(W_0)$ và trên $W$ tương ứng song ánh.
Nói cách khác, nếu $\pi'$ là hạn chế của $\pi$ lên $U'$, thì các $(\sh{O}_X|U_0)$-môđun
\[
  \pi_*(\sh{O}_V/\sh{J}^{n+1})|U_0
  \quad\hbox{và}\quad
  \pi'_*\bigl((\sh{O}_V/\sh{J}^{n+1})|U'\bigr)
\]
trùng nhau.
Vì $U'$ và $U_0$ là affine và các $U_0$ phủ $X$, từ đó ta kết luận \sref[I]{I.1.6.3} rằng $\pi_*(\sh{O}_V/\sh{J}^{n+1})$ là quasi-coherent, và chứng minh đối với $\pi_*(\sh{J}/\sh{J}^{n+1})$ hoàn toàn tương tự.

Cuối cùng, để chứng minh (iii), lưu ý rằng $\sh{S}$ chính là đại số đối xứng $\bb{S}_{\sh{O}_X}(\sh{L})$;
do đó $\vphi$ cảm sinh chính tắc một đồng cấu $\sh{O}_X$-đại số
\[
  \psi:\sh{S}\longrightarrow
  \varprojlim_n\pi_*(\sh{O}_V/\sh{J}^{n+1})
\]
\sref{II.1.7.4};
hơn nữa, đồng cấu này gửi $\sh{L}^{\otimes n}$ vào $\varprojlim_m\pi_*(\sh{J}^n/\sh{J}^{m+1})$, nên liên tục đối với các tôpô đang xét và vì thế mở rộng thành một đồng cấu
\[
  \widehat{\psi}:\widehat{\sh{S}}\longrightarrow
  \varprojlim_n\pi_*(\sh{O}_V/\sh{J}^{n+1}).
\]
Để thấy đây thực sự là một đẳng cấu, như trong chứng minh (i), ta có thể quy về trường hợp $X=\Spec(A)$ và $V=\Spec(B)$ là affine, với $\sh{J}=\widetilde{\mathfrak{J}}$, trong đó $\mathfrak{J}$ là một iđêan của $B$;
ứng với $\pi$ là một đơn ánh $A\to B$ đồng nhất $A$ với một vành con của $B$ bù trực tiếp cho $\mathfrak{J}$, và $\sh{L}$ (tương ứng, $\pi_*(\sh{O}_V/\sh{J}^{n+1})$) là $\sh{O}_X$-môđun quasi-coherent liên kết với $A$-môđun $L=\mathfrak{J}/\mathfrak{J}^2$ (tương ứng, $B/\mathfrak{J}^{n+1}$).
Vì $\sh{J}$ là một $\sh{O}_V$-môđun \emph{khả nghịch}, ta còn có thể giả sử $\mathfrak{J}=Bt$, trong đó $t$ không phải là một ước của không trong $B$.
Từ hệ thức $B=A\oplus Bt$, với mọi $n>0$ ta suy ra
\[
  B=A\oplus At\oplus At^2\oplus\cdots\oplus At^n\oplus Bt^{n+1}
\]
và do đó tồn tại một $A$-đẳng cấu chính tắc từ vành các chuỗi lũy thừa hình thức $A[[T]]$ lên $C=\varprojlim B/\mathfrak{J}^{n+1}$ gửi $T$ tới $t$.
Mặt khác, $L=A\overline{t}$, trong đó $\overline{t}$ là lớp của $t$ modulo $Bt^2$, và theo giả thiết đồng cấu $\vphi$ gửi $\overline{t}$ tới một phần tử $t'\in C$ đồng dư với $t$ modulo $Ct^2$.
Bằng quy nạp theo $n$, từ đó ta suy ra
\[
  A\oplus At'\oplus\cdots\oplus A(t')^n\oplus Ct^{n+1}
  =A\oplus At\oplus\cdots\oplus At^n\oplus Ct^{n+1},
\]
điều này chứng minh rằng đồng cấu $\widehat{\psi}$ thực sự tương ứng với một đẳng cấu từ $\prod_{n\geq0}L^{\otimes n}$ lên $C$.
\end{proof}

\begin{lemma}[8.10.2]
\label{II.8.10.2}
Dưới các giả thiết của Bổ đề~\sref{II.8.10.1}, cho $g:X'\to X$ là một cấu xạ,
\oldpage[II]{186}
đặt $V'=V\times_X X'$, và cho $\pi':V'\to X'$ và $g':V'\to V$ là các phép chiếu chính tắc, sao cho ta có sơ đồ giao hoán
\[
  \xymatrix{
    V
      \ar[d]_{\pi}
  & V'
      \ar[l]^{g'}
      \ar[d]^{\pi'}
  \\X
  & X'
      \ar[l]_{g}
  }
\]

Khi đó $j'=j\times1_{X'}$ là một $X'$-tiết diện của $V'$, đồng thời là một phép nhúng đóng, và $\sh{J}'=g^{'*}(\sh{J})\sh{O}_{V'}$ là iđêan quasi-coherent của $\sh{O}_{V'}$ xác định tiền lược đồ con đóng của $V'$ ứng với $j'$.
Hơn nữa,
\[
  \pi'_*(\sh{O}_{V'}/\sh{J}'^{\,n+1})
  =g^*(\pi_*(\sh{O}_V/\sh{J}^{n+1})).
\]
Cuối cùng, $\sh{J}'$ là một $\sh{O}_{V'}$-môđun đẳng cấu chính tắc với $g^{'*}(\sh{J})$ và đặc biệt là khả nghịch nếu $\sh{J}$ là một $\sh{O}_V$-môđun khả nghịch.
\end{lemma}

\begin{proof}
Việc $j'$ là một phép nhúng đóng suy ra từ \sref[I]{I.4.3.1}, và $j'$ là một $X'$-tiết diện của $V'$ do tính hàm tử của phép đổi cơ sở đối với các tiền lược đồ.
Hơn nữa, nếu $Z$ (tương ứng, $Z'$) là tiền lược đồ con đóng của $V$ (tương ứng, $V'$) ứng với $j$ (tương ứng, $j'$), thì $Z'=g^{'-1}(Z)$ \sref[I]{I.4.3.1}, và khẳng định thứ hai suy ra từ \sref[I]{I.4.4.5}.
Để chứng minh các khẳng định còn lại, như trong \sref{II.8.10.1}, ta thấy có thể quy về trường hợp $X$, $V$ và $X'$ (do đó cả $V'$) đều là affine;
ta giữ nguyên các ký hiệu của chứng minh \sref{II.8.10.1} và đặt $X'=\Spec(A')$.
Khi đó $V'=\Spec(B')$, trong đó $B'=B\otimes_A A'$, và $\sh{J}'=\widetilde{\mathfrak{J}'}$, trong đó $\mathfrak{J}'=\Im(\mathfrak{J}\otimes_A A')$.
Do đó,
\[
  B'/\mathfrak{J}'^{\,n+1}
  =(B/\mathfrak{J}^{n+1})\otimes_A A';
\]
hơn nữa, vì $\mathfrak{J}$ là một hạng tử trực tiếp của $B$ xét như $A$-môđun, nên $\mathfrak{J}\otimes_A A'$ là một hạng tử trực tiếp của $B'$ xét như $A'$-môđun và do đó được đồng nhất chính tắc với $\mathfrak{J}'$.
\end{proof}

\begin{corollary}[8.10.3]
\label{II.8.10.3}
Dưới các giả thiết của Bổ đề~\sref{II.8.10.1}, giả sử thêm rằng $\pi$ là kiểu hữu hạn và $\sh{J}$ là một $\sh{O}_V$-môđun khả nghịch.
Khi đó, với mọi $x\in X$, vành địa phương của thớ $\pi^{-1}(x)$ tại điểm $j(x)$ là một vành chính quy (do đó là miền nguyên) có chiều 1, vành hoàn thành của nó đẳng cấu với vành các chuỗi lũy thừa hình thức $\kres(x)[[T]]$ (trong đó $T$ là một biến);
hơn nữa, tồn tại đúng một thành phần bất khả quy của $\pi^{-1}(x)$ chứa $j(x)$.
\end{corollary}

\begin{proof}
Vì $\pi^{-1}(x)=V\times_X\Spec(\kres(x))$, \sref{II.8.10.2} cho phép ta quy về trường hợp $X$ là phổ của một trường $K$.
Vì $\pi$ là kiểu hữu hạn \sref[I]{I.6.3.4}[iv], khi đó $\sh{O}_{j(x)}$ là một vành Noether địa phương và do đó tách đối với tôpô $\mathfrak{m}_{j(x)}$-adic \sref[0]{0.7.3.5};
từ \sref{II.8.10.1}[ii và iii] suy ra rằng vành hoàn thành của nó đẳng cấu với $K[[T]]$, và do đó $\sh{O}_{j(x)}$ chính quy và có chiều 1 (\cite[p.~17-01, th.~1]{I-1});
cuối cùng, vì $\sh{O}_{j(x)}$ là miền nguyên, $j(x)$ thuộc đúng một trong các thành phần bất khả quy (hữu hạn) của $\pi^{-1}(x)$ \sref[I]{I.5.1.4}.
\end{proof}

\begin{corollary}[8.10.4]
\label{II.8.10.4}
Dưới các giả thiết của Bổ đề~\sref{II.8.10.1}, giả sử thêm rằng $\sh{J}$ là một $\sh{O}_V$-môđun khả nghịch.
Đặt $W=V\setmin j(X)$;
với mọi iđêan quasi-coherent $\sh{K}$ của $\sh{O}_X$, đặt $\sh{K}_V=\pi^*(\sh{K})\sh{O}_V$ và $\sh{K}_W=\sh{K}_V|W$.
Khi đó $\sh{K}_V$ là iđêan quasi-coherent lớn nhất của $\sh{O}_V$ có hạn chế trên $W$ bằng $\sh{K}_W$.
\end{corollary}

\begin{proof}
Thật vậy, như trong \sref{II.8.10.1}, ta thấy vấn đề mang tính địa phương trên $X$ và trên $V$;
do đó ta có thể giữ nguyên các ký hiệu của chứng minh \sref{II.8.10.1}, với $\mathfrak{J}=Bt$, trong đó $t$ không phải là một ước của không trong $B$.
Hơn nữa, $W=\Spec(B_t)$ và $\sh{K}=\widetilde{\mathfrak{K}}$, trong đó $\mathfrak{K}$ là một iđêan của $A$;
do đó $\pi^*(\sh{K})\sh{O}_V=(\mathfrak{K}.B)\supertilde$ \sref[I]{I.1.6.9}, $\sh{K}_W=(\mathfrak{K}.B_t)\supertilde$, và iđêan lớn nhất
\oldpage[II]{187}
của $B$ mà ảnh chính tắc trong $B_t$ là $\mathfrak{K}.B_t$ chính là ảnh ngược của $\mathfrak{K}.B_t$, nghĩa là tập hợp các $s\in B$ sao cho, với một số nguyên $n>0$, ta có $t^ns\in\mathfrak{K}.B$.
Ta phải chứng minh rằng hệ thức cuối cùng này kéo theo $s\in\mathfrak{K}.B$, hay tương đương rằng ảnh chính tắc của $t$ không phải là một ước của không trong $B/\mathfrak{K}B=(A/\mathfrak{K})\otimes_A B$, điều này suy ra từ \sref{II.8.10.2} áp dụng cho $X'=\Spec(A/\mathfrak{K})$.
\end{proof}

\begin{corollary}[8.10.5]
\label{II.8.10.5}
Giả sử các giả thiết của \sref{II.8.10.3} được thỏa mãn;
cho $W=V\setmin j(X)$, $x$ là một điểm của $X$, $\sh{K}$ là một iđêan quasi-coherent của $\sh{O}_X$, và $z$ là điểm tổng quát của thành phần bất khả quy của $\pi^{-1}(x)$ chứa $j(x)$ \sref{II.8.10.3}.
\begin{enumerate}
    \item[\rm{(i)}] Cho $g$ là một tiết diện của $\sh{O}_V$ trên $V$ sao cho $g|W$ là một tiết diện của $\sh{K}_W$ trên $W$ (theo các ký hiệu của \sref{II.8.10.4}).
        Khi đó $g$ là một tiết diện của $\sh{K}_V$;
        nếu thêm $g(z)\neq0$, và nếu, với mọi số nguyên $m>0$, ta ký hiệu bởi $g_m^x$ mầm tại điểm $x$ của ảnh chính tắc $g_m$ của $g$ trong $\Gamma(X,\pi_*(\sh{O}_V/\sh{J}^{m+1}))$, thì tồn tại một số nguyên $m>0$ sao cho ảnh của $g_m^x$ trong
        \[
            (\pi_*(\sh{O}_V/\sh{J}^{m+1}))_x \otimes_{\sh{O}_x} \kres(x)
        \]
        khác không.
    \item[\rm{(ii)}] Giả sử thêm rằng các điều kiện của \sref{II.8.10.1}[iii] được thỏa mãn.
        Khi đó, nếu tồn tại một tiết diện $g$ của $\sh{K}_V$ trên $V$ sao cho $g(z)\neq0$, thì tồn tại một số nguyên $n\geq0$ và một tiết diện $f$ của $\sh{K}.\sh{L}^{\otimes n}=\sh{K}\otimes\sh{L}^{\otimes n}\subset\sh{L}^{\otimes n}$ sao cho $f(x)\neq0$.
        Nếu $g$ là một tiết diện của $\sh{J}$, ta có thể lấy $n>0$.
\end{enumerate}
\end{corollary}

\begin{proof}
\begin{enumerate}
    \item[\rm{(i)}] Vì iđêan của $\sh{O}_W$ sinh bởi $g|W$ được chứa trong $\sh{K}_W$ theo giả thiết, nên iđêan của $\sh{O}_V$ sinh bởi $g$ được chứa trong $\sh{K}_V$ theo \sref{II.8.10.4}; nói cách khác, $g$ là một tiết diện của $\sh{K}_V$.
        Để chứng minh khẳng định thứ hai của (i), ta lại có thể giả sử $X$ và $V$ đều affine và giữ nguyên các ký hiệu của \sref{II.8.10.1};
        khi đó thớ $\pi^{-1}(x)$ là affine, với vành $B'=B\otimes_A\kres(x)$, và trong $B'$ tồn tại một phần tử $t'$ không phải là một ước của không sao cho $B'=\kres(x)\oplus B't'$.
        Vì $j(x)$ là một điểm chuyên biệt hóa của $z$ và $g(z)\neq0$, nên nhất thiết $g_{j(x)}\neq0$.
        Nhưng $\sh{O}_{j(x)}$ là một vành địa phương tách \sref{II.8.10.3}, và do đó nhúng vào vành hoàn thành của nó; vì thế ảnh của $g$ trong vành hoàn thành này khác không.
        Vành hoàn thành này lại đẳng cấu với $\varprojlim_n(B'/B't'^{\,n+1})$ \sref{II.8.10.3};
        nếu $g'=g\otimes1\in B'$, thì tồn tại một số nguyên $m$ sao cho $g'\notin B't'^{\,m+1}$, hay tương đương, ảnh $g'_m$ của $g'$ trong $B'/B't'^{\,m+1}$ khác không.
        Nhưng vì $g'_m$ chính là ảnh của $g_m^x$, khẳng định của ta được chứng minh.
    \item[\rm{(ii)}] Theo \sref{II.8.10.1}[iii], $\pi_*(\sh{O}_V/\sh{J}^{m+1})$ đẳng cấu với tổng trực tiếp của các $\sh{L}^{\otimes k}$ với $0\leq k\leq m$;
        ta ký hiệu bởi $f_k$ thành phần của phần tử trong $\bigoplus_{k=0}^m\Gamma(X,\sh{L}^{\otimes k})$ tương ứng với $g_m$ qua đẳng cấu này; $f_k$ là một tiết diện của $\sh{L}^{\otimes k}$ trên $X$.
        Chọn $m$ như trong (i), theo (i) tồn tại một chỉ số $k$ sao cho $f_k(x)\neq0$.
        Để thấy rằng $f_k$ là một tiết diện của $\sh{K}\sh{L}^{\otimes k}$, chỉ cần xét, như trên, trường hợp $X$ và $V$ đều affine, và điều này suy ra ngay từ việc $g\in\mathfrak{K}.B$ (theo các ký hiệu của \sref{II.8.10.4}).
        Khẳng định cuối cùng suy ra từ việc giả thiết $g\in\Gamma(V,\sh{J})$ kéo theo $f_0=0$.
\end{enumerate}
\end{proof}

\begin{env}[8.10.6]
\label{II.8.10.6}
\emph{Chứng minh \sref{II.8.9.4}.}
Vấn đề mang tính địa phương trên $Y$ \sref{II.4.6.4};
vì $\varepsilon$ là một $Y$-tiết diện, ta có thể thay $C$ bằng một lân cận mở affine $U$ của một điểm của $\varepsilon(Y)$ sao cho $\varepsilon(Y)\cap U$ đóng trong $U$.
Nói cách khác, ta có thể giả sử $C$ affine và $Y$ là một tiền lược đồ con đóng của $C$ (do đó cũng affine) được xác định bởi một bó quasi-coherent
\oldpage[II]{188}
$\sh{I}$ gồm các iđêan của $\sh{O}_C$.
Vì $p$ tách và quasi-compact, $X$ do đó là một lược đồ quasi-compact, và ta quy về việc chứng minh rằng $\sh{L}$ là \emph{ample} \sref{II.4.6.4}.
Theo tiêu chuẩn~\sref{II.4.5.2}[a], do đó ta phải chứng minh điều sau:
với mọi iđêan quasi-coherent $\sh{K}$ của $\sh{O}_X$ và mọi điểm $x\in X$ không thuộc giá của $\sh{O}_X/\sh{K}$, tồn tại một số nguyên $n>0$ và một tiết diện $f$ của $\sh{K}\otimes\sh{L}^{\otimes n}$ trên $X$ sao cho $f(x)\neq0$.

Để làm điều đó, đặt
\begin{align*}
  \sh{K}_V &= \pi^*(\sh{K})\sh{O}_V,\\
  \sh{K}_W &= \sh{K}_V|W,
\end{align*}
trong đó $W=V\setmin j(X)$;
vì hạn chế của $q$ lên $W$ là một phép nhúng mở quasi-compact vào $C$, từ \sref[I]{I.9.4.2} suy ra $\sh{K}_W$ là hạn chế lên $W$ của một iđêan quasi-coherent $\sh{K}'_V$ của $\sh{O}_V$ có dạng
\[
    \sh{K}'_V = q^*(\sh{K}_C)\sh{O}_V,
\]
trong đó $\sh{K}_C$ là một iđêan quasi-coherent của $\sh{O}_C$.
Hơn nữa, vì theo giả thiết $q^{-1}(Y)\subset j(X)$, và vì $Y$ được xác định bởi iđêan $\sh{I}$, hạn chế lên $W$ của $q^*(\sh{I})\sh{O}_V$ trùng với hạn chế của $\sh{O}_V$; do đó $\sh{K}_W$ cũng là hạn chế lên $W$ của $q^*(\sh{I}\sh{K}_C)\sh{O}_V$, và ta có thể giả sử rằng $\sh{K}_C\subset\sh{I}$, từ đó
\[
\label{II.8.10.6.1}
    \sh{K}'_V \subset q^*(\sh{I})\sh{O}_V \subset \sh{J}
\tag{8.10.6.1}
\]
có tính đến \sref[I]{I.4.4.6} và tính giao hoán của \sref{II.8.9.4.1}.
Hơn nữa, từ \sref{II.8.10.4} suy ra
\[
\label{II.8.10.6.2}
    \sh{K}'_V \subset \sh{K}_V.
\tag{8.10.6.2}
\]
Như vậy, từ \sref{II.8.10.3} suy ra $j(x)$ thuộc một thành phần bất khả quy duy nhất của $\pi^{-1}(x)$;
cho $z$ là điểm tổng quát của thành phần này, và đặt $z'=q(z)$.
Theo \sref{II.8.10.5}, chứng minh sẽ hoàn tất (có tính đến \sref{II.8.10.6.1} và \sref{II.8.10.6.2}) nếu ta chứng minh tồn tại một tiết diện $g$ của $\sh{K}'_V$ trên $V$ sao cho $g(z)\neq0$.
Nhưng theo giả thiết, hạn chế của $\sh{K}$ bằng hạn chế của $\sh{O}_X$ trong một lân cận mở của $x$;
mặt khác, từ \sref{II.8.10.3} suy ra $z\neq j(x)$, nên $z\in W$, và do đó $(\sh{K}_W)_z=\sh{O}_{V,z}$, từ đó theo định nghĩa $(\sh{K}_C)_{z'}=\sh{O}_{C,z'}$.
Vì $C$ affine, tồn tại một tiết diện $g'$ của $\sh{K}_C$ trên $C$ sao cho $g'(z')\neq0$; lấy $g$ là tiết diện của $\sh{K}'_V$ tương ứng chính tắc với $g'$, ta quả thực có $g(z)\neq0$, điều này hoàn tất chứng minh.
\end{env}

\begin{remark}[8.10.7]
\label{II.8.10.7}
Ta không biết liệu trong phát biểu của \sref{II.8.9.4}, điều kiện~(ii) có thừa hay không.
Dù sao, kết luận không còn đúng nếu không giả sử tồn tại một $Y$-cấu xạ $\pi:V\to X$ sao cho $\pi\circ j=1_X$;
ta chỉ ra ngắn gọn cách thực sự xây dựng một phản ví dụ, mà các chi tiết sẽ chỉ được trình bày về sau.
Lấy $Y=\Spec(k)$, trong đó $k$ là một trường, và $C=\Spec(A)$, trong đó $A=k[T_1,T_2]$, cùng với $Y$-tiết diện $\varepsilon$ tương ứng với đồng cấu tăng cường $A\to k$.
Ký hiệu $C'$ là lược đồ thu được từ $C$ bằng phép nổ điểm đóng $a=\varepsilon(Y)$ của $C$;
nếu $D$ là ảnh ngược của $a$ trong $C'$, ta xét trong $D$ một điểm đóng $b$,
\oldpage[II]{189}
và ký hiệu $V$ là lược đồ thu được từ $C'$ bằng phép nổ $b$;
$X$ là tiền lược đồ con đóng của $V$ được cho bởi ảnh ngược của $a$ qua cấu xạ cấu trúc $q:V\to C$.
Bây giờ ta chứng minh rằng $X$ là hợp của hai thành phần bất khả quy, $X_1$ và $X_2$, trong đó $X_1$ là ảnh ngược của $b$ trong $V$.
Hiển nhiên iđêan $\sh{J}$ của $\sh{O}_V$ xác định $X$ vẫn khả nghịch, nhưng có thể chứng minh rằng $j^*(\sh{J})=\sh{L}$ (trong đó $j$ là phép nhúng chính tắc $X\to V$) không ample, bằng cách xét «bậc» của ảnh ngược của $\sh{L}$ trên $X_1$, vốn sẽ $>0$ nếu $\sh{L}$ ample, nhưng (bằng một phép tính giao sơ cấp) có thể chứng minh rằng bậc này bằng $0$.
\end{remark}

\subsection{Tính duy nhất của các phép co}
\label{subsection:II.8.11}

\begin{lemma}[8.11.1]
\label{II.8.11.1}
Cho $U$ và $V$ là hai tiền lược đồ, và $h=(h_0,\lambda):U\to V$ là một cấu xạ toàn ánh.
Giả sử rằng:
\begin{enumerate}
  \item[$1^\circ$] $\lambda:\sh{O}_V\to h_*(\sh{O}_U)=(h_0)_*(\sh{O}_U)$ là một đẳng cấu.
  \item[$2^\circ$] Không gian nền của $V$ được đồng nhất với không gian thương của không gian nền của $U$ theo quan hệ
  \[
      R:h_0(x)=h_0(y)
  \]
  (điều kiện luôn được thỏa mãn khi cấu xạ $h$ là \emph{mở} hoặc \emph{đóng}, hay, \emph{a fortiori}, khi $h$ là \emph{riêng}).
\end{enumerate}
Khi đó, với mọi tiền lược đồ $W$, ánh xạ
\[
\label{II.8.11.1.1}
    \Hom(V,W)\longrightarrow\Hom(U,W)
\tag{8.11.1.1}
\]
gán cho mỗi cấu xạ $v=(v_0,\nu)$ từ $V$ vào $W$ cấu xạ $u=v\circ h=(u_0,\mu)$, là một song ánh từ $\Hom(V,W)$ lên tập hợp các $u$ sao cho $u_0$ không đổi trên mọi thớ $h_0^{-1}(x)$.
\end{lemma}

\begin{proof}
Hiển nhiên rằng, nếu $u=v\circ h$, do đó $u_0=v_0\circ h_0$, thì $u_0$ không đổi trên mọi tập $h_0^{-1}(x)$.
Ngược lại, nếu $u$ có tính chất này, ta sẽ chứng minh rằng tồn tại duy nhất một $v\in\Hom(V,W)$ sao cho $u=v\circ h$.
Sự tồn tại và tính duy nhất của ánh xạ liên tục $v_0:V\to W$ sao cho $u_0=v_0\circ h_0$ suy ra từ các giả thiết, vì $h_0$ được đồng nhất với ánh xạ chính tắc từ $U$ lên $U/R$.
Ta cũng có thể, bằng cách thay $V$ bởi một tiền lược đồ đẳng cấu nếu cần, giả sử rằng $\lambda$ là đồng nhất;
theo giả thiết, $\mu$ khi đó là một đồng cấu
\[
    \mu:\sh{O}_W\longrightarrow(u_0)_*(\sh{O}_U)
    =(v_0)_*((h_0)_*(\sh{O}_U))
\]
sao cho đồng cấu tương ứng $\mu^\sharp:u_0^*(\sh{O}_W)\to\sh{O}_U$ là \emph{địa phương} trên mỗi thớ.
Vì $(v_0)_*((h_0)_*(\sh{O}_U))=(v_0)_*(\sh{O}_V)$, ta nhất thiết có $\nu=\mu$, và mọi việc quy về chứng minh rằng đồng cấu tương ứng $\nu^\sharp:v_0^*(\sh{O}_W)\to\sh{O}_V$ là địa phương trên mỗi thớ.
Nhưng mọi $y\in V$ đều có dạng $h_0(x)$ với một $x\in U$ nào đó;
đặt $z=v_0(y)=u_0(x)$.
Khi đó, theo \sref[0]{0.3.5.5}, đồng cấu $\mu_x^\sharp$ phân tích thành
\[
    \mu_x^\sharp:\sh{O}_z
    \xrightarrow{\nu_y^\sharp}\sh{O}_y
    \xrightarrow{\lambda_x^\sharp}\sh{O}_x.
\]
Theo giả thiết, $\lambda_x^\sharp$ và $\mu_x^\sharp$ là các đồng cấu địa phương;
do đó, $\lambda_x^\sharp$ biến mọi phần tử khả nghịch của $\sh{O}_y$ thành một phần tử khả nghịch của $\sh{O}_x$;
nếu $\nu_y^\sharp$ biến một phần tử không khả nghịch của $\sh{O}_z$ thành một phần tử khả nghịch của $\sh{O}_y$, thì $\mu_x^\sharp$ sẽ biến phần tử đó của $\sh{O}_z$ thành một phần tử khả nghịch của $\sh{O}_x$, trái với giả thiết, từ đó suy ra bổ đề.
\end{proof}

\begin{corollary}[8.11.2]
\label{II.8.11.2}
Cho $U$ là một tiền lược đồ nguyên và $V$ là một tiền lược đồ chuẩn tắc;
khi đó mọi cấu xạ $h:U\to V$ đóng phổ dụng, song hữu tỉ và radicial đều là một đẳng cấu.
\end{corollary}

\begin{proof}
Nếu $h=(h_0,\lambda)$, thì từ các giả thiết suy ra $h_0$ đơn ánh và đóng, và $h_0(U)$
\oldpage[II]{190}
là trù mật trong $V$, vì thế $h_0$ là một \emph{đồng phôi} từ $U$ lên $V$.
Để chứng minh hệ quả, chỉ cần chứng minh rằng $\lambda:\sh{O}_V\to(h_0)_*(\sh{O}_U)$ là một đẳng cấu: khi đó ta có thể áp dụng \sref{II.8.11.1}, và bổ đề này sẽ chứng minh ánh xạ \sref{II.8.11.1.1} là song ánh (vì mỗi thớ $h_0^{-1}(x)$ chỉ gồm một điểm);
do đó, $h$ là một đẳng cấu.
Vì vấn đề hiển nhiên có tính địa phương trên $V$, ta có thể giả sử $V=\Spec(A)$ là affine, với $A$ là một vành nguyên và đóng nguyên \sref{II.8.8.6.1};
khi đó $h$ tương ứng \sref[I]{I.2.2.4} với một đồng cấu $\vphi:A\to\Gamma(U,\sh{O}_U)$, và mọi việc quy về chứng minh rằng $\vphi$ là một đẳng cấu.
Nếu $K$ là trường phân thức của $A$, thì theo giả thiết, $\Gamma(U,\sh{O}_U)$ cũng có trường phân thức là $K$, và $A$ là một vành con của $\Gamma(U,\sh{O}_U)$, với $\vphi$ là phép nhúng chính tắc \sref[I]{I.8.2.7}.
Vì cấu xạ $h$ thỏa mãn các giả thiết của \sref{II.7.3.11}, $\Gamma(U,\sh{O}_U)$ là một vành con của bao đóng nguyên của $A$ trong $K$, nên bằng $A$ theo giả thiết.
\end{proof}

\begin{remark}[8.11.3]
\label{II.8.11.3}
Ta sẽ thấy trong Chương~III \sref[III]{III.4.4.11} rằng, khi $V$ là một tiền lược đồ \emph{Noether địa phương}, mọi cấu xạ $h:U\to V$ riêng và quasi-hữu hạn (đặc biệt, mọi cấu xạ thỏa mãn các giả thiết của \sref{II.8.11.2}) nhất thiết là \emph{hữu hạn}.
Kết luận của \sref{II.8.11.2} khi ấy suy ra, trong trường hợp này, từ \sref{II.6.1.15}.
\end{remark}

\begin{env}[8.11.4]
\label{II.8.11.4}
Bây giờ ta sẽ thấy rằng, trong tiêu chuẩn Grauert \sref{II.8.9.1}, thường có thể khẳng định tiền lược đồ $C$ và «phép co» $q$ được xác định \emph{duy nhất về bản chất}.
\end{env}

\begin{lemma}[8.11.5]
\label{II.8.11.5}
Cho $Y$ là một tiền lược đồ, $p:X\to Y$ là một cấu xạ riêng, $\sh{L}$ là một $\sh{O}_X$-môđun khả nghịch ample đối với $p$, $C$ là một $Y$-tiền lược đồ, $\varepsilon:Y\to C$ là một $Y$-tiết diện, và $q:V=\bb{V}(\sh{L})\to C$ là một $Y$-cấu xạ, sao cho biểu đồ trong \sref{II.8.9.1.1} giao hoán.
Giả sử thêm rằng, nếu $p=(p_0,\theta)$, thì $\theta:\sh{O}_Y\to p_*(\sh{O}_X)$ là một đẳng cấu.
Đặt $\sh{S}'=\bigoplus_{n\geq0}p_*(\sh{L}^{\otimes n})$ và $C'=\Spec(\sh{S}')$, và cho $q':\bb{V}(\sh{L})\to C'$ là $Y$-cấu xạ chính tắc \sref{II.8.8.5}.
Khi đó tồn tại duy nhất một $Y$-cấu xạ $u:C'\to C$ sao cho $q=u\circ q'$.
\end{lemma}

\begin{proof}
Giả thiết về $\theta$ đặc biệt kéo theo rằng $p$ toàn ánh;
vì theo \sref{II.8.8.4}, hạn chế của $q'$ lên $\bb{V}(\sh{L})\setmin j(X)$ là một \emph{đẳng cấu} lên $C'\setmin\varepsilon'(Y)$ (trong đó $\varepsilon'$ là tiết diện đỉnh của $C'$), suy ra từ \sref{II.8.8.4} rằng $q'$ \emph{riêng và toàn ánh};
hơn nữa, theo \sref{II.8.8.6}, nếu đặt $q'=(q'_0,\tau)$, thì $\tau:\sh{O}_{C'}\to q'_*(\sh{O}_V)$ là một đẳng cấu.
Do đó ta ở trong các điều kiện áp dụng \sref{II.8.11.1}, và bổ đề sẽ được chứng minh nếu ta chứng minh rằng $q$ không đổi trên mọi thớ $q'{}^{-1}(z')$, với $z'\in C'$.
Điều kiện này hiển nhiên được thỏa mãn khi $z'\notin\varepsilon'(Y)$.
Nếu $z'\in\varepsilon'(Y)$, tồn tại duy nhất một $y\in Y$ sao cho $z'=\varepsilon'(y)$; và từ tính giao hoán của \sref{II.8.8.5.2}, cùng với việc $q'$ ánh xạ $\bb{V}(\sh{L})\setmin j(X)$ vào $C'\setmin\varepsilon'(Y)$, ta có $q'{}^{-1}(z')=j(p^{-1}(y))$;
tính giao hoán của biểu đồ \sref{II.8.9.1.1} khi đó chứng minh khẳng định của ta.
\end{proof}

\begin{corollary}[8.11.6]
\label{II.8.11.6}
Dưới các giả thiết của \sref{II.8.11.5}, giả sử thêm rằng $q$ riêng và hạn chế của $q$ lên $\bb{V}(\sh{L})\setmin j(X)$ là một đẳng cấu lên $C\setmin\varepsilon(Y)$.
Khi đó cấu xạ $u$ đóng phổ dụng, toàn ánh và radicial, và hạn chế của nó lên $C'\setmin\varepsilon'(Y)$ là một đẳng cấu lên $C\setmin\varepsilon(Y)$.
\end{corollary}

\begin{proof}
Vì $q'$ là một đẳng cấu từ $\bb{V}(\sh{L})\setmin j(X)$ lên $C'\setmin\varepsilon'(Y)$ \sref{II.8.8.4}, khẳng định cuối cùng suy ra ngay từ $q=u\circ q'$.
Hơn nữa, tính giao hoán của các biểu đồ
\oldpage[II]{191}
\sref{II.8.8.5.2} và \sref{II.8.9.1.1} cho thấy rằng hạn chế của $u$ lên tiền lược đồ con đóng $\varepsilon'(Y)$ của $C'$ là một đẳng cấu lên tiền lược đồ con đóng $\varepsilon(Y)$ của $C$; từ đó suy ra ngay rằng, với mọi $z'\in\varepsilon'(Y)$, nếu $z=u(z')$, thì $u$ xác định một đẳng cấu từ $\kres(z)$ lên $\kres(z')$.
Các nhận xét này chứng minh rằng $u$ song ánh và radicial;
hơn nữa, nếu $\psi$ và $\psi'$ là các cấu xạ cấu trúc $C\to Y$, $C'\to Y$, thì $\psi'=\psi\circ u$, và vì $\psi'$ tách \sref{II.1.2.4}, nên $u$ cũng tách \sref[I]{I.5.5.1}[v].
Mặt khác, ta đã thấy trong chứng minh Bổ đề~\sref{II.8.11.5} rằng $q'$ toàn ánh;
vì $q=u\circ q'$ riêng, cuối cùng từ \sref{II.5.4.3} và \sref{II.5.4.9} ta kết luận rằng $u$ đóng phổ dụng.
\end{proof}

\begin{proposition}[8.11.7]
\label{II.8.11.7}
Cho $Y$ là một tiền lược đồ, $X$ là một tiền lược đồ nguyên, $p:X\to Y$ là một cấu xạ riêng, $\sh{L}$ là một $\sh{O}_X$-môđun khả nghịch ample đối với $p$, $C$ là một $Y$-tiền lược đồ chuẩn tắc, $\varepsilon:Y\to C$ là một $Y$-tiết diện, và $q:V=\bb{V}(\sh{L})\to C$ là một $Y$-cấu xạ, sao cho sơ đồ \sref{II.8.9.1.1} giao hoán.
Giả sử thêm rằng $q$ là riêng và hạn chế của nó trên $V\setmin j(X)$ là một đẳng cấu lên $C\setmin\varepsilon(Y)$.
Giả sử thêm rằng, nếu $p=(p_0,\theta)$, thì $\theta:\sh{O}_Y\to p_*(\sh{O}_X)$ là một đẳng cấu.
Đặt $\sh{S}'=\bigoplus_{n\geq0}p_*(\sh{L}^{\otimes n})$ và $C'=\Spec(\sh{S}')$, và cho $q':\bb{V}(\sh{L})\to C'$ là $Y$-cấu xạ chính tắc \sref{II.8.8.5}.
Khi đó $Y$-cấu xạ duy nhất $u:C'\to C$ sao cho $q=u\circ q'$ là một đẳng cấu.
\end{proposition}

\begin{proof}
Từ \sref{II.8.8.6} suy ra $C'$ nguyên;
vì $u$ là một đồng phôi của các không gian nền $C'\to C$ (do $u$ song ánh và đóng theo \sref{II.8.11.6}), $C$ bất khả quy, do đó nguyên, và vì hạn chế của $u$ lên một tập mở khác rỗng của $C'$ là một đẳng cấu lên một tập mở của $C$, nên $u$ song hữu tỉ.
Vì $C$ được giả sử chuẩn tắc, chỉ cần áp dụng \sref{II.8.11.2} để thu được kết luận.
\end{proof}

\begin{remark}[8.11.8]
\label{II.8.11.8}
\begin{enumerate}
  \item[\rm{(i)}] Lưu ý rằng giả thiết $C$ chuẩn tắc kéo theo $X$ cũng chuẩn tắc.
  Thật vậy, $C'=\Spec(\sh{S}')$ khi đó chuẩn tắc vì đẳng cấu với $C$, và nguyên theo \sref{II.8.8.6}; do đó ta kết luận rằng $\Proj(\sh{S}')$ là \emph{chuẩn tắc}.
  Thật vậy, vấn đề mang tính địa phương trên $Y$;
  nếu $Y$ là affine, với $\sh{S}'=\widetilde{S'}$, thì vành $S'=\Gamma(C',\sh{O}_{C'})$ nguyên và đóng nguyên \sref{II.8.8.6.1}, nên với mọi phần tử thuần nhất $f\in S'_+$, vành phân bậc $S'_f$ nguyên và đóng nguyên \cite[t.~I, p.~257 và 261]{I-13}, và do đó vành $S'_{(f)}$ gồm các phần tử bậc 0 của nó cũng vậy, bởi vì giao của $S'_f$ với trường phân thức của $S'_{(f)}$ bằng $S'_{(f)}$;
  điều này chứng minh khẳng định của ta \sref{II.6.3.4}.
  Cuối cùng, vì $X$ đẳng cấu với một tiền lược đồ con mở của $\Proj(\sh{S}')$ \sref{II.8.8.1}, nên $X$ quả thực chuẩn tắc.
  Vì vậy có thể phát biểu Mệnh đề~\sref{II.8.11.7} dưới dạng sau:
  \emph{Nếu $X$ nguyên và chuẩn tắc, và $p=(p_0,\theta):X\to Y$ là một cấu xạ riêng sao cho $\theta:\sh{O}_Y\to p_*(\sh{O}_X)$ là một đẳng cấu, thì với mọi $\sh{O}_X$-môđun $\sh{L}$ ample đối với $p$, tồn tại đúng một cách co tiết diện không của $V=\bb{V}(\sh{L})$ sao cho thu được một $Y$-lược đồ chuẩn tắc $C$ và một $Y$-cấu xạ riêng $q:V\to C$.}
  \item[\rm{(ii)}] Khi $p$ riêng, giả thiết $p_*(\sh{O}_X)=\sh{O}_Y$ có thể được xem là một giả thiết phụ trợ, không thực sự hạn chế tính tổng quát của kết quả.
  Thật vậy, nếu giả thiết này không được thỏa mãn, chỉ cần thay $Y$ bằng $Y$-lược đồ $Y'=\Spec(p_*(\sh{O}_X))$ và xem $X$ như một $Y'$-lược đồ.
  Ta sẽ trở lại phương pháp tổng quát này trong Chương~III, \textsection~4.
\end{enumerate}
\end{remark}

\subsection{Các bó quasi-coherent trên các nón chiếu}
\label{subsection:II.8.12}

\begin{env}[8.12.1]
\label{II.8.12.1}
Giữ nguyên các giả thiết và ký hiệu của \sref{II.8.3.1}.
Cho $\sh{M}$ là một $\sh{S}$-môđun phân bậc quasi-coherent;
để tránh mọi nhầm lẫn, ta ký hiệu bởi $\widetilde{\sh{M}}$ $\sh{O}_C$-môđun
\oldpage[II]{192}
quasi-coherent liên kết với $\sh{M}$ \sref{II.1.4.3} khi $\sh{M}$ được xét như một $\sh{S}$-môđun \emph{không phân bậc}, và bởi $\shProj_0(\sh{M})$ $\sh{O}_X$-môđun quasi-coherent liên kết với $\sh{M}$, lần này $\sh{M}$ được xét như một $\sh{S}$-môđun phân bậc (nói cách khác, đó là $\sh{O}_X$-môđun được ký hiệu $\widetilde{\sh{M}}$ trong \sref{II.3.2.2}).
Ngoài ra, ta đặt
\[
\label{II.8.12.1.1}
  \sh{M}_X=\shProj(\sh{M})
  =\bigoplus_{n\in\bb{Z}}\shProj_0(\sh{M}(n));
\tag{8.12.1.1}
\]
vì $\sh{O}_X$-đại số phân bậc quasi-coherent $\sh{S}_X$ được định nghĩa bởi \sref{II.8.6.1.1}, $\shProj(\sh{M})$ được trang bị một cấu trúc $\sh{S}_X$-môđun phân bậc (quasi-coherent), nhờ các đồng cấu chính tắc \sref{II.3.2.6.1}
\[
\label{II.8.12.1.2}
  \sh{O}_X(m)\otimes_{\sh{O}_X}
  \shProj_0(\sh{M}(n))
  \to \shProj_0(\sh{S}(m)\otimes_{\sh{S}}\sh{M}(n))
  \to \shProj_0(\sh{M}(m+n)),
\tag{8.12.1.2}
\]
việc kiểm tra các tiên đề môđun được thực hiện nhờ sơ đồ giao hoán \sref{II.2.5.11.4}.

Nếu $Y=\Spec(A)$ là affine, $\sh{S}=\widetilde{S}$ và $\sh{M}=\widetilde{M}$, trong đó $S$ là một $A$-đại số phân bậc và $M$ là một $S$-môđun phân bậc, thì với mọi phần tử thuần nhất $f\in S_+$, ta có
\[
\label{II.8.12.1.3}
  \Gamma(X_f,\shProj(\widetilde{M}))=M_f
\tag{8.12.1.3}
\]
theo các định nghĩa và \sref{II.8.2.9.1}.

Bây giờ xét $\widehat{\sh{S}}$-môđun phân bậc quasi-coherent
\[
\label{II.8.12.1.4}
  \widehat{\sh{M}}
  =\sh{M}\otimes_{\sh{S}}\widehat{\sh{S}}
\tag{8.12.1.4}
\]
(trong đó $\widehat{\sh{S}}$ được định nghĩa bởi \sref{II.8.3.1.1});
từ đó thu được một $\sh{O}_{\widehat{C}}$-môđun phân bậc quasi-coherent $\shProj_0(\widehat{\sh{M}})$, mà ta cũng ký hiệu là
\[
\label{II.8.12.1.5}
  \sh{M}^{\square}=\shProj_0(\widehat{\sh{M}}).
\tag{8.12.1.5}
\]
Hiển nhiên \sref{II.3.2.4}, $\sh{M}^{\square}$ là một hàm tử cộng tính, chính xác theo $\sh{M}$, giao hoán với các tổng trực tiếp và các giới hạn quy nạp.
\end{env}

\begin{proposition}[8.12.2]
\label{II.8.12.2}
Với các ký hiệu của \sref{II.8.3.2}, ta có các đẳng cấu chính tắc có tính hàm tử
\[
\label{II.8.12.2.1}
  i^*(\sh{M}^{\square})\simto\widetilde{\sh{M}},
  \qquad
  j^*(\sh{M}^{\square})\simto\shProj_0(\sh{M}).
\tag{8.12.2.1}
\]
\end{proposition}

\begin{proof}
Thật vậy, $i^*(\sh{M}^{\square})$ được đồng nhất chính tắc với
$(\widehat{\sh{M}}/(\bb{z}-1)\widehat{\sh{M}})\supertilde$ trên
$\Spec(\widehat{\sh{S}}/(\bb{z}-1)\widehat{\sh{S}})$ theo \sref{II.3.2.3};
đẳng cấu thứ nhất trong \sref{II.8.12.2.1} khi đó suy ra ngay \sref{II.1.4.1} từ đẳng cấu chính tắc
$\widehat{\sh{M}}/(\bb{z}-1)\widehat{\sh{M}}\simto\sh{M}$.
Mặt khác, phép nhúng chính tắc $j:X\to\widehat{C}$ tương ứng với đồng cấu chính tắc $\widehat{\sh{S}}\to\sh{S}$ có hạt nhân $\bb{z}\widehat{\sh{S}}$ \sref{II.8.3.2};
đồng cấu thứ hai trong \sref{II.8.12.2.1} là trường hợp riêng của đồng cấu chính tắc \sref{II.3.5.2}[ii], bởi vì ở đây ta có
$\widehat{\sh{M}}\otimes_{\widehat{\sh{S}}}\sh{S}=\sh{M}$;
để kiểm tra rằng đó là một đẳng cấu, ta có thể quy về trường hợp $Y=\Spec(A)$ là affine, $\sh{S}=\widetilde{S}$ và $\sh{M}=\widetilde{M}$;
dựa vào \sref{II.2.8.8}, việc kiểm tra rằng với mọi $f$ thuần nhất trong $S_+$, đồng cấu trên, khi hạn chế lên $X_f$, là một đẳng cấu khi đó trở nên hiển nhiên.
\end{proof}
\oldpage[II]{193}
Theo lối lạm dụng ngôn ngữ, do tồn tại đẳng cấu thứ nhất trong \sref{II.8.12.2.1}, ta vẫn nói rằng $\sh{M}^{\square}$ là \emph{bao đóng xạ ảnh} của $\sh{O}_C$-môđun $\widetilde{\sh{M}}$ (ngầm hiểu rằng dữ liệu của $\sh{O}_C$-môđun $\widetilde{\sh{M}}$ bao gồm phân bậc của $\sh{S}$-môđun $\sh{M}$).

\begin{env}[8.12.3]
\label{II.8.12.3}
Với các ký hiệu của \sref{II.8.3.5}, ta có một đồng cấu chính tắc có tính hàm tử
\[
\label{II.8.12.3.1}
  p^*(\shProj_0(\sh{M}))
  \to\sh{M}^{\square}|\widehat{E}.
\tag{8.12.3.1}
\]
Thật vậy, đây là một trường hợp riêng của đồng cấu $\nu^{\sharp}$ được định nghĩa tổng quát trong \sref{II.3.5.6}.
Nếu $Y=\Spec(A)$ là affine, $\sh{S}=\widetilde{S}$ và $\sh{M}=\widetilde{M}$, thì dựa vào \sref{II.2.8.8}, ta thấy hạn chế của \sref{II.8.12.3.1} lên $p^{-1}(X_f)=\widehat{C}_f$ (với $f$ thuần nhất trong $S_+$) tương ứng với đồng cấu chính tắc
\[
\label{II.8.12.3.2}
  M_{(f)}\otimes_{S_{(f)}}S_f^{\leq}\to M_f^{\leq}
\tag{8.12.3.2}
\]
có tính đến \sref{II.8.2.3.2} và \sref{II.8.2.5.2}.
\end{env}

\begin{env}[8.12.4]
\label{II.8.12.4}
Giả sử các giả thiết của \sref{II.8.5.1} được thỏa mãn và giữ nguyên các ký hiệu ở đó.
Từ \sref{II.1.5.6} suy ra rằng, với mọi $\sh{S}$-môđun phân bậc quasi-coherent $\sh{M}$, một mặt ta có một đẳng cấu chính tắc
\[
\label{II.8.12.4.1}
  \Phi^*(\widetilde{\sh{M}})
  \simto
  (q^*(\sh{M})\otimes_{q^*(\sh{S})}\sh{S}')\supertilde
\tag{8.12.4.1}
\]
giữa các $\sh{O}_{C'}$-môđun;
mặt khác, \sref{II.3.5.6} kéo theo sự tồn tại của một $\Proj(\vphi)$-cấu xạ chính tắc
\[
\label{II.8.12.4.2}
  \shProj_0(\sh{M})\to
  \shProj_0(q^*(\sh{M})\otimes_{q^*(\sh{S})}\sh{S}')
  |G(\vphi)
\tag{8.12.4.2}
\]
và cũng của một $\widehat{\Phi}$-cấu xạ chính tắc
\[
\label{II.8.12.4.3}
  \shProj_0(\widehat{\sh{M}})\to
  \shProj_0(q^*(\widehat{\sh{M}})
  \otimes_{q^*(\widehat{\sh{S}})}\widehat{\sh{S}'})
  |G(\widehat{\vphi}).
\tag{8.12.4.3}
\]
\end{env}

\begin{env}[8.12.5]
\label{II.8.12.5}
Bây giờ xét bối cảnh của \sref{II.8.6.1}, với cùng các ký hiệu;
như vậy, trong những điều trên ta lấy $Y'=X$, cấu xạ $q:X\to Y$ là cấu xạ cấu trúc, và $\vphi$ là $q$-cấu xạ chính tắc \sref{II.8.6.1.2}.
Khi đó ta có một đẳng cấu chính tắc
\[
\label{II.8.12.5.1}
  q^*(\sh{M})\otimes_{q^*(\sh{S})}\sh{S}_X^{\geq}
  \simto\sh{M}_X^{\geq}
\tag{8.12.5.1}
\]
khi đặt
$\sh{M}_X^{\geq}=\bigoplus_{n\geq0}\shProj_0(\sh{M}(n))$.
Thật vậy, ta có thể quy về trường hợp $Y=\Spec(A)$ là affine, $\sh{S}=\widetilde{S}$ và $\sh{M}=\widetilde{M}$, rồi định nghĩa đẳng cấu \sref{II.8.12.5.1} trên mỗi tập mở affine $X_f$ (với $f$ thuần nhất trong $S_+$), bằng cách kiểm tra tính tương thích khi chuyển sang một bội thuần nhất của $f$.
Thế nhưng, hạn chế lên $X_f$ của vế trái trong \sref{II.8.12.5.1} là
$\widetilde{M}'=((M\otimes_A S_{(f)})\otimes_{S\otimes_A S_{(f)}}S_f^{\geq})\supertilde$
theo \sref{II.8.6.2.1};
vì ta có một đẳng cấu chính tắc từ $M\otimes_A S_{(f)}$ lên $M\otimes_S(S\otimes_A S_{(f)})$, nên suy ra một đẳng cấu từ $\widetilde{M}'$ lên $(M\otimes_S S_f^{\geq})\supertilde$, và môđun sau đẳng cấu chính tắc theo \sref{II.8.2.9.1} với hạn chế lên $X_f$ của vế phải trong \sref{II.8.12.5.1}, đồng thời thỏa mãn các điều kiện tương thích cần thiết.
\oldpage[II]{194}
Thay $\sh{M}$ bởi $\widehat{\sh{M}}$, $\sh{S}$ bởi $\widehat{\sh{S}}$ và $\sh{S}_X$ bởi $(\sh{S}_X^{\geq})^{\wedge}$ trong lập luận trên, ta cũng có một đẳng cấu chính tắc
\[
\label{II.8.12.5.2}
  q^*(\widehat{\sh{M}})
  \otimes_{q^*(\widehat{\sh{S}})}(\sh{S}_X^{\geq})^{\wedge}
  \simto(\sh{M}_X^{\geq})^{\wedge}.
\tag{8.12.5.2}
\]
Nhắc lại từ \sref{II.8.6.2} rằng cấu xạ cấu trúc $u:\Proj(\sh{S}_X^{\geq})\to X$ là một đẳng cấu, trước hết từ những điều trên ta suy ra một $u$-đẳng cấu chính tắc
\[
\label{II.8.12.5.3}
  \shProj_0\sh{M}
  \simto\shProj_0(\sh{M}_X^{\geq})
\tag{8.12.5.3}
\]
như một trường hợp riêng của \sref{II.8.12.4.2}.
Thật vậy, với các ký hiệu trong chứng minh \sref{II.8.6.2}, điều này quy về việc nhận thấy đồng cấu chính tắc
$M_{(f)}\otimes_{S_{(f)}}(S_f^{\geq})^{(d)}\to(M_f^{\geq})^{(d)}$ là một đẳng cấu khi $f\in S_d$, điều này hiển nhiên.

Thứ hai, đẳng cấu \sref{II.8.12.5.2} cho phép, bằng cách lần này áp dụng \sref{II.8.12.4.3} cho cấu xạ chính tắc $r=\Proj(\widehat{\alpha}):\widehat{C}_X\to\widehat{C}$, thu được một $r$-cấu xạ chính tắc
\[
\label{II.8.12.5.4}
  \sh{M}^{\square}\to(\sh{M}_X^{\geq})^{\square}.
\tag{8.12.5.4}
\]
Bây giờ nhắc lại từ \sref{II.8.6.2} rằng các hạn chế của $r$ lên các nón bỏ đỉnh $\widehat{E}_X$ và $E_X$ lần lượt là các \emph{đẳng cấu} lên $\widehat{E}$ và $E$.
Hơn nữa:
\end{env}

\begin{proposition}[8.12.6]
\label{II.8.12.6}
Các hạn chế lên $\widehat{E}_X$ và $E_X$ của $r$-cấu xạ chính tắc \sref{II.8.12.5.4} là các đẳng cấu
\[
\label{II.8.12.6.1}
  \sh{M}^{\square}|\widehat{E}
  \simto(\sh{M}_X^{\geq})^{\square}|\widehat{E}_X
\tag{8.12.6.1}
\]
\[
\label{II.8.12.6.2}
  \widetilde{\sh{M}}|E
  \simto(\sh{M}_X^{\geq})\supertilde|E_X.
\tag{8.12.6.2}
\]
\end{proposition}

\begin{proof}
Ta quy về trường hợp $Y$ affine như trong chứng minh \sref{II.8.6.2} và giữ nguyên các ký hiệu ở đó;
trở về các định nghĩa trong \sref{II.2.8.8}, ta phải chứng minh rằng đồng cấu chính tắc
\[
  \widehat{M}_{(f)}\otimes_{\widehat{S}_{(f)}}
  (S_f^{\geq})_{(f/1)}^{\wedge}
  \to(M\otimes_S S_f^{\geq})_{(f/1)}^{\wedge}
\]
là một đẳng cấu;
nhưng theo \sref{II.8.2.3.2} và \sref{II.8.2.5.2}, vế trái được đồng nhất chính tắc với
$M_f^{\leq}\otimes_{S_f^{\leq}}(S_f^{\geq})_{f/1}^{\leq}$, do đó với $M_f^{\leq}$ theo \sref{II.8.2.7.2}, còn vế phải được đồng nhất với $(M_f^{\geq})_{f/1}^{\leq}$, do đó cũng với $M_f^{\leq}$ theo \sref{II.8.2.9.2};
từ đó suy ra kết luận đối với \sref{II.8.12.6.1};
công thức \sref{II.8.12.6.2} khi đó suy ra từ \sref{II.8.12.6.1} và \sref{II.8.12.2.1}.
\end{proof}

\begin{corollary}[8.12.7]
\label{II.8.12.7}
Với các phép đồng nhất của \sref{II.8.6.3}, hạn chế của $(\sh{M}_X^{\geq})^{\square}$ lên $\widehat{E}_X$ được đồng nhất với $(\sh{M}_X^{\leq})\supertilde$, và hạn chế của $(\sh{M}_X^{\geq})^{\square}$ lên $E_X$ được đồng nhất với $\widetilde{\sh{M}}_X$.
\end{corollary}

\begin{proof}
Thật vậy, ta quy về trường hợp affine, và khẳng định suy ra từ việc đồng nhất $(M_f^{\geq})_{f/1}^{\leq}$ với $M_f^{\leq}$ và $(M_f^{\geq})_{f/1}$ với $M_f$ \sref{II.8.2.9.2}.
\end{proof}

\begin{proposition}[8.12.8]
\label{II.8.12.8}
Dưới các giả thiết của \sref{II.8.6.4}, đồng cấu chính tắc \sref{II.8.12.3.1} là một đẳng cấu.
\end{proposition}

\begin{proof}
Xét rằng $\Proj(\sh{S}_X^{\geq})\to X$ là một đẳng cấu \sref{II.8.6.2} và các
\oldpage[II]{195}
đẳng cấu \sref{II.8.12.5.3} và \sref{II.8.12.6.1}, ta quy về việc chứng minh mệnh đề tương ứng cho đồng cấu chính tắc
$p_X^*(\shProj_0(\sh{M}_X^{\geq}))
\to(\sh{M}_X^{\geq})^{\square}|\widehat{E}_X$;
nói cách khác, ta quy về trường hợp $\sh{S}_1$ là một $\sh{O}_Y$-môđun khả nghịch và $\sh{S}$ được sinh bởi $\sh{S}_1$.
Với các ký hiệu của \sref{II.8.12.3}, khi đó, với $f\in S_1$, ta có $S_f^{\leq}=S_{(f)}[1/f]$, và đồng cấu chính tắc
$M_{(f)}\otimes_{S_{(f)}}S_f^{\leq}\to M_f^{\leq}$ là một đẳng cấu theo định nghĩa của $M_f^{\leq}$.
\end{proof}

\begin{env}[8.12.9]
\label{II.8.12.9}
Bây giờ xét các $\sh{S}$-môđun quasi-coherent
\[
  \sh{M}_{[n]}=\bigoplus_{m\geq n}\sh{M}_m
\]
và (với các ký hiệu của \sref{II.8.7.2}) $\sh{S}^{\natural}$-môđun phân bậc quasi-coherent
\[
\label{II.8.12.9.1}
  \sh{M}^{\natural}
  =\left(\bigoplus_{n\geq0}\sh{M}_{[n]}\right)\supertilde.
\tag{8.12.9.1}
\]
Ta đã thấy trong \sref{II.8.7.3} rằng tồn tại một $C$-đẳng cấu chính tắc
$h:C_X\simto\Proj(\sh{S}^{\natural})$.
Hơn nữa:
\end{env}

\begin{proposition}[8.12.10]
\label{II.8.12.10}
Tồn tại một $h$-đẳng cấu chính tắc
\[
\label{II.8.12.10.1}
  \shProj_0(\sh{M}^{\natural})
  \simto\widetilde{\sh{M}}_X.
\tag{8.12.10.1}
\]
\end{proposition}

\begin{proof}
Ta lập luận như trong \sref{II.8.7.3}, lần này dùng sự tồn tại của di-đẳng cấu \sref{II.8.2.9.3} thay cho \sref{II.8.2.7.3}.
Ta để các chi tiết cho độc giả.
\end{proof}

\subsection{Các bao đóng xạ ảnh của các bó con và các lược đồ con đóng}
\label{subsection:II.8.13}

\begin{env}[8.13.1]
\label{II.8.13.1}
Với các giả thiết và ký hiệu như trong \sref{II.8.12.1}, xét một $\sh{S}$-môđun con quasi-coherent $\sh{N}$ của $\sh{M}$, \emph{không nhất thiết phân bậc}.
Khi đó ta có thể xét $\sh{O}_C$-môđun quasi-coherent $\widetilde{\sh{N}}$ liên kết với $\sh{N}$, là một $\sh{O}_C$-môđun con của $\widetilde{\sh{M}}$.
Mặt khác, ta đã thấy trong \sref{II.8.12.2.1} rằng $\widetilde{\sh{M}}$ được đồng nhất với hạn chế của $\sh{M}^{\square}$ lên $C$.
Vì phép nhúng chính tắc $i:C\to\widehat{C}$ là một cấu xạ affine \sref{II.8.3.2}, và \emph{a fortiori} quasi-compact, nên \emph{mở rộng chính tắc} $(\widetilde{\sh{N}})^-$, tức $\sh{O}_{\widehat{C}}$-môđun con lớn nhất được chứa trong $\sh{M}^{\square}$ và cảm sinh $\widetilde{\sh{N}}$ trên $C$, là một $\sh{O}_{\widehat{C}}$-môđun \emph{quasi-coherent} \sref[I]{I.9.4.2}.
Ta sẽ mô tả nó tường minh hơn bằng một $\widehat{\sh{S}}$-môđun phân bậc.
\end{env}

\begin{env}[8.13.2]
\label{II.8.13.2}
Để làm điều đó, với mỗi số nguyên $n\geq0$, xét đồng cấu
$\bigoplus_{i\leq n}\sh{M}_i\to\sh{M}$ mà, với mọi tập mở $U$ của $Y$, gửi họ
\[
  (s_i)\in\bigoplus_{i\leq n}\Gamma(U,\sh{M}_i)
\]
tới tiết diện $\sum_i s_i\in\Gamma(U,\sh{M})$.
Ký hiệu bởi $\sh{N}'_n$ ảnh ngược của $\sh{N}$ qua đồng cấu này; đó là một $\sh{O}_Y$-môđun con quasi-coherent của $\bigoplus_{i\leq n}\sh{M}_i$.
Bây giờ xét đồng cấu
$\bigoplus_{i\leq n}\sh{M}_i\to\widehat{\sh{M}}=\sh{M}[\bb{z}]$ gửi $(s_i)$ tới tiết diện
$\sum_{i\leq n}s_i\bb{z}^{n-i}\in\Gamma(U,\widehat{\sh{M}}_n)$, và cho $\sh{N}_n$ là ảnh của $\sh{N}'_n$ qua đồng cấu này;
ta kiểm tra ngay rằng
$\overline{\sh{N}}=\bigoplus_{n\geq0}\sh{N}_n$ là một $\widehat{\sh{S}}$-môđun con phân bậc (quasi-coherent) của $\widehat{\sh{M}}$;
ta nói rằng $\overline{\sh{N}}$ được suy ra từ $\sh{N}$ bằng \emph{phép thuần nhất hóa}, nhờ «biến thuần nhất hóa» $\bb{z}$.
Ta sẽ ký hiệu
\oldpage[II]{196}
rằng nếu $\sh{N}$ vốn đã là một $\sh{S}$-môđun con \emph{phân bậc} của $\sh{M}$, thì $\overline{\sh{N}}$ đồng nhất được với tổng trực tiếp của các thành phần $\widehat{\sh{N}}_n$ bậc $n\geq0$ trong $\widehat{\sh{N}}=\sh{N}[\bb{z}]$.
\end{env}

\begin{proposition}[8.13.3]
\label{II.8.13.3}
$\sh{O}_{\widehat{C}}$-môđun $\shProj_0(\overline{\sh{N}})$ là mở rộng chính tắc $(\widetilde{\sh{N}})^-$ của $\widetilde{\sh{N}}$ lên $\widehat{C}$.
\end{proposition}

\begin{proof}
Theo định nghĩa của mở rộng chính tắc \sref[I]{I.9.4.1}, vấn đề mang tính địa phương trên $Y$ và $\widehat{C}$.
Do đó, ta có thể giả sử ngay rằng $Y=\Spec(A)$ là affine, với $\sh{S}=\widetilde{S}$, $\sh{M}=\widetilde{M}$ và $\sh{N}=\widetilde{N}$, trong đó $N$ là một $S$-môđun con không nhất thiết phân bậc của $M$.
Hơn nữa, theo \sref{II.8.3.2.6}, $\widehat{C}$ là hợp của các tập mở affine $\widehat{C}_z=C$ và $\widehat{C}_f=\Spec(S_f^\leq)$ (với $f$ thuần nhất trong $S_+$).
Vì vậy, chỉ cần chứng minh rằng: (1) hạn chế của $\shProj_0(\overline{\sh{N}})$ lên $C$ là $\widetilde{\sh{N}}$; (2) hạn chế của $\shProj_0(\overline{\sh{N}})$ lên mỗi $\widehat{C}_f$ là mở rộng chính tắc của hạn chế của $\widetilde{\sh{N}}$ lên $C\cap\widehat{C}_f=\Spec(S_f)$ \sref{II.8.3.2.6}.
Đối với điểm thứ nhất, lưu ý rằng $\shProj_0(\overline{\sh{N}})|C$ đồng nhất được với $(\overline{N}_{(\bb{z})})\supertilde$ \sref{II.8.3.2.4};
nhưng $\overline{N}_{(\bb{z})}$ đồng nhất chính tắc được \sref{II.2.2.5} với ảnh của $\overline{N}$ trong $\widehat{M}/(\bb{z}-1)\widehat{M}$, và qua đẳng cấu chính tắc từ môđun sau đến $M$ \sref{II.8.2.5}, ảnh này đồng nhất được với $N$, theo định nghĩa của $\overline{N}$ trong \sref{II.8.13.2}.

Để chứng minh điểm thứ hai, lưu ý rằng đơn ánh $i:C\cap\widehat{C}_f\to\widehat{C}_f$ tương ứng với đơn ánh chính tắc $S_f^\leq\to S_f$ \sref{II.8.3.2.6};
ta cũng có $\Gamma(\widehat{C}_f,\sh{M}^\square)=M_f^\leq$, $\Gamma(\widehat{C}_f,i_*(\widetilde{\sh{N}}))=N_f$ và, theo \sref{II.8.12.2.1}, $\Gamma(\widehat{C}_f,i_*(i^*(\sh{M}^\square)))=M_f$.
Xét đến \sref[I]{I.9.4.2}, ta được quy về việc chứng minh rằng $\overline{N}_{(f)}\subset\widehat{M}_{(f)}=M_f^\leq$ đồng nhất chính tắc được với ảnh ngược của $N_f$ qua đơn ánh chính tắc $M_f^\leq\to M_f$.
Thật vậy, cho $d=\deg(f)>0$, và giả sử một phần tử $(\sum_{k\leq md}x_k)/f^m$ của $M_f$ (với $x_k\in M_k$) có dạng $y/f^m$, trong đó $y\in N$.
Bằng cách nhân $y$ và các $x_k$ với cùng một lũy thừa thích hợp $f^h$, ta có thể giả sử ngay rằng $\sum_{k\leq md}x_k=y$.
Nhưng trong phép đồng nhất ở \sref{II.8.2.5.2}, $(\sum_{k\leq md}x_k)/f^m$ tương ứng với $\sum_{k\leq md}x_k\bb{z}^{md-k}/f^m$, và đây quả là một phần tử của $\overline{N}_{(f)}$, vì $\sum_{k\leq md}x_k\in N$;
chiều ngược lại là hiển nhiên.
\end{proof}

\begin{remark}[8.13.4]
\label{II.8.13.4}
\begin{enumerate}
    \item[\rm{(i)}] Trường hợp áp dụng quan trọng nhất của \sref{II.8.13.3} là trường hợp $\sh{M}=\sh{S}$, khi đó $\widetilde{\sh{N}}$ là một bó iđêan quasi-coherent \emph{tùy ý} $\sh{J}$ của $\sh{O}_C$ \sref{II.1.4.3}, tương ứng song ánh với một \emph{tiền lược đồ con đóng} $Z$ của $C$.
        Khi đó, mở rộng chính tắc $\overline{\sh{J}}$ của $\sh{J}$ là bó iđêan quasi-coherent của $\sh{O}_{\widehat{C}}$ xác định \emph{bao đóng} $\overline{Z}$ của $Z$ trong $\widehat{C}$ \sref[I]{I.9.5.10};
        Mệnh đề~\sref{II.8.13.3} cho một cách chính tắc để xác định $\overline{Z}$ bằng một iđêan phân bậc trong $\widehat{\sh{S}}=\sh{S}[\bb{z}]$.
    \item[\rm{(ii)}] Để đơn giản, giả sử $Y$ affine và dùng các ký hiệu trong chứng minh của \sref{II.8.13.3}.
        Với mỗi $x\in N$ khác không, gọi $d(x)$ là bậc lớn nhất trong các bậc của những thành phần thuần nhất $x_i$ của $x$ trong $M$;
        theo định nghĩa, $\overline{N}$ là môđun con của $\widehat{M}$ gồm $0$ và các phần tử có dạng $h(x,k)=\bb{z}^k\sum_{i\leq d(x)}x_i\bb{z}^{d(x)-i}$ (với số nguyên $k\geq0$);
        do đó, như một môđun trên $\widehat{S}=S[\bb{z}]$, nó được sinh bởi các phần tử có dạng
        \[
            h(x,0) = \sum_{i\leq d(x)}x_i\bb{z}^{d(x)-i}.
        \]
\oldpage[II]{197}
        Ta nói $h(x,0)$ được suy ra từ $x$ bằng phép \emph{thuần nhất hóa} với «biến thuần nhất hóa» $\bb{z}$.
        Nhưng vì phép gán $x\mapsto h(x,0)$ không có tính cộng (và, \emph{a fortiori}, không $S$-tuyến tính), \emph{ta không nên lầm tưởng} (ngay cả khi $M=S$) rằng các $h(x,0)$ tạo thành một \emph{hệ sinh} của $\widehat{S}$-môđun phân bậc $\overline{N}$ khi cho $x$ chạy qua một \emph{hệ sinh} của $S$-môđun $N$.
        Tuy nhiên, điều này đúng khi $N$ là một $S$-môđun \emph{cyclic tự do} (trường hợp duy nhất được xét trong hình học đại số sơ cấp), vì nếu $t$ là một cơ sở của $N$, thì $h(t,0)$ sinh $\widehat{S}$-môđun $\overline{N}$.
\end{enumerate}
\end{remark}


\subsection{Bổ sung về các bó liên kết với $\mathcal{S}$-môđun phân bậc}
\label{subsection:II.8.14}

\begin{env}[8.14.1]
\label{II.8.14.1}
Cho $Y$ là một tiền lược đồ, $\sh{S}$ là một $\sh{O}_Y$-đại số quasi-coherent \emph{phân bậc dương}, $X=\Proj(\sh{S})$, và $q:X\to Y$ là cấu xạ cấu trúc (là cấu xạ tách, theo \sref{II.3.1.3}).
Với ký hiệu của \sref{II.8.12.1}, ta đã định nghĩa hàm tử $\sh{M}\mapsto\sh{M}_X=\shProj(\sh{M})$ từ phạm trù các $\sh{S}$-môđun phân bậc quasi-coherent đến phạm trù các $\sh{S}_X$-môđun phân bậc quasi-coherent;
hơn nữa, rõ ràng \sref{II.3.2.4} đây là một hàm tử \emph{cộng tính} và \emph{khớp}, giao hoán với các giới hạn quy nạp.

Ngoài ra, lưu ý rằng từ định nghĩa \sref{II.8.12.1.1} suy ra ngay
\[
\label{II.8.14.1.1}
    \shProj(\sh{M}(n)) = (\shProj(\sh{M}))(n)
    \quad\mbox{với mọi $n\in\bb{Z}$.}
\tag{8.14.1.1}
\]
\end{env}

\begin{env}[8.14.2]
\label{II.8.14.2}
Trước hết, ta sẽ mở rộng các đồng cấu chính tắc $\lambda$ và $\mu$, được định nghĩa trong \sref{II.3.2.6}, cho các $\sh{S}_X$-môđun có dạng $\shProj(\sh{M})$.
Để làm điều này, lưu ý rằng với mọi $m\in\bb{Z}$ và $n\in\bb{Z}$, theo \sref{II.2.1.2.1}, ta có một đồng cấu chính tắc của các $\sh{O}_X$-môđun
\[
\label{II.8.14.2.1}
    \lambda_{mn}:
    \shProj_0(\sh{M}(m))\otimes_{\sh{O}_X}\shProj_0(\sh{N}(n))
    \to
    \shProj_0((\sh{M}\otimes_{\sh{S}}\sh{N})(m+n)),
\tag{8.14.2.1}
\]
và tương tự, theo \sref{II.2.1.2.2}, một đồng cấu chính tắc của các $\sh{O}_X$-môđun
\[
\label{II.8.14.2.2}
    \mu_{mn}:
    \shProj_0((\shHom_{\sh{S}}(\sh{M},\sh{N}))(n-m))
    \to
    \shHom_{\sh{O}_X}(\shProj_0(\sh{M}(m)),\shProj_0(\sh{N}(n)))
\tag{8.14.2.2}
\]
với mọi $\sh{S}$-môđun phân bậc quasi-coherent $\sh{M}$ và $\sh{N}$.
Điều này cảm sinh một đồng cấu
\[
    \mu_k:
    \shProj_0((\shHom_{\sh{S}}(\sh{M},\sh{N}))(k))
    \to
    (\shHom_{\sh{S}_X}(\shProj(\sh{M}),\shProj(\sh{N})))_k
\]
bằng cách gửi mỗi $u\in\Gamma(U,\shProj_0((\shHom_{\sh{S}}(\sh{M},\sh{N}))(k)))$ đến đồng cấu $\mu_k(u)$ bậc $k$ của các $\bb{Z}$-môđun phân bậc $\Gamma(U,\shProj(\sh{M}))\to\Gamma(U,\shProj(\sh{N}))$ (trong đó $U$ mở trong $X$), sao cho trên mỗi $\Gamma(U,\shProj_0(\sh{M}(m)))$ nó trùng với $\mu_{m,m+k}(u)$;
hơn nữa, bằng cách trở lại định nghĩa của các $\mu_{mn}$ \sref{II.2.5.12.1}, ta thấy ngay $\mu_k(u)$ thực ra là một đồng cấu bậc $k$ của các $\Gamma(U,\sh{S}_X)$-môđun phân bậc, và các $\mu_k$ xác định một đồng cấu của các \emph{$\sh{S}_X$-môđun phân bậc}
\[
\label{II.8.14.2.3}
    \shProj(\shHom_{\sh{S}}(\sh{M},\sh{N}))
    \to
    \shHom_{\sh{S}_X}(\shProj(\sh{M}),\shProj(\sh{N})).
\tag{8.14.2.3}
\]

Tương tự, xét đến biểu đồ kết hợp \sref{II.2.5.11.4}, các đồng cấu \sref{II.8.14.2.1} cho một đồng cấu của các \emph{$\sh{S}_X$-môđun phân bậc}
\[
\label{II.8.14.2.4}
    \lambda:
    \shProj(\sh{M})\otimes_{\sh{S}_X}\shProj(\sh{N})
    \to
    \shProj(\sh{M}\otimes_{\sh{S}}\sh{N}).
\tag{8.14.2.4}
\]
\end{env}
\oldpage[II]{198}

\begin{proposition}[8.14.3]
\label{II.8.14.3}
Đồng cấu \sref{II.8.14.2.4} là song ánh;
đồng cấu \sref{II.8.14.2.3} cũng vậy khi $\sh{S}$-môđun phân bậc $\sh{M}$ có một trình bày hữu hạn \sref{II.3.1.1}.
\end{proposition}

\begin{proof}
Vấn đề hiển nhiên mang tính địa phương trên $X$ và $Y$;
do đó, ta có thể giả sử $Y=\Spec(A)$ là affine, với $\sh{S}=\widetilde{S}$, $\sh{M}=\widetilde{M}$ và $\sh{N}=\widetilde{N}$, trong đó $S$ là một $A$-đại số phân bậc dương, còn $M$ và $N$ là các $S$-môđun phân bậc.
Nếu $f$ là một phần tử thuần nhất của $S_+$, thì các đồng cấu \sref{II.8.14.2.1} và \sref{II.8.14.2.2}, khi hạn chế lên tập mở affine $D_+(f)$, tương ứng với các đồng cấu chính tắc \sref{II.2.5.11.1} và \sref{II.2.5.12.1}:
\begin{align*}
    M(m)_{(f)}\otimes_{S_{(f)}}N(n)_{(f)}
    &\to
    (M\otimes_S N)(m+n)_{(f)}
\\  (\Hom_S(M,N))(n-m)_{(f)}
    &\to
    \Hom_{S_{(f)}}(M(m)_{(f)},N(n)_{(f)}).
\end{align*}

Đối chiếu các định nghĩa của những đồng cấu này, ta thấy (có xét đến \sref{II.8.2.9.1}) rằng hạn chế của \sref{II.8.14.2.4} lên $D_+(f)$ tương ứng với đồng cấu chính tắc
\[
    M_f\otimes_{S_f}N_f \to (M\otimes_S N)_f
\]
được định nghĩa trong \sref[0]{0.1.3.4}, và ta biết đồng cấu sau là một đẳng cấu.
Tương tự, hạn chế của \sref{II.8.14.2.3} lên $D_+(f)$ tương ứng với đồng cấu chính tắc \sref[0]{0.1.3.5}
\[
    (\Hom_S(M,N))_f \to \Hom_{S_f}(M_f,N_f)
\]
ở đây ta lưu ý rằng, vì $M$ hữu hạn sinh, môđun $\Hom_S(M,N)$, tức tổng trực tiếp của các nhóm con gồm những đồng cấu \emph{thuần nhất} của các $S$-môđun \sref{II.2.1.2}, trùng với tập hợp \emph{tất cả} các đồng cấu $M\to N$ của các $S$-môđun.
Theo \sref[0]{0.1.3.5}, giả thiết rằng $M$ có một trình bày hữu hạn kéo theo đồng cấu chính tắc đang xét quả là một đẳng cấu.
\end{proof}

\begin{proposition}[8.14.4]
\label{II.8.14.4}
Nếu $U$ là một tập mở quasi-compact của $X$, thì tồn tại một số nguyên $d$ sao cho, với mọi số nguyên $n$ là bội của $d$, $\sh{O}_X(n)|U$ khả nghịch, với nghịch đảo là $\sh{O}_X(-n)|U$.
\end{proposition}

\begin{proof}
Vì $q(U)$ quasi-compact, nó được phủ bởi hữu hạn tập mở affine $V_i$, nên mỗi $x\in U$ nằm trong một tập mở affine có dạng $D_+(f)$, trong đó $f$ là một phần tử thuần nhất bậc $>0$ của một trong các vành $\Gamma(V_i,\sh{S})$.
Vì $U$ quasi-compact, ta có thể phủ nó bởi hữu hạn tập mở như vậy $D_+(f_j)$;
gọi $d$ là một bội chung của các bậc của $f_j$.
Theo \sref{II.2.5.17}, số $d$ này có tính chất cần tìm.
\end{proof}

\begin{env}[8.14.5]
\label{II.8.14.5}
Dưới các giả thiết và với ký hiệu của \sref{II.8.14.1}, trong \sref{II.3.3.2} ta đã định nghĩa các đồng cấu chính tắc của các $\sh{O}_Y$-môđun
\[
\label{II.8.14.5.1}
    \alpha_n: \sh{M}_n \to q_*(\shProj_0(\sh{M}(n)))
    \qquad (n\in\bb{Z}).
\tag{8.14.5.1}
\]

Mở rộng ký hiệu của \sref{II.3.3.1}, với mọi \emph{$\sh{S}_X$-môđun phân bậc $\sh{F}$}, ta đặt
\[
\label{II.8.14.5.2}
    \boldsymbol{\Gamma}_*(\sh{F}) = \bigoplus_{n\in\bb{Z}}q_*(\sh{F}_n).
\tag{8.14.5.2}
\]

Đặc biệt, $\boldsymbol{\Gamma}_*(\sh{S}_X)=\bigoplus_{n\in\bb{Z}}q_*(\sh{O}_X(n))$ là $\sh{O}_Y$-đại số phân bậc được ký hiệu bởi $\boldsymbol{\Gamma}_*(\sh{O}_X)$ trong \sref{II.3.3.1.2};
rõ ràng $\boldsymbol{\Gamma}_*(\sh{F})$ là một \emph{$\boldsymbol{\Gamma}_*(\sh{S}_X)$-môđun phân bậc} \sref[0]{0.4.2.2}.
Khi
\oldpage[II]{199}
ta lấy $\sh{M}=\sh{S}$ trong các đồng cấu \sref{II.8.14.5.1}, ta thu được đồng cấu của các $\sh{O}_Y$-đại số phân bậc
\[
\label{II.8.14.5.3}
    \alpha:\sh{S}\to\boldsymbol{\Gamma}_*(\sh{S}_X)
\tag{8.14.5.3}
\]
đã được định nghĩa trong \sref{II.3.3.2}, và đồng cấu này làm cho $\boldsymbol{\Gamma}_*(\sh{F})$ thành một \emph{$\sh{S}$-môđun phân bậc};
khi đó các đồng cấu \sref{II.8.14.5.1} xác định một đồng cấu (bậc $0$) của các \emph{$\sh{S}$-môđun phân bậc}
\[
\label{II.8.14.5.4}
    \alpha:\sh{M}\to\boldsymbol{\Gamma}_*(\shProj(\sh{M})).
\tag{8.14.5.4}
\]
\end{env}

\begin{env}[8.14.6]
\label{II.8.14.6}
Nói chung, với một $\sh{S}_X$-môđun phân bậc quasi-coherent $\sh{F}$, $\sh{S}$-môđun phân bậc $\boldsymbol{\Gamma}_*(\sh{F})$ không nhất thiết quasi-coherent.
Xét một tập mở $X'$ của $X$ sao cho hạn chế $q':X'\to Y$ của $q$ lên $X'$ là một cấu xạ \emph{quasi-compact}.
Vì $q'$ còn là cấu xạ tách, nên $q'_*(\sh{F}')$ là một $\sh{O}_Y$-môđun quasi-coherent với mọi $\sh{O}_{X'}$-môđun quasi-coherent $\sh{F}'$ \sref[I]{I.9.2.2}[b].
Ta đặt
\[
\label{II.8.14.6.1}
    \sh{S}_{X'} = \sh{S}_X|X' = \bigoplus_{n\in\bb{Z}}\sh{O}_X(n)|X'
\tag{8.14.6.1}
\]
và, với mọi $\sh{S}_{X'}$-môđun phân bậc $\sh{F}'$,
\[
\label{II.8.14.6.2}
    \boldsymbol{\Gamma}'_*(\sh{F}') = \bigoplus_{n\in\bb{Z}}q'_*(\sh{F}'_n).
\tag{8.14.6.2}
\]

Nhận xét trên cho thấy rằng nếu $\sh{F}'$ là một $\sh{S}_{X'}$-môđun quasi-coherent, thì $\boldsymbol{\Gamma}'_*(\sh{F}')$ là một $\sh{S}$-môđun phân bậc \emph{quasi-coherent} \sref[I]{I.9.6.1}.

Ta cũng lưu ý rằng đơn ánh chính tắc $j:X'\to X$ là \emph{quasi-compact}, vì $q'=q\circ j$ quasi-compact và $q$ tách \sref[I]{I.6.6.4}[v].
Khi đó $\sh{F}=j_*(\sh{F}')$ là một $\sh{S}_X$-môđun phân bậc quasi-coherent với mọi $\sh{S}_{X'}$-môđun phân bậc quasi-coherent $\sh{F}'$, và từ các định nghĩa trên suy ra
\[
\label{II.8.14.6.3}
    \boldsymbol{\Gamma}'_*(\sh{F}') = \boldsymbol{\Gamma}_*(\sh{F}).
\tag{8.14.6.3}
\]

Với cùng các giả thiết trên $X'$, đối với mỗi $\sh{S}$-môđun phân bậc quasi-coherent $\sh{M}$, ta đặt
\[
\label{II.8.14.6.4}
    \shProj'(\sh{M}) = \shProj(\sh{M})|X'
\tag{8.14.6.4}
\]
và đây là một $\sh{S}_{X'}$-môđun phân bậc quasi-coherent.
Đồng cấu chính tắc
\[
    \shProj(\sh{M}) \to j_*(\shProj'(\sh{M}))
\]
\sref[0]{0.4.4.3} do đó cho một đồng cấu chính tắc $\boldsymbol{\Gamma}_*(\shProj(\sh{M}))\to\boldsymbol{\Gamma}'_*(\shProj'(\sh{M}))$ của các $\sh{S}$-môđun phân bậc, và bằng cách hợp thành với \sref{II.8.14.5.4}, ta thu được một đồng cấu chính tắc, có tính hàm tử (bậc $0$), của các $\sh{S}$-môđun phân bậc quasi-coherent
\[
\label{II.8.14.6.5}
    \alpha': \sh{M} \to \boldsymbol{\Gamma}'_*(\shProj'(\sh{M})).
\tag{8.14.6.5}
\]
\end{env}

\begin{env}[8.14.7]
\label{II.8.14.7}
Giữ các giả thiết trên $X'$ của \sref{II.8.14.6}, cho $\sh{F}'$ là một \emph{$\sh{S}_{X'}$-môđun phân bậc quasi-coherent}, sao cho $\shProj'(\boldsymbol{\Gamma}'_*(\sh{F}'))$ cũng là một $\sh{S}_{X'}$-môđun phân bậc \emph{quasi-coherent}.
\oldpage[II]{200}
Ta sẽ định nghĩa một đồng cấu chính tắc, có tính hàm tử (bậc $0$), của các $\sh{S}_{X'}$-môđun phân bậc
\[
\label{II.8.14.7.1}
    \beta': \shProj'(\boldsymbol{\Gamma}'_*(\sh{F}')) \to \sh{F}'.
\tag{8.14.7.1}
\]

Trước hết, giả sử $Y=\Spec(A)$ là affine và $\sh{S}=\widetilde{S}$, trong đó $S$ là một $A$-đại số phân bậc dương;
khi đó $\boldsymbol{\Gamma}'_*(\sh{F}')=\widetilde{M}$, trong đó $M=\bigoplus_{n\in\bb{Z}}\Gamma(X',\sh{F}'_n)$ là một $S$-môđun phân bậc.
Cho $f\in S_d$ sao cho $D_+(f)\subset X'$;
theo định nghĩa \sref{II.2.6.2}, $\alpha_d(f)$ khi hạn chế lên $D_+(f)$ là tiết diện của $\sh{O}_X(d)$ trên $D_+(f)$ tương ứng với phần tử $f/1$ của $(S(d))_{(f)}$, và vì vậy khả nghịch;
do đó $\alpha_{nd}(f^n)$ cũng khả nghịch với mọi $n>0$.
Từ đây, ta kết luận ngay rằng đã định nghĩa được một $S_f$-đồng cấu (bậc $0$) của các môđun phân bậc $\beta_f:M_f\to\Gamma(D_+(f),\sh{F}')$ bằng cách gửi mỗi phần tử $z/f^n\in M_f$ (trong đó $z\in M$) đến tiết diện $(z|D_+(f))(\alpha_{nd}(f^n)|D_+(f))^{-1}$ của $\sh{F}'$ trên $D_+(f)$.
Hơn nữa, ta có một biểu đồ giao hoán tương ứng với \sref{II.2.6.4.1}, qua đó xác định được $\beta'$ trong trường hợp này.
Để chuyển sang trường hợp tổng quát, ta phải xét một $A$-đại số $A'$, $A'$-đại số phân bậc $S'=S\otimes_A A'$, và dùng biểu đồ giao hoán tương tự \sref{II.2.8.13.2};
ta để lại các chi tiết cho độc giả.
\end{env}

\begin{proposition}[8.14.8]
\label{II.8.14.8}
Nếu $X'$ là một tập mở của $X=\Proj(\sh{S})$ sao cho $q':X'\to Y$ quasi-compact, thì đồng cấu $\beta'$ được định nghĩa trong \sref{II.8.14.7} là song ánh.
\end{proposition}

\begin{proof}
Rõ ràng ta có thể giới hạn vào trường hợp $Y$ affine, và khi đó mọi việc quy về chứng minh (với ký hiệu của \sref{II.8.14.7}) rằng đồng cấu $\beta_f:M_f\to\Gamma(D_+(f),\sh{F}')$ là một đẳng cấu.
Nhưng thay $f$ bởi một lũy thừa của nó không làm thay đổi $D_+(f)$ cũng như $\beta_f$;
vì $X'$ \emph{quasi-compact} theo giả thiết, nên theo \sref{II.8.14.4}, ta luôn có thể giả sử bó $\sh{O}_X(d)$ \emph{khả nghịch}.
Vì $X'$ là một lược đồ (do $q'$ tách), mệnh đề khi đó chính là \sref[I]{I.9.3.1}.
\end{proof}

\begin{corollary}[8.14.9]
\label{II.8.14.9}
Dưới các giả thiết của \sref{II.8.14.8}, mọi $\sh{S}_{X'}$-môđun phân bậc quasi-coherent đều đẳng cấu với một $\sh{S}_{X'}$-môđun phân bậc có dạng $\shProj'(\sh{M})$, trong đó $\sh{M}$ là một $\sh{S}$-môđun phân bậc quasi-coherent.
Hơn nữa, nếu $\sh{F}'$ hữu hạn sinh và nếu ta giả sử $Y$ là một lược đồ quasi-compact hoặc một tiền lược đồ có không gian nền Noether, thì ta có thể giả sử $\sh{M}$ hữu hạn sinh.
\end{corollary}

\begin{proof}
Chứng minh dựa trên \sref{II.8.14.8} hoàn toàn theo cùng lập luận như chứng minh \sref{II.3.4.5} dựa trên \sref{II.3.4.4}; ta để lại các chi tiết cho độc giả.
\end{proof}

\begin{proposition}[8.14.10]
\label{II.8.14.10}
Dưới các giả thiết của \sref{II.8.14.7}, cho $\sh{M}$ là một $\sh{S}$-môđun phân bậc quasi-coherent và $\sh{F}'$ là một $\sh{S}_{X'}$-môđun phân bậc quasi-coherent;
các đồng cấu hợp thành
\[
\label{II.8.14.10.1}
    \shProj'(\sh{M})
    \xrightarrow{\shProj'(\alpha')}
    \shProj'(\boldsymbol{\Gamma}'_*(\shProj'(\sh{M})))
    \xrightarrow{\beta'}
    \shProj'(\sh{M})
\tag{8.14.10.1}
\]
\[
\label{II.8.14.10.2}
    \boldsymbol{\Gamma}'_*(\sh{F}')
    \xrightarrow{\alpha'}
    \boldsymbol{\Gamma}'_*(\shProj'(\boldsymbol{\Gamma}'_*(\sh{F}')))
    \xrightarrow{\boldsymbol{\Gamma}'_*(\beta')}
    \boldsymbol{\Gamma}'_*(\sh{F}')
\tag{8.14.10.2}
\]
là các đẳng cấu đồng nhất.
\end{proposition}

\begin{proof}
Vấn đề mang tính địa phương trên $Y$, và chứng minh được thực hiện như trong \sref{II.2.6.5};
ta để lại các chi tiết cho độc giả.
\end{proof}

\begin{remark}[8.14.11]
\label{II.8.14.11}
Trong chương~III \sref[III]{III.2.3.1}, ta sẽ thấy rằng, khi $Y$ \emph{Noether địa phương} và $\sh{S}$ là một $\sh{O}_Y$-đại số phân bậc quasi-coherent \emph{hữu hạn sinh} (trong trường hợp đó
\oldpage[II]{201}
ta có thể lấy $X'=X$), khi đó đồng cấu $\alpha$ \sref{II.8.14.5.4} là \textbf{(TN)}-\emph{song ánh} đối với mọi $\sh{S}$-môđun phân bậc quasi-coherent $\sh{M}$ thỏa mãn điều kiện~\textbf{(TF)}.
\end{remark}

\begin{remark}[8.14.12]
\label{II.8.14.12}
Tình huống được mô tả trong \sref{II.8.14.4} là một trường hợp riêng của tình huống sau đây.
Cho $X$ là một không gian vành hóa, và $\sh{S}$ là một $\sh{O}_X$-đại số phân bậc (theo cả bậc dương lẫn bậc âm);
giả sử tồn tại một số nguyên $d>0$ sao cho $\sh{S}_d$ và $\sh{S}_{-d}$ đều \emph{khả nghịch}, và đồng cấu chính tắc
\[
\label{II.8.14.12.1}
    \sh{S}_d\otimes_{\sh{O}_X}\sh{S}_{-d} \to \sh{O}_X
\tag{8.14.12.1}
\]
là một \emph{đẳng cấu} (do đó $\sh{S}_{-d}$ được đồng nhất với $\sh{S}_d^{-1}$).
Khi ấy ta nói $\sh{O}_X$-đại số phân bậc $\sh{S}$ là \emph{tuần hoàn}, \emph{với chu kỳ $d$}.
Thuật ngữ này xuất phát từ tính chất sau:
\emph{dưới các giả thiết trên, với mọi $\sh{S}$-môđun phân bậc $\sh{F}$, đồng cấu chính tắc}
\[
\label{II.8.14.12.2}
    \sh{S}_d\otimes_{\sh{O}_X}\sh{F}_n \to \sh{F}_{n+d}
\tag{8.14.12.2}
\]
\emph{là một đẳng cấu với mọi $n\in\bb{Z}$.}
Thật vậy, vấn đề có tính địa phương trên $X$, nên ta có thể giả sử $\sh{S}_d$ có một tiết diện \emph{khả nghịch} $s$ trên $X$, với nghịch đảo $s'$ là một tiết diện của $\sh{S}_{-d}$.
Khi đó, đồng cấu $\sh{F}_{n+d}\to\sh{S}_d\otimes\sh{F}_n$ gửi mỗi tiết diện $z\in\Gamma(U,\sh{F}_{n+d})$ tới tiết diện $(s|U)\otimes(s'|U)z$ của $\sh{S}_d\otimes\sh{F}_n$ trên $U$ là nghịch đảo của \sref{II.8.14.12.2}, và điều đó chứng minh khẳng định.
Từ đây, với mọi $k\in\bb{Z}$, ta thu được một đẳng cấu chính tắc
\[
    (\sh{S}_d)^{\otimes k}\otimes\sh{F}_n \simto \sh{F}_{n+kd}.
\]
Do đó, \emph{việc cho một $\sh{S}$-môđun phân bậc $\sh{F}$ tương đương với việc cho các $\sh{S}_0$-môđun $\sh{F}_i$ ($0\leq i\leq d-1$) cùng các đồng cấu chính tắc}
\[
    \sh{S}_i\otimes\sh{F}_j \to \sh{F}_{i+j}
    \qquad
    \mbox{với $0\leq i,j\leq d-1$}
\]
(quy ước $\sh{F}_{i+j}=\sh{S}_d\otimes_{\sh{S}_0}\sh{F}_{i+j-d}$ khi $i+j\geq d$).
Dĩ nhiên, để các đồng cấu này xác định tốt một cấu trúc $\sh{S}$-môđun trên tổng trực tiếp của các $(\sh{S}_d)^{\otimes k}\otimes\sh{F}_i$ ($k\in\bb{Z}$, $0\leq i\leq d-1$), chúng phải thỏa mãn một số điều kiện kết hợp mà ta sẽ không trình bày tường minh.

Trong trường hợp $d=1$ (chính là trường hợp được xét trong \sref{subsection:II.3.3}), ta có thể nói rằng phạm trù các $\sh{S}$-môđun phân bậc (tương ứng, các $\sh{S}$-môđun phân bậc quasi-coherent nếu $X$ là một tiền lược đồ và $\sh{S}$ là quasi-coherent) \emph{tương đương} với phạm trù các $\sh{S}_0$-môđun tùy ý (tương ứng, các $\sh{S}_0$-môđun quasi-coherent);
theo nghĩa đó, ta có thể xem các kết quả của mục này là sự tổng quát hóa các kết quả của §\hyperref[section:II.3]{3}.
Hơn nữa, dưới các điều kiện hữu hạn thích hợp, các kết quả của mục này (cùng với \sref{II.8.14.11}) theo một nghĩa nào đó quy việc nghiên cứu các đại số phân bậc quasi-coherent trên một tiền lược đồ, và các môđun phân bậc ``modulo \textbf{(TN)}'' trên những đại số ấy, về việc nghiên cứu trường hợp riêng trong đó các đại số đang xét là \emph{tuần hoàn} (và khi đó điều kiện~\textbf{(TN)} đối với $\sh{M}$ \sref{II.3.4.2} kéo theo $\sh{M}=0$).
\end{remark}

\begin{remark}[8.14.13]
\label{II.8.14.13}
Dưới các giả thiết của \sref{II.8.14.1}, cho $d$ là một số nguyên $>0$;
ta đã định nghĩa một $Y$-đẳng cấu chính tắc $h$ từ $X$ lên $X^{(d)}=\Proj(\sh{S}^{(d)})$ \sref{II.3.1.8}.
Với mọi
\oldpage[II]{202}
$\sh{S}$-môđun phân bậc quasi-coherent $\sh{M}$ và mọi số nguyên $k$ sao cho $0\leq k\leq d-1$, ta cũng có (với ký hiệu của \sref{II.3.1.1}) một $h$-đẳng cấu chính tắc
\[
\label{II.8.14.13.1}
    (\shProj(\sh{M}))^{(d,k)} \xleftarrow{\sim} \shProj(\sh{M}^{(d,k)}).
\tag{8.14.13.1}
\]

Trước hết giả sử $Y=\Spec(A)$ là affine, $\sh{S}=\widetilde{S}$ và $\sh{M}=\widetilde{M}$, trong đó $S$ là một $A$-đại số phân bậc dương và $M$ là một $S$-môđun phân bậc.
Với mọi $f\in S_e$ ($e>0$), ta biết rằng $h$ gửi $D_+(f)$ lên $D_+(f^d)$ và tương ứng với đẳng cấu chính tắc $S_{(f^d)}\to S_{(f)}$ \sref{II.2.2.2}.
Khi đó, hạn chế của \sref{II.8.14.13.1} lên $D_+(f^d)$ tương ứng với di-đẳng cấu chính tắc $M_{f^d}\to M_f$, được hạn chế trên các phần tử của $M_{f^d}$ có bậc đồng dư với $k$ modulo $d$.
Ta dành cho người đọc việc chứng minh rằng các đẳng cấu này tương thích với phép chuyển từ $f$ sang một bội thuần nhất $fg$, rồi chứng minh tính tương thích tương tự với phép chuyển từ $S$ sang một $A'$-đại số phân bậc $S'=S\otimes_A A'$, trong đó $A'$ là một $A$-đại số.
Đặc biệt, điều này cho ta một $h$-đẳng cấu
\[
\label{II.8.14.13.2}
    (\sh{S}^{(d)})_{X^{(d)}} \simto (\sh{S}_X)^{(d)}
\tag{8.14.13.2}
\]
tôn trọng các cấu trúc nhân ở cả nguồn lẫn đích và, nhờ \sref{II.8.14.13.1}, biến thành một $h$-di-đẳng cấu giữa một $(\sh{S}^{(d)})_{X^{(d)}}$-môđun phân bậc và một $(\sh{S}_X)^{(d)}$-môđun phân bậc.
Tương tự, ta có một $h$-đẳng cấu
\[
\label{II.8.14.13.3}
    \shProj_0(\sh{S}^{(d,k)}(n)) \simto \sh{O}_X(nd+k),
\tag{8.14.13.3}
\]
bổ sung cho kết quả của \sref{II.3.2.9}[ii].

Đẳng cấu trong \sref{II.8.14.13.1} lập tức cảm sinh một đẳng cấu các $\sh{S}^{(d)}$-môđun phân bậc
\[
\label{II.8.14.13.4}
    \boldsymbol{\Gamma}_*^{(d)}(\shProj(\sh{M}^{(d,k)})) \simto \boldsymbol{\Gamma}_*((\shProj(\sh{M}))^{(d,k)})
\tag{8.14.13.4}
\]
trong đó $\boldsymbol{\Gamma}_*^{(d)}$ tương ứng với cấu xạ cấu trúc $q^{(d)}:X^{(d)}\to Y$;
ta kiểm tra ngay được rằng đồng cấu chính tắc $\alpha$ \sref{II.8.14.5.4}, cùng đồng cấu tương tự $\alpha^{(d)}$ đối với $X^{(d)}$, làm cho biểu đồ sau giao hoán:
\[
\label{II.8.14.13.5}
    \xymatrix{
        &\sh{M}^{(d,k)} \ar[dl]_{\alpha^{(d)}} \ar[dr]^{\alpha}&
    \\  \boldsymbol{\Gamma}_*^{(d)}(\shProj(\sh{M}^{(d,k)})) \ar[rr]^{\sim}
        && \boldsymbol{\Gamma}_*((\shProj(\sh{M}))^{(d,k)})
    }
\tag{8.14.13.5}
\]
việc kiểm tra được thực hiện bằng cách giả sử $Y$ affine rồi tính các hạn chế trên các tập mở $D_+(f^d)$ và $D_+(f)$ của ảnh qua $\alpha^{(d)}$ và $\alpha$ của cùng một phần tử thuộc $M^{(d,k)}$ (với cùng ký hiệu như trên).
\end{remark}

\begin{proposition}[8.14.14]
\label{II.8.14.14}
Cho $Y$ là một tiền lược đồ quasi-compact, $\sh{S}$ là một $\sh{O}_Y$-đại số phân bậc quasi-coherent hữu hạn sinh, và $\sh{M}$ là một $\sh{S}$-môđun phân bậc quasi-coherent thỏa mãn điều kiện~\textbf{(TF)};
đặt $X=\Proj(\sh{S})$.
Khi đó $\sh{S}_X$ là một $\sh{O}_X$-đại số phân bậc tuần hoàn \sref{II.8.14.12}, và tồn tại một
\oldpage[II]{203}
chu kỳ $d$ của $\sh{S}_X$ sao cho các $(\shProj(\sh{M}))^{(d,k)}$ ($0\leq k\leq d-1$) là các $(\sh{S}_X)^{(d)}$-môđun kiểu hữu hạn.
\end{proposition}

\begin{proof}
Thật vậy, \sref{II.3.1.10} chứng minh rằng tồn tại một $d$ sao cho $\sh{S}^{(d)}$ được sinh bởi $\sh{S}_d=(\sh{S}^{(d)})_1$, mà $\sh{S}_d$ là một $\sh{S}_0$-môđun kiểu hữu hạn.
Để chứng minh khẳng định thứ nhất, nhờ \sref{II.8.14.13.2}, ta có thể quy về trường hợp $d=1$, và khi ấy mệnh đề suy ra từ \sref{II.3.2.7}.
Hơn nữa, khi lưu ý đến \sref{II.8.14.13.1}, khẳng định thứ hai là hệ quả của \sref{II.2.1.6}[iii] và \sref{II.3.4.3}.
\end{proof}
\clearpage
\phantomsection
\addcontentsline{toc}{section}{Tài liệu tham khảo}
\renewcommand{\refname}{Tài liệu tham khảo}
\noindent\textit{Ghi chú biên tập. Để có thể dùng độc lập, các mục [1], [9],
[11] và [13], được viện dẫn trong tập này, được chép lại ở đây từ thư mục
tài liệu tham khảo của loạt trước. Bản nguồn EGA~II chỉ in phần tiếp theo
[23]--[26] trên trang in~205.}
\par\medskip

\begin{thebibliography}{26}

\bibitem[1]{I-1}
H.~Cartan và C.~Chevalley,
\emph{Séminaire de l'École Normale Supérieure, 8th year (1955--56):
Géométrie algébrique}.

\bibitem[9]{I-9}
M.~Nagata,
\emph{A general theory of algebraic geometry over Dedekind domains},
Amer. Math. J. \textbf{78} (1956), 78--116, và \textbf{80} (1958),
382--420.

\bibitem[11]{I-11}
P.~Samuel,
\emph{Commutative algebra} (ghi chép bởi D.~Herzig), Cornell University,
1953.

\bibitem[13]{I-13}
P.~Samuel và O.~Zariski,
\emph{Commutative algebra}, 2 tập, Van Nostrand, New York, 1958--60.

\oldpage[II]{205}
\bibitem[23]{II-23}
C.~Chevalley,
\emph{Introduction to the theory of algebraic functions of one variable}.
Mathematical Surveys, số~6, American Mathematical Society, New York, 1951.

\bibitem[24]{II-24}
P.~Jaffard,
\emph{Les systèmes d'idéaux}, Paris (Dunod), 1960.

\bibitem[25]{II-25}
M.~Nagata,
\emph{On the derived normal rings of Noetherian integral domains},
Mem. Coll. Sci. Kyoto, Ser.~A, tập~XXIX (1955), tr.~293--303.

\bibitem[26]{II-26}
M.~Nagata,
\emph{Existence theorems for non projective complete algebraic varieties},
Ill. J. Math., tập~II (1958), tr.~490--498.

\end{thebibliography}
\clearpage
\oldpage[II]{206}
\clearpage
\oldpage[II]{207}
\phantomsection
\addcontentsline{toc}{section}{Chỉ mục ký hiệu}
\section*{Chỉ mục ký hiệu}

\begingroup
\small
\setlength{\parindent}{0pt}
\setlength{\parskip}{2pt}

$\sh{A}(X),\sh{A}(\sh{F}),\sh{A}(\sh{B})$
($X$ là một $S$-tiền lược đồ, $\sh{F}$ là một $\sh{O}_X$-môđun,
$\sh{B}$ là một $\sh{O}_X$-đại số): \hyperref[II.1.1.1]{1.1.1}.

$\sh{A}(h)$ ($h$ là một $S$-cấu xạ): \hyperref[II.1.1.2]{1.1.2}.

$\Spec(\sh{B})$ ($\sh{B}$ là một $\sh{O}_S$-đại số): \hyperref[II.1.3.1]{1.3.1}.

$\widetilde{\sh{M}}$ ($\sh{M}$ là một $\sh{B}$-môđun,
$\sh{B}$ là một $\sh{O}_S$-đại số): \hyperref[II.1.4.3]{1.4.3}.

$\bb{T}(E),\bb{S}(E),\bb{S}_A(E)$ ($E$ là một $A$-môđun): \hyperref[II.1.7.1]{1.7.1}.

$\bb{S}(\sh{E}),\bb{S}_{\sh{A}}(\sh{E})$
($\sh{E}$ là một $\sh{A}$-môđun): \hyperref[II.1.7.4]{1.7.4}.

$\bb{V}(\sh{E})$ ($\sh{E}$ là một $\sh{O}_S$-môđun quasi-coherent): \hyperref[II.1.7.8]{1.7.8}.

$S_n,S_+,M_n,S^{(d)},M^{(d,k)},M^{(d)}$
($S$ là một vành phân bậc, $M$ là một $S$-môđun phân bậc): \hyperref[II.2.1.1]{2.1.1}.

$M(n),S(n)$ ($S$ là một vành phân bậc, $M$ là một $S$-môđun phân bậc): \hyperref[II.2.1.1]{2.1.1}.

$\Hom_S(M,N)$ ($S$ là một vành phân bậc, $M,N$ là các $S$-môđun phân bậc): \hyperref[II.2.1.2]{2.1.2}.

$\mathfrak N_+,\mathfrak r_+(\mathfrak J)$
($\mathfrak J$ là một iđêan của $S_+$): \hyperref[II.2.1.10]{2.1.10}.

$S_{(f)},M_{(f)}$ ($f$ là một phần tử thuần nhất của $S_+$): \hyperref[II.2.2.1]{2.2.1}.

$S_{(T)},M_{(T)},S_{(\mathfrak p)},M_{(\mathfrak p)}$
($T$ là một tập con nhân của $S_+$ gồm các phần tử thuần nhất,
$\mathfrak p$ là một iđêan nguyên tố phân bậc của $S_+$): \hyperref[II.2.2.7]{2.2.7}.

$\Proj(S)$ ($S$ là một vành phân bậc): \hyperref[II.2.3.1]{2.3.1}.

$V_+(E)$ ($E$ là một tập con của $S_+$): \hyperref[II.2.3.2]{2.3.2}.

$D_+(f)$ ($f$ là một phần tử của $S$): \hyperref[II.2.3.3]{2.3.3}.

$\mathfrak j_+(Y)$ ($Y$ là một tập con của $\Proj(S)$): \hyperref[II.2.3.10]{2.3.10}.

$\widetilde M$ ($M$ là một $S$-môđun phân bậc): \hyperref[II.2.5.2]{2.5.2}.

$\widetilde u$ ($u$ là một đồng cấu bậc không): \hyperref[II.2.5.4]{2.5.4}.

$\sh{O}_X(n),\sh{F}(n)$
($X=\Proj(S)$, $\sh{F}$ là một $\sh{O}_X$-môđun): \hyperref[II.2.5.10]{2.5.10}.

$\lambda$ (đồng cấu chính tắc): \hyperref[II.2.5.11]{2.5.11}.

$\mu$ (đồng cấu chính tắc): \hyperref[II.2.5.12]{2.5.12}.

$X^{(d)}$ ($X=\Proj(S)$, $d$ là một số nguyên $>0$): \hyperref[II.2.5.16]{2.5.16}.

$\Gamma_*(\sh{F})$ ($\sh{F}$ là một $\sh{O}_X$-môđun,
$X=\Proj(S)$): \hyperref[II.2.6.1]{2.6.1}.

$\alpha_n,\alpha,\alpha_{\sh{M}}$ (các đồng cấu chính tắc): \hyperref[II.2.6.2]{2.6.2}.

$\beta,\beta_{\sh{F}}$ (các đồng cấu chính tắc): \hyperref[II.2.6.4]{2.6.4}.

$G(\vphi),{}^a\vphi$ (theo lạm dụng ngôn ngữ;
$\vphi$ là một đồng cấu các vành phân bậc): \hyperref[II.2.8.1]{2.8.1}.

$\Proj(\vphi)$ ($\vphi$ là một đồng cấu các vành phân bậc): \hyperref[II.2.8.2]{2.8.2}.

$\sh{S}^{(d)},\sh{M}^{(d,k)},\sh{M}^{(d)},\sh{M}(n)$
($\sh{S}$ là một $\sh{O}_Y$-đại số phân bậc,
$\sh{M}$ là một $\sh{S}$-môđun phân bậc): \hyperref[II.3.1.1]{3.1.1}.

$\Proj(\sh{S})$ ($\sh{S}$ là một $\sh{O}_Y$-đại số phân bậc quasi-coherent): \hyperref[II.3.1.3]{3.1.3}.

$X_f$ ($X=\Proj(\sh{S})$, $f$ là một tiết diện của $\sh{S}_d$ trên $Y$): \hyperref[II.3.1.4]{3.1.4}.

$\widetilde{\sh{M}}$ ($\sh{M}$ là một $\sh{S}$-môđun phân bậc): \hyperref[II.3.2.2]{3.2.2}.

$\sh{O}_X(n),\sh{F}(n)$
($X=\Proj(\sh{S})$, $\sh{F}$ là một $\sh{O}_X$-môđun): \hyperref[II.3.2.5]{3.2.5}.

$\lambda,\mu$ (các đồng cấu chính tắc): \hyperref[II.3.2.6]{3.2.6}.

$\Gamma_*(\sh{F})$ ($\sh{F}$ là một $\sh{O}_X$-môđun,
$X=\Proj(\sh{S})$): \hyperref[II.3.3.1]{3.3.1}.

$\alpha_n,\alpha,\alpha_{\sh{M}}$ (các đồng cấu chính tắc): \hyperref[II.3.3.2]{3.3.2}.
\oldpage[II]{208}
$\beta,\beta_{\sh{F}}$ (các đồng cấu chính tắc): \hyperref[II.3.3.4]{3.3.4}.

$G(\vphi),\Proj(\vphi)$
($\vphi$ là một đồng cấu các $\sh{O}_Y$-đại số phân bậc): \hyperref[II.3.5.1]{3.5.1}.

$\Proj(u)$ ($u$ là một $g$-cấu xạ từ một $\sh{O}_Y$-đại số phân bậc
đến một $\sh{O}_{Y'}$-đại số phân bậc): \hyperref[II.3.5.6]{3.5.6}.

$r_{\sh{L},\psi}$ ($\sh{L}$ là một $\sh{O}_X$-môđun khả nghịch): \hyperref[II.3.7.1]{3.7.1}.

$\bb{P}(\sh{E}),\bb{P}(E),\bb{P}_Y^n,\bb{P}_A^n$: \hyperref[II.4.1.1]{4.1.1}.

$\bb{P}(u)$ ($u$ là một đồng cấu toàn ánh các $\sh{O}_Y$-môđun): \hyperref[II.4.1.2]{4.1.2}.

$\Hyp_Y(X,\sh{E})$ ($\sh{E}$ là một $\sh{O}_Y$-môđun quasi-coherent): \hyperref[II.4.2.5]{4.2.5}.

$\varsigma$ (cấu xạ Segre): \hyperref[II.4.3.1]{4.3.1}.

$\sh{F}(n)$ ($\sh{F}$ là một $\sh{O}_X$-môđun, $X$ được trang bị một
$\sh{O}_X$-môđun khả nghịch): \hyperref[II.4.5.1]{4.5.1}.

$\varepsilon$ (tự đồng cấu đồng nhất): \hyperref[II.4.5.1]{4.5.1}.

$P(u,T),\sigma_i(u)$ ($u$ là một tự đồng cấu của môđun tự do có cơ sở hữu hạn): \hyperref[II.6.4.1]{6.4.1}.

$P(u,T),\sigma_i(u),\det u$
($u$ là một tự đồng cấu của môđun hữu hạn sinh trên một vành nguyên hoặc một
vành Noether rút gọn): \hyperref[II.6.4.2]{6.4.2} và \hyperref[II.6.4.7]{6.4.7}.

$\sigma_i(u),\det u$ ($u$ là một tự đồng cấu của $\sh{A}$-môđun tự do địa phương): \hyperref[II.6.4.8]{6.4.8}.

$\sigma_i,\det$ (các đồng cấu bó): \hyperref[II.6.4.8]{6.4.8}, \hyperref[II.6.4.9]{6.4.9} và \hyperref[II.6.4.10]{6.4.10}.

$\sigma_i(u),\det u$
($u$ là một tự đồng cấu của $\sh{O}_X$-môđun hữu hạn sinh trên một tiền lược đồ
nguyên địa phương hoặc một tiền lược đồ Noether địa phương và rút gọn):
\hyperref[II.6.4.9]{6.4.9} và \hyperref[II.6.4.10]{6.4.10}.

$N_{\sh{B}/\sh{A}}(f)$
($f$ là một tiết diện của $\sh{A}$-đại số $\sh{B}$): \hyperref[II.6.5.1]{6.5.1}.

$N_{\sh{B}/\sh{A}}(\sh{L}')$
($\sh{L}'$ là một $\sh{B}$-môđun khả nghịch): \hyperref[II.6.5.2]{6.5.2}.

$N_{\sh{B}/\sh{A}}(h')$
($h'$ là một đồng cấu các $\sh{B}$-môđun khả nghịch): \hyperref[II.6.5.3]{6.5.3}.

$N_{X'/X}(\sh{L}')$
($X'$ hữu hạn trên $X$, $\sh{L}'$ là một $\sh{O}_{X'}$-môđun khả nghịch): \hyperref[II.6.5.5]{6.5.5}.

$N_{X'/X}(h')$
($h'$ là một đồng cấu các $\sh{O}_{X'}$-môđun khả nghịch): \hyperref[II.6.5.5]{6.5.5}.

$\sh{S}_X$ ($X=\Proj(\sh{S})$, $\sh{S}$ là một tổng các iđêan của
$\sh{O}_Y$): \hyperref[II.8.1.5]{8.1.5}.

$S^\geq,S^\leq,S_f^\geq,S_f^\leq,
M^\geq,M^\leq,M_f^\geq,M_f^\leq$
($S$ là một vành phân bậc, $M$ là một $S$-môđun phân bậc,
$f$ là một phần tử thuần nhất của $S_+$): \hyperref[II.8.2.1]{8.2.1}.

$\widehat S$ ($S$ là một vành phân bậc): \hyperref[II.8.2.2]{8.2.2}.

$\widehat M$ ($M$ là một $S$-môđun phân bậc): \hyperref[II.8.2.4]{8.2.4}.

$S_{[n]},S^\natural,f^\natural$
($S$ là một vành phân bậc, $f$ thuần nhất trong $S$): \hyperref[II.8.2.6]{8.2.6}.

$M_{[n]},M^\natural$ ($M$ là một $S$-môđun phân bậc): \hyperref[II.8.2.8]{8.2.8}.

$\widehat{\sh{S}}$ ($\sh{S}$ là một $\sh{O}_Y$-đại số phân bậc): \hyperref[II.8.3.1]{8.3.1}.

$\bb{G}_m$ (lược đồ nhóm): \hyperref[II.8.3.9]{8.3.9}.

$\sh{S}_X$ ($X=\Proj(\sh{S})$, $\sh{S}$ là một $\sh{O}_Y$-đại số phân bậc): \hyperref[II.8.6.1]{8.6.1}.

$\sh{S}_X,C_X,\widehat C_X,E_X,\widehat E_X$
($X=\Proj(\sh{S})$, $\sh{S}$ là một $\sh{O}_Y$-đại số phân bậc
với $\sh{S}_0=\sh{O}_Y$): \hyperref[II.8.6.1]{8.6.1}.

$\sh{S}_{[n]},\sh{S}^{\natural}$
($\sh{S}$ là một $\sh{O}_Y$-đại số phân bậc): \hyperref[II.8.7.2]{8.7.2}.

$\shProj_0(\sh{M}),\shProj(\sh{M}),\sh{M}_X,
\widehat{\sh{M}},\sh{M}^{\square}$
($\sh{M}$ là một $\sh{S}$-môđun phân bậc): \hyperref[II.8.12.1]{8.12.1}.

$\sh{M}_X^\geq$
($X=\Proj(\sh{S})$, $\sh{M}$ là một $\sh{S}$-môđun phân bậc): \hyperref[II.8.12.5]{8.12.5}.

$\sh{M}_{[n]},\sh{M}^{\natural}$
($\sh{M}$ là một $\sh{S}$-môđun phân bậc): \hyperref[II.8.12.9]{8.12.9}.

$(\widetilde{\sh{N}})^-$
($\sh{N}$ là một $\sh{S}$-môđun con của một $\sh{S}$-môđun phân bậc): \hyperref[II.8.13.1]{8.13.1}.

$\overline{\sh{N}}$
($\sh{N}$ là một $\sh{S}$-môđun con của một $\sh{S}$-môđun phân bậc): \hyperref[II.8.13.2]{8.13.2}.

$\Gamma_*(\sh{F})$
($X=\Proj(\sh{S})$, $\sh{F}$ là một $\sh{S}_X$-môđun phân bậc): \hyperref[II.8.14.5]{8.14.5}.

$\sh{S}_{X'},\Gamma'_*(\sh{F}'),\shProj'(\sh{M}),
\alpha',\beta'$
($X'$ mở trong $X=\Proj(\sh{S})$, $\sh{F}'$ là một $\sh{S}_{X'}$-môđun phân bậc,
$\sh{M}$ là một $\sh{S}$-môđun phân bậc): \hyperref[II.8.14.6]{8.14.6}.

\endgroup

\clearpage
\oldpage[II]{209}
\phantomsection
\addcontentsline{toc}{section}{Chỉ mục thuật ngữ}
\section*{Chỉ mục thuật ngữ}

\begingroup
\small
\setlength{\parindent}{0pt}
\setlength{\parskip}{2pt}
\newcommand{\egaTerm}[2]{\textbf{#1}: #2.\par}

\egaTerm{Affine (cấu xạ)}{\hyperref[II.1.6.1]{1.6.1}}
\egaTerm{Affine (tiền lược đồ) trên một tiền lược đồ}{\hyperref[II.1.2.1]{1.2.1}}
\egaTerm{Đại số phân bậc trên một vành phân bậc}{\hyperref[II.2.1.2]{2.1.2}}
\egaTerm{Đại số đối xứng của một $A$-môđun}{\hyperref[II.1.7.1]{1.7.1}}
\egaTerm{$\sh{A}$-đại số đối xứng của một $\sh{A}$-môđun}{\hyperref[II.1.7.4]{1.7.4}}
\egaTerm{$\sh{O}_Y$-đại số phân bậc rút gọn cốt yếu}{\hyperref[II.3.1.12]{3.1.12}}
\egaTerm{$\sh{O}_Y$-đại số phân bậc nguyên}{\hyperref[II.3.1.12]{3.1.12}}
\egaTerm{Ample ($\sh{O}_X$-môđun khả nghịch)}{\hyperref[II.4.5.3]{4.5.3}}
\egaTerm{Ample tương đối với $f$; $f$-ample; $Y$-ample ($\sh{O}_X$-môđun khả nghịch)}{\hyperref[II.4.6.1]{4.6.1}}
\egaTerm{Vành phân bậc nguyên cốt yếu}{\hyperref[II.2.1.11]{2.1.11}}
\egaTerm{Vành phân bậc rút gọn cốt yếu}{\hyperref[II.2.1.10]{2.1.10}}
\egaTerm{Liên kết với một $\sh{B}$-môđun (bó)}{\hyperref[II.1.4.3]{1.4.3}}
\egaTerm{Liên kết với một $S$-môđun phân bậc (bó)}{\hyperref[II.2.5.3]{2.5.3}}
\egaTerm{Liên kết với một $\sh{S}$-môđun phân bậc (bó)}{\hyperref[II.3.2.2]{3.2.2}}
\egaTerm{Liên kết với một $\sh{O}_S$-đại số ($S$-lược đồ)}{\hyperref[II.1.3.1]{1.3.1}}

\egaTerm{Điều kiện (TF); điều kiện (TN)}{\hyperref[II.2.7.2]{2.7.2}}
\egaTerm{Điều kiện (TF); điều kiện (TN)}{\hyperref[II.3.4.2]{3.4.2}}
\egaTerm{Nón affine, nón xạ ảnh xác định bởi một $\sh{O}_Y$-đại số phân bậc}{\hyperref[II.8.3.1]{8.3.1}}
\egaTerm{Nón affine bỏ đỉnh, nón xạ ảnh bỏ đỉnh xác định bởi một $\sh{O}_Y$-đại số phân bậc}{\hyperref[II.8.3.4]{8.3.4}}
\egaTerm{Nón chiếu affine, nón chiếu xạ ảnh của một tiền lược đồ $\Proj(\sh{S})$}{\hyperref[II.8.3.1]{8.3.1}}
\egaTerm{Đường cong đại số trên một trường}{\hyperref[II.7.4.2]{7.4.2}}
\egaTerm{Đường cong đại số đầy đủ}{\hyperref[II.7.4.7]{7.4.7}}
\egaTerm{Tiêu chuẩn Serre}{\hyperref[II.5.2.1]{5.2.1}}

\egaTerm{Định thức của một tự đồng cấu của môđun hữu hạn sinh trên một vành nguyên}{\hyperref[II.6.4.2]{6.4.2}}
\egaTerm{Định thức của một tự đồng cấu của môđun hữu hạn sinh trên một vành Noether rút gọn}{\hyperref[II.6.4.7]{6.4.7}}
\egaTerm{Định thức của một tự đồng cấu của $\sh{A}$-môđun tự do địa phương}{\hyperref[II.6.4.8]{6.4.8}}
\egaTerm{Định thức của một tự đồng cấu của $\sh{O}_X$-môđun hữu hạn sinh trên một tiền lược đồ $X$ nguyên địa phương}{\hyperref[II.6.4.9]{6.4.9}}
\egaTerm{Định thức của một tự đồng cấu của $\sh{O}_X$-môđun hữu hạn sinh trên một tiền lược đồ $X$ Noether địa phương và rút gọn}{\hyperref[II.6.4.10]{6.4.10}}

\egaTerm{Tiền lược đồ thu được bằng phép thổi}{\hyperref[II.8.1.3]{8.1.3}}
\egaTerm{Nguyên (cấu xạ)}{\hyperref[II.6.1.1]{6.1.1}}
\egaTerm{Nguyên (tiền lược đồ) trên một tiền lược đồ}{\hyperref[II.6.1.1]{6.1.1}}
\egaTerm{Nguyên ($\sh{O}_S$-đại số quasi-coherent)}{\hyperref[II.6.1.2]{6.1.2}}
\egaTerm{Nguyên (tiết diện của một $\sh{A}$-đại số) trên $\sh{A}$}{\hyperref[II.6.3.1]{6.3.1}}

\egaTerm{Bó cơ bản của một phân thớ xạ ảnh}{\hyperref[II.4.1.1]{4.1.1}}
\egaTerm{Bao đóng nguyên của một bó vành $\sh{A}$ trong một $\sh{A}$-đại số}{\hyperref[II.6.3.2]{6.3.2}}
\egaTerm{Bao đóng nguyên của một tiền lược đồ $X$ tương đối với một $\sh{O}_X$-đại số}{\hyperref[II.6.3.4]{6.3.4}}
\egaTerm{Bao đóng xạ ảnh của một nón affine}{\hyperref[II.8.3.1]{8.3.1}}
\oldpage[II]{210}
\egaTerm{Thớ hình học hữu tỉ của một phân thớ xạ ảnh trên một mở rộng $K$ của $\kres(y)$}{\hyperref[II.4.2.6]{4.2.6}}
\egaTerm{Phân thớ xạ ảnh xác định bởi một $\sh{O}_Y$-môđun}{\hyperref[II.4.1.1]{4.1.1}}
\egaTerm{Phân thớ vectơ xác định bởi một $\sh{O}_S$-môđun}{\hyperref[II.1.7.8]{1.7.8}}
\egaTerm{Hữu hạn (cấu xạ)}{\hyperref[II.6.1.1]{6.1.1}}
\egaTerm{Hữu hạn (tiền lược đồ) trên một tiền lược đồ}{\hyperref[II.6.1.1]{6.1.1}}
\egaTerm{Hữu hạn ($\sh{O}_S$-đại số quasi-coherent)}{\hyperref[II.6.1.2]{6.1.2}}
\egaTerm{Các hàm đối xứng sơ cấp của một tự đồng cấu của môđun tự do địa phương}{\hyperref[II.6.4.1]{6.4.1}}
\egaTerm{Các hàm đối xứng sơ cấp của một tự đồng cấu của môđun hữu hạn sinh trên một vành nguyên}{\hyperref[II.6.4.2]{6.4.2}}
\egaTerm{Các hàm đối xứng sơ cấp của một tự đồng cấu của môđun hữu hạn sinh trên một vành Noether rút gọn}{\hyperref[II.6.4.7]{6.4.7}}
\egaTerm{Các hàm đối xứng sơ cấp của một tự đồng cấu của $\sh{A}$-môđun tự do địa phương có hạng không đổi}{\hyperref[II.6.4.8]{6.4.8}}
\egaTerm{Các hàm đối xứng sơ cấp của một tự đồng cấu của $\sh{O}_X$-môđun hữu hạn sinh trên một tiền lược đồ $X$ nguyên địa phương}{\hyperref[II.6.4.9]{6.4.9}}
\egaTerm{Các hàm đối xứng sơ cấp của một tự đồng cấu của $\sh{O}_X$-môđun hữu hạn sinh trên một tiền lược đồ $X$ Noether địa phương và rút gọn}{\hyperref[II.6.4.10]{6.4.10}}
\egaTerm{Các hàm đối xứng sơ cấp của một tiết diện của một $\sh{A}$-đại số}{\hyperref[II.6.5.1]{6.5.1}}

\egaTerm{Thuần nhất hóa một $\sh{S}$-môđun con của một $\sh{S}$-môđun phân bậc}{\hyperref[II.8.13.2]{8.13.2}}
\egaTerm{Đồng cấu các vành phân bậc}{\hyperref[II.2.1.2]{2.1.2}}
\egaTerm{Đồng cấu bậc $k$ của các môđun phân bậc}{\hyperref[II.2.1.2]{2.1.2}}

\egaTerm{Iđêan của $S_+$; iđêan nguyên tố phân bậc của $S_+$ ($S$ là một vành phân bậc)}{\hyperref[II.2.1.10]{2.1.10}}

\egaTerm{Bổ đề Chow}{\hyperref[II.5.6.1]{5.6.1}}
\egaTerm{Quỹ tích ở vô cực của một nón xạ ảnh}{\hyperref[II.8.3.3]{8.3.3}}

\egaTerm{Cấu xạ chính tắc $G(\varepsilon)\to\Proj(S)$, $S=\bigoplus_{n\geq0}\Gamma(X,\sh{L}^{\otimes n})$}{\hyperref[II.4.5.1]{4.5.1}}
\egaTerm{Cấu xạ chính tắc $X\to\Spec(\Gamma(X,\sh{O}_X))$}{\hyperref[II.5.1.1]{5.1.1}}
\egaTerm{Cấu xạ Segre}{\hyperref[II.4.3.1]{4.3.1}}

\egaTerm{Căn nil của $S_+$ ($S$ là một vành phân bậc)}{\hyperref[II.2.1.10]{2.1.10}}
\egaTerm{Căn nil của $\sh{S}_+$ ($\sh{S}$ là một $\sh{O}_Y$-đại số phân bậc quasi-coherent)}{\hyperref[II.3.1.12]{3.1.12}}
\egaTerm{Chuẩn của một tiết diện của một $\sh{A}$-đại số}{\hyperref[II.6.5.1]{6.5.1}}
\egaTerm{Chuẩn của một $\sh{B}$-môđun khả nghịch}{\hyperref[II.6.5.2]{6.5.2}}
\egaTerm{Chuẩn của một tiết diện của một $\sh{B}$-môđun khả nghịch}{\hyperref[II.6.5.3]{6.5.3}}
\egaTerm{Chuẩn của một $\sh{O}_{X'}$-môđun khả nghịch}{\hyperref[II.6.5.5]{6.5.5}}

\egaTerm{Tuần hoàn (đại số phân bậc)}{\hyperref[II.8.14.12]{8.14.12}}
\egaTerm{Đa thức đặc trưng của một tự đồng cấu của môđun hữu hạn sinh trên một vành nguyên}{\hyperref[II.6.4.2]{6.4.2}}
\egaTerm{Đa thức đặc trưng của một tự đồng cấu của môđun hữu hạn sinh trên một vành Noether rút gọn}{\hyperref[II.6.4.7]{6.4.7}}
\egaTerm{Tiền lược đồ các đỉnh của một nón affine (xạ ảnh)}{\hyperref[II.8.3.3]{8.3.3}}
\egaTerm{Trình bày hữu hạn (môđun phân bậc có một trình bày như vậy)}{\hyperref[II.2.1.1]{2.1.1}}
\egaTerm{Trình bày hữu hạn ($\sh{S}$-môđun phân bậc có một trình bày như vậy)}{\hyperref[II.3.1.1]{3.1.1}}
\egaTerm{Tích tensor phân bậc của hai môđun phân bậc}{\hyperref[II.2.1.2]{2.1.2}}
\egaTerm{Xạ ảnh (cấu xạ)}{\hyperref[II.5.5.2]{5.5.2}}
\egaTerm{Xạ ảnh (tiền lược đồ) trên một tiền lược đồ}{\hyperref[II.5.5.2]{5.5.2}}
\egaTerm{Riêng (cấu xạ)}{\hyperref[II.5.4.1]{5.4.1}}
\egaTerm{Riêng (tập con)}{\hyperref[II.5.4.10]{5.4.10}}
\egaTerm{Riêng (tiền lược đồ) trên một tiền lược đồ}{\hyperref[II.5.4.1]{5.4.1}}

\egaTerm{Quasi-affine (cấu xạ)}{\hyperref[II.5.1.1]{5.1.1}}
\egaTerm{Quasi-affine (tiền lược đồ) trên một tiền lược đồ}{\hyperref[II.5.1.1]{5.1.1}}
\egaTerm{Quasi-affine (lược đồ)}{\hyperref[II.5.1.1]{5.1.1}}
\oldpage[II]{211}
\markboth{CHƯƠNG II}{CHƯƠNG II}
\egaTerm{Cấu xạ quasi-hữu hạn}{\hyperref[II.6.2.3]{6.2.3}}
\egaTerm{Tiền lược đồ quasi-hữu hạn trên một tiền lược đồ}{\hyperref[II.6.2.3]{6.2.3}}
\egaTerm{Cấu xạ quasi-xạ ảnh}{\hyperref[II.5.3.1]{5.3.1}}
\egaTerm{Tiền lược đồ quasi-xạ ảnh trên một tiền lược đồ}{\hyperref[II.5.3.1]{5.3.1}}

\egaTerm{Căn của một iđêan của $S_+$ ($S$ là vành phân bậc)}{\hyperref[II.2.1.10]{2.1.10}}
\egaTerm{Phép co rút chính tắc của một nón xạ ảnh bỏ đỉnh lên quỹ tích của nó ở vô cực}{\hyperref[II.8.3.5]{8.3.5}}

\egaTerm{Tiết diện chính tắc của $\sh{O}_X(-1)$ ($X$ là tiền lược đồ thu được bằng phép thổi)}{\hyperref[II.8.1.9]{8.1.9}}
\egaTerm{Tiết diện không của một nón affine}{\hyperref[II.8.3.3]{8.3.3}}
\egaTerm{Tiết diện không của một phân thớ vectơ}{\hyperref[II.1.7.9]{1.7.9}}
\egaTerm{Tiết diện đỉnh của một nón affine (xạ ảnh)}{\hyperref[II.8.3.3]{8.3.3}}
\egaTerm{Phổ của một $\sh{O}_S$-đại số}{\hyperref[II.1.3.1]{1.3.1}}
\egaTerm{Phổ nguyên tố thuần nhất của một vành phân bậc}{\hyperref[II.2.3.1]{2.3.1}}
\egaTerm{Phổ thuần nhất của một $\sh{O}_S$-đại số phân bậc quasi-coherent}{\hyperref[II.3.1.3]{3.1.3}}

\egaTerm{Đồng cấu (TN)-đơn ánh, (TN)-toàn ánh hoặc (TN)-song ánh của các môđun phân bậc}{\hyperref[II.2.7.2]{2.7.2}}
\egaTerm{(TN)-đẳng cấu của các môđun phân bậc}{\hyperref[II.2.7.2]{2.7.2}}
\egaTerm{Đồng cấu (TN)-đơn ánh, (TN)-toàn ánh hoặc (TN)-song ánh của các vành phân bậc}{\hyperref[II.2.9.1]{2.9.1}}
\egaTerm{(TN)-đẳng cấu của các vành phân bậc}{\hyperref[II.2.9.1]{2.9.1}}
\egaTerm{Đồng cấu (TN)-đơn ánh, (TN)-toàn ánh hoặc (TN)-song ánh của các $\sh{S}$-môđun phân bậc}{\hyperref[II.3.4.2]{3.4.2}}
\egaTerm{(TN)-đẳng cấu của các $\sh{S}$-môđun phân bậc}{\hyperref[II.3.4.2]{3.4.2}}
\egaTerm{Đồng cấu (TN)-đơn ánh, (TN)-toàn ánh hoặc (TN)-song ánh của các $\sh{O}_Y$-đại số phân bậc}{\hyperref[II.3.6.1]{3.6.1}}
\egaTerm{(TN)-đẳng cấu của các $\sh{O}_Y$-đại số phân bậc}{\hyperref[II.3.6.1]{3.6.1}}
\egaTerm{Tôpô phổ trên $\Proj(S)$}{\hyperref[II.2.3.3]{2.3.3}}
\egaTerm{Rất ample đối với $q$, rất ample đối với $Y$ ($\sh{O}_X$-môđun khả nghịch)}{\hyperref[II.4.4.2]{4.4.2}}
\egaTerm{$\sh{S}$-môđun phân bậc kiểu hữu hạn}{\hyperref[II.3.1.1]{3.1.1}}

\egaTerm{Cấu xạ đóng phổ dụng}{\hyperref[II.5.4.9]{5.4.9}}

\endgroup

\clearpage
\oldpage[II]{212}
\thispagestyle{empty}
\mbox{}

\clearpage
\oldpage[II]{213}
\markboth{MỤC LỤC}{MỤC LỤC}
\phantomsection
\addcontentsline{toc}{section}{Mục lục nguyên bản}
\section*{Mục lục}
\thispagestyle{plain}

\begingroup
\small
\renewcommand{\arraystretch}{1.04}
\begin{longtable}{@{}p{0.09\textwidth}p{0.69\textwidth}r@{}}
\textbf{Đơn vị} & \textbf{Tiêu đề} & \textbf{Trang nguyên bản}\\
\hline
\textbf{II} & \textbf{Khảo sát toàn cục sơ cấp về một số lớp cấu xạ} & \textbf{5}\\
\hyperref[section:II.1]{\textsection1} & Các cấu xạ affine & 5\\
\hyperref[subsection:II.1.1]{1.1} & $S$-tiền lược đồ và $\sh{O}_S$-đại số & 5\\
\hyperref[subsection:II.1.2]{1.2} & Các tiền lược đồ affine trên một tiền lược đồ & 6\\
\hyperref[subsection:II.1.3]{1.3} & Tiền lược đồ affine trên $S$ liên kết với một $\sh{O}_S$-đại số & 8\\
\hyperref[subsection:II.1.4]{1.4} & Các bó quasi-coherent trên một tiền lược đồ affine trên $S$ & 9\\
\hyperref[subsection:II.1.5]{1.5} & Đổi tiền lược đồ cơ sở & 12\\
\hyperref[subsection:II.1.6]{1.6} & Các cấu xạ affine & 14\\
\hyperref[subsection:II.1.7]{1.7} & Phân thớ vectơ liên kết với một bó môđun & 14\\
\hyperref[section:II.2]{\textsection2} & Các phổ nguyên tố thuần nhất & 19\\
\hyperref[subsection:II.2.1]{2.1} & Đại cương về các vành và môđun phân bậc & 19\\
\hyperref[subsection:II.2.2]{2.2} & Các vành phân thức của một vành phân bậc & 23\\
\hyperref[subsection:II.2.3]{2.3} & Phổ nguyên tố thuần nhất của một vành phân bậc & 25\\
\hyperref[subsection:II.2.4]{2.4} & Cấu trúc lược đồ trên $\Proj(S)$ & 28\\
\hyperref[subsection:II.2.5]{2.5} & Bó liên kết với một môđun phân bậc & 30\\
\hyperref[subsection:II.2.6]{2.6} & $S$-môđun phân bậc liên kết với một bó trên $\Proj(S)$ & 36\\
\hyperref[subsection:II.2.7]{2.7} & Các điều kiện hữu hạn & 38\\
\hyperref[subsection:II.2.8]{2.8} & Tính hàm tử & 41\\
\hyperref[subsection:II.2.9]{2.9} & Các tiền lược đồ con đóng của lược đồ $\Proj(S)$ & 48\\
\hyperref[section:II.3]{\textsection3} & Phổ thuần nhất của một bó đại số phân bậc & 49\\
\hyperref[subsection:II.3.1]{3.1} & Phổ thuần nhất của một $\sh{O}_Y$-đại số phân bậc quasi-coherent & 49\\
\hyperref[subsection:II.3.2]{3.2} & Bó trên $\Proj(\sh{S})$ liên kết với một $\sh{S}$-môđun phân bậc & 54\\
\hyperref[subsection:II.3.3]{3.3} & $\sh{S}$-môđun phân bậc liên kết với một bó trên $\Proj(\sh{S})$ & 56\\
\hyperref[subsection:II.3.4]{3.4} & Các điều kiện hữu hạn & 59\\
\hyperref[subsection:II.3.5]{3.5} & Tính hàm tử & 61\\
\hyperref[subsection:II.3.6]{3.6} & Các tiền lược đồ con đóng của tiền lược đồ $\Proj(\sh{S})$ & 64\\
\hyperref[subsection:II.3.7]{3.7} & Các cấu xạ từ một tiền lược đồ vào một phổ thuần nhất & 65\\
\hyperref[subsection:II.3.8]{3.8} & Các tiêu chuẩn về phép nhúng vào một phổ thuần nhất & 69\\
\end{longtable}
\oldpage[II]{214}
\markboth{A. GROTHENDIECK --- Chương II}{A. GROTHENDIECK --- Chương II}
\begin{longtable}{@{}p{0.09\textwidth}p{0.69\textwidth}r@{}}
\hyperref[section:II.4]{\textsection4} & Các phân thớ xạ ảnh; các bó ample & 71\\
\hyperref[subsection:II.4.1]{4.1} & Định nghĩa các phân thớ xạ ảnh & 71\\
\hyperref[subsection:II.4.2]{4.2} & Các cấu xạ từ một tiền lược đồ vào một phân thớ xạ ảnh & 72\\
\hyperref[subsection:II.4.3]{4.3} & Cấu xạ Segre & 76\\
\hyperref[subsection:II.4.4]{4.4} & Các phép nhúng vào phân thớ xạ ảnh; các bó rất ample & 78\\
\hyperref[subsection:II.4.5]{4.5} & Các bó ample & 83\\
\hyperref[subsection:II.4.6]{4.6} & Các bó ample tương đối & 89\\
\hyperref[section:II.5]{\textsection5} & Các cấu xạ quasi-affine; các cấu xạ quasi-xạ ảnh; các cấu xạ riêng; các cấu xạ xạ ảnh & 94\\
\hyperref[subsection:II.5.1]{5.1} & Các cấu xạ quasi-affine & 94\\
\hyperref[subsection:II.5.2]{5.2} & Tiêu chuẩn Serre & 97\\
\hyperref[subsection:II.5.3]{5.3} & Các cấu xạ quasi-xạ ảnh & 99\\
\hyperref[subsection:II.5.4]{5.4} & Các cấu xạ riêng và các cấu xạ đóng phổ dụng & 100\\
\hyperref[subsection:II.5.5]{5.5} & Các cấu xạ xạ ảnh & 103\\
\hyperref[subsection:II.5.6]{5.6} & Bổ đề Chow & 106\\
\hyperref[section:II.6]{\textsection6} & Các cấu xạ nguyên và các cấu xạ hữu hạn & 110\\
\hyperref[subsection:II.6.1]{6.1} & Các tiền lược đồ nguyên trên một tiền lược đồ khác & 110\\
\hyperref[subsection:II.6.2]{6.2} & Các cấu xạ quasi-hữu hạn & 114\\
\hyperref[subsection:II.6.3]{6.3} & Bao đóng nguyên của một tiền lược đồ & 116\\
\hyperref[subsection:II.6.4]{6.4} & Định thức của một tự đồng cấu của một $\sh{O}_X$-môđun & 120\\
\hyperref[subsection:II.6.5]{6.5} & Chuẩn của một bó khả nghịch & 125\\
\hyperref[subsection:II.6.6]{6.6} & Ứng dụng: các tiêu chuẩn về tính ample & 130\\
\hyperref[subsection:II.6.7]{6.7} & Định lý Chevalley & 135\\
\hyperref[section:II.7]{\textsection7} & Các tiêu chuẩn định giá & 138\\
\hyperref[subsection:II.7.1]{7.1} & Ôn lại về các vành định giá & 138\\
\hyperref[subsection:II.7.2]{7.2} & Tiêu chuẩn định giá cho tính tách & 141\\
\hyperref[subsection:II.7.3]{7.3} & Tiêu chuẩn định giá cho tính riêng & 143\\
\hyperref[subsection:II.7.4]{7.4} & Đường cong đại số và trường hàm chiều 1 & 148\\
\hyperref[section:II.8]{\textsection8} & Các lược đồ thu được bằng phép thổi; các nón chiếu; bao đóng xạ ảnh & 152\\
\hyperref[subsection:II.8.1]{8.1} & Các tiền lược đồ thu được bằng phép thổi & 152\\
\hyperref[subsection:II.8.2]{8.2} & Các kết quả sơ bộ về phép địa phương hóa trong các vành phân bậc & 157\\
\hyperref[subsection:II.8.3]{8.3} & Các nón chiếu & 162\\
\hyperref[subsection:II.8.4]{8.4} & Bao đóng xạ ảnh của một phân thớ vectơ & 168\\
\hyperref[subsection:II.8.5]{8.5} & Tính hàm tử & 169\\
\hyperref[subsection:II.8.6]{8.6} & Một đẳng cấu chính tắc cho các nón bỏ đỉnh & 171\\
\end{longtable}
\oldpage[II]{215}
\markboth{MỤC LỤC}{MỤC LỤC}
\begin{longtable}{@{}p{0.09\textwidth}p{0.69\textwidth}r@{}}
\hyperref[subsection:II.8.7]{8.7} & Phép thổi các nón chiếu & 173\\
\hyperref[subsection:II.8.8]{8.8} & Các bó ample và các phép co & 177\\
\hyperref[subsection:II.8.9]{8.9} & Tiêu chuẩn về tính ample của Grauert: phát biểu & 182\\
\hyperref[subsection:II.8.10]{8.10} & Tiêu chuẩn về tính ample của Grauert: chứng minh & 184\\
\hyperref[subsection:II.8.11]{8.11} & Tính duy nhất của các phép co & 189\\
\hyperref[subsection:II.8.12]{8.12} & Các bó quasi-coherent trên các nón chiếu & 191\\
\hyperref[subsection:II.8.13]{8.13} & Bao đóng xạ ảnh của các bó con và các lược đồ con đóng & 195\\
\hyperref[subsection:II.8.14]{8.14} & Bổ sung về các bó liên kết với các $\sh{S}$-môđun phân bậc & 197\\
\multicolumn{2}{@{}l}{Thư mục (tiếp theo)} & 205\\
\multicolumn{2}{@{}l}{Chỉ mục ký hiệu} & 207\\
\multicolumn{2}{@{}l}{Chỉ mục thuật ngữ} & 209\\
\multicolumn{2}{@{}l}{Đính chính và bổ sung (Danh sách 1)} & 217\\
\end{longtable}

\begin{flushright}
\emph{Nhận ngày 22 tháng 7 năm 1960.}
\end{flushright}
\endgroup
\oldpage[II]{216}
\oldpage[II]{217}
\phantomsection
\addcontentsline{toc}{section}{Đính chính và bổ sung (Danh sách 1)}
\section*{Đính chính và bổ sung}
\begin{center}
\large (Danh sách 1)
\end{center}

\subsection*{Chương 0}

\paragraph{(0, 2.1.3).}
Ở dòng thứ 15 tính từ cuối tr.~21, thay «giao-phần» bằng «giao».

\paragraph{(0, 2.1.6).}
Ở dòng 13 của tr.~22, thay
\[
  Z_i\cap C\left(\bigcup_{j\ne i}Z_j\right)
  \quad\hbox{bằng}\quad
  \bigcup_{j\ne i}Z_j.
\]

\paragraph{(0, 3.1.6).}
Ở dòng thứ 17 tính từ cuối tr.~25, thay
$\Gamma(U,\sh{F})$ bằng $\Gamma(V,\sh{F})$.

\paragraph{(0, 4.1.1).}
Ở dòng 12 của tr.~36, thay $\sh{G}\to\sh{F}$ bằng
$\sh{B}\to\sh{A}$.

\paragraph{(0, 4.2.1).}
Ở dòng thứ 11 tính từ cuối tr.~39, thay «các bó» bằng «bó».

\paragraph{(0, 5.1.2).}
Sau dòng 14 của tr.~45, thêm:

Nếu $\sh{G}$ là một $\sh{O}_Y$-môđun sinh bởi các tiết diện của nó
trên $Y$, thì $f^*(\sh{G})$ sinh bởi các tiết diện của nó trên $X$, vì
$f^*$ là một hàm tử khớp phải.

\paragraph{(0, 5.2.4).}
Ở dòng thứ 18 tính từ cuối tr.~46, thay $f_U$ bằng $f_U^*$.

\paragraph{(0, 5.3.4).}
Trước dòng thứ 16 tính từ cuối tr.~47, thêm:

Nếu
\[
  \sh{A}\to\sh{B}\to\sh{C}\to\sh{D}\to\sh{E}
\]
là một dãy khớp các $\sh{O}_X$-môđun, và nếu
$\sh{A},\sh{B},\sh{D},\sh{E}$ là coherent, thì $\sh{C}$ là coherent.

\paragraph{(0, 5.4.1).}
Sau dòng 2 của tr.~49, thêm:

Giả sử $\sh{O}_X$ là coherent, và cho $\sh{F}$ là một
$\sh{O}_X$-môđun coherent. Nếu tại một điểm $x\in X$, $\sh{F}_x$ là
một $\sh{O}_x$-môđun tự do hạng $n$, thì tồn tại một lân cận $U$ của
$x$ sao cho $\sh{F}|U$ tự do địa phương hạng $n$; thật vậy,
$\sh{F}_x$ khi đó đẳng cấu với $\sh{O}_x^n$, và khẳng định suy ra từ
(0, 5.2.7).

\paragraph{(0, 6.6.2).}
Thay các dòng 10 và 11 của tr.~59 bằng:

Thật vậy, điều này suy ra từ (0, 6.4.1, d).

\paragraph{(0, 6.7.1).}
Trước dòng thứ 10 tính từ cuối tr.~59, thêm:

Nếu $f$ là một cấu xạ phẳng, ta còn nói rằng $X$ \emph{phẳng trên
$Y$}, hay \emph{$Y$-phẳng}.

\paragraph{(0, 7.2.9).}
Ở dòng thứ 11 tính từ cuối tr.~65, thay $M_n$ bằng $M_{n-1}$.

\subsection*{Chương I}

\paragraph{(I, 1.2.1).}
Ở dòng 18 của tr.~83, thay $K$ bằng $k$ (hai lần).

\paragraph{(I, 1.2.7).}
Ở dòng thứ 16 tính từ cuối tr.~84, thay $A$ bằng $A'$ (hai lần) và
$\rad$ bằng $\rad'$; ở dòng thứ 17 tính từ cuối, thay $X$ bằng $X'$.

\paragraph{(I, 1.7.3) và (I, 2.2.4).}
Theo một nhận xét của J.~Tate, các mệnh đề nêu dưới những số này có
thể được tổng quát hóa như sau.

\subsection*{1.8. Cấu xạ từ không gian vành địa phương vào lược đồ affine}
\label{ega2:addendum:subsection:I.1.8}

\begin{proposition}[1.8.1]
\label{ega2:addendum:I.1.8.1}
Cho $(S,\sh{O}_S)$ là một lược đồ affine và $(X,\sh{O}_X)$ là một
không gian vành địa phương. Khi đó có một song ánh chính tắc từ tập
các đồng cấu vành
\oldpage[II]{218}
\[
  \Gamma(S,\sh{O}_S)\to\Gamma(X,\sh{O}_X)
\]
vào tập các cấu xạ không gian vành
$(\psi,\theta):(X,\sh{O}_X)\to(S,\sh{O}_S)$ sao cho, với mọi
$x\in X$, $\theta_x^\sharp:\sh{O}_{\psi(x)}\to\sh{O}_x$ là một
đồng cấu địa phương.
\end{proposition}

\begin{proof}
Trước hết lưu ý rằng, nếu $(X,\sh{O}_X)$ và $(S,\sh{O}_S)$ là hai
không gian vành bất kỳ, một cấu xạ $(\psi,\theta)$ từ $(X,\sh{O}_X)$
vào $(S,\sh{O}_S)$ xác định chính tắc một đồng cấu vành
\[
  \Gamma(\theta):\Gamma(S,\sh{O}_S)\to\Gamma(X,\sh{O}_X),
\]
và do đó xác định một ánh xạ thứ nhất
\begin{equation}
\label{ega2:addendum:I.1.8.1.1}
  \rho:\Hom((X,\sh{O}_X),(S,\sh{O}_S))
  \to\Hom(\Gamma(S,\sh{O}_S),\Gamma(X,\sh{O}_X)).
\tag{1.8.1.1}
\end{equation}
Ngược lại, dưới các giả thiết của mệnh đề, đặt
$A=\Gamma(S,\sh{O}_S)$ và xét một đồng cấu vành
$\vphi:A\to\Gamma(X,\sh{O}_X)$. Với mọi $x\in X$, tập các $f\in A$
sao cho $\vphi(f)(x)=0$ là một iđêan nguyên tố của $A$, vì
$\sh{O}_x/\mathfrak{m}_x=\kres(x)$ là một trường. Do đó nó là một
phần tử của $S=\Spec(A)$, mà ta vẫn ký hiệu là ${}^a\vphi(x)$. Hơn
nữa, với mọi $f\in A$, theo định nghĩa (0, 5.5.2),
\[
  {}^a\vphi^{-1}(D(f))=X_f,
\]
điều này chứng minh rằng ${}^a\vphi$ là một ánh xạ liên tục $X\to S$.

Tiếp theo, ta xác định một đồng cấu $\sh{O}_S$-môđun
\[
  \widetilde{\vphi}:\sh{O}_S\to{}^a\vphi_*(\sh{O}_X).
\]
Với mọi $f\in A$, ta có
$\Gamma(D(f),\sh{O}_S)=A_f$ (I, 1.3.6); với $s/f\in A_f$, cho tương
ứng phần tử
\[
  (\vphi(s)|X_f)(\vphi(f)|X_f)^{-1}
  \in\Gamma(X_f,\sh{O}_X)
  =\Gamma(D(f),{}^a\vphi_*(\sh{O}_X)).
\]
Chuyển từ $D(f)$ sang $D(fg)$ cho thấy ngay rằng điều này xác định một
đồng cấu bó, và do đó một cấu xạ không gian vành
$({}^a\vphi,\widetilde{\vphi})$. Với cùng ký hiệu, và viết
$y={}^a\vphi(x)$, ta có (0, 3.7.1)
\[
  \widetilde{\vphi}_x^\sharp(s_y/f_y)
  =(\vphi(s)_x)(\vphi(f)_x)^{-1}.
\]
Vì $s_y\in\mathfrak{m}_y$ theo định nghĩa tương đương với
$\vphi(s)_x\in\mathfrak{m}_x$, đồng cấu
$\widetilde{\vphi}_x^\sharp:\sh{O}_y\to\sh{O}_x$ là địa phương. Như
vậy ta đã xác định một ánh xạ thứ hai
\[
  \sigma:\Hom(\Gamma(S,\sh{O}_S),\Gamma(X,\sh{O}_X))
  \to\mathfrak{L},
\]
trong đó $\mathfrak{L}$ là tập các cấu xạ
$(\psi,\theta):(X,\sh{O}_X)\to(S,\sh{O}_S)$ sao cho
$\theta_x^\sharp$ là địa phương với mọi $x\in X$.

Còn phải chứng minh rằng $\sigma$ và hạn chế của $\rho$ lên
$\mathfrak{L}$ là nghịch đảo của nhau. Định nghĩa của
$\widetilde{\vphi}$ cho ngay $\Gamma(\widetilde{\vphi})=\vphi$, nên
$\rho\circ\sigma$ là đồng nhất. Để thấy $\sigma\circ\rho$ là đồng
nhất, bắt đầu với $(\psi,\theta)\in\mathfrak{L}$ và đặt
$\vphi=\Gamma(\theta)$. Giả thiết về $\theta_x^\sharp$ cho, sau khi
chuyển sang các trường thặng dư, một đơn cấu
\[
  \theta^x:\kres(\psi(x))\to\kres(x)
\]
sao cho, với mọi $f\in A=\Gamma(S,\sh{O}_S)$,
\[
  \theta^x(f(\psi(x)))=\vphi(f)(x).
\]
Vì thế $f(\psi(x))=0$ tương đương với $\vphi(f)(x)=0$, điều này chứng
minh ${}^a\vphi=\psi$. Mặt khác, các định nghĩa làm cho biểu đồ
\[
  \xymatrix{
    A\ar[r]^\vphi\ar[d] &
    \Gamma(X,\sh{O}_X)\ar[d]\\
    A_{\psi(x)}\ar[r]^{\theta_x^\sharp} &
    \sh{O}_x
  }
\]
giao hoán; biểu đồ tương tự trong đó $\theta_x^\sharp$ được thay bởi
$\widetilde{\vphi}_x^\sharp$ cũng giao hoán. Từ (0, 1.2.4) suy ra
$\widetilde{\vphi}_x^\sharp=\theta_x^\sharp$, và do đó
$\widetilde{\vphi}=\theta$.
\end{proof}

\begin{env}[1.8.2]
\label{ega2:addendum:I.1.8.2}
Khi $(X,\sh{O}_X)$ và $(Y,\sh{O}_Y)$ là các không gian vành địa
phương, ta sẽ xét các cấu xạ
$(\psi,\theta):(X,\sh{O}_X)\to(Y,\sh{O}_Y)$ sao cho, với mọi
$x\in X$, $\theta_x^\sharp:\sh{O}_{\psi(x)}\to\sh{O}_x$ là một đồng
cấu địa phương.

Từ nay, khi nói đến một
\oldpage[II]{219}
\emph{cấu xạ các không gian vành địa phương}, ta luôn hiểu đó là một
cấu xạ thuộc loại này. Với định nghĩa cấu xạ này, các không gian vành
địa phương tạo thành một phạm trù. Đối với hai vật $X,Y$ của phạm trù
này, $\Hom(X,Y)$ do đó sẽ ký hiệu tập các cấu xạ không gian vành địa
phương từ $X$ vào $Y$ (tập được ký hiệu $\mathfrak{L}$ trong
\sref{I.1.8.1}). Khi cần xét tập các cấu xạ không gian vành từ $X$ vào
$Y$, ta viết $\Hom_{\mathrm{rs}}(X,Y)$ để tránh nhầm lẫn. Ánh xạ
\sref{I.1.8.1.1} do đó được viết
\begin{equation}
\label{ega2:addendum:I.1.8.2.1}
  \rho:\Hom_{\mathrm{rs}}(X,Y)
  \to\Hom(\Gamma(Y,\sh{O}_Y),\Gamma(X,\sh{O}_X)),
\tag{1.8.2.1}
\end{equation}
và hạn chế của nó
\begin{equation}
\label{ega2:addendum:I.1.8.2.2}
  \rho':\Hom(X,Y)
  \to\Hom(\Gamma(Y,\sh{O}_Y),\Gamma(X,\sh{O}_X))
\tag{1.8.2.2}
\end{equation}
có tính hàm tử theo $X$ và $Y$ trong phạm trù các không gian vành địa
phương.
\end{env}

\begin{corollary}[1.8.3]
\label{ega2:addendum:I.1.8.3}
Cho $(Y,\sh{O}_Y)$ là một không gian vành địa phương. Điều kiện cần và
đủ để $Y$ là một lược đồ affine là, với mọi không gian vành địa phương
$(X,\sh{O}_X)$, ánh xạ \sref{I.1.8.2.2} là song ánh.
\end{corollary}

\begin{proof}
Mệnh đề~\sref{I.1.8.1} cho thấy điều kiện là cần. Ngược lại, giả sử
điều kiện ấy đúng và đặt $A=\Gamma(Y,\sh{O}_Y)$. Từ giả thiết và
\sref{I.1.8.1}, các hàm tử
\[
  X\longmapsto\Hom(X,Y)
  \quad\hbox{và}\quad
  X\longmapsto\Hom(X,\Spec(A))
\]
từ các không gian vành địa phương vào các tập hợp là đẳng cấu. Điều
này kéo theo sự tồn tại của một đẳng cấu chính tắc
$Y\to\Spec(A)$ (xem 0, 8).
\end{proof}

\begin{env}[1.8.4]
\label{ega2:addendum:I.1.8.4}
Cho $S=\Spec(A)$ là một lược đồ affine. Ký hiệu bởi $(S',A')$ không
gian vành có không gian nền chỉ gồm một điểm và có bó cấu trúc $A'$ là
bó (tất yếu đơn) trên $S'$ xác định bởi vành $A$. Cho
$\pi:S\to S'$ là ánh xạ duy nhất. Với mọi tập mở $U$ của $S$, ánh xạ
chính tắc
\[
  \Gamma(S',A')=\Gamma(S,\sh{O}_S)\to\Gamma(U,\sh{O}_S)
\]
xác định một $\pi$-cấu xạ $\iota:A'\to\sh{O}_S$ của các bó vành. Như
vậy ta có một cấu xạ chính tắc của các không gian vành
\[
  i=(\pi,\iota):(S,\sh{O}_S)\to(S',A').
\]
Với mọi $A$-môđun $M$, ký hiệu bởi $M'$ bó đơn trên $S'$ xác định bởi
$M$, hiển nhiên là một $A'$-môđun. Rõ ràng
$i_*(\widetilde M)=M'$ (I, 1.3.7).
\end{env}

\begin{lemma}[1.8.5]
\label{ega2:addendum:I.1.8.5}
Với ký hiệu của \sref{I.1.8.4}, đối với mọi $A$-môđun $M$, đồng cấu
$\sh{O}_S$-môđun chính tắc, có tính hàm tử (0, 4.4.3.3)
\begin{equation}
\label{ega2:addendum:I.1.8.5.1}
  i^*(i_*(\widetilde M))\to\widetilde M
\tag{1.8.5.1}
\end{equation}
là một đẳng cấu.
\end{lemma}

\begin{proof}
Cả hai vế của \sref{I.1.8.5.1} đều khớp phải theo $M$ (hàm tử
$M\mapsto i_*(\widetilde M)$ hiển nhiên là khớp) và giao hoán với các
tổng trực tiếp. Xem $M$ như đối hạt nhân của một đồng cấu
$A^{(I)}\to A^{(J)}$, ta quy về trường hợp $M=A$, ở đó khẳng định là
hiển nhiên.
\end{proof}

\begin{corollary}[1.8.6]
\label{ega2:addendum:I.1.8.6}
Cho $(X,\sh{O}_X)$ là một không gian vành và $u:X\to S$ là một cấu xạ
\oldpage[II]{220}
không gian vành. Với mọi $A$-môđun $M$, theo ký hiệu của
\sref{I.1.8.4}, có một đẳng cấu chính tắc, có tính hàm tử, của các
$\sh{O}_X$-môđun
\begin{equation}
\label{ega2:addendum:I.1.8.6.1}
  u^*(\widetilde M)\isoto u^*(i^*(M')).
\tag{1.8.6.1}
\end{equation}
\end{corollary}

\begin{corollary}[1.8.7]
\label{ega2:addendum:I.1.8.7}
Dưới các giả thiết của \sref{I.1.8.6}, với mọi $A$-môđun $M$ và mọi
$\sh{O}_X$-môđun $\sh{F}$, có một đẳng cấu chính tắc, có tính hàm tử
theo $M$ và $\sh{F}$,
\begin{equation}
\label{ega2:addendum:I.1.8.7.1}
  \Hom_{\sh{O}_S}(\widetilde M,u_*(\sh{F}))
  \isoto\Hom_A(M,\Gamma(X,\sh{F})).
\tag{1.8.7.1}
\end{equation}
\end{corollary}

\begin{proof}
Theo (0, 4.4.3) và Bổ đề~\sref{I.1.8.5}, có một đẳng cấu chính tắc
của các song hàm tử
\[
  \Hom_{\sh{O}_S}(\widetilde M,u_*(\sh{F}))
  \isoto\Hom_{A'}(M',i_*(u_*(\sh{F}))).
\]
Vế phải hiển nhiên chính là $\Hom_A(M,\Gamma(X,\sh{F}))$. Dưới đồng
cấu chính tắc \sref{I.1.8.7.1}, mỗi $\sh{O}_S$-đồng cấu
$h:\widetilde M\to u_*(\sh{F})$ (tương đương, mỗi $u$-cấu xạ
$\widetilde M\to\sh{F}$) tương ứng với
\[
  \Gamma(h):M\to\Gamma(S,u_*(\sh{F}))=\Gamma(X,\sh{F}).
\]
\end{proof}

\begin{env}[1.8.8]
\label{ega2:addendum:I.1.8.8}
Với ký hiệu của \sref{I.1.8.4}, từ (0, 4.1.1) thấy rõ rằng mỗi cấu xạ
không gian vành $(\psi,\theta):X\to S'$ tương đương với việc cho một
đồng cấu vành $A\to\Gamma(X,\sh{O}_X)$. Do đó
Mệnh đề~\sref{I.1.8.1} còn có thể được hiểu là cho một song ánh chính
tắc
\[
  \Hom(X,S)\isoto\Hom(X,S'),
\]
trong đó vế phải được hiểu là các cấu xạ không gian vành, vì nói chung
$A$ không phải là một vành địa phương.

Tổng quát hơn, cho $X,Y$ là hai không gian vành địa phương và cho
$(Y',A')$ là không gian vành có không gian nền chỉ gồm một điểm và có
bó vành là bó đơn xác định bởi $\Gamma(Y,\sh{O}_Y)$. Khi đó
\sref{I.1.8.2.1} có thể được hiểu là một ánh xạ
\begin{equation}
\label{ega2:addendum:I.1.8.8.1}
  \rho:\Hom_{\mathrm{rs}}(X,Y)\to\Hom(X,Y').
\tag{1.8.8.1}
\end{equation}
Hệ quả~\sref{I.1.8.3} nói rằng, trong các không gian vành địa phương,
các lược đồ affine được đặc trưng là những không gian mà hạn chế của
$\rho$ lên $\Hom(X,Y)$,
\begin{equation}
\label{ega2:addendum:I.1.8.8.2}
  \rho':\Hom(X,Y)\to\Hom(X,Y'),
\tag{1.8.8.2}
\end{equation}
là song ánh với mọi không gian vành địa phương $X$. Trong một chương
sau, ta sẽ tổng quát hóa định nghĩa này, qua đó liên kết với một không
gian vành tùy ý $Z$ (chứ không chỉ với một không gian vành có không
gian nền chỉ gồm một điểm) một không gian vành địa phương mà ta vẫn
gọi là $\Spec(Z)$. Đây sẽ là điểm xuất phát của một lý thuyết
«tương đối» về các tiền lược đồ trên các không gian vành tùy ý, mở
rộng các kết quả của Chương~I.
\end{env}

\begin{env}[1.8.9]
\label{ega2:addendum:I.1.8.9}
Các cặp $(X,\sh{F})$, trong đó $X$ là một không gian vành địa phương
và $\sh{F}$ là một $\sh{O}_X$-môđun, tạo thành một phạm trù: một cấu
xạ $(X,\sh{F})\to(Y,\sh{G})$ là một cặp $(u,h)$ gồm một cấu xạ
\oldpage[II]{221}
giữa các không gian vành địa phương $u:X\to Y$ và một $u$-cấu xạ
$h:\sh{G}\to\sh{F}$ của các môđun; các
cấu xạ này (với $(X,\sh{F})$ và $(Y,\sh{G})$ cố định) tạo thành một tập
hợp, ký hiệu là $\Hom((X,\sh{F}),(Y,\sh{G}))$; ánh xạ
$(u,h)\mapsto(\rho'(u),\Gamma(h))$ là ánh xạ chính tắc
\begin{equation}
\label{ega2:addendum:I.1.8.9.1}
\begin{split}
  \Hom((X,\sh{F}),(Y,\sh{G}))\to
  \Hom\bigl(
    &(\Gamma(Y,\sh{O}_Y),\Gamma(Y,\sh{G})),\\
    &(\Gamma(X,\sh{O}_X),\Gamma(X,\sh{F}))
  \bigr),
\end{split}
\tag{1.8.9.1}
\end{equation}
có tính hàm tử theo $(X,\sh{F})$ và $(Y,\sh{G})$; vế phải là tập hợp
các di-đồng cấu tương ứng với các vành và môđun đang xét (0, 1.0.2).
\end{env}

\begin{corollary}[1.8.10]
\label{ega2:addendum:I.1.8.10}
Cho $Y$ là một không gian vành địa phương và $\sh{G}$ là một
$\sh{O}_Y$-môđun. Điều kiện cần và đủ để $Y$ là một lược đồ affine và
$\sh{G}$ là một $\sh{O}_Y$-môđun quasi-coherent là: với mọi cặp
$(X,\sh{F})$ gồm một không gian vành địa phương $X$ và một
$\sh{O}_X$-môđun $\sh{F}$, ánh xạ chính tắc \sref{I.1.8.9.1} là song
ánh.
\end{corollary}

Ta để bạn đọc khai triển lập luận, được mô phỏng theo lập luận của
\sref{I.1.8.3} và sử dụng \sref{I.1.8.1} cùng \sref{I.1.8.7}.

\begin{remark}[1.8.11]
\label{ega2:addendum:I.1.8.11}
Các mệnh đề (I, 1.7.3), (I, 1.7.4), (I, 2.2.4), cũng như định nghĩa
(I, 1.6.1), là các trường hợp riêng của \sref{I.1.8.1}; tương tự,
(I, 2.2.5) suy ra từ \sref{I.1.8.7}. Hệ quả~\sref{I.1.8.7} cũng cho
(I, 1.6.3) (và do đó (I, 1.6.4)) như một trường hợp riêng. Thật vậy,
nếu $X$ affine và $\Gamma(X,\sh{F})=N$, thì các hàm tử
\[
  M\longmapsto
  \Hom_{\sh{O}_S}(\widetilde M,u_*(\widetilde N))
  \quad\hbox{và}\quad
  M\longmapsto
  \Hom_{\sh{O}_S}(\widetilde M,(N_{[\vphi]})\supertilde)
\]
(trong đó $\vphi:A\to\Gamma(X,\sh{O}_X)$ tương ứng với $u$) đẳng cấu
theo \sref{I.1.8.7} và (I, 1.3.8). Cuối cùng, (I, 1.6.5) (và do đó
(I, 1.6.6)) suy ra từ \sref{I.1.8.6} và từ sự kiện rằng, với mọi
$f\in A$, các $A_f$-môđun
\[
  N'\otimes_{A'}A_f
  \quad\hbox{và}\quad
  (N'\otimes_{A'}A)_f
\]
(theo ký hiệu của (I, 1.6.5)) đẳng cấu chính tắc.
\end{remark}

\paragraph{Chèn sau (I, 3.2.8):}

\begin{env}[3.2.9]
\label{ega2:addendum:I.3.2.9}
Từ \sref{I.1.8.1} suy ra có thể làm chính xác hơn (I, 3.2.2) như sau.
Tiền lược đồ
\[
  Z=\Spec(B\otimes_A C)
\]
không chỉ là tích của $X=\Spec(B)$ và $Y=\Spec(C)$ trong phạm trù các
$S$-tiền lược đồ, mà còn trong phạm trù các không gian vành địa
phương trên $S$ (với định nghĩa $S$-cấu xạ được mô phỏng theo
(I, 2.5.2)). Chứng minh của (I, 3.2.6) cũng cho thấy rằng, với hai
$S$-tiền lược đồ tùy ý $X,Y$, tiền lược đồ $X\times_S Y$ không chỉ là
tích của chúng trong phạm trù các $S$-tiền lược đồ, mà còn trong phạm
trù các không gian vành địa phương trên tiền lược đồ $S$.
\end{env}

\paragraph{(I, 4.4.5).}
Ở dòng thứ 14 tính từ cuối tr.~126, thay $B$ bằng $A$; ở dòng thứ 13
tính từ cuối, thay «$A$-đại số» bằng «$B$-đại số».

\paragraph{(I, 5.3.5).}
Ở dòng thứ 2 tính từ cuối tr.~132, thay $Y\to S$ bằng $g:Y\to S$.

\paragraph{(I, 6.3.2.1).}
Ở dòng thứ 6 tính từ cuối tr.~144, đọc là:
\[
  D(g_i)\subset W.
\]

\paragraph{(I, 7.3.8).}
Thay văn bản hiện tại, vốn sai, bằng văn bản sau:

Cho $X,Y$ là hai tiền lược đồ nguyên; khi đó $\sh{R}(X)$ và
$\sh{R}(Y)$ lần lượt là một $\sh{O}_X$-môđun quasi-coherent và một
$\sh{O}_Y$-môđun quasi-coherent (I, 7.3.3). Nếu $f:X\to Y$ là một cấu
xạ trội, thì tồn tại một đồng cấu chính tắc của các $\sh{O}_X$-môđun
\oldpage[II]{222}
\begin{equation}
\label{ega2:addendum:I.7.3.8.1}
  \tau:f^*(\sh{R}(Y))\to\sh{R}(X).
\tag{7.3.8.1}
\end{equation}

Trước hết giả sử $X=\Spec(A)$ và $Y=\Spec(B)$ là affine, với $A,B$
là các miền nguyên. Khi đó $f$ tương ứng với một đồng cấu đơn ánh
$B\to A$, và đồng cấu này mở rộng thành một đơn cấu $L\to K$ từ trường
phân thức $L$ của $B$ vào trường phân thức $K$ của $A$. Khi ấy đồng
cấu \sref{I.7.3.8.1} tương ứng với đồng cấu chính tắc
\[
  L\otimes_B A\to K
\]
của (I, 1.6.5).

Trong trường hợp tổng quát, với mỗi cặp tập mở affine không rỗng
$U\subset X$, $V\subset Y$ sao cho $f(U)\subset V$, phép dựng trên xác
định một đồng cấu $\tau_{U,V}$. Nếu $U'\subset U$, $V'\subset V$ và
$f(U')\subset V'$, thì $\tau_{U,V}$ mở rộng $\tau_{U',V'}$; khẳng định
suy ra bằng phép dán. Nếu $x,y$ lần lượt là các điểm tổng quát của
$X,Y$, thì $f(x)=y$ và
\[
  \bigl(f^*(\sh{R}(Y))\bigr)_x
  =\sh{O}_y\otimes_{\sh{O}_y}\sh{O}_x
  =\sh{O}_x
\]
theo (0, 4.3.1), vì thế $\tau_x$ là một đẳng cấu.

\paragraph{(I, 9.5.2).}
Trong ba dòng cuối tr.~176 và năm dòng đầu tr.~177, thay mọi $X'$ bằng
$Y'$ và mọi $X''$ bằng $Y''$.

\paragraph{(I, 10.5.4).}
Ở dòng thứ 11 tính từ cuối tr.~187, thay (10.3.4) bằng (10.3.5).

\paragraph{(I, 10.6.3).}
Ở dòng thứ 3 tính từ cuối tr.~189, thay $X$ bằng $\mathfrak X$.

\paragraph{(I, 10.14.2).}
Ở dòng 5 tr.~210, thay «$\sh{O}_{\mathfrak X}$-môđun coherent» bằng
«iđêan coherent của $\sh{O}_{\mathfrak X}$»; ở dòng 7, thay
«các $\sh{O}_{\mathfrak X}$-môđun coherent» bằng «các iđêan coherent
của $\sh{O}_{\mathfrak X}$».
\end{document}
